\documentclass{article}
\usepackage[paperwidth=16cm,paperheight=14cm,margin=0cm]{geometry}
\usepackage[utf8]{inputenc}\usepackage[T1]{fontenc}
\usepackage{mathpazo}\usepackage{amsmath,amssymb}
\usepackage[most]{tcolorbox}
\pagestyle{empty}\setlength{\parindent}{0pt}

\definecolor{cnum}{HTML}{1F6FA8}
\definecolor{ccal}{HTML}{6A4C93}
\definecolor{cfra}{HTML}{2E7D6B}
\definecolor{cgeo}{HTML}{C0503A}
\definecolor{cmes}{HTML}{8A5A2B}
\definecolor{cpro}{HTML}{C77F0A}
\begin{document}
\begin{tcolorbox}[enhanced,width=16cm,height=14cm,colback=white,colframe=cnum,boxrule=2pt,arc=0mm,
left=6mm,right=6mm,top=0mm,bottom=4mm,sharp corners,
overlay={\fill[cnum] (frame.north west) rectangle ([yshift=-1.25cm]frame.north east);
\node[anchor=west,text=white,font=\sffamily\bfseries\LARGE] at ([xshift=5mm,yshift=-0.62cm]frame.north west) {Question 1};
\node[anchor=east,text=white,font=\sffamily\itshape\large] at ([xshift=-5mm,yshift=-0.62cm]frame.north east) {Nombres entiers et décimaux};}]
\vspace*{1.75cm}
\begin{minipage}[t][6.9cm][c]{\linewidth}\LARGE\centering
\begin{minipage}{0.95\linewidth}\raggedright\linespread{1.25}\selectfont Dans le nombre $345{,}672$, quel est le chiffre des centièmes~?\end{minipage}
\end{minipage}

\vfill
\begin{tcolorbox}[enhanced,nobeforeafter,width=7.1cm,height=2.1cm,valign=center,colback=cnum!6,colframe=cnum,arc=2.5mm,boxrule=0.9pt,left=3mm]{\Large\bfseries\color{cnum}a)}\quad\begin{minipage}[c]{5.6cm}\raggedright\Large $7$\end{minipage}\end{tcolorbox}\hfill\begin{tcolorbox}[enhanced,nobeforeafter,width=7.1cm,height=2.1cm,valign=center,colback=cnum!6,colframe=cnum,arc=2.5mm,boxrule=0.9pt,left=3mm]{\Large\bfseries\color{cnum}b)}\quad\begin{minipage}[c]{5.6cm}\raggedright\Large $4$\end{minipage}\end{tcolorbox}\par\vspace{2.5mm}
\begin{tcolorbox}[enhanced,nobeforeafter,width=7.1cm,height=2.1cm,valign=center,colback=cnum!6,colframe=cnum,arc=2.5mm,boxrule=0.9pt,left=3mm]{\Large\bfseries\color{cnum}c)}\quad\begin{minipage}[c]{5.6cm}\raggedright\Large $5$\end{minipage}\end{tcolorbox}\hfill\begin{tcolorbox}[enhanced,nobeforeafter,width=7.1cm,height=2.1cm,valign=center,colback=cnum!6,colframe=cnum,arc=2.5mm,boxrule=0.9pt,left=3mm]{\Large\bfseries\color{cnum}d)}\quad\begin{minipage}[c]{5.6cm}\raggedright\Large $3$\end{minipage}\end{tcolorbox}
\end{tcolorbox}
\newpage
\begin{tcolorbox}[enhanced,width=16cm,height=14cm,colback=white,colframe=cnum,boxrule=2pt,arc=0mm,
left=6mm,right=6mm,top=0mm,bottom=4mm,sharp corners,
overlay={\fill[cnum] (frame.north west) rectangle ([yshift=-1.25cm]frame.north east);
\node[anchor=west,text=white,font=\sffamily\bfseries\LARGE] at ([xshift=5mm,yshift=-0.62cm]frame.north west) {Question 2};
\node[anchor=east,text=white,font=\sffamily\itshape\large] at ([xshift=-5mm,yshift=-0.62cm]frame.north east) {Nombres entiers et décimaux};}]
\vspace*{1.75cm}
\begin{minipage}[t][6.9cm][c]{\linewidth}\LARGE\centering
\begin{minipage}{0.95\linewidth}\raggedright\linespread{1.25}\selectfont Dans le nombre $1\,208{,}53$, quel est le chiffre des dizaines~?\end{minipage}
\end{minipage}

\vfill
\begin{tcolorbox}[enhanced,nobeforeafter,width=7.1cm,height=2.1cm,valign=center,colback=cnum!6,colframe=cnum,arc=2.5mm,boxrule=0.9pt,left=3mm]{\Large\bfseries\color{cnum}a)}\quad\begin{minipage}[c]{5.6cm}\raggedright\Large $2$\end{minipage}\end{tcolorbox}\hfill\begin{tcolorbox}[enhanced,nobeforeafter,width=7.1cm,height=2.1cm,valign=center,colback=cnum!6,colframe=cnum,arc=2.5mm,boxrule=0.9pt,left=3mm]{\Large\bfseries\color{cnum}b)}\quad\begin{minipage}[c]{5.6cm}\raggedright\Large $0$\end{minipage}\end{tcolorbox}\par\vspace{2.5mm}
\begin{tcolorbox}[enhanced,nobeforeafter,width=7.1cm,height=2.1cm,valign=center,colback=cnum!6,colframe=cnum,arc=2.5mm,boxrule=0.9pt,left=3mm]{\Large\bfseries\color{cnum}c)}\quad\begin{minipage}[c]{5.6cm}\raggedright\Large $1$\end{minipage}\end{tcolorbox}\hfill\begin{tcolorbox}[enhanced,nobeforeafter,width=7.1cm,height=2.1cm,valign=center,colback=cnum!6,colframe=cnum,arc=2.5mm,boxrule=0.9pt,left=3mm]{\Large\bfseries\color{cnum}d)}\quad\begin{minipage}[c]{5.6cm}\raggedright\Large $8$\end{minipage}\end{tcolorbox}
\end{tcolorbox}
\newpage
\begin{tcolorbox}[enhanced,width=16cm,height=14cm,colback=white,colframe=cnum,boxrule=2pt,arc=0mm,
left=6mm,right=6mm,top=0mm,bottom=4mm,sharp corners,
overlay={\fill[cnum] (frame.north west) rectangle ([yshift=-1.25cm]frame.north east);
\node[anchor=west,text=white,font=\sffamily\bfseries\LARGE] at ([xshift=5mm,yshift=-0.62cm]frame.north west) {Question 3};
\node[anchor=east,text=white,font=\sffamily\itshape\large] at ([xshift=-5mm,yshift=-0.62cm]frame.north east) {Nombres entiers et décimaux};}]
\vspace*{1.75cm}
\begin{minipage}[t][6.9cm][c]{\linewidth}\LARGE\centering
\begin{minipage}{0.95\linewidth}\raggedright\linespread{1.25}\selectfont Quelle est l'écriture décimale de $\dfrac{347}{100}$~?\end{minipage}
\end{minipage}

\vfill
\begin{tcolorbox}[enhanced,nobeforeafter,width=7.1cm,height=2.1cm,valign=center,colback=cnum!6,colframe=cnum,arc=2.5mm,boxrule=0.9pt,left=3mm]{\Large\bfseries\color{cnum}a)}\quad\begin{minipage}[c]{5.6cm}\raggedright\Large $34{,}7$\end{minipage}\end{tcolorbox}\hfill\begin{tcolorbox}[enhanced,nobeforeafter,width=7.1cm,height=2.1cm,valign=center,colback=cnum!6,colframe=cnum,arc=2.5mm,boxrule=0.9pt,left=3mm]{\Large\bfseries\color{cnum}b)}\quad\begin{minipage}[c]{5.6cm}\raggedright\Large $0{,}347$\end{minipage}\end{tcolorbox}\par\vspace{2.5mm}
\begin{tcolorbox}[enhanced,nobeforeafter,width=7.1cm,height=2.1cm,valign=center,colback=cnum!6,colframe=cnum,arc=2.5mm,boxrule=0.9pt,left=3mm]{\Large\bfseries\color{cnum}c)}\quad\begin{minipage}[c]{5.6cm}\raggedright\Large $347$\end{minipage}\end{tcolorbox}\hfill\begin{tcolorbox}[enhanced,nobeforeafter,width=7.1cm,height=2.1cm,valign=center,colback=cnum!6,colframe=cnum,arc=2.5mm,boxrule=0.9pt,left=3mm]{\Large\bfseries\color{cnum}d)}\quad\begin{minipage}[c]{5.6cm}\raggedright\Large $3{,}47$\end{minipage}\end{tcolorbox}
\end{tcolorbox}
\newpage
\begin{tcolorbox}[enhanced,width=16cm,height=14cm,colback=white,colframe=cnum,boxrule=2pt,arc=0mm,
left=6mm,right=6mm,top=0mm,bottom=4mm,sharp corners,
overlay={\fill[cnum] (frame.north west) rectangle ([yshift=-1.25cm]frame.north east);
\node[anchor=west,text=white,font=\sffamily\bfseries\LARGE] at ([xshift=5mm,yshift=-0.62cm]frame.north west) {Question 4};
\node[anchor=east,text=white,font=\sffamily\itshape\large] at ([xshift=-5mm,yshift=-0.62cm]frame.north east) {Nombres entiers et décimaux};}]
\vspace*{1.75cm}
\begin{minipage}[t][6.9cm][c]{\linewidth}\LARGE\centering
\begin{minipage}{0.95\linewidth}\raggedright\linespread{1.25}\selectfont Quelle est l'écriture décimale de $\dfrac{52}{10}$~?\end{minipage}
\end{minipage}

\vfill
\begin{tcolorbox}[enhanced,nobeforeafter,width=7.1cm,height=2.1cm,valign=center,colback=cnum!6,colframe=cnum,arc=2.5mm,boxrule=0.9pt,left=3mm]{\Large\bfseries\color{cnum}a)}\quad\begin{minipage}[c]{5.6cm}\raggedright\Large $52$\end{minipage}\end{tcolorbox}\hfill\begin{tcolorbox}[enhanced,nobeforeafter,width=7.1cm,height=2.1cm,valign=center,colback=cnum!6,colframe=cnum,arc=2.5mm,boxrule=0.9pt,left=3mm]{\Large\bfseries\color{cnum}b)}\quad\begin{minipage}[c]{5.6cm}\raggedright\Large $520$\end{minipage}\end{tcolorbox}\par\vspace{2.5mm}
\begin{tcolorbox}[enhanced,nobeforeafter,width=7.1cm,height=2.1cm,valign=center,colback=cnum!6,colframe=cnum,arc=2.5mm,boxrule=0.9pt,left=3mm]{\Large\bfseries\color{cnum}c)}\quad\begin{minipage}[c]{5.6cm}\raggedright\Large $0{,}52$\end{minipage}\end{tcolorbox}\hfill\begin{tcolorbox}[enhanced,nobeforeafter,width=7.1cm,height=2.1cm,valign=center,colback=cnum!6,colframe=cnum,arc=2.5mm,boxrule=0.9pt,left=3mm]{\Large\bfseries\color{cnum}d)}\quad\begin{minipage}[c]{5.6cm}\raggedright\Large $5{,}2$\end{minipage}\end{tcolorbox}
\end{tcolorbox}
\newpage
\begin{tcolorbox}[enhanced,width=16cm,height=14cm,colback=white,colframe=cnum,boxrule=2pt,arc=0mm,
left=6mm,right=6mm,top=0mm,bottom=4mm,sharp corners,
overlay={\fill[cnum] (frame.north west) rectangle ([yshift=-1.25cm]frame.north east);
\node[anchor=west,text=white,font=\sffamily\bfseries\LARGE] at ([xshift=5mm,yshift=-0.62cm]frame.north west) {Question 5};
\node[anchor=east,text=white,font=\sffamily\itshape\large] at ([xshift=-5mm,yshift=-0.62cm]frame.north east) {Nombres entiers et décimaux};}]
\vspace*{1.75cm}
\begin{minipage}[t][6.9cm][c]{\linewidth}\LARGE\centering
\begin{minipage}{0.95\linewidth}\raggedright\linespread{1.25}\selectfont Quel est le plus grand de ces quatre nombres~: $3{,}5\ ;\ 3{,}45\ ;\ 3{,}405\ ;\ 3{,}054$~?\end{minipage}
\end{minipage}

\vfill
\begin{tcolorbox}[enhanced,nobeforeafter,width=7.1cm,height=2.1cm,valign=center,colback=cnum!6,colframe=cnum,arc=2.5mm,boxrule=0.9pt,left=3mm]{\Large\bfseries\color{cnum}a)}\quad\begin{minipage}[c]{5.6cm}\raggedright\Large $3{,}054$\end{minipage}\end{tcolorbox}\hfill\begin{tcolorbox}[enhanced,nobeforeafter,width=7.1cm,height=2.1cm,valign=center,colback=cnum!6,colframe=cnum,arc=2.5mm,boxrule=0.9pt,left=3mm]{\Large\bfseries\color{cnum}b)}\quad\begin{minipage}[c]{5.6cm}\raggedright\Large $3{,}45$\end{minipage}\end{tcolorbox}\par\vspace{2.5mm}
\begin{tcolorbox}[enhanced,nobeforeafter,width=7.1cm,height=2.1cm,valign=center,colback=cnum!6,colframe=cnum,arc=2.5mm,boxrule=0.9pt,left=3mm]{\Large\bfseries\color{cnum}c)}\quad\begin{minipage}[c]{5.6cm}\raggedright\Large $3{,}405$\end{minipage}\end{tcolorbox}\hfill\begin{tcolorbox}[enhanced,nobeforeafter,width=7.1cm,height=2.1cm,valign=center,colback=cnum!6,colframe=cnum,arc=2.5mm,boxrule=0.9pt,left=3mm]{\Large\bfseries\color{cnum}d)}\quad\begin{minipage}[c]{5.6cm}\raggedright\Large $3{,}5$\end{minipage}\end{tcolorbox}
\end{tcolorbox}
\newpage
\begin{tcolorbox}[enhanced,width=16cm,height=14cm,colback=white,colframe=cnum,boxrule=2pt,arc=0mm,
left=6mm,right=6mm,top=0mm,bottom=4mm,sharp corners,
overlay={\fill[cnum] (frame.north west) rectangle ([yshift=-1.25cm]frame.north east);
\node[anchor=west,text=white,font=\sffamily\bfseries\LARGE] at ([xshift=5mm,yshift=-0.62cm]frame.north west) {Question 6};
\node[anchor=east,text=white,font=\sffamily\itshape\large] at ([xshift=-5mm,yshift=-0.62cm]frame.north east) {Nombres entiers et décimaux};}]
\vspace*{1.75cm}
\begin{minipage}[t][6.9cm][c]{\linewidth}\LARGE\centering
\begin{minipage}{0.95\linewidth}\raggedright\linespread{1.25}\selectfont Quel est le plus petit de ces quatre nombres~: $12{,}09\ ;\ 12{,}1\ ;\ 12{,}089\ ;\ 12{,}01$~?\end{minipage}
\end{minipage}

\vfill
\begin{tcolorbox}[enhanced,nobeforeafter,width=7.1cm,height=2.1cm,valign=center,colback=cnum!6,colframe=cnum,arc=2.5mm,boxrule=0.9pt,left=3mm]{\Large\bfseries\color{cnum}a)}\quad\begin{minipage}[c]{5.6cm}\raggedright\Large $12{,}1$\end{minipage}\end{tcolorbox}\hfill\begin{tcolorbox}[enhanced,nobeforeafter,width=7.1cm,height=2.1cm,valign=center,colback=cnum!6,colframe=cnum,arc=2.5mm,boxrule=0.9pt,left=3mm]{\Large\bfseries\color{cnum}b)}\quad\begin{minipage}[c]{5.6cm}\raggedright\Large $12{,}09$\end{minipage}\end{tcolorbox}\par\vspace{2.5mm}
\begin{tcolorbox}[enhanced,nobeforeafter,width=7.1cm,height=2.1cm,valign=center,colback=cnum!6,colframe=cnum,arc=2.5mm,boxrule=0.9pt,left=3mm]{\Large\bfseries\color{cnum}c)}\quad\begin{minipage}[c]{5.6cm}\raggedright\Large $12{,}01$\end{minipage}\end{tcolorbox}\hfill\begin{tcolorbox}[enhanced,nobeforeafter,width=7.1cm,height=2.1cm,valign=center,colback=cnum!6,colframe=cnum,arc=2.5mm,boxrule=0.9pt,left=3mm]{\Large\bfseries\color{cnum}d)}\quad\begin{minipage}[c]{5.6cm}\raggedright\Large $12{,}089$\end{minipage}\end{tcolorbox}
\end{tcolorbox}
\newpage
\begin{tcolorbox}[enhanced,width=16cm,height=14cm,colback=white,colframe=cnum,boxrule=2pt,arc=0mm,
left=6mm,right=6mm,top=0mm,bottom=4mm,sharp corners,
overlay={\fill[cnum] (frame.north west) rectangle ([yshift=-1.25cm]frame.north east);
\node[anchor=west,text=white,font=\sffamily\bfseries\LARGE] at ([xshift=5mm,yshift=-0.62cm]frame.north west) {Question 7};
\node[anchor=east,text=white,font=\sffamily\itshape\large] at ([xshift=-5mm,yshift=-0.62cm]frame.north east) {Nombres entiers et décimaux};}]
\vspace*{1.75cm}
\begin{minipage}[t][6.9cm][c]{\linewidth}\LARGE\centering
\begin{minipage}{0.95\linewidth}\raggedright\linespread{1.25}\selectfont Quel est l'arrondi au dixième de $12{,}46$~?\end{minipage}
\end{minipage}

\vfill
\begin{tcolorbox}[enhanced,nobeforeafter,width=7.1cm,height=2.1cm,valign=center,colback=cnum!6,colframe=cnum,arc=2.5mm,boxrule=0.9pt,left=3mm]{\Large\bfseries\color{cnum}a)}\quad\begin{minipage}[c]{5.6cm}\raggedright\Large $12$\end{minipage}\end{tcolorbox}\hfill\begin{tcolorbox}[enhanced,nobeforeafter,width=7.1cm,height=2.1cm,valign=center,colback=cnum!6,colframe=cnum,arc=2.5mm,boxrule=0.9pt,left=3mm]{\Large\bfseries\color{cnum}b)}\quad\begin{minipage}[c]{5.6cm}\raggedright\Large $12{,}46$\end{minipage}\end{tcolorbox}\par\vspace{2.5mm}
\begin{tcolorbox}[enhanced,nobeforeafter,width=7.1cm,height=2.1cm,valign=center,colback=cnum!6,colframe=cnum,arc=2.5mm,boxrule=0.9pt,left=3mm]{\Large\bfseries\color{cnum}c)}\quad\begin{minipage}[c]{5.6cm}\raggedright\Large $12{,}4$\end{minipage}\end{tcolorbox}\hfill\begin{tcolorbox}[enhanced,nobeforeafter,width=7.1cm,height=2.1cm,valign=center,colback=cnum!6,colframe=cnum,arc=2.5mm,boxrule=0.9pt,left=3mm]{\Large\bfseries\color{cnum}d)}\quad\begin{minipage}[c]{5.6cm}\raggedright\Large $12{,}5$\end{minipage}\end{tcolorbox}
\end{tcolorbox}
\newpage
\begin{tcolorbox}[enhanced,width=16cm,height=14cm,colback=white,colframe=cnum,boxrule=2pt,arc=0mm,
left=6mm,right=6mm,top=0mm,bottom=4mm,sharp corners,
overlay={\fill[cnum] (frame.north west) rectangle ([yshift=-1.25cm]frame.north east);
\node[anchor=west,text=white,font=\sffamily\bfseries\LARGE] at ([xshift=5mm,yshift=-0.62cm]frame.north west) {Question 8};
\node[anchor=east,text=white,font=\sffamily\itshape\large] at ([xshift=-5mm,yshift=-0.62cm]frame.north east) {Nombres entiers et décimaux};}]
\vspace*{1.75cm}
\begin{minipage}[t][6.9cm][c]{\linewidth}\LARGE\centering
\begin{minipage}{0.95\linewidth}\raggedright\linespread{1.25}\selectfont Quel est l'arrondi au dixième de $7{,}83$~?\end{minipage}
\end{minipage}

\vfill
\begin{tcolorbox}[enhanced,nobeforeafter,width=7.1cm,height=2.1cm,valign=center,colback=cnum!6,colframe=cnum,arc=2.5mm,boxrule=0.9pt,left=3mm]{\Large\bfseries\color{cnum}a)}\quad\begin{minipage}[c]{5.6cm}\raggedright\Large $8$\end{minipage}\end{tcolorbox}\hfill\begin{tcolorbox}[enhanced,nobeforeafter,width=7.1cm,height=2.1cm,valign=center,colback=cnum!6,colframe=cnum,arc=2.5mm,boxrule=0.9pt,left=3mm]{\Large\bfseries\color{cnum}b)}\quad\begin{minipage}[c]{5.6cm}\raggedright\Large $7{,}9$\end{minipage}\end{tcolorbox}\par\vspace{2.5mm}
\begin{tcolorbox}[enhanced,nobeforeafter,width=7.1cm,height=2.1cm,valign=center,colback=cnum!6,colframe=cnum,arc=2.5mm,boxrule=0.9pt,left=3mm]{\Large\bfseries\color{cnum}c)}\quad\begin{minipage}[c]{5.6cm}\raggedright\Large $7{,}83$\end{minipage}\end{tcolorbox}\hfill\begin{tcolorbox}[enhanced,nobeforeafter,width=7.1cm,height=2.1cm,valign=center,colback=cnum!6,colframe=cnum,arc=2.5mm,boxrule=0.9pt,left=3mm]{\Large\bfseries\color{cnum}d)}\quad\begin{minipage}[c]{5.6cm}\raggedright\Large $7{,}8$\end{minipage}\end{tcolorbox}
\end{tcolorbox}
\newpage
\begin{tcolorbox}[enhanced,width=16cm,height=14cm,colback=white,colframe=cnum,boxrule=2pt,arc=0mm,
left=6mm,right=6mm,top=0mm,bottom=4mm,sharp corners,
overlay={\fill[cnum] (frame.north west) rectangle ([yshift=-1.25cm]frame.north east);
\node[anchor=west,text=white,font=\sffamily\bfseries\LARGE] at ([xshift=5mm,yshift=-0.62cm]frame.north west) {Question 9};
\node[anchor=east,text=white,font=\sffamily\itshape\large] at ([xshift=-5mm,yshift=-0.62cm]frame.north east) {Nombres entiers et décimaux};}]
\vspace*{1.75cm}
\begin{minipage}[t][6.9cm][c]{\linewidth}\LARGE\centering
\begin{minipage}{0.95\linewidth}\raggedright\linespread{1.25}\selectfont Que vaut $3{,}45\times100$~?\end{minipage}
\end{minipage}

\vfill
\begin{tcolorbox}[enhanced,nobeforeafter,width=7.1cm,height=2.1cm,valign=center,colback=cnum!6,colframe=cnum,arc=2.5mm,boxrule=0.9pt,left=3mm]{\Large\bfseries\color{cnum}a)}\quad\begin{minipage}[c]{5.6cm}\raggedright\Large $3\,450$\end{minipage}\end{tcolorbox}\hfill\begin{tcolorbox}[enhanced,nobeforeafter,width=7.1cm,height=2.1cm,valign=center,colback=cnum!6,colframe=cnum,arc=2.5mm,boxrule=0.9pt,left=3mm]{\Large\bfseries\color{cnum}b)}\quad\begin{minipage}[c]{5.6cm}\raggedright\Large $34{,}5$\end{minipage}\end{tcolorbox}\par\vspace{2.5mm}
\begin{tcolorbox}[enhanced,nobeforeafter,width=7.1cm,height=2.1cm,valign=center,colback=cnum!6,colframe=cnum,arc=2.5mm,boxrule=0.9pt,left=3mm]{\Large\bfseries\color{cnum}c)}\quad\begin{minipage}[c]{5.6cm}\raggedright\Large $345$\end{minipage}\end{tcolorbox}\hfill\begin{tcolorbox}[enhanced,nobeforeafter,width=7.1cm,height=2.1cm,valign=center,colback=cnum!6,colframe=cnum,arc=2.5mm,boxrule=0.9pt,left=3mm]{\Large\bfseries\color{cnum}d)}\quad\begin{minipage}[c]{5.6cm}\raggedright\Large $34\,500$\end{minipage}\end{tcolorbox}
\end{tcolorbox}
\newpage
\begin{tcolorbox}[enhanced,width=16cm,height=14cm,colback=white,colframe=cnum,boxrule=2pt,arc=0mm,
left=6mm,right=6mm,top=0mm,bottom=4mm,sharp corners,
overlay={\fill[cnum] (frame.north west) rectangle ([yshift=-1.25cm]frame.north east);
\node[anchor=west,text=white,font=\sffamily\bfseries\LARGE] at ([xshift=5mm,yshift=-0.62cm]frame.north west) {Question 10};
\node[anchor=east,text=white,font=\sffamily\itshape\large] at ([xshift=-5mm,yshift=-0.62cm]frame.north east) {Nombres entiers et décimaux};}]
\vspace*{1.75cm}
\begin{minipage}[t][6.9cm][c]{\linewidth}\LARGE\centering
\begin{minipage}{0.95\linewidth}\raggedright\linespread{1.25}\selectfont Que vaut $0{,}72\times10$~?\end{minipage}
\end{minipage}

\vfill
\begin{tcolorbox}[enhanced,nobeforeafter,width=7.1cm,height=2.1cm,valign=center,colback=cnum!6,colframe=cnum,arc=2.5mm,boxrule=0.9pt,left=3mm]{\Large\bfseries\color{cnum}a)}\quad\begin{minipage}[c]{5.6cm}\raggedright\Large $0{,}72$\end{minipage}\end{tcolorbox}\hfill\begin{tcolorbox}[enhanced,nobeforeafter,width=7.1cm,height=2.1cm,valign=center,colback=cnum!6,colframe=cnum,arc=2.5mm,boxrule=0.9pt,left=3mm]{\Large\bfseries\color{cnum}b)}\quad\begin{minipage}[c]{5.6cm}\raggedright\Large $0{,}072$\end{minipage}\end{tcolorbox}\par\vspace{2.5mm}
\begin{tcolorbox}[enhanced,nobeforeafter,width=7.1cm,height=2.1cm,valign=center,colback=cnum!6,colframe=cnum,arc=2.5mm,boxrule=0.9pt,left=3mm]{\Large\bfseries\color{cnum}c)}\quad\begin{minipage}[c]{5.6cm}\raggedright\Large $72$\end{minipage}\end{tcolorbox}\hfill\begin{tcolorbox}[enhanced,nobeforeafter,width=7.1cm,height=2.1cm,valign=center,colback=cnum!6,colframe=cnum,arc=2.5mm,boxrule=0.9pt,left=3mm]{\Large\bfseries\color{cnum}d)}\quad\begin{minipage}[c]{5.6cm}\raggedright\Large $7{,}2$\end{minipage}\end{tcolorbox}
\end{tcolorbox}
\newpage
\begin{tcolorbox}[enhanced,width=16cm,height=14cm,colback=white,colframe=cnum,boxrule=2pt,arc=0mm,
left=6mm,right=6mm,top=0mm,bottom=4mm,sharp corners,
overlay={\fill[cnum] (frame.north west) rectangle ([yshift=-1.25cm]frame.north east);
\node[anchor=west,text=white,font=\sffamily\bfseries\LARGE] at ([xshift=5mm,yshift=-0.62cm]frame.north west) {Question 11};
\node[anchor=east,text=white,font=\sffamily\itshape\large] at ([xshift=-5mm,yshift=-0.62cm]frame.north east) {Nombres entiers et décimaux};}]
\vspace*{1.75cm}
\begin{minipage}[t][6.9cm][c]{\linewidth}\LARGE\centering
\begin{minipage}{0.95\linewidth}\raggedright\linespread{1.25}\selectfont Que vaut $52\div100$~?\end{minipage}
\end{minipage}

\vfill
\begin{tcolorbox}[enhanced,nobeforeafter,width=7.1cm,height=2.1cm,valign=center,colback=cnum!6,colframe=cnum,arc=2.5mm,boxrule=0.9pt,left=3mm]{\Large\bfseries\color{cnum}a)}\quad\begin{minipage}[c]{5.6cm}\raggedright\Large $0{,}52$\end{minipage}\end{tcolorbox}\hfill\begin{tcolorbox}[enhanced,nobeforeafter,width=7.1cm,height=2.1cm,valign=center,colback=cnum!6,colframe=cnum,arc=2.5mm,boxrule=0.9pt,left=3mm]{\Large\bfseries\color{cnum}b)}\quad\begin{minipage}[c]{5.6cm}\raggedright\Large $0{,}052$\end{minipage}\end{tcolorbox}\par\vspace{2.5mm}
\begin{tcolorbox}[enhanced,nobeforeafter,width=7.1cm,height=2.1cm,valign=center,colback=cnum!6,colframe=cnum,arc=2.5mm,boxrule=0.9pt,left=3mm]{\Large\bfseries\color{cnum}c)}\quad\begin{minipage}[c]{5.6cm}\raggedright\Large $5\,200$\end{minipage}\end{tcolorbox}\hfill\begin{tcolorbox}[enhanced,nobeforeafter,width=7.1cm,height=2.1cm,valign=center,colback=cnum!6,colframe=cnum,arc=2.5mm,boxrule=0.9pt,left=3mm]{\Large\bfseries\color{cnum}d)}\quad\begin{minipage}[c]{5.6cm}\raggedright\Large $5{,}2$\end{minipage}\end{tcolorbox}
\end{tcolorbox}
\newpage
\begin{tcolorbox}[enhanced,width=16cm,height=14cm,colback=white,colframe=cnum,boxrule=2pt,arc=0mm,
left=6mm,right=6mm,top=0mm,bottom=4mm,sharp corners,
overlay={\fill[cnum] (frame.north west) rectangle ([yshift=-1.25cm]frame.north east);
\node[anchor=west,text=white,font=\sffamily\bfseries\LARGE] at ([xshift=5mm,yshift=-0.62cm]frame.north west) {Question 12};
\node[anchor=east,text=white,font=\sffamily\itshape\large] at ([xshift=-5mm,yshift=-0.62cm]frame.north east) {Nombres entiers et décimaux};}]
\vspace*{1.75cm}
\begin{minipage}[t][6.9cm][c]{\linewidth}\LARGE\centering
\begin{minipage}{0.95\linewidth}\raggedright\linespread{1.25}\selectfont Que vaut $7{,}5\div10$~?\end{minipage}
\end{minipage}

\vfill
\begin{tcolorbox}[enhanced,nobeforeafter,width=7.1cm,height=2.1cm,valign=center,colback=cnum!6,colframe=cnum,arc=2.5mm,boxrule=0.9pt,left=3mm]{\Large\bfseries\color{cnum}a)}\quad\begin{minipage}[c]{5.6cm}\raggedright\Large $7{,}5$\end{minipage}\end{tcolorbox}\hfill\begin{tcolorbox}[enhanced,nobeforeafter,width=7.1cm,height=2.1cm,valign=center,colback=cnum!6,colframe=cnum,arc=2.5mm,boxrule=0.9pt,left=3mm]{\Large\bfseries\color{cnum}b)}\quad\begin{minipage}[c]{5.6cm}\raggedright\Large $0{,}75$\end{minipage}\end{tcolorbox}\par\vspace{2.5mm}
\begin{tcolorbox}[enhanced,nobeforeafter,width=7.1cm,height=2.1cm,valign=center,colback=cnum!6,colframe=cnum,arc=2.5mm,boxrule=0.9pt,left=3mm]{\Large\bfseries\color{cnum}c)}\quad\begin{minipage}[c]{5.6cm}\raggedright\Large $75$\end{minipage}\end{tcolorbox}\hfill\begin{tcolorbox}[enhanced,nobeforeafter,width=7.1cm,height=2.1cm,valign=center,colback=cnum!6,colframe=cnum,arc=2.5mm,boxrule=0.9pt,left=3mm]{\Large\bfseries\color{cnum}d)}\quad\begin{minipage}[c]{5.6cm}\raggedright\Large $0{,}075$\end{minipage}\end{tcolorbox}
\end{tcolorbox}
\newpage
\begin{tcolorbox}[enhanced,width=16cm,height=14cm,colback=white,colframe=cnum,boxrule=2pt,arc=0mm,
left=6mm,right=6mm,top=0mm,bottom=4mm,sharp corners,
overlay={\fill[cnum] (frame.north west) rectangle ([yshift=-1.25cm]frame.north east);
\node[anchor=west,text=white,font=\sffamily\bfseries\LARGE] at ([xshift=5mm,yshift=-0.62cm]frame.north west) {Question 13};
\node[anchor=east,text=white,font=\sffamily\itshape\large] at ([xshift=-5mm,yshift=-0.62cm]frame.north east) {Nombres entiers et décimaux};}]
\vspace*{1.75cm}
\begin{minipage}[t][6.9cm][c]{\linewidth}\LARGE\centering
\begin{minipage}{0.95\linewidth}\raggedright\linespread{1.25}\selectfont Comment s'écrit en chiffres le nombre «~deux millions trois cent mille quarante~»~?\end{minipage}
\end{minipage}

\vfill
\begin{tcolorbox}[enhanced,nobeforeafter,width=7.1cm,height=2.1cm,valign=center,colback=cnum!6,colframe=cnum,arc=2.5mm,boxrule=0.9pt,left=3mm]{\Large\bfseries\color{cnum}a)}\quad\begin{minipage}[c]{5.6cm}\raggedright\Large $2\,030\,040$\end{minipage}\end{tcolorbox}\hfill\begin{tcolorbox}[enhanced,nobeforeafter,width=7.1cm,height=2.1cm,valign=center,colback=cnum!6,colframe=cnum,arc=2.5mm,boxrule=0.9pt,left=3mm]{\Large\bfseries\color{cnum}b)}\quad\begin{minipage}[c]{5.6cm}\raggedright\Large $230\,040$\end{minipage}\end{tcolorbox}\par\vspace{2.5mm}
\begin{tcolorbox}[enhanced,nobeforeafter,width=7.1cm,height=2.1cm,valign=center,colback=cnum!6,colframe=cnum,arc=2.5mm,boxrule=0.9pt,left=3mm]{\Large\bfseries\color{cnum}c)}\quad\begin{minipage}[c]{5.6cm}\raggedright\Large $2\,300\,040$\end{minipage}\end{tcolorbox}\hfill\begin{tcolorbox}[enhanced,nobeforeafter,width=7.1cm,height=2.1cm,valign=center,colback=cnum!6,colframe=cnum,arc=2.5mm,boxrule=0.9pt,left=3mm]{\Large\bfseries\color{cnum}d)}\quad\begin{minipage}[c]{5.6cm}\raggedright\Large $2\,300\,400$\end{minipage}\end{tcolorbox}
\end{tcolorbox}
\newpage
\begin{tcolorbox}[enhanced,width=16cm,height=14cm,colback=white,colframe=cnum,boxrule=2pt,arc=0mm,
left=6mm,right=6mm,top=0mm,bottom=4mm,sharp corners,
overlay={\fill[cnum] (frame.north west) rectangle ([yshift=-1.25cm]frame.north east);
\node[anchor=west,text=white,font=\sffamily\bfseries\LARGE] at ([xshift=5mm,yshift=-0.62cm]frame.north west) {Question 14};
\node[anchor=east,text=white,font=\sffamily\itshape\large] at ([xshift=-5mm,yshift=-0.62cm]frame.north east) {Nombres entiers et décimaux};}]
\vspace*{1.75cm}
\begin{minipage}[t][6.9cm][c]{\linewidth}\LARGE\centering
\begin{minipage}{0.95\linewidth}\raggedright\linespread{1.25}\selectfont Comment s'écrit en chiffres le nombre «~quinze mille six cent huit~»~?\end{minipage}
\end{minipage}

\vfill
\begin{tcolorbox}[enhanced,nobeforeafter,width=7.1cm,height=2.1cm,valign=center,colback=cnum!6,colframe=cnum,arc=2.5mm,boxrule=0.9pt,left=3mm]{\Large\bfseries\color{cnum}a)}\quad\begin{minipage}[c]{5.6cm}\raggedright\Large $1\,568$\end{minipage}\end{tcolorbox}\hfill\begin{tcolorbox}[enhanced,nobeforeafter,width=7.1cm,height=2.1cm,valign=center,colback=cnum!6,colframe=cnum,arc=2.5mm,boxrule=0.9pt,left=3mm]{\Large\bfseries\color{cnum}b)}\quad\begin{minipage}[c]{5.6cm}\raggedright\Large $15\,680$\end{minipage}\end{tcolorbox}\par\vspace{2.5mm}
\begin{tcolorbox}[enhanced,nobeforeafter,width=7.1cm,height=2.1cm,valign=center,colback=cnum!6,colframe=cnum,arc=2.5mm,boxrule=0.9pt,left=3mm]{\Large\bfseries\color{cnum}c)}\quad\begin{minipage}[c]{5.6cm}\raggedright\Large $15\,608$\end{minipage}\end{tcolorbox}\hfill\begin{tcolorbox}[enhanced,nobeforeafter,width=7.1cm,height=2.1cm,valign=center,colback=cnum!6,colframe=cnum,arc=2.5mm,boxrule=0.9pt,left=3mm]{\Large\bfseries\color{cnum}d)}\quad\begin{minipage}[c]{5.6cm}\raggedright\Large $156\,008$\end{minipage}\end{tcolorbox}
\end{tcolorbox}
\newpage
\begin{tcolorbox}[enhanced,width=16cm,height=14cm,colback=white,colframe=cnum,boxrule=2pt,arc=0mm,
left=6mm,right=6mm,top=0mm,bottom=4mm,sharp corners,
overlay={\fill[cnum] (frame.north west) rectangle ([yshift=-1.25cm]frame.north east);
\node[anchor=west,text=white,font=\sffamily\bfseries\LARGE] at ([xshift=5mm,yshift=-0.62cm]frame.north west) {Question 15};
\node[anchor=east,text=white,font=\sffamily\itshape\large] at ([xshift=-5mm,yshift=-0.62cm]frame.north east) {Nombres entiers et décimaux};}]
\vspace*{1.75cm}
\begin{minipage}[t][6.9cm][c]{\linewidth}\LARGE\centering
\begin{minipage}{0.95\linewidth}\raggedright\linespread{1.25}\selectfont Quel nombre est compris entre $2{,}3$ et $2{,}4$~?\end{minipage}
\end{minipage}

\vfill
\begin{tcolorbox}[enhanced,nobeforeafter,width=7.1cm,height=2.1cm,valign=center,colback=cnum!6,colframe=cnum,arc=2.5mm,boxrule=0.9pt,left=3mm]{\Large\bfseries\color{cnum}a)}\quad\begin{minipage}[c]{5.6cm}\raggedright\Large $2{,}5$\end{minipage}\end{tcolorbox}\hfill\begin{tcolorbox}[enhanced,nobeforeafter,width=7.1cm,height=2.1cm,valign=center,colback=cnum!6,colframe=cnum,arc=2.5mm,boxrule=0.9pt,left=3mm]{\Large\bfseries\color{cnum}b)}\quad\begin{minipage}[c]{5.6cm}\raggedright\Large $2{,}04$\end{minipage}\end{tcolorbox}\par\vspace{2.5mm}
\begin{tcolorbox}[enhanced,nobeforeafter,width=7.1cm,height=2.1cm,valign=center,colback=cnum!6,colframe=cnum,arc=2.5mm,boxrule=0.9pt,left=3mm]{\Large\bfseries\color{cnum}c)}\quad\begin{minipage}[c]{5.6cm}\raggedright\Large $2{,}29$\end{minipage}\end{tcolorbox}\hfill\begin{tcolorbox}[enhanced,nobeforeafter,width=7.1cm,height=2.1cm,valign=center,colback=cnum!6,colframe=cnum,arc=2.5mm,boxrule=0.9pt,left=3mm]{\Large\bfseries\color{cnum}d)}\quad\begin{minipage}[c]{5.6cm}\raggedright\Large $2{,}35$\end{minipage}\end{tcolorbox}
\end{tcolorbox}
\newpage
\begin{tcolorbox}[enhanced,width=16cm,height=14cm,colback=white,colframe=cnum,boxrule=2pt,arc=0mm,
left=6mm,right=6mm,top=0mm,bottom=4mm,sharp corners,
overlay={\fill[cnum] (frame.north west) rectangle ([yshift=-1.25cm]frame.north east);
\node[anchor=west,text=white,font=\sffamily\bfseries\LARGE] at ([xshift=5mm,yshift=-0.62cm]frame.north west) {Question 16};
\node[anchor=east,text=white,font=\sffamily\itshape\large] at ([xshift=-5mm,yshift=-0.62cm]frame.north east) {Nombres entiers et décimaux};}]
\vspace*{1.75cm}
\begin{minipage}[t][6.9cm][c]{\linewidth}\LARGE\centering
\begin{minipage}{0.95\linewidth}\raggedright\linespread{1.25}\selectfont Quel nombre est compris entre $7$ et $7{,}1$~?\end{minipage}
\end{minipage}

\vfill
\begin{tcolorbox}[enhanced,nobeforeafter,width=7.1cm,height=2.1cm,valign=center,colback=cnum!6,colframe=cnum,arc=2.5mm,boxrule=0.9pt,left=3mm]{\Large\bfseries\color{cnum}a)}\quad\begin{minipage}[c]{5.6cm}\raggedright\Large $7{,}15$\end{minipage}\end{tcolorbox}\hfill\begin{tcolorbox}[enhanced,nobeforeafter,width=7.1cm,height=2.1cm,valign=center,colback=cnum!6,colframe=cnum,arc=2.5mm,boxrule=0.9pt,left=3mm]{\Large\bfseries\color{cnum}b)}\quad\begin{minipage}[c]{5.6cm}\raggedright\Large $7{,}2$\end{minipage}\end{tcolorbox}\par\vspace{2.5mm}
\begin{tcolorbox}[enhanced,nobeforeafter,width=7.1cm,height=2.1cm,valign=center,colback=cnum!6,colframe=cnum,arc=2.5mm,boxrule=0.9pt,left=3mm]{\Large\bfseries\color{cnum}c)}\quad\begin{minipage}[c]{5.6cm}\raggedright\Large $7{,}05$\end{minipage}\end{tcolorbox}\hfill\begin{tcolorbox}[enhanced,nobeforeafter,width=7.1cm,height=2.1cm,valign=center,colback=cnum!6,colframe=cnum,arc=2.5mm,boxrule=0.9pt,left=3mm]{\Large\bfseries\color{cnum}d)}\quad\begin{minipage}[c]{5.6cm}\raggedright\Large $6{,}9$\end{minipage}\end{tcolorbox}
\end{tcolorbox}
\newpage
\begin{tcolorbox}[enhanced,width=16cm,height=14cm,colback=white,colframe=ccal,boxrule=2pt,arc=0mm,
left=6mm,right=6mm,top=0mm,bottom=4mm,sharp corners,
overlay={\fill[ccal] (frame.north west) rectangle ([yshift=-1.25cm]frame.north east);
\node[anchor=west,text=white,font=\sffamily\bfseries\LARGE] at ([xshift=5mm,yshift=-0.62cm]frame.north west) {Question 17};
\node[anchor=east,text=white,font=\sffamily\itshape\large] at ([xshift=-5mm,yshift=-0.62cm]frame.north east) {Calculs};}]
\vspace*{1.75cm}
\begin{minipage}[t][6.9cm][c]{\linewidth}\LARGE\centering
\begin{minipage}{0.95\linewidth}\raggedright\linespread{1.25}\selectfont Que vaut $12{,}5+7{,}85$~?\end{minipage}
\end{minipage}

\vfill
\begin{tcolorbox}[enhanced,nobeforeafter,width=7.1cm,height=2.1cm,valign=center,colback=ccal!6,colframe=ccal,arc=2.5mm,boxrule=0.9pt,left=3mm]{\Large\bfseries\color{ccal}a)}\quad\begin{minipage}[c]{5.6cm}\raggedright\Large $19{,}35$\end{minipage}\end{tcolorbox}\hfill\begin{tcolorbox}[enhanced,nobeforeafter,width=7.1cm,height=2.1cm,valign=center,colback=ccal!6,colframe=ccal,arc=2.5mm,boxrule=0.9pt,left=3mm]{\Large\bfseries\color{ccal}b)}\quad\begin{minipage}[c]{5.6cm}\raggedright\Large $20{,}35$\end{minipage}\end{tcolorbox}\par\vspace{2.5mm}
\begin{tcolorbox}[enhanced,nobeforeafter,width=7.1cm,height=2.1cm,valign=center,colback=ccal!6,colframe=ccal,arc=2.5mm,boxrule=0.9pt,left=3mm]{\Large\bfseries\color{ccal}c)}\quad\begin{minipage}[c]{5.6cm}\raggedright\Large $20{,}45$\end{minipage}\end{tcolorbox}\hfill\begin{tcolorbox}[enhanced,nobeforeafter,width=7.1cm,height=2.1cm,valign=center,colback=ccal!6,colframe=ccal,arc=2.5mm,boxrule=0.9pt,left=3mm]{\Large\bfseries\color{ccal}d)}\quad\begin{minipage}[c]{5.6cm}\raggedright\Large $21{,}35$\end{minipage}\end{tcolorbox}
\end{tcolorbox}
\newpage
\begin{tcolorbox}[enhanced,width=16cm,height=14cm,colback=white,colframe=ccal,boxrule=2pt,arc=0mm,
left=6mm,right=6mm,top=0mm,bottom=4mm,sharp corners,
overlay={\fill[ccal] (frame.north west) rectangle ([yshift=-1.25cm]frame.north east);
\node[anchor=west,text=white,font=\sffamily\bfseries\LARGE] at ([xshift=5mm,yshift=-0.62cm]frame.north west) {Question 18};
\node[anchor=east,text=white,font=\sffamily\itshape\large] at ([xshift=-5mm,yshift=-0.62cm]frame.north east) {Calculs};}]
\vspace*{1.75cm}
\begin{minipage}[t][6.9cm][c]{\linewidth}\LARGE\centering
\begin{minipage}{0.95\linewidth}\raggedright\linespread{1.25}\selectfont Que vaut $3{,}07+14{,}6$~?\end{minipage}
\end{minipage}

\vfill
\begin{tcolorbox}[enhanced,nobeforeafter,width=7.1cm,height=2.1cm,valign=center,colback=ccal!6,colframe=ccal,arc=2.5mm,boxrule=0.9pt,left=3mm]{\Large\bfseries\color{ccal}a)}\quad\begin{minipage}[c]{5.6cm}\raggedright\Large $16{,}67$\end{minipage}\end{tcolorbox}\hfill\begin{tcolorbox}[enhanced,nobeforeafter,width=7.1cm,height=2.1cm,valign=center,colback=ccal!6,colframe=ccal,arc=2.5mm,boxrule=0.9pt,left=3mm]{\Large\bfseries\color{ccal}b)}\quad\begin{minipage}[c]{5.6cm}\raggedright\Large $17{,}77$\end{minipage}\end{tcolorbox}\par\vspace{2.5mm}
\begin{tcolorbox}[enhanced,nobeforeafter,width=7.1cm,height=2.1cm,valign=center,colback=ccal!6,colframe=ccal,arc=2.5mm,boxrule=0.9pt,left=3mm]{\Large\bfseries\color{ccal}c)}\quad\begin{minipage}[c]{5.6cm}\raggedright\Large $18{,}67$\end{minipage}\end{tcolorbox}\hfill\begin{tcolorbox}[enhanced,nobeforeafter,width=7.1cm,height=2.1cm,valign=center,colback=ccal!6,colframe=ccal,arc=2.5mm,boxrule=0.9pt,left=3mm]{\Large\bfseries\color{ccal}d)}\quad\begin{minipage}[c]{5.6cm}\raggedright\Large $17{,}67$\end{minipage}\end{tcolorbox}
\end{tcolorbox}
\newpage
\begin{tcolorbox}[enhanced,width=16cm,height=14cm,colback=white,colframe=ccal,boxrule=2pt,arc=0mm,
left=6mm,right=6mm,top=0mm,bottom=4mm,sharp corners,
overlay={\fill[ccal] (frame.north west) rectangle ([yshift=-1.25cm]frame.north east);
\node[anchor=west,text=white,font=\sffamily\bfseries\LARGE] at ([xshift=5mm,yshift=-0.62cm]frame.north west) {Question 19};
\node[anchor=east,text=white,font=\sffamily\itshape\large] at ([xshift=-5mm,yshift=-0.62cm]frame.north east) {Calculs};}]
\vspace*{1.75cm}
\begin{minipage}[t][6.9cm][c]{\linewidth}\LARGE\centering
\begin{minipage}{0.95\linewidth}\raggedright\linespread{1.25}\selectfont Que vaut $20-6{,}35$~?\end{minipage}
\end{minipage}

\vfill
\begin{tcolorbox}[enhanced,nobeforeafter,width=7.1cm,height=2.1cm,valign=center,colback=ccal!6,colframe=ccal,arc=2.5mm,boxrule=0.9pt,left=3mm]{\Large\bfseries\color{ccal}a)}\quad\begin{minipage}[c]{5.6cm}\raggedright\Large $13{,}75$\end{minipage}\end{tcolorbox}\hfill\begin{tcolorbox}[enhanced,nobeforeafter,width=7.1cm,height=2.1cm,valign=center,colback=ccal!6,colframe=ccal,arc=2.5mm,boxrule=0.9pt,left=3mm]{\Large\bfseries\color{ccal}b)}\quad\begin{minipage}[c]{5.6cm}\raggedright\Large $14{,}65$\end{minipage}\end{tcolorbox}\par\vspace{2.5mm}
\begin{tcolorbox}[enhanced,nobeforeafter,width=7.1cm,height=2.1cm,valign=center,colback=ccal!6,colframe=ccal,arc=2.5mm,boxrule=0.9pt,left=3mm]{\Large\bfseries\color{ccal}c)}\quad\begin{minipage}[c]{5.6cm}\raggedright\Large $13{,}55$\end{minipage}\end{tcolorbox}\hfill\begin{tcolorbox}[enhanced,nobeforeafter,width=7.1cm,height=2.1cm,valign=center,colback=ccal!6,colframe=ccal,arc=2.5mm,boxrule=0.9pt,left=3mm]{\Large\bfseries\color{ccal}d)}\quad\begin{minipage}[c]{5.6cm}\raggedright\Large $13{,}65$\end{minipage}\end{tcolorbox}
\end{tcolorbox}
\newpage
\begin{tcolorbox}[enhanced,width=16cm,height=14cm,colback=white,colframe=ccal,boxrule=2pt,arc=0mm,
left=6mm,right=6mm,top=0mm,bottom=4mm,sharp corners,
overlay={\fill[ccal] (frame.north west) rectangle ([yshift=-1.25cm]frame.north east);
\node[anchor=west,text=white,font=\sffamily\bfseries\LARGE] at ([xshift=5mm,yshift=-0.62cm]frame.north west) {Question 20};
\node[anchor=east,text=white,font=\sffamily\itshape\large] at ([xshift=-5mm,yshift=-0.62cm]frame.north east) {Calculs};}]
\vspace*{1.75cm}
\begin{minipage}[t][6.9cm][c]{\linewidth}\LARGE\centering
\begin{minipage}{0.95\linewidth}\raggedright\linespread{1.25}\selectfont Que vaut $15{,}4-8{,}75$~?\end{minipage}
\end{minipage}

\vfill
\begin{tcolorbox}[enhanced,nobeforeafter,width=7.1cm,height=2.1cm,valign=center,colback=ccal!6,colframe=ccal,arc=2.5mm,boxrule=0.9pt,left=3mm]{\Large\bfseries\color{ccal}a)}\quad\begin{minipage}[c]{5.6cm}\raggedright\Large $6{,}55$\end{minipage}\end{tcolorbox}\hfill\begin{tcolorbox}[enhanced,nobeforeafter,width=7.1cm,height=2.1cm,valign=center,colback=ccal!6,colframe=ccal,arc=2.5mm,boxrule=0.9pt,left=3mm]{\Large\bfseries\color{ccal}b)}\quad\begin{minipage}[c]{5.6cm}\raggedright\Large $6{,}65$\end{minipage}\end{tcolorbox}\par\vspace{2.5mm}
\begin{tcolorbox}[enhanced,nobeforeafter,width=7.1cm,height=2.1cm,valign=center,colback=ccal!6,colframe=ccal,arc=2.5mm,boxrule=0.9pt,left=3mm]{\Large\bfseries\color{ccal}c)}\quad\begin{minipage}[c]{5.6cm}\raggedright\Large $7{,}65$\end{minipage}\end{tcolorbox}\hfill\begin{tcolorbox}[enhanced,nobeforeafter,width=7.1cm,height=2.1cm,valign=center,colback=ccal!6,colframe=ccal,arc=2.5mm,boxrule=0.9pt,left=3mm]{\Large\bfseries\color{ccal}d)}\quad\begin{minipage}[c]{5.6cm}\raggedright\Large $6{,}75$\end{minipage}\end{tcolorbox}
\end{tcolorbox}
\newpage
\begin{tcolorbox}[enhanced,width=16cm,height=14cm,colback=white,colframe=ccal,boxrule=2pt,arc=0mm,
left=6mm,right=6mm,top=0mm,bottom=4mm,sharp corners,
overlay={\fill[ccal] (frame.north west) rectangle ([yshift=-1.25cm]frame.north east);
\node[anchor=west,text=white,font=\sffamily\bfseries\LARGE] at ([xshift=5mm,yshift=-0.62cm]frame.north west) {Question 21};
\node[anchor=east,text=white,font=\sffamily\itshape\large] at ([xshift=-5mm,yshift=-0.62cm]frame.north east) {Calculs};}]
\vspace*{1.75cm}
\begin{minipage}[t][6.9cm][c]{\linewidth}\LARGE\centering
\begin{minipage}{0.95\linewidth}\raggedright\linespread{1.25}\selectfont Que vaut $2{,}5\times1{,}2$~?\end{minipage}
\end{minipage}

\vfill
\begin{tcolorbox}[enhanced,nobeforeafter,width=7.1cm,height=2.1cm,valign=center,colback=ccal!6,colframe=ccal,arc=2.5mm,boxrule=0.9pt,left=3mm]{\Large\bfseries\color{ccal}a)}\quad\begin{minipage}[c]{5.6cm}\raggedright\Large $3$\end{minipage}\end{tcolorbox}\hfill\begin{tcolorbox}[enhanced,nobeforeafter,width=7.1cm,height=2.1cm,valign=center,colback=ccal!6,colframe=ccal,arc=2.5mm,boxrule=0.9pt,left=3mm]{\Large\bfseries\color{ccal}b)}\quad\begin{minipage}[c]{5.6cm}\raggedright\Large $3{,}7$\end{minipage}\end{tcolorbox}\par\vspace{2.5mm}
\begin{tcolorbox}[enhanced,nobeforeafter,width=7.1cm,height=2.1cm,valign=center,colback=ccal!6,colframe=ccal,arc=2.5mm,boxrule=0.9pt,left=3mm]{\Large\bfseries\color{ccal}c)}\quad\begin{minipage}[c]{5.6cm}\raggedright\Large $0{,}3$\end{minipage}\end{tcolorbox}\hfill\begin{tcolorbox}[enhanced,nobeforeafter,width=7.1cm,height=2.1cm,valign=center,colback=ccal!6,colframe=ccal,arc=2.5mm,boxrule=0.9pt,left=3mm]{\Large\bfseries\color{ccal}d)}\quad\begin{minipage}[c]{5.6cm}\raggedright\Large $30$\end{minipage}\end{tcolorbox}
\end{tcolorbox}
\newpage
\begin{tcolorbox}[enhanced,width=16cm,height=14cm,colback=white,colframe=ccal,boxrule=2pt,arc=0mm,
left=6mm,right=6mm,top=0mm,bottom=4mm,sharp corners,
overlay={\fill[ccal] (frame.north west) rectangle ([yshift=-1.25cm]frame.north east);
\node[anchor=west,text=white,font=\sffamily\bfseries\LARGE] at ([xshift=5mm,yshift=-0.62cm]frame.north west) {Question 22};
\node[anchor=east,text=white,font=\sffamily\itshape\large] at ([xshift=-5mm,yshift=-0.62cm]frame.north east) {Calculs};}]
\vspace*{1.75cm}
\begin{minipage}[t][6.9cm][c]{\linewidth}\LARGE\centering
\begin{minipage}{0.95\linewidth}\raggedright\linespread{1.25}\selectfont Que vaut $0{,}4\times0{,}3$~?\end{minipage}
\end{minipage}

\vfill
\begin{tcolorbox}[enhanced,nobeforeafter,width=7.1cm,height=2.1cm,valign=center,colback=ccal!6,colframe=ccal,arc=2.5mm,boxrule=0.9pt,left=3mm]{\Large\bfseries\color{ccal}a)}\quad\begin{minipage}[c]{5.6cm}\raggedright\Large $0{,}7$\end{minipage}\end{tcolorbox}\hfill\begin{tcolorbox}[enhanced,nobeforeafter,width=7.1cm,height=2.1cm,valign=center,colback=ccal!6,colframe=ccal,arc=2.5mm,boxrule=0.9pt,left=3mm]{\Large\bfseries\color{ccal}b)}\quad\begin{minipage}[c]{5.6cm}\raggedright\Large $0{,}012$\end{minipage}\end{tcolorbox}\par\vspace{2.5mm}
\begin{tcolorbox}[enhanced,nobeforeafter,width=7.1cm,height=2.1cm,valign=center,colback=ccal!6,colframe=ccal,arc=2.5mm,boxrule=0.9pt,left=3mm]{\Large\bfseries\color{ccal}c)}\quad\begin{minipage}[c]{5.6cm}\raggedright\Large $1{,}2$\end{minipage}\end{tcolorbox}\hfill\begin{tcolorbox}[enhanced,nobeforeafter,width=7.1cm,height=2.1cm,valign=center,colback=ccal!6,colframe=ccal,arc=2.5mm,boxrule=0.9pt,left=3mm]{\Large\bfseries\color{ccal}d)}\quad\begin{minipage}[c]{5.6cm}\raggedright\Large $0{,}12$\end{minipage}\end{tcolorbox}
\end{tcolorbox}
\newpage
\begin{tcolorbox}[enhanced,width=16cm,height=14cm,colback=white,colframe=ccal,boxrule=2pt,arc=0mm,
left=6mm,right=6mm,top=0mm,bottom=4mm,sharp corners,
overlay={\fill[ccal] (frame.north west) rectangle ([yshift=-1.25cm]frame.north east);
\node[anchor=west,text=white,font=\sffamily\bfseries\LARGE] at ([xshift=5mm,yshift=-0.62cm]frame.north west) {Question 23};
\node[anchor=east,text=white,font=\sffamily\itshape\large] at ([xshift=-5mm,yshift=-0.62cm]frame.north east) {Calculs};}]
\vspace*{1.75cm}
\begin{minipage}[t][6.9cm][c]{\linewidth}\LARGE\centering
\begin{minipage}{0.95\linewidth}\raggedright\linespread{1.25}\selectfont Dans la division euclidienne de $157$ par $12$, quels sont le quotient $q$ et le reste $r$~?\end{minipage}
\end{minipage}

\vfill
\begin{tcolorbox}[enhanced,nobeforeafter,width=7.1cm,height=2.1cm,valign=center,colback=ccal!6,colframe=ccal,arc=2.5mm,boxrule=0.9pt,left=3mm]{\Large\bfseries\color{ccal}a)}\quad\begin{minipage}[c]{5.6cm}\raggedright\Large $q=13\text{ et }r=2$\end{minipage}\end{tcolorbox}\hfill\begin{tcolorbox}[enhanced,nobeforeafter,width=7.1cm,height=2.1cm,valign=center,colback=ccal!6,colframe=ccal,arc=2.5mm,boxrule=0.9pt,left=3mm]{\Large\bfseries\color{ccal}b)}\quad\begin{minipage}[c]{5.6cm}\raggedright\Large $q=13\text{ et }r=1$\end{minipage}\end{tcolorbox}\par\vspace{2.5mm}
\begin{tcolorbox}[enhanced,nobeforeafter,width=7.1cm,height=2.1cm,valign=center,colback=ccal!6,colframe=ccal,arc=2.5mm,boxrule=0.9pt,left=3mm]{\Large\bfseries\color{ccal}c)}\quad\begin{minipage}[c]{5.6cm}\raggedright\Large $q=12\text{ et }r=13$\end{minipage}\end{tcolorbox}\hfill\begin{tcolorbox}[enhanced,nobeforeafter,width=7.1cm,height=2.1cm,valign=center,colback=ccal!6,colframe=ccal,arc=2.5mm,boxrule=0.9pt,left=3mm]{\Large\bfseries\color{ccal}d)}\quad\begin{minipage}[c]{5.6cm}\raggedright\Large $q=14\text{ et }r=1$\end{minipage}\end{tcolorbox}
\end{tcolorbox}
\newpage
\begin{tcolorbox}[enhanced,width=16cm,height=14cm,colback=white,colframe=ccal,boxrule=2pt,arc=0mm,
left=6mm,right=6mm,top=0mm,bottom=4mm,sharp corners,
overlay={\fill[ccal] (frame.north west) rectangle ([yshift=-1.25cm]frame.north east);
\node[anchor=west,text=white,font=\sffamily\bfseries\LARGE] at ([xshift=5mm,yshift=-0.62cm]frame.north west) {Question 24};
\node[anchor=east,text=white,font=\sffamily\itshape\large] at ([xshift=-5mm,yshift=-0.62cm]frame.north east) {Calculs};}]
\vspace*{1.75cm}
\begin{minipage}[t][6.9cm][c]{\linewidth}\LARGE\centering
\begin{minipage}{0.95\linewidth}\raggedright\linespread{1.25}\selectfont Dans la division euclidienne de $100$ par $7$, quels sont le quotient $q$ et le reste $r$~?\end{minipage}
\end{minipage}

\vfill
\begin{tcolorbox}[enhanced,nobeforeafter,width=7.1cm,height=2.1cm,valign=center,colback=ccal!6,colframe=ccal,arc=2.5mm,boxrule=0.9pt,left=3mm]{\Large\bfseries\color{ccal}a)}\quad\begin{minipage}[c]{5.6cm}\raggedright\Large $q=13\text{ et }r=9$\end{minipage}\end{tcolorbox}\hfill\begin{tcolorbox}[enhanced,nobeforeafter,width=7.1cm,height=2.1cm,valign=center,colback=ccal!6,colframe=ccal,arc=2.5mm,boxrule=0.9pt,left=3mm]{\Large\bfseries\color{ccal}b)}\quad\begin{minipage}[c]{5.6cm}\raggedright\Large $q=14\text{ et }r=2$\end{minipage}\end{tcolorbox}\par\vspace{2.5mm}
\begin{tcolorbox}[enhanced,nobeforeafter,width=7.1cm,height=2.1cm,valign=center,colback=ccal!6,colframe=ccal,arc=2.5mm,boxrule=0.9pt,left=3mm]{\Large\bfseries\color{ccal}c)}\quad\begin{minipage}[c]{5.6cm}\raggedright\Large $q=14\text{ et }r=3$\end{minipage}\end{tcolorbox}\hfill\begin{tcolorbox}[enhanced,nobeforeafter,width=7.1cm,height=2.1cm,valign=center,colback=ccal!6,colframe=ccal,arc=2.5mm,boxrule=0.9pt,left=3mm]{\Large\bfseries\color{ccal}d)}\quad\begin{minipage}[c]{5.6cm}\raggedright\Large $q=15\text{ et }r=2$\end{minipage}\end{tcolorbox}
\end{tcolorbox}
\newpage
\begin{tcolorbox}[enhanced,width=16cm,height=14cm,colback=white,colframe=ccal,boxrule=2pt,arc=0mm,
left=6mm,right=6mm,top=0mm,bottom=4mm,sharp corners,
overlay={\fill[ccal] (frame.north west) rectangle ([yshift=-1.25cm]frame.north east);
\node[anchor=west,text=white,font=\sffamily\bfseries\LARGE] at ([xshift=5mm,yshift=-0.62cm]frame.north west) {Question 25};
\node[anchor=east,text=white,font=\sffamily\itshape\large] at ([xshift=-5mm,yshift=-0.62cm]frame.north east) {Calculs};}]
\vspace*{1.75cm}
\begin{minipage}[t][6.9cm][c]{\linewidth}\LARGE\centering
\begin{minipage}{0.95\linewidth}\raggedright\linespread{1.25}\selectfont Lequel de ces nombres est un multiple de $9$~?\end{minipage}
\end{minipage}

\vfill
\begin{tcolorbox}[enhanced,nobeforeafter,width=7.1cm,height=2.1cm,valign=center,colback=ccal!6,colframe=ccal,arc=2.5mm,boxrule=0.9pt,left=3mm]{\Large\bfseries\color{ccal}a)}\quad\begin{minipage}[c]{5.6cm}\raggedright\Large $109$\end{minipage}\end{tcolorbox}\hfill\begin{tcolorbox}[enhanced,nobeforeafter,width=7.1cm,height=2.1cm,valign=center,colback=ccal!6,colframe=ccal,arc=2.5mm,boxrule=0.9pt,left=3mm]{\Large\bfseries\color{ccal}b)}\quad\begin{minipage}[c]{5.6cm}\raggedright\Large $119$\end{minipage}\end{tcolorbox}\par\vspace{2.5mm}
\begin{tcolorbox}[enhanced,nobeforeafter,width=7.1cm,height=2.1cm,valign=center,colback=ccal!6,colframe=ccal,arc=2.5mm,boxrule=0.9pt,left=3mm]{\Large\bfseries\color{ccal}c)}\quad\begin{minipage}[c]{5.6cm}\raggedright\Large $129$\end{minipage}\end{tcolorbox}\hfill\begin{tcolorbox}[enhanced,nobeforeafter,width=7.1cm,height=2.1cm,valign=center,colback=ccal!6,colframe=ccal,arc=2.5mm,boxrule=0.9pt,left=3mm]{\Large\bfseries\color{ccal}d)}\quad\begin{minipage}[c]{5.6cm}\raggedright\Large $117$\end{minipage}\end{tcolorbox}
\end{tcolorbox}
\newpage
\begin{tcolorbox}[enhanced,width=16cm,height=14cm,colback=white,colframe=ccal,boxrule=2pt,arc=0mm,
left=6mm,right=6mm,top=0mm,bottom=4mm,sharp corners,
overlay={\fill[ccal] (frame.north west) rectangle ([yshift=-1.25cm]frame.north east);
\node[anchor=west,text=white,font=\sffamily\bfseries\LARGE] at ([xshift=5mm,yshift=-0.62cm]frame.north west) {Question 26};
\node[anchor=east,text=white,font=\sffamily\itshape\large] at ([xshift=-5mm,yshift=-0.62cm]frame.north east) {Calculs};}]
\vspace*{1.75cm}
\begin{minipage}[t][6.9cm][c]{\linewidth}\LARGE\centering
\begin{minipage}{0.95\linewidth}\raggedright\linespread{1.25}\selectfont Lequel de ces nombres est un multiple de $7$~?\end{minipage}
\end{minipage}

\vfill
\begin{tcolorbox}[enhanced,nobeforeafter,width=7.1cm,height=2.1cm,valign=center,colback=ccal!6,colframe=ccal,arc=2.5mm,boxrule=0.9pt,left=3mm]{\Large\bfseries\color{ccal}a)}\quad\begin{minipage}[c]{5.6cm}\raggedright\Large $97$\end{minipage}\end{tcolorbox}\hfill\begin{tcolorbox}[enhanced,nobeforeafter,width=7.1cm,height=2.1cm,valign=center,colback=ccal!6,colframe=ccal,arc=2.5mm,boxrule=0.9pt,left=3mm]{\Large\bfseries\color{ccal}b)}\quad\begin{minipage}[c]{5.6cm}\raggedright\Large $91$\end{minipage}\end{tcolorbox}\par\vspace{2.5mm}
\begin{tcolorbox}[enhanced,nobeforeafter,width=7.1cm,height=2.1cm,valign=center,colback=ccal!6,colframe=ccal,arc=2.5mm,boxrule=0.9pt,left=3mm]{\Large\bfseries\color{ccal}c)}\quad\begin{minipage}[c]{5.6cm}\raggedright\Large $87$\end{minipage}\end{tcolorbox}\hfill\begin{tcolorbox}[enhanced,nobeforeafter,width=7.1cm,height=2.1cm,valign=center,colback=ccal!6,colframe=ccal,arc=2.5mm,boxrule=0.9pt,left=3mm]{\Large\bfseries\color{ccal}d)}\quad\begin{minipage}[c]{5.6cm}\raggedright\Large $101$\end{minipage}\end{tcolorbox}
\end{tcolorbox}
\newpage
\begin{tcolorbox}[enhanced,width=16cm,height=14cm,colback=white,colframe=ccal,boxrule=2pt,arc=0mm,
left=6mm,right=6mm,top=0mm,bottom=4mm,sharp corners,
overlay={\fill[ccal] (frame.north west) rectangle ([yshift=-1.25cm]frame.north east);
\node[anchor=west,text=white,font=\sffamily\bfseries\LARGE] at ([xshift=5mm,yshift=-0.62cm]frame.north west) {Question 27};
\node[anchor=east,text=white,font=\sffamily\itshape\large] at ([xshift=-5mm,yshift=-0.62cm]frame.north east) {Calculs};}]
\vspace*{1.75cm}
\begin{minipage}[t][6.9cm][c]{\linewidth}\LARGE\centering
\begin{minipage}{0.95\linewidth}\raggedright\linespread{1.25}\selectfont Lequel de ces nombres est divisible par $3$~?\end{minipage}
\end{minipage}

\vfill
\begin{tcolorbox}[enhanced,nobeforeafter,width=7.1cm,height=2.1cm,valign=center,colback=ccal!6,colframe=ccal,arc=2.5mm,boxrule=0.9pt,left=3mm]{\Large\bfseries\color{ccal}a)}\quad\begin{minipage}[c]{5.6cm}\raggedright\Large $346$\end{minipage}\end{tcolorbox}\hfill\begin{tcolorbox}[enhanced,nobeforeafter,width=7.1cm,height=2.1cm,valign=center,colback=ccal!6,colframe=ccal,arc=2.5mm,boxrule=0.9pt,left=3mm]{\Large\bfseries\color{ccal}b)}\quad\begin{minipage}[c]{5.6cm}\raggedright\Large $342$\end{minipage}\end{tcolorbox}\par\vspace{2.5mm}
\begin{tcolorbox}[enhanced,nobeforeafter,width=7.1cm,height=2.1cm,valign=center,colback=ccal!6,colframe=ccal,arc=2.5mm,boxrule=0.9pt,left=3mm]{\Large\bfseries\color{ccal}c)}\quad\begin{minipage}[c]{5.6cm}\raggedright\Large $349$\end{minipage}\end{tcolorbox}\hfill\begin{tcolorbox}[enhanced,nobeforeafter,width=7.1cm,height=2.1cm,valign=center,colback=ccal!6,colframe=ccal,arc=2.5mm,boxrule=0.9pt,left=3mm]{\Large\bfseries\color{ccal}d)}\quad\begin{minipage}[c]{5.6cm}\raggedright\Large $341$\end{minipage}\end{tcolorbox}
\end{tcolorbox}
\newpage
\begin{tcolorbox}[enhanced,width=16cm,height=14cm,colback=white,colframe=ccal,boxrule=2pt,arc=0mm,
left=6mm,right=6mm,top=0mm,bottom=4mm,sharp corners,
overlay={\fill[ccal] (frame.north west) rectangle ([yshift=-1.25cm]frame.north east);
\node[anchor=west,text=white,font=\sffamily\bfseries\LARGE] at ([xshift=5mm,yshift=-0.62cm]frame.north west) {Question 28};
\node[anchor=east,text=white,font=\sffamily\itshape\large] at ([xshift=-5mm,yshift=-0.62cm]frame.north east) {Calculs};}]
\vspace*{1.75cm}
\begin{minipage}[t][6.9cm][c]{\linewidth}\LARGE\centering
\begin{minipage}{0.95\linewidth}\raggedright\linespread{1.25}\selectfont Lequel de ces nombres est divisible par $9$~?\end{minipage}
\end{minipage}

\vfill
\begin{tcolorbox}[enhanced,nobeforeafter,width=7.1cm,height=2.1cm,valign=center,colback=ccal!6,colframe=ccal,arc=2.5mm,boxrule=0.9pt,left=3mm]{\Large\bfseries\color{ccal}a)}\quad\begin{minipage}[c]{5.6cm}\raggedright\Large $4\,527$\end{minipage}\end{tcolorbox}\hfill\begin{tcolorbox}[enhanced,nobeforeafter,width=7.1cm,height=2.1cm,valign=center,colback=ccal!6,colframe=ccal,arc=2.5mm,boxrule=0.9pt,left=3mm]{\Large\bfseries\color{ccal}b)}\quad\begin{minipage}[c]{5.6cm}\raggedright\Large $4\,528$\end{minipage}\end{tcolorbox}\par\vspace{2.5mm}
\begin{tcolorbox}[enhanced,nobeforeafter,width=7.1cm,height=2.1cm,valign=center,colback=ccal!6,colframe=ccal,arc=2.5mm,boxrule=0.9pt,left=3mm]{\Large\bfseries\color{ccal}c)}\quad\begin{minipage}[c]{5.6cm}\raggedright\Large $4\,537$\end{minipage}\end{tcolorbox}\hfill\begin{tcolorbox}[enhanced,nobeforeafter,width=7.1cm,height=2.1cm,valign=center,colback=ccal!6,colframe=ccal,arc=2.5mm,boxrule=0.9pt,left=3mm]{\Large\bfseries\color{ccal}d)}\quad\begin{minipage}[c]{5.6cm}\raggedright\Large $4\,625$\end{minipage}\end{tcolorbox}
\end{tcolorbox}
\newpage
\begin{tcolorbox}[enhanced,width=16cm,height=14cm,colback=white,colframe=ccal,boxrule=2pt,arc=0mm,
left=6mm,right=6mm,top=0mm,bottom=4mm,sharp corners,
overlay={\fill[ccal] (frame.north west) rectangle ([yshift=-1.25cm]frame.north east);
\node[anchor=west,text=white,font=\sffamily\bfseries\LARGE] at ([xshift=5mm,yshift=-0.62cm]frame.north west) {Question 29};
\node[anchor=east,text=white,font=\sffamily\itshape\large] at ([xshift=-5mm,yshift=-0.62cm]frame.north east) {Calculs};}]
\vspace*{1.75cm}
\begin{minipage}[t][6.9cm][c]{\linewidth}\LARGE\centering
\begin{minipage}{0.95\linewidth}\raggedright\linespread{1.25}\selectfont Que vaut $5+3\times4$~?\end{minipage}
\end{minipage}

\vfill
\begin{tcolorbox}[enhanced,nobeforeafter,width=7.1cm,height=2.1cm,valign=center,colback=ccal!6,colframe=ccal,arc=2.5mm,boxrule=0.9pt,left=3mm]{\Large\bfseries\color{ccal}a)}\quad\begin{minipage}[c]{5.6cm}\raggedright\Large $32$\end{minipage}\end{tcolorbox}\hfill\begin{tcolorbox}[enhanced,nobeforeafter,width=7.1cm,height=2.1cm,valign=center,colback=ccal!6,colframe=ccal,arc=2.5mm,boxrule=0.9pt,left=3mm]{\Large\bfseries\color{ccal}b)}\quad\begin{minipage}[c]{5.6cm}\raggedright\Large $17$\end{minipage}\end{tcolorbox}\par\vspace{2.5mm}
\begin{tcolorbox}[enhanced,nobeforeafter,width=7.1cm,height=2.1cm,valign=center,colback=ccal!6,colframe=ccal,arc=2.5mm,boxrule=0.9pt,left=3mm]{\Large\bfseries\color{ccal}c)}\quad\begin{minipage}[c]{5.6cm}\raggedright\Large $12$\end{minipage}\end{tcolorbox}\hfill\begin{tcolorbox}[enhanced,nobeforeafter,width=7.1cm,height=2.1cm,valign=center,colback=ccal!6,colframe=ccal,arc=2.5mm,boxrule=0.9pt,left=3mm]{\Large\bfseries\color{ccal}d)}\quad\begin{minipage}[c]{5.6cm}\raggedright\Large $20$\end{minipage}\end{tcolorbox}
\end{tcolorbox}
\newpage
\begin{tcolorbox}[enhanced,width=16cm,height=14cm,colback=white,colframe=ccal,boxrule=2pt,arc=0mm,
left=6mm,right=6mm,top=0mm,bottom=4mm,sharp corners,
overlay={\fill[ccal] (frame.north west) rectangle ([yshift=-1.25cm]frame.north east);
\node[anchor=west,text=white,font=\sffamily\bfseries\LARGE] at ([xshift=5mm,yshift=-0.62cm]frame.north west) {Question 30};
\node[anchor=east,text=white,font=\sffamily\itshape\large] at ([xshift=-5mm,yshift=-0.62cm]frame.north east) {Calculs};}]
\vspace*{1.75cm}
\begin{minipage}[t][6.9cm][c]{\linewidth}\LARGE\centering
\begin{minipage}{0.95\linewidth}\raggedright\linespread{1.25}\selectfont Que vaut $20-4\times3$~?\end{minipage}
\end{minipage}

\vfill
\begin{tcolorbox}[enhanced,nobeforeafter,width=7.1cm,height=2.1cm,valign=center,colback=ccal!6,colframe=ccal,arc=2.5mm,boxrule=0.9pt,left=3mm]{\Large\bfseries\color{ccal}a)}\quad\begin{minipage}[c]{5.6cm}\raggedright\Large $48$\end{minipage}\end{tcolorbox}\hfill\begin{tcolorbox}[enhanced,nobeforeafter,width=7.1cm,height=2.1cm,valign=center,colback=ccal!6,colframe=ccal,arc=2.5mm,boxrule=0.9pt,left=3mm]{\Large\bfseries\color{ccal}b)}\quad\begin{minipage}[c]{5.6cm}\raggedright\Large $17$\end{minipage}\end{tcolorbox}\par\vspace{2.5mm}
\begin{tcolorbox}[enhanced,nobeforeafter,width=7.1cm,height=2.1cm,valign=center,colback=ccal!6,colframe=ccal,arc=2.5mm,boxrule=0.9pt,left=3mm]{\Large\bfseries\color{ccal}c)}\quad\begin{minipage}[c]{5.6cm}\raggedright\Large $8$\end{minipage}\end{tcolorbox}\hfill\begin{tcolorbox}[enhanced,nobeforeafter,width=7.1cm,height=2.1cm,valign=center,colback=ccal!6,colframe=ccal,arc=2.5mm,boxrule=0.9pt,left=3mm]{\Large\bfseries\color{ccal}d)}\quad\begin{minipage}[c]{5.6cm}\raggedright\Large $12$\end{minipage}\end{tcolorbox}
\end{tcolorbox}
\newpage
\begin{tcolorbox}[enhanced,width=16cm,height=14cm,colback=white,colframe=ccal,boxrule=2pt,arc=0mm,
left=6mm,right=6mm,top=0mm,bottom=4mm,sharp corners,
overlay={\fill[ccal] (frame.north west) rectangle ([yshift=-1.25cm]frame.north east);
\node[anchor=west,text=white,font=\sffamily\bfseries\LARGE] at ([xshift=5mm,yshift=-0.62cm]frame.north west) {Question 31};
\node[anchor=east,text=white,font=\sffamily\itshape\large] at ([xshift=-5mm,yshift=-0.62cm]frame.north east) {Calculs};}]
\vspace*{1.75cm}
\begin{minipage}[t][6.9cm][c]{\linewidth}\LARGE\centering
\begin{minipage}{0.95\linewidth}\raggedright\linespread{1.25}\selectfont Que vaut $7\div4$~?\end{minipage}
\end{minipage}

\vfill
\begin{tcolorbox}[enhanced,nobeforeafter,width=7.1cm,height=2.1cm,valign=center,colback=ccal!6,colframe=ccal,arc=2.5mm,boxrule=0.9pt,left=3mm]{\Large\bfseries\color{ccal}a)}\quad\begin{minipage}[c]{5.6cm}\raggedright\Large $1{,}75$\end{minipage}\end{tcolorbox}\hfill\begin{tcolorbox}[enhanced,nobeforeafter,width=7.1cm,height=2.1cm,valign=center,colback=ccal!6,colframe=ccal,arc=2.5mm,boxrule=0.9pt,left=3mm]{\Large\bfseries\color{ccal}b)}\quad\begin{minipage}[c]{5.6cm}\raggedright\Large $17{,}5$\end{minipage}\end{tcolorbox}\par\vspace{2.5mm}
\begin{tcolorbox}[enhanced,nobeforeafter,width=7.1cm,height=2.1cm,valign=center,colback=ccal!6,colframe=ccal,arc=2.5mm,boxrule=0.9pt,left=3mm]{\Large\bfseries\color{ccal}c)}\quad\begin{minipage}[c]{5.6cm}\raggedright\Large $\dfrac{4}{7}$\end{minipage}\end{tcolorbox}\hfill\begin{tcolorbox}[enhanced,nobeforeafter,width=7.1cm,height=2.1cm,valign=center,colback=ccal!6,colframe=ccal,arc=2.5mm,boxrule=0.9pt,left=3mm]{\Large\bfseries\color{ccal}d)}\quad\begin{minipage}[c]{5.6cm}\raggedright\Large $1{,}3$\end{minipage}\end{tcolorbox}
\end{tcolorbox}
\newpage
\begin{tcolorbox}[enhanced,width=16cm,height=14cm,colback=white,colframe=ccal,boxrule=2pt,arc=0mm,
left=6mm,right=6mm,top=0mm,bottom=4mm,sharp corners,
overlay={\fill[ccal] (frame.north west) rectangle ([yshift=-1.25cm]frame.north east);
\node[anchor=west,text=white,font=\sffamily\bfseries\LARGE] at ([xshift=5mm,yshift=-0.62cm]frame.north west) {Question 32};
\node[anchor=east,text=white,font=\sffamily\itshape\large] at ([xshift=-5mm,yshift=-0.62cm]frame.north east) {Calculs};}]
\vspace*{1.75cm}
\begin{minipage}[t][6.9cm][c]{\linewidth}\LARGE\centering
\begin{minipage}{0.95\linewidth}\raggedright\linespread{1.25}\selectfont Que vaut $9\div8$~?\end{minipage}
\end{minipage}

\vfill
\begin{tcolorbox}[enhanced,nobeforeafter,width=7.1cm,height=2.1cm,valign=center,colback=ccal!6,colframe=ccal,arc=2.5mm,boxrule=0.9pt,left=3mm]{\Large\bfseries\color{ccal}a)}\quad\begin{minipage}[c]{5.6cm}\raggedright\Large $1{,}1$\end{minipage}\end{tcolorbox}\hfill\begin{tcolorbox}[enhanced,nobeforeafter,width=7.1cm,height=2.1cm,valign=center,colback=ccal!6,colframe=ccal,arc=2.5mm,boxrule=0.9pt,left=3mm]{\Large\bfseries\color{ccal}b)}\quad\begin{minipage}[c]{5.6cm}\raggedright\Large $\dfrac{8}{9}$\end{minipage}\end{tcolorbox}\par\vspace{2.5mm}
\begin{tcolorbox}[enhanced,nobeforeafter,width=7.1cm,height=2.1cm,valign=center,colback=ccal!6,colframe=ccal,arc=2.5mm,boxrule=0.9pt,left=3mm]{\Large\bfseries\color{ccal}c)}\quad\begin{minipage}[c]{5.6cm}\raggedright\Large $11{,}25$\end{minipage}\end{tcolorbox}\hfill\begin{tcolorbox}[enhanced,nobeforeafter,width=7.1cm,height=2.1cm,valign=center,colback=ccal!6,colframe=ccal,arc=2.5mm,boxrule=0.9pt,left=3mm]{\Large\bfseries\color{ccal}d)}\quad\begin{minipage}[c]{5.6cm}\raggedright\Large $1{,}125$\end{minipage}\end{tcolorbox}
\end{tcolorbox}
\newpage
\begin{tcolorbox}[enhanced,width=16cm,height=14cm,colback=white,colframe=cfra,boxrule=2pt,arc=0mm,
left=6mm,right=6mm,top=0mm,bottom=4mm,sharp corners,
overlay={\fill[cfra] (frame.north west) rectangle ([yshift=-1.25cm]frame.north east);
\node[anchor=west,text=white,font=\sffamily\bfseries\LARGE] at ([xshift=5mm,yshift=-0.62cm]frame.north west) {Question 33};
\node[anchor=east,text=white,font=\sffamily\itshape\large] at ([xshift=-5mm,yshift=-0.62cm]frame.north east) {Fractions};}]
\vspace*{1.75cm}
\begin{minipage}[t][6.9cm][c]{\linewidth}\LARGE\centering
\begin{minipage}{0.95\linewidth}\raggedright\linespread{1.25}\selectfont Combien vaut $\dfrac{3}{4}$ de $20$~?\end{minipage}
\end{minipage}

\vfill
\begin{tcolorbox}[enhanced,nobeforeafter,width=7.1cm,height=2.1cm,valign=center,colback=cfra!6,colframe=cfra,arc=2.5mm,boxrule=0.9pt,left=3mm]{\Large\bfseries\color{cfra}a)}\quad\begin{minipage}[c]{5.6cm}\raggedright\Large $15$\end{minipage}\end{tcolorbox}\hfill\begin{tcolorbox}[enhanced,nobeforeafter,width=7.1cm,height=2.1cm,valign=center,colback=cfra!6,colframe=cfra,arc=2.5mm,boxrule=0.9pt,left=3mm]{\Large\bfseries\color{cfra}b)}\quad\begin{minipage}[c]{5.6cm}\raggedright\Large $60$\end{minipage}\end{tcolorbox}\par\vspace{2.5mm}
\begin{tcolorbox}[enhanced,nobeforeafter,width=7.1cm,height=2.1cm,valign=center,colback=cfra!6,colframe=cfra,arc=2.5mm,boxrule=0.9pt,left=3mm]{\Large\bfseries\color{cfra}c)}\quad\begin{minipage}[c]{5.6cm}\raggedright\Large $5$\end{minipage}\end{tcolorbox}\hfill\begin{tcolorbox}[enhanced,nobeforeafter,width=7.1cm,height=2.1cm,valign=center,colback=cfra!6,colframe=cfra,arc=2.5mm,boxrule=0.9pt,left=3mm]{\Large\bfseries\color{cfra}d)}\quad\begin{minipage}[c]{5.6cm}\raggedright\Large $16$\end{minipage}\end{tcolorbox}
\end{tcolorbox}
\newpage
\begin{tcolorbox}[enhanced,width=16cm,height=14cm,colback=white,colframe=cfra,boxrule=2pt,arc=0mm,
left=6mm,right=6mm,top=0mm,bottom=4mm,sharp corners,
overlay={\fill[cfra] (frame.north west) rectangle ([yshift=-1.25cm]frame.north east);
\node[anchor=west,text=white,font=\sffamily\bfseries\LARGE] at ([xshift=5mm,yshift=-0.62cm]frame.north west) {Question 34};
\node[anchor=east,text=white,font=\sffamily\itshape\large] at ([xshift=-5mm,yshift=-0.62cm]frame.north east) {Fractions};}]
\vspace*{1.75cm}
\begin{minipage}[t][6.9cm][c]{\linewidth}\LARGE\centering
\begin{minipage}{0.95\linewidth}\raggedright\linespread{1.25}\selectfont Combien vaut $\dfrac{2}{5}$ de $35$~?\end{minipage}
\end{minipage}

\vfill
\begin{tcolorbox}[enhanced,nobeforeafter,width=7.1cm,height=2.1cm,valign=center,colback=cfra!6,colframe=cfra,arc=2.5mm,boxrule=0.9pt,left=3mm]{\Large\bfseries\color{cfra}a)}\quad\begin{minipage}[c]{5.6cm}\raggedright\Large $21$\end{minipage}\end{tcolorbox}\hfill\begin{tcolorbox}[enhanced,nobeforeafter,width=7.1cm,height=2.1cm,valign=center,colback=cfra!6,colframe=cfra,arc=2.5mm,boxrule=0.9pt,left=3mm]{\Large\bfseries\color{cfra}b)}\quad\begin{minipage}[c]{5.6cm}\raggedright\Large $14$\end{minipage}\end{tcolorbox}\par\vspace{2.5mm}
\begin{tcolorbox}[enhanced,nobeforeafter,width=7.1cm,height=2.1cm,valign=center,colback=cfra!6,colframe=cfra,arc=2.5mm,boxrule=0.9pt,left=3mm]{\Large\bfseries\color{cfra}c)}\quad\begin{minipage}[c]{5.6cm}\raggedright\Large $70$\end{minipage}\end{tcolorbox}\hfill\begin{tcolorbox}[enhanced,nobeforeafter,width=7.1cm,height=2.1cm,valign=center,colback=cfra!6,colframe=cfra,arc=2.5mm,boxrule=0.9pt,left=3mm]{\Large\bfseries\color{cfra}d)}\quad\begin{minipage}[c]{5.6cm}\raggedright\Large $7$\end{minipage}\end{tcolorbox}
\end{tcolorbox}
\newpage
\begin{tcolorbox}[enhanced,width=16cm,height=14cm,colback=white,colframe=cfra,boxrule=2pt,arc=0mm,
left=6mm,right=6mm,top=0mm,bottom=4mm,sharp corners,
overlay={\fill[cfra] (frame.north west) rectangle ([yshift=-1.25cm]frame.north east);
\node[anchor=west,text=white,font=\sffamily\bfseries\LARGE] at ([xshift=5mm,yshift=-0.62cm]frame.north west) {Question 35};
\node[anchor=east,text=white,font=\sffamily\itshape\large] at ([xshift=-5mm,yshift=-0.62cm]frame.north east) {Fractions};}]
\vspace*{1.75cm}
\begin{minipage}[t][6.9cm][c]{\linewidth}\LARGE\centering
\begin{minipage}{0.95\linewidth}\raggedright\linespread{1.25}\selectfont Compléter~: $\dfrac{2}{3}=\dfrac{?}{12}$\end{minipage}
\end{minipage}

\vfill
\begin{tcolorbox}[enhanced,nobeforeafter,width=7.1cm,height=2.1cm,valign=center,colback=cfra!6,colframe=cfra,arc=2.5mm,boxrule=0.9pt,left=3mm]{\Large\bfseries\color{cfra}a)}\quad\begin{minipage}[c]{5.6cm}\raggedright\Large $24$\end{minipage}\end{tcolorbox}\hfill\begin{tcolorbox}[enhanced,nobeforeafter,width=7.1cm,height=2.1cm,valign=center,colback=cfra!6,colframe=cfra,arc=2.5mm,boxrule=0.9pt,left=3mm]{\Large\bfseries\color{cfra}b)}\quad\begin{minipage}[c]{5.6cm}\raggedright\Large $4$\end{minipage}\end{tcolorbox}\par\vspace{2.5mm}
\begin{tcolorbox}[enhanced,nobeforeafter,width=7.1cm,height=2.1cm,valign=center,colback=cfra!6,colframe=cfra,arc=2.5mm,boxrule=0.9pt,left=3mm]{\Large\bfseries\color{cfra}c)}\quad\begin{minipage}[c]{5.6cm}\raggedright\Large $11$\end{minipage}\end{tcolorbox}\hfill\begin{tcolorbox}[enhanced,nobeforeafter,width=7.1cm,height=2.1cm,valign=center,colback=cfra!6,colframe=cfra,arc=2.5mm,boxrule=0.9pt,left=3mm]{\Large\bfseries\color{cfra}d)}\quad\begin{minipage}[c]{5.6cm}\raggedright\Large $8$\end{minipage}\end{tcolorbox}
\end{tcolorbox}
\newpage
\begin{tcolorbox}[enhanced,width=16cm,height=14cm,colback=white,colframe=cfra,boxrule=2pt,arc=0mm,
left=6mm,right=6mm,top=0mm,bottom=4mm,sharp corners,
overlay={\fill[cfra] (frame.north west) rectangle ([yshift=-1.25cm]frame.north east);
\node[anchor=west,text=white,font=\sffamily\bfseries\LARGE] at ([xshift=5mm,yshift=-0.62cm]frame.north west) {Question 36};
\node[anchor=east,text=white,font=\sffamily\itshape\large] at ([xshift=-5mm,yshift=-0.62cm]frame.north east) {Fractions};}]
\vspace*{1.75cm}
\begin{minipage}[t][6.9cm][c]{\linewidth}\LARGE\centering
\begin{minipage}{0.95\linewidth}\raggedright\linespread{1.25}\selectfont Compléter~: $\dfrac{3}{4}=\dfrac{?}{20}$\end{minipage}
\end{minipage}

\vfill
\begin{tcolorbox}[enhanced,nobeforeafter,width=7.1cm,height=2.1cm,valign=center,colback=cfra!6,colframe=cfra,arc=2.5mm,boxrule=0.9pt,left=3mm]{\Large\bfseries\color{cfra}a)}\quad\begin{minipage}[c]{5.6cm}\raggedright\Large $5$\end{minipage}\end{tcolorbox}\hfill\begin{tcolorbox}[enhanced,nobeforeafter,width=7.1cm,height=2.1cm,valign=center,colback=cfra!6,colframe=cfra,arc=2.5mm,boxrule=0.9pt,left=3mm]{\Large\bfseries\color{cfra}b)}\quad\begin{minipage}[c]{5.6cm}\raggedright\Large $19$\end{minipage}\end{tcolorbox}\par\vspace{2.5mm}
\begin{tcolorbox}[enhanced,nobeforeafter,width=7.1cm,height=2.1cm,valign=center,colback=cfra!6,colframe=cfra,arc=2.5mm,boxrule=0.9pt,left=3mm]{\Large\bfseries\color{cfra}c)}\quad\begin{minipage}[c]{5.6cm}\raggedright\Large $60$\end{minipage}\end{tcolorbox}\hfill\begin{tcolorbox}[enhanced,nobeforeafter,width=7.1cm,height=2.1cm,valign=center,colback=cfra!6,colframe=cfra,arc=2.5mm,boxrule=0.9pt,left=3mm]{\Large\bfseries\color{cfra}d)}\quad\begin{minipage}[c]{5.6cm}\raggedright\Large $15$\end{minipage}\end{tcolorbox}
\end{tcolorbox}
\newpage
\begin{tcolorbox}[enhanced,width=16cm,height=14cm,colback=white,colframe=cfra,boxrule=2pt,arc=0mm,
left=6mm,right=6mm,top=0mm,bottom=4mm,sharp corners,
overlay={\fill[cfra] (frame.north west) rectangle ([yshift=-1.25cm]frame.north east);
\node[anchor=west,text=white,font=\sffamily\bfseries\LARGE] at ([xshift=5mm,yshift=-0.62cm]frame.north west) {Question 37};
\node[anchor=east,text=white,font=\sffamily\itshape\large] at ([xshift=-5mm,yshift=-0.62cm]frame.north east) {Fractions};}]
\vspace*{1.75cm}
\begin{minipage}[t][6.9cm][c]{\linewidth}\LARGE\centering
\begin{minipage}{0.95\linewidth}\raggedright\linespread{1.25}\selectfont Quelle est l'écriture décimale de $\dfrac{3}{4}$~?\end{minipage}
\end{minipage}

\vfill
\begin{tcolorbox}[enhanced,nobeforeafter,width=7.1cm,height=2.1cm,valign=center,colback=cfra!6,colframe=cfra,arc=2.5mm,boxrule=0.9pt,left=3mm]{\Large\bfseries\color{cfra}a)}\quad\begin{minipage}[c]{5.6cm}\raggedright\Large $\dfrac{4}{3}$\end{minipage}\end{tcolorbox}\hfill\begin{tcolorbox}[enhanced,nobeforeafter,width=7.1cm,height=2.1cm,valign=center,colback=cfra!6,colframe=cfra,arc=2.5mm,boxrule=0.9pt,left=3mm]{\Large\bfseries\color{cfra}b)}\quad\begin{minipage}[c]{5.6cm}\raggedright\Large $7{,}5$\end{minipage}\end{tcolorbox}\par\vspace{2.5mm}
\begin{tcolorbox}[enhanced,nobeforeafter,width=7.1cm,height=2.1cm,valign=center,colback=cfra!6,colframe=cfra,arc=2.5mm,boxrule=0.9pt,left=3mm]{\Large\bfseries\color{cfra}c)}\quad\begin{minipage}[c]{5.6cm}\raggedright\Large $3{,}4$\end{minipage}\end{tcolorbox}\hfill\begin{tcolorbox}[enhanced,nobeforeafter,width=7.1cm,height=2.1cm,valign=center,colback=cfra!6,colframe=cfra,arc=2.5mm,boxrule=0.9pt,left=3mm]{\Large\bfseries\color{cfra}d)}\quad\begin{minipage}[c]{5.6cm}\raggedright\Large $0{,}75$\end{minipage}\end{tcolorbox}
\end{tcolorbox}
\newpage
\begin{tcolorbox}[enhanced,width=16cm,height=14cm,colback=white,colframe=cfra,boxrule=2pt,arc=0mm,
left=6mm,right=6mm,top=0mm,bottom=4mm,sharp corners,
overlay={\fill[cfra] (frame.north west) rectangle ([yshift=-1.25cm]frame.north east);
\node[anchor=west,text=white,font=\sffamily\bfseries\LARGE] at ([xshift=5mm,yshift=-0.62cm]frame.north west) {Question 38};
\node[anchor=east,text=white,font=\sffamily\itshape\large] at ([xshift=-5mm,yshift=-0.62cm]frame.north east) {Fractions};}]
\vspace*{1.75cm}
\begin{minipage}[t][6.9cm][c]{\linewidth}\LARGE\centering
\begin{minipage}{0.95\linewidth}\raggedright\linespread{1.25}\selectfont Quelle est l'écriture décimale de $\dfrac{1}{2}$~?\end{minipage}
\end{minipage}

\vfill
\begin{tcolorbox}[enhanced,nobeforeafter,width=7.1cm,height=2.1cm,valign=center,colback=cfra!6,colframe=cfra,arc=2.5mm,boxrule=0.9pt,left=3mm]{\Large\bfseries\color{cfra}a)}\quad\begin{minipage}[c]{5.6cm}\raggedright\Large $1{,}2$\end{minipage}\end{tcolorbox}\hfill\begin{tcolorbox}[enhanced,nobeforeafter,width=7.1cm,height=2.1cm,valign=center,colback=cfra!6,colframe=cfra,arc=2.5mm,boxrule=0.9pt,left=3mm]{\Large\bfseries\color{cfra}b)}\quad\begin{minipage}[c]{5.6cm}\raggedright\Large $0{,}5$\end{minipage}\end{tcolorbox}\par\vspace{2.5mm}
\begin{tcolorbox}[enhanced,nobeforeafter,width=7.1cm,height=2.1cm,valign=center,colback=cfra!6,colframe=cfra,arc=2.5mm,boxrule=0.9pt,left=3mm]{\Large\bfseries\color{cfra}c)}\quad\begin{minipage}[c]{5.6cm}\raggedright\Large $5$\end{minipage}\end{tcolorbox}\hfill\begin{tcolorbox}[enhanced,nobeforeafter,width=7.1cm,height=2.1cm,valign=center,colback=cfra!6,colframe=cfra,arc=2.5mm,boxrule=0.9pt,left=3mm]{\Large\bfseries\color{cfra}d)}\quad\begin{minipage}[c]{5.6cm}\raggedright\Large $2$\end{minipage}\end{tcolorbox}
\end{tcolorbox}
\newpage
\begin{tcolorbox}[enhanced,width=16cm,height=14cm,colback=white,colframe=cfra,boxrule=2pt,arc=0mm,
left=6mm,right=6mm,top=0mm,bottom=4mm,sharp corners,
overlay={\fill[cfra] (frame.north west) rectangle ([yshift=-1.25cm]frame.north east);
\node[anchor=west,text=white,font=\sffamily\bfseries\LARGE] at ([xshift=5mm,yshift=-0.62cm]frame.north west) {Question 39};
\node[anchor=east,text=white,font=\sffamily\itshape\large] at ([xshift=-5mm,yshift=-0.62cm]frame.north east) {Fractions};}]
\vspace*{1.75cm}
\begin{minipage}[t][6.9cm][c]{\linewidth}\LARGE\centering
\begin{minipage}{0.95\linewidth}\raggedright\linespread{1.25}\selectfont Laquelle de ces fractions est plus grande que $1$~?\end{minipage}
\end{minipage}

\vfill
\begin{tcolorbox}[enhanced,nobeforeafter,width=7.1cm,height=2.1cm,valign=center,colback=cfra!6,colframe=cfra,arc=2.5mm,boxrule=0.9pt,left=3mm]{\Large\bfseries\color{cfra}a)}\quad\begin{minipage}[c]{5.6cm}\raggedright\Large $\dfrac{4}{4}$\end{minipage}\end{tcolorbox}\hfill\begin{tcolorbox}[enhanced,nobeforeafter,width=7.1cm,height=2.1cm,valign=center,colback=cfra!6,colframe=cfra,arc=2.5mm,boxrule=0.9pt,left=3mm]{\Large\bfseries\color{cfra}b)}\quad\begin{minipage}[c]{5.6cm}\raggedright\Large $\dfrac{5}{7}$\end{minipage}\end{tcolorbox}\par\vspace{2.5mm}
\begin{tcolorbox}[enhanced,nobeforeafter,width=7.1cm,height=2.1cm,valign=center,colback=cfra!6,colframe=cfra,arc=2.5mm,boxrule=0.9pt,left=3mm]{\Large\bfseries\color{cfra}c)}\quad\begin{minipage}[c]{5.6cm}\raggedright\Large $\dfrac{3}{4}$\end{minipage}\end{tcolorbox}\hfill\begin{tcolorbox}[enhanced,nobeforeafter,width=7.1cm,height=2.1cm,valign=center,colback=cfra!6,colframe=cfra,arc=2.5mm,boxrule=0.9pt,left=3mm]{\Large\bfseries\color{cfra}d)}\quad\begin{minipage}[c]{5.6cm}\raggedright\Large $\dfrac{7}{5}$\end{minipage}\end{tcolorbox}
\end{tcolorbox}
\newpage
\begin{tcolorbox}[enhanced,width=16cm,height=14cm,colback=white,colframe=cfra,boxrule=2pt,arc=0mm,
left=6mm,right=6mm,top=0mm,bottom=4mm,sharp corners,
overlay={\fill[cfra] (frame.north west) rectangle ([yshift=-1.25cm]frame.north east);
\node[anchor=west,text=white,font=\sffamily\bfseries\LARGE] at ([xshift=5mm,yshift=-0.62cm]frame.north west) {Question 40};
\node[anchor=east,text=white,font=\sffamily\itshape\large] at ([xshift=-5mm,yshift=-0.62cm]frame.north east) {Fractions};}]
\vspace*{1.75cm}
\begin{minipage}[t][6.9cm][c]{\linewidth}\LARGE\centering
\begin{minipage}{0.95\linewidth}\raggedright\linespread{1.25}\selectfont Laquelle de ces fractions est plus grande que $1$~?\end{minipage}
\end{minipage}

\vfill
\begin{tcolorbox}[enhanced,nobeforeafter,width=7.1cm,height=2.1cm,valign=center,colback=cfra!6,colframe=cfra,arc=2.5mm,boxrule=0.9pt,left=3mm]{\Large\bfseries\color{cfra}a)}\quad\begin{minipage}[c]{5.6cm}\raggedright\Large $\dfrac{9}{8}$\end{minipage}\end{tcolorbox}\hfill\begin{tcolorbox}[enhanced,nobeforeafter,width=7.1cm,height=2.1cm,valign=center,colback=cfra!6,colframe=cfra,arc=2.5mm,boxrule=0.9pt,left=3mm]{\Large\bfseries\color{cfra}b)}\quad\begin{minipage}[c]{5.6cm}\raggedright\Large $\dfrac{8}{9}$\end{minipage}\end{tcolorbox}\par\vspace{2.5mm}
\begin{tcolorbox}[enhanced,nobeforeafter,width=7.1cm,height=2.1cm,valign=center,colback=cfra!6,colframe=cfra,arc=2.5mm,boxrule=0.9pt,left=3mm]{\Large\bfseries\color{cfra}c)}\quad\begin{minipage}[c]{5.6cm}\raggedright\Large $\dfrac{1}{2}$\end{minipage}\end{tcolorbox}\hfill\begin{tcolorbox}[enhanced,nobeforeafter,width=7.1cm,height=2.1cm,valign=center,colback=cfra!6,colframe=cfra,arc=2.5mm,boxrule=0.9pt,left=3mm]{\Large\bfseries\color{cfra}d)}\quad\begin{minipage}[c]{5.6cm}\raggedright\Large $\dfrac{6}{6}$\end{minipage}\end{tcolorbox}
\end{tcolorbox}
\newpage
\begin{tcolorbox}[enhanced,width=16cm,height=14cm,colback=white,colframe=cfra,boxrule=2pt,arc=0mm,
left=6mm,right=6mm,top=0mm,bottom=4mm,sharp corners,
overlay={\fill[cfra] (frame.north west) rectangle ([yshift=-1.25cm]frame.north east);
\node[anchor=west,text=white,font=\sffamily\bfseries\LARGE] at ([xshift=5mm,yshift=-0.62cm]frame.north west) {Question 41};
\node[anchor=east,text=white,font=\sffamily\itshape\large] at ([xshift=-5mm,yshift=-0.62cm]frame.north east) {Fractions};}]
\vspace*{1.75cm}
\begin{minipage}[t][6.9cm][c]{\linewidth}\LARGE\centering
\begin{minipage}{0.95\linewidth}\raggedright\linespread{1.25}\selectfont Quel nombre, multiplié par $7$, donne $3$~?\end{minipage}
\end{minipage}

\vfill
\begin{tcolorbox}[enhanced,nobeforeafter,width=7.1cm,height=2.1cm,valign=center,colback=cfra!6,colframe=cfra,arc=2.5mm,boxrule=0.9pt,left=3mm]{\Large\bfseries\color{cfra}a)}\quad\begin{minipage}[c]{5.6cm}\raggedright\Large $4$\end{minipage}\end{tcolorbox}\hfill\begin{tcolorbox}[enhanced,nobeforeafter,width=7.1cm,height=2.1cm,valign=center,colback=cfra!6,colframe=cfra,arc=2.5mm,boxrule=0.9pt,left=3mm]{\Large\bfseries\color{cfra}b)}\quad\begin{minipage}[c]{5.6cm}\raggedright\Large $\dfrac{3}{7}$\end{minipage}\end{tcolorbox}\par\vspace{2.5mm}
\begin{tcolorbox}[enhanced,nobeforeafter,width=7.1cm,height=2.1cm,valign=center,colback=cfra!6,colframe=cfra,arc=2.5mm,boxrule=0.9pt,left=3mm]{\Large\bfseries\color{cfra}c)}\quad\begin{minipage}[c]{5.6cm}\raggedright\Large $\dfrac{7}{3}$\end{minipage}\end{tcolorbox}\hfill\begin{tcolorbox}[enhanced,nobeforeafter,width=7.1cm,height=2.1cm,valign=center,colback=cfra!6,colframe=cfra,arc=2.5mm,boxrule=0.9pt,left=3mm]{\Large\bfseries\color{cfra}d)}\quad\begin{minipage}[c]{5.6cm}\raggedright\Large $21$\end{minipage}\end{tcolorbox}
\end{tcolorbox}
\newpage
\begin{tcolorbox}[enhanced,width=16cm,height=14cm,colback=white,colframe=cfra,boxrule=2pt,arc=0mm,
left=6mm,right=6mm,top=0mm,bottom=4mm,sharp corners,
overlay={\fill[cfra] (frame.north west) rectangle ([yshift=-1.25cm]frame.north east);
\node[anchor=west,text=white,font=\sffamily\bfseries\LARGE] at ([xshift=5mm,yshift=-0.62cm]frame.north west) {Question 42};
\node[anchor=east,text=white,font=\sffamily\itshape\large] at ([xshift=-5mm,yshift=-0.62cm]frame.north east) {Fractions};}]
\vspace*{1.75cm}
\begin{minipage}[t][6.9cm][c]{\linewidth}\LARGE\centering
\begin{minipage}{0.95\linewidth}\raggedright\linespread{1.25}\selectfont Quel nombre, multiplié par $4$, donne $5$~?\end{minipage}
\end{minipage}

\vfill
\begin{tcolorbox}[enhanced,nobeforeafter,width=7.1cm,height=2.1cm,valign=center,colback=cfra!6,colframe=cfra,arc=2.5mm,boxrule=0.9pt,left=3mm]{\Large\bfseries\color{cfra}a)}\quad\begin{minipage}[c]{5.6cm}\raggedright\Large $\dfrac{4}{5}$\end{minipage}\end{tcolorbox}\hfill\begin{tcolorbox}[enhanced,nobeforeafter,width=7.1cm,height=2.1cm,valign=center,colback=cfra!6,colframe=cfra,arc=2.5mm,boxrule=0.9pt,left=3mm]{\Large\bfseries\color{cfra}b)}\quad\begin{minipage}[c]{5.6cm}\raggedright\Large $\dfrac{5}{4}$\end{minipage}\end{tcolorbox}\par\vspace{2.5mm}
\begin{tcolorbox}[enhanced,nobeforeafter,width=7.1cm,height=2.1cm,valign=center,colback=cfra!6,colframe=cfra,arc=2.5mm,boxrule=0.9pt,left=3mm]{\Large\bfseries\color{cfra}c)}\quad\begin{minipage}[c]{5.6cm}\raggedright\Large $1$\end{minipage}\end{tcolorbox}\hfill\begin{tcolorbox}[enhanced,nobeforeafter,width=7.1cm,height=2.1cm,valign=center,colback=cfra!6,colframe=cfra,arc=2.5mm,boxrule=0.9pt,left=3mm]{\Large\bfseries\color{cfra}d)}\quad\begin{minipage}[c]{5.6cm}\raggedright\Large $20$\end{minipage}\end{tcolorbox}
\end{tcolorbox}
\newpage
\begin{tcolorbox}[enhanced,width=16cm,height=14cm,colback=white,colframe=cfra,boxrule=2pt,arc=0mm,
left=6mm,right=6mm,top=0mm,bottom=4mm,sharp corners,
overlay={\fill[cfra] (frame.north west) rectangle ([yshift=-1.25cm]frame.north east);
\node[anchor=west,text=white,font=\sffamily\bfseries\LARGE] at ([xshift=5mm,yshift=-0.62cm]frame.north west) {Question 43};
\node[anchor=east,text=white,font=\sffamily\itshape\large] at ([xshift=-5mm,yshift=-0.62cm]frame.north east) {Fractions};}]
\vspace*{1.75cm}
\begin{minipage}[t][6.9cm][c]{\linewidth}\LARGE\centering
\begin{minipage}{0.95\linewidth}\raggedright\linespread{1.25}\selectfont Que vaut $\dfrac{2}{7}+\dfrac{3}{7}$~?\end{minipage}
\end{minipage}

\vfill
\begin{tcolorbox}[enhanced,nobeforeafter,width=7.1cm,height=2.1cm,valign=center,colback=cfra!6,colframe=cfra,arc=2.5mm,boxrule=0.9pt,left=3mm]{\Large\bfseries\color{cfra}a)}\quad\begin{minipage}[c]{5.6cm}\raggedright\Large $\dfrac{5}{14}$\end{minipage}\end{tcolorbox}\hfill\begin{tcolorbox}[enhanced,nobeforeafter,width=7.1cm,height=2.1cm,valign=center,colback=cfra!6,colframe=cfra,arc=2.5mm,boxrule=0.9pt,left=3mm]{\Large\bfseries\color{cfra}b)}\quad\begin{minipage}[c]{5.6cm}\raggedright\Large $\dfrac{6}{7}$\end{minipage}\end{tcolorbox}\par\vspace{2.5mm}
\begin{tcolorbox}[enhanced,nobeforeafter,width=7.1cm,height=2.1cm,valign=center,colback=cfra!6,colframe=cfra,arc=2.5mm,boxrule=0.9pt,left=3mm]{\Large\bfseries\color{cfra}c)}\quad\begin{minipage}[c]{5.6cm}\raggedright\Large $\dfrac{5}{7}$\end{minipage}\end{tcolorbox}\hfill\begin{tcolorbox}[enhanced,nobeforeafter,width=7.1cm,height=2.1cm,valign=center,colback=cfra!6,colframe=cfra,arc=2.5mm,boxrule=0.9pt,left=3mm]{\Large\bfseries\color{cfra}d)}\quad\begin{minipage}[c]{5.6cm}\raggedright\Large $\dfrac{5}{49}$\end{minipage}\end{tcolorbox}
\end{tcolorbox}
\newpage
\begin{tcolorbox}[enhanced,width=16cm,height=14cm,colback=white,colframe=cfra,boxrule=2pt,arc=0mm,
left=6mm,right=6mm,top=0mm,bottom=4mm,sharp corners,
overlay={\fill[cfra] (frame.north west) rectangle ([yshift=-1.25cm]frame.north east);
\node[anchor=west,text=white,font=\sffamily\bfseries\LARGE] at ([xshift=5mm,yshift=-0.62cm]frame.north west) {Question 44};
\node[anchor=east,text=white,font=\sffamily\itshape\large] at ([xshift=-5mm,yshift=-0.62cm]frame.north east) {Fractions};}]
\vspace*{1.75cm}
\begin{minipage}[t][6.9cm][c]{\linewidth}\LARGE\centering
\begin{minipage}{0.95\linewidth}\raggedright\linespread{1.25}\selectfont Que vaut $\dfrac{4}{9}+\dfrac{3}{9}$~?\end{minipage}
\end{minipage}

\vfill
\begin{tcolorbox}[enhanced,nobeforeafter,width=7.1cm,height=2.1cm,valign=center,colback=cfra!6,colframe=cfra,arc=2.5mm,boxrule=0.9pt,left=3mm]{\Large\bfseries\color{cfra}a)}\quad\begin{minipage}[c]{5.6cm}\raggedright\Large $\dfrac{7}{81}$\end{minipage}\end{tcolorbox}\hfill\begin{tcolorbox}[enhanced,nobeforeafter,width=7.1cm,height=2.1cm,valign=center,colback=cfra!6,colframe=cfra,arc=2.5mm,boxrule=0.9pt,left=3mm]{\Large\bfseries\color{cfra}b)}\quad\begin{minipage}[c]{5.6cm}\raggedright\Large $\dfrac{12}{9}$\end{minipage}\end{tcolorbox}\par\vspace{2.5mm}
\begin{tcolorbox}[enhanced,nobeforeafter,width=7.1cm,height=2.1cm,valign=center,colback=cfra!6,colframe=cfra,arc=2.5mm,boxrule=0.9pt,left=3mm]{\Large\bfseries\color{cfra}c)}\quad\begin{minipage}[c]{5.6cm}\raggedright\Large $\dfrac{7}{18}$\end{minipage}\end{tcolorbox}\hfill\begin{tcolorbox}[enhanced,nobeforeafter,width=7.1cm,height=2.1cm,valign=center,colback=cfra!6,colframe=cfra,arc=2.5mm,boxrule=0.9pt,left=3mm]{\Large\bfseries\color{cfra}d)}\quad\begin{minipage}[c]{5.6cm}\raggedright\Large $\dfrac{7}{9}$\end{minipage}\end{tcolorbox}
\end{tcolorbox}
\newpage
\begin{tcolorbox}[enhanced,width=16cm,height=14cm,colback=white,colframe=cfra,boxrule=2pt,arc=0mm,
left=6mm,right=6mm,top=0mm,bottom=4mm,sharp corners,
overlay={\fill[cfra] (frame.north west) rectangle ([yshift=-1.25cm]frame.north east);
\node[anchor=west,text=white,font=\sffamily\bfseries\LARGE] at ([xshift=5mm,yshift=-0.62cm]frame.north west) {Question 45};
\node[anchor=east,text=white,font=\sffamily\itshape\large] at ([xshift=-5mm,yshift=-0.62cm]frame.north east) {Fractions};}]
\vspace*{1.75cm}
\begin{minipage}[t][6.9cm][c]{\linewidth}\LARGE\centering
\begin{minipage}{0.95\linewidth}\raggedright\linespread{1.25}\selectfont Un gâteau est coupé en $8$ parts égales. On mange $3$ parts. Quelle fraction du gâteau reste-t-il~?\end{minipage}
\end{minipage}

\vfill
\begin{tcolorbox}[enhanced,nobeforeafter,width=7.1cm,height=2.1cm,valign=center,colback=cfra!6,colframe=cfra,arc=2.5mm,boxrule=0.9pt,left=3mm]{\Large\bfseries\color{cfra}a)}\quad\begin{minipage}[c]{5.6cm}\raggedright\Large $\dfrac{3}{8}$\end{minipage}\end{tcolorbox}\hfill\begin{tcolorbox}[enhanced,nobeforeafter,width=7.1cm,height=2.1cm,valign=center,colback=cfra!6,colframe=cfra,arc=2.5mm,boxrule=0.9pt,left=3mm]{\Large\bfseries\color{cfra}b)}\quad\begin{minipage}[c]{5.6cm}\raggedright\Large $\dfrac{5}{3}$\end{minipage}\end{tcolorbox}\par\vspace{2.5mm}
\begin{tcolorbox}[enhanced,nobeforeafter,width=7.1cm,height=2.1cm,valign=center,colback=cfra!6,colframe=cfra,arc=2.5mm,boxrule=0.9pt,left=3mm]{\Large\bfseries\color{cfra}c)}\quad\begin{minipage}[c]{5.6cm}\raggedright\Large $\dfrac{3}{5}$\end{minipage}\end{tcolorbox}\hfill\begin{tcolorbox}[enhanced,nobeforeafter,width=7.1cm,height=2.1cm,valign=center,colback=cfra!6,colframe=cfra,arc=2.5mm,boxrule=0.9pt,left=3mm]{\Large\bfseries\color{cfra}d)}\quad\begin{minipage}[c]{5.6cm}\raggedright\Large $\dfrac{5}{8}$\end{minipage}\end{tcolorbox}
\end{tcolorbox}
\newpage
\begin{tcolorbox}[enhanced,width=16cm,height=14cm,colback=white,colframe=cfra,boxrule=2pt,arc=0mm,
left=6mm,right=6mm,top=0mm,bottom=4mm,sharp corners,
overlay={\fill[cfra] (frame.north west) rectangle ([yshift=-1.25cm]frame.north east);
\node[anchor=west,text=white,font=\sffamily\bfseries\LARGE] at ([xshift=5mm,yshift=-0.62cm]frame.north west) {Question 46};
\node[anchor=east,text=white,font=\sffamily\itshape\large] at ([xshift=-5mm,yshift=-0.62cm]frame.north east) {Fractions};}]
\vspace*{1.75cm}
\begin{minipage}[t][6.9cm][c]{\linewidth}\LARGE\centering
\begin{minipage}{0.95\linewidth}\raggedright\linespread{1.25}\selectfont Un gâteau est coupé en $6$ parts égales. On mange la part. Quelle fraction du gâteau reste-t-il~?\end{minipage}
\end{minipage}

\vfill
\begin{tcolorbox}[enhanced,nobeforeafter,width=7.1cm,height=2.1cm,valign=center,colback=cfra!6,colframe=cfra,arc=2.5mm,boxrule=0.9pt,left=3mm]{\Large\bfseries\color{cfra}a)}\quad\begin{minipage}[c]{5.6cm}\raggedright\Large $\dfrac{1}{5}$\end{minipage}\end{tcolorbox}\hfill\begin{tcolorbox}[enhanced,nobeforeafter,width=7.1cm,height=2.1cm,valign=center,colback=cfra!6,colframe=cfra,arc=2.5mm,boxrule=0.9pt,left=3mm]{\Large\bfseries\color{cfra}b)}\quad\begin{minipage}[c]{5.6cm}\raggedright\Large $\dfrac{1}{6}$\end{minipage}\end{tcolorbox}\par\vspace{2.5mm}
\begin{tcolorbox}[enhanced,nobeforeafter,width=7.1cm,height=2.1cm,valign=center,colback=cfra!6,colframe=cfra,arc=2.5mm,boxrule=0.9pt,left=3mm]{\Large\bfseries\color{cfra}c)}\quad\begin{minipage}[c]{5.6cm}\raggedright\Large $\dfrac{5}{6}$\end{minipage}\end{tcolorbox}\hfill\begin{tcolorbox}[enhanced,nobeforeafter,width=7.1cm,height=2.1cm,valign=center,colback=cfra!6,colframe=cfra,arc=2.5mm,boxrule=0.9pt,left=3mm]{\Large\bfseries\color{cfra}d)}\quad\begin{minipage}[c]{5.6cm}\raggedright\Large $\dfrac{5}{1}$\end{minipage}\end{tcolorbox}
\end{tcolorbox}
\newpage
\begin{tcolorbox}[enhanced,width=16cm,height=14cm,colback=white,colframe=cfra,boxrule=2pt,arc=0mm,
left=6mm,right=6mm,top=0mm,bottom=4mm,sharp corners,
overlay={\fill[cfra] (frame.north west) rectangle ([yshift=-1.25cm]frame.north east);
\node[anchor=west,text=white,font=\sffamily\bfseries\LARGE] at ([xshift=5mm,yshift=-0.62cm]frame.north west) {Question 47};
\node[anchor=east,text=white,font=\sffamily\itshape\large] at ([xshift=-5mm,yshift=-0.62cm]frame.north east) {Fractions};}]
\vspace*{1.75cm}
\begin{minipage}[t][6.9cm][c]{\linewidth}\LARGE\centering
\begin{minipage}{0.95\linewidth}\raggedright\linespread{1.25}\selectfont Entre quels nombres entiers consécutifs se trouve $\dfrac{17}{5}$~?\end{minipage}
\end{minipage}

\vfill
\begin{tcolorbox}[enhanced,nobeforeafter,width=7.1cm,height=2.1cm,valign=center,colback=cfra!6,colframe=cfra,arc=2.5mm,boxrule=0.9pt,left=3mm]{\Large\bfseries\color{cfra}a)}\quad\begin{minipage}[c]{5.6cm}\raggedright\Large $4\text{ et }5$\end{minipage}\end{tcolorbox}\hfill\begin{tcolorbox}[enhanced,nobeforeafter,width=7.1cm,height=2.1cm,valign=center,colback=cfra!6,colframe=cfra,arc=2.5mm,boxrule=0.9pt,left=3mm]{\Large\bfseries\color{cfra}b)}\quad\begin{minipage}[c]{5.6cm}\raggedright\Large $2\text{ et }3$\end{minipage}\end{tcolorbox}\par\vspace{2.5mm}
\begin{tcolorbox}[enhanced,nobeforeafter,width=7.1cm,height=2.1cm,valign=center,colback=cfra!6,colframe=cfra,arc=2.5mm,boxrule=0.9pt,left=3mm]{\Large\bfseries\color{cfra}c)}\quad\begin{minipage}[c]{5.6cm}\raggedright\Large $5\text{ et }6$\end{minipage}\end{tcolorbox}\hfill\begin{tcolorbox}[enhanced,nobeforeafter,width=7.1cm,height=2.1cm,valign=center,colback=cfra!6,colframe=cfra,arc=2.5mm,boxrule=0.9pt,left=3mm]{\Large\bfseries\color{cfra}d)}\quad\begin{minipage}[c]{5.6cm}\raggedright\Large $3\text{ et }4$\end{minipage}\end{tcolorbox}
\end{tcolorbox}
\newpage
\begin{tcolorbox}[enhanced,width=16cm,height=14cm,colback=white,colframe=cfra,boxrule=2pt,arc=0mm,
left=6mm,right=6mm,top=0mm,bottom=4mm,sharp corners,
overlay={\fill[cfra] (frame.north west) rectangle ([yshift=-1.25cm]frame.north east);
\node[anchor=west,text=white,font=\sffamily\bfseries\LARGE] at ([xshift=5mm,yshift=-0.62cm]frame.north west) {Question 48};
\node[anchor=east,text=white,font=\sffamily\itshape\large] at ([xshift=-5mm,yshift=-0.62cm]frame.north east) {Fractions};}]
\vspace*{1.75cm}
\begin{minipage}[t][6.9cm][c]{\linewidth}\LARGE\centering
\begin{minipage}{0.95\linewidth}\raggedright\linespread{1.25}\selectfont Entre quels nombres entiers consécutifs se trouve $\dfrac{23}{4}$~?\end{minipage}
\end{minipage}

\vfill
\begin{tcolorbox}[enhanced,nobeforeafter,width=7.1cm,height=2.1cm,valign=center,colback=cfra!6,colframe=cfra,arc=2.5mm,boxrule=0.9pt,left=3mm]{\Large\bfseries\color{cfra}a)}\quad\begin{minipage}[c]{5.6cm}\raggedright\Large $5\text{ et }6$\end{minipage}\end{tcolorbox}\hfill\begin{tcolorbox}[enhanced,nobeforeafter,width=7.1cm,height=2.1cm,valign=center,colback=cfra!6,colframe=cfra,arc=2.5mm,boxrule=0.9pt,left=3mm]{\Large\bfseries\color{cfra}b)}\quad\begin{minipage}[c]{5.6cm}\raggedright\Large $6\text{ et }7$\end{minipage}\end{tcolorbox}\par\vspace{2.5mm}
\begin{tcolorbox}[enhanced,nobeforeafter,width=7.1cm,height=2.1cm,valign=center,colback=cfra!6,colframe=cfra,arc=2.5mm,boxrule=0.9pt,left=3mm]{\Large\bfseries\color{cfra}c)}\quad\begin{minipage}[c]{5.6cm}\raggedright\Large $7\text{ et }8$\end{minipage}\end{tcolorbox}\hfill\begin{tcolorbox}[enhanced,nobeforeafter,width=7.1cm,height=2.1cm,valign=center,colback=cfra!6,colframe=cfra,arc=2.5mm,boxrule=0.9pt,left=3mm]{\Large\bfseries\color{cfra}d)}\quad\begin{minipage}[c]{5.6cm}\raggedright\Large $4\text{ et }5$\end{minipage}\end{tcolorbox}
\end{tcolorbox}
\newpage
\begin{tcolorbox}[enhanced,width=16cm,height=14cm,colback=white,colframe=cgeo,boxrule=2pt,arc=0mm,
left=6mm,right=6mm,top=0mm,bottom=4mm,sharp corners,
overlay={\fill[cgeo] (frame.north west) rectangle ([yshift=-1.25cm]frame.north east);
\node[anchor=west,text=white,font=\sffamily\bfseries\LARGE] at ([xshift=5mm,yshift=-0.62cm]frame.north west) {Question 49};
\node[anchor=east,text=white,font=\sffamily\itshape\large] at ([xshift=-5mm,yshift=-0.62cm]frame.north east) {Géométrie};}]
\vspace*{1.75cm}
\begin{minipage}[t][6.9cm][c]{\linewidth}\LARGE\centering
\begin{minipage}{0.95\linewidth}\raggedright\linespread{1.25}\selectfont Un angle mesure $135^\circ$. Quelle est sa nature~?\end{minipage}
\end{minipage}

\vfill
\begin{tcolorbox}[enhanced,nobeforeafter,width=7.1cm,height=2.1cm,valign=center,colback=cgeo!6,colframe=cgeo,arc=2.5mm,boxrule=0.9pt,left=3mm]{\Large\bfseries\color{cgeo}a)}\quad\begin{minipage}[c]{5.6cm}\raggedright\Large plat\end{minipage}\end{tcolorbox}\hfill\begin{tcolorbox}[enhanced,nobeforeafter,width=7.1cm,height=2.1cm,valign=center,colback=cgeo!6,colframe=cgeo,arc=2.5mm,boxrule=0.9pt,left=3mm]{\Large\bfseries\color{cgeo}b)}\quad\begin{minipage}[c]{5.6cm}\raggedright\Large obtus\end{minipage}\end{tcolorbox}\par\vspace{2.5mm}
\begin{tcolorbox}[enhanced,nobeforeafter,width=7.1cm,height=2.1cm,valign=center,colback=cgeo!6,colframe=cgeo,arc=2.5mm,boxrule=0.9pt,left=3mm]{\Large\bfseries\color{cgeo}c)}\quad\begin{minipage}[c]{5.6cm}\raggedright\Large aigu\end{minipage}\end{tcolorbox}\hfill\begin{tcolorbox}[enhanced,nobeforeafter,width=7.1cm,height=2.1cm,valign=center,colback=cgeo!6,colframe=cgeo,arc=2.5mm,boxrule=0.9pt,left=3mm]{\Large\bfseries\color{cgeo}d)}\quad\begin{minipage}[c]{5.6cm}\raggedright\Large droit\end{minipage}\end{tcolorbox}
\end{tcolorbox}
\newpage
\begin{tcolorbox}[enhanced,width=16cm,height=14cm,colback=white,colframe=cgeo,boxrule=2pt,arc=0mm,
left=6mm,right=6mm,top=0mm,bottom=4mm,sharp corners,
overlay={\fill[cgeo] (frame.north west) rectangle ([yshift=-1.25cm]frame.north east);
\node[anchor=west,text=white,font=\sffamily\bfseries\LARGE] at ([xshift=5mm,yshift=-0.62cm]frame.north west) {Question 50};
\node[anchor=east,text=white,font=\sffamily\itshape\large] at ([xshift=-5mm,yshift=-0.62cm]frame.north east) {Géométrie};}]
\vspace*{1.75cm}
\begin{minipage}[t][6.9cm][c]{\linewidth}\LARGE\centering
\begin{minipage}{0.95\linewidth}\raggedright\linespread{1.25}\selectfont Un angle mesure $42^\circ$. Quelle est sa nature~?\end{minipage}
\end{minipage}

\vfill
\begin{tcolorbox}[enhanced,nobeforeafter,width=7.1cm,height=2.1cm,valign=center,colback=cgeo!6,colframe=cgeo,arc=2.5mm,boxrule=0.9pt,left=3mm]{\Large\bfseries\color{cgeo}a)}\quad\begin{minipage}[c]{5.6cm}\raggedright\Large aigu\end{minipage}\end{tcolorbox}\hfill\begin{tcolorbox}[enhanced,nobeforeafter,width=7.1cm,height=2.1cm,valign=center,colback=cgeo!6,colframe=cgeo,arc=2.5mm,boxrule=0.9pt,left=3mm]{\Large\bfseries\color{cgeo}b)}\quad\begin{minipage}[c]{5.6cm}\raggedright\Large plat\end{minipage}\end{tcolorbox}\par\vspace{2.5mm}
\begin{tcolorbox}[enhanced,nobeforeafter,width=7.1cm,height=2.1cm,valign=center,colback=cgeo!6,colframe=cgeo,arc=2.5mm,boxrule=0.9pt,left=3mm]{\Large\bfseries\color{cgeo}c)}\quad\begin{minipage}[c]{5.6cm}\raggedright\Large droit\end{minipage}\end{tcolorbox}\hfill\begin{tcolorbox}[enhanced,nobeforeafter,width=7.1cm,height=2.1cm,valign=center,colback=cgeo!6,colframe=cgeo,arc=2.5mm,boxrule=0.9pt,left=3mm]{\Large\bfseries\color{cgeo}d)}\quad\begin{minipage}[c]{5.6cm}\raggedright\Large obtus\end{minipage}\end{tcolorbox}
\end{tcolorbox}
\newpage
\begin{tcolorbox}[enhanced,width=16cm,height=14cm,colback=white,colframe=cgeo,boxrule=2pt,arc=0mm,
left=6mm,right=6mm,top=0mm,bottom=4mm,sharp corners,
overlay={\fill[cgeo] (frame.north west) rectangle ([yshift=-1.25cm]frame.north east);
\node[anchor=west,text=white,font=\sffamily\bfseries\LARGE] at ([xshift=5mm,yshift=-0.62cm]frame.north west) {Question 51};
\node[anchor=east,text=white,font=\sffamily\itshape\large] at ([xshift=-5mm,yshift=-0.62cm]frame.north east) {Géométrie};}]
\vspace*{1.75cm}
\begin{minipage}[t][6.9cm][c]{\linewidth}\LARGE\centering
\begin{minipage}{0.95\linewidth}\raggedright\linespread{1.25}\selectfont L'angle droit $\widehat{xOy}$ est partagé par la demi-droite $[Oz)$ en deux angles. $\widehat{xOz}=35^\circ$. Combien mesure $\widehat{zOy}$~?\end{minipage}
\end{minipage}

\vfill
\begin{tcolorbox}[enhanced,nobeforeafter,width=7.1cm,height=2.1cm,valign=center,colback=cgeo!6,colframe=cgeo,arc=2.5mm,boxrule=0.9pt,left=3mm]{\Large\bfseries\color{cgeo}a)}\quad\begin{minipage}[c]{5.6cm}\raggedright\Large $145^\circ$\end{minipage}\end{tcolorbox}\hfill\begin{tcolorbox}[enhanced,nobeforeafter,width=7.1cm,height=2.1cm,valign=center,colback=cgeo!6,colframe=cgeo,arc=2.5mm,boxrule=0.9pt,left=3mm]{\Large\bfseries\color{cgeo}b)}\quad\begin{minipage}[c]{5.6cm}\raggedright\Large $55^\circ$\end{minipage}\end{tcolorbox}\par\vspace{2.5mm}
\begin{tcolorbox}[enhanced,nobeforeafter,width=7.1cm,height=2.1cm,valign=center,colback=cgeo!6,colframe=cgeo,arc=2.5mm,boxrule=0.9pt,left=3mm]{\Large\bfseries\color{cgeo}c)}\quad\begin{minipage}[c]{5.6cm}\raggedright\Large $125^\circ$\end{minipage}\end{tcolorbox}\hfill\begin{tcolorbox}[enhanced,nobeforeafter,width=7.1cm,height=2.1cm,valign=center,colback=cgeo!6,colframe=cgeo,arc=2.5mm,boxrule=0.9pt,left=3mm]{\Large\bfseries\color{cgeo}d)}\quad\begin{minipage}[c]{5.6cm}\raggedright\Large $35^\circ$\end{minipage}\end{tcolorbox}
\end{tcolorbox}
\newpage
\begin{tcolorbox}[enhanced,width=16cm,height=14cm,colback=white,colframe=cgeo,boxrule=2pt,arc=0mm,
left=6mm,right=6mm,top=0mm,bottom=4mm,sharp corners,
overlay={\fill[cgeo] (frame.north west) rectangle ([yshift=-1.25cm]frame.north east);
\node[anchor=west,text=white,font=\sffamily\bfseries\LARGE] at ([xshift=5mm,yshift=-0.62cm]frame.north west) {Question 52};
\node[anchor=east,text=white,font=\sffamily\itshape\large] at ([xshift=-5mm,yshift=-0.62cm]frame.north east) {Géométrie};}]
\vspace*{1.75cm}
\begin{minipage}[t][6.9cm][c]{\linewidth}\LARGE\centering
\begin{minipage}{0.95\linewidth}\raggedright\linespread{1.25}\selectfont L'angle droit $\widehat{xOy}$ est partagé par la demi-droite $[Oz)$ en deux angles. $\widehat{xOz}=28^\circ$. Combien mesure $\widehat{zOy}$~?\end{minipage}
\end{minipage}

\vfill
\begin{tcolorbox}[enhanced,nobeforeafter,width=7.1cm,height=2.1cm,valign=center,colback=cgeo!6,colframe=cgeo,arc=2.5mm,boxrule=0.9pt,left=3mm]{\Large\bfseries\color{cgeo}a)}\quad\begin{minipage}[c]{5.6cm}\raggedright\Large $118^\circ$\end{minipage}\end{tcolorbox}\hfill\begin{tcolorbox}[enhanced,nobeforeafter,width=7.1cm,height=2.1cm,valign=center,colback=cgeo!6,colframe=cgeo,arc=2.5mm,boxrule=0.9pt,left=3mm]{\Large\bfseries\color{cgeo}b)}\quad\begin{minipage}[c]{5.6cm}\raggedright\Large $28^\circ$\end{minipage}\end{tcolorbox}\par\vspace{2.5mm}
\begin{tcolorbox}[enhanced,nobeforeafter,width=7.1cm,height=2.1cm,valign=center,colback=cgeo!6,colframe=cgeo,arc=2.5mm,boxrule=0.9pt,left=3mm]{\Large\bfseries\color{cgeo}c)}\quad\begin{minipage}[c]{5.6cm}\raggedright\Large $152^\circ$\end{minipage}\end{tcolorbox}\hfill\begin{tcolorbox}[enhanced,nobeforeafter,width=7.1cm,height=2.1cm,valign=center,colback=cgeo!6,colframe=cgeo,arc=2.5mm,boxrule=0.9pt,left=3mm]{\Large\bfseries\color{cgeo}d)}\quad\begin{minipage}[c]{5.6cm}\raggedright\Large $62^\circ$\end{minipage}\end{tcolorbox}
\end{tcolorbox}
\newpage
\begin{tcolorbox}[enhanced,width=16cm,height=14cm,colback=white,colframe=cgeo,boxrule=2pt,arc=0mm,
left=6mm,right=6mm,top=0mm,bottom=4mm,sharp corners,
overlay={\fill[cgeo] (frame.north west) rectangle ([yshift=-1.25cm]frame.north east);
\node[anchor=west,text=white,font=\sffamily\bfseries\LARGE] at ([xshift=5mm,yshift=-0.62cm]frame.north west) {Question 53};
\node[anchor=east,text=white,font=\sffamily\itshape\large] at ([xshift=-5mm,yshift=-0.62cm]frame.north east) {Géométrie};}]
\vspace*{1.75cm}
\begin{minipage}[t][6.9cm][c]{\linewidth}\LARGE\centering
\begin{minipage}{0.95\linewidth}\raggedright\linespread{1.25}\selectfont Les points $A$, $O$ et $B$ sont alignés dans cet ordre et $\widehat{AOC}=120^\circ$. Combien mesure $\widehat{COB}$~?\end{minipage}
\end{minipage}

\vfill
\begin{tcolorbox}[enhanced,nobeforeafter,width=7.1cm,height=2.1cm,valign=center,colback=cgeo!6,colframe=cgeo,arc=2.5mm,boxrule=0.9pt,left=3mm]{\Large\bfseries\color{cgeo}a)}\quad\begin{minipage}[c]{5.6cm}\raggedright\Large $60^\circ$\end{minipage}\end{tcolorbox}\hfill\begin{tcolorbox}[enhanced,nobeforeafter,width=7.1cm,height=2.1cm,valign=center,colback=cgeo!6,colframe=cgeo,arc=2.5mm,boxrule=0.9pt,left=3mm]{\Large\bfseries\color{cgeo}b)}\quad\begin{minipage}[c]{5.6cm}\raggedright\Large $240^\circ$\end{minipage}\end{tcolorbox}\par\vspace{2.5mm}
\begin{tcolorbox}[enhanced,nobeforeafter,width=7.1cm,height=2.1cm,valign=center,colback=cgeo!6,colframe=cgeo,arc=2.5mm,boxrule=0.9pt,left=3mm]{\Large\bfseries\color{cgeo}c)}\quad\begin{minipage}[c]{5.6cm}\raggedright\Large $30^\circ$\end{minipage}\end{tcolorbox}\hfill\begin{tcolorbox}[enhanced,nobeforeafter,width=7.1cm,height=2.1cm,valign=center,colback=cgeo!6,colframe=cgeo,arc=2.5mm,boxrule=0.9pt,left=3mm]{\Large\bfseries\color{cgeo}d)}\quad\begin{minipage}[c]{5.6cm}\raggedright\Large $120^\circ$\end{minipage}\end{tcolorbox}
\end{tcolorbox}
\newpage
\begin{tcolorbox}[enhanced,width=16cm,height=14cm,colback=white,colframe=cgeo,boxrule=2pt,arc=0mm,
left=6mm,right=6mm,top=0mm,bottom=4mm,sharp corners,
overlay={\fill[cgeo] (frame.north west) rectangle ([yshift=-1.25cm]frame.north east);
\node[anchor=west,text=white,font=\sffamily\bfseries\LARGE] at ([xshift=5mm,yshift=-0.62cm]frame.north west) {Question 54};
\node[anchor=east,text=white,font=\sffamily\itshape\large] at ([xshift=-5mm,yshift=-0.62cm]frame.north east) {Géométrie};}]
\vspace*{1.75cm}
\begin{minipage}[t][6.9cm][c]{\linewidth}\LARGE\centering
\begin{minipage}{0.95\linewidth}\raggedright\linespread{1.25}\selectfont Les points $A$, $O$ et $B$ sont alignés dans cet ordre et $\widehat{AOC}=65^\circ$. Combien mesure $\widehat{COB}$~?\end{minipage}
\end{minipage}

\vfill
\begin{tcolorbox}[enhanced,nobeforeafter,width=7.1cm,height=2.1cm,valign=center,colback=cgeo!6,colframe=cgeo,arc=2.5mm,boxrule=0.9pt,left=3mm]{\Large\bfseries\color{cgeo}a)}\quad\begin{minipage}[c]{5.6cm}\raggedright\Large $295^\circ$\end{minipage}\end{tcolorbox}\hfill\begin{tcolorbox}[enhanced,nobeforeafter,width=7.1cm,height=2.1cm,valign=center,colback=cgeo!6,colframe=cgeo,arc=2.5mm,boxrule=0.9pt,left=3mm]{\Large\bfseries\color{cgeo}b)}\quad\begin{minipage}[c]{5.6cm}\raggedright\Large $25^\circ$\end{minipage}\end{tcolorbox}\par\vspace{2.5mm}
\begin{tcolorbox}[enhanced,nobeforeafter,width=7.1cm,height=2.1cm,valign=center,colback=cgeo!6,colframe=cgeo,arc=2.5mm,boxrule=0.9pt,left=3mm]{\Large\bfseries\color{cgeo}c)}\quad\begin{minipage}[c]{5.6cm}\raggedright\Large $65^\circ$\end{minipage}\end{tcolorbox}\hfill\begin{tcolorbox}[enhanced,nobeforeafter,width=7.1cm,height=2.1cm,valign=center,colback=cgeo!6,colframe=cgeo,arc=2.5mm,boxrule=0.9pt,left=3mm]{\Large\bfseries\color{cgeo}d)}\quad\begin{minipage}[c]{5.6cm}\raggedright\Large $115^\circ$\end{minipage}\end{tcolorbox}
\end{tcolorbox}
\newpage
\begin{tcolorbox}[enhanced,width=16cm,height=14cm,colback=white,colframe=cgeo,boxrule=2pt,arc=0mm,
left=6mm,right=6mm,top=0mm,bottom=4mm,sharp corners,
overlay={\fill[cgeo] (frame.north west) rectangle ([yshift=-1.25cm]frame.north east);
\node[anchor=west,text=white,font=\sffamily\bfseries\LARGE] at ([xshift=5mm,yshift=-0.62cm]frame.north west) {Question 55};
\node[anchor=east,text=white,font=\sffamily\itshape\large] at ([xshift=-5mm,yshift=-0.62cm]frame.north east) {Géométrie};}]
\vspace*{1.75cm}
\begin{minipage}[t][6.9cm][c]{\linewidth}\LARGE\centering
\begin{minipage}{0.95\linewidth}\raggedright\linespread{1.25}\selectfont Combien d'axes de symétrie possède un carré~?\end{minipage}
\end{minipage}

\vfill
\begin{tcolorbox}[enhanced,nobeforeafter,width=7.1cm,height=2.1cm,valign=center,colback=cgeo!6,colframe=cgeo,arc=2.5mm,boxrule=0.9pt,left=3mm]{\Large\bfseries\color{cgeo}a)}\quad\begin{minipage}[c]{5.6cm}\raggedright\Large 2\end{minipage}\end{tcolorbox}\hfill\begin{tcolorbox}[enhanced,nobeforeafter,width=7.1cm,height=2.1cm,valign=center,colback=cgeo!6,colframe=cgeo,arc=2.5mm,boxrule=0.9pt,left=3mm]{\Large\bfseries\color{cgeo}b)}\quad\begin{minipage}[c]{5.6cm}\raggedright\Large 1\end{minipage}\end{tcolorbox}\par\vspace{2.5mm}
\begin{tcolorbox}[enhanced,nobeforeafter,width=7.1cm,height=2.1cm,valign=center,colback=cgeo!6,colframe=cgeo,arc=2.5mm,boxrule=0.9pt,left=3mm]{\Large\bfseries\color{cgeo}c)}\quad\begin{minipage}[c]{5.6cm}\raggedright\Large 3\end{minipage}\end{tcolorbox}\hfill\begin{tcolorbox}[enhanced,nobeforeafter,width=7.1cm,height=2.1cm,valign=center,colback=cgeo!6,colframe=cgeo,arc=2.5mm,boxrule=0.9pt,left=3mm]{\Large\bfseries\color{cgeo}d)}\quad\begin{minipage}[c]{5.6cm}\raggedright\Large 4\end{minipage}\end{tcolorbox}
\end{tcolorbox}
\newpage
\begin{tcolorbox}[enhanced,width=16cm,height=14cm,colback=white,colframe=cgeo,boxrule=2pt,arc=0mm,
left=6mm,right=6mm,top=0mm,bottom=4mm,sharp corners,
overlay={\fill[cgeo] (frame.north west) rectangle ([yshift=-1.25cm]frame.north east);
\node[anchor=west,text=white,font=\sffamily\bfseries\LARGE] at ([xshift=5mm,yshift=-0.62cm]frame.north west) {Question 56};
\node[anchor=east,text=white,font=\sffamily\itshape\large] at ([xshift=-5mm,yshift=-0.62cm]frame.north east) {Géométrie};}]
\vspace*{1.75cm}
\begin{minipage}[t][6.9cm][c]{\linewidth}\LARGE\centering
\begin{minipage}{0.95\linewidth}\raggedright\linespread{1.25}\selectfont Combien d'axes de symétrie possède un rectangle qui n'est pas un carré~?\end{minipage}
\end{minipage}

\vfill
\begin{tcolorbox}[enhanced,nobeforeafter,width=7.1cm,height=2.1cm,valign=center,colback=cgeo!6,colframe=cgeo,arc=2.5mm,boxrule=0.9pt,left=3mm]{\Large\bfseries\color{cgeo}a)}\quad\begin{minipage}[c]{5.6cm}\raggedright\Large 1\end{minipage}\end{tcolorbox}\hfill\begin{tcolorbox}[enhanced,nobeforeafter,width=7.1cm,height=2.1cm,valign=center,colback=cgeo!6,colframe=cgeo,arc=2.5mm,boxrule=0.9pt,left=3mm]{\Large\bfseries\color{cgeo}b)}\quad\begin{minipage}[c]{5.6cm}\raggedright\Large 4\end{minipage}\end{tcolorbox}\par\vspace{2.5mm}
\begin{tcolorbox}[enhanced,nobeforeafter,width=7.1cm,height=2.1cm,valign=center,colback=cgeo!6,colframe=cgeo,arc=2.5mm,boxrule=0.9pt,left=3mm]{\Large\bfseries\color{cgeo}c)}\quad\begin{minipage}[c]{5.6cm}\raggedright\Large 2\end{minipage}\end{tcolorbox}\hfill\begin{tcolorbox}[enhanced,nobeforeafter,width=7.1cm,height=2.1cm,valign=center,colback=cgeo!6,colframe=cgeo,arc=2.5mm,boxrule=0.9pt,left=3mm]{\Large\bfseries\color{cgeo}d)}\quad\begin{minipage}[c]{5.6cm}\raggedright\Large 3\end{minipage}\end{tcolorbox}
\end{tcolorbox}
\newpage
\begin{tcolorbox}[enhanced,width=16cm,height=14cm,colback=white,colframe=cgeo,boxrule=2pt,arc=0mm,
left=6mm,right=6mm,top=0mm,bottom=4mm,sharp corners,
overlay={\fill[cgeo] (frame.north west) rectangle ([yshift=-1.25cm]frame.north east);
\node[anchor=west,text=white,font=\sffamily\bfseries\LARGE] at ([xshift=5mm,yshift=-0.62cm]frame.north west) {Question 57};
\node[anchor=east,text=white,font=\sffamily\itshape\large] at ([xshift=-5mm,yshift=-0.62cm]frame.north east) {Géométrie};}]
\vspace*{1.75cm}
\begin{minipage}[t][6.9cm][c]{\linewidth}\LARGE\centering
\begin{minipage}{0.95\linewidth}\raggedright\linespread{1.25}\selectfont Un cercle a un diamètre de $9$ cm. Quelle est la longueur de son rayon~?\end{minipage}
\end{minipage}

\vfill
\begin{tcolorbox}[enhanced,nobeforeafter,width=7.1cm,height=2.1cm,valign=center,colback=cgeo!6,colframe=cgeo,arc=2.5mm,boxrule=0.9pt,left=3mm]{\Large\bfseries\color{cgeo}a)}\quad\begin{minipage}[c]{5.6cm}\raggedright\Large $11~\text{cm}$\end{minipage}\end{tcolorbox}\hfill\begin{tcolorbox}[enhanced,nobeforeafter,width=7.1cm,height=2.1cm,valign=center,colback=cgeo!6,colframe=cgeo,arc=2.5mm,boxrule=0.9pt,left=3mm]{\Large\bfseries\color{cgeo}b)}\quad\begin{minipage}[c]{5.6cm}\raggedright\Large $4{,}5~\text{cm}$\end{minipage}\end{tcolorbox}\par\vspace{2.5mm}
\begin{tcolorbox}[enhanced,nobeforeafter,width=7.1cm,height=2.1cm,valign=center,colback=cgeo!6,colframe=cgeo,arc=2.5mm,boxrule=0.9pt,left=3mm]{\Large\bfseries\color{cgeo}c)}\quad\begin{minipage}[c]{5.6cm}\raggedright\Large $9~\text{cm}$\end{minipage}\end{tcolorbox}\hfill\begin{tcolorbox}[enhanced,nobeforeafter,width=7.1cm,height=2.1cm,valign=center,colback=cgeo!6,colframe=cgeo,arc=2.5mm,boxrule=0.9pt,left=3mm]{\Large\bfseries\color{cgeo}d)}\quad\begin{minipage}[c]{5.6cm}\raggedright\Large $18~\text{cm}$\end{minipage}\end{tcolorbox}
\end{tcolorbox}
\newpage
\begin{tcolorbox}[enhanced,width=16cm,height=14cm,colback=white,colframe=cgeo,boxrule=2pt,arc=0mm,
left=6mm,right=6mm,top=0mm,bottom=4mm,sharp corners,
overlay={\fill[cgeo] (frame.north west) rectangle ([yshift=-1.25cm]frame.north east);
\node[anchor=west,text=white,font=\sffamily\bfseries\LARGE] at ([xshift=5mm,yshift=-0.62cm]frame.north west) {Question 58};
\node[anchor=east,text=white,font=\sffamily\itshape\large] at ([xshift=-5mm,yshift=-0.62cm]frame.north east) {Géométrie};}]
\vspace*{1.75cm}
\begin{minipage}[t][6.9cm][c]{\linewidth}\LARGE\centering
\begin{minipage}{0.95\linewidth}\raggedright\linespread{1.25}\selectfont Un cercle a un rayon de $6$ cm. Quelle est la longueur de son diamètre~?\end{minipage}
\end{minipage}

\vfill
\begin{tcolorbox}[enhanced,nobeforeafter,width=7.1cm,height=2.1cm,valign=center,colback=cgeo!6,colframe=cgeo,arc=2.5mm,boxrule=0.9pt,left=3mm]{\Large\bfseries\color{cgeo}a)}\quad\begin{minipage}[c]{5.6cm}\raggedright\Large $8~\text{cm}$\end{minipage}\end{tcolorbox}\hfill\begin{tcolorbox}[enhanced,nobeforeafter,width=7.1cm,height=2.1cm,valign=center,colback=cgeo!6,colframe=cgeo,arc=2.5mm,boxrule=0.9pt,left=3mm]{\Large\bfseries\color{cgeo}b)}\quad\begin{minipage}[c]{5.6cm}\raggedright\Large $3~\text{cm}$\end{minipage}\end{tcolorbox}\par\vspace{2.5mm}
\begin{tcolorbox}[enhanced,nobeforeafter,width=7.1cm,height=2.1cm,valign=center,colback=cgeo!6,colframe=cgeo,arc=2.5mm,boxrule=0.9pt,left=3mm]{\Large\bfseries\color{cgeo}c)}\quad\begin{minipage}[c]{5.6cm}\raggedright\Large $6~\text{cm}$\end{minipage}\end{tcolorbox}\hfill\begin{tcolorbox}[enhanced,nobeforeafter,width=7.1cm,height=2.1cm,valign=center,colback=cgeo!6,colframe=cgeo,arc=2.5mm,boxrule=0.9pt,left=3mm]{\Large\bfseries\color{cgeo}d)}\quad\begin{minipage}[c]{5.6cm}\raggedright\Large $12~\text{cm}$\end{minipage}\end{tcolorbox}
\end{tcolorbox}
\newpage
\begin{tcolorbox}[enhanced,width=16cm,height=14cm,colback=white,colframe=cgeo,boxrule=2pt,arc=0mm,
left=6mm,right=6mm,top=0mm,bottom=4mm,sharp corners,
overlay={\fill[cgeo] (frame.north west) rectangle ([yshift=-1.25cm]frame.north east);
\node[anchor=west,text=white,font=\sffamily\bfseries\LARGE] at ([xshift=5mm,yshift=-0.62cm]frame.north west) {Question 59};
\node[anchor=east,text=white,font=\sffamily\itshape\large] at ([xshift=-5mm,yshift=-0.62cm]frame.north east) {Géométrie};}]
\vspace*{1.75cm}
\begin{minipage}[t][6.9cm][c]{\linewidth}\LARGE\centering
\begin{minipage}{0.95\linewidth}\raggedright\linespread{1.25}\selectfont La droite $(d_1)$ est perpendiculaire à $(d)$ et la droite $(d_2)$ est perpendiculaire à $(d)$. Que peut-on dire des droites $(d_1)$ et $(d_2)$~?\end{minipage}
\end{minipage}

\vfill
\begin{tcolorbox}[enhanced,nobeforeafter,width=7.1cm,height=2.1cm,valign=center,colback=cgeo!6,colframe=cgeo,arc=2.5mm,boxrule=0.9pt,left=3mm]{\Large\bfseries\color{cgeo}a)}\quad\begin{minipage}[c]{5.6cm}\raggedright\Large perpendiculaires\end{minipage}\end{tcolorbox}\hfill\begin{tcolorbox}[enhanced,nobeforeafter,width=7.1cm,height=2.1cm,valign=center,colback=cgeo!6,colframe=cgeo,arc=2.5mm,boxrule=0.9pt,left=3mm]{\Large\bfseries\color{cgeo}b)}\quad\begin{minipage}[c]{5.6cm}\raggedright\Large on ne peut pas savoir\end{minipage}\end{tcolorbox}\par\vspace{2.5mm}
\begin{tcolorbox}[enhanced,nobeforeafter,width=7.1cm,height=2.1cm,valign=center,colback=cgeo!6,colframe=cgeo,arc=2.5mm,boxrule=0.9pt,left=3mm]{\Large\bfseries\color{cgeo}c)}\quad\begin{minipage}[c]{5.6cm}\raggedright\large sécantes, non perpendiculaires\end{minipage}\end{tcolorbox}\hfill\begin{tcolorbox}[enhanced,nobeforeafter,width=7.1cm,height=2.1cm,valign=center,colback=cgeo!6,colframe=cgeo,arc=2.5mm,boxrule=0.9pt,left=3mm]{\Large\bfseries\color{cgeo}d)}\quad\begin{minipage}[c]{5.6cm}\raggedright\Large parallèles\end{minipage}\end{tcolorbox}
\end{tcolorbox}
\newpage
\begin{tcolorbox}[enhanced,width=16cm,height=14cm,colback=white,colframe=cgeo,boxrule=2pt,arc=0mm,
left=6mm,right=6mm,top=0mm,bottom=4mm,sharp corners,
overlay={\fill[cgeo] (frame.north west) rectangle ([yshift=-1.25cm]frame.north east);
\node[anchor=west,text=white,font=\sffamily\bfseries\LARGE] at ([xshift=5mm,yshift=-0.62cm]frame.north west) {Question 60};
\node[anchor=east,text=white,font=\sffamily\itshape\large] at ([xshift=-5mm,yshift=-0.62cm]frame.north east) {Géométrie};}]
\vspace*{1.75cm}
\begin{minipage}[t][6.9cm][c]{\linewidth}\LARGE\centering
\begin{minipage}{0.95\linewidth}\raggedright\linespread{1.25}\selectfont La droite $(a)$ est parallèle à $(c)$ et la droite $(b)$ est perpendiculaire à $(c)$. Que peut-on dire des droites $(a)$ et $(b)$~?\end{minipage}
\end{minipage}

\vfill
\begin{tcolorbox}[enhanced,nobeforeafter,width=7.1cm,height=2.1cm,valign=center,colback=cgeo!6,colframe=cgeo,arc=2.5mm,boxrule=0.9pt,left=3mm]{\Large\bfseries\color{cgeo}a)}\quad\begin{minipage}[c]{5.6cm}\raggedright\Large on ne peut pas savoir\end{minipage}\end{tcolorbox}\hfill\begin{tcolorbox}[enhanced,nobeforeafter,width=7.1cm,height=2.1cm,valign=center,colback=cgeo!6,colframe=cgeo,arc=2.5mm,boxrule=0.9pt,left=3mm]{\Large\bfseries\color{cgeo}b)}\quad\begin{minipage}[c]{5.6cm}\raggedright\large sécantes, non perpendiculaires\end{minipage}\end{tcolorbox}\par\vspace{2.5mm}
\begin{tcolorbox}[enhanced,nobeforeafter,width=7.1cm,height=2.1cm,valign=center,colback=cgeo!6,colframe=cgeo,arc=2.5mm,boxrule=0.9pt,left=3mm]{\Large\bfseries\color{cgeo}c)}\quad\begin{minipage}[c]{5.6cm}\raggedright\Large parallèles\end{minipage}\end{tcolorbox}\hfill\begin{tcolorbox}[enhanced,nobeforeafter,width=7.1cm,height=2.1cm,valign=center,colback=cgeo!6,colframe=cgeo,arc=2.5mm,boxrule=0.9pt,left=3mm]{\Large\bfseries\color{cgeo}d)}\quad\begin{minipage}[c]{5.6cm}\raggedright\Large perpendiculaires\end{minipage}\end{tcolorbox}
\end{tcolorbox}
\newpage
\begin{tcolorbox}[enhanced,width=16cm,height=14cm,colback=white,colframe=cgeo,boxrule=2pt,arc=0mm,
left=6mm,right=6mm,top=0mm,bottom=4mm,sharp corners,
overlay={\fill[cgeo] (frame.north west) rectangle ([yshift=-1.25cm]frame.north east);
\node[anchor=west,text=white,font=\sffamily\bfseries\LARGE] at ([xshift=5mm,yshift=-0.62cm]frame.north west) {Question 61};
\node[anchor=east,text=white,font=\sffamily\itshape\large] at ([xshift=-5mm,yshift=-0.62cm]frame.north east) {Géométrie};}]
\vspace*{1.75cm}
\begin{minipage}[t][6.9cm][c]{\linewidth}\LARGE\centering
\begin{minipage}{0.95\linewidth}\raggedright\linespread{1.25}\selectfont Un triangle a des côtés de $5$ cm, $5$ cm et $7$ cm. Quelle est sa nature~?\end{minipage}
\end{minipage}

\vfill
\begin{tcolorbox}[enhanced,nobeforeafter,width=7.1cm,height=2.1cm,valign=center,colback=cgeo!6,colframe=cgeo,arc=2.5mm,boxrule=0.9pt,left=3mm]{\Large\bfseries\color{cgeo}a)}\quad\begin{minipage}[c]{5.6cm}\raggedright\Large rectangle\end{minipage}\end{tcolorbox}\hfill\begin{tcolorbox}[enhanced,nobeforeafter,width=7.1cm,height=2.1cm,valign=center,colback=cgeo!6,colframe=cgeo,arc=2.5mm,boxrule=0.9pt,left=3mm]{\Large\bfseries\color{cgeo}b)}\quad\begin{minipage}[c]{5.6cm}\raggedright\Large équilatéral\end{minipage}\end{tcolorbox}\par\vspace{2.5mm}
\begin{tcolorbox}[enhanced,nobeforeafter,width=7.1cm,height=2.1cm,valign=center,colback=cgeo!6,colframe=cgeo,arc=2.5mm,boxrule=0.9pt,left=3mm]{\Large\bfseries\color{cgeo}c)}\quad\begin{minipage}[c]{5.6cm}\raggedright\Large isocèle\end{minipage}\end{tcolorbox}\hfill\begin{tcolorbox}[enhanced,nobeforeafter,width=7.1cm,height=2.1cm,valign=center,colback=cgeo!6,colframe=cgeo,arc=2.5mm,boxrule=0.9pt,left=3mm]{\Large\bfseries\color{cgeo}d)}\quad\begin{minipage}[c]{5.6cm}\raggedright\Large quelconque\end{minipage}\end{tcolorbox}
\end{tcolorbox}
\newpage
\begin{tcolorbox}[enhanced,width=16cm,height=14cm,colback=white,colframe=cgeo,boxrule=2pt,arc=0mm,
left=6mm,right=6mm,top=0mm,bottom=4mm,sharp corners,
overlay={\fill[cgeo] (frame.north west) rectangle ([yshift=-1.25cm]frame.north east);
\node[anchor=west,text=white,font=\sffamily\bfseries\LARGE] at ([xshift=5mm,yshift=-0.62cm]frame.north west) {Question 62};
\node[anchor=east,text=white,font=\sffamily\itshape\large] at ([xshift=-5mm,yshift=-0.62cm]frame.north east) {Géométrie};}]
\vspace*{1.75cm}
\begin{minipage}[t][6.9cm][c]{\linewidth}\LARGE\centering
\begin{minipage}{0.95\linewidth}\raggedright\linespread{1.25}\selectfont Un triangle a trois côtés de $6$ cm. Quelle est sa nature~?\end{minipage}
\end{minipage}

\vfill
\begin{tcolorbox}[enhanced,nobeforeafter,width=7.1cm,height=2.1cm,valign=center,colback=cgeo!6,colframe=cgeo,arc=2.5mm,boxrule=0.9pt,left=3mm]{\Large\bfseries\color{cgeo}a)}\quad\begin{minipage}[c]{5.6cm}\raggedright\Large isocèle\end{minipage}\end{tcolorbox}\hfill\begin{tcolorbox}[enhanced,nobeforeafter,width=7.1cm,height=2.1cm,valign=center,colback=cgeo!6,colframe=cgeo,arc=2.5mm,boxrule=0.9pt,left=3mm]{\Large\bfseries\color{cgeo}b)}\quad\begin{minipage}[c]{5.6cm}\raggedright\Large équilatéral\end{minipage}\end{tcolorbox}\par\vspace{2.5mm}
\begin{tcolorbox}[enhanced,nobeforeafter,width=7.1cm,height=2.1cm,valign=center,colback=cgeo!6,colframe=cgeo,arc=2.5mm,boxrule=0.9pt,left=3mm]{\Large\bfseries\color{cgeo}c)}\quad\begin{minipage}[c]{5.6cm}\raggedright\Large quelconque\end{minipage}\end{tcolorbox}\hfill\begin{tcolorbox}[enhanced,nobeforeafter,width=7.1cm,height=2.1cm,valign=center,colback=cgeo!6,colframe=cgeo,arc=2.5mm,boxrule=0.9pt,left=3mm]{\Large\bfseries\color{cgeo}d)}\quad\begin{minipage}[c]{5.6cm}\raggedright\Large rectangle\end{minipage}\end{tcolorbox}
\end{tcolorbox}
\newpage
\begin{tcolorbox}[enhanced,width=16cm,height=14cm,colback=white,colframe=cgeo,boxrule=2pt,arc=0mm,
left=6mm,right=6mm,top=0mm,bottom=4mm,sharp corners,
overlay={\fill[cgeo] (frame.north west) rectangle ([yshift=-1.25cm]frame.north east);
\node[anchor=west,text=white,font=\sffamily\bfseries\LARGE] at ([xshift=5mm,yshift=-0.62cm]frame.north west) {Question 63};
\node[anchor=east,text=white,font=\sffamily\itshape\large] at ([xshift=-5mm,yshift=-0.62cm]frame.north east) {Géométrie};}]
\vspace*{1.75cm}
\begin{minipage}[t][6.9cm][c]{\linewidth}\LARGE\centering
\begin{minipage}{0.95\linewidth}\raggedright\linespread{1.25}\selectfont Un quadrilatère a quatre angles droits et quatre côtés de même longueur. Quelle est sa nature la plus précise~?\end{minipage}
\end{minipage}

\vfill
\begin{tcolorbox}[enhanced,nobeforeafter,width=7.1cm,height=2.1cm,valign=center,colback=cgeo!6,colframe=cgeo,arc=2.5mm,boxrule=0.9pt,left=3mm]{\Large\bfseries\color{cgeo}a)}\quad\begin{minipage}[c]{5.6cm}\raggedright\Large un rectangle\end{minipage}\end{tcolorbox}\hfill\begin{tcolorbox}[enhanced,nobeforeafter,width=7.1cm,height=2.1cm,valign=center,colback=cgeo!6,colframe=cgeo,arc=2.5mm,boxrule=0.9pt,left=3mm]{\Large\bfseries\color{cgeo}b)}\quad\begin{minipage}[c]{5.6cm}\raggedright\Large un losange\end{minipage}\end{tcolorbox}\par\vspace{2.5mm}
\begin{tcolorbox}[enhanced,nobeforeafter,width=7.1cm,height=2.1cm,valign=center,colback=cgeo!6,colframe=cgeo,arc=2.5mm,boxrule=0.9pt,left=3mm]{\Large\bfseries\color{cgeo}c)}\quad\begin{minipage}[c]{5.6cm}\raggedright\Large un carré\end{minipage}\end{tcolorbox}\hfill\begin{tcolorbox}[enhanced,nobeforeafter,width=7.1cm,height=2.1cm,valign=center,colback=cgeo!6,colframe=cgeo,arc=2.5mm,boxrule=0.9pt,left=3mm]{\Large\bfseries\color{cgeo}d)}\quad\begin{minipage}[c]{5.6cm}\raggedright\Large un triangle\end{minipage}\end{tcolorbox}
\end{tcolorbox}
\newpage
\begin{tcolorbox}[enhanced,width=16cm,height=14cm,colback=white,colframe=cgeo,boxrule=2pt,arc=0mm,
left=6mm,right=6mm,top=0mm,bottom=4mm,sharp corners,
overlay={\fill[cgeo] (frame.north west) rectangle ([yshift=-1.25cm]frame.north east);
\node[anchor=west,text=white,font=\sffamily\bfseries\LARGE] at ([xshift=5mm,yshift=-0.62cm]frame.north west) {Question 64};
\node[anchor=east,text=white,font=\sffamily\itshape\large] at ([xshift=-5mm,yshift=-0.62cm]frame.north east) {Géométrie};}]
\vspace*{1.75cm}
\begin{minipage}[t][6.9cm][c]{\linewidth}\LARGE\centering
\begin{minipage}{0.95\linewidth}\raggedright\linespread{1.25}\selectfont Un quadrilatère a quatre angles droits. Quelle est sa nature la plus précise~?\end{minipage}
\end{minipage}

\vfill
\begin{tcolorbox}[enhanced,nobeforeafter,width=7.1cm,height=2.1cm,valign=center,colback=cgeo!6,colframe=cgeo,arc=2.5mm,boxrule=0.9pt,left=3mm]{\Large\bfseries\color{cgeo}a)}\quad\begin{minipage}[c]{5.6cm}\raggedright\Large un losange\end{minipage}\end{tcolorbox}\hfill\begin{tcolorbox}[enhanced,nobeforeafter,width=7.1cm,height=2.1cm,valign=center,colback=cgeo!6,colframe=cgeo,arc=2.5mm,boxrule=0.9pt,left=3mm]{\Large\bfseries\color{cgeo}b)}\quad\begin{minipage}[c]{5.6cm}\raggedright\Large un rectangle\end{minipage}\end{tcolorbox}\par\vspace{2.5mm}
\begin{tcolorbox}[enhanced,nobeforeafter,width=7.1cm,height=2.1cm,valign=center,colback=cgeo!6,colframe=cgeo,arc=2.5mm,boxrule=0.9pt,left=3mm]{\Large\bfseries\color{cgeo}c)}\quad\begin{minipage}[c]{5.6cm}\raggedright\Large un carré\end{minipage}\end{tcolorbox}\hfill\begin{tcolorbox}[enhanced,nobeforeafter,width=7.1cm,height=2.1cm,valign=center,colback=cgeo!6,colframe=cgeo,arc=2.5mm,boxrule=0.9pt,left=3mm]{\Large\bfseries\color{cgeo}d)}\quad\begin{minipage}[c]{5.6cm}\raggedright\Large un triangle\end{minipage}\end{tcolorbox}
\end{tcolorbox}
\newpage
\begin{tcolorbox}[enhanced,width=16cm,height=14cm,colback=white,colframe=cmes,boxrule=2pt,arc=0mm,
left=6mm,right=6mm,top=0mm,bottom=4mm,sharp corners,
overlay={\fill[cmes] (frame.north west) rectangle ([yshift=-1.25cm]frame.north east);
\node[anchor=west,text=white,font=\sffamily\bfseries\LARGE] at ([xshift=5mm,yshift=-0.62cm]frame.north west) {Question 65};
\node[anchor=east,text=white,font=\sffamily\itshape\large] at ([xshift=-5mm,yshift=-0.62cm]frame.north east) {Grandeurs et mesures};}]
\vspace*{1.75cm}
\begin{minipage}[t][6.9cm][c]{\linewidth}\LARGE\centering
\begin{minipage}{0.95\linewidth}\raggedright\linespread{1.25}\selectfont Un rectangle a une longueur de $8$ cm et une largeur de $5$ cm. Quel est son périmètre~?\end{minipage}
\end{minipage}

\vfill
\begin{tcolorbox}[enhanced,nobeforeafter,width=7.1cm,height=2.1cm,valign=center,colback=cmes!6,colframe=cmes,arc=2.5mm,boxrule=0.9pt,left=3mm]{\Large\bfseries\color{cmes}a)}\quad\begin{minipage}[c]{5.6cm}\raggedright\Large $13~\text{cm}$\end{minipage}\end{tcolorbox}\hfill\begin{tcolorbox}[enhanced,nobeforeafter,width=7.1cm,height=2.1cm,valign=center,colback=cmes!6,colframe=cmes,arc=2.5mm,boxrule=0.9pt,left=3mm]{\Large\bfseries\color{cmes}b)}\quad\begin{minipage}[c]{5.6cm}\raggedright\Large $40~\text{cm}$\end{minipage}\end{tcolorbox}\par\vspace{2.5mm}
\begin{tcolorbox}[enhanced,nobeforeafter,width=7.1cm,height=2.1cm,valign=center,colback=cmes!6,colframe=cmes,arc=2.5mm,boxrule=0.9pt,left=3mm]{\Large\bfseries\color{cmes}c)}\quad\begin{minipage}[c]{5.6cm}\raggedright\Large $26~\text{cm}$\end{minipage}\end{tcolorbox}\hfill\begin{tcolorbox}[enhanced,nobeforeafter,width=7.1cm,height=2.1cm,valign=center,colback=cmes!6,colframe=cmes,arc=2.5mm,boxrule=0.9pt,left=3mm]{\Large\bfseries\color{cmes}d)}\quad\begin{minipage}[c]{5.6cm}\raggedright\Large $21~\text{cm}$\end{minipage}\end{tcolorbox}
\end{tcolorbox}
\newpage
\begin{tcolorbox}[enhanced,width=16cm,height=14cm,colback=white,colframe=cmes,boxrule=2pt,arc=0mm,
left=6mm,right=6mm,top=0mm,bottom=4mm,sharp corners,
overlay={\fill[cmes] (frame.north west) rectangle ([yshift=-1.25cm]frame.north east);
\node[anchor=west,text=white,font=\sffamily\bfseries\LARGE] at ([xshift=5mm,yshift=-0.62cm]frame.north west) {Question 66};
\node[anchor=east,text=white,font=\sffamily\itshape\large] at ([xshift=-5mm,yshift=-0.62cm]frame.north east) {Grandeurs et mesures};}]
\vspace*{1.75cm}
\begin{minipage}[t][6.9cm][c]{\linewidth}\LARGE\centering
\begin{minipage}{0.95\linewidth}\raggedright\linespread{1.25}\selectfont Un rectangle a une longueur de $12$ cm et une largeur de $7$ cm. Quel est son périmètre~?\end{minipage}
\end{minipage}

\vfill
\begin{tcolorbox}[enhanced,nobeforeafter,width=7.1cm,height=2.1cm,valign=center,colback=cmes!6,colframe=cmes,arc=2.5mm,boxrule=0.9pt,left=3mm]{\Large\bfseries\color{cmes}a)}\quad\begin{minipage}[c]{5.6cm}\raggedright\Large $84~\text{cm}$\end{minipage}\end{tcolorbox}\hfill\begin{tcolorbox}[enhanced,nobeforeafter,width=7.1cm,height=2.1cm,valign=center,colback=cmes!6,colframe=cmes,arc=2.5mm,boxrule=0.9pt,left=3mm]{\Large\bfseries\color{cmes}b)}\quad\begin{minipage}[c]{5.6cm}\raggedright\Large $31~\text{cm}$\end{minipage}\end{tcolorbox}\par\vspace{2.5mm}
\begin{tcolorbox}[enhanced,nobeforeafter,width=7.1cm,height=2.1cm,valign=center,colback=cmes!6,colframe=cmes,arc=2.5mm,boxrule=0.9pt,left=3mm]{\Large\bfseries\color{cmes}c)}\quad\begin{minipage}[c]{5.6cm}\raggedright\Large $38~\text{cm}$\end{minipage}\end{tcolorbox}\hfill\begin{tcolorbox}[enhanced,nobeforeafter,width=7.1cm,height=2.1cm,valign=center,colback=cmes!6,colframe=cmes,arc=2.5mm,boxrule=0.9pt,left=3mm]{\Large\bfseries\color{cmes}d)}\quad\begin{minipage}[c]{5.6cm}\raggedright\Large $19~\text{cm}$\end{minipage}\end{tcolorbox}
\end{tcolorbox}
\newpage
\begin{tcolorbox}[enhanced,width=16cm,height=14cm,colback=white,colframe=cmes,boxrule=2pt,arc=0mm,
left=6mm,right=6mm,top=0mm,bottom=4mm,sharp corners,
overlay={\fill[cmes] (frame.north west) rectangle ([yshift=-1.25cm]frame.north east);
\node[anchor=west,text=white,font=\sffamily\bfseries\LARGE] at ([xshift=5mm,yshift=-0.62cm]frame.north west) {Question 67};
\node[anchor=east,text=white,font=\sffamily\itshape\large] at ([xshift=-5mm,yshift=-0.62cm]frame.north east) {Grandeurs et mesures};}]
\vspace*{1.75cm}
\begin{minipage}[t][6.9cm][c]{\linewidth}\LARGE\centering
\begin{minipage}{0.95\linewidth}\raggedright\linespread{1.25}\selectfont Un rectangle a une longueur de $8$ cm et une largeur de $5$ cm. Quelle est son aire~?\end{minipage}
\end{minipage}

\vfill
\begin{tcolorbox}[enhanced,nobeforeafter,width=7.1cm,height=2.1cm,valign=center,colback=cmes!6,colframe=cmes,arc=2.5mm,boxrule=0.9pt,left=3mm]{\Large\bfseries\color{cmes}a)}\quad\begin{minipage}[c]{5.6cm}\raggedright\Large $40~\text{cm}^2$\end{minipage}\end{tcolorbox}\hfill\begin{tcolorbox}[enhanced,nobeforeafter,width=7.1cm,height=2.1cm,valign=center,colback=cmes!6,colframe=cmes,arc=2.5mm,boxrule=0.9pt,left=3mm]{\Large\bfseries\color{cmes}b)}\quad\begin{minipage}[c]{5.6cm}\raggedright\Large $13~\text{cm}^2$\end{minipage}\end{tcolorbox}\par\vspace{2.5mm}
\begin{tcolorbox}[enhanced,nobeforeafter,width=7.1cm,height=2.1cm,valign=center,colback=cmes!6,colframe=cmes,arc=2.5mm,boxrule=0.9pt,left=3mm]{\Large\bfseries\color{cmes}c)}\quad\begin{minipage}[c]{5.6cm}\raggedright\Large $20~\text{cm}^2$\end{minipage}\end{tcolorbox}\hfill\begin{tcolorbox}[enhanced,nobeforeafter,width=7.1cm,height=2.1cm,valign=center,colback=cmes!6,colframe=cmes,arc=2.5mm,boxrule=0.9pt,left=3mm]{\Large\bfseries\color{cmes}d)}\quad\begin{minipage}[c]{5.6cm}\raggedright\Large $26~\text{cm}^2$\end{minipage}\end{tcolorbox}
\end{tcolorbox}
\newpage
\begin{tcolorbox}[enhanced,width=16cm,height=14cm,colback=white,colframe=cmes,boxrule=2pt,arc=0mm,
left=6mm,right=6mm,top=0mm,bottom=4mm,sharp corners,
overlay={\fill[cmes] (frame.north west) rectangle ([yshift=-1.25cm]frame.north east);
\node[anchor=west,text=white,font=\sffamily\bfseries\LARGE] at ([xshift=5mm,yshift=-0.62cm]frame.north west) {Question 68};
\node[anchor=east,text=white,font=\sffamily\itshape\large] at ([xshift=-5mm,yshift=-0.62cm]frame.north east) {Grandeurs et mesures};}]
\vspace*{1.75cm}
\begin{minipage}[t][6.9cm][c]{\linewidth}\LARGE\centering
\begin{minipage}{0.95\linewidth}\raggedright\linespread{1.25}\selectfont Un rectangle a une longueur de $12$ cm et une largeur de $7$ cm. Quelle est son aire~?\end{minipage}
\end{minipage}

\vfill
\begin{tcolorbox}[enhanced,nobeforeafter,width=7.1cm,height=2.1cm,valign=center,colback=cmes!6,colframe=cmes,arc=2.5mm,boxrule=0.9pt,left=3mm]{\Large\bfseries\color{cmes}a)}\quad\begin{minipage}[c]{5.6cm}\raggedright\Large $19~\text{cm}^2$\end{minipage}\end{tcolorbox}\hfill\begin{tcolorbox}[enhanced,nobeforeafter,width=7.1cm,height=2.1cm,valign=center,colback=cmes!6,colframe=cmes,arc=2.5mm,boxrule=0.9pt,left=3mm]{\Large\bfseries\color{cmes}b)}\quad\begin{minipage}[c]{5.6cm}\raggedright\Large $38~\text{cm}^2$\end{minipage}\end{tcolorbox}\par\vspace{2.5mm}
\begin{tcolorbox}[enhanced,nobeforeafter,width=7.1cm,height=2.1cm,valign=center,colback=cmes!6,colframe=cmes,arc=2.5mm,boxrule=0.9pt,left=3mm]{\Large\bfseries\color{cmes}c)}\quad\begin{minipage}[c]{5.6cm}\raggedright\Large $84~\text{cm}^2$\end{minipage}\end{tcolorbox}\hfill\begin{tcolorbox}[enhanced,nobeforeafter,width=7.1cm,height=2.1cm,valign=center,colback=cmes!6,colframe=cmes,arc=2.5mm,boxrule=0.9pt,left=3mm]{\Large\bfseries\color{cmes}d)}\quad\begin{minipage}[c]{5.6cm}\raggedright\Large $42~\text{cm}^2$\end{minipage}\end{tcolorbox}
\end{tcolorbox}
\newpage
\begin{tcolorbox}[enhanced,width=16cm,height=14cm,colback=white,colframe=cmes,boxrule=2pt,arc=0mm,
left=6mm,right=6mm,top=0mm,bottom=4mm,sharp corners,
overlay={\fill[cmes] (frame.north west) rectangle ([yshift=-1.25cm]frame.north east);
\node[anchor=west,text=white,font=\sffamily\bfseries\LARGE] at ([xshift=5mm,yshift=-0.62cm]frame.north west) {Question 69};
\node[anchor=east,text=white,font=\sffamily\itshape\large] at ([xshift=-5mm,yshift=-0.62cm]frame.north east) {Grandeurs et mesures};}]
\vspace*{1.75cm}
\begin{minipage}[t][6.9cm][c]{\linewidth}\LARGE\centering
\begin{minipage}{0.95\linewidth}\raggedright\linespread{1.25}\selectfont Un triangle est rectangle ; les deux côtés de l'angle droit mesurent $6$ cm et $4$ cm. Quelle est son aire~?\end{minipage}
\end{minipage}

\vfill
\begin{tcolorbox}[enhanced,nobeforeafter,width=7.1cm,height=2.1cm,valign=center,colback=cmes!6,colframe=cmes,arc=2.5mm,boxrule=0.9pt,left=3mm]{\Large\bfseries\color{cmes}a)}\quad\begin{minipage}[c]{5.6cm}\raggedright\Large $12~\text{cm}^2$\end{minipage}\end{tcolorbox}\hfill\begin{tcolorbox}[enhanced,nobeforeafter,width=7.1cm,height=2.1cm,valign=center,colback=cmes!6,colframe=cmes,arc=2.5mm,boxrule=0.9pt,left=3mm]{\Large\bfseries\color{cmes}b)}\quad\begin{minipage}[c]{5.6cm}\raggedright\Large $10~\text{cm}^2$\end{minipage}\end{tcolorbox}\par\vspace{2.5mm}
\begin{tcolorbox}[enhanced,nobeforeafter,width=7.1cm,height=2.1cm,valign=center,colback=cmes!6,colframe=cmes,arc=2.5mm,boxrule=0.9pt,left=3mm]{\Large\bfseries\color{cmes}c)}\quad\begin{minipage}[c]{5.6cm}\raggedright\Large $24~\text{cm}^2$\end{minipage}\end{tcolorbox}\hfill\begin{tcolorbox}[enhanced,nobeforeafter,width=7.1cm,height=2.1cm,valign=center,colback=cmes!6,colframe=cmes,arc=2.5mm,boxrule=0.9pt,left=3mm]{\Large\bfseries\color{cmes}d)}\quad\begin{minipage}[c]{5.6cm}\raggedright\Large $20~\text{cm}^2$\end{minipage}\end{tcolorbox}
\end{tcolorbox}
\newpage
\begin{tcolorbox}[enhanced,width=16cm,height=14cm,colback=white,colframe=cmes,boxrule=2pt,arc=0mm,
left=6mm,right=6mm,top=0mm,bottom=4mm,sharp corners,
overlay={\fill[cmes] (frame.north west) rectangle ([yshift=-1.25cm]frame.north east);
\node[anchor=west,text=white,font=\sffamily\bfseries\LARGE] at ([xshift=5mm,yshift=-0.62cm]frame.north west) {Question 70};
\node[anchor=east,text=white,font=\sffamily\itshape\large] at ([xshift=-5mm,yshift=-0.62cm]frame.north east) {Grandeurs et mesures};}]
\vspace*{1.75cm}
\begin{minipage}[t][6.9cm][c]{\linewidth}\LARGE\centering
\begin{minipage}{0.95\linewidth}\raggedright\linespread{1.25}\selectfont Un triangle est rectangle ; les deux côtés de l'angle droit mesurent $10$ cm et $7$ cm. Quelle est son aire~?\end{minipage}
\end{minipage}

\vfill
\begin{tcolorbox}[enhanced,nobeforeafter,width=7.1cm,height=2.1cm,valign=center,colback=cmes!6,colframe=cmes,arc=2.5mm,boxrule=0.9pt,left=3mm]{\Large\bfseries\color{cmes}a)}\quad\begin{minipage}[c]{5.6cm}\raggedright\Large $70~\text{cm}^2$\end{minipage}\end{tcolorbox}\hfill\begin{tcolorbox}[enhanced,nobeforeafter,width=7.1cm,height=2.1cm,valign=center,colback=cmes!6,colframe=cmes,arc=2.5mm,boxrule=0.9pt,left=3mm]{\Large\bfseries\color{cmes}b)}\quad\begin{minipage}[c]{5.6cm}\raggedright\Large $35~\text{cm}^2$\end{minipage}\end{tcolorbox}\par\vspace{2.5mm}
\begin{tcolorbox}[enhanced,nobeforeafter,width=7.1cm,height=2.1cm,valign=center,colback=cmes!6,colframe=cmes,arc=2.5mm,boxrule=0.9pt,left=3mm]{\Large\bfseries\color{cmes}c)}\quad\begin{minipage}[c]{5.6cm}\raggedright\Large $34~\text{cm}^2$\end{minipage}\end{tcolorbox}\hfill\begin{tcolorbox}[enhanced,nobeforeafter,width=7.1cm,height=2.1cm,valign=center,colback=cmes!6,colframe=cmes,arc=2.5mm,boxrule=0.9pt,left=3mm]{\Large\bfseries\color{cmes}d)}\quad\begin{minipage}[c]{5.6cm}\raggedright\Large $17~\text{cm}^2$\end{minipage}\end{tcolorbox}
\end{tcolorbox}
\newpage
\begin{tcolorbox}[enhanced,width=16cm,height=14cm,colback=white,colframe=cmes,boxrule=2pt,arc=0mm,
left=6mm,right=6mm,top=0mm,bottom=4mm,sharp corners,
overlay={\fill[cmes] (frame.north west) rectangle ([yshift=-1.25cm]frame.north east);
\node[anchor=west,text=white,font=\sffamily\bfseries\LARGE] at ([xshift=5mm,yshift=-0.62cm]frame.north west) {Question 71};
\node[anchor=east,text=white,font=\sffamily\itshape\large] at ([xshift=-5mm,yshift=-0.62cm]frame.north east) {Grandeurs et mesures};}]
\vspace*{1.75cm}
\begin{minipage}[t][6.9cm][c]{\linewidth}\LARGE\centering
\begin{minipage}{0.95\linewidth}\raggedright\linespread{1.25}\selectfont Compléter~: $3{,}5~\text{m}=\;?\;\text{cm}$\end{minipage}
\end{minipage}

\vfill
\begin{tcolorbox}[enhanced,nobeforeafter,width=7.1cm,height=2.1cm,valign=center,colback=cmes!6,colframe=cmes,arc=2.5mm,boxrule=0.9pt,left=3mm]{\Large\bfseries\color{cmes}a)}\quad\begin{minipage}[c]{5.6cm}\raggedright\Large $350~\text{cm}$\end{minipage}\end{tcolorbox}\hfill\begin{tcolorbox}[enhanced,nobeforeafter,width=7.1cm,height=2.1cm,valign=center,colback=cmes!6,colframe=cmes,arc=2.5mm,boxrule=0.9pt,left=3mm]{\Large\bfseries\color{cmes}b)}\quad\begin{minipage}[c]{5.6cm}\raggedright\Large $3\,500~\text{cm}$\end{minipage}\end{tcolorbox}\par\vspace{2.5mm}
\begin{tcolorbox}[enhanced,nobeforeafter,width=7.1cm,height=2.1cm,valign=center,colback=cmes!6,colframe=cmes,arc=2.5mm,boxrule=0.9pt,left=3mm]{\Large\bfseries\color{cmes}c)}\quad\begin{minipage}[c]{5.6cm}\raggedright\Large $35\,000~\text{cm}$\end{minipage}\end{tcolorbox}\hfill\begin{tcolorbox}[enhanced,nobeforeafter,width=7.1cm,height=2.1cm,valign=center,colback=cmes!6,colframe=cmes,arc=2.5mm,boxrule=0.9pt,left=3mm]{\Large\bfseries\color{cmes}d)}\quad\begin{minipage}[c]{5.6cm}\raggedright\Large $35~\text{cm}$\end{minipage}\end{tcolorbox}
\end{tcolorbox}
\newpage
\begin{tcolorbox}[enhanced,width=16cm,height=14cm,colback=white,colframe=cmes,boxrule=2pt,arc=0mm,
left=6mm,right=6mm,top=0mm,bottom=4mm,sharp corners,
overlay={\fill[cmes] (frame.north west) rectangle ([yshift=-1.25cm]frame.north east);
\node[anchor=west,text=white,font=\sffamily\bfseries\LARGE] at ([xshift=5mm,yshift=-0.62cm]frame.north west) {Question 72};
\node[anchor=east,text=white,font=\sffamily\itshape\large] at ([xshift=-5mm,yshift=-0.62cm]frame.north east) {Grandeurs et mesures};}]
\vspace*{1.75cm}
\begin{minipage}[t][6.9cm][c]{\linewidth}\LARGE\centering
\begin{minipage}{0.95\linewidth}\raggedright\linespread{1.25}\selectfont Compléter~: $1\,250~\text{m}=\;?\;\text{km}$\end{minipage}
\end{minipage}

\vfill
\begin{tcolorbox}[enhanced,nobeforeafter,width=7.1cm,height=2.1cm,valign=center,colback=cmes!6,colframe=cmes,arc=2.5mm,boxrule=0.9pt,left=3mm]{\Large\bfseries\color{cmes}a)}\quad\begin{minipage}[c]{5.6cm}\raggedright\Large $0{,}125~\text{km}$\end{minipage}\end{tcolorbox}\hfill\begin{tcolorbox}[enhanced,nobeforeafter,width=7.1cm,height=2.1cm,valign=center,colback=cmes!6,colframe=cmes,arc=2.5mm,boxrule=0.9pt,left=3mm]{\Large\bfseries\color{cmes}b)}\quad\begin{minipage}[c]{5.6cm}\raggedright\Large $125~\text{km}$\end{minipage}\end{tcolorbox}\par\vspace{2.5mm}
\begin{tcolorbox}[enhanced,nobeforeafter,width=7.1cm,height=2.1cm,valign=center,colback=cmes!6,colframe=cmes,arc=2.5mm,boxrule=0.9pt,left=3mm]{\Large\bfseries\color{cmes}c)}\quad\begin{minipage}[c]{5.6cm}\raggedright\Large $12{,}5~\text{km}$\end{minipage}\end{tcolorbox}\hfill\begin{tcolorbox}[enhanced,nobeforeafter,width=7.1cm,height=2.1cm,valign=center,colback=cmes!6,colframe=cmes,arc=2.5mm,boxrule=0.9pt,left=3mm]{\Large\bfseries\color{cmes}d)}\quad\begin{minipage}[c]{5.6cm}\raggedright\Large $1{,}25~\text{km}$\end{minipage}\end{tcolorbox}
\end{tcolorbox}
\newpage
\begin{tcolorbox}[enhanced,width=16cm,height=14cm,colback=white,colframe=cmes,boxrule=2pt,arc=0mm,
left=6mm,right=6mm,top=0mm,bottom=4mm,sharp corners,
overlay={\fill[cmes] (frame.north west) rectangle ([yshift=-1.25cm]frame.north east);
\node[anchor=west,text=white,font=\sffamily\bfseries\LARGE] at ([xshift=5mm,yshift=-0.62cm]frame.north west) {Question 73};
\node[anchor=east,text=white,font=\sffamily\itshape\large] at ([xshift=-5mm,yshift=-0.62cm]frame.north east) {Grandeurs et mesures};}]
\vspace*{1.75cm}
\begin{minipage}[t][6.9cm][c]{\linewidth}\LARGE\centering
\begin{minipage}{0.95\linewidth}\raggedright\linespread{1.25}\selectfont Compléter~: $2{,}5~\text{kg}=\;?\;\text{g}$\end{minipage}
\end{minipage}

\vfill
\begin{tcolorbox}[enhanced,nobeforeafter,width=7.1cm,height=2.1cm,valign=center,colback=cmes!6,colframe=cmes,arc=2.5mm,boxrule=0.9pt,left=3mm]{\Large\bfseries\color{cmes}a)}\quad\begin{minipage}[c]{5.6cm}\raggedright\Large $250\,000~\text{g}$\end{minipage}\end{tcolorbox}\hfill\begin{tcolorbox}[enhanced,nobeforeafter,width=7.1cm,height=2.1cm,valign=center,colback=cmes!6,colframe=cmes,arc=2.5mm,boxrule=0.9pt,left=3mm]{\Large\bfseries\color{cmes}b)}\quad\begin{minipage}[c]{5.6cm}\raggedright\Large $250~\text{g}$\end{minipage}\end{tcolorbox}\par\vspace{2.5mm}
\begin{tcolorbox}[enhanced,nobeforeafter,width=7.1cm,height=2.1cm,valign=center,colback=cmes!6,colframe=cmes,arc=2.5mm,boxrule=0.9pt,left=3mm]{\Large\bfseries\color{cmes}c)}\quad\begin{minipage}[c]{5.6cm}\raggedright\Large $2\,500~\text{g}$\end{minipage}\end{tcolorbox}\hfill\begin{tcolorbox}[enhanced,nobeforeafter,width=7.1cm,height=2.1cm,valign=center,colback=cmes!6,colframe=cmes,arc=2.5mm,boxrule=0.9pt,left=3mm]{\Large\bfseries\color{cmes}d)}\quad\begin{minipage}[c]{5.6cm}\raggedright\Large $25\,000~\text{g}$\end{minipage}\end{tcolorbox}
\end{tcolorbox}
\newpage
\begin{tcolorbox}[enhanced,width=16cm,height=14cm,colback=white,colframe=cmes,boxrule=2pt,arc=0mm,
left=6mm,right=6mm,top=0mm,bottom=4mm,sharp corners,
overlay={\fill[cmes] (frame.north west) rectangle ([yshift=-1.25cm]frame.north east);
\node[anchor=west,text=white,font=\sffamily\bfseries\LARGE] at ([xshift=5mm,yshift=-0.62cm]frame.north west) {Question 74};
\node[anchor=east,text=white,font=\sffamily\itshape\large] at ([xshift=-5mm,yshift=-0.62cm]frame.north east) {Grandeurs et mesures};}]
\vspace*{1.75cm}
\begin{minipage}[t][6.9cm][c]{\linewidth}\LARGE\centering
\begin{minipage}{0.95\linewidth}\raggedright\linespread{1.25}\selectfont Compléter~: $750~\text{g}=\;?\;\text{kg}$\end{minipage}
\end{minipage}

\vfill
\begin{tcolorbox}[enhanced,nobeforeafter,width=7.1cm,height=2.1cm,valign=center,colback=cmes!6,colframe=cmes,arc=2.5mm,boxrule=0.9pt,left=3mm]{\Large\bfseries\color{cmes}a)}\quad\begin{minipage}[c]{5.6cm}\raggedright\Large $0{,}075~\text{kg}$\end{minipage}\end{tcolorbox}\hfill\begin{tcolorbox}[enhanced,nobeforeafter,width=7.1cm,height=2.1cm,valign=center,colback=cmes!6,colframe=cmes,arc=2.5mm,boxrule=0.9pt,left=3mm]{\Large\bfseries\color{cmes}b)}\quad\begin{minipage}[c]{5.6cm}\raggedright\Large $0{,}75~\text{kg}$\end{minipage}\end{tcolorbox}\par\vspace{2.5mm}
\begin{tcolorbox}[enhanced,nobeforeafter,width=7.1cm,height=2.1cm,valign=center,colback=cmes!6,colframe=cmes,arc=2.5mm,boxrule=0.9pt,left=3mm]{\Large\bfseries\color{cmes}c)}\quad\begin{minipage}[c]{5.6cm}\raggedright\Large $75~\text{kg}$\end{minipage}\end{tcolorbox}\hfill\begin{tcolorbox}[enhanced,nobeforeafter,width=7.1cm,height=2.1cm,valign=center,colback=cmes!6,colframe=cmes,arc=2.5mm,boxrule=0.9pt,left=3mm]{\Large\bfseries\color{cmes}d)}\quad\begin{minipage}[c]{5.6cm}\raggedright\Large $7{,}5~\text{kg}$\end{minipage}\end{tcolorbox}
\end{tcolorbox}
\newpage
\begin{tcolorbox}[enhanced,width=16cm,height=14cm,colback=white,colframe=cmes,boxrule=2pt,arc=0mm,
left=6mm,right=6mm,top=0mm,bottom=4mm,sharp corners,
overlay={\fill[cmes] (frame.north west) rectangle ([yshift=-1.25cm]frame.north east);
\node[anchor=west,text=white,font=\sffamily\bfseries\LARGE] at ([xshift=5mm,yshift=-0.62cm]frame.north west) {Question 75};
\node[anchor=east,text=white,font=\sffamily\itshape\large] at ([xshift=-5mm,yshift=-0.62cm]frame.north east) {Grandeurs et mesures};}]
\vspace*{1.75cm}
\begin{minipage}[t][6.9cm][c]{\linewidth}\LARGE\centering
\begin{minipage}{0.95\linewidth}\raggedright\linespread{1.25}\selectfont Combien de minutes y a-t-il dans 1 h 45 min~?\end{minipage}
\end{minipage}

\vfill
\begin{tcolorbox}[enhanced,nobeforeafter,width=7.1cm,height=2.1cm,valign=center,colback=cmes!6,colframe=cmes,arc=2.5mm,boxrule=0.9pt,left=3mm]{\Large\bfseries\color{cmes}a)}\quad\begin{minipage}[c]{5.6cm}\raggedright\Large $105~\text{min}$\end{minipage}\end{tcolorbox}\hfill\begin{tcolorbox}[enhanced,nobeforeafter,width=7.1cm,height=2.1cm,valign=center,colback=cmes!6,colframe=cmes,arc=2.5mm,boxrule=0.9pt,left=3mm]{\Large\bfseries\color{cmes}b)}\quad\begin{minipage}[c]{5.6cm}\raggedright\Large $145~\text{min}$\end{minipage}\end{tcolorbox}\par\vspace{2.5mm}
\begin{tcolorbox}[enhanced,nobeforeafter,width=7.1cm,height=2.1cm,valign=center,colback=cmes!6,colframe=cmes,arc=2.5mm,boxrule=0.9pt,left=3mm]{\Large\bfseries\color{cmes}c)}\quad\begin{minipage}[c]{5.6cm}\raggedright\Large $46~\text{min}$\end{minipage}\end{tcolorbox}\hfill\begin{tcolorbox}[enhanced,nobeforeafter,width=7.1cm,height=2.1cm,valign=center,colback=cmes!6,colframe=cmes,arc=2.5mm,boxrule=0.9pt,left=3mm]{\Large\bfseries\color{cmes}d)}\quad\begin{minipage}[c]{5.6cm}\raggedright\Large $15~\text{min}$\end{minipage}\end{tcolorbox}
\end{tcolorbox}
\newpage
\begin{tcolorbox}[enhanced,width=16cm,height=14cm,colback=white,colframe=cmes,boxrule=2pt,arc=0mm,
left=6mm,right=6mm,top=0mm,bottom=4mm,sharp corners,
overlay={\fill[cmes] (frame.north west) rectangle ([yshift=-1.25cm]frame.north east);
\node[anchor=west,text=white,font=\sffamily\bfseries\LARGE] at ([xshift=5mm,yshift=-0.62cm]frame.north west) {Question 76};
\node[anchor=east,text=white,font=\sffamily\itshape\large] at ([xshift=-5mm,yshift=-0.62cm]frame.north east) {Grandeurs et mesures};}]
\vspace*{1.75cm}
\begin{minipage}[t][6.9cm][c]{\linewidth}\LARGE\centering
\begin{minipage}{0.95\linewidth}\raggedright\linespread{1.25}\selectfont Combien de minutes y a-t-il dans 2 h 30 min~?\end{minipage}
\end{minipage}

\vfill
\begin{tcolorbox}[enhanced,nobeforeafter,width=7.1cm,height=2.1cm,valign=center,colback=cmes!6,colframe=cmes,arc=2.5mm,boxrule=0.9pt,left=3mm]{\Large\bfseries\color{cmes}a)}\quad\begin{minipage}[c]{5.6cm}\raggedright\Large $32~\text{min}$\end{minipage}\end{tcolorbox}\hfill\begin{tcolorbox}[enhanced,nobeforeafter,width=7.1cm,height=2.1cm,valign=center,colback=cmes!6,colframe=cmes,arc=2.5mm,boxrule=0.9pt,left=3mm]{\Large\bfseries\color{cmes}b)}\quad\begin{minipage}[c]{5.6cm}\raggedright\Large $150~\text{min}$\end{minipage}\end{tcolorbox}\par\vspace{2.5mm}
\begin{tcolorbox}[enhanced,nobeforeafter,width=7.1cm,height=2.1cm,valign=center,colback=cmes!6,colframe=cmes,arc=2.5mm,boxrule=0.9pt,left=3mm]{\Large\bfseries\color{cmes}c)}\quad\begin{minipage}[c]{5.6cm}\raggedright\Large $230~\text{min}$\end{minipage}\end{tcolorbox}\hfill\begin{tcolorbox}[enhanced,nobeforeafter,width=7.1cm,height=2.1cm,valign=center,colback=cmes!6,colframe=cmes,arc=2.5mm,boxrule=0.9pt,left=3mm]{\Large\bfseries\color{cmes}d)}\quad\begin{minipage}[c]{5.6cm}\raggedright\Large $90~\text{min}$\end{minipage}\end{tcolorbox}
\end{tcolorbox}
\newpage
\begin{tcolorbox}[enhanced,width=16cm,height=14cm,colback=white,colframe=cmes,boxrule=2pt,arc=0mm,
left=6mm,right=6mm,top=0mm,bottom=4mm,sharp corners,
overlay={\fill[cmes] (frame.north west) rectangle ([yshift=-1.25cm]frame.north east);
\node[anchor=west,text=white,font=\sffamily\bfseries\LARGE] at ([xshift=5mm,yshift=-0.62cm]frame.north west) {Question 77};
\node[anchor=east,text=white,font=\sffamily\itshape\large] at ([xshift=-5mm,yshift=-0.62cm]frame.north east) {Grandeurs et mesures};}]
\vspace*{1.75cm}
\begin{minipage}[t][6.9cm][c]{\linewidth}\LARGE\centering
\begin{minipage}{0.95\linewidth}\raggedright\linespread{1.25}\selectfont En prenant $\pi\approx3{,}14$, quel est le périmètre d'un cercle de diamètre $10$ cm~?\end{minipage}
\end{minipage}

\vfill
\begin{tcolorbox}[enhanced,nobeforeafter,width=7.1cm,height=2.1cm,valign=center,colback=cmes!6,colframe=cmes,arc=2.5mm,boxrule=0.9pt,left=3mm]{\Large\bfseries\color{cmes}a)}\quad\begin{minipage}[c]{5.6cm}\raggedright\Large $62{,}8~\text{cm}$\end{minipage}\end{tcolorbox}\hfill\begin{tcolorbox}[enhanced,nobeforeafter,width=7.1cm,height=2.1cm,valign=center,colback=cmes!6,colframe=cmes,arc=2.5mm,boxrule=0.9pt,left=3mm]{\Large\bfseries\color{cmes}b)}\quad\begin{minipage}[c]{5.6cm}\raggedright\Large $15{,}7~\text{cm}$\end{minipage}\end{tcolorbox}\par\vspace{2.5mm}
\begin{tcolorbox}[enhanced,nobeforeafter,width=7.1cm,height=2.1cm,valign=center,colback=cmes!6,colframe=cmes,arc=2.5mm,boxrule=0.9pt,left=3mm]{\Large\bfseries\color{cmes}c)}\quad\begin{minipage}[c]{5.6cm}\raggedright\Large $31{,}4~\text{cm}$\end{minipage}\end{tcolorbox}\hfill\begin{tcolorbox}[enhanced,nobeforeafter,width=7.1cm,height=2.1cm,valign=center,colback=cmes!6,colframe=cmes,arc=2.5mm,boxrule=0.9pt,left=3mm]{\Large\bfseries\color{cmes}d)}\quad\begin{minipage}[c]{5.6cm}\raggedright\Large $78{,}5~\text{cm}$\end{minipage}\end{tcolorbox}
\end{tcolorbox}
\newpage
\begin{tcolorbox}[enhanced,width=16cm,height=14cm,colback=white,colframe=cmes,boxrule=2pt,arc=0mm,
left=6mm,right=6mm,top=0mm,bottom=4mm,sharp corners,
overlay={\fill[cmes] (frame.north west) rectangle ([yshift=-1.25cm]frame.north east);
\node[anchor=west,text=white,font=\sffamily\bfseries\LARGE] at ([xshift=5mm,yshift=-0.62cm]frame.north west) {Question 78};
\node[anchor=east,text=white,font=\sffamily\itshape\large] at ([xshift=-5mm,yshift=-0.62cm]frame.north east) {Grandeurs et mesures};}]
\vspace*{1.75cm}
\begin{minipage}[t][6.9cm][c]{\linewidth}\LARGE\centering
\begin{minipage}{0.95\linewidth}\raggedright\linespread{1.25}\selectfont En prenant $\pi\approx3{,}14$, quel est le périmètre d'un cercle de diamètre $4$ cm~?\end{minipage}
\end{minipage}

\vfill
\begin{tcolorbox}[enhanced,nobeforeafter,width=7.1cm,height=2.1cm,valign=center,colback=cmes!6,colframe=cmes,arc=2.5mm,boxrule=0.9pt,left=3mm]{\Large\bfseries\color{cmes}a)}\quad\begin{minipage}[c]{5.6cm}\raggedright\Large $6{,}28~\text{cm}$\end{minipage}\end{tcolorbox}\hfill\begin{tcolorbox}[enhanced,nobeforeafter,width=7.1cm,height=2.1cm,valign=center,colback=cmes!6,colframe=cmes,arc=2.5mm,boxrule=0.9pt,left=3mm]{\Large\bfseries\color{cmes}b)}\quad\begin{minipage}[c]{5.6cm}\raggedright\Large $12~\text{cm}$\end{minipage}\end{tcolorbox}\par\vspace{2.5mm}
\begin{tcolorbox}[enhanced,nobeforeafter,width=7.1cm,height=2.1cm,valign=center,colback=cmes!6,colframe=cmes,arc=2.5mm,boxrule=0.9pt,left=3mm]{\Large\bfseries\color{cmes}c)}\quad\begin{minipage}[c]{5.6cm}\raggedright\Large $12{,}56~\text{cm}$\end{minipage}\end{tcolorbox}\hfill\begin{tcolorbox}[enhanced,nobeforeafter,width=7.1cm,height=2.1cm,valign=center,colback=cmes!6,colframe=cmes,arc=2.5mm,boxrule=0.9pt,left=3mm]{\Large\bfseries\color{cmes}d)}\quad\begin{minipage}[c]{5.6cm}\raggedright\Large $25{,}12~\text{cm}$\end{minipage}\end{tcolorbox}
\end{tcolorbox}
\newpage
\begin{tcolorbox}[enhanced,width=16cm,height=14cm,colback=white,colframe=cmes,boxrule=2pt,arc=0mm,
left=6mm,right=6mm,top=0mm,bottom=4mm,sharp corners,
overlay={\fill[cmes] (frame.north west) rectangle ([yshift=-1.25cm]frame.north east);
\node[anchor=west,text=white,font=\sffamily\bfseries\LARGE] at ([xshift=5mm,yshift=-0.62cm]frame.north west) {Question 79};
\node[anchor=east,text=white,font=\sffamily\itshape\large] at ([xshift=-5mm,yshift=-0.62cm]frame.north east) {Grandeurs et mesures};}]
\vspace*{1.75cm}
\begin{minipage}[t][6.9cm][c]{\linewidth}\LARGE\centering
\begin{minipage}{0.95\linewidth}\raggedright\linespread{1.25}\selectfont Un pavé droit mesure $5$ cm de longueur, $4$ cm de largeur et $3$ cm de hauteur. Quel est son volume~?\end{minipage}
\end{minipage}

\vfill
\begin{tcolorbox}[enhanced,nobeforeafter,width=7.1cm,height=2.1cm,valign=center,colback=cmes!6,colframe=cmes,arc=2.5mm,boxrule=0.9pt,left=3mm]{\Large\bfseries\color{cmes}a)}\quad\begin{minipage}[c]{5.6cm}\raggedright\Large $12~\text{cm}^3$\end{minipage}\end{tcolorbox}\hfill\begin{tcolorbox}[enhanced,nobeforeafter,width=7.1cm,height=2.1cm,valign=center,colback=cmes!6,colframe=cmes,arc=2.5mm,boxrule=0.9pt,left=3mm]{\Large\bfseries\color{cmes}b)}\quad\begin{minipage}[c]{5.6cm}\raggedright\Large $60~\text{cm}^3$\end{minipage}\end{tcolorbox}\par\vspace{2.5mm}
\begin{tcolorbox}[enhanced,nobeforeafter,width=7.1cm,height=2.1cm,valign=center,colback=cmes!6,colframe=cmes,arc=2.5mm,boxrule=0.9pt,left=3mm]{\Large\bfseries\color{cmes}c)}\quad\begin{minipage}[c]{5.6cm}\raggedright\Large $20~\text{cm}^3$\end{minipage}\end{tcolorbox}\hfill\begin{tcolorbox}[enhanced,nobeforeafter,width=7.1cm,height=2.1cm,valign=center,colback=cmes!6,colframe=cmes,arc=2.5mm,boxrule=0.9pt,left=3mm]{\Large\bfseries\color{cmes}d)}\quad\begin{minipage}[c]{5.6cm}\raggedright\Large $94~\text{cm}^3$\end{minipage}\end{tcolorbox}
\end{tcolorbox}
\newpage
\begin{tcolorbox}[enhanced,width=16cm,height=14cm,colback=white,colframe=cmes,boxrule=2pt,arc=0mm,
left=6mm,right=6mm,top=0mm,bottom=4mm,sharp corners,
overlay={\fill[cmes] (frame.north west) rectangle ([yshift=-1.25cm]frame.north east);
\node[anchor=west,text=white,font=\sffamily\bfseries\LARGE] at ([xshift=5mm,yshift=-0.62cm]frame.north west) {Question 80};
\node[anchor=east,text=white,font=\sffamily\itshape\large] at ([xshift=-5mm,yshift=-0.62cm]frame.north east) {Grandeurs et mesures};}]
\vspace*{1.75cm}
\begin{minipage}[t][6.9cm][c]{\linewidth}\LARGE\centering
\begin{minipage}{0.95\linewidth}\raggedright\linespread{1.25}\selectfont Un pavé droit mesure $10$ cm de longueur, $2$ cm de largeur et $6$ cm de hauteur. Quel est son volume~?\end{minipage}
\end{minipage}

\vfill
\begin{tcolorbox}[enhanced,nobeforeafter,width=7.1cm,height=2.1cm,valign=center,colback=cmes!6,colframe=cmes,arc=2.5mm,boxrule=0.9pt,left=3mm]{\Large\bfseries\color{cmes}a)}\quad\begin{minipage}[c]{5.6cm}\raggedright\Large $20~\text{cm}^3$\end{minipage}\end{tcolorbox}\hfill\begin{tcolorbox}[enhanced,nobeforeafter,width=7.1cm,height=2.1cm,valign=center,colback=cmes!6,colframe=cmes,arc=2.5mm,boxrule=0.9pt,left=3mm]{\Large\bfseries\color{cmes}b)}\quad\begin{minipage}[c]{5.6cm}\raggedright\Large $184~\text{cm}^3$\end{minipage}\end{tcolorbox}\par\vspace{2.5mm}
\begin{tcolorbox}[enhanced,nobeforeafter,width=7.1cm,height=2.1cm,valign=center,colback=cmes!6,colframe=cmes,arc=2.5mm,boxrule=0.9pt,left=3mm]{\Large\bfseries\color{cmes}c)}\quad\begin{minipage}[c]{5.6cm}\raggedright\Large $120~\text{cm}^3$\end{minipage}\end{tcolorbox}\hfill\begin{tcolorbox}[enhanced,nobeforeafter,width=7.1cm,height=2.1cm,valign=center,colback=cmes!6,colframe=cmes,arc=2.5mm,boxrule=0.9pt,left=3mm]{\Large\bfseries\color{cmes}d)}\quad\begin{minipage}[c]{5.6cm}\raggedright\Large $18~\text{cm}^3$\end{minipage}\end{tcolorbox}
\end{tcolorbox}
\newpage
\begin{tcolorbox}[enhanced,width=16cm,height=14cm,colback=white,colframe=cpro,boxrule=2pt,arc=0mm,
left=6mm,right=6mm,top=0mm,bottom=4mm,sharp corners,
overlay={\fill[cpro] (frame.north west) rectangle ([yshift=-1.25cm]frame.north east);
\node[anchor=west,text=white,font=\sffamily\bfseries\LARGE] at ([xshift=5mm,yshift=-0.62cm]frame.north west) {Question 81};
\node[anchor=east,text=white,font=\sffamily\itshape\large] at ([xshift=-5mm,yshift=-0.62cm]frame.north east) {Proportionnalité et données};}]
\vspace*{1.75cm}
\begin{minipage}[t][6.9cm][c]{\linewidth}\LARGE\centering
\begin{minipage}{0.95\linewidth}\raggedright\linespread{1.25}\selectfont Pour $4$ personnes, il faut $300$ g de farine. Combien en faut-il pour $6$ personnes~?\end{minipage}
\end{minipage}

\vfill
\begin{tcolorbox}[enhanced,nobeforeafter,width=7.1cm,height=2.1cm,valign=center,colback=cpro!6,colframe=cpro,arc=2.5mm,boxrule=0.9pt,left=3mm]{\Large\bfseries\color{cpro}a)}\quad\begin{minipage}[c]{5.6cm}\raggedright\Large $75~\text{g}$\end{minipage}\end{tcolorbox}\hfill\begin{tcolorbox}[enhanced,nobeforeafter,width=7.1cm,height=2.1cm,valign=center,colback=cpro!6,colframe=cpro,arc=2.5mm,boxrule=0.9pt,left=3mm]{\Large\bfseries\color{cpro}b)}\quad\begin{minipage}[c]{5.6cm}\raggedright\Large $302~\text{g}$\end{minipage}\end{tcolorbox}\par\vspace{2.5mm}
\begin{tcolorbox}[enhanced,nobeforeafter,width=7.1cm,height=2.1cm,valign=center,colback=cpro!6,colframe=cpro,arc=2.5mm,boxrule=0.9pt,left=3mm]{\Large\bfseries\color{cpro}c)}\quad\begin{minipage}[c]{5.6cm}\raggedright\Large $1\,800~\text{g}$\end{minipage}\end{tcolorbox}\hfill\begin{tcolorbox}[enhanced,nobeforeafter,width=7.1cm,height=2.1cm,valign=center,colback=cpro!6,colframe=cpro,arc=2.5mm,boxrule=0.9pt,left=3mm]{\Large\bfseries\color{cpro}d)}\quad\begin{minipage}[c]{5.6cm}\raggedright\Large $450~\text{g}$\end{minipage}\end{tcolorbox}
\end{tcolorbox}
\newpage
\begin{tcolorbox}[enhanced,width=16cm,height=14cm,colback=white,colframe=cpro,boxrule=2pt,arc=0mm,
left=6mm,right=6mm,top=0mm,bottom=4mm,sharp corners,
overlay={\fill[cpro] (frame.north west) rectangle ([yshift=-1.25cm]frame.north east);
\node[anchor=west,text=white,font=\sffamily\bfseries\LARGE] at ([xshift=5mm,yshift=-0.62cm]frame.north west) {Question 82};
\node[anchor=east,text=white,font=\sffamily\itshape\large] at ([xshift=-5mm,yshift=-0.62cm]frame.north east) {Proportionnalité et données};}]
\vspace*{1.75cm}
\begin{minipage}[t][6.9cm][c]{\linewidth}\LARGE\centering
\begin{minipage}{0.95\linewidth}\raggedright\linespread{1.25}\selectfont Pour $5$ personnes, il faut $250$ g de farine. Combien en faut-il pour $2$ personnes~?\end{minipage}
\end{minipage}

\vfill
\begin{tcolorbox}[enhanced,nobeforeafter,width=7.1cm,height=2.1cm,valign=center,colback=cpro!6,colframe=cpro,arc=2.5mm,boxrule=0.9pt,left=3mm]{\Large\bfseries\color{cpro}a)}\quad\begin{minipage}[c]{5.6cm}\raggedright\Large $100~\text{g}$\end{minipage}\end{tcolorbox}\hfill\begin{tcolorbox}[enhanced,nobeforeafter,width=7.1cm,height=2.1cm,valign=center,colback=cpro!6,colframe=cpro,arc=2.5mm,boxrule=0.9pt,left=3mm]{\Large\bfseries\color{cpro}b)}\quad\begin{minipage}[c]{5.6cm}\raggedright\Large $50~\text{g}$\end{minipage}\end{tcolorbox}\par\vspace{2.5mm}
\begin{tcolorbox}[enhanced,nobeforeafter,width=7.1cm,height=2.1cm,valign=center,colback=cpro!6,colframe=cpro,arc=2.5mm,boxrule=0.9pt,left=3mm]{\Large\bfseries\color{cpro}c)}\quad\begin{minipage}[c]{5.6cm}\raggedright\Large $247~\text{g}$\end{minipage}\end{tcolorbox}\hfill\begin{tcolorbox}[enhanced,nobeforeafter,width=7.1cm,height=2.1cm,valign=center,colback=cpro!6,colframe=cpro,arc=2.5mm,boxrule=0.9pt,left=3mm]{\Large\bfseries\color{cpro}d)}\quad\begin{minipage}[c]{5.6cm}\raggedright\Large $500~\text{g}$\end{minipage}\end{tcolorbox}
\end{tcolorbox}
\newpage
\begin{tcolorbox}[enhanced,width=16cm,height=14cm,colback=white,colframe=cpro,boxrule=2pt,arc=0mm,
left=6mm,right=6mm,top=0mm,bottom=4mm,sharp corners,
overlay={\fill[cpro] (frame.north west) rectangle ([yshift=-1.25cm]frame.north east);
\node[anchor=west,text=white,font=\sffamily\bfseries\LARGE] at ([xshift=5mm,yshift=-0.62cm]frame.north west) {Question 83};
\node[anchor=east,text=white,font=\sffamily\itshape\large] at ([xshift=-5mm,yshift=-0.62cm]frame.north east) {Proportionnalité et données};}]
\vspace*{1.75cm}
\begin{minipage}[t][6.9cm][c]{\linewidth}\LARGE\centering
\begin{minipage}{0.95\linewidth}\raggedright\linespread{1.25}\selectfont $5$ cahiers coûtent $1\,250$ F. Combien coûte un cahier~?\end{minipage}
\end{minipage}

\vfill
\begin{tcolorbox}[enhanced,nobeforeafter,width=7.1cm,height=2.1cm,valign=center,colback=cpro!6,colframe=cpro,arc=2.5mm,boxrule=0.9pt,left=3mm]{\Large\bfseries\color{cpro}a)}\quad\begin{minipage}[c]{5.6cm}\raggedright\Large $250~\text{F}$\end{minipage}\end{tcolorbox}\hfill\begin{tcolorbox}[enhanced,nobeforeafter,width=7.1cm,height=2.1cm,valign=center,colback=cpro!6,colframe=cpro,arc=2.5mm,boxrule=0.9pt,left=3mm]{\Large\bfseries\color{cpro}b)}\quad\begin{minipage}[c]{5.6cm}\raggedright\Large $1\,245~\text{F}$\end{minipage}\end{tcolorbox}\par\vspace{2.5mm}
\begin{tcolorbox}[enhanced,nobeforeafter,width=7.1cm,height=2.1cm,valign=center,colback=cpro!6,colframe=cpro,arc=2.5mm,boxrule=0.9pt,left=3mm]{\Large\bfseries\color{cpro}c)}\quad\begin{minipage}[c]{5.6cm}\raggedright\Large $500~\text{F}$\end{minipage}\end{tcolorbox}\hfill\begin{tcolorbox}[enhanced,nobeforeafter,width=7.1cm,height=2.1cm,valign=center,colback=cpro!6,colframe=cpro,arc=2.5mm,boxrule=0.9pt,left=3mm]{\Large\bfseries\color{cpro}d)}\quad\begin{minipage}[c]{5.6cm}\raggedright\Large $255~\text{F}$\end{minipage}\end{tcolorbox}
\end{tcolorbox}
\newpage
\begin{tcolorbox}[enhanced,width=16cm,height=14cm,colback=white,colframe=cpro,boxrule=2pt,arc=0mm,
left=6mm,right=6mm,top=0mm,bottom=4mm,sharp corners,
overlay={\fill[cpro] (frame.north west) rectangle ([yshift=-1.25cm]frame.north east);
\node[anchor=west,text=white,font=\sffamily\bfseries\LARGE] at ([xshift=5mm,yshift=-0.62cm]frame.north west) {Question 84};
\node[anchor=east,text=white,font=\sffamily\itshape\large] at ([xshift=-5mm,yshift=-0.62cm]frame.north east) {Proportionnalité et données};}]
\vspace*{1.75cm}
\begin{minipage}[t][6.9cm][c]{\linewidth}\LARGE\centering
\begin{minipage}{0.95\linewidth}\raggedright\linespread{1.25}\selectfont $4$ cahiers coûtent $1\,100$ F. Combien coûte un cahier~?\end{minipage}
\end{minipage}

\vfill
\begin{tcolorbox}[enhanced,nobeforeafter,width=7.1cm,height=2.1cm,valign=center,colback=cpro!6,colframe=cpro,arc=2.5mm,boxrule=0.9pt,left=3mm]{\Large\bfseries\color{cpro}a)}\quad\begin{minipage}[c]{5.6cm}\raggedright\Large $550~\text{F}$\end{minipage}\end{tcolorbox}\hfill\begin{tcolorbox}[enhanced,nobeforeafter,width=7.1cm,height=2.1cm,valign=center,colback=cpro!6,colframe=cpro,arc=2.5mm,boxrule=0.9pt,left=3mm]{\Large\bfseries\color{cpro}b)}\quad\begin{minipage}[c]{5.6cm}\raggedright\Large $1\,096~\text{F}$\end{minipage}\end{tcolorbox}\par\vspace{2.5mm}
\begin{tcolorbox}[enhanced,nobeforeafter,width=7.1cm,height=2.1cm,valign=center,colback=cpro!6,colframe=cpro,arc=2.5mm,boxrule=0.9pt,left=3mm]{\Large\bfseries\color{cpro}c)}\quad\begin{minipage}[c]{5.6cm}\raggedright\Large $275~\text{F}$\end{minipage}\end{tcolorbox}\hfill\begin{tcolorbox}[enhanced,nobeforeafter,width=7.1cm,height=2.1cm,valign=center,colback=cpro!6,colframe=cpro,arc=2.5mm,boxrule=0.9pt,left=3mm]{\Large\bfseries\color{cpro}d)}\quad\begin{minipage}[c]{5.6cm}\raggedright\Large $279~\text{F}$\end{minipage}\end{tcolorbox}
\end{tcolorbox}
\newpage
\begin{tcolorbox}[enhanced,width=16cm,height=14cm,colback=white,colframe=cpro,boxrule=2pt,arc=0mm,
left=6mm,right=6mm,top=0mm,bottom=4mm,sharp corners,
overlay={\fill[cpro] (frame.north west) rectangle ([yshift=-1.25cm]frame.north east);
\node[anchor=west,text=white,font=\sffamily\bfseries\LARGE] at ([xshift=5mm,yshift=-0.62cm]frame.north west) {Question 85};
\node[anchor=east,text=white,font=\sffamily\itshape\large] at ([xshift=-5mm,yshift=-0.62cm]frame.north east) {Proportionnalité et données};}]
\vspace*{1.75cm}
\begin{minipage}[t][6.9cm][c]{\linewidth}\LARGE\centering
\begin{minipage}{0.95\linewidth}\raggedright\linespread{1.25}\selectfont Que vaut $50\,\%$ de $84$~?\end{minipage}
\end{minipage}

\vfill
\begin{tcolorbox}[enhanced,nobeforeafter,width=7.1cm,height=2.1cm,valign=center,colback=cpro!6,colframe=cpro,arc=2.5mm,boxrule=0.9pt,left=3mm]{\Large\bfseries\color{cpro}a)}\quad\begin{minipage}[c]{5.6cm}\raggedright\Large $84$\end{minipage}\end{tcolorbox}\hfill\begin{tcolorbox}[enhanced,nobeforeafter,width=7.1cm,height=2.1cm,valign=center,colback=cpro!6,colframe=cpro,arc=2.5mm,boxrule=0.9pt,left=3mm]{\Large\bfseries\color{cpro}b)}\quad\begin{minipage}[c]{5.6cm}\raggedright\Large $4{,}2$\end{minipage}\end{tcolorbox}\par\vspace{2.5mm}
\begin{tcolorbox}[enhanced,nobeforeafter,width=7.1cm,height=2.1cm,valign=center,colback=cpro!6,colframe=cpro,arc=2.5mm,boxrule=0.9pt,left=3mm]{\Large\bfseries\color{cpro}c)}\quad\begin{minipage}[c]{5.6cm}\raggedright\Large $52$\end{minipage}\end{tcolorbox}\hfill\begin{tcolorbox}[enhanced,nobeforeafter,width=7.1cm,height=2.1cm,valign=center,colback=cpro!6,colframe=cpro,arc=2.5mm,boxrule=0.9pt,left=3mm]{\Large\bfseries\color{cpro}d)}\quad\begin{minipage}[c]{5.6cm}\raggedright\Large $42$\end{minipage}\end{tcolorbox}
\end{tcolorbox}
\newpage
\begin{tcolorbox}[enhanced,width=16cm,height=14cm,colback=white,colframe=cpro,boxrule=2pt,arc=0mm,
left=6mm,right=6mm,top=0mm,bottom=4mm,sharp corners,
overlay={\fill[cpro] (frame.north west) rectangle ([yshift=-1.25cm]frame.north east);
\node[anchor=west,text=white,font=\sffamily\bfseries\LARGE] at ([xshift=5mm,yshift=-0.62cm]frame.north west) {Question 86};
\node[anchor=east,text=white,font=\sffamily\itshape\large] at ([xshift=-5mm,yshift=-0.62cm]frame.north east) {Proportionnalité et données};}]
\vspace*{1.75cm}
\begin{minipage}[t][6.9cm][c]{\linewidth}\LARGE\centering
\begin{minipage}{0.95\linewidth}\raggedright\linespread{1.25}\selectfont Que vaut $25\,\%$ de $120$~?\end{minipage}
\end{minipage}

\vfill
\begin{tcolorbox}[enhanced,nobeforeafter,width=7.1cm,height=2.1cm,valign=center,colback=cpro!6,colframe=cpro,arc=2.5mm,boxrule=0.9pt,left=3mm]{\Large\bfseries\color{cpro}a)}\quad\begin{minipage}[c]{5.6cm}\raggedright\Large $40$\end{minipage}\end{tcolorbox}\hfill\begin{tcolorbox}[enhanced,nobeforeafter,width=7.1cm,height=2.1cm,valign=center,colback=cpro!6,colframe=cpro,arc=2.5mm,boxrule=0.9pt,left=3mm]{\Large\bfseries\color{cpro}b)}\quad\begin{minipage}[c]{5.6cm}\raggedright\Large $30$\end{minipage}\end{tcolorbox}\par\vspace{2.5mm}
\begin{tcolorbox}[enhanced,nobeforeafter,width=7.1cm,height=2.1cm,valign=center,colback=cpro!6,colframe=cpro,arc=2.5mm,boxrule=0.9pt,left=3mm]{\Large\bfseries\color{cpro}c)}\quad\begin{minipage}[c]{5.6cm}\raggedright\Large $90$\end{minipage}\end{tcolorbox}\hfill\begin{tcolorbox}[enhanced,nobeforeafter,width=7.1cm,height=2.1cm,valign=center,colback=cpro!6,colframe=cpro,arc=2.5mm,boxrule=0.9pt,left=3mm]{\Large\bfseries\color{cpro}d)}\quad\begin{minipage}[c]{5.6cm}\raggedright\Large $4{,}8$\end{minipage}\end{tcolorbox}
\end{tcolorbox}
\newpage
\begin{tcolorbox}[enhanced,width=16cm,height=14cm,colback=white,colframe=cpro,boxrule=2pt,arc=0mm,
left=6mm,right=6mm,top=0mm,bottom=4mm,sharp corners,
overlay={\fill[cpro] (frame.north west) rectangle ([yshift=-1.25cm]frame.north east);
\node[anchor=west,text=white,font=\sffamily\bfseries\LARGE] at ([xshift=5mm,yshift=-0.62cm]frame.north west) {Question 87};
\node[anchor=east,text=white,font=\sffamily\itshape\large] at ([xshift=-5mm,yshift=-0.62cm]frame.north east) {Proportionnalité et données};}]
\vspace*{1.75cm}
\begin{minipage}[t][6.9cm][c]{\linewidth}\LARGE\centering
\begin{minipage}{0.95\linewidth}\raggedright\linespread{1.25}\selectfont Sur un plan à l'échelle $\dfrac{1}{100}$, une longueur mesure $5$ cm. Quelle est la longueur réelle~?\end{minipage}
\end{minipage}

\vfill
\begin{tcolorbox}[enhanced,nobeforeafter,width=7.1cm,height=2.1cm,valign=center,colback=cpro!6,colframe=cpro,arc=2.5mm,boxrule=0.9pt,left=3mm]{\Large\bfseries\color{cpro}a)}\quad\begin{minipage}[c]{5.6cm}\raggedright\Large $5~\text{m}$\end{minipage}\end{tcolorbox}\hfill\begin{tcolorbox}[enhanced,nobeforeafter,width=7.1cm,height=2.1cm,valign=center,colback=cpro!6,colframe=cpro,arc=2.5mm,boxrule=0.9pt,left=3mm]{\Large\bfseries\color{cpro}b)}\quad\begin{minipage}[c]{5.6cm}\raggedright\Large $50~\text{m}$\end{minipage}\end{tcolorbox}\par\vspace{2.5mm}
\begin{tcolorbox}[enhanced,nobeforeafter,width=7.1cm,height=2.1cm,valign=center,colback=cpro!6,colframe=cpro,arc=2.5mm,boxrule=0.9pt,left=3mm]{\Large\bfseries\color{cpro}c)}\quad\begin{minipage}[c]{5.6cm}\raggedright\Large $0{,}5~\text{m}$\end{minipage}\end{tcolorbox}\hfill\begin{tcolorbox}[enhanced,nobeforeafter,width=7.1cm,height=2.1cm,valign=center,colback=cpro!6,colframe=cpro,arc=2.5mm,boxrule=0.9pt,left=3mm]{\Large\bfseries\color{cpro}d)}\quad\begin{minipage}[c]{5.6cm}\raggedright\Large $500~\text{m}$\end{minipage}\end{tcolorbox}
\end{tcolorbox}
\newpage
\begin{tcolorbox}[enhanced,width=16cm,height=14cm,colback=white,colframe=cpro,boxrule=2pt,arc=0mm,
left=6mm,right=6mm,top=0mm,bottom=4mm,sharp corners,
overlay={\fill[cpro] (frame.north west) rectangle ([yshift=-1.25cm]frame.north east);
\node[anchor=west,text=white,font=\sffamily\bfseries\LARGE] at ([xshift=5mm,yshift=-0.62cm]frame.north west) {Question 88};
\node[anchor=east,text=white,font=\sffamily\itshape\large] at ([xshift=-5mm,yshift=-0.62cm]frame.north east) {Proportionnalité et données};}]
\vspace*{1.75cm}
\begin{minipage}[t][6.9cm][c]{\linewidth}\LARGE\centering
\begin{minipage}{0.95\linewidth}\raggedright\linespread{1.25}\selectfont Sur un plan à l'échelle $\dfrac{1}{1\,000}$, une longueur mesure $3$ cm. Quelle est la longueur réelle~?\end{minipage}
\end{minipage}

\vfill
\begin{tcolorbox}[enhanced,nobeforeafter,width=7.1cm,height=2.1cm,valign=center,colback=cpro!6,colframe=cpro,arc=2.5mm,boxrule=0.9pt,left=3mm]{\Large\bfseries\color{cpro}a)}\quad\begin{minipage}[c]{5.6cm}\raggedright\Large $300~\text{m}$\end{minipage}\end{tcolorbox}\hfill\begin{tcolorbox}[enhanced,nobeforeafter,width=7.1cm,height=2.1cm,valign=center,colback=cpro!6,colframe=cpro,arc=2.5mm,boxrule=0.9pt,left=3mm]{\Large\bfseries\color{cpro}b)}\quad\begin{minipage}[c]{5.6cm}\raggedright\Large $3\,000~\text{m}$\end{minipage}\end{tcolorbox}\par\vspace{2.5mm}
\begin{tcolorbox}[enhanced,nobeforeafter,width=7.1cm,height=2.1cm,valign=center,colback=cpro!6,colframe=cpro,arc=2.5mm,boxrule=0.9pt,left=3mm]{\Large\bfseries\color{cpro}c)}\quad\begin{minipage}[c]{5.6cm}\raggedright\Large $30~\text{m}$\end{minipage}\end{tcolorbox}\hfill\begin{tcolorbox}[enhanced,nobeforeafter,width=7.1cm,height=2.1cm,valign=center,colback=cpro!6,colframe=cpro,arc=2.5mm,boxrule=0.9pt,left=3mm]{\Large\bfseries\color{cpro}d)}\quad\begin{minipage}[c]{5.6cm}\raggedright\Large $3~\text{m}$\end{minipage}\end{tcolorbox}
\end{tcolorbox}
\newpage
\begin{tcolorbox}[enhanced,width=16cm,height=14cm,colback=white,colframe=cpro,boxrule=2pt,arc=0mm,
left=6mm,right=6mm,top=0mm,bottom=4mm,sharp corners,
overlay={\fill[cpro] (frame.north west) rectangle ([yshift=-1.25cm]frame.north east);
\node[anchor=west,text=white,font=\sffamily\bfseries\LARGE] at ([xshift=5mm,yshift=-0.62cm]frame.north west) {Question 89};
\node[anchor=east,text=white,font=\sffamily\itshape\large] at ([xshift=-5mm,yshift=-0.62cm]frame.north east) {Proportionnalité et données};}]
\vspace*{1.75cm}
\begin{minipage}[t][6.9cm][c]{\linewidth}\Large\centering
\begin{minipage}{0.95\linewidth}\raggedright\linespread{1.25}\selectfont Ce tableau est un tableau de proportionnalité~: $\begin{array}{|c|c|c|}\hline 2 & 3 & 5\\\hline 6 & 9 & 15\\\hline\end{array}$. Par quel nombre multiplie-t-on la 1re ligne pour obtenir la 2e~?\end{minipage}
\end{minipage}

\vfill
\begin{tcolorbox}[enhanced,nobeforeafter,width=7.1cm,height=2.1cm,valign=center,colback=cpro!6,colframe=cpro,arc=2.5mm,boxrule=0.9pt,left=3mm]{\Large\bfseries\color{cpro}a)}\quad\begin{minipage}[c]{5.6cm}\raggedright\Large $\dfrac{1}{3}$\end{minipage}\end{tcolorbox}\hfill\begin{tcolorbox}[enhanced,nobeforeafter,width=7.1cm,height=2.1cm,valign=center,colback=cpro!6,colframe=cpro,arc=2.5mm,boxrule=0.9pt,left=3mm]{\Large\bfseries\color{cpro}b)}\quad\begin{minipage}[c]{5.6cm}\raggedright\Large $3$\end{minipage}\end{tcolorbox}\par\vspace{2.5mm}
\begin{tcolorbox}[enhanced,nobeforeafter,width=7.1cm,height=2.1cm,valign=center,colback=cpro!6,colframe=cpro,arc=2.5mm,boxrule=0.9pt,left=3mm]{\Large\bfseries\color{cpro}c)}\quad\begin{minipage}[c]{5.6cm}\raggedright\Large $4$\end{minipage}\end{tcolorbox}\hfill\begin{tcolorbox}[enhanced,nobeforeafter,width=7.1cm,height=2.1cm,valign=center,colback=cpro!6,colframe=cpro,arc=2.5mm,boxrule=0.9pt,left=3mm]{\Large\bfseries\color{cpro}d)}\quad\begin{minipage}[c]{5.6cm}\raggedright\Large $6$\end{minipage}\end{tcolorbox}
\end{tcolorbox}
\newpage
\begin{tcolorbox}[enhanced,width=16cm,height=14cm,colback=white,colframe=cpro,boxrule=2pt,arc=0mm,
left=6mm,right=6mm,top=0mm,bottom=4mm,sharp corners,
overlay={\fill[cpro] (frame.north west) rectangle ([yshift=-1.25cm]frame.north east);
\node[anchor=west,text=white,font=\sffamily\bfseries\LARGE] at ([xshift=5mm,yshift=-0.62cm]frame.north west) {Question 90};
\node[anchor=east,text=white,font=\sffamily\itshape\large] at ([xshift=-5mm,yshift=-0.62cm]frame.north east) {Proportionnalité et données};}]
\vspace*{1.75cm}
\begin{minipage}[t][6.9cm][c]{\linewidth}\Large\centering
\begin{minipage}{0.95\linewidth}\raggedright\linespread{1.25}\selectfont Ce tableau est un tableau de proportionnalité~: $\begin{array}{|c|c|c|}\hline 1 & 5 & 7\\\hline 4 & 20 & 28\\\hline\end{array}$. Par quel nombre multiplie-t-on la 1re ligne pour obtenir la 2e~?\end{minipage}
\end{minipage}

\vfill
\begin{tcolorbox}[enhanced,nobeforeafter,width=7.1cm,height=2.1cm,valign=center,colback=cpro!6,colframe=cpro,arc=2.5mm,boxrule=0.9pt,left=3mm]{\Large\bfseries\color{cpro}a)}\quad\begin{minipage}[c]{5.6cm}\raggedright\Large $5$\end{minipage}\end{tcolorbox}\hfill\begin{tcolorbox}[enhanced,nobeforeafter,width=7.1cm,height=2.1cm,valign=center,colback=cpro!6,colframe=cpro,arc=2.5mm,boxrule=0.9pt,left=3mm]{\Large\bfseries\color{cpro}b)}\quad\begin{minipage}[c]{5.6cm}\raggedright\Large $4$\end{minipage}\end{tcolorbox}\par\vspace{2.5mm}
\begin{tcolorbox}[enhanced,nobeforeafter,width=7.1cm,height=2.1cm,valign=center,colback=cpro!6,colframe=cpro,arc=2.5mm,boxrule=0.9pt,left=3mm]{\Large\bfseries\color{cpro}c)}\quad\begin{minipage}[c]{5.6cm}\raggedright\Large $3$\end{minipage}\end{tcolorbox}\hfill\begin{tcolorbox}[enhanced,nobeforeafter,width=7.1cm,height=2.1cm,valign=center,colback=cpro!6,colframe=cpro,arc=2.5mm,boxrule=0.9pt,left=3mm]{\Large\bfseries\color{cpro}d)}\quad\begin{minipage}[c]{5.6cm}\raggedright\Large $\dfrac{1}{4}$\end{minipage}\end{tcolorbox}
\end{tcolorbox}
\newpage
\begin{tcolorbox}[enhanced,width=16cm,height=14cm,colback=white,colframe=cpro,boxrule=2pt,arc=0mm,
left=6mm,right=6mm,top=0mm,bottom=4mm,sharp corners,
overlay={\fill[cpro] (frame.north west) rectangle ([yshift=-1.25cm]frame.north east);
\node[anchor=west,text=white,font=\sffamily\bfseries\LARGE] at ([xshift=5mm,yshift=-0.62cm]frame.north west) {Question 91};
\node[anchor=east,text=white,font=\sffamily\itshape\large] at ([xshift=-5mm,yshift=-0.62cm]frame.north east) {Proportionnalité et données};}]
\vspace*{1.75cm}
\begin{minipage}[t][6.9cm][c]{\linewidth}\LARGE\centering
\begin{minipage}{0.95\linewidth}\raggedright\linespread{1.25}\selectfont Dans une classe, les élèves viennent au collège ainsi~: $8$ à pied, $12$ en bus et $5$ à vélo. Combien y a-t-il d'élèves dans la classe~?\end{minipage}
\end{minipage}

\vfill
\begin{tcolorbox}[enhanced,nobeforeafter,width=7.1cm,height=2.1cm,valign=center,colback=cpro!6,colframe=cpro,arc=2.5mm,boxrule=0.9pt,left=3mm]{\Large\bfseries\color{cpro}a)}\quad\begin{minipage}[c]{5.6cm}\raggedright\Large $25$\end{minipage}\end{tcolorbox}\hfill\begin{tcolorbox}[enhanced,nobeforeafter,width=7.1cm,height=2.1cm,valign=center,colback=cpro!6,colframe=cpro,arc=2.5mm,boxrule=0.9pt,left=3mm]{\Large\bfseries\color{cpro}b)}\quad\begin{minipage}[c]{5.6cm}\raggedright\Large $20$\end{minipage}\end{tcolorbox}\par\vspace{2.5mm}
\begin{tcolorbox}[enhanced,nobeforeafter,width=7.1cm,height=2.1cm,valign=center,colback=cpro!6,colframe=cpro,arc=2.5mm,boxrule=0.9pt,left=3mm]{\Large\bfseries\color{cpro}c)}\quad\begin{minipage}[c]{5.6cm}\raggedright\Large $12$\end{minipage}\end{tcolorbox}\hfill\begin{tcolorbox}[enhanced,nobeforeafter,width=7.1cm,height=2.1cm,valign=center,colback=cpro!6,colframe=cpro,arc=2.5mm,boxrule=0.9pt,left=3mm]{\Large\bfseries\color{cpro}d)}\quad\begin{minipage}[c]{5.6cm}\raggedright\Large $26$\end{minipage}\end{tcolorbox}
\end{tcolorbox}
\newpage
\begin{tcolorbox}[enhanced,width=16cm,height=14cm,colback=white,colframe=cpro,boxrule=2pt,arc=0mm,
left=6mm,right=6mm,top=0mm,bottom=4mm,sharp corners,
overlay={\fill[cpro] (frame.north west) rectangle ([yshift=-1.25cm]frame.north east);
\node[anchor=west,text=white,font=\sffamily\bfseries\LARGE] at ([xshift=5mm,yshift=-0.62cm]frame.north west) {Question 92};
\node[anchor=east,text=white,font=\sffamily\itshape\large] at ([xshift=-5mm,yshift=-0.62cm]frame.north east) {Proportionnalité et données};}]
\vspace*{1.75cm}
\begin{minipage}[t][6.9cm][c]{\linewidth}\LARGE\centering
\begin{minipage}{0.95\linewidth}\raggedright\linespread{1.25}\selectfont Dans une classe, les élèves viennent au collège ainsi~: $14$ à pied, $9$ en bus et $6$ à vélo. Combien y a-t-il d'élèves dans la classe~?\end{minipage}
\end{minipage}

\vfill
\begin{tcolorbox}[enhanced,nobeforeafter,width=7.1cm,height=2.1cm,valign=center,colback=cpro!6,colframe=cpro,arc=2.5mm,boxrule=0.9pt,left=3mm]{\Large\bfseries\color{cpro}a)}\quad\begin{minipage}[c]{5.6cm}\raggedright\Large $14$\end{minipage}\end{tcolorbox}\hfill\begin{tcolorbox}[enhanced,nobeforeafter,width=7.1cm,height=2.1cm,valign=center,colback=cpro!6,colframe=cpro,arc=2.5mm,boxrule=0.9pt,left=3mm]{\Large\bfseries\color{cpro}b)}\quad\begin{minipage}[c]{5.6cm}\raggedright\Large $29$\end{minipage}\end{tcolorbox}\par\vspace{2.5mm}
\begin{tcolorbox}[enhanced,nobeforeafter,width=7.1cm,height=2.1cm,valign=center,colback=cpro!6,colframe=cpro,arc=2.5mm,boxrule=0.9pt,left=3mm]{\Large\bfseries\color{cpro}c)}\quad\begin{minipage}[c]{5.6cm}\raggedright\Large $30$\end{minipage}\end{tcolorbox}\hfill\begin{tcolorbox}[enhanced,nobeforeafter,width=7.1cm,height=2.1cm,valign=center,colback=cpro!6,colframe=cpro,arc=2.5mm,boxrule=0.9pt,left=3mm]{\Large\bfseries\color{cpro}d)}\quad\begin{minipage}[c]{5.6cm}\raggedright\Large $23$\end{minipage}\end{tcolorbox}
\end{tcolorbox}
\newpage
\begin{tcolorbox}[enhanced,width=16cm,height=14cm,colback=white,colframe=cpro,boxrule=2pt,arc=0mm,
left=6mm,right=6mm,top=0mm,bottom=4mm,sharp corners,
overlay={\fill[cpro] (frame.north west) rectangle ([yshift=-1.25cm]frame.north east);
\node[anchor=west,text=white,font=\sffamily\bfseries\LARGE] at ([xshift=5mm,yshift=-0.62cm]frame.north west) {Question 93};
\node[anchor=east,text=white,font=\sffamily\itshape\large] at ([xshift=-5mm,yshift=-0.62cm]frame.north east) {Proportionnalité et données};}]
\vspace*{1.75cm}
\begin{minipage}[t][6.9cm][c]{\linewidth}\LARGE\centering
\begin{minipage}{0.95\linewidth}\raggedright\linespread{1.25}\selectfont Ce tableau est-il un tableau de proportionnalité~? $\begin{array}{|c|c|}\hline 2 & 5\\\hline 6 & 15\\\hline\end{array}$\end{minipage}
\end{minipage}

\vfill
\begin{tcolorbox}[enhanced,nobeforeafter,width=7.1cm,height=2.1cm,valign=center,colback=cpro!6,colframe=cpro,arc=2.5mm,boxrule=0.9pt,left=3mm]{\Large\bfseries\color{cpro}a)}\quad\begin{minipage}[c]{5.6cm}\raggedright\Large non\end{minipage}\end{tcolorbox}\hfill\begin{tcolorbox}[enhanced,nobeforeafter,width=7.1cm,height=2.1cm,valign=center,colback=cpro!6,colframe=cpro,arc=2.5mm,boxrule=0.9pt,left=3mm]{\Large\bfseries\color{cpro}b)}\quad\begin{minipage}[c]{5.6cm}\raggedright\Large oui\end{minipage}\end{tcolorbox}\par\vspace{2.5mm}
\begin{tcolorbox}[enhanced,nobeforeafter,width=7.1cm,height=2.1cm,valign=center,colback=cpro!6,colframe=cpro,arc=2.5mm,boxrule=0.9pt,left=3mm]{\Large\bfseries\color{cpro}c)}\quad\begin{minipage}[c]{5.6cm}\raggedright\Large il manque une colonne\end{minipage}\end{tcolorbox}\hfill\begin{tcolorbox}[enhanced,nobeforeafter,width=7.1cm,height=2.1cm,valign=center,colback=cpro!6,colframe=cpro,arc=2.5mm,boxrule=0.9pt,left=3mm]{\Large\bfseries\color{cpro}d)}\quad\begin{minipage}[c]{5.6cm}\raggedright\Large on ne peut pas savoir\end{minipage}\end{tcolorbox}
\end{tcolorbox}
\newpage
\begin{tcolorbox}[enhanced,width=16cm,height=14cm,colback=white,colframe=cpro,boxrule=2pt,arc=0mm,
left=6mm,right=6mm,top=0mm,bottom=4mm,sharp corners,
overlay={\fill[cpro] (frame.north west) rectangle ([yshift=-1.25cm]frame.north east);
\node[anchor=west,text=white,font=\sffamily\bfseries\LARGE] at ([xshift=5mm,yshift=-0.62cm]frame.north west) {Question 94};
\node[anchor=east,text=white,font=\sffamily\itshape\large] at ([xshift=-5mm,yshift=-0.62cm]frame.north east) {Proportionnalité et données};}]
\vspace*{1.75cm}
\begin{minipage}[t][6.9cm][c]{\linewidth}\LARGE\centering
\begin{minipage}{0.95\linewidth}\raggedright\linespread{1.25}\selectfont Ce tableau est-il un tableau de proportionnalité~? $\begin{array}{|c|c|}\hline 3 & 4\\\hline 9 & 16\\\hline\end{array}$\end{minipage}
\end{minipage}

\vfill
\begin{tcolorbox}[enhanced,nobeforeafter,width=7.1cm,height=2.1cm,valign=center,colback=cpro!6,colframe=cpro,arc=2.5mm,boxrule=0.9pt,left=3mm]{\Large\bfseries\color{cpro}a)}\quad\begin{minipage}[c]{5.6cm}\raggedright\Large non\end{minipage}\end{tcolorbox}\hfill\begin{tcolorbox}[enhanced,nobeforeafter,width=7.1cm,height=2.1cm,valign=center,colback=cpro!6,colframe=cpro,arc=2.5mm,boxrule=0.9pt,left=3mm]{\Large\bfseries\color{cpro}b)}\quad\begin{minipage}[c]{5.6cm}\raggedright\Large il manque une colonne\end{minipage}\end{tcolorbox}\par\vspace{2.5mm}
\begin{tcolorbox}[enhanced,nobeforeafter,width=7.1cm,height=2.1cm,valign=center,colback=cpro!6,colframe=cpro,arc=2.5mm,boxrule=0.9pt,left=3mm]{\Large\bfseries\color{cpro}c)}\quad\begin{minipage}[c]{5.6cm}\raggedright\Large oui\end{minipage}\end{tcolorbox}\hfill\begin{tcolorbox}[enhanced,nobeforeafter,width=7.1cm,height=2.1cm,valign=center,colback=cpro!6,colframe=cpro,arc=2.5mm,boxrule=0.9pt,left=3mm]{\Large\bfseries\color{cpro}d)}\quad\begin{minipage}[c]{5.6cm}\raggedright\Large on ne peut pas savoir\end{minipage}\end{tcolorbox}
\end{tcolorbox}
\newpage
\begin{tcolorbox}[enhanced,width=16cm,height=14cm,colback=white,colframe=cpro,boxrule=2pt,arc=0mm,
left=6mm,right=6mm,top=0mm,bottom=4mm,sharp corners,
overlay={\fill[cpro] (frame.north west) rectangle ([yshift=-1.25cm]frame.north east);
\node[anchor=west,text=white,font=\sffamily\bfseries\LARGE] at ([xshift=5mm,yshift=-0.62cm]frame.north west) {Question 95};
\node[anchor=east,text=white,font=\sffamily\itshape\large] at ([xshift=-5mm,yshift=-0.62cm]frame.north east) {Proportionnalité et données};}]
\vspace*{1.75cm}
\begin{minipage}[t][6.9cm][c]{\linewidth}\LARGE\centering
\begin{minipage}{0.95\linewidth}\raggedright\linespread{1.25}\selectfont Un robinet remplit $12$ L en $3$ min. Quel volume remplit-il en $5$ min, au même débit~?\end{minipage}
\end{minipage}

\vfill
\begin{tcolorbox}[enhanced,nobeforeafter,width=7.1cm,height=2.1cm,valign=center,colback=cpro!6,colframe=cpro,arc=2.5mm,boxrule=0.9pt,left=3mm]{\Large\bfseries\color{cpro}a)}\quad\begin{minipage}[c]{5.6cm}\raggedright\Large $60~\text{L}$\end{minipage}\end{tcolorbox}\hfill\begin{tcolorbox}[enhanced,nobeforeafter,width=7.1cm,height=2.1cm,valign=center,colback=cpro!6,colframe=cpro,arc=2.5mm,boxrule=0.9pt,left=3mm]{\Large\bfseries\color{cpro}b)}\quad\begin{minipage}[c]{5.6cm}\raggedright\Large $14~\text{L}$\end{minipage}\end{tcolorbox}\par\vspace{2.5mm}
\begin{tcolorbox}[enhanced,nobeforeafter,width=7.1cm,height=2.1cm,valign=center,colback=cpro!6,colframe=cpro,arc=2.5mm,boxrule=0.9pt,left=3mm]{\Large\bfseries\color{cpro}c)}\quad\begin{minipage}[c]{5.6cm}\raggedright\Large $4~\text{L}$\end{minipage}\end{tcolorbox}\hfill\begin{tcolorbox}[enhanced,nobeforeafter,width=7.1cm,height=2.1cm,valign=center,colback=cpro!6,colframe=cpro,arc=2.5mm,boxrule=0.9pt,left=3mm]{\Large\bfseries\color{cpro}d)}\quad\begin{minipage}[c]{5.6cm}\raggedright\Large $20~\text{L}$\end{minipage}\end{tcolorbox}
\end{tcolorbox}
\newpage
\begin{tcolorbox}[enhanced,width=16cm,height=14cm,colback=white,colframe=cpro,boxrule=2pt,arc=0mm,
left=6mm,right=6mm,top=0mm,bottom=4mm,sharp corners,
overlay={\fill[cpro] (frame.north west) rectangle ([yshift=-1.25cm]frame.north east);
\node[anchor=west,text=white,font=\sffamily\bfseries\LARGE] at ([xshift=5mm,yshift=-0.62cm]frame.north west) {Question 96};
\node[anchor=east,text=white,font=\sffamily\itshape\large] at ([xshift=-5mm,yshift=-0.62cm]frame.north east) {Proportionnalité et données};}]
\vspace*{1.75cm}
\begin{minipage}[t][6.9cm][c]{\linewidth}\LARGE\centering
\begin{minipage}{0.95\linewidth}\raggedright\linespread{1.25}\selectfont Un robinet remplit $25$ L en $5$ min. Quel volume remplit-il en $8$ min, au même débit~?\end{minipage}
\end{minipage}

\vfill
\begin{tcolorbox}[enhanced,nobeforeafter,width=7.1cm,height=2.1cm,valign=center,colback=cpro!6,colframe=cpro,arc=2.5mm,boxrule=0.9pt,left=3mm]{\Large\bfseries\color{cpro}a)}\quad\begin{minipage}[c]{5.6cm}\raggedright\Large $200~\text{L}$\end{minipage}\end{tcolorbox}\hfill\begin{tcolorbox}[enhanced,nobeforeafter,width=7.1cm,height=2.1cm,valign=center,colback=cpro!6,colframe=cpro,arc=2.5mm,boxrule=0.9pt,left=3mm]{\Large\bfseries\color{cpro}b)}\quad\begin{minipage}[c]{5.6cm}\raggedright\Large $40~\text{L}$\end{minipage}\end{tcolorbox}\par\vspace{2.5mm}
\begin{tcolorbox}[enhanced,nobeforeafter,width=7.1cm,height=2.1cm,valign=center,colback=cpro!6,colframe=cpro,arc=2.5mm,boxrule=0.9pt,left=3mm]{\Large\bfseries\color{cpro}c)}\quad\begin{minipage}[c]{5.6cm}\raggedright\Large $5~\text{L}$\end{minipage}\end{tcolorbox}\hfill\begin{tcolorbox}[enhanced,nobeforeafter,width=7.1cm,height=2.1cm,valign=center,colback=cpro!6,colframe=cpro,arc=2.5mm,boxrule=0.9pt,left=3mm]{\Large\bfseries\color{cpro}d)}\quad\begin{minipage}[c]{5.6cm}\raggedright\Large $28~\text{L}$\end{minipage}\end{tcolorbox}
\end{tcolorbox}
\newpage
\begin{tcolorbox}[enhanced,width=16cm,height=14cm,colback=white,colframe=cnum,boxrule=2pt,arc=0mm,
left=6mm,right=6mm,top=0mm,bottom=4mm,sharp corners,
overlay={\fill[cnum] (frame.north west) rectangle ([yshift=-1.25cm]frame.north east);
\node[anchor=west,text=white,font=\sffamily\bfseries\LARGE] at ([xshift=5mm,yshift=-0.62cm]frame.north west) {Question 1 B};
\node[anchor=east,text=white,font=\sffamily\itshape\large] at ([xshift=-5mm,yshift=-0.62cm]frame.north east) {Nombres entiers et décimaux};}]
\vspace*{1.75cm}
\begin{minipage}[t][6.9cm][c]{\linewidth}\LARGE\centering
\begin{minipage}{0.95\linewidth}\raggedright\linespread{1.25}\selectfont Dans le nombre $67{,}094$, quel est le chiffre des millièmes~?\end{minipage}
\end{minipage}

\vfill
\begin{tcolorbox}[enhanced,nobeforeafter,width=7.1cm,height=2.1cm,valign=center,colback=cnum!6,colframe=cnum,arc=2.5mm,boxrule=0.9pt,left=3mm]{\Large\bfseries\color{cnum}a)}\quad\begin{minipage}[c]{5.6cm}\raggedright\Large $6$\end{minipage}\end{tcolorbox}\hfill\begin{tcolorbox}[enhanced,nobeforeafter,width=7.1cm,height=2.1cm,valign=center,colback=cnum!6,colframe=cnum,arc=2.5mm,boxrule=0.9pt,left=3mm]{\Large\bfseries\color{cnum}b)}\quad\begin{minipage}[c]{5.6cm}\raggedright\Large $7$\end{minipage}\end{tcolorbox}\par\vspace{2.5mm}
\begin{tcolorbox}[enhanced,nobeforeafter,width=7.1cm,height=2.1cm,valign=center,colback=cnum!6,colframe=cnum,arc=2.5mm,boxrule=0.9pt,left=3mm]{\Large\bfseries\color{cnum}c)}\quad\begin{minipage}[c]{5.6cm}\raggedright\Large $0$\end{minipage}\end{tcolorbox}\hfill\begin{tcolorbox}[enhanced,nobeforeafter,width=7.1cm,height=2.1cm,valign=center,colback=cnum!6,colframe=cnum,arc=2.5mm,boxrule=0.9pt,left=3mm]{\Large\bfseries\color{cnum}d)}\quad\begin{minipage}[c]{5.6cm}\raggedright\Large $4$\end{minipage}\end{tcolorbox}
\end{tcolorbox}
\newpage
\begin{tcolorbox}[enhanced,width=16cm,height=14cm,colback=white,colframe=cnum,boxrule=2pt,arc=0mm,
left=6mm,right=6mm,top=0mm,bottom=4mm,sharp corners,
overlay={\fill[cnum] (frame.north west) rectangle ([yshift=-1.25cm]frame.north east);
\node[anchor=west,text=white,font=\sffamily\bfseries\LARGE] at ([xshift=5mm,yshift=-0.62cm]frame.north west) {Question 2 B};
\node[anchor=east,text=white,font=\sffamily\itshape\large] at ([xshift=-5mm,yshift=-0.62cm]frame.north east) {Nombres entiers et décimaux};}]
\vspace*{1.75cm}
\begin{minipage}[t][6.9cm][c]{\linewidth}\LARGE\centering
\begin{minipage}{0.95\linewidth}\raggedright\linespread{1.25}\selectfont Dans le nombre $5\,431{,}8$, quel est le chiffre des centaines~?\end{minipage}
\end{minipage}

\vfill
\begin{tcolorbox}[enhanced,nobeforeafter,width=7.1cm,height=2.1cm,valign=center,colback=cnum!6,colframe=cnum,arc=2.5mm,boxrule=0.9pt,left=3mm]{\Large\bfseries\color{cnum}a)}\quad\begin{minipage}[c]{5.6cm}\raggedright\Large $4$\end{minipage}\end{tcolorbox}\hfill\begin{tcolorbox}[enhanced,nobeforeafter,width=7.1cm,height=2.1cm,valign=center,colback=cnum!6,colframe=cnum,arc=2.5mm,boxrule=0.9pt,left=3mm]{\Large\bfseries\color{cnum}b)}\quad\begin{minipage}[c]{5.6cm}\raggedright\Large $3$\end{minipage}\end{tcolorbox}\par\vspace{2.5mm}
\begin{tcolorbox}[enhanced,nobeforeafter,width=7.1cm,height=2.1cm,valign=center,colback=cnum!6,colframe=cnum,arc=2.5mm,boxrule=0.9pt,left=3mm]{\Large\bfseries\color{cnum}c)}\quad\begin{minipage}[c]{5.6cm}\raggedright\Large $5$\end{minipage}\end{tcolorbox}\hfill\begin{tcolorbox}[enhanced,nobeforeafter,width=7.1cm,height=2.1cm,valign=center,colback=cnum!6,colframe=cnum,arc=2.5mm,boxrule=0.9pt,left=3mm]{\Large\bfseries\color{cnum}d)}\quad\begin{minipage}[c]{5.6cm}\raggedright\Large $1$\end{minipage}\end{tcolorbox}
\end{tcolorbox}
\newpage
\begin{tcolorbox}[enhanced,width=16cm,height=14cm,colback=white,colframe=cnum,boxrule=2pt,arc=0mm,
left=6mm,right=6mm,top=0mm,bottom=4mm,sharp corners,
overlay={\fill[cnum] (frame.north west) rectangle ([yshift=-1.25cm]frame.north east);
\node[anchor=west,text=white,font=\sffamily\bfseries\LARGE] at ([xshift=5mm,yshift=-0.62cm]frame.north west) {Question 3 B};
\node[anchor=east,text=white,font=\sffamily\itshape\large] at ([xshift=-5mm,yshift=-0.62cm]frame.north east) {Nombres entiers et décimaux};}]
\vspace*{1.75cm}
\begin{minipage}[t][6.9cm][c]{\linewidth}\LARGE\centering
\begin{minipage}{0.95\linewidth}\raggedright\linespread{1.25}\selectfont Quelle est l'écriture décimale de $\dfrac{8}{100}$~?\end{minipage}
\end{minipage}

\vfill
\begin{tcolorbox}[enhanced,nobeforeafter,width=7.1cm,height=2.1cm,valign=center,colback=cnum!6,colframe=cnum,arc=2.5mm,boxrule=0.9pt,left=3mm]{\Large\bfseries\color{cnum}a)}\quad\begin{minipage}[c]{5.6cm}\raggedright\Large $0{,}8$\end{minipage}\end{tcolorbox}\hfill\begin{tcolorbox}[enhanced,nobeforeafter,width=7.1cm,height=2.1cm,valign=center,colback=cnum!6,colframe=cnum,arc=2.5mm,boxrule=0.9pt,left=3mm]{\Large\bfseries\color{cnum}b)}\quad\begin{minipage}[c]{5.6cm}\raggedright\Large $0{,}008$\end{minipage}\end{tcolorbox}\par\vspace{2.5mm}
\begin{tcolorbox}[enhanced,nobeforeafter,width=7.1cm,height=2.1cm,valign=center,colback=cnum!6,colframe=cnum,arc=2.5mm,boxrule=0.9pt,left=3mm]{\Large\bfseries\color{cnum}c)}\quad\begin{minipage}[c]{5.6cm}\raggedright\Large $0{,}08$\end{minipage}\end{tcolorbox}\hfill\begin{tcolorbox}[enhanced,nobeforeafter,width=7.1cm,height=2.1cm,valign=center,colback=cnum!6,colframe=cnum,arc=2.5mm,boxrule=0.9pt,left=3mm]{\Large\bfseries\color{cnum}d)}\quad\begin{minipage}[c]{5.6cm}\raggedright\Large $8$\end{minipage}\end{tcolorbox}
\end{tcolorbox}
\newpage
\begin{tcolorbox}[enhanced,width=16cm,height=14cm,colback=white,colframe=cnum,boxrule=2pt,arc=0mm,
left=6mm,right=6mm,top=0mm,bottom=4mm,sharp corners,
overlay={\fill[cnum] (frame.north west) rectangle ([yshift=-1.25cm]frame.north east);
\node[anchor=west,text=white,font=\sffamily\bfseries\LARGE] at ([xshift=5mm,yshift=-0.62cm]frame.north west) {Question 4 B};
\node[anchor=east,text=white,font=\sffamily\itshape\large] at ([xshift=-5mm,yshift=-0.62cm]frame.north east) {Nombres entiers et décimaux};}]
\vspace*{1.75cm}
\begin{minipage}[t][6.9cm][c]{\linewidth}\LARGE\centering
\begin{minipage}{0.95\linewidth}\raggedright\linespread{1.25}\selectfont Quelle est l'écriture décimale de $\dfrac{1205}{1000}$~?\end{minipage}
\end{minipage}

\vfill
\begin{tcolorbox}[enhanced,nobeforeafter,width=7.1cm,height=2.1cm,valign=center,colback=cnum!6,colframe=cnum,arc=2.5mm,boxrule=0.9pt,left=3mm]{\Large\bfseries\color{cnum}a)}\quad\begin{minipage}[c]{5.6cm}\raggedright\Large $120{,}5$\end{minipage}\end{tcolorbox}\hfill\begin{tcolorbox}[enhanced,nobeforeafter,width=7.1cm,height=2.1cm,valign=center,colback=cnum!6,colframe=cnum,arc=2.5mm,boxrule=0.9pt,left=3mm]{\Large\bfseries\color{cnum}b)}\quad\begin{minipage}[c]{5.6cm}\raggedright\Large $1{,}205$\end{minipage}\end{tcolorbox}\par\vspace{2.5mm}
\begin{tcolorbox}[enhanced,nobeforeafter,width=7.1cm,height=2.1cm,valign=center,colback=cnum!6,colframe=cnum,arc=2.5mm,boxrule=0.9pt,left=3mm]{\Large\bfseries\color{cnum}c)}\quad\begin{minipage}[c]{5.6cm}\raggedright\Large $0{,}1205$\end{minipage}\end{tcolorbox}\hfill\begin{tcolorbox}[enhanced,nobeforeafter,width=7.1cm,height=2.1cm,valign=center,colback=cnum!6,colframe=cnum,arc=2.5mm,boxrule=0.9pt,left=3mm]{\Large\bfseries\color{cnum}d)}\quad\begin{minipage}[c]{5.6cm}\raggedright\Large $12{,}05$\end{minipage}\end{tcolorbox}
\end{tcolorbox}
\newpage
\begin{tcolorbox}[enhanced,width=16cm,height=14cm,colback=white,colframe=cnum,boxrule=2pt,arc=0mm,
left=6mm,right=6mm,top=0mm,bottom=4mm,sharp corners,
overlay={\fill[cnum] (frame.north west) rectangle ([yshift=-1.25cm]frame.north east);
\node[anchor=west,text=white,font=\sffamily\bfseries\LARGE] at ([xshift=5mm,yshift=-0.62cm]frame.north west) {Question 5 B};
\node[anchor=east,text=white,font=\sffamily\itshape\large] at ([xshift=-5mm,yshift=-0.62cm]frame.north east) {Nombres entiers et décimaux};}]
\vspace*{1.75cm}
\begin{minipage}[t][6.9cm][c]{\linewidth}\LARGE\centering
\begin{minipage}{0.95\linewidth}\raggedright\linespread{1.25}\selectfont Quel est le plus grand de ces quatre nombres~: $0{,}7\ ;\ 0{,}69\ ;\ 0{,}075\ ;\ 0{,}701$~?\end{minipage}
\end{minipage}

\vfill
\begin{tcolorbox}[enhanced,nobeforeafter,width=7.1cm,height=2.1cm,valign=center,colback=cnum!6,colframe=cnum,arc=2.5mm,boxrule=0.9pt,left=3mm]{\Large\bfseries\color{cnum}a)}\quad\begin{minipage}[c]{5.6cm}\raggedright\Large $0{,}701$\end{minipage}\end{tcolorbox}\hfill\begin{tcolorbox}[enhanced,nobeforeafter,width=7.1cm,height=2.1cm,valign=center,colback=cnum!6,colframe=cnum,arc=2.5mm,boxrule=0.9pt,left=3mm]{\Large\bfseries\color{cnum}b)}\quad\begin{minipage}[c]{5.6cm}\raggedright\Large $0{,}69$\end{minipage}\end{tcolorbox}\par\vspace{2.5mm}
\begin{tcolorbox}[enhanced,nobeforeafter,width=7.1cm,height=2.1cm,valign=center,colback=cnum!6,colframe=cnum,arc=2.5mm,boxrule=0.9pt,left=3mm]{\Large\bfseries\color{cnum}c)}\quad\begin{minipage}[c]{5.6cm}\raggedright\Large $0{,}075$\end{minipage}\end{tcolorbox}\hfill\begin{tcolorbox}[enhanced,nobeforeafter,width=7.1cm,height=2.1cm,valign=center,colback=cnum!6,colframe=cnum,arc=2.5mm,boxrule=0.9pt,left=3mm]{\Large\bfseries\color{cnum}d)}\quad\begin{minipage}[c]{5.6cm}\raggedright\Large $0{,}7$\end{minipage}\end{tcolorbox}
\end{tcolorbox}
\newpage
\begin{tcolorbox}[enhanced,width=16cm,height=14cm,colback=white,colframe=cnum,boxrule=2pt,arc=0mm,
left=6mm,right=6mm,top=0mm,bottom=4mm,sharp corners,
overlay={\fill[cnum] (frame.north west) rectangle ([yshift=-1.25cm]frame.north east);
\node[anchor=west,text=white,font=\sffamily\bfseries\LARGE] at ([xshift=5mm,yshift=-0.62cm]frame.north west) {Question 6 B};
\node[anchor=east,text=white,font=\sffamily\itshape\large] at ([xshift=-5mm,yshift=-0.62cm]frame.north east) {Nombres entiers et décimaux};}]
\vspace*{1.75cm}
\begin{minipage}[t][6.9cm][c]{\linewidth}\LARGE\centering
\begin{minipage}{0.95\linewidth}\raggedright\linespread{1.25}\selectfont Quel est le plus petit de ces quatre nombres~: $5{,}03\ ;\ 5{,}3\ ;\ 5{,}003\ ;\ 5{,}29$~?\end{minipage}
\end{minipage}

\vfill
\begin{tcolorbox}[enhanced,nobeforeafter,width=7.1cm,height=2.1cm,valign=center,colback=cnum!6,colframe=cnum,arc=2.5mm,boxrule=0.9pt,left=3mm]{\Large\bfseries\color{cnum}a)}\quad\begin{minipage}[c]{5.6cm}\raggedright\Large $5{,}3$\end{minipage}\end{tcolorbox}\hfill\begin{tcolorbox}[enhanced,nobeforeafter,width=7.1cm,height=2.1cm,valign=center,colback=cnum!6,colframe=cnum,arc=2.5mm,boxrule=0.9pt,left=3mm]{\Large\bfseries\color{cnum}b)}\quad\begin{minipage}[c]{5.6cm}\raggedright\Large $5{,}29$\end{minipage}\end{tcolorbox}\par\vspace{2.5mm}
\begin{tcolorbox}[enhanced,nobeforeafter,width=7.1cm,height=2.1cm,valign=center,colback=cnum!6,colframe=cnum,arc=2.5mm,boxrule=0.9pt,left=3mm]{\Large\bfseries\color{cnum}c)}\quad\begin{minipage}[c]{5.6cm}\raggedright\Large $5{,}03$\end{minipage}\end{tcolorbox}\hfill\begin{tcolorbox}[enhanced,nobeforeafter,width=7.1cm,height=2.1cm,valign=center,colback=cnum!6,colframe=cnum,arc=2.5mm,boxrule=0.9pt,left=3mm]{\Large\bfseries\color{cnum}d)}\quad\begin{minipage}[c]{5.6cm}\raggedright\Large $5{,}003$\end{minipage}\end{tcolorbox}
\end{tcolorbox}
\newpage
\begin{tcolorbox}[enhanced,width=16cm,height=14cm,colback=white,colframe=cnum,boxrule=2pt,arc=0mm,
left=6mm,right=6mm,top=0mm,bottom=4mm,sharp corners,
overlay={\fill[cnum] (frame.north west) rectangle ([yshift=-1.25cm]frame.north east);
\node[anchor=west,text=white,font=\sffamily\bfseries\LARGE] at ([xshift=5mm,yshift=-0.62cm]frame.north west) {Question 7 B};
\node[anchor=east,text=white,font=\sffamily\itshape\large] at ([xshift=-5mm,yshift=-0.62cm]frame.north east) {Nombres entiers et décimaux};}]
\vspace*{1.75cm}
\begin{minipage}[t][6.9cm][c]{\linewidth}\LARGE\centering
\begin{minipage}{0.95\linewidth}\raggedright\linespread{1.25}\selectfont Quel est l'arrondi au dixième de $0{,}95$~?\end{minipage}
\end{minipage}

\vfill
\begin{tcolorbox}[enhanced,nobeforeafter,width=7.1cm,height=2.1cm,valign=center,colback=cnum!6,colframe=cnum,arc=2.5mm,boxrule=0.9pt,left=3mm]{\Large\bfseries\color{cnum}a)}\quad\begin{minipage}[c]{5.6cm}\raggedright\Large $1$\end{minipage}\end{tcolorbox}\hfill\begin{tcolorbox}[enhanced,nobeforeafter,width=7.1cm,height=2.1cm,valign=center,colback=cnum!6,colframe=cnum,arc=2.5mm,boxrule=0.9pt,left=3mm]{\Large\bfseries\color{cnum}b)}\quad\begin{minipage}[c]{5.6cm}\raggedright\Large $1{,}2$\end{minipage}\end{tcolorbox}\par\vspace{2.5mm}
\begin{tcolorbox}[enhanced,nobeforeafter,width=7.1cm,height=2.1cm,valign=center,colback=cnum!6,colframe=cnum,arc=2.5mm,boxrule=0.9pt,left=3mm]{\Large\bfseries\color{cnum}c)}\quad\begin{minipage}[c]{5.6cm}\raggedright\Large $0{,}9$\end{minipage}\end{tcolorbox}\hfill\begin{tcolorbox}[enhanced,nobeforeafter,width=7.1cm,height=2.1cm,valign=center,colback=cnum!6,colframe=cnum,arc=2.5mm,boxrule=0.9pt,left=3mm]{\Large\bfseries\color{cnum}d)}\quad\begin{minipage}[c]{5.6cm}\raggedright\Large $0{,}95$\end{minipage}\end{tcolorbox}
\end{tcolorbox}
\newpage
\begin{tcolorbox}[enhanced,width=16cm,height=14cm,colback=white,colframe=cnum,boxrule=2pt,arc=0mm,
left=6mm,right=6mm,top=0mm,bottom=4mm,sharp corners,
overlay={\fill[cnum] (frame.north west) rectangle ([yshift=-1.25cm]frame.north east);
\node[anchor=west,text=white,font=\sffamily\bfseries\LARGE] at ([xshift=5mm,yshift=-0.62cm]frame.north west) {Question 8 B};
\node[anchor=east,text=white,font=\sffamily\itshape\large] at ([xshift=-5mm,yshift=-0.62cm]frame.north east) {Nombres entiers et décimaux};}]
\vspace*{1.75cm}
\begin{minipage}[t][6.9cm][c]{\linewidth}\LARGE\centering
\begin{minipage}{0.95\linewidth}\raggedright\linespread{1.25}\selectfont Quel est l'arrondi au dixième de $23{,}149$~?\end{minipage}
\end{minipage}

\vfill
\begin{tcolorbox}[enhanced,nobeforeafter,width=7.1cm,height=2.1cm,valign=center,colback=cnum!6,colframe=cnum,arc=2.5mm,boxrule=0.9pt,left=3mm]{\Large\bfseries\color{cnum}a)}\quad\begin{minipage}[c]{5.6cm}\raggedright\Large $23$\end{minipage}\end{tcolorbox}\hfill\begin{tcolorbox}[enhanced,nobeforeafter,width=7.1cm,height=2.1cm,valign=center,colback=cnum!6,colframe=cnum,arc=2.5mm,boxrule=0.9pt,left=3mm]{\Large\bfseries\color{cnum}b)}\quad\begin{minipage}[c]{5.6cm}\raggedright\Large $23{,}1$\end{minipage}\end{tcolorbox}\par\vspace{2.5mm}
\begin{tcolorbox}[enhanced,nobeforeafter,width=7.1cm,height=2.1cm,valign=center,colback=cnum!6,colframe=cnum,arc=2.5mm,boxrule=0.9pt,left=3mm]{\Large\bfseries\color{cnum}c)}\quad\begin{minipage}[c]{5.6cm}\raggedright\Large $23{,}15$\end{minipage}\end{tcolorbox}\hfill\begin{tcolorbox}[enhanced,nobeforeafter,width=7.1cm,height=2.1cm,valign=center,colback=cnum!6,colframe=cnum,arc=2.5mm,boxrule=0.9pt,left=3mm]{\Large\bfseries\color{cnum}d)}\quad\begin{minipage}[c]{5.6cm}\raggedright\Large $23{,}2$\end{minipage}\end{tcolorbox}
\end{tcolorbox}
\newpage
\begin{tcolorbox}[enhanced,width=16cm,height=14cm,colback=white,colframe=cnum,boxrule=2pt,arc=0mm,
left=6mm,right=6mm,top=0mm,bottom=4mm,sharp corners,
overlay={\fill[cnum] (frame.north west) rectangle ([yshift=-1.25cm]frame.north east);
\node[anchor=west,text=white,font=\sffamily\bfseries\LARGE] at ([xshift=5mm,yshift=-0.62cm]frame.north west) {Question 9 B};
\node[anchor=east,text=white,font=\sffamily\itshape\large] at ([xshift=-5mm,yshift=-0.62cm]frame.north east) {Nombres entiers et décimaux};}]
\vspace*{1.75cm}
\begin{minipage}[t][6.9cm][c]{\linewidth}\LARGE\centering
\begin{minipage}{0.95\linewidth}\raggedright\linespread{1.25}\selectfont Que vaut $5{,}3\times1\,000$~?\end{minipage}
\end{minipage}

\vfill
\begin{tcolorbox}[enhanced,nobeforeafter,width=7.1cm,height=2.1cm,valign=center,colback=cnum!6,colframe=cnum,arc=2.5mm,boxrule=0.9pt,left=3mm]{\Large\bfseries\color{cnum}a)}\quad\begin{minipage}[c]{5.6cm}\raggedright\Large $530\,000$\end{minipage}\end{tcolorbox}\hfill\begin{tcolorbox}[enhanced,nobeforeafter,width=7.1cm,height=2.1cm,valign=center,colback=cnum!6,colframe=cnum,arc=2.5mm,boxrule=0.9pt,left=3mm]{\Large\bfseries\color{cnum}b)}\quad\begin{minipage}[c]{5.6cm}\raggedright\Large $53\,000$\end{minipage}\end{tcolorbox}\par\vspace{2.5mm}
\begin{tcolorbox}[enhanced,nobeforeafter,width=7.1cm,height=2.1cm,valign=center,colback=cnum!6,colframe=cnum,arc=2.5mm,boxrule=0.9pt,left=3mm]{\Large\bfseries\color{cnum}c)}\quad\begin{minipage}[c]{5.6cm}\raggedright\Large $530$\end{minipage}\end{tcolorbox}\hfill\begin{tcolorbox}[enhanced,nobeforeafter,width=7.1cm,height=2.1cm,valign=center,colback=cnum!6,colframe=cnum,arc=2.5mm,boxrule=0.9pt,left=3mm]{\Large\bfseries\color{cnum}d)}\quad\begin{minipage}[c]{5.6cm}\raggedright\Large $5\,300$\end{minipage}\end{tcolorbox}
\end{tcolorbox}
\newpage
\begin{tcolorbox}[enhanced,width=16cm,height=14cm,colback=white,colframe=cnum,boxrule=2pt,arc=0mm,
left=6mm,right=6mm,top=0mm,bottom=4mm,sharp corners,
overlay={\fill[cnum] (frame.north west) rectangle ([yshift=-1.25cm]frame.north east);
\node[anchor=west,text=white,font=\sffamily\bfseries\LARGE] at ([xshift=5mm,yshift=-0.62cm]frame.north west) {Question 10 B};
\node[anchor=east,text=white,font=\sffamily\itshape\large] at ([xshift=-5mm,yshift=-0.62cm]frame.north east) {Nombres entiers et décimaux};}]
\vspace*{1.75cm}
\begin{minipage}[t][6.9cm][c]{\linewidth}\LARGE\centering
\begin{minipage}{0.95\linewidth}\raggedright\linespread{1.25}\selectfont Que vaut $0{,}068\times100$~?\end{minipage}
\end{minipage}

\vfill
\begin{tcolorbox}[enhanced,nobeforeafter,width=7.1cm,height=2.1cm,valign=center,colback=cnum!6,colframe=cnum,arc=2.5mm,boxrule=0.9pt,left=3mm]{\Large\bfseries\color{cnum}a)}\quad\begin{minipage}[c]{5.6cm}\raggedright\Large $680$\end{minipage}\end{tcolorbox}\hfill\begin{tcolorbox}[enhanced,nobeforeafter,width=7.1cm,height=2.1cm,valign=center,colback=cnum!6,colframe=cnum,arc=2.5mm,boxrule=0.9pt,left=3mm]{\Large\bfseries\color{cnum}b)}\quad\begin{minipage}[c]{5.6cm}\raggedright\Large $68$\end{minipage}\end{tcolorbox}\par\vspace{2.5mm}
\begin{tcolorbox}[enhanced,nobeforeafter,width=7.1cm,height=2.1cm,valign=center,colback=cnum!6,colframe=cnum,arc=2.5mm,boxrule=0.9pt,left=3mm]{\Large\bfseries\color{cnum}c)}\quad\begin{minipage}[c]{5.6cm}\raggedright\Large $0{,}68$\end{minipage}\end{tcolorbox}\hfill\begin{tcolorbox}[enhanced,nobeforeafter,width=7.1cm,height=2.1cm,valign=center,colback=cnum!6,colframe=cnum,arc=2.5mm,boxrule=0.9pt,left=3mm]{\Large\bfseries\color{cnum}d)}\quad\begin{minipage}[c]{5.6cm}\raggedright\Large $6{,}8$\end{minipage}\end{tcolorbox}
\end{tcolorbox}
\newpage
\begin{tcolorbox}[enhanced,width=16cm,height=14cm,colback=white,colframe=cnum,boxrule=2pt,arc=0mm,
left=6mm,right=6mm,top=0mm,bottom=4mm,sharp corners,
overlay={\fill[cnum] (frame.north west) rectangle ([yshift=-1.25cm]frame.north east);
\node[anchor=west,text=white,font=\sffamily\bfseries\LARGE] at ([xshift=5mm,yshift=-0.62cm]frame.north west) {Question 11 B};
\node[anchor=east,text=white,font=\sffamily\itshape\large] at ([xshift=-5mm,yshift=-0.62cm]frame.north east) {Nombres entiers et décimaux};}]
\vspace*{1.75cm}
\begin{minipage}[t][6.9cm][c]{\linewidth}\LARGE\centering
\begin{minipage}{0.95\linewidth}\raggedright\linespread{1.25}\selectfont Que vaut $340\div1\,000$~?\end{minipage}
\end{minipage}

\vfill
\begin{tcolorbox}[enhanced,nobeforeafter,width=7.1cm,height=2.1cm,valign=center,colback=cnum!6,colframe=cnum,arc=2.5mm,boxrule=0.9pt,left=3mm]{\Large\bfseries\color{cnum}a)}\quad\begin{minipage}[c]{5.6cm}\raggedright\Large $340\,000$\end{minipage}\end{tcolorbox}\hfill\begin{tcolorbox}[enhanced,nobeforeafter,width=7.1cm,height=2.1cm,valign=center,colback=cnum!6,colframe=cnum,arc=2.5mm,boxrule=0.9pt,left=3mm]{\Large\bfseries\color{cnum}b)}\quad\begin{minipage}[c]{5.6cm}\raggedright\Large $0{,}34$\end{minipage}\end{tcolorbox}\par\vspace{2.5mm}
\begin{tcolorbox}[enhanced,nobeforeafter,width=7.1cm,height=2.1cm,valign=center,colback=cnum!6,colframe=cnum,arc=2.5mm,boxrule=0.9pt,left=3mm]{\Large\bfseries\color{cnum}c)}\quad\begin{minipage}[c]{5.6cm}\raggedright\Large $3{,}4$\end{minipage}\end{tcolorbox}\hfill\begin{tcolorbox}[enhanced,nobeforeafter,width=7.1cm,height=2.1cm,valign=center,colback=cnum!6,colframe=cnum,arc=2.5mm,boxrule=0.9pt,left=3mm]{\Large\bfseries\color{cnum}d)}\quad\begin{minipage}[c]{5.6cm}\raggedright\Large $0{,}034$\end{minipage}\end{tcolorbox}
\end{tcolorbox}
\newpage
\begin{tcolorbox}[enhanced,width=16cm,height=14cm,colback=white,colframe=cnum,boxrule=2pt,arc=0mm,
left=6mm,right=6mm,top=0mm,bottom=4mm,sharp corners,
overlay={\fill[cnum] (frame.north west) rectangle ([yshift=-1.25cm]frame.north east);
\node[anchor=west,text=white,font=\sffamily\bfseries\LARGE] at ([xshift=5mm,yshift=-0.62cm]frame.north west) {Question 12 B};
\node[anchor=east,text=white,font=\sffamily\itshape\large] at ([xshift=-5mm,yshift=-0.62cm]frame.north east) {Nombres entiers et décimaux};}]
\vspace*{1.75cm}
\begin{minipage}[t][6.9cm][c]{\linewidth}\LARGE\centering
\begin{minipage}{0.95\linewidth}\raggedright\linespread{1.25}\selectfont Que vaut $6\div100$~?\end{minipage}
\end{minipage}

\vfill
\begin{tcolorbox}[enhanced,nobeforeafter,width=7.1cm,height=2.1cm,valign=center,colback=cnum!6,colframe=cnum,arc=2.5mm,boxrule=0.9pt,left=3mm]{\Large\bfseries\color{cnum}a)}\quad\begin{minipage}[c]{5.6cm}\raggedright\Large $0{,}006$\end{minipage}\end{tcolorbox}\hfill\begin{tcolorbox}[enhanced,nobeforeafter,width=7.1cm,height=2.1cm,valign=center,colback=cnum!6,colframe=cnum,arc=2.5mm,boxrule=0.9pt,left=3mm]{\Large\bfseries\color{cnum}b)}\quad\begin{minipage}[c]{5.6cm}\raggedright\Large $0{,}6$\end{minipage}\end{tcolorbox}\par\vspace{2.5mm}
\begin{tcolorbox}[enhanced,nobeforeafter,width=7.1cm,height=2.1cm,valign=center,colback=cnum!6,colframe=cnum,arc=2.5mm,boxrule=0.9pt,left=3mm]{\Large\bfseries\color{cnum}c)}\quad\begin{minipage}[c]{5.6cm}\raggedright\Large $0{,}06$\end{minipage}\end{tcolorbox}\hfill\begin{tcolorbox}[enhanced,nobeforeafter,width=7.1cm,height=2.1cm,valign=center,colback=cnum!6,colframe=cnum,arc=2.5mm,boxrule=0.9pt,left=3mm]{\Large\bfseries\color{cnum}d)}\quad\begin{minipage}[c]{5.6cm}\raggedright\Large $600$\end{minipage}\end{tcolorbox}
\end{tcolorbox}
\newpage
\begin{tcolorbox}[enhanced,width=16cm,height=14cm,colback=white,colframe=cnum,boxrule=2pt,arc=0mm,
left=6mm,right=6mm,top=0mm,bottom=4mm,sharp corners,
overlay={\fill[cnum] (frame.north west) rectangle ([yshift=-1.25cm]frame.north east);
\node[anchor=west,text=white,font=\sffamily\bfseries\LARGE] at ([xshift=5mm,yshift=-0.62cm]frame.north west) {Question 13 B};
\node[anchor=east,text=white,font=\sffamily\itshape\large] at ([xshift=-5mm,yshift=-0.62cm]frame.north east) {Nombres entiers et décimaux};}]
\vspace*{1.75cm}
\begin{minipage}[t][6.9cm][c]{\linewidth}\LARGE\centering
\begin{minipage}{0.95\linewidth}\raggedright\linespread{1.25}\selectfont Comment s'écrit en chiffres le nombre «~trois cent quatre mille deux cents~»~?\end{minipage}
\end{minipage}

\vfill
\begin{tcolorbox}[enhanced,nobeforeafter,width=7.1cm,height=2.1cm,valign=center,colback=cnum!6,colframe=cnum,arc=2.5mm,boxrule=0.9pt,left=3mm]{\Large\bfseries\color{cnum}a)}\quad\begin{minipage}[c]{5.6cm}\raggedright\Large $304\,200$\end{minipage}\end{tcolorbox}\hfill\begin{tcolorbox}[enhanced,nobeforeafter,width=7.1cm,height=2.1cm,valign=center,colback=cnum!6,colframe=cnum,arc=2.5mm,boxrule=0.9pt,left=3mm]{\Large\bfseries\color{cnum}b)}\quad\begin{minipage}[c]{5.6cm}\raggedright\Large $304\,020$\end{minipage}\end{tcolorbox}\par\vspace{2.5mm}
\begin{tcolorbox}[enhanced,nobeforeafter,width=7.1cm,height=2.1cm,valign=center,colback=cnum!6,colframe=cnum,arc=2.5mm,boxrule=0.9pt,left=3mm]{\Large\bfseries\color{cnum}c)}\quad\begin{minipage}[c]{5.6cm}\raggedright\Large $300\,402$\end{minipage}\end{tcolorbox}\hfill\begin{tcolorbox}[enhanced,nobeforeafter,width=7.1cm,height=2.1cm,valign=center,colback=cnum!6,colframe=cnum,arc=2.5mm,boxrule=0.9pt,left=3mm]{\Large\bfseries\color{cnum}d)}\quad\begin{minipage}[c]{5.6cm}\raggedright\Large $34\,200$\end{minipage}\end{tcolorbox}
\end{tcolorbox}
\newpage
\begin{tcolorbox}[enhanced,width=16cm,height=14cm,colback=white,colframe=cnum,boxrule=2pt,arc=0mm,
left=6mm,right=6mm,top=0mm,bottom=4mm,sharp corners,
overlay={\fill[cnum] (frame.north west) rectangle ([yshift=-1.25cm]frame.north east);
\node[anchor=west,text=white,font=\sffamily\bfseries\LARGE] at ([xshift=5mm,yshift=-0.62cm]frame.north west) {Question 14 B};
\node[anchor=east,text=white,font=\sffamily\itshape\large] at ([xshift=-5mm,yshift=-0.62cm]frame.north east) {Nombres entiers et décimaux};}]
\vspace*{1.75cm}
\begin{minipage}[t][6.9cm][c]{\linewidth}\LARGE\centering
\begin{minipage}{0.95\linewidth}\raggedright\linespread{1.25}\selectfont Comment s'écrit en chiffres le nombre «~sept millions vingt mille cinq~»~?\end{minipage}
\end{minipage}

\vfill
\begin{tcolorbox}[enhanced,nobeforeafter,width=7.1cm,height=2.1cm,valign=center,colback=cnum!6,colframe=cnum,arc=2.5mm,boxrule=0.9pt,left=3mm]{\Large\bfseries\color{cnum}a)}\quad\begin{minipage}[c]{5.6cm}\raggedright\Large $7\,020\,050$\end{minipage}\end{tcolorbox}\hfill\begin{tcolorbox}[enhanced,nobeforeafter,width=7.1cm,height=2.1cm,valign=center,colback=cnum!6,colframe=cnum,arc=2.5mm,boxrule=0.9pt,left=3mm]{\Large\bfseries\color{cnum}b)}\quad\begin{minipage}[c]{5.6cm}\raggedright\Large $720\,005$\end{minipage}\end{tcolorbox}\par\vspace{2.5mm}
\begin{tcolorbox}[enhanced,nobeforeafter,width=7.1cm,height=2.1cm,valign=center,colback=cnum!6,colframe=cnum,arc=2.5mm,boxrule=0.9pt,left=3mm]{\Large\bfseries\color{cnum}c)}\quad\begin{minipage}[c]{5.6cm}\raggedright\Large $7\,020\,005$\end{minipage}\end{tcolorbox}\hfill\begin{tcolorbox}[enhanced,nobeforeafter,width=7.1cm,height=2.1cm,valign=center,colback=cnum!6,colframe=cnum,arc=2.5mm,boxrule=0.9pt,left=3mm]{\Large\bfseries\color{cnum}d)}\quad\begin{minipage}[c]{5.6cm}\raggedright\Large $7\,200\,005$\end{minipage}\end{tcolorbox}
\end{tcolorbox}
\newpage
\begin{tcolorbox}[enhanced,width=16cm,height=14cm,colback=white,colframe=cnum,boxrule=2pt,arc=0mm,
left=6mm,right=6mm,top=0mm,bottom=4mm,sharp corners,
overlay={\fill[cnum] (frame.north west) rectangle ([yshift=-1.25cm]frame.north east);
\node[anchor=west,text=white,font=\sffamily\bfseries\LARGE] at ([xshift=5mm,yshift=-0.62cm]frame.north west) {Question 15 B};
\node[anchor=east,text=white,font=\sffamily\itshape\large] at ([xshift=-5mm,yshift=-0.62cm]frame.north east) {Nombres entiers et décimaux};}]
\vspace*{1.75cm}
\begin{minipage}[t][6.9cm][c]{\linewidth}\LARGE\centering
\begin{minipage}{0.95\linewidth}\raggedright\linespread{1.25}\selectfont Quel nombre est compris entre $0{,}5$ et $0{,}6$~?\end{minipage}
\end{minipage}

\vfill
\begin{tcolorbox}[enhanced,nobeforeafter,width=7.1cm,height=2.1cm,valign=center,colback=cnum!6,colframe=cnum,arc=2.5mm,boxrule=0.9pt,left=3mm]{\Large\bfseries\color{cnum}a)}\quad\begin{minipage}[c]{5.6cm}\raggedright\Large $0{,}61$\end{minipage}\end{tcolorbox}\hfill\begin{tcolorbox}[enhanced,nobeforeafter,width=7.1cm,height=2.1cm,valign=center,colback=cnum!6,colframe=cnum,arc=2.5mm,boxrule=0.9pt,left=3mm]{\Large\bfseries\color{cnum}b)}\quad\begin{minipage}[c]{5.6cm}\raggedright\Large $0{,}49$\end{minipage}\end{tcolorbox}\par\vspace{2.5mm}
\begin{tcolorbox}[enhanced,nobeforeafter,width=7.1cm,height=2.1cm,valign=center,colback=cnum!6,colframe=cnum,arc=2.5mm,boxrule=0.9pt,left=3mm]{\Large\bfseries\color{cnum}c)}\quad\begin{minipage}[c]{5.6cm}\raggedright\Large $0{,}065$\end{minipage}\end{tcolorbox}\hfill\begin{tcolorbox}[enhanced,nobeforeafter,width=7.1cm,height=2.1cm,valign=center,colback=cnum!6,colframe=cnum,arc=2.5mm,boxrule=0.9pt,left=3mm]{\Large\bfseries\color{cnum}d)}\quad\begin{minipage}[c]{5.6cm}\raggedright\Large $0{,}58$\end{minipage}\end{tcolorbox}
\end{tcolorbox}
\newpage
\begin{tcolorbox}[enhanced,width=16cm,height=14cm,colback=white,colframe=cnum,boxrule=2pt,arc=0mm,
left=6mm,right=6mm,top=0mm,bottom=4mm,sharp corners,
overlay={\fill[cnum] (frame.north west) rectangle ([yshift=-1.25cm]frame.north east);
\node[anchor=west,text=white,font=\sffamily\bfseries\LARGE] at ([xshift=5mm,yshift=-0.62cm]frame.north west) {Question 16 B};
\node[anchor=east,text=white,font=\sffamily\itshape\large] at ([xshift=-5mm,yshift=-0.62cm]frame.north east) {Nombres entiers et décimaux};}]
\vspace*{1.75cm}
\begin{minipage}[t][6.9cm][c]{\linewidth}\LARGE\centering
\begin{minipage}{0.95\linewidth}\raggedright\linespread{1.25}\selectfont Quel nombre est compris entre $12{,}09$ et $12{,}1$~?\end{minipage}
\end{minipage}

\vfill
\begin{tcolorbox}[enhanced,nobeforeafter,width=7.1cm,height=2.1cm,valign=center,colback=cnum!6,colframe=cnum,arc=2.5mm,boxrule=0.9pt,left=3mm]{\Large\bfseries\color{cnum}a)}\quad\begin{minipage}[c]{5.6cm}\raggedright\Large $12{,}095$\end{minipage}\end{tcolorbox}\hfill\begin{tcolorbox}[enhanced,nobeforeafter,width=7.1cm,height=2.1cm,valign=center,colback=cnum!6,colframe=cnum,arc=2.5mm,boxrule=0.9pt,left=3mm]{\Large\bfseries\color{cnum}b)}\quad\begin{minipage}[c]{5.6cm}\raggedright\Large $12{,}08$\end{minipage}\end{tcolorbox}\par\vspace{2.5mm}
\begin{tcolorbox}[enhanced,nobeforeafter,width=7.1cm,height=2.1cm,valign=center,colback=cnum!6,colframe=cnum,arc=2.5mm,boxrule=0.9pt,left=3mm]{\Large\bfseries\color{cnum}c)}\quad\begin{minipage}[c]{5.6cm}\raggedright\Large $12{,}11$\end{minipage}\end{tcolorbox}\hfill\begin{tcolorbox}[enhanced,nobeforeafter,width=7.1cm,height=2.1cm,valign=center,colback=cnum!6,colframe=cnum,arc=2.5mm,boxrule=0.9pt,left=3mm]{\Large\bfseries\color{cnum}d)}\quad\begin{minipage}[c]{5.6cm}\raggedright\Large $12{,}9$\end{minipage}\end{tcolorbox}
\end{tcolorbox}
\newpage
\begin{tcolorbox}[enhanced,width=16cm,height=14cm,colback=white,colframe=ccal,boxrule=2pt,arc=0mm,
left=6mm,right=6mm,top=0mm,bottom=4mm,sharp corners,
overlay={\fill[ccal] (frame.north west) rectangle ([yshift=-1.25cm]frame.north east);
\node[anchor=west,text=white,font=\sffamily\bfseries\LARGE] at ([xshift=5mm,yshift=-0.62cm]frame.north west) {Question 17 B};
\node[anchor=east,text=white,font=\sffamily\itshape\large] at ([xshift=-5mm,yshift=-0.62cm]frame.north east) {Calculs};}]
\vspace*{1.75cm}
\begin{minipage}[t][6.9cm][c]{\linewidth}\LARGE\centering
\begin{minipage}{0.95\linewidth}\raggedright\linespread{1.25}\selectfont Que vaut $0{,}95+2{,}4$~?\end{minipage}
\end{minipage}

\vfill
\begin{tcolorbox}[enhanced,nobeforeafter,width=7.1cm,height=2.1cm,valign=center,colback=ccal!6,colframe=ccal,arc=2.5mm,boxrule=0.9pt,left=3mm]{\Large\bfseries\color{ccal}a)}\quad\begin{minipage}[c]{5.6cm}\raggedright\Large $4{,}35$\end{minipage}\end{tcolorbox}\hfill\begin{tcolorbox}[enhanced,nobeforeafter,width=7.1cm,height=2.1cm,valign=center,colback=ccal!6,colframe=ccal,arc=2.5mm,boxrule=0.9pt,left=3mm]{\Large\bfseries\color{ccal}b)}\quad\begin{minipage}[c]{5.6cm}\raggedright\Large $2{,}35$\end{minipage}\end{tcolorbox}\par\vspace{2.5mm}
\begin{tcolorbox}[enhanced,nobeforeafter,width=7.1cm,height=2.1cm,valign=center,colback=ccal!6,colframe=ccal,arc=2.5mm,boxrule=0.9pt,left=3mm]{\Large\bfseries\color{ccal}c)}\quad\begin{minipage}[c]{5.6cm}\raggedright\Large $3{,}35$\end{minipage}\end{tcolorbox}\hfill\begin{tcolorbox}[enhanced,nobeforeafter,width=7.1cm,height=2.1cm,valign=center,colback=ccal!6,colframe=ccal,arc=2.5mm,boxrule=0.9pt,left=3mm]{\Large\bfseries\color{ccal}d)}\quad\begin{minipage}[c]{5.6cm}\raggedright\Large $3{,}45$\end{minipage}\end{tcolorbox}
\end{tcolorbox}
\newpage
\begin{tcolorbox}[enhanced,width=16cm,height=14cm,colback=white,colframe=ccal,boxrule=2pt,arc=0mm,
left=6mm,right=6mm,top=0mm,bottom=4mm,sharp corners,
overlay={\fill[ccal] (frame.north west) rectangle ([yshift=-1.25cm]frame.north east);
\node[anchor=west,text=white,font=\sffamily\bfseries\LARGE] at ([xshift=5mm,yshift=-0.62cm]frame.north west) {Question 18 B};
\node[anchor=east,text=white,font=\sffamily\itshape\large] at ([xshift=-5mm,yshift=-0.62cm]frame.north east) {Calculs};}]
\vspace*{1.75cm}
\begin{minipage}[t][6.9cm][c]{\linewidth}\LARGE\centering
\begin{minipage}{0.95\linewidth}\raggedright\linespread{1.25}\selectfont Que vaut $25{,}3+4{,}75$~?\end{minipage}
\end{minipage}

\vfill
\begin{tcolorbox}[enhanced,nobeforeafter,width=7.1cm,height=2.1cm,valign=center,colback=ccal!6,colframe=ccal,arc=2.5mm,boxrule=0.9pt,left=3mm]{\Large\bfseries\color{ccal}a)}\quad\begin{minipage}[c]{5.6cm}\raggedright\Large $31{,}05$\end{minipage}\end{tcolorbox}\hfill\begin{tcolorbox}[enhanced,nobeforeafter,width=7.1cm,height=2.1cm,valign=center,colback=ccal!6,colframe=ccal,arc=2.5mm,boxrule=0.9pt,left=3mm]{\Large\bfseries\color{ccal}b)}\quad\begin{minipage}[c]{5.6cm}\raggedright\Large $30{,}15$\end{minipage}\end{tcolorbox}\par\vspace{2.5mm}
\begin{tcolorbox}[enhanced,nobeforeafter,width=7.1cm,height=2.1cm,valign=center,colback=ccal!6,colframe=ccal,arc=2.5mm,boxrule=0.9pt,left=3mm]{\Large\bfseries\color{ccal}c)}\quad\begin{minipage}[c]{5.6cm}\raggedright\Large $30{,}05$\end{minipage}\end{tcolorbox}\hfill\begin{tcolorbox}[enhanced,nobeforeafter,width=7.1cm,height=2.1cm,valign=center,colback=ccal!6,colframe=ccal,arc=2.5mm,boxrule=0.9pt,left=3mm]{\Large\bfseries\color{ccal}d)}\quad\begin{minipage}[c]{5.6cm}\raggedright\Large $29{,}05$\end{minipage}\end{tcolorbox}
\end{tcolorbox}
\newpage
\begin{tcolorbox}[enhanced,width=16cm,height=14cm,colback=white,colframe=ccal,boxrule=2pt,arc=0mm,
left=6mm,right=6mm,top=0mm,bottom=4mm,sharp corners,
overlay={\fill[ccal] (frame.north west) rectangle ([yshift=-1.25cm]frame.north east);
\node[anchor=west,text=white,font=\sffamily\bfseries\LARGE] at ([xshift=5mm,yshift=-0.62cm]frame.north west) {Question 19 B};
\node[anchor=east,text=white,font=\sffamily\itshape\large] at ([xshift=-5mm,yshift=-0.62cm]frame.north east) {Calculs};}]
\vspace*{1.75cm}
\begin{minipage}[t][6.9cm][c]{\linewidth}\LARGE\centering
\begin{minipage}{0.95\linewidth}\raggedright\linespread{1.25}\selectfont Que vaut $7{,}02-3{,}5$~?\end{minipage}
\end{minipage}

\vfill
\begin{tcolorbox}[enhanced,nobeforeafter,width=7.1cm,height=2.1cm,valign=center,colback=ccal!6,colframe=ccal,arc=2.5mm,boxrule=0.9pt,left=3mm]{\Large\bfseries\color{ccal}a)}\quad\begin{minipage}[c]{5.6cm}\raggedright\Large $4{,}52$\end{minipage}\end{tcolorbox}\hfill\begin{tcolorbox}[enhanced,nobeforeafter,width=7.1cm,height=2.1cm,valign=center,colback=ccal!6,colframe=ccal,arc=2.5mm,boxrule=0.9pt,left=3mm]{\Large\bfseries\color{ccal}b)}\quad\begin{minipage}[c]{5.6cm}\raggedright\Large $3{,}62$\end{minipage}\end{tcolorbox}\par\vspace{2.5mm}
\begin{tcolorbox}[enhanced,nobeforeafter,width=7.1cm,height=2.1cm,valign=center,colback=ccal!6,colframe=ccal,arc=2.5mm,boxrule=0.9pt,left=3mm]{\Large\bfseries\color{ccal}c)}\quad\begin{minipage}[c]{5.6cm}\raggedright\Large $3{,}42$\end{minipage}\end{tcolorbox}\hfill\begin{tcolorbox}[enhanced,nobeforeafter,width=7.1cm,height=2.1cm,valign=center,colback=ccal!6,colframe=ccal,arc=2.5mm,boxrule=0.9pt,left=3mm]{\Large\bfseries\color{ccal}d)}\quad\begin{minipage}[c]{5.6cm}\raggedright\Large $3{,}52$\end{minipage}\end{tcolorbox}
\end{tcolorbox}
\newpage
\begin{tcolorbox}[enhanced,width=16cm,height=14cm,colback=white,colframe=ccal,boxrule=2pt,arc=0mm,
left=6mm,right=6mm,top=0mm,bottom=4mm,sharp corners,
overlay={\fill[ccal] (frame.north west) rectangle ([yshift=-1.25cm]frame.north east);
\node[anchor=west,text=white,font=\sffamily\bfseries\LARGE] at ([xshift=5mm,yshift=-0.62cm]frame.north west) {Question 20 B};
\node[anchor=east,text=white,font=\sffamily\itshape\large] at ([xshift=-5mm,yshift=-0.62cm]frame.north east) {Calculs};}]
\vspace*{1.75cm}
\begin{minipage}[t][6.9cm][c]{\linewidth}\LARGE\centering
\begin{minipage}{0.95\linewidth}\raggedright\linespread{1.25}\selectfont Que vaut $100-45{,}8$~?\end{minipage}
\end{minipage}

\vfill
\begin{tcolorbox}[enhanced,nobeforeafter,width=7.1cm,height=2.1cm,valign=center,colback=ccal!6,colframe=ccal,arc=2.5mm,boxrule=0.9pt,left=3mm]{\Large\bfseries\color{ccal}a)}\quad\begin{minipage}[c]{5.6cm}\raggedright\Large $55{,}2$\end{minipage}\end{tcolorbox}\hfill\begin{tcolorbox}[enhanced,nobeforeafter,width=7.1cm,height=2.1cm,valign=center,colback=ccal!6,colframe=ccal,arc=2.5mm,boxrule=0.9pt,left=3mm]{\Large\bfseries\color{ccal}b)}\quad\begin{minipage}[c]{5.6cm}\raggedright\Large $54{,}3$\end{minipage}\end{tcolorbox}\par\vspace{2.5mm}
\begin{tcolorbox}[enhanced,nobeforeafter,width=7.1cm,height=2.1cm,valign=center,colback=ccal!6,colframe=ccal,arc=2.5mm,boxrule=0.9pt,left=3mm]{\Large\bfseries\color{ccal}c)}\quad\begin{minipage}[c]{5.6cm}\raggedright\Large $54{,}1$\end{minipage}\end{tcolorbox}\hfill\begin{tcolorbox}[enhanced,nobeforeafter,width=7.1cm,height=2.1cm,valign=center,colback=ccal!6,colframe=ccal,arc=2.5mm,boxrule=0.9pt,left=3mm]{\Large\bfseries\color{ccal}d)}\quad\begin{minipage}[c]{5.6cm}\raggedright\Large $54{,}2$\end{minipage}\end{tcolorbox}
\end{tcolorbox}
\newpage
\begin{tcolorbox}[enhanced,width=16cm,height=14cm,colback=white,colframe=ccal,boxrule=2pt,arc=0mm,
left=6mm,right=6mm,top=0mm,bottom=4mm,sharp corners,
overlay={\fill[ccal] (frame.north west) rectangle ([yshift=-1.25cm]frame.north east);
\node[anchor=west,text=white,font=\sffamily\bfseries\LARGE] at ([xshift=5mm,yshift=-0.62cm]frame.north west) {Question 21 B};
\node[anchor=east,text=white,font=\sffamily\itshape\large] at ([xshift=-5mm,yshift=-0.62cm]frame.north east) {Calculs};}]
\vspace*{1.75cm}
\begin{minipage}[t][6.9cm][c]{\linewidth}\LARGE\centering
\begin{minipage}{0.95\linewidth}\raggedright\linespread{1.25}\selectfont Que vaut $3{,}6\times5$~?\end{minipage}
\end{minipage}

\vfill
\begin{tcolorbox}[enhanced,nobeforeafter,width=7.1cm,height=2.1cm,valign=center,colback=ccal!6,colframe=ccal,arc=2.5mm,boxrule=0.9pt,left=3mm]{\Large\bfseries\color{ccal}a)}\quad\begin{minipage}[c]{5.6cm}\raggedright\Large $18$\end{minipage}\end{tcolorbox}\hfill\begin{tcolorbox}[enhanced,nobeforeafter,width=7.1cm,height=2.1cm,valign=center,colback=ccal!6,colframe=ccal,arc=2.5mm,boxrule=0.9pt,left=3mm]{\Large\bfseries\color{ccal}b)}\quad\begin{minipage}[c]{5.6cm}\raggedright\Large $180$\end{minipage}\end{tcolorbox}\par\vspace{2.5mm}
\begin{tcolorbox}[enhanced,nobeforeafter,width=7.1cm,height=2.1cm,valign=center,colback=ccal!6,colframe=ccal,arc=2.5mm,boxrule=0.9pt,left=3mm]{\Large\bfseries\color{ccal}c)}\quad\begin{minipage}[c]{5.6cm}\raggedright\Large $1{,}8$\end{minipage}\end{tcolorbox}\hfill\begin{tcolorbox}[enhanced,nobeforeafter,width=7.1cm,height=2.1cm,valign=center,colback=ccal!6,colframe=ccal,arc=2.5mm,boxrule=0.9pt,left=3mm]{\Large\bfseries\color{ccal}d)}\quad\begin{minipage}[c]{5.6cm}\raggedright\Large $8{,}6$\end{minipage}\end{tcolorbox}
\end{tcolorbox}
\newpage
\begin{tcolorbox}[enhanced,width=16cm,height=14cm,colback=white,colframe=ccal,boxrule=2pt,arc=0mm,
left=6mm,right=6mm,top=0mm,bottom=4mm,sharp corners,
overlay={\fill[ccal] (frame.north west) rectangle ([yshift=-1.25cm]frame.north east);
\node[anchor=west,text=white,font=\sffamily\bfseries\LARGE] at ([xshift=5mm,yshift=-0.62cm]frame.north west) {Question 22 B};
\node[anchor=east,text=white,font=\sffamily\itshape\large] at ([xshift=-5mm,yshift=-0.62cm]frame.north east) {Calculs};}]
\vspace*{1.75cm}
\begin{minipage}[t][6.9cm][c]{\linewidth}\LARGE\centering
\begin{minipage}{0.95\linewidth}\raggedright\linespread{1.25}\selectfont Que vaut $1{,}5\times1{,}5$~?\end{minipage}
\end{minipage}

\vfill
\begin{tcolorbox}[enhanced,nobeforeafter,width=7.1cm,height=2.1cm,valign=center,colback=ccal!6,colframe=ccal,arc=2.5mm,boxrule=0.9pt,left=3mm]{\Large\bfseries\color{ccal}a)}\quad\begin{minipage}[c]{5.6cm}\raggedright\Large $2{,}25$\end{minipage}\end{tcolorbox}\hfill\begin{tcolorbox}[enhanced,nobeforeafter,width=7.1cm,height=2.1cm,valign=center,colback=ccal!6,colframe=ccal,arc=2.5mm,boxrule=0.9pt,left=3mm]{\Large\bfseries\color{ccal}b)}\quad\begin{minipage}[c]{5.6cm}\raggedright\Large $0{,}225$\end{minipage}\end{tcolorbox}\par\vspace{2.5mm}
\begin{tcolorbox}[enhanced,nobeforeafter,width=7.1cm,height=2.1cm,valign=center,colback=ccal!6,colframe=ccal,arc=2.5mm,boxrule=0.9pt,left=3mm]{\Large\bfseries\color{ccal}c)}\quad\begin{minipage}[c]{5.6cm}\raggedright\Large $3$\end{minipage}\end{tcolorbox}\hfill\begin{tcolorbox}[enhanced,nobeforeafter,width=7.1cm,height=2.1cm,valign=center,colback=ccal!6,colframe=ccal,arc=2.5mm,boxrule=0.9pt,left=3mm]{\Large\bfseries\color{ccal}d)}\quad\begin{minipage}[c]{5.6cm}\raggedright\Large $22{,}5$\end{minipage}\end{tcolorbox}
\end{tcolorbox}
\newpage
\begin{tcolorbox}[enhanced,width=16cm,height=14cm,colback=white,colframe=ccal,boxrule=2pt,arc=0mm,
left=6mm,right=6mm,top=0mm,bottom=4mm,sharp corners,
overlay={\fill[ccal] (frame.north west) rectangle ([yshift=-1.25cm]frame.north east);
\node[anchor=west,text=white,font=\sffamily\bfseries\LARGE] at ([xshift=5mm,yshift=-0.62cm]frame.north west) {Question 23 B};
\node[anchor=east,text=white,font=\sffamily\itshape\large] at ([xshift=-5mm,yshift=-0.62cm]frame.north east) {Calculs};}]
\vspace*{1.75cm}
\begin{minipage}[t][6.9cm][c]{\linewidth}\LARGE\centering
\begin{minipage}{0.95\linewidth}\raggedright\linespread{1.25}\selectfont Dans la division euclidienne de $245$ par $9$, quels sont le quotient $q$ et le reste $r$~?\end{minipage}
\end{minipage}

\vfill
\begin{tcolorbox}[enhanced,nobeforeafter,width=7.1cm,height=2.1cm,valign=center,colback=ccal!6,colframe=ccal,arc=2.5mm,boxrule=0.9pt,left=3mm]{\Large\bfseries\color{ccal}a)}\quad\begin{minipage}[c]{5.6cm}\raggedright\Large $q=28\text{ et }r=2$\end{minipage}\end{tcolorbox}\hfill\begin{tcolorbox}[enhanced,nobeforeafter,width=7.1cm,height=2.1cm,valign=center,colback=ccal!6,colframe=ccal,arc=2.5mm,boxrule=0.9pt,left=3mm]{\Large\bfseries\color{ccal}b)}\quad\begin{minipage}[c]{5.6cm}\raggedright\Large $q=27\text{ et }r=3$\end{minipage}\end{tcolorbox}\par\vspace{2.5mm}
\begin{tcolorbox}[enhanced,nobeforeafter,width=7.1cm,height=2.1cm,valign=center,colback=ccal!6,colframe=ccal,arc=2.5mm,boxrule=0.9pt,left=3mm]{\Large\bfseries\color{ccal}c)}\quad\begin{minipage}[c]{5.6cm}\raggedright\Large $q=26\text{ et }r=11$\end{minipage}\end{tcolorbox}\hfill\begin{tcolorbox}[enhanced,nobeforeafter,width=7.1cm,height=2.1cm,valign=center,colback=ccal!6,colframe=ccal,arc=2.5mm,boxrule=0.9pt,left=3mm]{\Large\bfseries\color{ccal}d)}\quad\begin{minipage}[c]{5.6cm}\raggedright\Large $q=27\text{ et }r=2$\end{minipage}\end{tcolorbox}
\end{tcolorbox}
\newpage
\begin{tcolorbox}[enhanced,width=16cm,height=14cm,colback=white,colframe=ccal,boxrule=2pt,arc=0mm,
left=6mm,right=6mm,top=0mm,bottom=4mm,sharp corners,
overlay={\fill[ccal] (frame.north west) rectangle ([yshift=-1.25cm]frame.north east);
\node[anchor=west,text=white,font=\sffamily\bfseries\LARGE] at ([xshift=5mm,yshift=-0.62cm]frame.north west) {Question 24 B};
\node[anchor=east,text=white,font=\sffamily\itshape\large] at ([xshift=-5mm,yshift=-0.62cm]frame.north east) {Calculs};}]
\vspace*{1.75cm}
\begin{minipage}[t][6.9cm][c]{\linewidth}\LARGE\centering
\begin{minipage}{0.95\linewidth}\raggedright\linespread{1.25}\selectfont Dans la division euclidienne de $83$ par $6$, quels sont le quotient $q$ et le reste $r$~?\end{minipage}
\end{minipage}

\vfill
\begin{tcolorbox}[enhanced,nobeforeafter,width=7.1cm,height=2.1cm,valign=center,colback=ccal!6,colframe=ccal,arc=2.5mm,boxrule=0.9pt,left=3mm]{\Large\bfseries\color{ccal}a)}\quad\begin{minipage}[c]{5.6cm}\raggedright\Large $q=13\text{ et }r=6$\end{minipage}\end{tcolorbox}\hfill\begin{tcolorbox}[enhanced,nobeforeafter,width=7.1cm,height=2.1cm,valign=center,colback=ccal!6,colframe=ccal,arc=2.5mm,boxrule=0.9pt,left=3mm]{\Large\bfseries\color{ccal}b)}\quad\begin{minipage}[c]{5.6cm}\raggedright\Large $q=12\text{ et }r=11$\end{minipage}\end{tcolorbox}\par\vspace{2.5mm}
\begin{tcolorbox}[enhanced,nobeforeafter,width=7.1cm,height=2.1cm,valign=center,colback=ccal!6,colframe=ccal,arc=2.5mm,boxrule=0.9pt,left=3mm]{\Large\bfseries\color{ccal}c)}\quad\begin{minipage}[c]{5.6cm}\raggedright\Large $q=14\text{ et }r=5$\end{minipage}\end{tcolorbox}\hfill\begin{tcolorbox}[enhanced,nobeforeafter,width=7.1cm,height=2.1cm,valign=center,colback=ccal!6,colframe=ccal,arc=2.5mm,boxrule=0.9pt,left=3mm]{\Large\bfseries\color{ccal}d)}\quad\begin{minipage}[c]{5.6cm}\raggedright\Large $q=13\text{ et }r=5$\end{minipage}\end{tcolorbox}
\end{tcolorbox}
\newpage
\begin{tcolorbox}[enhanced,width=16cm,height=14cm,colback=white,colframe=ccal,boxrule=2pt,arc=0mm,
left=6mm,right=6mm,top=0mm,bottom=4mm,sharp corners,
overlay={\fill[ccal] (frame.north west) rectangle ([yshift=-1.25cm]frame.north east);
\node[anchor=west,text=white,font=\sffamily\bfseries\LARGE] at ([xshift=5mm,yshift=-0.62cm]frame.north west) {Question 25 B};
\node[anchor=east,text=white,font=\sffamily\itshape\large] at ([xshift=-5mm,yshift=-0.62cm]frame.north east) {Calculs};}]
\vspace*{1.75cm}
\begin{minipage}[t][6.9cm][c]{\linewidth}\LARGE\centering
\begin{minipage}{0.95\linewidth}\raggedright\linespread{1.25}\selectfont Lequel de ces nombres est un multiple de $6$~?\end{minipage}
\end{minipage}

\vfill
\begin{tcolorbox}[enhanced,nobeforeafter,width=7.1cm,height=2.1cm,valign=center,colback=ccal!6,colframe=ccal,arc=2.5mm,boxrule=0.9pt,left=3mm]{\Large\bfseries\color{ccal}a)}\quad\begin{minipage}[c]{5.6cm}\raggedright\Large $84$\end{minipage}\end{tcolorbox}\hfill\begin{tcolorbox}[enhanced,nobeforeafter,width=7.1cm,height=2.1cm,valign=center,colback=ccal!6,colframe=ccal,arc=2.5mm,boxrule=0.9pt,left=3mm]{\Large\bfseries\color{ccal}b)}\quad\begin{minipage}[c]{5.6cm}\raggedright\Large $86$\end{minipage}\end{tcolorbox}\par\vspace{2.5mm}
\begin{tcolorbox}[enhanced,nobeforeafter,width=7.1cm,height=2.1cm,valign=center,colback=ccal!6,colframe=ccal,arc=2.5mm,boxrule=0.9pt,left=3mm]{\Large\bfseries\color{ccal}c)}\quad\begin{minipage}[c]{5.6cm}\raggedright\Large $76$\end{minipage}\end{tcolorbox}\hfill\begin{tcolorbox}[enhanced,nobeforeafter,width=7.1cm,height=2.1cm,valign=center,colback=ccal!6,colframe=ccal,arc=2.5mm,boxrule=0.9pt,left=3mm]{\Large\bfseries\color{ccal}d)}\quad\begin{minipage}[c]{5.6cm}\raggedright\Large $64$\end{minipage}\end{tcolorbox}
\end{tcolorbox}
\newpage
\begin{tcolorbox}[enhanced,width=16cm,height=14cm,colback=white,colframe=ccal,boxrule=2pt,arc=0mm,
left=6mm,right=6mm,top=0mm,bottom=4mm,sharp corners,
overlay={\fill[ccal] (frame.north west) rectangle ([yshift=-1.25cm]frame.north east);
\node[anchor=west,text=white,font=\sffamily\bfseries\LARGE] at ([xshift=5mm,yshift=-0.62cm]frame.north west) {Question 26 B};
\node[anchor=east,text=white,font=\sffamily\itshape\large] at ([xshift=-5mm,yshift=-0.62cm]frame.north east) {Calculs};}]
\vspace*{1.75cm}
\begin{minipage}[t][6.9cm][c]{\linewidth}\LARGE\centering
\begin{minipage}{0.95\linewidth}\raggedright\linespread{1.25}\selectfont Lequel de ces nombres est un multiple de $8$~?\end{minipage}
\end{minipage}

\vfill
\begin{tcolorbox}[enhanced,nobeforeafter,width=7.1cm,height=2.1cm,valign=center,colback=ccal!6,colframe=ccal,arc=2.5mm,boxrule=0.9pt,left=3mm]{\Large\bfseries\color{ccal}a)}\quad\begin{minipage}[c]{5.6cm}\raggedright\Large $98$\end{minipage}\end{tcolorbox}\hfill\begin{tcolorbox}[enhanced,nobeforeafter,width=7.1cm,height=2.1cm,valign=center,colback=ccal!6,colframe=ccal,arc=2.5mm,boxrule=0.9pt,left=3mm]{\Large\bfseries\color{ccal}b)}\quad\begin{minipage}[c]{5.6cm}\raggedright\Large $86$\end{minipage}\end{tcolorbox}\par\vspace{2.5mm}
\begin{tcolorbox}[enhanced,nobeforeafter,width=7.1cm,height=2.1cm,valign=center,colback=ccal!6,colframe=ccal,arc=2.5mm,boxrule=0.9pt,left=3mm]{\Large\bfseries\color{ccal}c)}\quad\begin{minipage}[c]{5.6cm}\raggedright\Large $106$\end{minipage}\end{tcolorbox}\hfill\begin{tcolorbox}[enhanced,nobeforeafter,width=7.1cm,height=2.1cm,valign=center,colback=ccal!6,colframe=ccal,arc=2.5mm,boxrule=0.9pt,left=3mm]{\Large\bfseries\color{ccal}d)}\quad\begin{minipage}[c]{5.6cm}\raggedright\Large $96$\end{minipage}\end{tcolorbox}
\end{tcolorbox}
\newpage
\begin{tcolorbox}[enhanced,width=16cm,height=14cm,colback=white,colframe=ccal,boxrule=2pt,arc=0mm,
left=6mm,right=6mm,top=0mm,bottom=4mm,sharp corners,
overlay={\fill[ccal] (frame.north west) rectangle ([yshift=-1.25cm]frame.north east);
\node[anchor=west,text=white,font=\sffamily\bfseries\LARGE] at ([xshift=5mm,yshift=-0.62cm]frame.north west) {Question 27 B};
\node[anchor=east,text=white,font=\sffamily\itshape\large] at ([xshift=-5mm,yshift=-0.62cm]frame.north east) {Calculs};}]
\vspace*{1.75cm}
\begin{minipage}[t][6.9cm][c]{\linewidth}\LARGE\centering
\begin{minipage}{0.95\linewidth}\raggedright\linespread{1.25}\selectfont Lequel de ces nombres est divisible par $5$~?\end{minipage}
\end{minipage}

\vfill
\begin{tcolorbox}[enhanced,nobeforeafter,width=7.1cm,height=2.1cm,valign=center,colback=ccal!6,colframe=ccal,arc=2.5mm,boxrule=0.9pt,left=3mm]{\Large\bfseries\color{ccal}a)}\quad\begin{minipage}[c]{5.6cm}\raggedright\Large $1\,352$\end{minipage}\end{tcolorbox}\hfill\begin{tcolorbox}[enhanced,nobeforeafter,width=7.1cm,height=2.1cm,valign=center,colback=ccal!6,colframe=ccal,arc=2.5mm,boxrule=0.9pt,left=3mm]{\Large\bfseries\color{ccal}b)}\quad\begin{minipage}[c]{5.6cm}\raggedright\Large $1\,235$\end{minipage}\end{tcolorbox}\par\vspace{2.5mm}
\begin{tcolorbox}[enhanced,nobeforeafter,width=7.1cm,height=2.1cm,valign=center,colback=ccal!6,colframe=ccal,arc=2.5mm,boxrule=0.9pt,left=3mm]{\Large\bfseries\color{ccal}c)}\quad\begin{minipage}[c]{5.6cm}\raggedright\Large $1\,253$\end{minipage}\end{tcolorbox}\hfill\begin{tcolorbox}[enhanced,nobeforeafter,width=7.1cm,height=2.1cm,valign=center,colback=ccal!6,colframe=ccal,arc=2.5mm,boxrule=0.9pt,left=3mm]{\Large\bfseries\color{ccal}d)}\quad\begin{minipage}[c]{5.6cm}\raggedright\Large $1\,532$\end{minipage}\end{tcolorbox}
\end{tcolorbox}
\newpage
\begin{tcolorbox}[enhanced,width=16cm,height=14cm,colback=white,colframe=ccal,boxrule=2pt,arc=0mm,
left=6mm,right=6mm,top=0mm,bottom=4mm,sharp corners,
overlay={\fill[ccal] (frame.north west) rectangle ([yshift=-1.25cm]frame.north east);
\node[anchor=west,text=white,font=\sffamily\bfseries\LARGE] at ([xshift=5mm,yshift=-0.62cm]frame.north west) {Question 28 B};
\node[anchor=east,text=white,font=\sffamily\itshape\large] at ([xshift=-5mm,yshift=-0.62cm]frame.north east) {Calculs};}]
\vspace*{1.75cm}
\begin{minipage}[t][6.9cm][c]{\linewidth}\LARGE\centering
\begin{minipage}{0.95\linewidth}\raggedright\linespread{1.25}\selectfont Lequel de ces nombres est divisible par $2$~?\end{minipage}
\end{minipage}

\vfill
\begin{tcolorbox}[enhanced,nobeforeafter,width=7.1cm,height=2.1cm,valign=center,colback=ccal!6,colframe=ccal,arc=2.5mm,boxrule=0.9pt,left=3mm]{\Large\bfseries\color{ccal}a)}\quad\begin{minipage}[c]{5.6cm}\raggedright\Large $7\,835$\end{minipage}\end{tcolorbox}\hfill\begin{tcolorbox}[enhanced,nobeforeafter,width=7.1cm,height=2.1cm,valign=center,colback=ccal!6,colframe=ccal,arc=2.5mm,boxrule=0.9pt,left=3mm]{\Large\bfseries\color{ccal}b)}\quad\begin{minipage}[c]{5.6cm}\raggedright\Large $5\,783$\end{minipage}\end{tcolorbox}\par\vspace{2.5mm}
\begin{tcolorbox}[enhanced,nobeforeafter,width=7.1cm,height=2.1cm,valign=center,colback=ccal!6,colframe=ccal,arc=2.5mm,boxrule=0.9pt,left=3mm]{\Large\bfseries\color{ccal}c)}\quad\begin{minipage}[c]{5.6cm}\raggedright\Large $3\,587$\end{minipage}\end{tcolorbox}\hfill\begin{tcolorbox}[enhanced,nobeforeafter,width=7.1cm,height=2.1cm,valign=center,colback=ccal!6,colframe=ccal,arc=2.5mm,boxrule=0.9pt,left=3mm]{\Large\bfseries\color{ccal}d)}\quad\begin{minipage}[c]{5.6cm}\raggedright\Large $3\,578$\end{minipage}\end{tcolorbox}
\end{tcolorbox}
\newpage
\begin{tcolorbox}[enhanced,width=16cm,height=14cm,colback=white,colframe=ccal,boxrule=2pt,arc=0mm,
left=6mm,right=6mm,top=0mm,bottom=4mm,sharp corners,
overlay={\fill[ccal] (frame.north west) rectangle ([yshift=-1.25cm]frame.north east);
\node[anchor=west,text=white,font=\sffamily\bfseries\LARGE] at ([xshift=5mm,yshift=-0.62cm]frame.north west) {Question 29 B};
\node[anchor=east,text=white,font=\sffamily\itshape\large] at ([xshift=-5mm,yshift=-0.62cm]frame.north east) {Calculs};}]
\vspace*{1.75cm}
\begin{minipage}[t][6.9cm][c]{\linewidth}\LARGE\centering
\begin{minipage}{0.95\linewidth}\raggedright\linespread{1.25}\selectfont Que vaut $(7+5)\times2$~?\end{minipage}
\end{minipage}

\vfill
\begin{tcolorbox}[enhanced,nobeforeafter,width=7.1cm,height=2.1cm,valign=center,colback=ccal!6,colframe=ccal,arc=2.5mm,boxrule=0.9pt,left=3mm]{\Large\bfseries\color{ccal}a)}\quad\begin{minipage}[c]{5.6cm}\raggedright\Large $24$\end{minipage}\end{tcolorbox}\hfill\begin{tcolorbox}[enhanced,nobeforeafter,width=7.1cm,height=2.1cm,valign=center,colback=ccal!6,colframe=ccal,arc=2.5mm,boxrule=0.9pt,left=3mm]{\Large\bfseries\color{ccal}b)}\quad\begin{minipage}[c]{5.6cm}\raggedright\Large $14$\end{minipage}\end{tcolorbox}\par\vspace{2.5mm}
\begin{tcolorbox}[enhanced,nobeforeafter,width=7.1cm,height=2.1cm,valign=center,colback=ccal!6,colframe=ccal,arc=2.5mm,boxrule=0.9pt,left=3mm]{\Large\bfseries\color{ccal}c)}\quad\begin{minipage}[c]{5.6cm}\raggedright\Large $19$\end{minipage}\end{tcolorbox}\hfill\begin{tcolorbox}[enhanced,nobeforeafter,width=7.1cm,height=2.1cm,valign=center,colback=ccal!6,colframe=ccal,arc=2.5mm,boxrule=0.9pt,left=3mm]{\Large\bfseries\color{ccal}d)}\quad\begin{minipage}[c]{5.6cm}\raggedright\Large $17$\end{minipage}\end{tcolorbox}
\end{tcolorbox}
\newpage
\begin{tcolorbox}[enhanced,width=16cm,height=14cm,colback=white,colframe=ccal,boxrule=2pt,arc=0mm,
left=6mm,right=6mm,top=0mm,bottom=4mm,sharp corners,
overlay={\fill[ccal] (frame.north west) rectangle ([yshift=-1.25cm]frame.north east);
\node[anchor=west,text=white,font=\sffamily\bfseries\LARGE] at ([xshift=5mm,yshift=-0.62cm]frame.north west) {Question 30 B};
\node[anchor=east,text=white,font=\sffamily\itshape\large] at ([xshift=-5mm,yshift=-0.62cm]frame.north east) {Calculs};}]
\vspace*{1.75cm}
\begin{minipage}[t][6.9cm][c]{\linewidth}\LARGE\centering
\begin{minipage}{0.95\linewidth}\raggedright\linespread{1.25}\selectfont Que vaut $18\div3+6$~?\end{minipage}
\end{minipage}

\vfill
\begin{tcolorbox}[enhanced,nobeforeafter,width=7.1cm,height=2.1cm,valign=center,colback=ccal!6,colframe=ccal,arc=2.5mm,boxrule=0.9pt,left=3mm]{\Large\bfseries\color{ccal}a)}\quad\begin{minipage}[c]{5.6cm}\raggedright\Large $12$\end{minipage}\end{tcolorbox}\hfill\begin{tcolorbox}[enhanced,nobeforeafter,width=7.1cm,height=2.1cm,valign=center,colback=ccal!6,colframe=ccal,arc=2.5mm,boxrule=0.9pt,left=3mm]{\Large\bfseries\color{ccal}b)}\quad\begin{minipage}[c]{5.6cm}\raggedright\Large $9$\end{minipage}\end{tcolorbox}\par\vspace{2.5mm}
\begin{tcolorbox}[enhanced,nobeforeafter,width=7.1cm,height=2.1cm,valign=center,colback=ccal!6,colframe=ccal,arc=2.5mm,boxrule=0.9pt,left=3mm]{\Large\bfseries\color{ccal}c)}\quad\begin{minipage}[c]{5.6cm}\raggedright\Large $27$\end{minipage}\end{tcolorbox}\hfill\begin{tcolorbox}[enhanced,nobeforeafter,width=7.1cm,height=2.1cm,valign=center,colback=ccal!6,colframe=ccal,arc=2.5mm,boxrule=0.9pt,left=3mm]{\Large\bfseries\color{ccal}d)}\quad\begin{minipage}[c]{5.6cm}\raggedright\Large $2$\end{minipage}\end{tcolorbox}
\end{tcolorbox}
\newpage
\begin{tcolorbox}[enhanced,width=16cm,height=14cm,colback=white,colframe=ccal,boxrule=2pt,arc=0mm,
left=6mm,right=6mm,top=0mm,bottom=4mm,sharp corners,
overlay={\fill[ccal] (frame.north west) rectangle ([yshift=-1.25cm]frame.north east);
\node[anchor=west,text=white,font=\sffamily\bfseries\LARGE] at ([xshift=5mm,yshift=-0.62cm]frame.north west) {Question 31 B};
\node[anchor=east,text=white,font=\sffamily\itshape\large] at ([xshift=-5mm,yshift=-0.62cm]frame.north east) {Calculs};}]
\vspace*{1.75cm}
\begin{minipage}[t][6.9cm][c]{\linewidth}\LARGE\centering
\begin{minipage}{0.95\linewidth}\raggedright\linespread{1.25}\selectfont Que vaut $15\div6$~?\end{minipage}
\end{minipage}

\vfill
\begin{tcolorbox}[enhanced,nobeforeafter,width=7.1cm,height=2.1cm,valign=center,colback=ccal!6,colframe=ccal,arc=2.5mm,boxrule=0.9pt,left=3mm]{\Large\bfseries\color{ccal}a)}\quad\begin{minipage}[c]{5.6cm}\raggedright\Large $2{,}5$\end{minipage}\end{tcolorbox}\hfill\begin{tcolorbox}[enhanced,nobeforeafter,width=7.1cm,height=2.1cm,valign=center,colback=ccal!6,colframe=ccal,arc=2.5mm,boxrule=0.9pt,left=3mm]{\Large\bfseries\color{ccal}b)}\quad\begin{minipage}[c]{5.6cm}\raggedright\Large $0{,}4$\end{minipage}\end{tcolorbox}\par\vspace{2.5mm}
\begin{tcolorbox}[enhanced,nobeforeafter,width=7.1cm,height=2.1cm,valign=center,colback=ccal!6,colframe=ccal,arc=2.5mm,boxrule=0.9pt,left=3mm]{\Large\bfseries\color{ccal}c)}\quad\begin{minipage}[c]{5.6cm}\raggedright\Large $2{,}3$\end{minipage}\end{tcolorbox}\hfill\begin{tcolorbox}[enhanced,nobeforeafter,width=7.1cm,height=2.1cm,valign=center,colback=ccal!6,colframe=ccal,arc=2.5mm,boxrule=0.9pt,left=3mm]{\Large\bfseries\color{ccal}d)}\quad\begin{minipage}[c]{5.6cm}\raggedright\Large $25$\end{minipage}\end{tcolorbox}
\end{tcolorbox}
\newpage
\begin{tcolorbox}[enhanced,width=16cm,height=14cm,colback=white,colframe=ccal,boxrule=2pt,arc=0mm,
left=6mm,right=6mm,top=0mm,bottom=4mm,sharp corners,
overlay={\fill[ccal] (frame.north west) rectangle ([yshift=-1.25cm]frame.north east);
\node[anchor=west,text=white,font=\sffamily\bfseries\LARGE] at ([xshift=5mm,yshift=-0.62cm]frame.north west) {Question 32 B};
\node[anchor=east,text=white,font=\sffamily\itshape\large] at ([xshift=-5mm,yshift=-0.62cm]frame.north east) {Calculs};}]
\vspace*{1.75cm}
\begin{minipage}[t][6.9cm][c]{\linewidth}\LARGE\centering
\begin{minipage}{0.95\linewidth}\raggedright\linespread{1.25}\selectfont Que vaut $3\div5$~?\end{minipage}
\end{minipage}

\vfill
\begin{tcolorbox}[enhanced,nobeforeafter,width=7.1cm,height=2.1cm,valign=center,colback=ccal!6,colframe=ccal,arc=2.5mm,boxrule=0.9pt,left=3mm]{\Large\bfseries\color{ccal}a)}\quad\begin{minipage}[c]{5.6cm}\raggedright\Large $6$\end{minipage}\end{tcolorbox}\hfill\begin{tcolorbox}[enhanced,nobeforeafter,width=7.1cm,height=2.1cm,valign=center,colback=ccal!6,colframe=ccal,arc=2.5mm,boxrule=0.9pt,left=3mm]{\Large\bfseries\color{ccal}b)}\quad\begin{minipage}[c]{5.6cm}\raggedright\Large $0{,}6$\end{minipage}\end{tcolorbox}\par\vspace{2.5mm}
\begin{tcolorbox}[enhanced,nobeforeafter,width=7.1cm,height=2.1cm,valign=center,colback=ccal!6,colframe=ccal,arc=2.5mm,boxrule=0.9pt,left=3mm]{\Large\bfseries\color{ccal}c)}\quad\begin{minipage}[c]{5.6cm}\raggedright\Large $\dfrac{5}{3}$\end{minipage}\end{tcolorbox}\hfill\begin{tcolorbox}[enhanced,nobeforeafter,width=7.1cm,height=2.1cm,valign=center,colback=ccal!6,colframe=ccal,arc=2.5mm,boxrule=0.9pt,left=3mm]{\Large\bfseries\color{ccal}d)}\quad\begin{minipage}[c]{5.6cm}\raggedright\Large $0{,}3$\end{minipage}\end{tcolorbox}
\end{tcolorbox}
\newpage
\begin{tcolorbox}[enhanced,width=16cm,height=14cm,colback=white,colframe=cfra,boxrule=2pt,arc=0mm,
left=6mm,right=6mm,top=0mm,bottom=4mm,sharp corners,
overlay={\fill[cfra] (frame.north west) rectangle ([yshift=-1.25cm]frame.north east);
\node[anchor=west,text=white,font=\sffamily\bfseries\LARGE] at ([xshift=5mm,yshift=-0.62cm]frame.north west) {Question 33 B};
\node[anchor=east,text=white,font=\sffamily\itshape\large] at ([xshift=-5mm,yshift=-0.62cm]frame.north east) {Fractions};}]
\vspace*{1.75cm}
\begin{minipage}[t][6.9cm][c]{\linewidth}\LARGE\centering
\begin{minipage}{0.95\linewidth}\raggedright\linespread{1.25}\selectfont Combien vaut $\dfrac{5}{6}$ de $42$~?\end{minipage}
\end{minipage}

\vfill
\begin{tcolorbox}[enhanced,nobeforeafter,width=7.1cm,height=2.1cm,valign=center,colback=cfra!6,colframe=cfra,arc=2.5mm,boxrule=0.9pt,left=3mm]{\Large\bfseries\color{cfra}a)}\quad\begin{minipage}[c]{5.6cm}\raggedright\Large $210$\end{minipage}\end{tcolorbox}\hfill\begin{tcolorbox}[enhanced,nobeforeafter,width=7.1cm,height=2.1cm,valign=center,colback=cfra!6,colframe=cfra,arc=2.5mm,boxrule=0.9pt,left=3mm]{\Large\bfseries\color{cfra}b)}\quad\begin{minipage}[c]{5.6cm}\raggedright\Large $7$\end{minipage}\end{tcolorbox}\par\vspace{2.5mm}
\begin{tcolorbox}[enhanced,nobeforeafter,width=7.1cm,height=2.1cm,valign=center,colback=cfra!6,colframe=cfra,arc=2.5mm,boxrule=0.9pt,left=3mm]{\Large\bfseries\color{cfra}c)}\quad\begin{minipage}[c]{5.6cm}\raggedright\Large $36$\end{minipage}\end{tcolorbox}\hfill\begin{tcolorbox}[enhanced,nobeforeafter,width=7.1cm,height=2.1cm,valign=center,colback=cfra!6,colframe=cfra,arc=2.5mm,boxrule=0.9pt,left=3mm]{\Large\bfseries\color{cfra}d)}\quad\begin{minipage}[c]{5.6cm}\raggedright\Large $35$\end{minipage}\end{tcolorbox}
\end{tcolorbox}
\newpage
\begin{tcolorbox}[enhanced,width=16cm,height=14cm,colback=white,colframe=cfra,boxrule=2pt,arc=0mm,
left=6mm,right=6mm,top=0mm,bottom=4mm,sharp corners,
overlay={\fill[cfra] (frame.north west) rectangle ([yshift=-1.25cm]frame.north east);
\node[anchor=west,text=white,font=\sffamily\bfseries\LARGE] at ([xshift=5mm,yshift=-0.62cm]frame.north west) {Question 34 B};
\node[anchor=east,text=white,font=\sffamily\itshape\large] at ([xshift=-5mm,yshift=-0.62cm]frame.north east) {Fractions};}]
\vspace*{1.75cm}
\begin{minipage}[t][6.9cm][c]{\linewidth}\LARGE\centering
\begin{minipage}{0.95\linewidth}\raggedright\linespread{1.25}\selectfont Combien vaut $\dfrac{3}{8}$ de $40$~?\end{minipage}
\end{minipage}

\vfill
\begin{tcolorbox}[enhanced,nobeforeafter,width=7.1cm,height=2.1cm,valign=center,colback=cfra!6,colframe=cfra,arc=2.5mm,boxrule=0.9pt,left=3mm]{\Large\bfseries\color{cfra}a)}\quad\begin{minipage}[c]{5.6cm}\raggedright\Large $15$\end{minipage}\end{tcolorbox}\hfill\begin{tcolorbox}[enhanced,nobeforeafter,width=7.1cm,height=2.1cm,valign=center,colback=cfra!6,colframe=cfra,arc=2.5mm,boxrule=0.9pt,left=3mm]{\Large\bfseries\color{cfra}b)}\quad\begin{minipage}[c]{5.6cm}\raggedright\Large $5$\end{minipage}\end{tcolorbox}\par\vspace{2.5mm}
\begin{tcolorbox}[enhanced,nobeforeafter,width=7.1cm,height=2.1cm,valign=center,colback=cfra!6,colframe=cfra,arc=2.5mm,boxrule=0.9pt,left=3mm]{\Large\bfseries\color{cfra}c)}\quad\begin{minipage}[c]{5.6cm}\raggedright\Large $25$\end{minipage}\end{tcolorbox}\hfill\begin{tcolorbox}[enhanced,nobeforeafter,width=7.1cm,height=2.1cm,valign=center,colback=cfra!6,colframe=cfra,arc=2.5mm,boxrule=0.9pt,left=3mm]{\Large\bfseries\color{cfra}d)}\quad\begin{minipage}[c]{5.6cm}\raggedright\Large $120$\end{minipage}\end{tcolorbox}
\end{tcolorbox}
\newpage
\begin{tcolorbox}[enhanced,width=16cm,height=14cm,colback=white,colframe=cfra,boxrule=2pt,arc=0mm,
left=6mm,right=6mm,top=0mm,bottom=4mm,sharp corners,
overlay={\fill[cfra] (frame.north west) rectangle ([yshift=-1.25cm]frame.north east);
\node[anchor=west,text=white,font=\sffamily\bfseries\LARGE] at ([xshift=5mm,yshift=-0.62cm]frame.north west) {Question 35 B};
\node[anchor=east,text=white,font=\sffamily\itshape\large] at ([xshift=-5mm,yshift=-0.62cm]frame.north east) {Fractions};}]
\vspace*{1.75cm}
\begin{minipage}[t][6.9cm][c]{\linewidth}\LARGE\centering
\begin{minipage}{0.95\linewidth}\raggedright\linespread{1.25}\selectfont Compléter~: $\dfrac{5}{6}=\dfrac{?}{18}$\end{minipage}
\end{minipage}

\vfill
\begin{tcolorbox}[enhanced,nobeforeafter,width=7.1cm,height=2.1cm,valign=center,colback=cfra!6,colframe=cfra,arc=2.5mm,boxrule=0.9pt,left=3mm]{\Large\bfseries\color{cfra}a)}\quad\begin{minipage}[c]{5.6cm}\raggedright\Large $17$\end{minipage}\end{tcolorbox}\hfill\begin{tcolorbox}[enhanced,nobeforeafter,width=7.1cm,height=2.1cm,valign=center,colback=cfra!6,colframe=cfra,arc=2.5mm,boxrule=0.9pt,left=3mm]{\Large\bfseries\color{cfra}b)}\quad\begin{minipage}[c]{5.6cm}\raggedright\Large $90$\end{minipage}\end{tcolorbox}\par\vspace{2.5mm}
\begin{tcolorbox}[enhanced,nobeforeafter,width=7.1cm,height=2.1cm,valign=center,colback=cfra!6,colframe=cfra,arc=2.5mm,boxrule=0.9pt,left=3mm]{\Large\bfseries\color{cfra}c)}\quad\begin{minipage}[c]{5.6cm}\raggedright\Large $3$\end{minipage}\end{tcolorbox}\hfill\begin{tcolorbox}[enhanced,nobeforeafter,width=7.1cm,height=2.1cm,valign=center,colback=cfra!6,colframe=cfra,arc=2.5mm,boxrule=0.9pt,left=3mm]{\Large\bfseries\color{cfra}d)}\quad\begin{minipage}[c]{5.6cm}\raggedright\Large $15$\end{minipage}\end{tcolorbox}
\end{tcolorbox}
\newpage
\begin{tcolorbox}[enhanced,width=16cm,height=14cm,colback=white,colframe=cfra,boxrule=2pt,arc=0mm,
left=6mm,right=6mm,top=0mm,bottom=4mm,sharp corners,
overlay={\fill[cfra] (frame.north west) rectangle ([yshift=-1.25cm]frame.north east);
\node[anchor=west,text=white,font=\sffamily\bfseries\LARGE] at ([xshift=5mm,yshift=-0.62cm]frame.north west) {Question 36 B};
\node[anchor=east,text=white,font=\sffamily\itshape\large] at ([xshift=-5mm,yshift=-0.62cm]frame.north east) {Fractions};}]
\vspace*{1.75cm}
\begin{minipage}[t][6.9cm][c]{\linewidth}\LARGE\centering
\begin{minipage}{0.95\linewidth}\raggedright\linespread{1.25}\selectfont Compléter~: $\dfrac{1}{7}=\dfrac{?}{21}$\end{minipage}
\end{minipage}

\vfill
\begin{tcolorbox}[enhanced,nobeforeafter,width=7.1cm,height=2.1cm,valign=center,colback=cfra!6,colframe=cfra,arc=2.5mm,boxrule=0.9pt,left=3mm]{\Large\bfseries\color{cfra}a)}\quad\begin{minipage}[c]{5.6cm}\raggedright\Large $15$\end{minipage}\end{tcolorbox}\hfill\begin{tcolorbox}[enhanced,nobeforeafter,width=7.1cm,height=2.1cm,valign=center,colback=cfra!6,colframe=cfra,arc=2.5mm,boxrule=0.9pt,left=3mm]{\Large\bfseries\color{cfra}b)}\quad\begin{minipage}[c]{5.6cm}\raggedright\Large $4$\end{minipage}\end{tcolorbox}\par\vspace{2.5mm}
\begin{tcolorbox}[enhanced,nobeforeafter,width=7.1cm,height=2.1cm,valign=center,colback=cfra!6,colframe=cfra,arc=2.5mm,boxrule=0.9pt,left=3mm]{\Large\bfseries\color{cfra}c)}\quad\begin{minipage}[c]{5.6cm}\raggedright\Large $3$\end{minipage}\end{tcolorbox}\hfill\begin{tcolorbox}[enhanced,nobeforeafter,width=7.1cm,height=2.1cm,valign=center,colback=cfra!6,colframe=cfra,arc=2.5mm,boxrule=0.9pt,left=3mm]{\Large\bfseries\color{cfra}d)}\quad\begin{minipage}[c]{5.6cm}\raggedright\Large $21$\end{minipage}\end{tcolorbox}
\end{tcolorbox}
\newpage
\begin{tcolorbox}[enhanced,width=16cm,height=14cm,colback=white,colframe=cfra,boxrule=2pt,arc=0mm,
left=6mm,right=6mm,top=0mm,bottom=4mm,sharp corners,
overlay={\fill[cfra] (frame.north west) rectangle ([yshift=-1.25cm]frame.north east);
\node[anchor=west,text=white,font=\sffamily\bfseries\LARGE] at ([xshift=5mm,yshift=-0.62cm]frame.north west) {Question 37 B};
\node[anchor=east,text=white,font=\sffamily\itshape\large] at ([xshift=-5mm,yshift=-0.62cm]frame.north east) {Fractions};}]
\vspace*{1.75cm}
\begin{minipage}[t][6.9cm][c]{\linewidth}\LARGE\centering
\begin{minipage}{0.95\linewidth}\raggedright\linespread{1.25}\selectfont Quelle est l'écriture décimale de $\dfrac{7}{5}$~?\end{minipage}
\end{minipage}

\vfill
\begin{tcolorbox}[enhanced,nobeforeafter,width=7.1cm,height=2.1cm,valign=center,colback=cfra!6,colframe=cfra,arc=2.5mm,boxrule=0.9pt,left=3mm]{\Large\bfseries\color{cfra}a)}\quad\begin{minipage}[c]{5.6cm}\raggedright\Large $14$\end{minipage}\end{tcolorbox}\hfill\begin{tcolorbox}[enhanced,nobeforeafter,width=7.1cm,height=2.1cm,valign=center,colback=cfra!6,colframe=cfra,arc=2.5mm,boxrule=0.9pt,left=3mm]{\Large\bfseries\color{cfra}b)}\quad\begin{minipage}[c]{5.6cm}\raggedright\Large $\dfrac{5}{7}$\end{minipage}\end{tcolorbox}\par\vspace{2.5mm}
\begin{tcolorbox}[enhanced,nobeforeafter,width=7.1cm,height=2.1cm,valign=center,colback=cfra!6,colframe=cfra,arc=2.5mm,boxrule=0.9pt,left=3mm]{\Large\bfseries\color{cfra}c)}\quad\begin{minipage}[c]{5.6cm}\raggedright\Large $7{,}5$\end{minipage}\end{tcolorbox}\hfill\begin{tcolorbox}[enhanced,nobeforeafter,width=7.1cm,height=2.1cm,valign=center,colback=cfra!6,colframe=cfra,arc=2.5mm,boxrule=0.9pt,left=3mm]{\Large\bfseries\color{cfra}d)}\quad\begin{minipage}[c]{5.6cm}\raggedright\Large $1{,}4$\end{minipage}\end{tcolorbox}
\end{tcolorbox}
\newpage
\begin{tcolorbox}[enhanced,width=16cm,height=14cm,colback=white,colframe=cfra,boxrule=2pt,arc=0mm,
left=6mm,right=6mm,top=0mm,bottom=4mm,sharp corners,
overlay={\fill[cfra] (frame.north west) rectangle ([yshift=-1.25cm]frame.north east);
\node[anchor=west,text=white,font=\sffamily\bfseries\LARGE] at ([xshift=5mm,yshift=-0.62cm]frame.north west) {Question 38 B};
\node[anchor=east,text=white,font=\sffamily\itshape\large] at ([xshift=-5mm,yshift=-0.62cm]frame.north east) {Fractions};}]
\vspace*{1.75cm}
\begin{minipage}[t][6.9cm][c]{\linewidth}\LARGE\centering
\begin{minipage}{0.95\linewidth}\raggedright\linespread{1.25}\selectfont Quelle est l'écriture décimale de $\dfrac{1}{8}$~?\end{minipage}
\end{minipage}

\vfill
\begin{tcolorbox}[enhanced,nobeforeafter,width=7.1cm,height=2.1cm,valign=center,colback=cfra!6,colframe=cfra,arc=2.5mm,boxrule=0.9pt,left=3mm]{\Large\bfseries\color{cfra}a)}\quad\begin{minipage}[c]{5.6cm}\raggedright\Large $8$\end{minipage}\end{tcolorbox}\hfill\begin{tcolorbox}[enhanced,nobeforeafter,width=7.1cm,height=2.1cm,valign=center,colback=cfra!6,colframe=cfra,arc=2.5mm,boxrule=0.9pt,left=3mm]{\Large\bfseries\color{cfra}b)}\quad\begin{minipage}[c]{5.6cm}\raggedright\Large $1{,}25$\end{minipage}\end{tcolorbox}\par\vspace{2.5mm}
\begin{tcolorbox}[enhanced,nobeforeafter,width=7.1cm,height=2.1cm,valign=center,colback=cfra!6,colframe=cfra,arc=2.5mm,boxrule=0.9pt,left=3mm]{\Large\bfseries\color{cfra}c)}\quad\begin{minipage}[c]{5.6cm}\raggedright\Large $0{,}125$\end{minipage}\end{tcolorbox}\hfill\begin{tcolorbox}[enhanced,nobeforeafter,width=7.1cm,height=2.1cm,valign=center,colback=cfra!6,colframe=cfra,arc=2.5mm,boxrule=0.9pt,left=3mm]{\Large\bfseries\color{cfra}d)}\quad\begin{minipage}[c]{5.6cm}\raggedright\Large $1{,}8$\end{minipage}\end{tcolorbox}
\end{tcolorbox}
\newpage
\begin{tcolorbox}[enhanced,width=16cm,height=14cm,colback=white,colframe=cfra,boxrule=2pt,arc=0mm,
left=6mm,right=6mm,top=0mm,bottom=4mm,sharp corners,
overlay={\fill[cfra] (frame.north west) rectangle ([yshift=-1.25cm]frame.north east);
\node[anchor=west,text=white,font=\sffamily\bfseries\LARGE] at ([xshift=5mm,yshift=-0.62cm]frame.north west) {Question 39 B};
\node[anchor=east,text=white,font=\sffamily\itshape\large] at ([xshift=-5mm,yshift=-0.62cm]frame.north east) {Fractions};}]
\vspace*{1.75cm}
\begin{minipage}[t][6.9cm][c]{\linewidth}\LARGE\centering
\begin{minipage}{0.95\linewidth}\raggedright\linespread{1.25}\selectfont Laquelle de ces fractions est plus grande que $1$~?\end{minipage}
\end{minipage}

\vfill
\begin{tcolorbox}[enhanced,nobeforeafter,width=7.1cm,height=2.1cm,valign=center,colback=cfra!6,colframe=cfra,arc=2.5mm,boxrule=0.9pt,left=3mm]{\Large\bfseries\color{cfra}a)}\quad\begin{minipage}[c]{5.6cm}\raggedright\Large $\dfrac{5}{6}$\end{minipage}\end{tcolorbox}\hfill\begin{tcolorbox}[enhanced,nobeforeafter,width=7.1cm,height=2.1cm,valign=center,colback=cfra!6,colframe=cfra,arc=2.5mm,boxrule=0.9pt,left=3mm]{\Large\bfseries\color{cfra}b)}\quad\begin{minipage}[c]{5.6cm}\raggedright\Large $\dfrac{11}{6}$\end{minipage}\end{tcolorbox}\par\vspace{2.5mm}
\begin{tcolorbox}[enhanced,nobeforeafter,width=7.1cm,height=2.1cm,valign=center,colback=cfra!6,colframe=cfra,arc=2.5mm,boxrule=0.9pt,left=3mm]{\Large\bfseries\color{cfra}c)}\quad\begin{minipage}[c]{5.6cm}\raggedright\Large $\dfrac{3}{3}$\end{minipage}\end{tcolorbox}\hfill\begin{tcolorbox}[enhanced,nobeforeafter,width=7.1cm,height=2.1cm,valign=center,colback=cfra!6,colframe=cfra,arc=2.5mm,boxrule=0.9pt,left=3mm]{\Large\bfseries\color{cfra}d)}\quad\begin{minipage}[c]{5.6cm}\raggedright\Large $\dfrac{6}{11}$\end{minipage}\end{tcolorbox}
\end{tcolorbox}
\newpage
\begin{tcolorbox}[enhanced,width=16cm,height=14cm,colback=white,colframe=cfra,boxrule=2pt,arc=0mm,
left=6mm,right=6mm,top=0mm,bottom=4mm,sharp corners,
overlay={\fill[cfra] (frame.north west) rectangle ([yshift=-1.25cm]frame.north east);
\node[anchor=west,text=white,font=\sffamily\bfseries\LARGE] at ([xshift=5mm,yshift=-0.62cm]frame.north west) {Question 40 B};
\node[anchor=east,text=white,font=\sffamily\itshape\large] at ([xshift=-5mm,yshift=-0.62cm]frame.north east) {Fractions};}]
\vspace*{1.75cm}
\begin{minipage}[t][6.9cm][c]{\linewidth}\LARGE\centering
\begin{minipage}{0.95\linewidth}\raggedright\linespread{1.25}\selectfont Laquelle de ces fractions est plus grande que $1$~?\end{minipage}
\end{minipage}

\vfill
\begin{tcolorbox}[enhanced,nobeforeafter,width=7.1cm,height=2.1cm,valign=center,colback=cfra!6,colframe=cfra,arc=2.5mm,boxrule=0.9pt,left=3mm]{\Large\bfseries\color{cfra}a)}\quad\begin{minipage}[c]{5.6cm}\raggedright\Large $\dfrac{5}{3}$\end{minipage}\end{tcolorbox}\hfill\begin{tcolorbox}[enhanced,nobeforeafter,width=7.1cm,height=2.1cm,valign=center,colback=cfra!6,colframe=cfra,arc=2.5mm,boxrule=0.9pt,left=3mm]{\Large\bfseries\color{cfra}b)}\quad\begin{minipage}[c]{5.6cm}\raggedright\Large $\dfrac{3}{5}$\end{minipage}\end{tcolorbox}\par\vspace{2.5mm}
\begin{tcolorbox}[enhanced,nobeforeafter,width=7.1cm,height=2.1cm,valign=center,colback=cfra!6,colframe=cfra,arc=2.5mm,boxrule=0.9pt,left=3mm]{\Large\bfseries\color{cfra}c)}\quad\begin{minipage}[c]{5.6cm}\raggedright\Large $\dfrac{2}{3}$\end{minipage}\end{tcolorbox}\hfill\begin{tcolorbox}[enhanced,nobeforeafter,width=7.1cm,height=2.1cm,valign=center,colback=cfra!6,colframe=cfra,arc=2.5mm,boxrule=0.9pt,left=3mm]{\Large\bfseries\color{cfra}d)}\quad\begin{minipage}[c]{5.6cm}\raggedright\Large $\dfrac{7}{7}$\end{minipage}\end{tcolorbox}
\end{tcolorbox}
\newpage
\begin{tcolorbox}[enhanced,width=16cm,height=14cm,colback=white,colframe=cfra,boxrule=2pt,arc=0mm,
left=6mm,right=6mm,top=0mm,bottom=4mm,sharp corners,
overlay={\fill[cfra] (frame.north west) rectangle ([yshift=-1.25cm]frame.north east);
\node[anchor=west,text=white,font=\sffamily\bfseries\LARGE] at ([xshift=5mm,yshift=-0.62cm]frame.north west) {Question 41 B};
\node[anchor=east,text=white,font=\sffamily\itshape\large] at ([xshift=-5mm,yshift=-0.62cm]frame.north east) {Fractions};}]
\vspace*{1.75cm}
\begin{minipage}[t][6.9cm][c]{\linewidth}\LARGE\centering
\begin{minipage}{0.95\linewidth}\raggedright\linespread{1.25}\selectfont Quel nombre, multiplié par $9$, donne $2$~?\end{minipage}
\end{minipage}

\vfill
\begin{tcolorbox}[enhanced,nobeforeafter,width=7.1cm,height=2.1cm,valign=center,colback=cfra!6,colframe=cfra,arc=2.5mm,boxrule=0.9pt,left=3mm]{\Large\bfseries\color{cfra}a)}\quad\begin{minipage}[c]{5.6cm}\raggedright\Large $18$\end{minipage}\end{tcolorbox}\hfill\begin{tcolorbox}[enhanced,nobeforeafter,width=7.1cm,height=2.1cm,valign=center,colback=cfra!6,colframe=cfra,arc=2.5mm,boxrule=0.9pt,left=3mm]{\Large\bfseries\color{cfra}b)}\quad\begin{minipage}[c]{5.6cm}\raggedright\Large $7$\end{minipage}\end{tcolorbox}\par\vspace{2.5mm}
\begin{tcolorbox}[enhanced,nobeforeafter,width=7.1cm,height=2.1cm,valign=center,colback=cfra!6,colframe=cfra,arc=2.5mm,boxrule=0.9pt,left=3mm]{\Large\bfseries\color{cfra}c)}\quad\begin{minipage}[c]{5.6cm}\raggedright\Large $\dfrac{9}{2}$\end{minipage}\end{tcolorbox}\hfill\begin{tcolorbox}[enhanced,nobeforeafter,width=7.1cm,height=2.1cm,valign=center,colback=cfra!6,colframe=cfra,arc=2.5mm,boxrule=0.9pt,left=3mm]{\Large\bfseries\color{cfra}d)}\quad\begin{minipage}[c]{5.6cm}\raggedright\Large $\dfrac{2}{9}$\end{minipage}\end{tcolorbox}
\end{tcolorbox}
\newpage
\begin{tcolorbox}[enhanced,width=16cm,height=14cm,colback=white,colframe=cfra,boxrule=2pt,arc=0mm,
left=6mm,right=6mm,top=0mm,bottom=4mm,sharp corners,
overlay={\fill[cfra] (frame.north west) rectangle ([yshift=-1.25cm]frame.north east);
\node[anchor=west,text=white,font=\sffamily\bfseries\LARGE] at ([xshift=5mm,yshift=-0.62cm]frame.north west) {Question 42 B};
\node[anchor=east,text=white,font=\sffamily\itshape\large] at ([xshift=-5mm,yshift=-0.62cm]frame.north east) {Fractions};}]
\vspace*{1.75cm}
\begin{minipage}[t][6.9cm][c]{\linewidth}\LARGE\centering
\begin{minipage}{0.95\linewidth}\raggedright\linespread{1.25}\selectfont Quel nombre, multiplié par $3$, donne $8$~?\end{minipage}
\end{minipage}

\vfill
\begin{tcolorbox}[enhanced,nobeforeafter,width=7.1cm,height=2.1cm,valign=center,colback=cfra!6,colframe=cfra,arc=2.5mm,boxrule=0.9pt,left=3mm]{\Large\bfseries\color{cfra}a)}\quad\begin{minipage}[c]{5.6cm}\raggedright\Large $24$\end{minipage}\end{tcolorbox}\hfill\begin{tcolorbox}[enhanced,nobeforeafter,width=7.1cm,height=2.1cm,valign=center,colback=cfra!6,colframe=cfra,arc=2.5mm,boxrule=0.9pt,left=3mm]{\Large\bfseries\color{cfra}b)}\quad\begin{minipage}[c]{5.6cm}\raggedright\Large $\dfrac{8}{3}$\end{minipage}\end{tcolorbox}\par\vspace{2.5mm}
\begin{tcolorbox}[enhanced,nobeforeafter,width=7.1cm,height=2.1cm,valign=center,colback=cfra!6,colframe=cfra,arc=2.5mm,boxrule=0.9pt,left=3mm]{\Large\bfseries\color{cfra}c)}\quad\begin{minipage}[c]{5.6cm}\raggedright\Large $5$\end{minipage}\end{tcolorbox}\hfill\begin{tcolorbox}[enhanced,nobeforeafter,width=7.1cm,height=2.1cm,valign=center,colback=cfra!6,colframe=cfra,arc=2.5mm,boxrule=0.9pt,left=3mm]{\Large\bfseries\color{cfra}d)}\quad\begin{minipage}[c]{5.6cm}\raggedright\Large $\dfrac{3}{8}$\end{minipage}\end{tcolorbox}
\end{tcolorbox}
\newpage
\begin{tcolorbox}[enhanced,width=16cm,height=14cm,colback=white,colframe=cfra,boxrule=2pt,arc=0mm,
left=6mm,right=6mm,top=0mm,bottom=4mm,sharp corners,
overlay={\fill[cfra] (frame.north west) rectangle ([yshift=-1.25cm]frame.north east);
\node[anchor=west,text=white,font=\sffamily\bfseries\LARGE] at ([xshift=5mm,yshift=-0.62cm]frame.north west) {Question 43 B};
\node[anchor=east,text=white,font=\sffamily\itshape\large] at ([xshift=-5mm,yshift=-0.62cm]frame.north east) {Fractions};}]
\vspace*{1.75cm}
\begin{minipage}[t][6.9cm][c]{\linewidth}\LARGE\centering
\begin{minipage}{0.95\linewidth}\raggedright\linespread{1.25}\selectfont Que vaut $\dfrac{1}{10}+\dfrac{6}{10}$~?\end{minipage}
\end{minipage}

\vfill
\begin{tcolorbox}[enhanced,nobeforeafter,width=7.1cm,height=2.1cm,valign=center,colback=cfra!6,colframe=cfra,arc=2.5mm,boxrule=0.9pt,left=3mm]{\Large\bfseries\color{cfra}a)}\quad\begin{minipage}[c]{5.6cm}\raggedright\Large $\dfrac{7}{10}$\end{minipage}\end{tcolorbox}\hfill\begin{tcolorbox}[enhanced,nobeforeafter,width=7.1cm,height=2.1cm,valign=center,colback=cfra!6,colframe=cfra,arc=2.5mm,boxrule=0.9pt,left=3mm]{\Large\bfseries\color{cfra}b)}\quad\begin{minipage}[c]{5.6cm}\raggedright\Large $\dfrac{7}{20}$\end{minipage}\end{tcolorbox}\par\vspace{2.5mm}
\begin{tcolorbox}[enhanced,nobeforeafter,width=7.1cm,height=2.1cm,valign=center,colback=cfra!6,colframe=cfra,arc=2.5mm,boxrule=0.9pt,left=3mm]{\Large\bfseries\color{cfra}c)}\quad\begin{minipage}[c]{5.6cm}\raggedright\Large $\dfrac{6}{10}$\end{minipage}\end{tcolorbox}\hfill\begin{tcolorbox}[enhanced,nobeforeafter,width=7.1cm,height=2.1cm,valign=center,colback=cfra!6,colframe=cfra,arc=2.5mm,boxrule=0.9pt,left=3mm]{\Large\bfseries\color{cfra}d)}\quad\begin{minipage}[c]{5.6cm}\raggedright\Large $\dfrac{7}{100}$\end{minipage}\end{tcolorbox}
\end{tcolorbox}
\newpage
\begin{tcolorbox}[enhanced,width=16cm,height=14cm,colback=white,colframe=cfra,boxrule=2pt,arc=0mm,
left=6mm,right=6mm,top=0mm,bottom=4mm,sharp corners,
overlay={\fill[cfra] (frame.north west) rectangle ([yshift=-1.25cm]frame.north east);
\node[anchor=west,text=white,font=\sffamily\bfseries\LARGE] at ([xshift=5mm,yshift=-0.62cm]frame.north west) {Question 44 B};
\node[anchor=east,text=white,font=\sffamily\itshape\large] at ([xshift=-5mm,yshift=-0.62cm]frame.north east) {Fractions};}]
\vspace*{1.75cm}
\begin{minipage}[t][6.9cm][c]{\linewidth}\LARGE\centering
\begin{minipage}{0.95\linewidth}\raggedright\linespread{1.25}\selectfont Que vaut $\dfrac{5}{11}+\dfrac{2}{11}$~?\end{minipage}
\end{minipage}

\vfill
\begin{tcolorbox}[enhanced,nobeforeafter,width=7.1cm,height=2.1cm,valign=center,colback=cfra!6,colframe=cfra,arc=2.5mm,boxrule=0.9pt,left=3mm]{\Large\bfseries\color{cfra}a)}\quad\begin{minipage}[c]{5.6cm}\raggedright\Large $\dfrac{7}{22}$\end{minipage}\end{tcolorbox}\hfill\begin{tcolorbox}[enhanced,nobeforeafter,width=7.1cm,height=2.1cm,valign=center,colback=cfra!6,colframe=cfra,arc=2.5mm,boxrule=0.9pt,left=3mm]{\Large\bfseries\color{cfra}b)}\quad\begin{minipage}[c]{5.6cm}\raggedright\Large $\dfrac{10}{11}$\end{minipage}\end{tcolorbox}\par\vspace{2.5mm}
\begin{tcolorbox}[enhanced,nobeforeafter,width=7.1cm,height=2.1cm,valign=center,colback=cfra!6,colframe=cfra,arc=2.5mm,boxrule=0.9pt,left=3mm]{\Large\bfseries\color{cfra}c)}\quad\begin{minipage}[c]{5.6cm}\raggedright\Large $\dfrac{7}{11}$\end{minipage}\end{tcolorbox}\hfill\begin{tcolorbox}[enhanced,nobeforeafter,width=7.1cm,height=2.1cm,valign=center,colback=cfra!6,colframe=cfra,arc=2.5mm,boxrule=0.9pt,left=3mm]{\Large\bfseries\color{cfra}d)}\quad\begin{minipage}[c]{5.6cm}\raggedright\Large $\dfrac{7}{121}$\end{minipage}\end{tcolorbox}
\end{tcolorbox}
\newpage
\begin{tcolorbox}[enhanced,width=16cm,height=14cm,colback=white,colframe=cfra,boxrule=2pt,arc=0mm,
left=6mm,right=6mm,top=0mm,bottom=4mm,sharp corners,
overlay={\fill[cfra] (frame.north west) rectangle ([yshift=-1.25cm]frame.north east);
\node[anchor=west,text=white,font=\sffamily\bfseries\LARGE] at ([xshift=5mm,yshift=-0.62cm]frame.north west) {Question 45 B};
\node[anchor=east,text=white,font=\sffamily\itshape\large] at ([xshift=-5mm,yshift=-0.62cm]frame.north east) {Fractions};}]
\vspace*{1.75cm}
\begin{minipage}[t][6.9cm][c]{\linewidth}\LARGE\centering
\begin{minipage}{0.95\linewidth}\raggedright\linespread{1.25}\selectfont Un gâteau est coupé en $10$ parts égales. On mange $7$ parts. Quelle fraction du gâteau reste-t-il~?\end{minipage}
\end{minipage}

\vfill
\begin{tcolorbox}[enhanced,nobeforeafter,width=7.1cm,height=2.1cm,valign=center,colback=cfra!6,colframe=cfra,arc=2.5mm,boxrule=0.9pt,left=3mm]{\Large\bfseries\color{cfra}a)}\quad\begin{minipage}[c]{5.6cm}\raggedright\Large $\dfrac{7}{10}$\end{minipage}\end{tcolorbox}\hfill\begin{tcolorbox}[enhanced,nobeforeafter,width=7.1cm,height=2.1cm,valign=center,colback=cfra!6,colframe=cfra,arc=2.5mm,boxrule=0.9pt,left=3mm]{\Large\bfseries\color{cfra}b)}\quad\begin{minipage}[c]{5.6cm}\raggedright\Large $\dfrac{7}{3}$\end{minipage}\end{tcolorbox}\par\vspace{2.5mm}
\begin{tcolorbox}[enhanced,nobeforeafter,width=7.1cm,height=2.1cm,valign=center,colback=cfra!6,colframe=cfra,arc=2.5mm,boxrule=0.9pt,left=3mm]{\Large\bfseries\color{cfra}c)}\quad\begin{minipage}[c]{5.6cm}\raggedright\Large $\dfrac{3}{7}$\end{minipage}\end{tcolorbox}\hfill\begin{tcolorbox}[enhanced,nobeforeafter,width=7.1cm,height=2.1cm,valign=center,colback=cfra!6,colframe=cfra,arc=2.5mm,boxrule=0.9pt,left=3mm]{\Large\bfseries\color{cfra}d)}\quad\begin{minipage}[c]{5.6cm}\raggedright\Large $\dfrac{3}{10}$\end{minipage}\end{tcolorbox}
\end{tcolorbox}
\newpage
\begin{tcolorbox}[enhanced,width=16cm,height=14cm,colback=white,colframe=cfra,boxrule=2pt,arc=0mm,
left=6mm,right=6mm,top=0mm,bottom=4mm,sharp corners,
overlay={\fill[cfra] (frame.north west) rectangle ([yshift=-1.25cm]frame.north east);
\node[anchor=west,text=white,font=\sffamily\bfseries\LARGE] at ([xshift=5mm,yshift=-0.62cm]frame.north west) {Question 46 B};
\node[anchor=east,text=white,font=\sffamily\itshape\large] at ([xshift=-5mm,yshift=-0.62cm]frame.north east) {Fractions};}]
\vspace*{1.75cm}
\begin{minipage}[t][6.9cm][c]{\linewidth}\LARGE\centering
\begin{minipage}{0.95\linewidth}\raggedright\linespread{1.25}\selectfont Un gâteau est coupé en $12$ parts égales. On mange $5$ parts. Quelle fraction du gâteau reste-t-il~?\end{minipage}
\end{minipage}

\vfill
\begin{tcolorbox}[enhanced,nobeforeafter,width=7.1cm,height=2.1cm,valign=center,colback=cfra!6,colframe=cfra,arc=2.5mm,boxrule=0.9pt,left=3mm]{\Large\bfseries\color{cfra}a)}\quad\begin{minipage}[c]{5.6cm}\raggedright\Large $\dfrac{7}{12}$\end{minipage}\end{tcolorbox}\hfill\begin{tcolorbox}[enhanced,nobeforeafter,width=7.1cm,height=2.1cm,valign=center,colback=cfra!6,colframe=cfra,arc=2.5mm,boxrule=0.9pt,left=3mm]{\Large\bfseries\color{cfra}b)}\quad\begin{minipage}[c]{5.6cm}\raggedright\Large $\dfrac{7}{5}$\end{minipage}\end{tcolorbox}\par\vspace{2.5mm}
\begin{tcolorbox}[enhanced,nobeforeafter,width=7.1cm,height=2.1cm,valign=center,colback=cfra!6,colframe=cfra,arc=2.5mm,boxrule=0.9pt,left=3mm]{\Large\bfseries\color{cfra}c)}\quad\begin{minipage}[c]{5.6cm}\raggedright\Large $\dfrac{5}{7}$\end{minipage}\end{tcolorbox}\hfill\begin{tcolorbox}[enhanced,nobeforeafter,width=7.1cm,height=2.1cm,valign=center,colback=cfra!6,colframe=cfra,arc=2.5mm,boxrule=0.9pt,left=3mm]{\Large\bfseries\color{cfra}d)}\quad\begin{minipage}[c]{5.6cm}\raggedright\Large $\dfrac{5}{12}$\end{minipage}\end{tcolorbox}
\end{tcolorbox}
\newpage
\begin{tcolorbox}[enhanced,width=16cm,height=14cm,colback=white,colframe=cfra,boxrule=2pt,arc=0mm,
left=6mm,right=6mm,top=0mm,bottom=4mm,sharp corners,
overlay={\fill[cfra] (frame.north west) rectangle ([yshift=-1.25cm]frame.north east);
\node[anchor=west,text=white,font=\sffamily\bfseries\LARGE] at ([xshift=5mm,yshift=-0.62cm]frame.north west) {Question 47 B};
\node[anchor=east,text=white,font=\sffamily\itshape\large] at ([xshift=-5mm,yshift=-0.62cm]frame.north east) {Fractions};}]
\vspace*{1.75cm}
\begin{minipage}[t][6.9cm][c]{\linewidth}\LARGE\centering
\begin{minipage}{0.95\linewidth}\raggedright\linespread{1.25}\selectfont Entre quels nombres entiers consécutifs se trouve $\dfrac{9}{2}$~?\end{minipage}
\end{minipage}

\vfill
\begin{tcolorbox}[enhanced,nobeforeafter,width=7.1cm,height=2.1cm,valign=center,colback=cfra!6,colframe=cfra,arc=2.5mm,boxrule=0.9pt,left=3mm]{\Large\bfseries\color{cfra}a)}\quad\begin{minipage}[c]{5.6cm}\raggedright\Large $4\text{ et }5$\end{minipage}\end{tcolorbox}\hfill\begin{tcolorbox}[enhanced,nobeforeafter,width=7.1cm,height=2.1cm,valign=center,colback=cfra!6,colframe=cfra,arc=2.5mm,boxrule=0.9pt,left=3mm]{\Large\bfseries\color{cfra}b)}\quad\begin{minipage}[c]{5.6cm}\raggedright\Large $6\text{ et }7$\end{minipage}\end{tcolorbox}\par\vspace{2.5mm}
\begin{tcolorbox}[enhanced,nobeforeafter,width=7.1cm,height=2.1cm,valign=center,colback=cfra!6,colframe=cfra,arc=2.5mm,boxrule=0.9pt,left=3mm]{\Large\bfseries\color{cfra}c)}\quad\begin{minipage}[c]{5.6cm}\raggedright\Large $5\text{ et }6$\end{minipage}\end{tcolorbox}\hfill\begin{tcolorbox}[enhanced,nobeforeafter,width=7.1cm,height=2.1cm,valign=center,colback=cfra!6,colframe=cfra,arc=2.5mm,boxrule=0.9pt,left=3mm]{\Large\bfseries\color{cfra}d)}\quad\begin{minipage}[c]{5.6cm}\raggedright\Large $3\text{ et }4$\end{minipage}\end{tcolorbox}
\end{tcolorbox}
\newpage
\begin{tcolorbox}[enhanced,width=16cm,height=14cm,colback=white,colframe=cfra,boxrule=2pt,arc=0mm,
left=6mm,right=6mm,top=0mm,bottom=4mm,sharp corners,
overlay={\fill[cfra] (frame.north west) rectangle ([yshift=-1.25cm]frame.north east);
\node[anchor=west,text=white,font=\sffamily\bfseries\LARGE] at ([xshift=5mm,yshift=-0.62cm]frame.north west) {Question 48 B};
\node[anchor=east,text=white,font=\sffamily\itshape\large] at ([xshift=-5mm,yshift=-0.62cm]frame.north east) {Fractions};}]
\vspace*{1.75cm}
\begin{minipage}[t][6.9cm][c]{\linewidth}\LARGE\centering
\begin{minipage}{0.95\linewidth}\raggedright\linespread{1.25}\selectfont Entre quels nombres entiers consécutifs se trouve $\dfrac{13}{6}$~?\end{minipage}
\end{minipage}

\vfill
\begin{tcolorbox}[enhanced,nobeforeafter,width=7.1cm,height=2.1cm,valign=center,colback=cfra!6,colframe=cfra,arc=2.5mm,boxrule=0.9pt,left=3mm]{\Large\bfseries\color{cfra}a)}\quad\begin{minipage}[c]{5.6cm}\raggedright\Large $3\text{ et }4$\end{minipage}\end{tcolorbox}\hfill\begin{tcolorbox}[enhanced,nobeforeafter,width=7.1cm,height=2.1cm,valign=center,colback=cfra!6,colframe=cfra,arc=2.5mm,boxrule=0.9pt,left=3mm]{\Large\bfseries\color{cfra}b)}\quad\begin{minipage}[c]{5.6cm}\raggedright\Large $4\text{ et }5$\end{minipage}\end{tcolorbox}\par\vspace{2.5mm}
\begin{tcolorbox}[enhanced,nobeforeafter,width=7.1cm,height=2.1cm,valign=center,colback=cfra!6,colframe=cfra,arc=2.5mm,boxrule=0.9pt,left=3mm]{\Large\bfseries\color{cfra}c)}\quad\begin{minipage}[c]{5.6cm}\raggedright\Large $2\text{ et }3$\end{minipage}\end{tcolorbox}\hfill\begin{tcolorbox}[enhanced,nobeforeafter,width=7.1cm,height=2.1cm,valign=center,colback=cfra!6,colframe=cfra,arc=2.5mm,boxrule=0.9pt,left=3mm]{\Large\bfseries\color{cfra}d)}\quad\begin{minipage}[c]{5.6cm}\raggedright\Large $1\text{ et }2$\end{minipage}\end{tcolorbox}
\end{tcolorbox}
\newpage
\begin{tcolorbox}[enhanced,width=16cm,height=14cm,colback=white,colframe=cgeo,boxrule=2pt,arc=0mm,
left=6mm,right=6mm,top=0mm,bottom=4mm,sharp corners,
overlay={\fill[cgeo] (frame.north west) rectangle ([yshift=-1.25cm]frame.north east);
\node[anchor=west,text=white,font=\sffamily\bfseries\LARGE] at ([xshift=5mm,yshift=-0.62cm]frame.north west) {Question 49 B};
\node[anchor=east,text=white,font=\sffamily\itshape\large] at ([xshift=-5mm,yshift=-0.62cm]frame.north east) {Géométrie};}]
\vspace*{1.75cm}
\begin{minipage}[t][6.9cm][c]{\linewidth}\LARGE\centering
\begin{minipage}{0.95\linewidth}\raggedright\linespread{1.25}\selectfont Un angle mesure $90^\circ$. Quelle est sa nature~?\end{minipage}
\end{minipage}

\vfill
\begin{tcolorbox}[enhanced,nobeforeafter,width=7.1cm,height=2.1cm,valign=center,colback=cgeo!6,colframe=cgeo,arc=2.5mm,boxrule=0.9pt,left=3mm]{\Large\bfseries\color{cgeo}a)}\quad\begin{minipage}[c]{5.6cm}\raggedright\Large aigu\end{minipage}\end{tcolorbox}\hfill\begin{tcolorbox}[enhanced,nobeforeafter,width=7.1cm,height=2.1cm,valign=center,colback=cgeo!6,colframe=cgeo,arc=2.5mm,boxrule=0.9pt,left=3mm]{\Large\bfseries\color{cgeo}b)}\quad\begin{minipage}[c]{5.6cm}\raggedright\Large plat\end{minipage}\end{tcolorbox}\par\vspace{2.5mm}
\begin{tcolorbox}[enhanced,nobeforeafter,width=7.1cm,height=2.1cm,valign=center,colback=cgeo!6,colframe=cgeo,arc=2.5mm,boxrule=0.9pt,left=3mm]{\Large\bfseries\color{cgeo}c)}\quad\begin{minipage}[c]{5.6cm}\raggedright\Large droit\end{minipage}\end{tcolorbox}\hfill\begin{tcolorbox}[enhanced,nobeforeafter,width=7.1cm,height=2.1cm,valign=center,colback=cgeo!6,colframe=cgeo,arc=2.5mm,boxrule=0.9pt,left=3mm]{\Large\bfseries\color{cgeo}d)}\quad\begin{minipage}[c]{5.6cm}\raggedright\Large obtus\end{minipage}\end{tcolorbox}
\end{tcolorbox}
\newpage
\begin{tcolorbox}[enhanced,width=16cm,height=14cm,colback=white,colframe=cgeo,boxrule=2pt,arc=0mm,
left=6mm,right=6mm,top=0mm,bottom=4mm,sharp corners,
overlay={\fill[cgeo] (frame.north west) rectangle ([yshift=-1.25cm]frame.north east);
\node[anchor=west,text=white,font=\sffamily\bfseries\LARGE] at ([xshift=5mm,yshift=-0.62cm]frame.north west) {Question 50 B};
\node[anchor=east,text=white,font=\sffamily\itshape\large] at ([xshift=-5mm,yshift=-0.62cm]frame.north east) {Géométrie};}]
\vspace*{1.75cm}
\begin{minipage}[t][6.9cm][c]{\linewidth}\LARGE\centering
\begin{minipage}{0.95\linewidth}\raggedright\linespread{1.25}\selectfont Un angle mesure $180^\circ$. Quelle est sa nature~?\end{minipage}
\end{minipage}

\vfill
\begin{tcolorbox}[enhanced,nobeforeafter,width=7.1cm,height=2.1cm,valign=center,colback=cgeo!6,colframe=cgeo,arc=2.5mm,boxrule=0.9pt,left=3mm]{\Large\bfseries\color{cgeo}a)}\quad\begin{minipage}[c]{5.6cm}\raggedright\Large droit\end{minipage}\end{tcolorbox}\hfill\begin{tcolorbox}[enhanced,nobeforeafter,width=7.1cm,height=2.1cm,valign=center,colback=cgeo!6,colframe=cgeo,arc=2.5mm,boxrule=0.9pt,left=3mm]{\Large\bfseries\color{cgeo}b)}\quad\begin{minipage}[c]{5.6cm}\raggedright\Large plat\end{minipage}\end{tcolorbox}\par\vspace{2.5mm}
\begin{tcolorbox}[enhanced,nobeforeafter,width=7.1cm,height=2.1cm,valign=center,colback=cgeo!6,colframe=cgeo,arc=2.5mm,boxrule=0.9pt,left=3mm]{\Large\bfseries\color{cgeo}c)}\quad\begin{minipage}[c]{5.6cm}\raggedright\Large obtus\end{minipage}\end{tcolorbox}\hfill\begin{tcolorbox}[enhanced,nobeforeafter,width=7.1cm,height=2.1cm,valign=center,colback=cgeo!6,colframe=cgeo,arc=2.5mm,boxrule=0.9pt,left=3mm]{\Large\bfseries\color{cgeo}d)}\quad\begin{minipage}[c]{5.6cm}\raggedright\Large aigu\end{minipage}\end{tcolorbox}
\end{tcolorbox}
\newpage
\begin{tcolorbox}[enhanced,width=16cm,height=14cm,colback=white,colframe=cgeo,boxrule=2pt,arc=0mm,
left=6mm,right=6mm,top=0mm,bottom=4mm,sharp corners,
overlay={\fill[cgeo] (frame.north west) rectangle ([yshift=-1.25cm]frame.north east);
\node[anchor=west,text=white,font=\sffamily\bfseries\LARGE] at ([xshift=5mm,yshift=-0.62cm]frame.north west) {Question 51 B};
\node[anchor=east,text=white,font=\sffamily\itshape\large] at ([xshift=-5mm,yshift=-0.62cm]frame.north east) {Géométrie};}]
\vspace*{1.75cm}
\begin{minipage}[t][6.9cm][c]{\linewidth}\LARGE\centering
\begin{minipage}{0.95\linewidth}\raggedright\linespread{1.25}\selectfont L'angle droit $\widehat{xOy}$ est partagé par la demi-droite $[Oz)$ en deux angles. $\widehat{xOz}=72^\circ$. Combien mesure $\widehat{zOy}$~?\end{minipage}
\end{minipage}

\vfill
\begin{tcolorbox}[enhanced,nobeforeafter,width=7.1cm,height=2.1cm,valign=center,colback=cgeo!6,colframe=cgeo,arc=2.5mm,boxrule=0.9pt,left=3mm]{\Large\bfseries\color{cgeo}a)}\quad\begin{minipage}[c]{5.6cm}\raggedright\Large $18^\circ$\end{minipage}\end{tcolorbox}\hfill\begin{tcolorbox}[enhanced,nobeforeafter,width=7.1cm,height=2.1cm,valign=center,colback=cgeo!6,colframe=cgeo,arc=2.5mm,boxrule=0.9pt,left=3mm]{\Large\bfseries\color{cgeo}b)}\quad\begin{minipage}[c]{5.6cm}\raggedright\Large $108^\circ$\end{minipage}\end{tcolorbox}\par\vspace{2.5mm}
\begin{tcolorbox}[enhanced,nobeforeafter,width=7.1cm,height=2.1cm,valign=center,colback=cgeo!6,colframe=cgeo,arc=2.5mm,boxrule=0.9pt,left=3mm]{\Large\bfseries\color{cgeo}c)}\quad\begin{minipage}[c]{5.6cm}\raggedright\Large $162^\circ$\end{minipage}\end{tcolorbox}\hfill\begin{tcolorbox}[enhanced,nobeforeafter,width=7.1cm,height=2.1cm,valign=center,colback=cgeo!6,colframe=cgeo,arc=2.5mm,boxrule=0.9pt,left=3mm]{\Large\bfseries\color{cgeo}d)}\quad\begin{minipage}[c]{5.6cm}\raggedright\Large $72^\circ$\end{minipage}\end{tcolorbox}
\end{tcolorbox}
\newpage
\begin{tcolorbox}[enhanced,width=16cm,height=14cm,colback=white,colframe=cgeo,boxrule=2pt,arc=0mm,
left=6mm,right=6mm,top=0mm,bottom=4mm,sharp corners,
overlay={\fill[cgeo] (frame.north west) rectangle ([yshift=-1.25cm]frame.north east);
\node[anchor=west,text=white,font=\sffamily\bfseries\LARGE] at ([xshift=5mm,yshift=-0.62cm]frame.north west) {Question 52 B};
\node[anchor=east,text=white,font=\sffamily\itshape\large] at ([xshift=-5mm,yshift=-0.62cm]frame.north east) {Géométrie};}]
\vspace*{1.75cm}
\begin{minipage}[t][6.9cm][c]{\linewidth}\LARGE\centering
\begin{minipage}{0.95\linewidth}\raggedright\linespread{1.25}\selectfont L'angle droit $\widehat{xOy}$ est partagé par la demi-droite $[Oz)$ en deux angles. $\widehat{xOz}=45^\circ$. Combien mesure $\widehat{zOy}$~?\end{minipage}
\end{minipage}

\vfill
\begin{tcolorbox}[enhanced,nobeforeafter,width=7.1cm,height=2.1cm,valign=center,colback=cgeo!6,colframe=cgeo,arc=2.5mm,boxrule=0.9pt,left=3mm]{\Large\bfseries\color{cgeo}a)}\quad\begin{minipage}[c]{5.6cm}\raggedright\Large $135^\circ$\end{minipage}\end{tcolorbox}\hfill\begin{tcolorbox}[enhanced,nobeforeafter,width=7.1cm,height=2.1cm,valign=center,colback=cgeo!6,colframe=cgeo,arc=2.5mm,boxrule=0.9pt,left=3mm]{\Large\bfseries\color{cgeo}b)}\quad\begin{minipage}[c]{5.6cm}\raggedright\Large $90^\circ$\end{minipage}\end{tcolorbox}\par\vspace{2.5mm}
\begin{tcolorbox}[enhanced,nobeforeafter,width=7.1cm,height=2.1cm,valign=center,colback=cgeo!6,colframe=cgeo,arc=2.5mm,boxrule=0.9pt,left=3mm]{\Large\bfseries\color{cgeo}c)}\quad\begin{minipage}[c]{5.6cm}\raggedright\Large $45^\circ$\end{minipage}\end{tcolorbox}\hfill\begin{tcolorbox}[enhanced,nobeforeafter,width=7.1cm,height=2.1cm,valign=center,colback=cgeo!6,colframe=cgeo,arc=2.5mm,boxrule=0.9pt,left=3mm]{\Large\bfseries\color{cgeo}d)}\quad\begin{minipage}[c]{5.6cm}\raggedright\Large $55^\circ$\end{minipage}\end{tcolorbox}
\end{tcolorbox}
\newpage
\begin{tcolorbox}[enhanced,width=16cm,height=14cm,colback=white,colframe=cgeo,boxrule=2pt,arc=0mm,
left=6mm,right=6mm,top=0mm,bottom=4mm,sharp corners,
overlay={\fill[cgeo] (frame.north west) rectangle ([yshift=-1.25cm]frame.north east);
\node[anchor=west,text=white,font=\sffamily\bfseries\LARGE] at ([xshift=5mm,yshift=-0.62cm]frame.north west) {Question 53 B};
\node[anchor=east,text=white,font=\sffamily\itshape\large] at ([xshift=-5mm,yshift=-0.62cm]frame.north east) {Géométrie};}]
\vspace*{1.75cm}
\begin{minipage}[t][6.9cm][c]{\linewidth}\LARGE\centering
\begin{minipage}{0.95\linewidth}\raggedright\linespread{1.25}\selectfont Les points $A$, $O$ et $B$ sont alignés dans cet ordre et $\widehat{AOC}=38^\circ$. Combien mesure $\widehat{COB}$~?\end{minipage}
\end{minipage}

\vfill
\begin{tcolorbox}[enhanced,nobeforeafter,width=7.1cm,height=2.1cm,valign=center,colback=cgeo!6,colframe=cgeo,arc=2.5mm,boxrule=0.9pt,left=3mm]{\Large\bfseries\color{cgeo}a)}\quad\begin{minipage}[c]{5.6cm}\raggedright\Large $52^\circ$\end{minipage}\end{tcolorbox}\hfill\begin{tcolorbox}[enhanced,nobeforeafter,width=7.1cm,height=2.1cm,valign=center,colback=cgeo!6,colframe=cgeo,arc=2.5mm,boxrule=0.9pt,left=3mm]{\Large\bfseries\color{cgeo}b)}\quad\begin{minipage}[c]{5.6cm}\raggedright\Large $322^\circ$\end{minipage}\end{tcolorbox}\par\vspace{2.5mm}
\begin{tcolorbox}[enhanced,nobeforeafter,width=7.1cm,height=2.1cm,valign=center,colback=cgeo!6,colframe=cgeo,arc=2.5mm,boxrule=0.9pt,left=3mm]{\Large\bfseries\color{cgeo}c)}\quad\begin{minipage}[c]{5.6cm}\raggedright\Large $38^\circ$\end{minipage}\end{tcolorbox}\hfill\begin{tcolorbox}[enhanced,nobeforeafter,width=7.1cm,height=2.1cm,valign=center,colback=cgeo!6,colframe=cgeo,arc=2.5mm,boxrule=0.9pt,left=3mm]{\Large\bfseries\color{cgeo}d)}\quad\begin{minipage}[c]{5.6cm}\raggedright\Large $142^\circ$\end{minipage}\end{tcolorbox}
\end{tcolorbox}
\newpage
\begin{tcolorbox}[enhanced,width=16cm,height=14cm,colback=white,colframe=cgeo,boxrule=2pt,arc=0mm,
left=6mm,right=6mm,top=0mm,bottom=4mm,sharp corners,
overlay={\fill[cgeo] (frame.north west) rectangle ([yshift=-1.25cm]frame.north east);
\node[anchor=west,text=white,font=\sffamily\bfseries\LARGE] at ([xshift=5mm,yshift=-0.62cm]frame.north west) {Question 54 B};
\node[anchor=east,text=white,font=\sffamily\itshape\large] at ([xshift=-5mm,yshift=-0.62cm]frame.north east) {Géométrie};}]
\vspace*{1.75cm}
\begin{minipage}[t][6.9cm][c]{\linewidth}\LARGE\centering
\begin{minipage}{0.95\linewidth}\raggedright\linespread{1.25}\selectfont Les points $A$, $O$ et $B$ sont alignés dans cet ordre et $\widehat{AOC}=145^\circ$. Combien mesure $\widehat{COB}$~?\end{minipage}
\end{minipage}

\vfill
\begin{tcolorbox}[enhanced,nobeforeafter,width=7.1cm,height=2.1cm,valign=center,colback=cgeo!6,colframe=cgeo,arc=2.5mm,boxrule=0.9pt,left=3mm]{\Large\bfseries\color{cgeo}a)}\quad\begin{minipage}[c]{5.6cm}\raggedright\Large $215^\circ$\end{minipage}\end{tcolorbox}\hfill\begin{tcolorbox}[enhanced,nobeforeafter,width=7.1cm,height=2.1cm,valign=center,colback=cgeo!6,colframe=cgeo,arc=2.5mm,boxrule=0.9pt,left=3mm]{\Large\bfseries\color{cgeo}b)}\quad\begin{minipage}[c]{5.6cm}\raggedright\Large $55^\circ$\end{minipage}\end{tcolorbox}\par\vspace{2.5mm}
\begin{tcolorbox}[enhanced,nobeforeafter,width=7.1cm,height=2.1cm,valign=center,colback=cgeo!6,colframe=cgeo,arc=2.5mm,boxrule=0.9pt,left=3mm]{\Large\bfseries\color{cgeo}c)}\quad\begin{minipage}[c]{5.6cm}\raggedright\Large $35^\circ$\end{minipage}\end{tcolorbox}\hfill\begin{tcolorbox}[enhanced,nobeforeafter,width=7.1cm,height=2.1cm,valign=center,colback=cgeo!6,colframe=cgeo,arc=2.5mm,boxrule=0.9pt,left=3mm]{\Large\bfseries\color{cgeo}d)}\quad\begin{minipage}[c]{5.6cm}\raggedright\Large $145^\circ$\end{minipage}\end{tcolorbox}
\end{tcolorbox}
\newpage
\begin{tcolorbox}[enhanced,width=16cm,height=14cm,colback=white,colframe=cgeo,boxrule=2pt,arc=0mm,
left=6mm,right=6mm,top=0mm,bottom=4mm,sharp corners,
overlay={\fill[cgeo] (frame.north west) rectangle ([yshift=-1.25cm]frame.north east);
\node[anchor=west,text=white,font=\sffamily\bfseries\LARGE] at ([xshift=5mm,yshift=-0.62cm]frame.north west) {Question 55 B};
\node[anchor=east,text=white,font=\sffamily\itshape\large] at ([xshift=-5mm,yshift=-0.62cm]frame.north east) {Géométrie};}]
\vspace*{1.75cm}
\begin{minipage}[t][6.9cm][c]{\linewidth}\LARGE\centering
\begin{minipage}{0.95\linewidth}\raggedright\linespread{1.25}\selectfont Combien d'axes de symétrie possède un triangle équilatéral~?\end{minipage}
\end{minipage}

\vfill
\begin{tcolorbox}[enhanced,nobeforeafter,width=7.1cm,height=2.1cm,valign=center,colback=cgeo!6,colframe=cgeo,arc=2.5mm,boxrule=0.9pt,left=3mm]{\Large\bfseries\color{cgeo}a)}\quad\begin{minipage}[c]{5.6cm}\raggedright\Large 2\end{minipage}\end{tcolorbox}\hfill\begin{tcolorbox}[enhanced,nobeforeafter,width=7.1cm,height=2.1cm,valign=center,colback=cgeo!6,colframe=cgeo,arc=2.5mm,boxrule=0.9pt,left=3mm]{\Large\bfseries\color{cgeo}b)}\quad\begin{minipage}[c]{5.6cm}\raggedright\Large 3\end{minipage}\end{tcolorbox}\par\vspace{2.5mm}
\begin{tcolorbox}[enhanced,nobeforeafter,width=7.1cm,height=2.1cm,valign=center,colback=cgeo!6,colframe=cgeo,arc=2.5mm,boxrule=0.9pt,left=3mm]{\Large\bfseries\color{cgeo}c)}\quad\begin{minipage}[c]{5.6cm}\raggedright\Large 1\end{minipage}\end{tcolorbox}\hfill\begin{tcolorbox}[enhanced,nobeforeafter,width=7.1cm,height=2.1cm,valign=center,colback=cgeo!6,colframe=cgeo,arc=2.5mm,boxrule=0.9pt,left=3mm]{\Large\bfseries\color{cgeo}d)}\quad\begin{minipage}[c]{5.6cm}\raggedright\Large 4\end{minipage}\end{tcolorbox}
\end{tcolorbox}
\newpage
\begin{tcolorbox}[enhanced,width=16cm,height=14cm,colback=white,colframe=cgeo,boxrule=2pt,arc=0mm,
left=6mm,right=6mm,top=0mm,bottom=4mm,sharp corners,
overlay={\fill[cgeo] (frame.north west) rectangle ([yshift=-1.25cm]frame.north east);
\node[anchor=west,text=white,font=\sffamily\bfseries\LARGE] at ([xshift=5mm,yshift=-0.62cm]frame.north west) {Question 56 B};
\node[anchor=east,text=white,font=\sffamily\itshape\large] at ([xshift=-5mm,yshift=-0.62cm]frame.north east) {Géométrie};}]
\vspace*{1.75cm}
\begin{minipage}[t][6.9cm][c]{\linewidth}\LARGE\centering
\begin{minipage}{0.95\linewidth}\raggedright\linespread{1.25}\selectfont Combien d'axes de symétrie possède un losange qui n'est pas un carré~?\end{minipage}
\end{minipage}

\vfill
\begin{tcolorbox}[enhanced,nobeforeafter,width=7.1cm,height=2.1cm,valign=center,colback=cgeo!6,colframe=cgeo,arc=2.5mm,boxrule=0.9pt,left=3mm]{\Large\bfseries\color{cgeo}a)}\quad\begin{minipage}[c]{5.6cm}\raggedright\Large 3\end{minipage}\end{tcolorbox}\hfill\begin{tcolorbox}[enhanced,nobeforeafter,width=7.1cm,height=2.1cm,valign=center,colback=cgeo!6,colframe=cgeo,arc=2.5mm,boxrule=0.9pt,left=3mm]{\Large\bfseries\color{cgeo}b)}\quad\begin{minipage}[c]{5.6cm}\raggedright\Large 1\end{minipage}\end{tcolorbox}\par\vspace{2.5mm}
\begin{tcolorbox}[enhanced,nobeforeafter,width=7.1cm,height=2.1cm,valign=center,colback=cgeo!6,colframe=cgeo,arc=2.5mm,boxrule=0.9pt,left=3mm]{\Large\bfseries\color{cgeo}c)}\quad\begin{minipage}[c]{5.6cm}\raggedright\Large 2\end{minipage}\end{tcolorbox}\hfill\begin{tcolorbox}[enhanced,nobeforeafter,width=7.1cm,height=2.1cm,valign=center,colback=cgeo!6,colframe=cgeo,arc=2.5mm,boxrule=0.9pt,left=3mm]{\Large\bfseries\color{cgeo}d)}\quad\begin{minipage}[c]{5.6cm}\raggedright\Large 4\end{minipage}\end{tcolorbox}
\end{tcolorbox}
\newpage
\begin{tcolorbox}[enhanced,width=16cm,height=14cm,colback=white,colframe=cgeo,boxrule=2pt,arc=0mm,
left=6mm,right=6mm,top=0mm,bottom=4mm,sharp corners,
overlay={\fill[cgeo] (frame.north west) rectangle ([yshift=-1.25cm]frame.north east);
\node[anchor=west,text=white,font=\sffamily\bfseries\LARGE] at ([xshift=5mm,yshift=-0.62cm]frame.north west) {Question 57 B};
\node[anchor=east,text=white,font=\sffamily\itshape\large] at ([xshift=-5mm,yshift=-0.62cm]frame.north east) {Géométrie};}]
\vspace*{1.75cm}
\begin{minipage}[t][6.9cm][c]{\linewidth}\LARGE\centering
\begin{minipage}{0.95\linewidth}\raggedright\linespread{1.25}\selectfont Un cercle a un diamètre de $15$ cm. Quelle est la longueur de son rayon~?\end{minipage}
\end{minipage}

\vfill
\begin{tcolorbox}[enhanced,nobeforeafter,width=7.1cm,height=2.1cm,valign=center,colback=cgeo!6,colframe=cgeo,arc=2.5mm,boxrule=0.9pt,left=3mm]{\Large\bfseries\color{cgeo}a)}\quad\begin{minipage}[c]{5.6cm}\raggedright\Large $17~\text{cm}$\end{minipage}\end{tcolorbox}\hfill\begin{tcolorbox}[enhanced,nobeforeafter,width=7.1cm,height=2.1cm,valign=center,colback=cgeo!6,colframe=cgeo,arc=2.5mm,boxrule=0.9pt,left=3mm]{\Large\bfseries\color{cgeo}b)}\quad\begin{minipage}[c]{5.6cm}\raggedright\Large $15~\text{cm}$\end{minipage}\end{tcolorbox}\par\vspace{2.5mm}
\begin{tcolorbox}[enhanced,nobeforeafter,width=7.1cm,height=2.1cm,valign=center,colback=cgeo!6,colframe=cgeo,arc=2.5mm,boxrule=0.9pt,left=3mm]{\Large\bfseries\color{cgeo}c)}\quad\begin{minipage}[c]{5.6cm}\raggedright\Large $30~\text{cm}$\end{minipage}\end{tcolorbox}\hfill\begin{tcolorbox}[enhanced,nobeforeafter,width=7.1cm,height=2.1cm,valign=center,colback=cgeo!6,colframe=cgeo,arc=2.5mm,boxrule=0.9pt,left=3mm]{\Large\bfseries\color{cgeo}d)}\quad\begin{minipage}[c]{5.6cm}\raggedright\Large $7{,}5~\text{cm}$\end{minipage}\end{tcolorbox}
\end{tcolorbox}
\newpage
\begin{tcolorbox}[enhanced,width=16cm,height=14cm,colback=white,colframe=cgeo,boxrule=2pt,arc=0mm,
left=6mm,right=6mm,top=0mm,bottom=4mm,sharp corners,
overlay={\fill[cgeo] (frame.north west) rectangle ([yshift=-1.25cm]frame.north east);
\node[anchor=west,text=white,font=\sffamily\bfseries\LARGE] at ([xshift=5mm,yshift=-0.62cm]frame.north west) {Question 58 B};
\node[anchor=east,text=white,font=\sffamily\itshape\large] at ([xshift=-5mm,yshift=-0.62cm]frame.north east) {Géométrie};}]
\vspace*{1.75cm}
\begin{minipage}[t][6.9cm][c]{\linewidth}\LARGE\centering
\begin{minipage}{0.95\linewidth}\raggedright\linespread{1.25}\selectfont Un cercle a un rayon de $3{,}5$ cm. Quelle est la longueur de son diamètre~?\end{minipage}
\end{minipage}

\vfill
\begin{tcolorbox}[enhanced,nobeforeafter,width=7.1cm,height=2.1cm,valign=center,colback=cgeo!6,colframe=cgeo,arc=2.5mm,boxrule=0.9pt,left=3mm]{\Large\bfseries\color{cgeo}a)}\quad\begin{minipage}[c]{5.6cm}\raggedright\Large $5{,}5~\text{cm}$\end{minipage}\end{tcolorbox}\hfill\begin{tcolorbox}[enhanced,nobeforeafter,width=7.1cm,height=2.1cm,valign=center,colback=cgeo!6,colframe=cgeo,arc=2.5mm,boxrule=0.9pt,left=3mm]{\Large\bfseries\color{cgeo}b)}\quad\begin{minipage}[c]{5.6cm}\raggedright\Large $7~\text{cm}$\end{minipage}\end{tcolorbox}\par\vspace{2.5mm}
\begin{tcolorbox}[enhanced,nobeforeafter,width=7.1cm,height=2.1cm,valign=center,colback=cgeo!6,colframe=cgeo,arc=2.5mm,boxrule=0.9pt,left=3mm]{\Large\bfseries\color{cgeo}c)}\quad\begin{minipage}[c]{5.6cm}\raggedright\Large $1{,}75~\text{cm}$\end{minipage}\end{tcolorbox}\hfill\begin{tcolorbox}[enhanced,nobeforeafter,width=7.1cm,height=2.1cm,valign=center,colback=cgeo!6,colframe=cgeo,arc=2.5mm,boxrule=0.9pt,left=3mm]{\Large\bfseries\color{cgeo}d)}\quad\begin{minipage}[c]{5.6cm}\raggedright\Large $3{,}5~\text{cm}$\end{minipage}\end{tcolorbox}
\end{tcolorbox}
\newpage
\begin{tcolorbox}[enhanced,width=16cm,height=14cm,colback=white,colframe=cgeo,boxrule=2pt,arc=0mm,
left=6mm,right=6mm,top=0mm,bottom=4mm,sharp corners,
overlay={\fill[cgeo] (frame.north west) rectangle ([yshift=-1.25cm]frame.north east);
\node[anchor=west,text=white,font=\sffamily\bfseries\LARGE] at ([xshift=5mm,yshift=-0.62cm]frame.north west) {Question 59 B};
\node[anchor=east,text=white,font=\sffamily\itshape\large] at ([xshift=-5mm,yshift=-0.62cm]frame.north east) {Géométrie};}]
\vspace*{1.75cm}
\begin{minipage}[t][6.9cm][c]{\linewidth}\LARGE\centering
\begin{minipage}{0.95\linewidth}\raggedright\linespread{1.25}\selectfont La droite $(u)$ est parallèle à $(w)$ et la droite $(v)$ est parallèle à $(w)$. Que peut-on dire des droites $(u)$ et $(v)$~?\end{minipage}
\end{minipage}

\vfill
\begin{tcolorbox}[enhanced,nobeforeafter,width=7.1cm,height=2.1cm,valign=center,colback=cgeo!6,colframe=cgeo,arc=2.5mm,boxrule=0.9pt,left=3mm]{\Large\bfseries\color{cgeo}a)}\quad\begin{minipage}[c]{5.6cm}\raggedright\Large on ne peut pas savoir\end{minipage}\end{tcolorbox}\hfill\begin{tcolorbox}[enhanced,nobeforeafter,width=7.1cm,height=2.1cm,valign=center,colback=cgeo!6,colframe=cgeo,arc=2.5mm,boxrule=0.9pt,left=3mm]{\Large\bfseries\color{cgeo}b)}\quad\begin{minipage}[c]{5.6cm}\raggedright\Large parallèles\end{minipage}\end{tcolorbox}\par\vspace{2.5mm}
\begin{tcolorbox}[enhanced,nobeforeafter,width=7.1cm,height=2.1cm,valign=center,colback=cgeo!6,colframe=cgeo,arc=2.5mm,boxrule=0.9pt,left=3mm]{\Large\bfseries\color{cgeo}c)}\quad\begin{minipage}[c]{5.6cm}\raggedright\Large perpendiculaires\end{minipage}\end{tcolorbox}\hfill\begin{tcolorbox}[enhanced,nobeforeafter,width=7.1cm,height=2.1cm,valign=center,colback=cgeo!6,colframe=cgeo,arc=2.5mm,boxrule=0.9pt,left=3mm]{\Large\bfseries\color{cgeo}d)}\quad\begin{minipage}[c]{5.6cm}\raggedright\large sécantes, non perpendiculaires\end{minipage}\end{tcolorbox}
\end{tcolorbox}
\newpage
\begin{tcolorbox}[enhanced,width=16cm,height=14cm,colback=white,colframe=cgeo,boxrule=2pt,arc=0mm,
left=6mm,right=6mm,top=0mm,bottom=4mm,sharp corners,
overlay={\fill[cgeo] (frame.north west) rectangle ([yshift=-1.25cm]frame.north east);
\node[anchor=west,text=white,font=\sffamily\bfseries\LARGE] at ([xshift=5mm,yshift=-0.62cm]frame.north west) {Question 60 B};
\node[anchor=east,text=white,font=\sffamily\itshape\large] at ([xshift=-5mm,yshift=-0.62cm]frame.north east) {Géométrie};}]
\vspace*{1.75cm}
\begin{minipage}[t][6.9cm][c]{\linewidth}\LARGE\centering
\begin{minipage}{0.95\linewidth}\raggedright\linespread{1.25}\selectfont La droite $(e)$ est perpendiculaire à $(g)$ et la droite $(f)$ est parallèle à $(g)$. Que peut-on dire des droites $(e)$ et $(f)$~?\end{minipage}
\end{minipage}

\vfill
\begin{tcolorbox}[enhanced,nobeforeafter,width=7.1cm,height=2.1cm,valign=center,colback=cgeo!6,colframe=cgeo,arc=2.5mm,boxrule=0.9pt,left=3mm]{\Large\bfseries\color{cgeo}a)}\quad\begin{minipage}[c]{5.6cm}\raggedright\large sécantes, non perpendiculaires\end{minipage}\end{tcolorbox}\hfill\begin{tcolorbox}[enhanced,nobeforeafter,width=7.1cm,height=2.1cm,valign=center,colback=cgeo!6,colframe=cgeo,arc=2.5mm,boxrule=0.9pt,left=3mm]{\Large\bfseries\color{cgeo}b)}\quad\begin{minipage}[c]{5.6cm}\raggedright\Large perpendiculaires\end{minipage}\end{tcolorbox}\par\vspace{2.5mm}
\begin{tcolorbox}[enhanced,nobeforeafter,width=7.1cm,height=2.1cm,valign=center,colback=cgeo!6,colframe=cgeo,arc=2.5mm,boxrule=0.9pt,left=3mm]{\Large\bfseries\color{cgeo}c)}\quad\begin{minipage}[c]{5.6cm}\raggedright\Large on ne peut pas savoir\end{minipage}\end{tcolorbox}\hfill\begin{tcolorbox}[enhanced,nobeforeafter,width=7.1cm,height=2.1cm,valign=center,colback=cgeo!6,colframe=cgeo,arc=2.5mm,boxrule=0.9pt,left=3mm]{\Large\bfseries\color{cgeo}d)}\quad\begin{minipage}[c]{5.6cm}\raggedright\Large parallèles\end{minipage}\end{tcolorbox}
\end{tcolorbox}
\newpage
\begin{tcolorbox}[enhanced,width=16cm,height=14cm,colback=white,colframe=cgeo,boxrule=2pt,arc=0mm,
left=6mm,right=6mm,top=0mm,bottom=4mm,sharp corners,
overlay={\fill[cgeo] (frame.north west) rectangle ([yshift=-1.25cm]frame.north east);
\node[anchor=west,text=white,font=\sffamily\bfseries\LARGE] at ([xshift=5mm,yshift=-0.62cm]frame.north west) {Question 61 B};
\node[anchor=east,text=white,font=\sffamily\itshape\large] at ([xshift=-5mm,yshift=-0.62cm]frame.north east) {Géométrie};}]
\vspace*{1.75cm}
\begin{minipage}[t][6.9cm][c]{\linewidth}\LARGE\centering
\begin{minipage}{0.95\linewidth}\raggedright\linespread{1.25}\selectfont Un triangle a un angle de $90^\circ$. Quelle est sa nature~?\end{minipage}
\end{minipage}

\vfill
\begin{tcolorbox}[enhanced,nobeforeafter,width=7.1cm,height=2.1cm,valign=center,colback=cgeo!6,colframe=cgeo,arc=2.5mm,boxrule=0.9pt,left=3mm]{\Large\bfseries\color{cgeo}a)}\quad\begin{minipage}[c]{5.6cm}\raggedright\Large rectangle\end{minipage}\end{tcolorbox}\hfill\begin{tcolorbox}[enhanced,nobeforeafter,width=7.1cm,height=2.1cm,valign=center,colback=cgeo!6,colframe=cgeo,arc=2.5mm,boxrule=0.9pt,left=3mm]{\Large\bfseries\color{cgeo}b)}\quad\begin{minipage}[c]{5.6cm}\raggedright\Large quelconque\end{minipage}\end{tcolorbox}\par\vspace{2.5mm}
\begin{tcolorbox}[enhanced,nobeforeafter,width=7.1cm,height=2.1cm,valign=center,colback=cgeo!6,colframe=cgeo,arc=2.5mm,boxrule=0.9pt,left=3mm]{\Large\bfseries\color{cgeo}c)}\quad\begin{minipage}[c]{5.6cm}\raggedright\Large équilatéral\end{minipage}\end{tcolorbox}\hfill\begin{tcolorbox}[enhanced,nobeforeafter,width=7.1cm,height=2.1cm,valign=center,colback=cgeo!6,colframe=cgeo,arc=2.5mm,boxrule=0.9pt,left=3mm]{\Large\bfseries\color{cgeo}d)}\quad\begin{minipage}[c]{5.6cm}\raggedright\Large isocèle\end{minipage}\end{tcolorbox}
\end{tcolorbox}
\newpage
\begin{tcolorbox}[enhanced,width=16cm,height=14cm,colback=white,colframe=cgeo,boxrule=2pt,arc=0mm,
left=6mm,right=6mm,top=0mm,bottom=4mm,sharp corners,
overlay={\fill[cgeo] (frame.north west) rectangle ([yshift=-1.25cm]frame.north east);
\node[anchor=west,text=white,font=\sffamily\bfseries\LARGE] at ([xshift=5mm,yshift=-0.62cm]frame.north west) {Question 62 B};
\node[anchor=east,text=white,font=\sffamily\itshape\large] at ([xshift=-5mm,yshift=-0.62cm]frame.north east) {Géométrie};}]
\vspace*{1.75cm}
\begin{minipage}[t][6.9cm][c]{\linewidth}\LARGE\centering
\begin{minipage}{0.95\linewidth}\raggedright\linespread{1.25}\selectfont Un triangle a des côtés de $4$ cm, $6$ cm et $7$ cm, et aucun angle droit. Quelle est sa nature~?\end{minipage}
\end{minipage}

\vfill
\begin{tcolorbox}[enhanced,nobeforeafter,width=7.1cm,height=2.1cm,valign=center,colback=cgeo!6,colframe=cgeo,arc=2.5mm,boxrule=0.9pt,left=3mm]{\Large\bfseries\color{cgeo}a)}\quad\begin{minipage}[c]{5.6cm}\raggedright\Large isocèle\end{minipage}\end{tcolorbox}\hfill\begin{tcolorbox}[enhanced,nobeforeafter,width=7.1cm,height=2.1cm,valign=center,colback=cgeo!6,colframe=cgeo,arc=2.5mm,boxrule=0.9pt,left=3mm]{\Large\bfseries\color{cgeo}b)}\quad\begin{minipage}[c]{5.6cm}\raggedright\Large quelconque\end{minipage}\end{tcolorbox}\par\vspace{2.5mm}
\begin{tcolorbox}[enhanced,nobeforeafter,width=7.1cm,height=2.1cm,valign=center,colback=cgeo!6,colframe=cgeo,arc=2.5mm,boxrule=0.9pt,left=3mm]{\Large\bfseries\color{cgeo}c)}\quad\begin{minipage}[c]{5.6cm}\raggedright\Large rectangle\end{minipage}\end{tcolorbox}\hfill\begin{tcolorbox}[enhanced,nobeforeafter,width=7.1cm,height=2.1cm,valign=center,colback=cgeo!6,colframe=cgeo,arc=2.5mm,boxrule=0.9pt,left=3mm]{\Large\bfseries\color{cgeo}d)}\quad\begin{minipage}[c]{5.6cm}\raggedright\Large équilatéral\end{minipage}\end{tcolorbox}
\end{tcolorbox}
\newpage
\begin{tcolorbox}[enhanced,width=16cm,height=14cm,colback=white,colframe=cgeo,boxrule=2pt,arc=0mm,
left=6mm,right=6mm,top=0mm,bottom=4mm,sharp corners,
overlay={\fill[cgeo] (frame.north west) rectangle ([yshift=-1.25cm]frame.north east);
\node[anchor=west,text=white,font=\sffamily\bfseries\LARGE] at ([xshift=5mm,yshift=-0.62cm]frame.north west) {Question 63 B};
\node[anchor=east,text=white,font=\sffamily\itshape\large] at ([xshift=-5mm,yshift=-0.62cm]frame.north east) {Géométrie};}]
\vspace*{1.75cm}
\begin{minipage}[t][6.9cm][c]{\linewidth}\LARGE\centering
\begin{minipage}{0.95\linewidth}\raggedright\linespread{1.25}\selectfont Un quadrilatère a quatre côtés de même longueur. Quelle est sa nature la plus précise~?\end{minipage}
\end{minipage}

\vfill
\begin{tcolorbox}[enhanced,nobeforeafter,width=7.1cm,height=2.1cm,valign=center,colback=cgeo!6,colframe=cgeo,arc=2.5mm,boxrule=0.9pt,left=3mm]{\Large\bfseries\color{cgeo}a)}\quad\begin{minipage}[c]{5.6cm}\raggedright\Large un carré\end{minipage}\end{tcolorbox}\hfill\begin{tcolorbox}[enhanced,nobeforeafter,width=7.1cm,height=2.1cm,valign=center,colback=cgeo!6,colframe=cgeo,arc=2.5mm,boxrule=0.9pt,left=3mm]{\Large\bfseries\color{cgeo}b)}\quad\begin{minipage}[c]{5.6cm}\raggedright\Large un triangle\end{minipage}\end{tcolorbox}\par\vspace{2.5mm}
\begin{tcolorbox}[enhanced,nobeforeafter,width=7.1cm,height=2.1cm,valign=center,colback=cgeo!6,colframe=cgeo,arc=2.5mm,boxrule=0.9pt,left=3mm]{\Large\bfseries\color{cgeo}c)}\quad\begin{minipage}[c]{5.6cm}\raggedright\Large un losange\end{minipage}\end{tcolorbox}\hfill\begin{tcolorbox}[enhanced,nobeforeafter,width=7.1cm,height=2.1cm,valign=center,colback=cgeo!6,colframe=cgeo,arc=2.5mm,boxrule=0.9pt,left=3mm]{\Large\bfseries\color{cgeo}d)}\quad\begin{minipage}[c]{5.6cm}\raggedright\Large un rectangle\end{minipage}\end{tcolorbox}
\end{tcolorbox}
\newpage
\begin{tcolorbox}[enhanced,width=16cm,height=14cm,colback=white,colframe=cgeo,boxrule=2pt,arc=0mm,
left=6mm,right=6mm,top=0mm,bottom=4mm,sharp corners,
overlay={\fill[cgeo] (frame.north west) rectangle ([yshift=-1.25cm]frame.north east);
\node[anchor=west,text=white,font=\sffamily\bfseries\LARGE] at ([xshift=5mm,yshift=-0.62cm]frame.north west) {Question 64 B};
\node[anchor=east,text=white,font=\sffamily\itshape\large] at ([xshift=-5mm,yshift=-0.62cm]frame.north east) {Géométrie};}]
\vspace*{1.75cm}
\begin{minipage}[t][6.9cm][c]{\linewidth}\LARGE\centering
\begin{minipage}{0.95\linewidth}\raggedright\linespread{1.25}\selectfont Un quadrilatère a quatre côtés de même longueur et un angle droit. Quelle est sa nature la plus précise~?\end{minipage}
\end{minipage}

\vfill
\begin{tcolorbox}[enhanced,nobeforeafter,width=7.1cm,height=2.1cm,valign=center,colback=cgeo!6,colframe=cgeo,arc=2.5mm,boxrule=0.9pt,left=3mm]{\Large\bfseries\color{cgeo}a)}\quad\begin{minipage}[c]{5.6cm}\raggedright\Large un carré\end{minipage}\end{tcolorbox}\hfill\begin{tcolorbox}[enhanced,nobeforeafter,width=7.1cm,height=2.1cm,valign=center,colback=cgeo!6,colframe=cgeo,arc=2.5mm,boxrule=0.9pt,left=3mm]{\Large\bfseries\color{cgeo}b)}\quad\begin{minipage}[c]{5.6cm}\raggedright\Large un rectangle\end{minipage}\end{tcolorbox}\par\vspace{2.5mm}
\begin{tcolorbox}[enhanced,nobeforeafter,width=7.1cm,height=2.1cm,valign=center,colback=cgeo!6,colframe=cgeo,arc=2.5mm,boxrule=0.9pt,left=3mm]{\Large\bfseries\color{cgeo}c)}\quad\begin{minipage}[c]{5.6cm}\raggedright\Large un losange\end{minipage}\end{tcolorbox}\hfill\begin{tcolorbox}[enhanced,nobeforeafter,width=7.1cm,height=2.1cm,valign=center,colback=cgeo!6,colframe=cgeo,arc=2.5mm,boxrule=0.9pt,left=3mm]{\Large\bfseries\color{cgeo}d)}\quad\begin{minipage}[c]{5.6cm}\raggedright\Large un triangle\end{minipage}\end{tcolorbox}
\end{tcolorbox}
\newpage
\begin{tcolorbox}[enhanced,width=16cm,height=14cm,colback=white,colframe=cmes,boxrule=2pt,arc=0mm,
left=6mm,right=6mm,top=0mm,bottom=4mm,sharp corners,
overlay={\fill[cmes] (frame.north west) rectangle ([yshift=-1.25cm]frame.north east);
\node[anchor=west,text=white,font=\sffamily\bfseries\LARGE] at ([xshift=5mm,yshift=-0.62cm]frame.north west) {Question 65 B};
\node[anchor=east,text=white,font=\sffamily\itshape\large] at ([xshift=-5mm,yshift=-0.62cm]frame.north east) {Grandeurs et mesures};}]
\vspace*{1.75cm}
\begin{minipage}[t][6.9cm][c]{\linewidth}\LARGE\centering
\begin{minipage}{0.95\linewidth}\raggedright\linespread{1.25}\selectfont Un rectangle a une longueur de $6{,}5$ cm et une largeur de $4$ cm. Quel est son périmètre~?\end{minipage}
\end{minipage}

\vfill
\begin{tcolorbox}[enhanced,nobeforeafter,width=7.1cm,height=2.1cm,valign=center,colback=cmes!6,colframe=cmes,arc=2.5mm,boxrule=0.9pt,left=3mm]{\Large\bfseries\color{cmes}a)}\quad\begin{minipage}[c]{5.6cm}\raggedright\Large $26~\text{cm}$\end{minipage}\end{tcolorbox}\hfill\begin{tcolorbox}[enhanced,nobeforeafter,width=7.1cm,height=2.1cm,valign=center,colback=cmes!6,colframe=cmes,arc=2.5mm,boxrule=0.9pt,left=3mm]{\Large\bfseries\color{cmes}b)}\quad\begin{minipage}[c]{5.6cm}\raggedright\Large $10{,}5~\text{cm}$\end{minipage}\end{tcolorbox}\par\vspace{2.5mm}
\begin{tcolorbox}[enhanced,nobeforeafter,width=7.1cm,height=2.1cm,valign=center,colback=cmes!6,colframe=cmes,arc=2.5mm,boxrule=0.9pt,left=3mm]{\Large\bfseries\color{cmes}c)}\quad\begin{minipage}[c]{5.6cm}\raggedright\Large $21~\text{cm}$\end{minipage}\end{tcolorbox}\hfill\begin{tcolorbox}[enhanced,nobeforeafter,width=7.1cm,height=2.1cm,valign=center,colback=cmes!6,colframe=cmes,arc=2.5mm,boxrule=0.9pt,left=3mm]{\Large\bfseries\color{cmes}d)}\quad\begin{minipage}[c]{5.6cm}\raggedright\Large $17~\text{cm}$\end{minipage}\end{tcolorbox}
\end{tcolorbox}
\newpage
\begin{tcolorbox}[enhanced,width=16cm,height=14cm,colback=white,colframe=cmes,boxrule=2pt,arc=0mm,
left=6mm,right=6mm,top=0mm,bottom=4mm,sharp corners,
overlay={\fill[cmes] (frame.north west) rectangle ([yshift=-1.25cm]frame.north east);
\node[anchor=west,text=white,font=\sffamily\bfseries\LARGE] at ([xshift=5mm,yshift=-0.62cm]frame.north west) {Question 66 B};
\node[anchor=east,text=white,font=\sffamily\itshape\large] at ([xshift=-5mm,yshift=-0.62cm]frame.north east) {Grandeurs et mesures};}]
\vspace*{1.75cm}
\begin{minipage}[t][6.9cm][c]{\linewidth}\LARGE\centering
\begin{minipage}{0.95\linewidth}\raggedright\linespread{1.25}\selectfont Un rectangle a une longueur de $15$ cm et une largeur de $9$ cm. Quel est son périmètre~?\end{minipage}
\end{minipage}

\vfill
\begin{tcolorbox}[enhanced,nobeforeafter,width=7.1cm,height=2.1cm,valign=center,colback=cmes!6,colframe=cmes,arc=2.5mm,boxrule=0.9pt,left=3mm]{\Large\bfseries\color{cmes}a)}\quad\begin{minipage}[c]{5.6cm}\raggedright\Large $48~\text{cm}$\end{minipage}\end{tcolorbox}\hfill\begin{tcolorbox}[enhanced,nobeforeafter,width=7.1cm,height=2.1cm,valign=center,colback=cmes!6,colframe=cmes,arc=2.5mm,boxrule=0.9pt,left=3mm]{\Large\bfseries\color{cmes}b)}\quad\begin{minipage}[c]{5.6cm}\raggedright\Large $135~\text{cm}$\end{minipage}\end{tcolorbox}\par\vspace{2.5mm}
\begin{tcolorbox}[enhanced,nobeforeafter,width=7.1cm,height=2.1cm,valign=center,colback=cmes!6,colframe=cmes,arc=2.5mm,boxrule=0.9pt,left=3mm]{\Large\bfseries\color{cmes}c)}\quad\begin{minipage}[c]{5.6cm}\raggedright\Large $39~\text{cm}$\end{minipage}\end{tcolorbox}\hfill\begin{tcolorbox}[enhanced,nobeforeafter,width=7.1cm,height=2.1cm,valign=center,colback=cmes!6,colframe=cmes,arc=2.5mm,boxrule=0.9pt,left=3mm]{\Large\bfseries\color{cmes}d)}\quad\begin{minipage}[c]{5.6cm}\raggedright\Large $24~\text{cm}$\end{minipage}\end{tcolorbox}
\end{tcolorbox}
\newpage
\begin{tcolorbox}[enhanced,width=16cm,height=14cm,colback=white,colframe=cmes,boxrule=2pt,arc=0mm,
left=6mm,right=6mm,top=0mm,bottom=4mm,sharp corners,
overlay={\fill[cmes] (frame.north west) rectangle ([yshift=-1.25cm]frame.north east);
\node[anchor=west,text=white,font=\sffamily\bfseries\LARGE] at ([xshift=5mm,yshift=-0.62cm]frame.north west) {Question 67 B};
\node[anchor=east,text=white,font=\sffamily\itshape\large] at ([xshift=-5mm,yshift=-0.62cm]frame.north east) {Grandeurs et mesures};}]
\vspace*{1.75cm}
\begin{minipage}[t][6.9cm][c]{\linewidth}\LARGE\centering
\begin{minipage}{0.95\linewidth}\raggedright\linespread{1.25}\selectfont Un rectangle a une longueur de $6{,}5$ cm et une largeur de $4$ cm. Quelle est son aire~?\end{minipage}
\end{minipage}

\vfill
\begin{tcolorbox}[enhanced,nobeforeafter,width=7.1cm,height=2.1cm,valign=center,colback=cmes!6,colframe=cmes,arc=2.5mm,boxrule=0.9pt,left=3mm]{\Large\bfseries\color{cmes}a)}\quad\begin{minipage}[c]{5.6cm}\raggedright\Large $21~\text{cm}^2$\end{minipage}\end{tcolorbox}\hfill\begin{tcolorbox}[enhanced,nobeforeafter,width=7.1cm,height=2.1cm,valign=center,colback=cmes!6,colframe=cmes,arc=2.5mm,boxrule=0.9pt,left=3mm]{\Large\bfseries\color{cmes}b)}\quad\begin{minipage}[c]{5.6cm}\raggedright\Large $26~\text{cm}^2$\end{minipage}\end{tcolorbox}\par\vspace{2.5mm}
\begin{tcolorbox}[enhanced,nobeforeafter,width=7.1cm,height=2.1cm,valign=center,colback=cmes!6,colframe=cmes,arc=2.5mm,boxrule=0.9pt,left=3mm]{\Large\bfseries\color{cmes}c)}\quad\begin{minipage}[c]{5.6cm}\raggedright\Large $10{,}5~\text{cm}^2$\end{minipage}\end{tcolorbox}\hfill\begin{tcolorbox}[enhanced,nobeforeafter,width=7.1cm,height=2.1cm,valign=center,colback=cmes!6,colframe=cmes,arc=2.5mm,boxrule=0.9pt,left=3mm]{\Large\bfseries\color{cmes}d)}\quad\begin{minipage}[c]{5.6cm}\raggedright\Large $13~\text{cm}^2$\end{minipage}\end{tcolorbox}
\end{tcolorbox}
\newpage
\begin{tcolorbox}[enhanced,width=16cm,height=14cm,colback=white,colframe=cmes,boxrule=2pt,arc=0mm,
left=6mm,right=6mm,top=0mm,bottom=4mm,sharp corners,
overlay={\fill[cmes] (frame.north west) rectangle ([yshift=-1.25cm]frame.north east);
\node[anchor=west,text=white,font=\sffamily\bfseries\LARGE] at ([xshift=5mm,yshift=-0.62cm]frame.north west) {Question 68 B};
\node[anchor=east,text=white,font=\sffamily\itshape\large] at ([xshift=-5mm,yshift=-0.62cm]frame.north east) {Grandeurs et mesures};}]
\vspace*{1.75cm}
\begin{minipage}[t][6.9cm][c]{\linewidth}\LARGE\centering
\begin{minipage}{0.95\linewidth}\raggedright\linespread{1.25}\selectfont Un rectangle a une longueur de $15$ cm et une largeur de $9$ cm. Quelle est son aire~?\end{minipage}
\end{minipage}

\vfill
\begin{tcolorbox}[enhanced,nobeforeafter,width=7.1cm,height=2.1cm,valign=center,colback=cmes!6,colframe=cmes,arc=2.5mm,boxrule=0.9pt,left=3mm]{\Large\bfseries\color{cmes}a)}\quad\begin{minipage}[c]{5.6cm}\raggedright\Large $135~\text{cm}^2$\end{minipage}\end{tcolorbox}\hfill\begin{tcolorbox}[enhanced,nobeforeafter,width=7.1cm,height=2.1cm,valign=center,colback=cmes!6,colframe=cmes,arc=2.5mm,boxrule=0.9pt,left=3mm]{\Large\bfseries\color{cmes}b)}\quad\begin{minipage}[c]{5.6cm}\raggedright\Large $24~\text{cm}^2$\end{minipage}\end{tcolorbox}\par\vspace{2.5mm}
\begin{tcolorbox}[enhanced,nobeforeafter,width=7.1cm,height=2.1cm,valign=center,colback=cmes!6,colframe=cmes,arc=2.5mm,boxrule=0.9pt,left=3mm]{\Large\bfseries\color{cmes}c)}\quad\begin{minipage}[c]{5.6cm}\raggedright\Large $67{,}5~\text{cm}^2$\end{minipage}\end{tcolorbox}\hfill\begin{tcolorbox}[enhanced,nobeforeafter,width=7.1cm,height=2.1cm,valign=center,colback=cmes!6,colframe=cmes,arc=2.5mm,boxrule=0.9pt,left=3mm]{\Large\bfseries\color{cmes}d)}\quad\begin{minipage}[c]{5.6cm}\raggedright\Large $48~\text{cm}^2$\end{minipage}\end{tcolorbox}
\end{tcolorbox}
\newpage
\begin{tcolorbox}[enhanced,width=16cm,height=14cm,colback=white,colframe=cmes,boxrule=2pt,arc=0mm,
left=6mm,right=6mm,top=0mm,bottom=4mm,sharp corners,
overlay={\fill[cmes] (frame.north west) rectangle ([yshift=-1.25cm]frame.north east);
\node[anchor=west,text=white,font=\sffamily\bfseries\LARGE] at ([xshift=5mm,yshift=-0.62cm]frame.north west) {Question 69 B};
\node[anchor=east,text=white,font=\sffamily\itshape\large] at ([xshift=-5mm,yshift=-0.62cm]frame.north east) {Grandeurs et mesures};}]
\vspace*{1.75cm}
\begin{minipage}[t][6.9cm][c]{\linewidth}\LARGE\centering
\begin{minipage}{0.95\linewidth}\raggedright\linespread{1.25}\selectfont Un triangle est rectangle ; les deux côtés de l'angle droit mesurent $8$ cm et $5$ cm. Quelle est son aire~?\end{minipage}
\end{minipage}

\vfill
\begin{tcolorbox}[enhanced,nobeforeafter,width=7.1cm,height=2.1cm,valign=center,colback=cmes!6,colframe=cmes,arc=2.5mm,boxrule=0.9pt,left=3mm]{\Large\bfseries\color{cmes}a)}\quad\begin{minipage}[c]{5.6cm}\raggedright\Large $26~\text{cm}^2$\end{minipage}\end{tcolorbox}\hfill\begin{tcolorbox}[enhanced,nobeforeafter,width=7.1cm,height=2.1cm,valign=center,colback=cmes!6,colframe=cmes,arc=2.5mm,boxrule=0.9pt,left=3mm]{\Large\bfseries\color{cmes}b)}\quad\begin{minipage}[c]{5.6cm}\raggedright\Large $40~\text{cm}^2$\end{minipage}\end{tcolorbox}\par\vspace{2.5mm}
\begin{tcolorbox}[enhanced,nobeforeafter,width=7.1cm,height=2.1cm,valign=center,colback=cmes!6,colframe=cmes,arc=2.5mm,boxrule=0.9pt,left=3mm]{\Large\bfseries\color{cmes}c)}\quad\begin{minipage}[c]{5.6cm}\raggedright\Large $13~\text{cm}^2$\end{minipage}\end{tcolorbox}\hfill\begin{tcolorbox}[enhanced,nobeforeafter,width=7.1cm,height=2.1cm,valign=center,colback=cmes!6,colframe=cmes,arc=2.5mm,boxrule=0.9pt,left=3mm]{\Large\bfseries\color{cmes}d)}\quad\begin{minipage}[c]{5.6cm}\raggedright\Large $20~\text{cm}^2$\end{minipage}\end{tcolorbox}
\end{tcolorbox}
\newpage
\begin{tcolorbox}[enhanced,width=16cm,height=14cm,colback=white,colframe=cmes,boxrule=2pt,arc=0mm,
left=6mm,right=6mm,top=0mm,bottom=4mm,sharp corners,
overlay={\fill[cmes] (frame.north west) rectangle ([yshift=-1.25cm]frame.north east);
\node[anchor=west,text=white,font=\sffamily\bfseries\LARGE] at ([xshift=5mm,yshift=-0.62cm]frame.north west) {Question 70 B};
\node[anchor=east,text=white,font=\sffamily\itshape\large] at ([xshift=-5mm,yshift=-0.62cm]frame.north east) {Grandeurs et mesures};}]
\vspace*{1.75cm}
\begin{minipage}[t][6.9cm][c]{\linewidth}\LARGE\centering
\begin{minipage}{0.95\linewidth}\raggedright\linespread{1.25}\selectfont Un triangle est rectangle ; les deux côtés de l'angle droit mesurent $12$ cm et $9$ cm. Quelle est son aire~?\end{minipage}
\end{minipage}

\vfill
\begin{tcolorbox}[enhanced,nobeforeafter,width=7.1cm,height=2.1cm,valign=center,colback=cmes!6,colframe=cmes,arc=2.5mm,boxrule=0.9pt,left=3mm]{\Large\bfseries\color{cmes}a)}\quad\begin{minipage}[c]{5.6cm}\raggedright\Large $54~\text{cm}^2$\end{minipage}\end{tcolorbox}\hfill\begin{tcolorbox}[enhanced,nobeforeafter,width=7.1cm,height=2.1cm,valign=center,colback=cmes!6,colframe=cmes,arc=2.5mm,boxrule=0.9pt,left=3mm]{\Large\bfseries\color{cmes}b)}\quad\begin{minipage}[c]{5.6cm}\raggedright\Large $42~\text{cm}^2$\end{minipage}\end{tcolorbox}\par\vspace{2.5mm}
\begin{tcolorbox}[enhanced,nobeforeafter,width=7.1cm,height=2.1cm,valign=center,colback=cmes!6,colframe=cmes,arc=2.5mm,boxrule=0.9pt,left=3mm]{\Large\bfseries\color{cmes}c)}\quad\begin{minipage}[c]{5.6cm}\raggedright\Large $21~\text{cm}^2$\end{minipage}\end{tcolorbox}\hfill\begin{tcolorbox}[enhanced,nobeforeafter,width=7.1cm,height=2.1cm,valign=center,colback=cmes!6,colframe=cmes,arc=2.5mm,boxrule=0.9pt,left=3mm]{\Large\bfseries\color{cmes}d)}\quad\begin{minipage}[c]{5.6cm}\raggedright\Large $108~\text{cm}^2$\end{minipage}\end{tcolorbox}
\end{tcolorbox}
\newpage
\begin{tcolorbox}[enhanced,width=16cm,height=14cm,colback=white,colframe=cmes,boxrule=2pt,arc=0mm,
left=6mm,right=6mm,top=0mm,bottom=4mm,sharp corners,
overlay={\fill[cmes] (frame.north west) rectangle ([yshift=-1.25cm]frame.north east);
\node[anchor=west,text=white,font=\sffamily\bfseries\LARGE] at ([xshift=5mm,yshift=-0.62cm]frame.north west) {Question 71 B};
\node[anchor=east,text=white,font=\sffamily\itshape\large] at ([xshift=-5mm,yshift=-0.62cm]frame.north east) {Grandeurs et mesures};}]
\vspace*{1.75cm}
\begin{minipage}[t][6.9cm][c]{\linewidth}\LARGE\centering
\begin{minipage}{0.95\linewidth}\raggedright\linespread{1.25}\selectfont Compléter~: $42~\text{mm}=\;?\;\text{cm}$\end{minipage}
\end{minipage}

\vfill
\begin{tcolorbox}[enhanced,nobeforeafter,width=7.1cm,height=2.1cm,valign=center,colback=cmes!6,colframe=cmes,arc=2.5mm,boxrule=0.9pt,left=3mm]{\Large\bfseries\color{cmes}a)}\quad\begin{minipage}[c]{5.6cm}\raggedright\Large $42~\text{cm}$\end{minipage}\end{tcolorbox}\hfill\begin{tcolorbox}[enhanced,nobeforeafter,width=7.1cm,height=2.1cm,valign=center,colback=cmes!6,colframe=cmes,arc=2.5mm,boxrule=0.9pt,left=3mm]{\Large\bfseries\color{cmes}b)}\quad\begin{minipage}[c]{5.6cm}\raggedright\Large $0{,}42~\text{cm}$\end{minipage}\end{tcolorbox}\par\vspace{2.5mm}
\begin{tcolorbox}[enhanced,nobeforeafter,width=7.1cm,height=2.1cm,valign=center,colback=cmes!6,colframe=cmes,arc=2.5mm,boxrule=0.9pt,left=3mm]{\Large\bfseries\color{cmes}c)}\quad\begin{minipage}[c]{5.6cm}\raggedright\Large $420~\text{cm}$\end{minipage}\end{tcolorbox}\hfill\begin{tcolorbox}[enhanced,nobeforeafter,width=7.1cm,height=2.1cm,valign=center,colback=cmes!6,colframe=cmes,arc=2.5mm,boxrule=0.9pt,left=3mm]{\Large\bfseries\color{cmes}d)}\quad\begin{minipage}[c]{5.6cm}\raggedright\Large $4{,}2~\text{cm}$\end{minipage}\end{tcolorbox}
\end{tcolorbox}
\newpage
\begin{tcolorbox}[enhanced,width=16cm,height=14cm,colback=white,colframe=cmes,boxrule=2pt,arc=0mm,
left=6mm,right=6mm,top=0mm,bottom=4mm,sharp corners,
overlay={\fill[cmes] (frame.north west) rectangle ([yshift=-1.25cm]frame.north east);
\node[anchor=west,text=white,font=\sffamily\bfseries\LARGE] at ([xshift=5mm,yshift=-0.62cm]frame.north west) {Question 72 B};
\node[anchor=east,text=white,font=\sffamily\itshape\large] at ([xshift=-5mm,yshift=-0.62cm]frame.north east) {Grandeurs et mesures};}]
\vspace*{1.75cm}
\begin{minipage}[t][6.9cm][c]{\linewidth}\LARGE\centering
\begin{minipage}{0.95\linewidth}\raggedright\linespread{1.25}\selectfont Compléter~: $0{,}8~\text{km}=\;?\;\text{m}$\end{minipage}
\end{minipage}

\vfill
\begin{tcolorbox}[enhanced,nobeforeafter,width=7.1cm,height=2.1cm,valign=center,colback=cmes!6,colframe=cmes,arc=2.5mm,boxrule=0.9pt,left=3mm]{\Large\bfseries\color{cmes}a)}\quad\begin{minipage}[c]{5.6cm}\raggedright\Large $8\,000~\text{m}$\end{minipage}\end{tcolorbox}\hfill\begin{tcolorbox}[enhanced,nobeforeafter,width=7.1cm,height=2.1cm,valign=center,colback=cmes!6,colframe=cmes,arc=2.5mm,boxrule=0.9pt,left=3mm]{\Large\bfseries\color{cmes}b)}\quad\begin{minipage}[c]{5.6cm}\raggedright\Large $80~\text{m}$\end{minipage}\end{tcolorbox}\par\vspace{2.5mm}
\begin{tcolorbox}[enhanced,nobeforeafter,width=7.1cm,height=2.1cm,valign=center,colback=cmes!6,colframe=cmes,arc=2.5mm,boxrule=0.9pt,left=3mm]{\Large\bfseries\color{cmes}c)}\quad\begin{minipage}[c]{5.6cm}\raggedright\Large $80\,000~\text{m}$\end{minipage}\end{tcolorbox}\hfill\begin{tcolorbox}[enhanced,nobeforeafter,width=7.1cm,height=2.1cm,valign=center,colback=cmes!6,colframe=cmes,arc=2.5mm,boxrule=0.9pt,left=3mm]{\Large\bfseries\color{cmes}d)}\quad\begin{minipage}[c]{5.6cm}\raggedright\Large $800~\text{m}$\end{minipage}\end{tcolorbox}
\end{tcolorbox}
\newpage
\begin{tcolorbox}[enhanced,width=16cm,height=14cm,colback=white,colframe=cmes,boxrule=2pt,arc=0mm,
left=6mm,right=6mm,top=0mm,bottom=4mm,sharp corners,
overlay={\fill[cmes] (frame.north west) rectangle ([yshift=-1.25cm]frame.north east);
\node[anchor=west,text=white,font=\sffamily\bfseries\LARGE] at ([xshift=5mm,yshift=-0.62cm]frame.north west) {Question 73 B};
\node[anchor=east,text=white,font=\sffamily\itshape\large] at ([xshift=-5mm,yshift=-0.62cm]frame.north east) {Grandeurs et mesures};}]
\vspace*{1.75cm}
\begin{minipage}[t][6.9cm][c]{\linewidth}\LARGE\centering
\begin{minipage}{0.95\linewidth}\raggedright\linespread{1.25}\selectfont Compléter~: $45~\text{g}=\;?\;\text{mg}$\end{minipage}
\end{minipage}

\vfill
\begin{tcolorbox}[enhanced,nobeforeafter,width=7.1cm,height=2.1cm,valign=center,colback=cmes!6,colframe=cmes,arc=2.5mm,boxrule=0.9pt,left=3mm]{\Large\bfseries\color{cmes}a)}\quad\begin{minipage}[c]{5.6cm}\raggedright\Large $450\,000~\text{mg}$\end{minipage}\end{tcolorbox}\hfill\begin{tcolorbox}[enhanced,nobeforeafter,width=7.1cm,height=2.1cm,valign=center,colback=cmes!6,colframe=cmes,arc=2.5mm,boxrule=0.9pt,left=3mm]{\Large\bfseries\color{cmes}b)}\quad\begin{minipage}[c]{5.6cm}\raggedright\Large $45\,000~\text{mg}$\end{minipage}\end{tcolorbox}\par\vspace{2.5mm}
\begin{tcolorbox}[enhanced,nobeforeafter,width=7.1cm,height=2.1cm,valign=center,colback=cmes!6,colframe=cmes,arc=2.5mm,boxrule=0.9pt,left=3mm]{\Large\bfseries\color{cmes}c)}\quad\begin{minipage}[c]{5.6cm}\raggedright\Large $4\,500~\text{mg}$\end{minipage}\end{tcolorbox}\hfill\begin{tcolorbox}[enhanced,nobeforeafter,width=7.1cm,height=2.1cm,valign=center,colback=cmes!6,colframe=cmes,arc=2.5mm,boxrule=0.9pt,left=3mm]{\Large\bfseries\color{cmes}d)}\quad\begin{minipage}[c]{5.6cm}\raggedright\Large $4\,500\,000~\text{mg}$\end{minipage}\end{tcolorbox}
\end{tcolorbox}
\newpage
\begin{tcolorbox}[enhanced,width=16cm,height=14cm,colback=white,colframe=cmes,boxrule=2pt,arc=0mm,
left=6mm,right=6mm,top=0mm,bottom=4mm,sharp corners,
overlay={\fill[cmes] (frame.north west) rectangle ([yshift=-1.25cm]frame.north east);
\node[anchor=west,text=white,font=\sffamily\bfseries\LARGE] at ([xshift=5mm,yshift=-0.62cm]frame.north west) {Question 74 B};
\node[anchor=east,text=white,font=\sffamily\itshape\large] at ([xshift=-5mm,yshift=-0.62cm]frame.north east) {Grandeurs et mesures};}]
\vspace*{1.75cm}
\begin{minipage}[t][6.9cm][c]{\linewidth}\LARGE\centering
\begin{minipage}{0.95\linewidth}\raggedright\linespread{1.25}\selectfont Compléter~: $3\,500~\text{g}=\;?\;\text{kg}$\end{minipage}
\end{minipage}

\vfill
\begin{tcolorbox}[enhanced,nobeforeafter,width=7.1cm,height=2.1cm,valign=center,colback=cmes!6,colframe=cmes,arc=2.5mm,boxrule=0.9pt,left=3mm]{\Large\bfseries\color{cmes}a)}\quad\begin{minipage}[c]{5.6cm}\raggedright\Large $3{,}5~\text{kg}$\end{minipage}\end{tcolorbox}\hfill\begin{tcolorbox}[enhanced,nobeforeafter,width=7.1cm,height=2.1cm,valign=center,colback=cmes!6,colframe=cmes,arc=2.5mm,boxrule=0.9pt,left=3mm]{\Large\bfseries\color{cmes}b)}\quad\begin{minipage}[c]{5.6cm}\raggedright\Large $350~\text{kg}$\end{minipage}\end{tcolorbox}\par\vspace{2.5mm}
\begin{tcolorbox}[enhanced,nobeforeafter,width=7.1cm,height=2.1cm,valign=center,colback=cmes!6,colframe=cmes,arc=2.5mm,boxrule=0.9pt,left=3mm]{\Large\bfseries\color{cmes}c)}\quad\begin{minipage}[c]{5.6cm}\raggedright\Large $0{,}35~\text{kg}$\end{minipage}\end{tcolorbox}\hfill\begin{tcolorbox}[enhanced,nobeforeafter,width=7.1cm,height=2.1cm,valign=center,colback=cmes!6,colframe=cmes,arc=2.5mm,boxrule=0.9pt,left=3mm]{\Large\bfseries\color{cmes}d)}\quad\begin{minipage}[c]{5.6cm}\raggedright\Large $35~\text{kg}$\end{minipage}\end{tcolorbox}
\end{tcolorbox}
\newpage
\begin{tcolorbox}[enhanced,width=16cm,height=14cm,colback=white,colframe=cmes,boxrule=2pt,arc=0mm,
left=6mm,right=6mm,top=0mm,bottom=4mm,sharp corners,
overlay={\fill[cmes] (frame.north west) rectangle ([yshift=-1.25cm]frame.north east);
\node[anchor=west,text=white,font=\sffamily\bfseries\LARGE] at ([xshift=5mm,yshift=-0.62cm]frame.north west) {Question 75 B};
\node[anchor=east,text=white,font=\sffamily\itshape\large] at ([xshift=-5mm,yshift=-0.62cm]frame.north east) {Grandeurs et mesures};}]
\vspace*{1.75cm}
\begin{minipage}[t][6.9cm][c]{\linewidth}\LARGE\centering
\begin{minipage}{0.95\linewidth}\raggedright\linespread{1.25}\selectfont Combien de minutes y a-t-il dans 3 h 15 min~?\end{minipage}
\end{minipage}

\vfill
\begin{tcolorbox}[enhanced,nobeforeafter,width=7.1cm,height=2.1cm,valign=center,colback=cmes!6,colframe=cmes,arc=2.5mm,boxrule=0.9pt,left=3mm]{\Large\bfseries\color{cmes}a)}\quad\begin{minipage}[c]{5.6cm}\raggedright\Large $195~\text{min}$\end{minipage}\end{tcolorbox}\hfill\begin{tcolorbox}[enhanced,nobeforeafter,width=7.1cm,height=2.1cm,valign=center,colback=cmes!6,colframe=cmes,arc=2.5mm,boxrule=0.9pt,left=3mm]{\Large\bfseries\color{cmes}b)}\quad\begin{minipage}[c]{5.6cm}\raggedright\Large $315~\text{min}$\end{minipage}\end{tcolorbox}\par\vspace{2.5mm}
\begin{tcolorbox}[enhanced,nobeforeafter,width=7.1cm,height=2.1cm,valign=center,colback=cmes!6,colframe=cmes,arc=2.5mm,boxrule=0.9pt,left=3mm]{\Large\bfseries\color{cmes}c)}\quad\begin{minipage}[c]{5.6cm}\raggedright\Large $18~\text{min}$\end{minipage}\end{tcolorbox}\hfill\begin{tcolorbox}[enhanced,nobeforeafter,width=7.1cm,height=2.1cm,valign=center,colback=cmes!6,colframe=cmes,arc=2.5mm,boxrule=0.9pt,left=3mm]{\Large\bfseries\color{cmes}d)}\quad\begin{minipage}[c]{5.6cm}\raggedright\Large $165~\text{min}$\end{minipage}\end{tcolorbox}
\end{tcolorbox}
\newpage
\begin{tcolorbox}[enhanced,width=16cm,height=14cm,colback=white,colframe=cmes,boxrule=2pt,arc=0mm,
left=6mm,right=6mm,top=0mm,bottom=4mm,sharp corners,
overlay={\fill[cmes] (frame.north west) rectangle ([yshift=-1.25cm]frame.north east);
\node[anchor=west,text=white,font=\sffamily\bfseries\LARGE] at ([xshift=5mm,yshift=-0.62cm]frame.north west) {Question 76 B};
\node[anchor=east,text=white,font=\sffamily\itshape\large] at ([xshift=-5mm,yshift=-0.62cm]frame.north east) {Grandeurs et mesures};}]
\vspace*{1.75cm}
\begin{minipage}[t][6.9cm][c]{\linewidth}\LARGE\centering
\begin{minipage}{0.95\linewidth}\raggedright\linespread{1.25}\selectfont Combien de minutes y a-t-il dans 1 h 5 min~?\end{minipage}
\end{minipage}

\vfill
\begin{tcolorbox}[enhanced,nobeforeafter,width=7.1cm,height=2.1cm,valign=center,colback=cmes!6,colframe=cmes,arc=2.5mm,boxrule=0.9pt,left=3mm]{\Large\bfseries\color{cmes}a)}\quad\begin{minipage}[c]{5.6cm}\raggedright\Large $105~\text{min}$\end{minipage}\end{tcolorbox}\hfill\begin{tcolorbox}[enhanced,nobeforeafter,width=7.1cm,height=2.1cm,valign=center,colback=cmes!6,colframe=cmes,arc=2.5mm,boxrule=0.9pt,left=3mm]{\Large\bfseries\color{cmes}b)}\quad\begin{minipage}[c]{5.6cm}\raggedright\Large $65~\text{min}$\end{minipage}\end{tcolorbox}\par\vspace{2.5mm}
\begin{tcolorbox}[enhanced,nobeforeafter,width=7.1cm,height=2.1cm,valign=center,colback=cmes!6,colframe=cmes,arc=2.5mm,boxrule=0.9pt,left=3mm]{\Large\bfseries\color{cmes}c)}\quad\begin{minipage}[c]{5.6cm}\raggedright\Large $55~\text{min}$\end{minipage}\end{tcolorbox}\hfill\begin{tcolorbox}[enhanced,nobeforeafter,width=7.1cm,height=2.1cm,valign=center,colback=cmes!6,colframe=cmes,arc=2.5mm,boxrule=0.9pt,left=3mm]{\Large\bfseries\color{cmes}d)}\quad\begin{minipage}[c]{5.6cm}\raggedright\Large $6~\text{min}$\end{minipage}\end{tcolorbox}
\end{tcolorbox}
\newpage
\begin{tcolorbox}[enhanced,width=16cm,height=14cm,colback=white,colframe=cmes,boxrule=2pt,arc=0mm,
left=6mm,right=6mm,top=0mm,bottom=4mm,sharp corners,
overlay={\fill[cmes] (frame.north west) rectangle ([yshift=-1.25cm]frame.north east);
\node[anchor=west,text=white,font=\sffamily\bfseries\LARGE] at ([xshift=5mm,yshift=-0.62cm]frame.north west) {Question 77 B};
\node[anchor=east,text=white,font=\sffamily\itshape\large] at ([xshift=-5mm,yshift=-0.62cm]frame.north east) {Grandeurs et mesures};}]
\vspace*{1.75cm}
\begin{minipage}[t][6.9cm][c]{\linewidth}\LARGE\centering
\begin{minipage}{0.95\linewidth}\raggedright\linespread{1.25}\selectfont En prenant $\pi\approx3{,}14$, quel est le périmètre d'un cercle de diamètre $20$ cm~?\end{minipage}
\end{minipage}

\vfill
\begin{tcolorbox}[enhanced,nobeforeafter,width=7.1cm,height=2.1cm,valign=center,colback=cmes!6,colframe=cmes,arc=2.5mm,boxrule=0.9pt,left=3mm]{\Large\bfseries\color{cmes}a)}\quad\begin{minipage}[c]{5.6cm}\raggedright\Large $31{,}4~\text{cm}$\end{minipage}\end{tcolorbox}\hfill\begin{tcolorbox}[enhanced,nobeforeafter,width=7.1cm,height=2.1cm,valign=center,colback=cmes!6,colframe=cmes,arc=2.5mm,boxrule=0.9pt,left=3mm]{\Large\bfseries\color{cmes}b)}\quad\begin{minipage}[c]{5.6cm}\raggedright\Large $314~\text{cm}$\end{minipage}\end{tcolorbox}\par\vspace{2.5mm}
\begin{tcolorbox}[enhanced,nobeforeafter,width=7.1cm,height=2.1cm,valign=center,colback=cmes!6,colframe=cmes,arc=2.5mm,boxrule=0.9pt,left=3mm]{\Large\bfseries\color{cmes}c)}\quad\begin{minipage}[c]{5.6cm}\raggedright\Large $62{,}8~\text{cm}$\end{minipage}\end{tcolorbox}\hfill\begin{tcolorbox}[enhanced,nobeforeafter,width=7.1cm,height=2.1cm,valign=center,colback=cmes!6,colframe=cmes,arc=2.5mm,boxrule=0.9pt,left=3mm]{\Large\bfseries\color{cmes}d)}\quad\begin{minipage}[c]{5.6cm}\raggedright\Large $125{,}6~\text{cm}$\end{minipage}\end{tcolorbox}
\end{tcolorbox}
\newpage
\begin{tcolorbox}[enhanced,width=16cm,height=14cm,colback=white,colframe=cmes,boxrule=2pt,arc=0mm,
left=6mm,right=6mm,top=0mm,bottom=4mm,sharp corners,
overlay={\fill[cmes] (frame.north west) rectangle ([yshift=-1.25cm]frame.north east);
\node[anchor=west,text=white,font=\sffamily\bfseries\LARGE] at ([xshift=5mm,yshift=-0.62cm]frame.north west) {Question 78 B};
\node[anchor=east,text=white,font=\sffamily\itshape\large] at ([xshift=-5mm,yshift=-0.62cm]frame.north east) {Grandeurs et mesures};}]
\vspace*{1.75cm}
\begin{minipage}[t][6.9cm][c]{\linewidth}\LARGE\centering
\begin{minipage}{0.95\linewidth}\raggedright\linespread{1.25}\selectfont En prenant $\pi\approx3{,}14$, quel est le périmètre d'un cercle de diamètre $6$ cm~?\end{minipage}
\end{minipage}

\vfill
\begin{tcolorbox}[enhanced,nobeforeafter,width=7.1cm,height=2.1cm,valign=center,colback=cmes!6,colframe=cmes,arc=2.5mm,boxrule=0.9pt,left=3mm]{\Large\bfseries\color{cmes}a)}\quad\begin{minipage}[c]{5.6cm}\raggedright\Large $28{,}26~\text{cm}$\end{minipage}\end{tcolorbox}\hfill\begin{tcolorbox}[enhanced,nobeforeafter,width=7.1cm,height=2.1cm,valign=center,colback=cmes!6,colframe=cmes,arc=2.5mm,boxrule=0.9pt,left=3mm]{\Large\bfseries\color{cmes}b)}\quad\begin{minipage}[c]{5.6cm}\raggedright\Large $18{,}84~\text{cm}$\end{minipage}\end{tcolorbox}\par\vspace{2.5mm}
\begin{tcolorbox}[enhanced,nobeforeafter,width=7.1cm,height=2.1cm,valign=center,colback=cmes!6,colframe=cmes,arc=2.5mm,boxrule=0.9pt,left=3mm]{\Large\bfseries\color{cmes}c)}\quad\begin{minipage}[c]{5.6cm}\raggedright\Large $37{,}68~\text{cm}$\end{minipage}\end{tcolorbox}\hfill\begin{tcolorbox}[enhanced,nobeforeafter,width=7.1cm,height=2.1cm,valign=center,colback=cmes!6,colframe=cmes,arc=2.5mm,boxrule=0.9pt,left=3mm]{\Large\bfseries\color{cmes}d)}\quad\begin{minipage}[c]{5.6cm}\raggedright\Large $9{,}42~\text{cm}$\end{minipage}\end{tcolorbox}
\end{tcolorbox}
\newpage
\begin{tcolorbox}[enhanced,width=16cm,height=14cm,colback=white,colframe=cmes,boxrule=2pt,arc=0mm,
left=6mm,right=6mm,top=0mm,bottom=4mm,sharp corners,
overlay={\fill[cmes] (frame.north west) rectangle ([yshift=-1.25cm]frame.north east);
\node[anchor=west,text=white,font=\sffamily\bfseries\LARGE] at ([xshift=5mm,yshift=-0.62cm]frame.north west) {Question 79 B};
\node[anchor=east,text=white,font=\sffamily\itshape\large] at ([xshift=-5mm,yshift=-0.62cm]frame.north east) {Grandeurs et mesures};}]
\vspace*{1.75cm}
\begin{minipage}[t][6.9cm][c]{\linewidth}\LARGE\centering
\begin{minipage}{0.95\linewidth}\raggedright\linespread{1.25}\selectfont Un pavé droit mesure $8$ cm de longueur, $5$ cm de largeur et $2$ cm de hauteur. Quel est son volume~?\end{minipage}
\end{minipage}

\vfill
\begin{tcolorbox}[enhanced,nobeforeafter,width=7.1cm,height=2.1cm,valign=center,colback=cmes!6,colframe=cmes,arc=2.5mm,boxrule=0.9pt,left=3mm]{\Large\bfseries\color{cmes}a)}\quad\begin{minipage}[c]{5.6cm}\raggedright\Large $80~\text{cm}^3$\end{minipage}\end{tcolorbox}\hfill\begin{tcolorbox}[enhanced,nobeforeafter,width=7.1cm,height=2.1cm,valign=center,colback=cmes!6,colframe=cmes,arc=2.5mm,boxrule=0.9pt,left=3mm]{\Large\bfseries\color{cmes}b)}\quad\begin{minipage}[c]{5.6cm}\raggedright\Large $132~\text{cm}^3$\end{minipage}\end{tcolorbox}\par\vspace{2.5mm}
\begin{tcolorbox}[enhanced,nobeforeafter,width=7.1cm,height=2.1cm,valign=center,colback=cmes!6,colframe=cmes,arc=2.5mm,boxrule=0.9pt,left=3mm]{\Large\bfseries\color{cmes}c)}\quad\begin{minipage}[c]{5.6cm}\raggedright\Large $40~\text{cm}^3$\end{minipage}\end{tcolorbox}\hfill\begin{tcolorbox}[enhanced,nobeforeafter,width=7.1cm,height=2.1cm,valign=center,colback=cmes!6,colframe=cmes,arc=2.5mm,boxrule=0.9pt,left=3mm]{\Large\bfseries\color{cmes}d)}\quad\begin{minipage}[c]{5.6cm}\raggedright\Large $15~\text{cm}^3$\end{minipage}\end{tcolorbox}
\end{tcolorbox}
\newpage
\begin{tcolorbox}[enhanced,width=16cm,height=14cm,colback=white,colframe=cmes,boxrule=2pt,arc=0mm,
left=6mm,right=6mm,top=0mm,bottom=4mm,sharp corners,
overlay={\fill[cmes] (frame.north west) rectangle ([yshift=-1.25cm]frame.north east);
\node[anchor=west,text=white,font=\sffamily\bfseries\LARGE] at ([xshift=5mm,yshift=-0.62cm]frame.north west) {Question 80 B};
\node[anchor=east,text=white,font=\sffamily\itshape\large] at ([xshift=-5mm,yshift=-0.62cm]frame.north east) {Grandeurs et mesures};}]
\vspace*{1.75cm}
\begin{minipage}[t][6.9cm][c]{\linewidth}\LARGE\centering
\begin{minipage}{0.95\linewidth}\raggedright\linespread{1.25}\selectfont Un pavé droit mesure $7$ cm de longueur, $3$ cm de largeur et $4$ cm de hauteur. Quel est son volume~?\end{minipage}
\end{minipage}

\vfill
\begin{tcolorbox}[enhanced,nobeforeafter,width=7.1cm,height=2.1cm,valign=center,colback=cmes!6,colframe=cmes,arc=2.5mm,boxrule=0.9pt,left=3mm]{\Large\bfseries\color{cmes}a)}\quad\begin{minipage}[c]{5.6cm}\raggedright\Large $21~\text{cm}^3$\end{minipage}\end{tcolorbox}\hfill\begin{tcolorbox}[enhanced,nobeforeafter,width=7.1cm,height=2.1cm,valign=center,colback=cmes!6,colframe=cmes,arc=2.5mm,boxrule=0.9pt,left=3mm]{\Large\bfseries\color{cmes}b)}\quad\begin{minipage}[c]{5.6cm}\raggedright\Large $122~\text{cm}^3$\end{minipage}\end{tcolorbox}\par\vspace{2.5mm}
\begin{tcolorbox}[enhanced,nobeforeafter,width=7.1cm,height=2.1cm,valign=center,colback=cmes!6,colframe=cmes,arc=2.5mm,boxrule=0.9pt,left=3mm]{\Large\bfseries\color{cmes}c)}\quad\begin{minipage}[c]{5.6cm}\raggedright\Large $84~\text{cm}^3$\end{minipage}\end{tcolorbox}\hfill\begin{tcolorbox}[enhanced,nobeforeafter,width=7.1cm,height=2.1cm,valign=center,colback=cmes!6,colframe=cmes,arc=2.5mm,boxrule=0.9pt,left=3mm]{\Large\bfseries\color{cmes}d)}\quad\begin{minipage}[c]{5.6cm}\raggedright\Large $14~\text{cm}^3$\end{minipage}\end{tcolorbox}
\end{tcolorbox}
\newpage
\begin{tcolorbox}[enhanced,width=16cm,height=14cm,colback=white,colframe=cpro,boxrule=2pt,arc=0mm,
left=6mm,right=6mm,top=0mm,bottom=4mm,sharp corners,
overlay={\fill[cpro] (frame.north west) rectangle ([yshift=-1.25cm]frame.north east);
\node[anchor=west,text=white,font=\sffamily\bfseries\LARGE] at ([xshift=5mm,yshift=-0.62cm]frame.north west) {Question 81 B};
\node[anchor=east,text=white,font=\sffamily\itshape\large] at ([xshift=-5mm,yshift=-0.62cm]frame.north east) {Proportionnalité et données};}]
\vspace*{1.75cm}
\begin{minipage}[t][6.9cm][c]{\linewidth}\LARGE\centering
\begin{minipage}{0.95\linewidth}\raggedright\linespread{1.25}\selectfont Pour $3$ personnes, il faut $180$ g de farine. Combien en faut-il pour $9$ personnes~?\end{minipage}
\end{minipage}

\vfill
\begin{tcolorbox}[enhanced,nobeforeafter,width=7.1cm,height=2.1cm,valign=center,colback=cpro!6,colframe=cpro,arc=2.5mm,boxrule=0.9pt,left=3mm]{\Large\bfseries\color{cpro}a)}\quad\begin{minipage}[c]{5.6cm}\raggedright\Large $186~\text{g}$\end{minipage}\end{tcolorbox}\hfill\begin{tcolorbox}[enhanced,nobeforeafter,width=7.1cm,height=2.1cm,valign=center,colback=cpro!6,colframe=cpro,arc=2.5mm,boxrule=0.9pt,left=3mm]{\Large\bfseries\color{cpro}b)}\quad\begin{minipage}[c]{5.6cm}\raggedright\Large $540~\text{g}$\end{minipage}\end{tcolorbox}\par\vspace{2.5mm}
\begin{tcolorbox}[enhanced,nobeforeafter,width=7.1cm,height=2.1cm,valign=center,colback=cpro!6,colframe=cpro,arc=2.5mm,boxrule=0.9pt,left=3mm]{\Large\bfseries\color{cpro}c)}\quad\begin{minipage}[c]{5.6cm}\raggedright\Large $1\,620~\text{g}$\end{minipage}\end{tcolorbox}\hfill\begin{tcolorbox}[enhanced,nobeforeafter,width=7.1cm,height=2.1cm,valign=center,colback=cpro!6,colframe=cpro,arc=2.5mm,boxrule=0.9pt,left=3mm]{\Large\bfseries\color{cpro}d)}\quad\begin{minipage}[c]{5.6cm}\raggedright\Large $60~\text{g}$\end{minipage}\end{tcolorbox}
\end{tcolorbox}
\newpage
\begin{tcolorbox}[enhanced,width=16cm,height=14cm,colback=white,colframe=cpro,boxrule=2pt,arc=0mm,
left=6mm,right=6mm,top=0mm,bottom=4mm,sharp corners,
overlay={\fill[cpro] (frame.north west) rectangle ([yshift=-1.25cm]frame.north east);
\node[anchor=west,text=white,font=\sffamily\bfseries\LARGE] at ([xshift=5mm,yshift=-0.62cm]frame.north west) {Question 82 B};
\node[anchor=east,text=white,font=\sffamily\itshape\large] at ([xshift=-5mm,yshift=-0.62cm]frame.north east) {Proportionnalité et données};}]
\vspace*{1.75cm}
\begin{minipage}[t][6.9cm][c]{\linewidth}\LARGE\centering
\begin{minipage}{0.95\linewidth}\raggedright\linespread{1.25}\selectfont Pour $8$ personnes, il faut $400$ g de farine. Combien en faut-il pour $6$ personnes~?\end{minipage}
\end{minipage}

\vfill
\begin{tcolorbox}[enhanced,nobeforeafter,width=7.1cm,height=2.1cm,valign=center,colback=cpro!6,colframe=cpro,arc=2.5mm,boxrule=0.9pt,left=3mm]{\Large\bfseries\color{cpro}a)}\quad\begin{minipage}[c]{5.6cm}\raggedright\Large $2\,400~\text{g}$\end{minipage}\end{tcolorbox}\hfill\begin{tcolorbox}[enhanced,nobeforeafter,width=7.1cm,height=2.1cm,valign=center,colback=cpro!6,colframe=cpro,arc=2.5mm,boxrule=0.9pt,left=3mm]{\Large\bfseries\color{cpro}b)}\quad\begin{minipage}[c]{5.6cm}\raggedright\Large $300~\text{g}$\end{minipage}\end{tcolorbox}\par\vspace{2.5mm}
\begin{tcolorbox}[enhanced,nobeforeafter,width=7.1cm,height=2.1cm,valign=center,colback=cpro!6,colframe=cpro,arc=2.5mm,boxrule=0.9pt,left=3mm]{\Large\bfseries\color{cpro}c)}\quad\begin{minipage}[c]{5.6cm}\raggedright\Large $50~\text{g}$\end{minipage}\end{tcolorbox}\hfill\begin{tcolorbox}[enhanced,nobeforeafter,width=7.1cm,height=2.1cm,valign=center,colback=cpro!6,colframe=cpro,arc=2.5mm,boxrule=0.9pt,left=3mm]{\Large\bfseries\color{cpro}d)}\quad\begin{minipage}[c]{5.6cm}\raggedright\Large $398~\text{g}$\end{minipage}\end{tcolorbox}
\end{tcolorbox}
\newpage
\begin{tcolorbox}[enhanced,width=16cm,height=14cm,colback=white,colframe=cpro,boxrule=2pt,arc=0mm,
left=6mm,right=6mm,top=0mm,bottom=4mm,sharp corners,
overlay={\fill[cpro] (frame.north west) rectangle ([yshift=-1.25cm]frame.north east);
\node[anchor=west,text=white,font=\sffamily\bfseries\LARGE] at ([xshift=5mm,yshift=-0.62cm]frame.north west) {Question 83 B};
\node[anchor=east,text=white,font=\sffamily\itshape\large] at ([xshift=-5mm,yshift=-0.62cm]frame.north east) {Proportionnalité et données};}]
\vspace*{1.75cm}
\begin{minipage}[t][6.9cm][c]{\linewidth}\LARGE\centering
\begin{minipage}{0.95\linewidth}\raggedright\linespread{1.25}\selectfont $3$ cahiers coûtent $450$ F. Combien coûte un cahier~?\end{minipage}
\end{minipage}

\vfill
\begin{tcolorbox}[enhanced,nobeforeafter,width=7.1cm,height=2.1cm,valign=center,colback=cpro!6,colframe=cpro,arc=2.5mm,boxrule=0.9pt,left=3mm]{\Large\bfseries\color{cpro}a)}\quad\begin{minipage}[c]{5.6cm}\raggedright\Large $150~\text{F}$\end{minipage}\end{tcolorbox}\hfill\begin{tcolorbox}[enhanced,nobeforeafter,width=7.1cm,height=2.1cm,valign=center,colback=cpro!6,colframe=cpro,arc=2.5mm,boxrule=0.9pt,left=3mm]{\Large\bfseries\color{cpro}b)}\quad\begin{minipage}[c]{5.6cm}\raggedright\Large $447~\text{F}$\end{minipage}\end{tcolorbox}\par\vspace{2.5mm}
\begin{tcolorbox}[enhanced,nobeforeafter,width=7.1cm,height=2.1cm,valign=center,colback=cpro!6,colframe=cpro,arc=2.5mm,boxrule=0.9pt,left=3mm]{\Large\bfseries\color{cpro}c)}\quad\begin{minipage}[c]{5.6cm}\raggedright\Large $153~\text{F}$\end{minipage}\end{tcolorbox}\hfill\begin{tcolorbox}[enhanced,nobeforeafter,width=7.1cm,height=2.1cm,valign=center,colback=cpro!6,colframe=cpro,arc=2.5mm,boxrule=0.9pt,left=3mm]{\Large\bfseries\color{cpro}d)}\quad\begin{minipage}[c]{5.6cm}\raggedright\Large $300~\text{F}$\end{minipage}\end{tcolorbox}
\end{tcolorbox}
\newpage
\begin{tcolorbox}[enhanced,width=16cm,height=14cm,colback=white,colframe=cpro,boxrule=2pt,arc=0mm,
left=6mm,right=6mm,top=0mm,bottom=4mm,sharp corners,
overlay={\fill[cpro] (frame.north west) rectangle ([yshift=-1.25cm]frame.north east);
\node[anchor=west,text=white,font=\sffamily\bfseries\LARGE] at ([xshift=5mm,yshift=-0.62cm]frame.north west) {Question 84 B};
\node[anchor=east,text=white,font=\sffamily\itshape\large] at ([xshift=-5mm,yshift=-0.62cm]frame.north east) {Proportionnalité et données};}]
\vspace*{1.75cm}
\begin{minipage}[t][6.9cm][c]{\linewidth}\LARGE\centering
\begin{minipage}{0.95\linewidth}\raggedright\linespread{1.25}\selectfont $6$ cahiers coûtent $840$ F. Combien coûte un cahier~?\end{minipage}
\end{minipage}

\vfill
\begin{tcolorbox}[enhanced,nobeforeafter,width=7.1cm,height=2.1cm,valign=center,colback=cpro!6,colframe=cpro,arc=2.5mm,boxrule=0.9pt,left=3mm]{\Large\bfseries\color{cpro}a)}\quad\begin{minipage}[c]{5.6cm}\raggedright\Large $146~\text{F}$\end{minipage}\end{tcolorbox}\hfill\begin{tcolorbox}[enhanced,nobeforeafter,width=7.1cm,height=2.1cm,valign=center,colback=cpro!6,colframe=cpro,arc=2.5mm,boxrule=0.9pt,left=3mm]{\Large\bfseries\color{cpro}b)}\quad\begin{minipage}[c]{5.6cm}\raggedright\Large $280~\text{F}$\end{minipage}\end{tcolorbox}\par\vspace{2.5mm}
\begin{tcolorbox}[enhanced,nobeforeafter,width=7.1cm,height=2.1cm,valign=center,colback=cpro!6,colframe=cpro,arc=2.5mm,boxrule=0.9pt,left=3mm]{\Large\bfseries\color{cpro}c)}\quad\begin{minipage}[c]{5.6cm}\raggedright\Large $140~\text{F}$\end{minipage}\end{tcolorbox}\hfill\begin{tcolorbox}[enhanced,nobeforeafter,width=7.1cm,height=2.1cm,valign=center,colback=cpro!6,colframe=cpro,arc=2.5mm,boxrule=0.9pt,left=3mm]{\Large\bfseries\color{cpro}d)}\quad\begin{minipage}[c]{5.6cm}\raggedright\Large $834~\text{F}$\end{minipage}\end{tcolorbox}
\end{tcolorbox}
\newpage
\begin{tcolorbox}[enhanced,width=16cm,height=14cm,colback=white,colframe=cpro,boxrule=2pt,arc=0mm,
left=6mm,right=6mm,top=0mm,bottom=4mm,sharp corners,
overlay={\fill[cpro] (frame.north west) rectangle ([yshift=-1.25cm]frame.north east);
\node[anchor=west,text=white,font=\sffamily\bfseries\LARGE] at ([xshift=5mm,yshift=-0.62cm]frame.north west) {Question 85 B};
\node[anchor=east,text=white,font=\sffamily\itshape\large] at ([xshift=-5mm,yshift=-0.62cm]frame.north east) {Proportionnalité et données};}]
\vspace*{1.75cm}
\begin{minipage}[t][6.9cm][c]{\linewidth}\LARGE\centering
\begin{minipage}{0.95\linewidth}\raggedright\linespread{1.25}\selectfont Que vaut $10\,\%$ de $350$~?\end{minipage}
\end{minipage}

\vfill
\begin{tcolorbox}[enhanced,nobeforeafter,width=7.1cm,height=2.1cm,valign=center,colback=cpro!6,colframe=cpro,arc=2.5mm,boxrule=0.9pt,left=3mm]{\Large\bfseries\color{cpro}a)}\quad\begin{minipage}[c]{5.6cm}\raggedright\Large $3{,}5$\end{minipage}\end{tcolorbox}\hfill\begin{tcolorbox}[enhanced,nobeforeafter,width=7.1cm,height=2.1cm,valign=center,colback=cpro!6,colframe=cpro,arc=2.5mm,boxrule=0.9pt,left=3mm]{\Large\bfseries\color{cpro}b)}\quad\begin{minipage}[c]{5.6cm}\raggedright\Large $315$\end{minipage}\end{tcolorbox}\par\vspace{2.5mm}
\begin{tcolorbox}[enhanced,nobeforeafter,width=7.1cm,height=2.1cm,valign=center,colback=cpro!6,colframe=cpro,arc=2.5mm,boxrule=0.9pt,left=3mm]{\Large\bfseries\color{cpro}c)}\quad\begin{minipage}[c]{5.6cm}\raggedright\Large $45$\end{minipage}\end{tcolorbox}\hfill\begin{tcolorbox}[enhanced,nobeforeafter,width=7.1cm,height=2.1cm,valign=center,colback=cpro!6,colframe=cpro,arc=2.5mm,boxrule=0.9pt,left=3mm]{\Large\bfseries\color{cpro}d)}\quad\begin{minipage}[c]{5.6cm}\raggedright\Large $35$\end{minipage}\end{tcolorbox}
\end{tcolorbox}
\newpage
\begin{tcolorbox}[enhanced,width=16cm,height=14cm,colback=white,colframe=cpro,boxrule=2pt,arc=0mm,
left=6mm,right=6mm,top=0mm,bottom=4mm,sharp corners,
overlay={\fill[cpro] (frame.north west) rectangle ([yshift=-1.25cm]frame.north east);
\node[anchor=west,text=white,font=\sffamily\bfseries\LARGE] at ([xshift=5mm,yshift=-0.62cm]frame.north west) {Question 86 B};
\node[anchor=east,text=white,font=\sffamily\itshape\large] at ([xshift=-5mm,yshift=-0.62cm]frame.north east) {Proportionnalité et données};}]
\vspace*{1.75cm}
\begin{minipage}[t][6.9cm][c]{\linewidth}\LARGE\centering
\begin{minipage}{0.95\linewidth}\raggedright\linespread{1.25}\selectfont Que vaut $20\,\%$ de $45$~?\end{minipage}
\end{minipage}

\vfill
\begin{tcolorbox}[enhanced,nobeforeafter,width=7.1cm,height=2.1cm,valign=center,colback=cpro!6,colframe=cpro,arc=2.5mm,boxrule=0.9pt,left=3mm]{\Large\bfseries\color{cpro}a)}\quad\begin{minipage}[c]{5.6cm}\raggedright\Large $36$\end{minipage}\end{tcolorbox}\hfill\begin{tcolorbox}[enhanced,nobeforeafter,width=7.1cm,height=2.1cm,valign=center,colback=cpro!6,colframe=cpro,arc=2.5mm,boxrule=0.9pt,left=3mm]{\Large\bfseries\color{cpro}b)}\quad\begin{minipage}[c]{5.6cm}\raggedright\Large $2{,}25$\end{minipage}\end{tcolorbox}\par\vspace{2.5mm}
\begin{tcolorbox}[enhanced,nobeforeafter,width=7.1cm,height=2.1cm,valign=center,colback=cpro!6,colframe=cpro,arc=2.5mm,boxrule=0.9pt,left=3mm]{\Large\bfseries\color{cpro}c)}\quad\begin{minipage}[c]{5.6cm}\raggedright\Large $9$\end{minipage}\end{tcolorbox}\hfill\begin{tcolorbox}[enhanced,nobeforeafter,width=7.1cm,height=2.1cm,valign=center,colback=cpro!6,colframe=cpro,arc=2.5mm,boxrule=0.9pt,left=3mm]{\Large\bfseries\color{cpro}d)}\quad\begin{minipage}[c]{5.6cm}\raggedright\Large $19$\end{minipage}\end{tcolorbox}
\end{tcolorbox}
\newpage
\begin{tcolorbox}[enhanced,width=16cm,height=14cm,colback=white,colframe=cpro,boxrule=2pt,arc=0mm,
left=6mm,right=6mm,top=0mm,bottom=4mm,sharp corners,
overlay={\fill[cpro] (frame.north west) rectangle ([yshift=-1.25cm]frame.north east);
\node[anchor=west,text=white,font=\sffamily\bfseries\LARGE] at ([xshift=5mm,yshift=-0.62cm]frame.north west) {Question 87 B};
\node[anchor=east,text=white,font=\sffamily\itshape\large] at ([xshift=-5mm,yshift=-0.62cm]frame.north east) {Proportionnalité et données};}]
\vspace*{1.75cm}
\begin{minipage}[t][6.9cm][c]{\linewidth}\LARGE\centering
\begin{minipage}{0.95\linewidth}\raggedright\linespread{1.25}\selectfont Sur un plan à l'échelle $\dfrac{1}{200}$, une longueur mesure $4$ cm. Quelle est la longueur réelle~?\end{minipage}
\end{minipage}

\vfill
\begin{tcolorbox}[enhanced,nobeforeafter,width=7.1cm,height=2.1cm,valign=center,colback=cpro!6,colframe=cpro,arc=2.5mm,boxrule=0.9pt,left=3mm]{\Large\bfseries\color{cpro}a)}\quad\begin{minipage}[c]{5.6cm}\raggedright\Large $0{,}8~\text{m}$\end{minipage}\end{tcolorbox}\hfill\begin{tcolorbox}[enhanced,nobeforeafter,width=7.1cm,height=2.1cm,valign=center,colback=cpro!6,colframe=cpro,arc=2.5mm,boxrule=0.9pt,left=3mm]{\Large\bfseries\color{cpro}b)}\quad\begin{minipage}[c]{5.6cm}\raggedright\Large $800~\text{m}$\end{minipage}\end{tcolorbox}\par\vspace{2.5mm}
\begin{tcolorbox}[enhanced,nobeforeafter,width=7.1cm,height=2.1cm,valign=center,colback=cpro!6,colframe=cpro,arc=2.5mm,boxrule=0.9pt,left=3mm]{\Large\bfseries\color{cpro}c)}\quad\begin{minipage}[c]{5.6cm}\raggedright\Large $8~\text{m}$\end{minipage}\end{tcolorbox}\hfill\begin{tcolorbox}[enhanced,nobeforeafter,width=7.1cm,height=2.1cm,valign=center,colback=cpro!6,colframe=cpro,arc=2.5mm,boxrule=0.9pt,left=3mm]{\Large\bfseries\color{cpro}d)}\quad\begin{minipage}[c]{5.6cm}\raggedright\Large $80~\text{m}$\end{minipage}\end{tcolorbox}
\end{tcolorbox}
\newpage
\begin{tcolorbox}[enhanced,width=16cm,height=14cm,colback=white,colframe=cpro,boxrule=2pt,arc=0mm,
left=6mm,right=6mm,top=0mm,bottom=4mm,sharp corners,
overlay={\fill[cpro] (frame.north west) rectangle ([yshift=-1.25cm]frame.north east);
\node[anchor=west,text=white,font=\sffamily\bfseries\LARGE] at ([xshift=5mm,yshift=-0.62cm]frame.north west) {Question 88 B};
\node[anchor=east,text=white,font=\sffamily\itshape\large] at ([xshift=-5mm,yshift=-0.62cm]frame.north east) {Proportionnalité et données};}]
\vspace*{1.75cm}
\begin{minipage}[t][6.9cm][c]{\linewidth}\LARGE\centering
\begin{minipage}{0.95\linewidth}\raggedright\linespread{1.25}\selectfont Sur un plan à l'échelle $\dfrac{1}{50}$, une longueur mesure $8$ cm. Quelle est la longueur réelle~?\end{minipage}
\end{minipage}

\vfill
\begin{tcolorbox}[enhanced,nobeforeafter,width=7.1cm,height=2.1cm,valign=center,colback=cpro!6,colframe=cpro,arc=2.5mm,boxrule=0.9pt,left=3mm]{\Large\bfseries\color{cpro}a)}\quad\begin{minipage}[c]{5.6cm}\raggedright\Large $400~\text{m}$\end{minipage}\end{tcolorbox}\hfill\begin{tcolorbox}[enhanced,nobeforeafter,width=7.1cm,height=2.1cm,valign=center,colback=cpro!6,colframe=cpro,arc=2.5mm,boxrule=0.9pt,left=3mm]{\Large\bfseries\color{cpro}b)}\quad\begin{minipage}[c]{5.6cm}\raggedright\Large $0{,}4~\text{m}$\end{minipage}\end{tcolorbox}\par\vspace{2.5mm}
\begin{tcolorbox}[enhanced,nobeforeafter,width=7.1cm,height=2.1cm,valign=center,colback=cpro!6,colframe=cpro,arc=2.5mm,boxrule=0.9pt,left=3mm]{\Large\bfseries\color{cpro}c)}\quad\begin{minipage}[c]{5.6cm}\raggedright\Large $4~\text{m}$\end{minipage}\end{tcolorbox}\hfill\begin{tcolorbox}[enhanced,nobeforeafter,width=7.1cm,height=2.1cm,valign=center,colback=cpro!6,colframe=cpro,arc=2.5mm,boxrule=0.9pt,left=3mm]{\Large\bfseries\color{cpro}d)}\quad\begin{minipage}[c]{5.6cm}\raggedright\Large $40~\text{m}$\end{minipage}\end{tcolorbox}
\end{tcolorbox}
\newpage
\begin{tcolorbox}[enhanced,width=16cm,height=14cm,colback=white,colframe=cpro,boxrule=2pt,arc=0mm,
left=6mm,right=6mm,top=0mm,bottom=4mm,sharp corners,
overlay={\fill[cpro] (frame.north west) rectangle ([yshift=-1.25cm]frame.north east);
\node[anchor=west,text=white,font=\sffamily\bfseries\LARGE] at ([xshift=5mm,yshift=-0.62cm]frame.north west) {Question 89 B};
\node[anchor=east,text=white,font=\sffamily\itshape\large] at ([xshift=-5mm,yshift=-0.62cm]frame.north east) {Proportionnalité et données};}]
\vspace*{1.75cm}
\begin{minipage}[t][6.9cm][c]{\linewidth}\Large\centering
\begin{minipage}{0.95\linewidth}\raggedright\linespread{1.25}\selectfont Ce tableau est un tableau de proportionnalité~: $\begin{array}{|c|c|c|}\hline 2 & 4 & 9\\\hline 12 & 24 & 54\\\hline\end{array}$. Par quel nombre multiplie-t-on la 1re ligne pour obtenir la 2e~?\end{minipage}
\end{minipage}

\vfill
\begin{tcolorbox}[enhanced,nobeforeafter,width=7.1cm,height=2.1cm,valign=center,colback=cpro!6,colframe=cpro,arc=2.5mm,boxrule=0.9pt,left=3mm]{\Large\bfseries\color{cpro}a)}\quad\begin{minipage}[c]{5.6cm}\raggedright\Large $6$\end{minipage}\end{tcolorbox}\hfill\begin{tcolorbox}[enhanced,nobeforeafter,width=7.1cm,height=2.1cm,valign=center,colback=cpro!6,colframe=cpro,arc=2.5mm,boxrule=0.9pt,left=3mm]{\Large\bfseries\color{cpro}b)}\quad\begin{minipage}[c]{5.6cm}\raggedright\Large $7$\end{minipage}\end{tcolorbox}\par\vspace{2.5mm}
\begin{tcolorbox}[enhanced,nobeforeafter,width=7.1cm,height=2.1cm,valign=center,colback=cpro!6,colframe=cpro,arc=2.5mm,boxrule=0.9pt,left=3mm]{\Large\bfseries\color{cpro}c)}\quad\begin{minipage}[c]{5.6cm}\raggedright\Large $12$\end{minipage}\end{tcolorbox}\hfill\begin{tcolorbox}[enhanced,nobeforeafter,width=7.1cm,height=2.1cm,valign=center,colback=cpro!6,colframe=cpro,arc=2.5mm,boxrule=0.9pt,left=3mm]{\Large\bfseries\color{cpro}d)}\quad\begin{minipage}[c]{5.6cm}\raggedright\Large $\dfrac{1}{6}$\end{minipage}\end{tcolorbox}
\end{tcolorbox}
\newpage
\begin{tcolorbox}[enhanced,width=16cm,height=14cm,colback=white,colframe=cpro,boxrule=2pt,arc=0mm,
left=6mm,right=6mm,top=0mm,bottom=4mm,sharp corners,
overlay={\fill[cpro] (frame.north west) rectangle ([yshift=-1.25cm]frame.north east);
\node[anchor=west,text=white,font=\sffamily\bfseries\LARGE] at ([xshift=5mm,yshift=-0.62cm]frame.north west) {Question 90 B};
\node[anchor=east,text=white,font=\sffamily\itshape\large] at ([xshift=-5mm,yshift=-0.62cm]frame.north east) {Proportionnalité et données};}]
\vspace*{1.75cm}
\begin{minipage}[t][6.9cm][c]{\linewidth}\Large\centering
\begin{minipage}{0.95\linewidth}\raggedright\linespread{1.25}\selectfont Ce tableau est un tableau de proportionnalité~: $\begin{array}{|c|c|c|}\hline 2 & 4 & 6\\\hline 5 & 10 & 15\\\hline\end{array}$. Par quel nombre multiplie-t-on la 1re ligne pour obtenir la 2e~?\end{minipage}
\end{minipage}

\vfill
\begin{tcolorbox}[enhanced,nobeforeafter,width=7.1cm,height=2.1cm,valign=center,colback=cpro!6,colframe=cpro,arc=2.5mm,boxrule=0.9pt,left=3mm]{\Large\bfseries\color{cpro}a)}\quad\begin{minipage}[c]{5.6cm}\raggedright\Large $\dfrac{2}{5}$\end{minipage}\end{tcolorbox}\hfill\begin{tcolorbox}[enhanced,nobeforeafter,width=7.1cm,height=2.1cm,valign=center,colback=cpro!6,colframe=cpro,arc=2.5mm,boxrule=0.9pt,left=3mm]{\Large\bfseries\color{cpro}b)}\quad\begin{minipage}[c]{5.6cm}\raggedright\Large $3{,}5$\end{minipage}\end{tcolorbox}\par\vspace{2.5mm}
\begin{tcolorbox}[enhanced,nobeforeafter,width=7.1cm,height=2.1cm,valign=center,colback=cpro!6,colframe=cpro,arc=2.5mm,boxrule=0.9pt,left=3mm]{\Large\bfseries\color{cpro}c)}\quad\begin{minipage}[c]{5.6cm}\raggedright\Large $5$\end{minipage}\end{tcolorbox}\hfill\begin{tcolorbox}[enhanced,nobeforeafter,width=7.1cm,height=2.1cm,valign=center,colback=cpro!6,colframe=cpro,arc=2.5mm,boxrule=0.9pt,left=3mm]{\Large\bfseries\color{cpro}d)}\quad\begin{minipage}[c]{5.6cm}\raggedright\Large $2{,}5$\end{minipage}\end{tcolorbox}
\end{tcolorbox}
\newpage
\begin{tcolorbox}[enhanced,width=16cm,height=14cm,colback=white,colframe=cpro,boxrule=2pt,arc=0mm,
left=6mm,right=6mm,top=0mm,bottom=4mm,sharp corners,
overlay={\fill[cpro] (frame.north west) rectangle ([yshift=-1.25cm]frame.north east);
\node[anchor=west,text=white,font=\sffamily\bfseries\LARGE] at ([xshift=5mm,yshift=-0.62cm]frame.north west) {Question 91 B};
\node[anchor=east,text=white,font=\sffamily\itshape\large] at ([xshift=-5mm,yshift=-0.62cm]frame.north east) {Proportionnalité et données};}]
\vspace*{1.75cm}
\begin{minipage}[t][6.9cm][c]{\linewidth}\Large\centering
\begin{minipage}{0.95\linewidth}\raggedright\linespread{1.25}\selectfont Dans une classe, les élèves viennent au collège ainsi~: $7$ à pied, $7$ en bus, $11$ à vélo et $3$ en voiture. Combien y a-t-il d'élèves dans la classe~?\end{minipage}
\end{minipage}

\vfill
\begin{tcolorbox}[enhanced,nobeforeafter,width=7.1cm,height=2.1cm,valign=center,colback=cpro!6,colframe=cpro,arc=2.5mm,boxrule=0.9pt,left=3mm]{\Large\bfseries\color{cpro}a)}\quad\begin{minipage}[c]{5.6cm}\raggedright\Large $25$\end{minipage}\end{tcolorbox}\hfill\begin{tcolorbox}[enhanced,nobeforeafter,width=7.1cm,height=2.1cm,valign=center,colback=cpro!6,colframe=cpro,arc=2.5mm,boxrule=0.9pt,left=3mm]{\Large\bfseries\color{cpro}b)}\quad\begin{minipage}[c]{5.6cm}\raggedright\Large $11$\end{minipage}\end{tcolorbox}\par\vspace{2.5mm}
\begin{tcolorbox}[enhanced,nobeforeafter,width=7.1cm,height=2.1cm,valign=center,colback=cpro!6,colframe=cpro,arc=2.5mm,boxrule=0.9pt,left=3mm]{\Large\bfseries\color{cpro}c)}\quad\begin{minipage}[c]{5.6cm}\raggedright\Large $28$\end{minipage}\end{tcolorbox}\hfill\begin{tcolorbox}[enhanced,nobeforeafter,width=7.1cm,height=2.1cm,valign=center,colback=cpro!6,colframe=cpro,arc=2.5mm,boxrule=0.9pt,left=3mm]{\Large\bfseries\color{cpro}d)}\quad\begin{minipage}[c]{5.6cm}\raggedright\Large $29$\end{minipage}\end{tcolorbox}
\end{tcolorbox}
\newpage
\begin{tcolorbox}[enhanced,width=16cm,height=14cm,colback=white,colframe=cpro,boxrule=2pt,arc=0mm,
left=6mm,right=6mm,top=0mm,bottom=4mm,sharp corners,
overlay={\fill[cpro] (frame.north west) rectangle ([yshift=-1.25cm]frame.north east);
\node[anchor=west,text=white,font=\sffamily\bfseries\LARGE] at ([xshift=5mm,yshift=-0.62cm]frame.north west) {Question 92 B};
\node[anchor=east,text=white,font=\sffamily\itshape\large] at ([xshift=-5mm,yshift=-0.62cm]frame.north east) {Proportionnalité et données};}]
\vspace*{1.75cm}
\begin{minipage}[t][6.9cm][c]{\linewidth}\LARGE\centering
\begin{minipage}{0.95\linewidth}\raggedright\linespread{1.25}\selectfont Dans une classe, les élèves viennent au collège ainsi~: $20$ à pied, $15$ en bus et $8$ à vélo. Combien y a-t-il d'élèves dans la classe~?\end{minipage}
\end{minipage}

\vfill
\begin{tcolorbox}[enhanced,nobeforeafter,width=7.1cm,height=2.1cm,valign=center,colback=cpro!6,colframe=cpro,arc=2.5mm,boxrule=0.9pt,left=3mm]{\Large\bfseries\color{cpro}a)}\quad\begin{minipage}[c]{5.6cm}\raggedright\Large $43$\end{minipage}\end{tcolorbox}\hfill\begin{tcolorbox}[enhanced,nobeforeafter,width=7.1cm,height=2.1cm,valign=center,colback=cpro!6,colframe=cpro,arc=2.5mm,boxrule=0.9pt,left=3mm]{\Large\bfseries\color{cpro}b)}\quad\begin{minipage}[c]{5.6cm}\raggedright\Large $20$\end{minipage}\end{tcolorbox}\par\vspace{2.5mm}
\begin{tcolorbox}[enhanced,nobeforeafter,width=7.1cm,height=2.1cm,valign=center,colback=cpro!6,colframe=cpro,arc=2.5mm,boxrule=0.9pt,left=3mm]{\Large\bfseries\color{cpro}c)}\quad\begin{minipage}[c]{5.6cm}\raggedright\Large $35$\end{minipage}\end{tcolorbox}\hfill\begin{tcolorbox}[enhanced,nobeforeafter,width=7.1cm,height=2.1cm,valign=center,colback=cpro!6,colframe=cpro,arc=2.5mm,boxrule=0.9pt,left=3mm]{\Large\bfseries\color{cpro}d)}\quad\begin{minipage}[c]{5.6cm}\raggedright\Large $44$\end{minipage}\end{tcolorbox}
\end{tcolorbox}
\newpage
\begin{tcolorbox}[enhanced,width=16cm,height=14cm,colback=white,colframe=cpro,boxrule=2pt,arc=0mm,
left=6mm,right=6mm,top=0mm,bottom=4mm,sharp corners,
overlay={\fill[cpro] (frame.north west) rectangle ([yshift=-1.25cm]frame.north east);
\node[anchor=west,text=white,font=\sffamily\bfseries\LARGE] at ([xshift=5mm,yshift=-0.62cm]frame.north west) {Question 93 B};
\node[anchor=east,text=white,font=\sffamily\itshape\large] at ([xshift=-5mm,yshift=-0.62cm]frame.north east) {Proportionnalité et données};}]
\vspace*{1.75cm}
\begin{minipage}[t][6.9cm][c]{\linewidth}\LARGE\centering
\begin{minipage}{0.95\linewidth}\raggedright\linespread{1.25}\selectfont Ce tableau est-il un tableau de proportionnalité~? $\begin{array}{|c|c|}\hline 4 & 10\\\hline 10 & 25\\\hline\end{array}$\end{minipage}
\end{minipage}

\vfill
\begin{tcolorbox}[enhanced,nobeforeafter,width=7.1cm,height=2.1cm,valign=center,colback=cpro!6,colframe=cpro,arc=2.5mm,boxrule=0.9pt,left=3mm]{\Large\bfseries\color{cpro}a)}\quad\begin{minipage}[c]{5.6cm}\raggedright\Large il manque une colonne\end{minipage}\end{tcolorbox}\hfill\begin{tcolorbox}[enhanced,nobeforeafter,width=7.1cm,height=2.1cm,valign=center,colback=cpro!6,colframe=cpro,arc=2.5mm,boxrule=0.9pt,left=3mm]{\Large\bfseries\color{cpro}b)}\quad\begin{minipage}[c]{5.6cm}\raggedright\Large on ne peut pas savoir\end{minipage}\end{tcolorbox}\par\vspace{2.5mm}
\begin{tcolorbox}[enhanced,nobeforeafter,width=7.1cm,height=2.1cm,valign=center,colback=cpro!6,colframe=cpro,arc=2.5mm,boxrule=0.9pt,left=3mm]{\Large\bfseries\color{cpro}c)}\quad\begin{minipage}[c]{5.6cm}\raggedright\Large oui\end{minipage}\end{tcolorbox}\hfill\begin{tcolorbox}[enhanced,nobeforeafter,width=7.1cm,height=2.1cm,valign=center,colback=cpro!6,colframe=cpro,arc=2.5mm,boxrule=0.9pt,left=3mm]{\Large\bfseries\color{cpro}d)}\quad\begin{minipage}[c]{5.6cm}\raggedright\Large non\end{minipage}\end{tcolorbox}
\end{tcolorbox}
\newpage
\begin{tcolorbox}[enhanced,width=16cm,height=14cm,colback=white,colframe=cpro,boxrule=2pt,arc=0mm,
left=6mm,right=6mm,top=0mm,bottom=4mm,sharp corners,
overlay={\fill[cpro] (frame.north west) rectangle ([yshift=-1.25cm]frame.north east);
\node[anchor=west,text=white,font=\sffamily\bfseries\LARGE] at ([xshift=5mm,yshift=-0.62cm]frame.north west) {Question 94 B};
\node[anchor=east,text=white,font=\sffamily\itshape\large] at ([xshift=-5mm,yshift=-0.62cm]frame.north east) {Proportionnalité et données};}]
\vspace*{1.75cm}
\begin{minipage}[t][6.9cm][c]{\linewidth}\LARGE\centering
\begin{minipage}{0.95\linewidth}\raggedright\linespread{1.25}\selectfont Ce tableau est-il un tableau de proportionnalité~? $\begin{array}{|c|c|}\hline 5 & 7\\\hline 15 & 20\\\hline\end{array}$\end{minipage}
\end{minipage}

\vfill
\begin{tcolorbox}[enhanced,nobeforeafter,width=7.1cm,height=2.1cm,valign=center,colback=cpro!6,colframe=cpro,arc=2.5mm,boxrule=0.9pt,left=3mm]{\Large\bfseries\color{cpro}a)}\quad\begin{minipage}[c]{5.6cm}\raggedright\Large non\end{minipage}\end{tcolorbox}\hfill\begin{tcolorbox}[enhanced,nobeforeafter,width=7.1cm,height=2.1cm,valign=center,colback=cpro!6,colframe=cpro,arc=2.5mm,boxrule=0.9pt,left=3mm]{\Large\bfseries\color{cpro}b)}\quad\begin{minipage}[c]{5.6cm}\raggedright\Large on ne peut pas savoir\end{minipage}\end{tcolorbox}\par\vspace{2.5mm}
\begin{tcolorbox}[enhanced,nobeforeafter,width=7.1cm,height=2.1cm,valign=center,colback=cpro!6,colframe=cpro,arc=2.5mm,boxrule=0.9pt,left=3mm]{\Large\bfseries\color{cpro}c)}\quad\begin{minipage}[c]{5.6cm}\raggedright\Large il manque une colonne\end{minipage}\end{tcolorbox}\hfill\begin{tcolorbox}[enhanced,nobeforeafter,width=7.1cm,height=2.1cm,valign=center,colback=cpro!6,colframe=cpro,arc=2.5mm,boxrule=0.9pt,left=3mm]{\Large\bfseries\color{cpro}d)}\quad\begin{minipage}[c]{5.6cm}\raggedright\Large oui\end{minipage}\end{tcolorbox}
\end{tcolorbox}
\newpage
\begin{tcolorbox}[enhanced,width=16cm,height=14cm,colback=white,colframe=cpro,boxrule=2pt,arc=0mm,
left=6mm,right=6mm,top=0mm,bottom=4mm,sharp corners,
overlay={\fill[cpro] (frame.north west) rectangle ([yshift=-1.25cm]frame.north east);
\node[anchor=west,text=white,font=\sffamily\bfseries\LARGE] at ([xshift=5mm,yshift=-0.62cm]frame.north west) {Question 95 B};
\node[anchor=east,text=white,font=\sffamily\itshape\large] at ([xshift=-5mm,yshift=-0.62cm]frame.north east) {Proportionnalité et données};}]
\vspace*{1.75cm}
\begin{minipage}[t][6.9cm][c]{\linewidth}\LARGE\centering
\begin{minipage}{0.95\linewidth}\raggedright\linespread{1.25}\selectfont Un robinet remplit $18$ L en $6$ min. Quel volume remplit-il en $10$ min, au même débit~?\end{minipage}
\end{minipage}

\vfill
\begin{tcolorbox}[enhanced,nobeforeafter,width=7.1cm,height=2.1cm,valign=center,colback=cpro!6,colframe=cpro,arc=2.5mm,boxrule=0.9pt,left=3mm]{\Large\bfseries\color{cpro}a)}\quad\begin{minipage}[c]{5.6cm}\raggedright\Large $30~\text{L}$\end{minipage}\end{tcolorbox}\hfill\begin{tcolorbox}[enhanced,nobeforeafter,width=7.1cm,height=2.1cm,valign=center,colback=cpro!6,colframe=cpro,arc=2.5mm,boxrule=0.9pt,left=3mm]{\Large\bfseries\color{cpro}b)}\quad\begin{minipage}[c]{5.6cm}\raggedright\Large $180~\text{L}$\end{minipage}\end{tcolorbox}\par\vspace{2.5mm}
\begin{tcolorbox}[enhanced,nobeforeafter,width=7.1cm,height=2.1cm,valign=center,colback=cpro!6,colframe=cpro,arc=2.5mm,boxrule=0.9pt,left=3mm]{\Large\bfseries\color{cpro}c)}\quad\begin{minipage}[c]{5.6cm}\raggedright\Large $22~\text{L}$\end{minipage}\end{tcolorbox}\hfill\begin{tcolorbox}[enhanced,nobeforeafter,width=7.1cm,height=2.1cm,valign=center,colback=cpro!6,colframe=cpro,arc=2.5mm,boxrule=0.9pt,left=3mm]{\Large\bfseries\color{cpro}d)}\quad\begin{minipage}[c]{5.6cm}\raggedright\Large $3~\text{L}$\end{minipage}\end{tcolorbox}
\end{tcolorbox}
\newpage
\begin{tcolorbox}[enhanced,width=16cm,height=14cm,colback=white,colframe=cpro,boxrule=2pt,arc=0mm,
left=6mm,right=6mm,top=0mm,bottom=4mm,sharp corners,
overlay={\fill[cpro] (frame.north west) rectangle ([yshift=-1.25cm]frame.north east);
\node[anchor=west,text=white,font=\sffamily\bfseries\LARGE] at ([xshift=5mm,yshift=-0.62cm]frame.north west) {Question 96 B};
\node[anchor=east,text=white,font=\sffamily\itshape\large] at ([xshift=-5mm,yshift=-0.62cm]frame.north east) {Proportionnalité et données};}]
\vspace*{1.75cm}
\begin{minipage}[t][6.9cm][c]{\linewidth}\LARGE\centering
\begin{minipage}{0.95\linewidth}\raggedright\linespread{1.25}\selectfont Un robinet remplit $9$ L en $2$ min. Quel volume remplit-il en $7$ min, au même débit~?\end{minipage}
\end{minipage}

\vfill
\begin{tcolorbox}[enhanced,nobeforeafter,width=7.1cm,height=2.1cm,valign=center,colback=cpro!6,colframe=cpro,arc=2.5mm,boxrule=0.9pt,left=3mm]{\Large\bfseries\color{cpro}a)}\quad\begin{minipage}[c]{5.6cm}\raggedright\Large $4{,}5~\text{L}$\end{minipage}\end{tcolorbox}\hfill\begin{tcolorbox}[enhanced,nobeforeafter,width=7.1cm,height=2.1cm,valign=center,colback=cpro!6,colframe=cpro,arc=2.5mm,boxrule=0.9pt,left=3mm]{\Large\bfseries\color{cpro}b)}\quad\begin{minipage}[c]{5.6cm}\raggedright\Large $14~\text{L}$\end{minipage}\end{tcolorbox}\par\vspace{2.5mm}
\begin{tcolorbox}[enhanced,nobeforeafter,width=7.1cm,height=2.1cm,valign=center,colback=cpro!6,colframe=cpro,arc=2.5mm,boxrule=0.9pt,left=3mm]{\Large\bfseries\color{cpro}c)}\quad\begin{minipage}[c]{5.6cm}\raggedright\Large $31{,}5~\text{L}$\end{minipage}\end{tcolorbox}\hfill\begin{tcolorbox}[enhanced,nobeforeafter,width=7.1cm,height=2.1cm,valign=center,colback=cpro!6,colframe=cpro,arc=2.5mm,boxrule=0.9pt,left=3mm]{\Large\bfseries\color{cpro}d)}\quad\begin{minipage}[c]{5.6cm}\raggedright\Large $63~\text{L}$\end{minipage}\end{tcolorbox}
\end{tcolorbox}
\newpage
\begin{tcolorbox}[enhanced,width=16cm,height=14cm,colback=white,colframe=cnum,boxrule=2pt,arc=0mm,
left=6mm,right=6mm,top=0mm,bottom=4mm,sharp corners,
overlay={\fill[cnum] (frame.north west) rectangle ([yshift=-1.25cm]frame.north east);
\node[anchor=west,text=white,font=\sffamily\bfseries\LARGE] at ([xshift=5mm,yshift=-0.62cm]frame.north west) {Question 1 C};
\node[anchor=east,text=white,font=\sffamily\itshape\large] at ([xshift=-5mm,yshift=-0.62cm]frame.north east) {Nombres entiers et décimaux};}]
\vspace*{1.75cm}
\begin{minipage}[t][6.9cm][c]{\linewidth}\LARGE\centering
\begin{minipage}{0.95\linewidth}\raggedright\linespread{1.25}\selectfont Dans le nombre $902{,}375$, quel est le chiffre des dixièmes~?\end{minipage}
\end{minipage}

\vfill
\begin{tcolorbox}[enhanced,nobeforeafter,width=7.1cm,height=2.1cm,valign=center,colback=cnum!6,colframe=cnum,arc=2.5mm,boxrule=0.9pt,left=3mm]{\Large\bfseries\color{cnum}a)}\quad\begin{minipage}[c]{5.6cm}\raggedright\Large $9$\end{minipage}\end{tcolorbox}\hfill\begin{tcolorbox}[enhanced,nobeforeafter,width=7.1cm,height=2.1cm,valign=center,colback=cnum!6,colframe=cnum,arc=2.5mm,boxrule=0.9pt,left=3mm]{\Large\bfseries\color{cnum}b)}\quad\begin{minipage}[c]{5.6cm}\raggedright\Large $2$\end{minipage}\end{tcolorbox}\par\vspace{2.5mm}
\begin{tcolorbox}[enhanced,nobeforeafter,width=7.1cm,height=2.1cm,valign=center,colback=cnum!6,colframe=cnum,arc=2.5mm,boxrule=0.9pt,left=3mm]{\Large\bfseries\color{cnum}c)}\quad\begin{minipage}[c]{5.6cm}\raggedright\Large $0$\end{minipage}\end{tcolorbox}\hfill\begin{tcolorbox}[enhanced,nobeforeafter,width=7.1cm,height=2.1cm,valign=center,colback=cnum!6,colframe=cnum,arc=2.5mm,boxrule=0.9pt,left=3mm]{\Large\bfseries\color{cnum}d)}\quad\begin{minipage}[c]{5.6cm}\raggedright\Large $3$\end{minipage}\end{tcolorbox}
\end{tcolorbox}
\newpage
\begin{tcolorbox}[enhanced,width=16cm,height=14cm,colback=white,colframe=cnum,boxrule=2pt,arc=0mm,
left=6mm,right=6mm,top=0mm,bottom=4mm,sharp corners,
overlay={\fill[cnum] (frame.north west) rectangle ([yshift=-1.25cm]frame.north east);
\node[anchor=west,text=white,font=\sffamily\bfseries\LARGE] at ([xshift=5mm,yshift=-0.62cm]frame.north west) {Question 2 C};
\node[anchor=east,text=white,font=\sffamily\itshape\large] at ([xshift=-5mm,yshift=-0.62cm]frame.north east) {Nombres entiers et décimaux};}]
\vspace*{1.75cm}
\begin{minipage}[t][6.9cm][c]{\linewidth}\LARGE\centering
\begin{minipage}{0.95\linewidth}\raggedright\linespread{1.25}\selectfont Dans le nombre $18\,064{,}2$, quel est le chiffre des milliers~?\end{minipage}
\end{minipage}

\vfill
\begin{tcolorbox}[enhanced,nobeforeafter,width=7.1cm,height=2.1cm,valign=center,colback=cnum!6,colframe=cnum,arc=2.5mm,boxrule=0.9pt,left=3mm]{\Large\bfseries\color{cnum}a)}\quad\begin{minipage}[c]{5.6cm}\raggedright\Large $1$\end{minipage}\end{tcolorbox}\hfill\begin{tcolorbox}[enhanced,nobeforeafter,width=7.1cm,height=2.1cm,valign=center,colback=cnum!6,colframe=cnum,arc=2.5mm,boxrule=0.9pt,left=3mm]{\Large\bfseries\color{cnum}b)}\quad\begin{minipage}[c]{5.6cm}\raggedright\Large $6$\end{minipage}\end{tcolorbox}\par\vspace{2.5mm}
\begin{tcolorbox}[enhanced,nobeforeafter,width=7.1cm,height=2.1cm,valign=center,colback=cnum!6,colframe=cnum,arc=2.5mm,boxrule=0.9pt,left=3mm]{\Large\bfseries\color{cnum}c)}\quad\begin{minipage}[c]{5.6cm}\raggedright\Large $0$\end{minipage}\end{tcolorbox}\hfill\begin{tcolorbox}[enhanced,nobeforeafter,width=7.1cm,height=2.1cm,valign=center,colback=cnum!6,colframe=cnum,arc=2.5mm,boxrule=0.9pt,left=3mm]{\Large\bfseries\color{cnum}d)}\quad\begin{minipage}[c]{5.6cm}\raggedright\Large $8$\end{minipage}\end{tcolorbox}
\end{tcolorbox}
\newpage
\begin{tcolorbox}[enhanced,width=16cm,height=14cm,colback=white,colframe=cnum,boxrule=2pt,arc=0mm,
left=6mm,right=6mm,top=0mm,bottom=4mm,sharp corners,
overlay={\fill[cnum] (frame.north west) rectangle ([yshift=-1.25cm]frame.north east);
\node[anchor=west,text=white,font=\sffamily\bfseries\LARGE] at ([xshift=5mm,yshift=-0.62cm]frame.north west) {Question 3 C};
\node[anchor=east,text=white,font=\sffamily\itshape\large] at ([xshift=-5mm,yshift=-0.62cm]frame.north east) {Nombres entiers et décimaux};}]
\vspace*{1.75cm}
\begin{minipage}[t][6.9cm][c]{\linewidth}\LARGE\centering
\begin{minipage}{0.95\linewidth}\raggedright\linespread{1.25}\selectfont Quelle est l'écriture décimale de $\dfrac{63}{1000}$~?\end{minipage}
\end{minipage}

\vfill
\begin{tcolorbox}[enhanced,nobeforeafter,width=7.1cm,height=2.1cm,valign=center,colback=cnum!6,colframe=cnum,arc=2.5mm,boxrule=0.9pt,left=3mm]{\Large\bfseries\color{cnum}a)}\quad\begin{minipage}[c]{5.6cm}\raggedright\Large $0{,}063$\end{minipage}\end{tcolorbox}\hfill\begin{tcolorbox}[enhanced,nobeforeafter,width=7.1cm,height=2.1cm,valign=center,colback=cnum!6,colframe=cnum,arc=2.5mm,boxrule=0.9pt,left=3mm]{\Large\bfseries\color{cnum}b)}\quad\begin{minipage}[c]{5.6cm}\raggedright\Large $6{,}3$\end{minipage}\end{tcolorbox}\par\vspace{2.5mm}
\begin{tcolorbox}[enhanced,nobeforeafter,width=7.1cm,height=2.1cm,valign=center,colback=cnum!6,colframe=cnum,arc=2.5mm,boxrule=0.9pt,left=3mm]{\Large\bfseries\color{cnum}c)}\quad\begin{minipage}[c]{5.6cm}\raggedright\Large $0{,}0063$\end{minipage}\end{tcolorbox}\hfill\begin{tcolorbox}[enhanced,nobeforeafter,width=7.1cm,height=2.1cm,valign=center,colback=cnum!6,colframe=cnum,arc=2.5mm,boxrule=0.9pt,left=3mm]{\Large\bfseries\color{cnum}d)}\quad\begin{minipage}[c]{5.6cm}\raggedright\Large $0{,}63$\end{minipage}\end{tcolorbox}
\end{tcolorbox}
\newpage
\begin{tcolorbox}[enhanced,width=16cm,height=14cm,colback=white,colframe=cnum,boxrule=2pt,arc=0mm,
left=6mm,right=6mm,top=0mm,bottom=4mm,sharp corners,
overlay={\fill[cnum] (frame.north west) rectangle ([yshift=-1.25cm]frame.north east);
\node[anchor=west,text=white,font=\sffamily\bfseries\LARGE] at ([xshift=5mm,yshift=-0.62cm]frame.north west) {Question 4 C};
\node[anchor=east,text=white,font=\sffamily\itshape\large] at ([xshift=-5mm,yshift=-0.62cm]frame.north east) {Nombres entiers et décimaux};}]
\vspace*{1.75cm}
\begin{minipage}[t][6.9cm][c]{\linewidth}\LARGE\centering
\begin{minipage}{0.95\linewidth}\raggedright\linespread{1.25}\selectfont Quelle est l'écriture décimale de $\dfrac{9}{10}$~?\end{minipage}
\end{minipage}

\vfill
\begin{tcolorbox}[enhanced,nobeforeafter,width=7.1cm,height=2.1cm,valign=center,colback=cnum!6,colframe=cnum,arc=2.5mm,boxrule=0.9pt,left=3mm]{\Large\bfseries\color{cnum}a)}\quad\begin{minipage}[c]{5.6cm}\raggedright\Large $0{,}9$\end{minipage}\end{tcolorbox}\hfill\begin{tcolorbox}[enhanced,nobeforeafter,width=7.1cm,height=2.1cm,valign=center,colback=cnum!6,colframe=cnum,arc=2.5mm,boxrule=0.9pt,left=3mm]{\Large\bfseries\color{cnum}b)}\quad\begin{minipage}[c]{5.6cm}\raggedright\Large $0{,}09$\end{minipage}\end{tcolorbox}\par\vspace{2.5mm}
\begin{tcolorbox}[enhanced,nobeforeafter,width=7.1cm,height=2.1cm,valign=center,colback=cnum!6,colframe=cnum,arc=2.5mm,boxrule=0.9pt,left=3mm]{\Large\bfseries\color{cnum}c)}\quad\begin{minipage}[c]{5.6cm}\raggedright\Large $90$\end{minipage}\end{tcolorbox}\hfill\begin{tcolorbox}[enhanced,nobeforeafter,width=7.1cm,height=2.1cm,valign=center,colback=cnum!6,colframe=cnum,arc=2.5mm,boxrule=0.9pt,left=3mm]{\Large\bfseries\color{cnum}d)}\quad\begin{minipage}[c]{5.6cm}\raggedright\Large $9$\end{minipage}\end{tcolorbox}
\end{tcolorbox}
\newpage
\begin{tcolorbox}[enhanced,width=16cm,height=14cm,colback=white,colframe=cnum,boxrule=2pt,arc=0mm,
left=6mm,right=6mm,top=0mm,bottom=4mm,sharp corners,
overlay={\fill[cnum] (frame.north west) rectangle ([yshift=-1.25cm]frame.north east);
\node[anchor=west,text=white,font=\sffamily\bfseries\LARGE] at ([xshift=5mm,yshift=-0.62cm]frame.north west) {Question 5 C};
\node[anchor=east,text=white,font=\sffamily\itshape\large] at ([xshift=-5mm,yshift=-0.62cm]frame.north east) {Nombres entiers et décimaux};}]
\vspace*{1.75cm}
\begin{minipage}[t][6.9cm][c]{\linewidth}\LARGE\centering
\begin{minipage}{0.95\linewidth}\raggedright\linespread{1.25}\selectfont Quel est le plus grand de ces quatre nombres~: $2{,}25\ ;\ 2{,}5\ ;\ 2{,}205\ ;\ 2{,}052$~?\end{minipage}
\end{minipage}

\vfill
\begin{tcolorbox}[enhanced,nobeforeafter,width=7.1cm,height=2.1cm,valign=center,colback=cnum!6,colframe=cnum,arc=2.5mm,boxrule=0.9pt,left=3mm]{\Large\bfseries\color{cnum}a)}\quad\begin{minipage}[c]{5.6cm}\raggedright\Large $2{,}052$\end{minipage}\end{tcolorbox}\hfill\begin{tcolorbox}[enhanced,nobeforeafter,width=7.1cm,height=2.1cm,valign=center,colback=cnum!6,colframe=cnum,arc=2.5mm,boxrule=0.9pt,left=3mm]{\Large\bfseries\color{cnum}b)}\quad\begin{minipage}[c]{5.6cm}\raggedright\Large $2{,}5$\end{minipage}\end{tcolorbox}\par\vspace{2.5mm}
\begin{tcolorbox}[enhanced,nobeforeafter,width=7.1cm,height=2.1cm,valign=center,colback=cnum!6,colframe=cnum,arc=2.5mm,boxrule=0.9pt,left=3mm]{\Large\bfseries\color{cnum}c)}\quad\begin{minipage}[c]{5.6cm}\raggedright\Large $2{,}205$\end{minipage}\end{tcolorbox}\hfill\begin{tcolorbox}[enhanced,nobeforeafter,width=7.1cm,height=2.1cm,valign=center,colback=cnum!6,colframe=cnum,arc=2.5mm,boxrule=0.9pt,left=3mm]{\Large\bfseries\color{cnum}d)}\quad\begin{minipage}[c]{5.6cm}\raggedright\Large $2{,}25$\end{minipage}\end{tcolorbox}
\end{tcolorbox}
\newpage
\begin{tcolorbox}[enhanced,width=16cm,height=14cm,colback=white,colframe=cnum,boxrule=2pt,arc=0mm,
left=6mm,right=6mm,top=0mm,bottom=4mm,sharp corners,
overlay={\fill[cnum] (frame.north west) rectangle ([yshift=-1.25cm]frame.north east);
\node[anchor=west,text=white,font=\sffamily\bfseries\LARGE] at ([xshift=5mm,yshift=-0.62cm]frame.north west) {Question 6 C};
\node[anchor=east,text=white,font=\sffamily\itshape\large] at ([xshift=-5mm,yshift=-0.62cm]frame.north east) {Nombres entiers et décimaux};}]
\vspace*{1.75cm}
\begin{minipage}[t][6.9cm][c]{\linewidth}\LARGE\centering
\begin{minipage}{0.95\linewidth}\raggedright\linespread{1.25}\selectfont Quel est le plus petit de ces quatre nombres~: $8{,}9\ ;\ 8{,}89\ ;\ 8{,}099\ ;\ 8{,}91$~?\end{minipage}
\end{minipage}

\vfill
\begin{tcolorbox}[enhanced,nobeforeafter,width=7.1cm,height=2.1cm,valign=center,colback=cnum!6,colframe=cnum,arc=2.5mm,boxrule=0.9pt,left=3mm]{\Large\bfseries\color{cnum}a)}\quad\begin{minipage}[c]{5.6cm}\raggedright\Large $8{,}9$\end{minipage}\end{tcolorbox}\hfill\begin{tcolorbox}[enhanced,nobeforeafter,width=7.1cm,height=2.1cm,valign=center,colback=cnum!6,colframe=cnum,arc=2.5mm,boxrule=0.9pt,left=3mm]{\Large\bfseries\color{cnum}b)}\quad\begin{minipage}[c]{5.6cm}\raggedright\Large $8{,}91$\end{minipage}\end{tcolorbox}\par\vspace{2.5mm}
\begin{tcolorbox}[enhanced,nobeforeafter,width=7.1cm,height=2.1cm,valign=center,colback=cnum!6,colframe=cnum,arc=2.5mm,boxrule=0.9pt,left=3mm]{\Large\bfseries\color{cnum}c)}\quad\begin{minipage}[c]{5.6cm}\raggedright\Large $8{,}099$\end{minipage}\end{tcolorbox}\hfill\begin{tcolorbox}[enhanced,nobeforeafter,width=7.1cm,height=2.1cm,valign=center,colback=cnum!6,colframe=cnum,arc=2.5mm,boxrule=0.9pt,left=3mm]{\Large\bfseries\color{cnum}d)}\quad\begin{minipage}[c]{5.6cm}\raggedright\Large $8{,}89$\end{minipage}\end{tcolorbox}
\end{tcolorbox}
\newpage
\begin{tcolorbox}[enhanced,width=16cm,height=14cm,colback=white,colframe=cnum,boxrule=2pt,arc=0mm,
left=6mm,right=6mm,top=0mm,bottom=4mm,sharp corners,
overlay={\fill[cnum] (frame.north west) rectangle ([yshift=-1.25cm]frame.north east);
\node[anchor=west,text=white,font=\sffamily\bfseries\LARGE] at ([xshift=5mm,yshift=-0.62cm]frame.north west) {Question 7 C};
\node[anchor=east,text=white,font=\sffamily\itshape\large] at ([xshift=-5mm,yshift=-0.62cm]frame.north east) {Nombres entiers et décimaux};}]
\vspace*{1.75cm}
\begin{minipage}[t][6.9cm][c]{\linewidth}\LARGE\centering
\begin{minipage}{0.95\linewidth}\raggedright\linespread{1.25}\selectfont Quel est l'arrondi au dixième de $5{,}05$~?\end{minipage}
\end{minipage}

\vfill
\begin{tcolorbox}[enhanced,nobeforeafter,width=7.1cm,height=2.1cm,valign=center,colback=cnum!6,colframe=cnum,arc=2.5mm,boxrule=0.9pt,left=3mm]{\Large\bfseries\color{cnum}a)}\quad\begin{minipage}[c]{5.6cm}\raggedright\Large $5{,}1$\end{minipage}\end{tcolorbox}\hfill\begin{tcolorbox}[enhanced,nobeforeafter,width=7.1cm,height=2.1cm,valign=center,colback=cnum!6,colframe=cnum,arc=2.5mm,boxrule=0.9pt,left=3mm]{\Large\bfseries\color{cnum}b)}\quad\begin{minipage}[c]{5.6cm}\raggedright\Large $5{,}3$\end{minipage}\end{tcolorbox}\par\vspace{2.5mm}
\begin{tcolorbox}[enhanced,nobeforeafter,width=7.1cm,height=2.1cm,valign=center,colback=cnum!6,colframe=cnum,arc=2.5mm,boxrule=0.9pt,left=3mm]{\Large\bfseries\color{cnum}c)}\quad\begin{minipage}[c]{5.6cm}\raggedright\Large $5{,}05$\end{minipage}\end{tcolorbox}\hfill\begin{tcolorbox}[enhanced,nobeforeafter,width=7.1cm,height=2.1cm,valign=center,colback=cnum!6,colframe=cnum,arc=2.5mm,boxrule=0.9pt,left=3mm]{\Large\bfseries\color{cnum}d)}\quad\begin{minipage}[c]{5.6cm}\raggedright\Large $5$\end{minipage}\end{tcolorbox}
\end{tcolorbox}
\newpage
\begin{tcolorbox}[enhanced,width=16cm,height=14cm,colback=white,colframe=cnum,boxrule=2pt,arc=0mm,
left=6mm,right=6mm,top=0mm,bottom=4mm,sharp corners,
overlay={\fill[cnum] (frame.north west) rectangle ([yshift=-1.25cm]frame.north east);
\node[anchor=west,text=white,font=\sffamily\bfseries\LARGE] at ([xshift=5mm,yshift=-0.62cm]frame.north west) {Question 8 C};
\node[anchor=east,text=white,font=\sffamily\itshape\large] at ([xshift=-5mm,yshift=-0.62cm]frame.north east) {Nombres entiers et décimaux};}]
\vspace*{1.75cm}
\begin{minipage}[t][6.9cm][c]{\linewidth}\LARGE\centering
\begin{minipage}{0.95\linewidth}\raggedright\linespread{1.25}\selectfont Quel est l'arrondi au dixième de $9{,}98$~?\end{minipage}
\end{minipage}

\vfill
\begin{tcolorbox}[enhanced,nobeforeafter,width=7.1cm,height=2.1cm,valign=center,colback=cnum!6,colframe=cnum,arc=2.5mm,boxrule=0.9pt,left=3mm]{\Large\bfseries\color{cnum}a)}\quad\begin{minipage}[c]{5.6cm}\raggedright\Large $9{,}98$\end{minipage}\end{tcolorbox}\hfill\begin{tcolorbox}[enhanced,nobeforeafter,width=7.1cm,height=2.1cm,valign=center,colback=cnum!6,colframe=cnum,arc=2.5mm,boxrule=0.9pt,left=3mm]{\Large\bfseries\color{cnum}b)}\quad\begin{minipage}[c]{5.6cm}\raggedright\Large $9{,}9$\end{minipage}\end{tcolorbox}\par\vspace{2.5mm}
\begin{tcolorbox}[enhanced,nobeforeafter,width=7.1cm,height=2.1cm,valign=center,colback=cnum!6,colframe=cnum,arc=2.5mm,boxrule=0.9pt,left=3mm]{\Large\bfseries\color{cnum}c)}\quad\begin{minipage}[c]{5.6cm}\raggedright\Large $10$\end{minipage}\end{tcolorbox}\hfill\begin{tcolorbox}[enhanced,nobeforeafter,width=7.1cm,height=2.1cm,valign=center,colback=cnum!6,colframe=cnum,arc=2.5mm,boxrule=0.9pt,left=3mm]{\Large\bfseries\color{cnum}d)}\quad\begin{minipage}[c]{5.6cm}\raggedright\Large $10{,}2$\end{minipage}\end{tcolorbox}
\end{tcolorbox}
\newpage
\begin{tcolorbox}[enhanced,width=16cm,height=14cm,colback=white,colframe=cnum,boxrule=2pt,arc=0mm,
left=6mm,right=6mm,top=0mm,bottom=4mm,sharp corners,
overlay={\fill[cnum] (frame.north west) rectangle ([yshift=-1.25cm]frame.north east);
\node[anchor=west,text=white,font=\sffamily\bfseries\LARGE] at ([xshift=5mm,yshift=-0.62cm]frame.north west) {Question 9 C};
\node[anchor=east,text=white,font=\sffamily\itshape\large] at ([xshift=-5mm,yshift=-0.62cm]frame.north east) {Nombres entiers et décimaux};}]
\vspace*{1.75cm}
\begin{minipage}[t][6.9cm][c]{\linewidth}\LARGE\centering
\begin{minipage}{0.95\linewidth}\raggedright\linespread{1.25}\selectfont Que vaut $12{,}5\times10$~?\end{minipage}
\end{minipage}

\vfill
\begin{tcolorbox}[enhanced,nobeforeafter,width=7.1cm,height=2.1cm,valign=center,colback=cnum!6,colframe=cnum,arc=2.5mm,boxrule=0.9pt,left=3mm]{\Large\bfseries\color{cnum}a)}\quad\begin{minipage}[c]{5.6cm}\raggedright\Large $125$\end{minipage}\end{tcolorbox}\hfill\begin{tcolorbox}[enhanced,nobeforeafter,width=7.1cm,height=2.1cm,valign=center,colback=cnum!6,colframe=cnum,arc=2.5mm,boxrule=0.9pt,left=3mm]{\Large\bfseries\color{cnum}b)}\quad\begin{minipage}[c]{5.6cm}\raggedright\Large $12{,}5$\end{minipage}\end{tcolorbox}\par\vspace{2.5mm}
\begin{tcolorbox}[enhanced,nobeforeafter,width=7.1cm,height=2.1cm,valign=center,colback=cnum!6,colframe=cnum,arc=2.5mm,boxrule=0.9pt,left=3mm]{\Large\bfseries\color{cnum}c)}\quad\begin{minipage}[c]{5.6cm}\raggedright\Large $1\,250$\end{minipage}\end{tcolorbox}\hfill\begin{tcolorbox}[enhanced,nobeforeafter,width=7.1cm,height=2.1cm,valign=center,colback=cnum!6,colframe=cnum,arc=2.5mm,boxrule=0.9pt,left=3mm]{\Large\bfseries\color{cnum}d)}\quad\begin{minipage}[c]{5.6cm}\raggedright\Large $1{,}25$\end{minipage}\end{tcolorbox}
\end{tcolorbox}
\newpage
\begin{tcolorbox}[enhanced,width=16cm,height=14cm,colback=white,colframe=cnum,boxrule=2pt,arc=0mm,
left=6mm,right=6mm,top=0mm,bottom=4mm,sharp corners,
overlay={\fill[cnum] (frame.north west) rectangle ([yshift=-1.25cm]frame.north east);
\node[anchor=west,text=white,font=\sffamily\bfseries\LARGE] at ([xshift=5mm,yshift=-0.62cm]frame.north west) {Question 10 C};
\node[anchor=east,text=white,font=\sffamily\itshape\large] at ([xshift=-5mm,yshift=-0.62cm]frame.north east) {Nombres entiers et décimaux};}]
\vspace*{1.75cm}
\begin{minipage}[t][6.9cm][c]{\linewidth}\LARGE\centering
\begin{minipage}{0.95\linewidth}\raggedright\linespread{1.25}\selectfont Que vaut $0{,}4\times1\,000$~?\end{minipage}
\end{minipage}

\vfill
\begin{tcolorbox}[enhanced,nobeforeafter,width=7.1cm,height=2.1cm,valign=center,colback=cnum!6,colframe=cnum,arc=2.5mm,boxrule=0.9pt,left=3mm]{\Large\bfseries\color{cnum}a)}\quad\begin{minipage}[c]{5.6cm}\raggedright\Large $40\,000$\end{minipage}\end{tcolorbox}\hfill\begin{tcolorbox}[enhanced,nobeforeafter,width=7.1cm,height=2.1cm,valign=center,colback=cnum!6,colframe=cnum,arc=2.5mm,boxrule=0.9pt,left=3mm]{\Large\bfseries\color{cnum}b)}\quad\begin{minipage}[c]{5.6cm}\raggedright\Large $400$\end{minipage}\end{tcolorbox}\par\vspace{2.5mm}
\begin{tcolorbox}[enhanced,nobeforeafter,width=7.1cm,height=2.1cm,valign=center,colback=cnum!6,colframe=cnum,arc=2.5mm,boxrule=0.9pt,left=3mm]{\Large\bfseries\color{cnum}c)}\quad\begin{minipage}[c]{5.6cm}\raggedright\Large $40$\end{minipage}\end{tcolorbox}\hfill\begin{tcolorbox}[enhanced,nobeforeafter,width=7.1cm,height=2.1cm,valign=center,colback=cnum!6,colframe=cnum,arc=2.5mm,boxrule=0.9pt,left=3mm]{\Large\bfseries\color{cnum}d)}\quad\begin{minipage}[c]{5.6cm}\raggedright\Large $4\,000$\end{minipage}\end{tcolorbox}
\end{tcolorbox}
\newpage
\begin{tcolorbox}[enhanced,width=16cm,height=14cm,colback=white,colframe=cnum,boxrule=2pt,arc=0mm,
left=6mm,right=6mm,top=0mm,bottom=4mm,sharp corners,
overlay={\fill[cnum] (frame.north west) rectangle ([yshift=-1.25cm]frame.north east);
\node[anchor=west,text=white,font=\sffamily\bfseries\LARGE] at ([xshift=5mm,yshift=-0.62cm]frame.north west) {Question 11 C};
\node[anchor=east,text=white,font=\sffamily\itshape\large] at ([xshift=-5mm,yshift=-0.62cm]frame.north east) {Nombres entiers et décimaux};}]
\vspace*{1.75cm}
\begin{minipage}[t][6.9cm][c]{\linewidth}\LARGE\centering
\begin{minipage}{0.95\linewidth}\raggedright\linespread{1.25}\selectfont Que vaut $12{,}8\div10$~?\end{minipage}
\end{minipage}

\vfill
\begin{tcolorbox}[enhanced,nobeforeafter,width=7.1cm,height=2.1cm,valign=center,colback=cnum!6,colframe=cnum,arc=2.5mm,boxrule=0.9pt,left=3mm]{\Large\bfseries\color{cnum}a)}\quad\begin{minipage}[c]{5.6cm}\raggedright\Large $12{,}8$\end{minipage}\end{tcolorbox}\hfill\begin{tcolorbox}[enhanced,nobeforeafter,width=7.1cm,height=2.1cm,valign=center,colback=cnum!6,colframe=cnum,arc=2.5mm,boxrule=0.9pt,left=3mm]{\Large\bfseries\color{cnum}b)}\quad\begin{minipage}[c]{5.6cm}\raggedright\Large $128$\end{minipage}\end{tcolorbox}\par\vspace{2.5mm}
\begin{tcolorbox}[enhanced,nobeforeafter,width=7.1cm,height=2.1cm,valign=center,colback=cnum!6,colframe=cnum,arc=2.5mm,boxrule=0.9pt,left=3mm]{\Large\bfseries\color{cnum}c)}\quad\begin{minipage}[c]{5.6cm}\raggedright\Large $0{,}128$\end{minipage}\end{tcolorbox}\hfill\begin{tcolorbox}[enhanced,nobeforeafter,width=7.1cm,height=2.1cm,valign=center,colback=cnum!6,colframe=cnum,arc=2.5mm,boxrule=0.9pt,left=3mm]{\Large\bfseries\color{cnum}d)}\quad\begin{minipage}[c]{5.6cm}\raggedright\Large $1{,}28$\end{minipage}\end{tcolorbox}
\end{tcolorbox}
\newpage
\begin{tcolorbox}[enhanced,width=16cm,height=14cm,colback=white,colframe=cnum,boxrule=2pt,arc=0mm,
left=6mm,right=6mm,top=0mm,bottom=4mm,sharp corners,
overlay={\fill[cnum] (frame.north west) rectangle ([yshift=-1.25cm]frame.north east);
\node[anchor=west,text=white,font=\sffamily\bfseries\LARGE] at ([xshift=5mm,yshift=-0.62cm]frame.north west) {Question 12 C};
\node[anchor=east,text=white,font=\sffamily\itshape\large] at ([xshift=-5mm,yshift=-0.62cm]frame.north east) {Nombres entiers et décimaux};}]
\vspace*{1.75cm}
\begin{minipage}[t][6.9cm][c]{\linewidth}\LARGE\centering
\begin{minipage}{0.95\linewidth}\raggedright\linespread{1.25}\selectfont Que vaut $905\div1\,000$~?\end{minipage}
\end{minipage}

\vfill
\begin{tcolorbox}[enhanced,nobeforeafter,width=7.1cm,height=2.1cm,valign=center,colback=cnum!6,colframe=cnum,arc=2.5mm,boxrule=0.9pt,left=3mm]{\Large\bfseries\color{cnum}a)}\quad\begin{minipage}[c]{5.6cm}\raggedright\Large $0{,}905$\end{minipage}\end{tcolorbox}\hfill\begin{tcolorbox}[enhanced,nobeforeafter,width=7.1cm,height=2.1cm,valign=center,colback=cnum!6,colframe=cnum,arc=2.5mm,boxrule=0.9pt,left=3mm]{\Large\bfseries\color{cnum}b)}\quad\begin{minipage}[c]{5.6cm}\raggedright\Large $905\,000$\end{minipage}\end{tcolorbox}\par\vspace{2.5mm}
\begin{tcolorbox}[enhanced,nobeforeafter,width=7.1cm,height=2.1cm,valign=center,colback=cnum!6,colframe=cnum,arc=2.5mm,boxrule=0.9pt,left=3mm]{\Large\bfseries\color{cnum}c)}\quad\begin{minipage}[c]{5.6cm}\raggedright\Large $0{,}0905$\end{minipage}\end{tcolorbox}\hfill\begin{tcolorbox}[enhanced,nobeforeafter,width=7.1cm,height=2.1cm,valign=center,colback=cnum!6,colframe=cnum,arc=2.5mm,boxrule=0.9pt,left=3mm]{\Large\bfseries\color{cnum}d)}\quad\begin{minipage}[c]{5.6cm}\raggedright\Large $9{,}05$\end{minipage}\end{tcolorbox}
\end{tcolorbox}
\newpage
\begin{tcolorbox}[enhanced,width=16cm,height=14cm,colback=white,colframe=cnum,boxrule=2pt,arc=0mm,
left=6mm,right=6mm,top=0mm,bottom=4mm,sharp corners,
overlay={\fill[cnum] (frame.north west) rectangle ([yshift=-1.25cm]frame.north east);
\node[anchor=west,text=white,font=\sffamily\bfseries\LARGE] at ([xshift=5mm,yshift=-0.62cm]frame.north west) {Question 13 C};
\node[anchor=east,text=white,font=\sffamily\itshape\large] at ([xshift=-5mm,yshift=-0.62cm]frame.north east) {Nombres entiers et décimaux};}]
\vspace*{1.75cm}
\begin{minipage}[t][6.9cm][c]{\linewidth}\LARGE\centering
\begin{minipage}{0.95\linewidth}\raggedright\linespread{1.25}\selectfont Comment s'écrit en chiffres le nombre «~quatre-vingt-dix mille trois cent un~»~?\end{minipage}
\end{minipage}

\vfill
\begin{tcolorbox}[enhanced,nobeforeafter,width=7.1cm,height=2.1cm,valign=center,colback=cnum!6,colframe=cnum,arc=2.5mm,boxrule=0.9pt,left=3mm]{\Large\bfseries\color{cnum}a)}\quad\begin{minipage}[c]{5.6cm}\raggedright\Large $900\,301$\end{minipage}\end{tcolorbox}\hfill\begin{tcolorbox}[enhanced,nobeforeafter,width=7.1cm,height=2.1cm,valign=center,colback=cnum!6,colframe=cnum,arc=2.5mm,boxrule=0.9pt,left=3mm]{\Large\bfseries\color{cnum}b)}\quad\begin{minipage}[c]{5.6cm}\raggedright\Large $90\,031$\end{minipage}\end{tcolorbox}\par\vspace{2.5mm}
\begin{tcolorbox}[enhanced,nobeforeafter,width=7.1cm,height=2.1cm,valign=center,colback=cnum!6,colframe=cnum,arc=2.5mm,boxrule=0.9pt,left=3mm]{\Large\bfseries\color{cnum}c)}\quad\begin{minipage}[c]{5.6cm}\raggedright\Large $9\,301$\end{minipage}\end{tcolorbox}\hfill\begin{tcolorbox}[enhanced,nobeforeafter,width=7.1cm,height=2.1cm,valign=center,colback=cnum!6,colframe=cnum,arc=2.5mm,boxrule=0.9pt,left=3mm]{\Large\bfseries\color{cnum}d)}\quad\begin{minipage}[c]{5.6cm}\raggedright\Large $90\,301$\end{minipage}\end{tcolorbox}
\end{tcolorbox}
\newpage
\begin{tcolorbox}[enhanced,width=16cm,height=14cm,colback=white,colframe=cnum,boxrule=2pt,arc=0mm,
left=6mm,right=6mm,top=0mm,bottom=4mm,sharp corners,
overlay={\fill[cnum] (frame.north west) rectangle ([yshift=-1.25cm]frame.north east);
\node[anchor=west,text=white,font=\sffamily\bfseries\LARGE] at ([xshift=5mm,yshift=-0.62cm]frame.north west) {Question 14 C};
\node[anchor=east,text=white,font=\sffamily\itshape\large] at ([xshift=-5mm,yshift=-0.62cm]frame.north east) {Nombres entiers et décimaux};}]
\vspace*{1.75cm}
\begin{minipage}[t][6.9cm][c]{\linewidth}\LARGE\centering
\begin{minipage}{0.95\linewidth}\raggedright\linespread{1.25}\selectfont Comment s'écrit en chiffres le nombre «~six cent mille soixante-dix~»~?\end{minipage}
\end{minipage}

\vfill
\begin{tcolorbox}[enhanced,nobeforeafter,width=7.1cm,height=2.1cm,valign=center,colback=cnum!6,colframe=cnum,arc=2.5mm,boxrule=0.9pt,left=3mm]{\Large\bfseries\color{cnum}a)}\quad\begin{minipage}[c]{5.6cm}\raggedright\Large $60\,070$\end{minipage}\end{tcolorbox}\hfill\begin{tcolorbox}[enhanced,nobeforeafter,width=7.1cm,height=2.1cm,valign=center,colback=cnum!6,colframe=cnum,arc=2.5mm,boxrule=0.9pt,left=3mm]{\Large\bfseries\color{cnum}b)}\quad\begin{minipage}[c]{5.6cm}\raggedright\Large $6\,000\,070$\end{minipage}\end{tcolorbox}\par\vspace{2.5mm}
\begin{tcolorbox}[enhanced,nobeforeafter,width=7.1cm,height=2.1cm,valign=center,colback=cnum!6,colframe=cnum,arc=2.5mm,boxrule=0.9pt,left=3mm]{\Large\bfseries\color{cnum}c)}\quad\begin{minipage}[c]{5.6cm}\raggedright\Large $600\,070$\end{minipage}\end{tcolorbox}\hfill\begin{tcolorbox}[enhanced,nobeforeafter,width=7.1cm,height=2.1cm,valign=center,colback=cnum!6,colframe=cnum,arc=2.5mm,boxrule=0.9pt,left=3mm]{\Large\bfseries\color{cnum}d)}\quad\begin{minipage}[c]{5.6cm}\raggedright\Large $600\,700$\end{minipage}\end{tcolorbox}
\end{tcolorbox}
\newpage
\begin{tcolorbox}[enhanced,width=16cm,height=14cm,colback=white,colframe=cnum,boxrule=2pt,arc=0mm,
left=6mm,right=6mm,top=0mm,bottom=4mm,sharp corners,
overlay={\fill[cnum] (frame.north west) rectangle ([yshift=-1.25cm]frame.north east);
\node[anchor=west,text=white,font=\sffamily\bfseries\LARGE] at ([xshift=5mm,yshift=-0.62cm]frame.north west) {Question 15 C};
\node[anchor=east,text=white,font=\sffamily\itshape\large] at ([xshift=-5mm,yshift=-0.62cm]frame.north east) {Nombres entiers et décimaux};}]
\vspace*{1.75cm}
\begin{minipage}[t][6.9cm][c]{\linewidth}\LARGE\centering
\begin{minipage}{0.95\linewidth}\raggedright\linespread{1.25}\selectfont Quel nombre est compris entre $3$ et $3{,}01$~?\end{minipage}
\end{minipage}

\vfill
\begin{tcolorbox}[enhanced,nobeforeafter,width=7.1cm,height=2.1cm,valign=center,colback=cnum!6,colframe=cnum,arc=2.5mm,boxrule=0.9pt,left=3mm]{\Large\bfseries\color{cnum}a)}\quad\begin{minipage}[c]{5.6cm}\raggedright\Large $2{,}99$\end{minipage}\end{tcolorbox}\hfill\begin{tcolorbox}[enhanced,nobeforeafter,width=7.1cm,height=2.1cm,valign=center,colback=cnum!6,colframe=cnum,arc=2.5mm,boxrule=0.9pt,left=3mm]{\Large\bfseries\color{cnum}b)}\quad\begin{minipage}[c]{5.6cm}\raggedright\Large $3{,}005$\end{minipage}\end{tcolorbox}\par\vspace{2.5mm}
\begin{tcolorbox}[enhanced,nobeforeafter,width=7.1cm,height=2.1cm,valign=center,colback=cnum!6,colframe=cnum,arc=2.5mm,boxrule=0.9pt,left=3mm]{\Large\bfseries\color{cnum}c)}\quad\begin{minipage}[c]{5.6cm}\raggedright\Large $3{,}015$\end{minipage}\end{tcolorbox}\hfill\begin{tcolorbox}[enhanced,nobeforeafter,width=7.1cm,height=2.1cm,valign=center,colback=cnum!6,colframe=cnum,arc=2.5mm,boxrule=0.9pt,left=3mm]{\Large\bfseries\color{cnum}d)}\quad\begin{minipage}[c]{5.6cm}\raggedright\Large $3{,}1$\end{minipage}\end{tcolorbox}
\end{tcolorbox}
\newpage
\begin{tcolorbox}[enhanced,width=16cm,height=14cm,colback=white,colframe=cnum,boxrule=2pt,arc=0mm,
left=6mm,right=6mm,top=0mm,bottom=4mm,sharp corners,
overlay={\fill[cnum] (frame.north west) rectangle ([yshift=-1.25cm]frame.north east);
\node[anchor=west,text=white,font=\sffamily\bfseries\LARGE] at ([xshift=5mm,yshift=-0.62cm]frame.north west) {Question 16 C};
\node[anchor=east,text=white,font=\sffamily\itshape\large] at ([xshift=-5mm,yshift=-0.62cm]frame.north east) {Nombres entiers et décimaux};}]
\vspace*{1.75cm}
\begin{minipage}[t][6.9cm][c]{\linewidth}\LARGE\centering
\begin{minipage}{0.95\linewidth}\raggedright\linespread{1.25}\selectfont Quel nombre est compris entre $9{,}9$ et $10$~?\end{minipage}
\end{minipage}

\vfill
\begin{tcolorbox}[enhanced,nobeforeafter,width=7.1cm,height=2.1cm,valign=center,colback=cnum!6,colframe=cnum,arc=2.5mm,boxrule=0.9pt,left=3mm]{\Large\bfseries\color{cnum}a)}\quad\begin{minipage}[c]{5.6cm}\raggedright\Large $9{,}85$\end{minipage}\end{tcolorbox}\hfill\begin{tcolorbox}[enhanced,nobeforeafter,width=7.1cm,height=2.1cm,valign=center,colback=cnum!6,colframe=cnum,arc=2.5mm,boxrule=0.9pt,left=3mm]{\Large\bfseries\color{cnum}b)}\quad\begin{minipage}[c]{5.6cm}\raggedright\Large $9{,}95$\end{minipage}\end{tcolorbox}\par\vspace{2.5mm}
\begin{tcolorbox}[enhanced,nobeforeafter,width=7.1cm,height=2.1cm,valign=center,colback=cnum!6,colframe=cnum,arc=2.5mm,boxrule=0.9pt,left=3mm]{\Large\bfseries\color{cnum}c)}\quad\begin{minipage}[c]{5.6cm}\raggedright\Large $10{,}5$\end{minipage}\end{tcolorbox}\hfill\begin{tcolorbox}[enhanced,nobeforeafter,width=7.1cm,height=2.1cm,valign=center,colback=cnum!6,colframe=cnum,arc=2.5mm,boxrule=0.9pt,left=3mm]{\Large\bfseries\color{cnum}d)}\quad\begin{minipage}[c]{5.6cm}\raggedright\Large $9{,}09$\end{minipage}\end{tcolorbox}
\end{tcolorbox}
\newpage
\begin{tcolorbox}[enhanced,width=16cm,height=14cm,colback=white,colframe=ccal,boxrule=2pt,arc=0mm,
left=6mm,right=6mm,top=0mm,bottom=4mm,sharp corners,
overlay={\fill[ccal] (frame.north west) rectangle ([yshift=-1.25cm]frame.north east);
\node[anchor=west,text=white,font=\sffamily\bfseries\LARGE] at ([xshift=5mm,yshift=-0.62cm]frame.north west) {Question 17 C};
\node[anchor=east,text=white,font=\sffamily\itshape\large] at ([xshift=-5mm,yshift=-0.62cm]frame.north east) {Calculs};}]
\vspace*{1.75cm}
\begin{minipage}[t][6.9cm][c]{\linewidth}\LARGE\centering
\begin{minipage}{0.95\linewidth}\raggedright\linespread{1.25}\selectfont Que vaut $8{,}6+9{,}48$~?\end{minipage}
\end{minipage}

\vfill
\begin{tcolorbox}[enhanced,nobeforeafter,width=7.1cm,height=2.1cm,valign=center,colback=ccal!6,colframe=ccal,arc=2.5mm,boxrule=0.9pt,left=3mm]{\Large\bfseries\color{ccal}a)}\quad\begin{minipage}[c]{5.6cm}\raggedright\Large $17{,}08$\end{minipage}\end{tcolorbox}\hfill\begin{tcolorbox}[enhanced,nobeforeafter,width=7.1cm,height=2.1cm,valign=center,colback=ccal!6,colframe=ccal,arc=2.5mm,boxrule=0.9pt,left=3mm]{\Large\bfseries\color{ccal}b)}\quad\begin{minipage}[c]{5.6cm}\raggedright\Large $18{,}08$\end{minipage}\end{tcolorbox}\par\vspace{2.5mm}
\begin{tcolorbox}[enhanced,nobeforeafter,width=7.1cm,height=2.1cm,valign=center,colback=ccal!6,colframe=ccal,arc=2.5mm,boxrule=0.9pt,left=3mm]{\Large\bfseries\color{ccal}c)}\quad\begin{minipage}[c]{5.6cm}\raggedright\Large $19{,}08$\end{minipage}\end{tcolorbox}\hfill\begin{tcolorbox}[enhanced,nobeforeafter,width=7.1cm,height=2.1cm,valign=center,colback=ccal!6,colframe=ccal,arc=2.5mm,boxrule=0.9pt,left=3mm]{\Large\bfseries\color{ccal}d)}\quad\begin{minipage}[c]{5.6cm}\raggedright\Large $18{,}18$\end{minipage}\end{tcolorbox}
\end{tcolorbox}
\newpage
\begin{tcolorbox}[enhanced,width=16cm,height=14cm,colback=white,colframe=ccal,boxrule=2pt,arc=0mm,
left=6mm,right=6mm,top=0mm,bottom=4mm,sharp corners,
overlay={\fill[ccal] (frame.north west) rectangle ([yshift=-1.25cm]frame.north east);
\node[anchor=west,text=white,font=\sffamily\bfseries\LARGE] at ([xshift=5mm,yshift=-0.62cm]frame.north west) {Question 18 C};
\node[anchor=east,text=white,font=\sffamily\itshape\large] at ([xshift=-5mm,yshift=-0.62cm]frame.north east) {Calculs};}]
\vspace*{1.75cm}
\begin{minipage}[t][6.9cm][c]{\linewidth}\LARGE\centering
\begin{minipage}{0.95\linewidth}\raggedright\linespread{1.25}\selectfont Que vaut $100{,}5+0{,}75$~?\end{minipage}
\end{minipage}

\vfill
\begin{tcolorbox}[enhanced,nobeforeafter,width=7.1cm,height=2.1cm,valign=center,colback=ccal!6,colframe=ccal,arc=2.5mm,boxrule=0.9pt,left=3mm]{\Large\bfseries\color{ccal}a)}\quad\begin{minipage}[c]{5.6cm}\raggedright\Large $101{,}25$\end{minipage}\end{tcolorbox}\hfill\begin{tcolorbox}[enhanced,nobeforeafter,width=7.1cm,height=2.1cm,valign=center,colback=ccal!6,colframe=ccal,arc=2.5mm,boxrule=0.9pt,left=3mm]{\Large\bfseries\color{ccal}b)}\quad\begin{minipage}[c]{5.6cm}\raggedright\Large $101{,}35$\end{minipage}\end{tcolorbox}\par\vspace{2.5mm}
\begin{tcolorbox}[enhanced,nobeforeafter,width=7.1cm,height=2.1cm,valign=center,colback=ccal!6,colframe=ccal,arc=2.5mm,boxrule=0.9pt,left=3mm]{\Large\bfseries\color{ccal}c)}\quad\begin{minipage}[c]{5.6cm}\raggedright\Large $100{,}25$\end{minipage}\end{tcolorbox}\hfill\begin{tcolorbox}[enhanced,nobeforeafter,width=7.1cm,height=2.1cm,valign=center,colback=ccal!6,colframe=ccal,arc=2.5mm,boxrule=0.9pt,left=3mm]{\Large\bfseries\color{ccal}d)}\quad\begin{minipage}[c]{5.6cm}\raggedright\Large $102{,}25$\end{minipage}\end{tcolorbox}
\end{tcolorbox}
\newpage
\begin{tcolorbox}[enhanced,width=16cm,height=14cm,colback=white,colframe=ccal,boxrule=2pt,arc=0mm,
left=6mm,right=6mm,top=0mm,bottom=4mm,sharp corners,
overlay={\fill[ccal] (frame.north west) rectangle ([yshift=-1.25cm]frame.north east);
\node[anchor=west,text=white,font=\sffamily\bfseries\LARGE] at ([xshift=5mm,yshift=-0.62cm]frame.north west) {Question 19 C};
\node[anchor=east,text=white,font=\sffamily\itshape\large] at ([xshift=-5mm,yshift=-0.62cm]frame.north east) {Calculs};}]
\vspace*{1.75cm}
\begin{minipage}[t][6.9cm][c]{\linewidth}\LARGE\centering
\begin{minipage}{0.95\linewidth}\raggedright\linespread{1.25}\selectfont Que vaut $9{,}3-2{,}47$~?\end{minipage}
\end{minipage}

\vfill
\begin{tcolorbox}[enhanced,nobeforeafter,width=7.1cm,height=2.1cm,valign=center,colback=ccal!6,colframe=ccal,arc=2.5mm,boxrule=0.9pt,left=3mm]{\Large\bfseries\color{ccal}a)}\quad\begin{minipage}[c]{5.6cm}\raggedright\Large $7{,}83$\end{minipage}\end{tcolorbox}\hfill\begin{tcolorbox}[enhanced,nobeforeafter,width=7.1cm,height=2.1cm,valign=center,colback=ccal!6,colframe=ccal,arc=2.5mm,boxrule=0.9pt,left=3mm]{\Large\bfseries\color{ccal}b)}\quad\begin{minipage}[c]{5.6cm}\raggedright\Large $6{,}83$\end{minipage}\end{tcolorbox}\par\vspace{2.5mm}
\begin{tcolorbox}[enhanced,nobeforeafter,width=7.1cm,height=2.1cm,valign=center,colback=ccal!6,colframe=ccal,arc=2.5mm,boxrule=0.9pt,left=3mm]{\Large\bfseries\color{ccal}c)}\quad\begin{minipage}[c]{5.6cm}\raggedright\Large $6{,}93$\end{minipage}\end{tcolorbox}\hfill\begin{tcolorbox}[enhanced,nobeforeafter,width=7.1cm,height=2.1cm,valign=center,colback=ccal!6,colframe=ccal,arc=2.5mm,boxrule=0.9pt,left=3mm]{\Large\bfseries\color{ccal}d)}\quad\begin{minipage}[c]{5.6cm}\raggedright\Large $6{,}73$\end{minipage}\end{tcolorbox}
\end{tcolorbox}
\newpage
\begin{tcolorbox}[enhanced,width=16cm,height=14cm,colback=white,colframe=ccal,boxrule=2pt,arc=0mm,
left=6mm,right=6mm,top=0mm,bottom=4mm,sharp corners,
overlay={\fill[ccal] (frame.north west) rectangle ([yshift=-1.25cm]frame.north east);
\node[anchor=west,text=white,font=\sffamily\bfseries\LARGE] at ([xshift=5mm,yshift=-0.62cm]frame.north west) {Question 20 C};
\node[anchor=east,text=white,font=\sffamily\itshape\large] at ([xshift=-5mm,yshift=-0.62cm]frame.north east) {Calculs};}]
\vspace*{1.75cm}
\begin{minipage}[t][6.9cm][c]{\linewidth}\LARGE\centering
\begin{minipage}{0.95\linewidth}\raggedright\linespread{1.25}\selectfont Que vaut $50{,}5-12{,}75$~?\end{minipage}
\end{minipage}

\vfill
\begin{tcolorbox}[enhanced,nobeforeafter,width=7.1cm,height=2.1cm,valign=center,colback=ccal!6,colframe=ccal,arc=2.5mm,boxrule=0.9pt,left=3mm]{\Large\bfseries\color{ccal}a)}\quad\begin{minipage}[c]{5.6cm}\raggedright\Large $37{,}65$\end{minipage}\end{tcolorbox}\hfill\begin{tcolorbox}[enhanced,nobeforeafter,width=7.1cm,height=2.1cm,valign=center,colback=ccal!6,colframe=ccal,arc=2.5mm,boxrule=0.9pt,left=3mm]{\Large\bfseries\color{ccal}b)}\quad\begin{minipage}[c]{5.6cm}\raggedright\Large $38{,}75$\end{minipage}\end{tcolorbox}\par\vspace{2.5mm}
\begin{tcolorbox}[enhanced,nobeforeafter,width=7.1cm,height=2.1cm,valign=center,colback=ccal!6,colframe=ccal,arc=2.5mm,boxrule=0.9pt,left=3mm]{\Large\bfseries\color{ccal}c)}\quad\begin{minipage}[c]{5.6cm}\raggedright\Large $37{,}75$\end{minipage}\end{tcolorbox}\hfill\begin{tcolorbox}[enhanced,nobeforeafter,width=7.1cm,height=2.1cm,valign=center,colback=ccal!6,colframe=ccal,arc=2.5mm,boxrule=0.9pt,left=3mm]{\Large\bfseries\color{ccal}d)}\quad\begin{minipage}[c]{5.6cm}\raggedright\Large $37{,}85$\end{minipage}\end{tcolorbox}
\end{tcolorbox}
\newpage
\begin{tcolorbox}[enhanced,width=16cm,height=14cm,colback=white,colframe=ccal,boxrule=2pt,arc=0mm,
left=6mm,right=6mm,top=0mm,bottom=4mm,sharp corners,
overlay={\fill[ccal] (frame.north west) rectangle ([yshift=-1.25cm]frame.north east);
\node[anchor=west,text=white,font=\sffamily\bfseries\LARGE] at ([xshift=5mm,yshift=-0.62cm]frame.north west) {Question 21 C};
\node[anchor=east,text=white,font=\sffamily\itshape\large] at ([xshift=-5mm,yshift=-0.62cm]frame.north east) {Calculs};}]
\vspace*{1.75cm}
\begin{minipage}[t][6.9cm][c]{\linewidth}\LARGE\centering
\begin{minipage}{0.95\linewidth}\raggedright\linespread{1.25}\selectfont Que vaut $0{,}25\times8$~?\end{minipage}
\end{minipage}

\vfill
\begin{tcolorbox}[enhanced,nobeforeafter,width=7.1cm,height=2.1cm,valign=center,colback=ccal!6,colframe=ccal,arc=2.5mm,boxrule=0.9pt,left=3mm]{\Large\bfseries\color{ccal}a)}\quad\begin{minipage}[c]{5.6cm}\raggedright\Large $0{,}2$\end{minipage}\end{tcolorbox}\hfill\begin{tcolorbox}[enhanced,nobeforeafter,width=7.1cm,height=2.1cm,valign=center,colback=ccal!6,colframe=ccal,arc=2.5mm,boxrule=0.9pt,left=3mm]{\Large\bfseries\color{ccal}b)}\quad\begin{minipage}[c]{5.6cm}\raggedright\Large $2$\end{minipage}\end{tcolorbox}\par\vspace{2.5mm}
\begin{tcolorbox}[enhanced,nobeforeafter,width=7.1cm,height=2.1cm,valign=center,colback=ccal!6,colframe=ccal,arc=2.5mm,boxrule=0.9pt,left=3mm]{\Large\bfseries\color{ccal}c)}\quad\begin{minipage}[c]{5.6cm}\raggedright\Large $8{,}25$\end{minipage}\end{tcolorbox}\hfill\begin{tcolorbox}[enhanced,nobeforeafter,width=7.1cm,height=2.1cm,valign=center,colback=ccal!6,colframe=ccal,arc=2.5mm,boxrule=0.9pt,left=3mm]{\Large\bfseries\color{ccal}d)}\quad\begin{minipage}[c]{5.6cm}\raggedright\Large $20$\end{minipage}\end{tcolorbox}
\end{tcolorbox}
\newpage
\begin{tcolorbox}[enhanced,width=16cm,height=14cm,colback=white,colframe=ccal,boxrule=2pt,arc=0mm,
left=6mm,right=6mm,top=0mm,bottom=4mm,sharp corners,
overlay={\fill[ccal] (frame.north west) rectangle ([yshift=-1.25cm]frame.north east);
\node[anchor=west,text=white,font=\sffamily\bfseries\LARGE] at ([xshift=5mm,yshift=-0.62cm]frame.north west) {Question 22 C};
\node[anchor=east,text=white,font=\sffamily\itshape\large] at ([xshift=-5mm,yshift=-0.62cm]frame.north east) {Calculs};}]
\vspace*{1.75cm}
\begin{minipage}[t][6.9cm][c]{\linewidth}\LARGE\centering
\begin{minipage}{0.95\linewidth}\raggedright\linespread{1.25}\selectfont Que vaut $4{,}2\times0{,}5$~?\end{minipage}
\end{minipage}

\vfill
\begin{tcolorbox}[enhanced,nobeforeafter,width=7.1cm,height=2.1cm,valign=center,colback=ccal!6,colframe=ccal,arc=2.5mm,boxrule=0.9pt,left=3mm]{\Large\bfseries\color{ccal}a)}\quad\begin{minipage}[c]{5.6cm}\raggedright\Large $2{,}1$\end{minipage}\end{tcolorbox}\hfill\begin{tcolorbox}[enhanced,nobeforeafter,width=7.1cm,height=2.1cm,valign=center,colback=ccal!6,colframe=ccal,arc=2.5mm,boxrule=0.9pt,left=3mm]{\Large\bfseries\color{ccal}b)}\quad\begin{minipage}[c]{5.6cm}\raggedright\Large $4{,}7$\end{minipage}\end{tcolorbox}\par\vspace{2.5mm}
\begin{tcolorbox}[enhanced,nobeforeafter,width=7.1cm,height=2.1cm,valign=center,colback=ccal!6,colframe=ccal,arc=2.5mm,boxrule=0.9pt,left=3mm]{\Large\bfseries\color{ccal}c)}\quad\begin{minipage}[c]{5.6cm}\raggedright\Large $21$\end{minipage}\end{tcolorbox}\hfill\begin{tcolorbox}[enhanced,nobeforeafter,width=7.1cm,height=2.1cm,valign=center,colback=ccal!6,colframe=ccal,arc=2.5mm,boxrule=0.9pt,left=3mm]{\Large\bfseries\color{ccal}d)}\quad\begin{minipage}[c]{5.6cm}\raggedright\Large $0{,}21$\end{minipage}\end{tcolorbox}
\end{tcolorbox}
\newpage
\begin{tcolorbox}[enhanced,width=16cm,height=14cm,colback=white,colframe=ccal,boxrule=2pt,arc=0mm,
left=6mm,right=6mm,top=0mm,bottom=4mm,sharp corners,
overlay={\fill[ccal] (frame.north west) rectangle ([yshift=-1.25cm]frame.north east);
\node[anchor=west,text=white,font=\sffamily\bfseries\LARGE] at ([xshift=5mm,yshift=-0.62cm]frame.north west) {Question 23 C};
\node[anchor=east,text=white,font=\sffamily\itshape\large] at ([xshift=-5mm,yshift=-0.62cm]frame.north east) {Calculs};}]
\vspace*{1.75cm}
\begin{minipage}[t][6.9cm][c]{\linewidth}\LARGE\centering
\begin{minipage}{0.95\linewidth}\raggedright\linespread{1.25}\selectfont Dans la division euclidienne de $500$ par $13$, quels sont le quotient $q$ et le reste $r$~?\end{minipage}
\end{minipage}

\vfill
\begin{tcolorbox}[enhanced,nobeforeafter,width=7.1cm,height=2.1cm,valign=center,colback=ccal!6,colframe=ccal,arc=2.5mm,boxrule=0.9pt,left=3mm]{\Large\bfseries\color{ccal}a)}\quad\begin{minipage}[c]{5.6cm}\raggedright\Large $q=39\text{ et }r=6$\end{minipage}\end{tcolorbox}\hfill\begin{tcolorbox}[enhanced,nobeforeafter,width=7.1cm,height=2.1cm,valign=center,colback=ccal!6,colframe=ccal,arc=2.5mm,boxrule=0.9pt,left=3mm]{\Large\bfseries\color{ccal}b)}\quad\begin{minipage}[c]{5.6cm}\raggedright\Large $q=37\text{ et }r=19$\end{minipage}\end{tcolorbox}\par\vspace{2.5mm}
\begin{tcolorbox}[enhanced,nobeforeafter,width=7.1cm,height=2.1cm,valign=center,colback=ccal!6,colframe=ccal,arc=2.5mm,boxrule=0.9pt,left=3mm]{\Large\bfseries\color{ccal}c)}\quad\begin{minipage}[c]{5.6cm}\raggedright\Large $q=38\text{ et }r=6$\end{minipage}\end{tcolorbox}\hfill\begin{tcolorbox}[enhanced,nobeforeafter,width=7.1cm,height=2.1cm,valign=center,colback=ccal!6,colframe=ccal,arc=2.5mm,boxrule=0.9pt,left=3mm]{\Large\bfseries\color{ccal}d)}\quad\begin{minipage}[c]{5.6cm}\raggedright\Large $q=38\text{ et }r=7$\end{minipage}\end{tcolorbox}
\end{tcolorbox}
\newpage
\begin{tcolorbox}[enhanced,width=16cm,height=14cm,colback=white,colframe=ccal,boxrule=2pt,arc=0mm,
left=6mm,right=6mm,top=0mm,bottom=4mm,sharp corners,
overlay={\fill[ccal] (frame.north west) rectangle ([yshift=-1.25cm]frame.north east);
\node[anchor=west,text=white,font=\sffamily\bfseries\LARGE] at ([xshift=5mm,yshift=-0.62cm]frame.north west) {Question 24 C};
\node[anchor=east,text=white,font=\sffamily\itshape\large] at ([xshift=-5mm,yshift=-0.62cm]frame.north east) {Calculs};}]
\vspace*{1.75cm}
\begin{minipage}[t][6.9cm][c]{\linewidth}\LARGE\centering
\begin{minipage}{0.95\linewidth}\raggedright\linespread{1.25}\selectfont Dans la division euclidienne de $91$ par $4$, quels sont le quotient $q$ et le reste $r$~?\end{minipage}
\end{minipage}

\vfill
\begin{tcolorbox}[enhanced,nobeforeafter,width=7.1cm,height=2.1cm,valign=center,colback=ccal!6,colframe=ccal,arc=2.5mm,boxrule=0.9pt,left=3mm]{\Large\bfseries\color{ccal}a)}\quad\begin{minipage}[c]{5.6cm}\raggedright\Large $q=23\text{ et }r=3$\end{minipage}\end{tcolorbox}\hfill\begin{tcolorbox}[enhanced,nobeforeafter,width=7.1cm,height=2.1cm,valign=center,colback=ccal!6,colframe=ccal,arc=2.5mm,boxrule=0.9pt,left=3mm]{\Large\bfseries\color{ccal}b)}\quad\begin{minipage}[c]{5.6cm}\raggedright\Large $q=22\text{ et }r=4$\end{minipage}\end{tcolorbox}\par\vspace{2.5mm}
\begin{tcolorbox}[enhanced,nobeforeafter,width=7.1cm,height=2.1cm,valign=center,colback=ccal!6,colframe=ccal,arc=2.5mm,boxrule=0.9pt,left=3mm]{\Large\bfseries\color{ccal}c)}\quad\begin{minipage}[c]{5.6cm}\raggedright\Large $q=22\text{ et }r=3$\end{minipage}\end{tcolorbox}\hfill\begin{tcolorbox}[enhanced,nobeforeafter,width=7.1cm,height=2.1cm,valign=center,colback=ccal!6,colframe=ccal,arc=2.5mm,boxrule=0.9pt,left=3mm]{\Large\bfseries\color{ccal}d)}\quad\begin{minipage}[c]{5.6cm}\raggedright\Large $q=21\text{ et }r=7$\end{minipage}\end{tcolorbox}
\end{tcolorbox}
\newpage
\begin{tcolorbox}[enhanced,width=16cm,height=14cm,colback=white,colframe=ccal,boxrule=2pt,arc=0mm,
left=6mm,right=6mm,top=0mm,bottom=4mm,sharp corners,
overlay={\fill[ccal] (frame.north west) rectangle ([yshift=-1.25cm]frame.north east);
\node[anchor=west,text=white,font=\sffamily\bfseries\LARGE] at ([xshift=5mm,yshift=-0.62cm]frame.north west) {Question 25 C};
\node[anchor=east,text=white,font=\sffamily\itshape\large] at ([xshift=-5mm,yshift=-0.62cm]frame.north east) {Calculs};}]
\vspace*{1.75cm}
\begin{minipage}[t][6.9cm][c]{\linewidth}\LARGE\centering
\begin{minipage}{0.95\linewidth}\raggedright\linespread{1.25}\selectfont Lequel de ces nombres est un multiple de $12$~?\end{minipage}
\end{minipage}

\vfill
\begin{tcolorbox}[enhanced,nobeforeafter,width=7.1cm,height=2.1cm,valign=center,colback=ccal!6,colframe=ccal,arc=2.5mm,boxrule=0.9pt,left=3mm]{\Large\bfseries\color{ccal}a)}\quad\begin{minipage}[c]{5.6cm}\raggedright\Large $132$\end{minipage}\end{tcolorbox}\hfill\begin{tcolorbox}[enhanced,nobeforeafter,width=7.1cm,height=2.1cm,valign=center,colback=ccal!6,colframe=ccal,arc=2.5mm,boxrule=0.9pt,left=3mm]{\Large\bfseries\color{ccal}b)}\quad\begin{minipage}[c]{5.6cm}\raggedright\Large $142$\end{minipage}\end{tcolorbox}\par\vspace{2.5mm}
\begin{tcolorbox}[enhanced,nobeforeafter,width=7.1cm,height=2.1cm,valign=center,colback=ccal!6,colframe=ccal,arc=2.5mm,boxrule=0.9pt,left=3mm]{\Large\bfseries\color{ccal}c)}\quad\begin{minipage}[c]{5.6cm}\raggedright\Large $126$\end{minipage}\end{tcolorbox}\hfill\begin{tcolorbox}[enhanced,nobeforeafter,width=7.1cm,height=2.1cm,valign=center,colback=ccal!6,colframe=ccal,arc=2.5mm,boxrule=0.9pt,left=3mm]{\Large\bfseries\color{ccal}d)}\quad\begin{minipage}[c]{5.6cm}\raggedright\Large $122$\end{minipage}\end{tcolorbox}
\end{tcolorbox}
\newpage
\begin{tcolorbox}[enhanced,width=16cm,height=14cm,colback=white,colframe=ccal,boxrule=2pt,arc=0mm,
left=6mm,right=6mm,top=0mm,bottom=4mm,sharp corners,
overlay={\fill[ccal] (frame.north west) rectangle ([yshift=-1.25cm]frame.north east);
\node[anchor=west,text=white,font=\sffamily\bfseries\LARGE] at ([xshift=5mm,yshift=-0.62cm]frame.north west) {Question 26 C};
\node[anchor=east,text=white,font=\sffamily\itshape\large] at ([xshift=-5mm,yshift=-0.62cm]frame.north east) {Calculs};}]
\vspace*{1.75cm}
\begin{minipage}[t][6.9cm][c]{\linewidth}\LARGE\centering
\begin{minipage}{0.95\linewidth}\raggedright\linespread{1.25}\selectfont Lequel de ces nombres est un multiple de $4$~?\end{minipage}
\end{minipage}

\vfill
\begin{tcolorbox}[enhanced,nobeforeafter,width=7.1cm,height=2.1cm,valign=center,colback=ccal!6,colframe=ccal,arc=2.5mm,boxrule=0.9pt,left=3mm]{\Large\bfseries\color{ccal}a)}\quad\begin{minipage}[c]{5.6cm}\raggedright\Large $226$\end{minipage}\end{tcolorbox}\hfill\begin{tcolorbox}[enhanced,nobeforeafter,width=7.1cm,height=2.1cm,valign=center,colback=ccal!6,colframe=ccal,arc=2.5mm,boxrule=0.9pt,left=3mm]{\Large\bfseries\color{ccal}b)}\quad\begin{minipage}[c]{5.6cm}\raggedright\Large $242$\end{minipage}\end{tcolorbox}\par\vspace{2.5mm}
\begin{tcolorbox}[enhanced,nobeforeafter,width=7.1cm,height=2.1cm,valign=center,colback=ccal!6,colframe=ccal,arc=2.5mm,boxrule=0.9pt,left=3mm]{\Large\bfseries\color{ccal}c)}\quad\begin{minipage}[c]{5.6cm}\raggedright\Large $238$\end{minipage}\end{tcolorbox}\hfill\begin{tcolorbox}[enhanced,nobeforeafter,width=7.1cm,height=2.1cm,valign=center,colback=ccal!6,colframe=ccal,arc=2.5mm,boxrule=0.9pt,left=3mm]{\Large\bfseries\color{ccal}d)}\quad\begin{minipage}[c]{5.6cm}\raggedright\Large $236$\end{minipage}\end{tcolorbox}
\end{tcolorbox}
\newpage
\begin{tcolorbox}[enhanced,width=16cm,height=14cm,colback=white,colframe=ccal,boxrule=2pt,arc=0mm,
left=6mm,right=6mm,top=0mm,bottom=4mm,sharp corners,
overlay={\fill[ccal] (frame.north west) rectangle ([yshift=-1.25cm]frame.north east);
\node[anchor=west,text=white,font=\sffamily\bfseries\LARGE] at ([xshift=5mm,yshift=-0.62cm]frame.north west) {Question 27 C};
\node[anchor=east,text=white,font=\sffamily\itshape\large] at ([xshift=-5mm,yshift=-0.62cm]frame.north east) {Calculs};}]
\vspace*{1.75cm}
\begin{minipage}[t][6.9cm][c]{\linewidth}\LARGE\centering
\begin{minipage}{0.95\linewidth}\raggedright\linespread{1.25}\selectfont Lequel de ces nombres est divisible par $10$~?\end{minipage}
\end{minipage}

\vfill
\begin{tcolorbox}[enhanced,nobeforeafter,width=7.1cm,height=2.1cm,valign=center,colback=ccal!6,colframe=ccal,arc=2.5mm,boxrule=0.9pt,left=3mm]{\Large\bfseries\color{ccal}a)}\quad\begin{minipage}[c]{5.6cm}\raggedright\Large $4\,506$\end{minipage}\end{tcolorbox}\hfill\begin{tcolorbox}[enhanced,nobeforeafter,width=7.1cm,height=2.1cm,valign=center,colback=ccal!6,colframe=ccal,arc=2.5mm,boxrule=0.9pt,left=3mm]{\Large\bfseries\color{ccal}b)}\quad\begin{minipage}[c]{5.6cm}\raggedright\Large $6\,045$\end{minipage}\end{tcolorbox}\par\vspace{2.5mm}
\begin{tcolorbox}[enhanced,nobeforeafter,width=7.1cm,height=2.1cm,valign=center,colback=ccal!6,colframe=ccal,arc=2.5mm,boxrule=0.9pt,left=3mm]{\Large\bfseries\color{ccal}c)}\quad\begin{minipage}[c]{5.6cm}\raggedright\Large $4\,560$\end{minipage}\end{tcolorbox}\hfill\begin{tcolorbox}[enhanced,nobeforeafter,width=7.1cm,height=2.1cm,valign=center,colback=ccal!6,colframe=ccal,arc=2.5mm,boxrule=0.9pt,left=3mm]{\Large\bfseries\color{ccal}d)}\quad\begin{minipage}[c]{5.6cm}\raggedright\Large $5\,046$\end{minipage}\end{tcolorbox}
\end{tcolorbox}
\newpage
\begin{tcolorbox}[enhanced,width=16cm,height=14cm,colback=white,colframe=ccal,boxrule=2pt,arc=0mm,
left=6mm,right=6mm,top=0mm,bottom=4mm,sharp corners,
overlay={\fill[ccal] (frame.north west) rectangle ([yshift=-1.25cm]frame.north east);
\node[anchor=west,text=white,font=\sffamily\bfseries\LARGE] at ([xshift=5mm,yshift=-0.62cm]frame.north west) {Question 28 C};
\node[anchor=east,text=white,font=\sffamily\itshape\large] at ([xshift=-5mm,yshift=-0.62cm]frame.north east) {Calculs};}]
\vspace*{1.75cm}
\begin{minipage}[t][6.9cm][c]{\linewidth}\LARGE\centering
\begin{minipage}{0.95\linewidth}\raggedright\linespread{1.25}\selectfont Lequel de ces nombres est divisible par $3$~?\end{minipage}
\end{minipage}

\vfill
\begin{tcolorbox}[enhanced,nobeforeafter,width=7.1cm,height=2.1cm,valign=center,colback=ccal!6,colframe=ccal,arc=2.5mm,boxrule=0.9pt,left=3mm]{\Large\bfseries\color{ccal}a)}\quad\begin{minipage}[c]{5.6cm}\raggedright\Large $1\,411$\end{minipage}\end{tcolorbox}\hfill\begin{tcolorbox}[enhanced,nobeforeafter,width=7.1cm,height=2.1cm,valign=center,colback=ccal!6,colframe=ccal,arc=2.5mm,boxrule=0.9pt,left=3mm]{\Large\bfseries\color{ccal}b)}\quad\begin{minipage}[c]{5.6cm}\raggedright\Large $1\,114$\end{minipage}\end{tcolorbox}\par\vspace{2.5mm}
\begin{tcolorbox}[enhanced,nobeforeafter,width=7.1cm,height=2.1cm,valign=center,colback=ccal!6,colframe=ccal,arc=2.5mm,boxrule=0.9pt,left=3mm]{\Large\bfseries\color{ccal}c)}\quad\begin{minipage}[c]{5.6cm}\raggedright\Large $1\,113$\end{minipage}\end{tcolorbox}\hfill\begin{tcolorbox}[enhanced,nobeforeafter,width=7.1cm,height=2.1cm,valign=center,colback=ccal!6,colframe=ccal,arc=2.5mm,boxrule=0.9pt,left=3mm]{\Large\bfseries\color{ccal}d)}\quad\begin{minipage}[c]{5.6cm}\raggedright\Large $1\,141$\end{minipage}\end{tcolorbox}
\end{tcolorbox}
\newpage
\begin{tcolorbox}[enhanced,width=16cm,height=14cm,colback=white,colframe=ccal,boxrule=2pt,arc=0mm,
left=6mm,right=6mm,top=0mm,bottom=4mm,sharp corners,
overlay={\fill[ccal] (frame.north west) rectangle ([yshift=-1.25cm]frame.north east);
\node[anchor=west,text=white,font=\sffamily\bfseries\LARGE] at ([xshift=5mm,yshift=-0.62cm]frame.north west) {Question 29 C};
\node[anchor=east,text=white,font=\sffamily\itshape\large] at ([xshift=-5mm,yshift=-0.62cm]frame.north east) {Calculs};}]
\vspace*{1.75cm}
\begin{minipage}[t][6.9cm][c]{\linewidth}\LARGE\centering
\begin{minipage}{0.95\linewidth}\raggedright\linespread{1.25}\selectfont Que vaut $30-(8+2)$~?\end{minipage}
\end{minipage}

\vfill
\begin{tcolorbox}[enhanced,nobeforeafter,width=7.1cm,height=2.1cm,valign=center,colback=ccal!6,colframe=ccal,arc=2.5mm,boxrule=0.9pt,left=3mm]{\Large\bfseries\color{ccal}a)}\quad\begin{minipage}[c]{5.6cm}\raggedright\Large $40$\end{minipage}\end{tcolorbox}\hfill\begin{tcolorbox}[enhanced,nobeforeafter,width=7.1cm,height=2.1cm,valign=center,colback=ccal!6,colframe=ccal,arc=2.5mm,boxrule=0.9pt,left=3mm]{\Large\bfseries\color{ccal}b)}\quad\begin{minipage}[c]{5.6cm}\raggedright\Large $20$\end{minipage}\end{tcolorbox}\par\vspace{2.5mm}
\begin{tcolorbox}[enhanced,nobeforeafter,width=7.1cm,height=2.1cm,valign=center,colback=ccal!6,colframe=ccal,arc=2.5mm,boxrule=0.9pt,left=3mm]{\Large\bfseries\color{ccal}c)}\quad\begin{minipage}[c]{5.6cm}\raggedright\Large $24$\end{minipage}\end{tcolorbox}\hfill\begin{tcolorbox}[enhanced,nobeforeafter,width=7.1cm,height=2.1cm,valign=center,colback=ccal!6,colframe=ccal,arc=2.5mm,boxrule=0.9pt,left=3mm]{\Large\bfseries\color{ccal}d)}\quad\begin{minipage}[c]{5.6cm}\raggedright\Large $38$\end{minipage}\end{tcolorbox}
\end{tcolorbox}
\newpage
\begin{tcolorbox}[enhanced,width=16cm,height=14cm,colback=white,colframe=ccal,boxrule=2pt,arc=0mm,
left=6mm,right=6mm,top=0mm,bottom=4mm,sharp corners,
overlay={\fill[ccal] (frame.north west) rectangle ([yshift=-1.25cm]frame.north east);
\node[anchor=west,text=white,font=\sffamily\bfseries\LARGE] at ([xshift=5mm,yshift=-0.62cm]frame.north west) {Question 30 C};
\node[anchor=east,text=white,font=\sffamily\itshape\large] at ([xshift=-5mm,yshift=-0.62cm]frame.north east) {Calculs};}]
\vspace*{1.75cm}
\begin{minipage}[t][6.9cm][c]{\linewidth}\LARGE\centering
\begin{minipage}{0.95\linewidth}\raggedright\linespread{1.25}\selectfont Que vaut $4\times(10-7)$~?\end{minipage}
\end{minipage}

\vfill
\begin{tcolorbox}[enhanced,nobeforeafter,width=7.1cm,height=2.1cm,valign=center,colback=ccal!6,colframe=ccal,arc=2.5mm,boxrule=0.9pt,left=3mm]{\Large\bfseries\color{ccal}a)}\quad\begin{minipage}[c]{5.6cm}\raggedright\Large $33$\end{minipage}\end{tcolorbox}\hfill\begin{tcolorbox}[enhanced,nobeforeafter,width=7.1cm,height=2.1cm,valign=center,colback=ccal!6,colframe=ccal,arc=2.5mm,boxrule=0.9pt,left=3mm]{\Large\bfseries\color{ccal}b)}\quad\begin{minipage}[c]{5.6cm}\raggedright\Large $47$\end{minipage}\end{tcolorbox}\par\vspace{2.5mm}
\begin{tcolorbox}[enhanced,nobeforeafter,width=7.1cm,height=2.1cm,valign=center,colback=ccal!6,colframe=ccal,arc=2.5mm,boxrule=0.9pt,left=3mm]{\Large\bfseries\color{ccal}c)}\quad\begin{minipage}[c]{5.6cm}\raggedright\Large $12$\end{minipage}\end{tcolorbox}\hfill\begin{tcolorbox}[enhanced,nobeforeafter,width=7.1cm,height=2.1cm,valign=center,colback=ccal!6,colframe=ccal,arc=2.5mm,boxrule=0.9pt,left=3mm]{\Large\bfseries\color{ccal}d)}\quad\begin{minipage}[c]{5.6cm}\raggedright\Large $30$\end{minipage}\end{tcolorbox}
\end{tcolorbox}
\newpage
\begin{tcolorbox}[enhanced,width=16cm,height=14cm,colback=white,colframe=ccal,boxrule=2pt,arc=0mm,
left=6mm,right=6mm,top=0mm,bottom=4mm,sharp corners,
overlay={\fill[ccal] (frame.north west) rectangle ([yshift=-1.25cm]frame.north east);
\node[anchor=west,text=white,font=\sffamily\bfseries\LARGE] at ([xshift=5mm,yshift=-0.62cm]frame.north west) {Question 31 C};
\node[anchor=east,text=white,font=\sffamily\itshape\large] at ([xshift=-5mm,yshift=-0.62cm]frame.north east) {Calculs};}]
\vspace*{1.75cm}
\begin{minipage}[t][6.9cm][c]{\linewidth}\LARGE\centering
\begin{minipage}{0.95\linewidth}\raggedright\linespread{1.25}\selectfont Que vaut $21\div4$~?\end{minipage}
\end{minipage}

\vfill
\begin{tcolorbox}[enhanced,nobeforeafter,width=7.1cm,height=2.1cm,valign=center,colback=ccal!6,colframe=ccal,arc=2.5mm,boxrule=0.9pt,left=3mm]{\Large\bfseries\color{ccal}a)}\quad\begin{minipage}[c]{5.6cm}\raggedright\Large $5{,}1$\end{minipage}\end{tcolorbox}\hfill\begin{tcolorbox}[enhanced,nobeforeafter,width=7.1cm,height=2.1cm,valign=center,colback=ccal!6,colframe=ccal,arc=2.5mm,boxrule=0.9pt,left=3mm]{\Large\bfseries\color{ccal}b)}\quad\begin{minipage}[c]{5.6cm}\raggedright\Large $52{,}5$\end{minipage}\end{tcolorbox}\par\vspace{2.5mm}
\begin{tcolorbox}[enhanced,nobeforeafter,width=7.1cm,height=2.1cm,valign=center,colback=ccal!6,colframe=ccal,arc=2.5mm,boxrule=0.9pt,left=3mm]{\Large\bfseries\color{ccal}c)}\quad\begin{minipage}[c]{5.6cm}\raggedright\Large $5{,}25$\end{minipage}\end{tcolorbox}\hfill\begin{tcolorbox}[enhanced,nobeforeafter,width=7.1cm,height=2.1cm,valign=center,colback=ccal!6,colframe=ccal,arc=2.5mm,boxrule=0.9pt,left=3mm]{\Large\bfseries\color{ccal}d)}\quad\begin{minipage}[c]{5.6cm}\raggedright\Large $\dfrac{4}{21}$\end{minipage}\end{tcolorbox}
\end{tcolorbox}
\newpage
\begin{tcolorbox}[enhanced,width=16cm,height=14cm,colback=white,colframe=ccal,boxrule=2pt,arc=0mm,
left=6mm,right=6mm,top=0mm,bottom=4mm,sharp corners,
overlay={\fill[ccal] (frame.north west) rectangle ([yshift=-1.25cm]frame.north east);
\node[anchor=west,text=white,font=\sffamily\bfseries\LARGE] at ([xshift=5mm,yshift=-0.62cm]frame.north west) {Question 32 C};
\node[anchor=east,text=white,font=\sffamily\itshape\large] at ([xshift=-5mm,yshift=-0.62cm]frame.north east) {Calculs};}]
\vspace*{1.75cm}
\begin{minipage}[t][6.9cm][c]{\linewidth}\LARGE\centering
\begin{minipage}{0.95\linewidth}\raggedright\linespread{1.25}\selectfont Que vaut $13\div8$~?\end{minipage}
\end{minipage}

\vfill
\begin{tcolorbox}[enhanced,nobeforeafter,width=7.1cm,height=2.1cm,valign=center,colback=ccal!6,colframe=ccal,arc=2.5mm,boxrule=0.9pt,left=3mm]{\Large\bfseries\color{ccal}a)}\quad\begin{minipage}[c]{5.6cm}\raggedright\Large $\dfrac{8}{13}$\end{minipage}\end{tcolorbox}\hfill\begin{tcolorbox}[enhanced,nobeforeafter,width=7.1cm,height=2.1cm,valign=center,colback=ccal!6,colframe=ccal,arc=2.5mm,boxrule=0.9pt,left=3mm]{\Large\bfseries\color{ccal}b)}\quad\begin{minipage}[c]{5.6cm}\raggedright\Large $1{,}5$\end{minipage}\end{tcolorbox}\par\vspace{2.5mm}
\begin{tcolorbox}[enhanced,nobeforeafter,width=7.1cm,height=2.1cm,valign=center,colback=ccal!6,colframe=ccal,arc=2.5mm,boxrule=0.9pt,left=3mm]{\Large\bfseries\color{ccal}c)}\quad\begin{minipage}[c]{5.6cm}\raggedright\Large $16{,}25$\end{minipage}\end{tcolorbox}\hfill\begin{tcolorbox}[enhanced,nobeforeafter,width=7.1cm,height=2.1cm,valign=center,colback=ccal!6,colframe=ccal,arc=2.5mm,boxrule=0.9pt,left=3mm]{\Large\bfseries\color{ccal}d)}\quad\begin{minipage}[c]{5.6cm}\raggedright\Large $1{,}625$\end{minipage}\end{tcolorbox}
\end{tcolorbox}
\newpage
\begin{tcolorbox}[enhanced,width=16cm,height=14cm,colback=white,colframe=cfra,boxrule=2pt,arc=0mm,
left=6mm,right=6mm,top=0mm,bottom=4mm,sharp corners,
overlay={\fill[cfra] (frame.north west) rectangle ([yshift=-1.25cm]frame.north east);
\node[anchor=west,text=white,font=\sffamily\bfseries\LARGE] at ([xshift=5mm,yshift=-0.62cm]frame.north west) {Question 33 C};
\node[anchor=east,text=white,font=\sffamily\itshape\large] at ([xshift=-5mm,yshift=-0.62cm]frame.north east) {Fractions};}]
\vspace*{1.75cm}
\begin{minipage}[t][6.9cm][c]{\linewidth}\LARGE\centering
\begin{minipage}{0.95\linewidth}\raggedright\linespread{1.25}\selectfont Combien vaut $\dfrac{2}{3}$ de $27$~?\end{minipage}
\end{minipage}

\vfill
\begin{tcolorbox}[enhanced,nobeforeafter,width=7.1cm,height=2.1cm,valign=center,colback=cfra!6,colframe=cfra,arc=2.5mm,boxrule=0.9pt,left=3mm]{\Large\bfseries\color{cfra}a)}\quad\begin{minipage}[c]{5.6cm}\raggedright\Large $54$\end{minipage}\end{tcolorbox}\hfill\begin{tcolorbox}[enhanced,nobeforeafter,width=7.1cm,height=2.1cm,valign=center,colback=cfra!6,colframe=cfra,arc=2.5mm,boxrule=0.9pt,left=3mm]{\Large\bfseries\color{cfra}b)}\quad\begin{minipage}[c]{5.6cm}\raggedright\Large $19$\end{minipage}\end{tcolorbox}\par\vspace{2.5mm}
\begin{tcolorbox}[enhanced,nobeforeafter,width=7.1cm,height=2.1cm,valign=center,colback=cfra!6,colframe=cfra,arc=2.5mm,boxrule=0.9pt,left=3mm]{\Large\bfseries\color{cfra}c)}\quad\begin{minipage}[c]{5.6cm}\raggedright\Large $18$\end{minipage}\end{tcolorbox}\hfill\begin{tcolorbox}[enhanced,nobeforeafter,width=7.1cm,height=2.1cm,valign=center,colback=cfra!6,colframe=cfra,arc=2.5mm,boxrule=0.9pt,left=3mm]{\Large\bfseries\color{cfra}d)}\quad\begin{minipage}[c]{5.6cm}\raggedright\Large $9$\end{minipage}\end{tcolorbox}
\end{tcolorbox}
\newpage
\begin{tcolorbox}[enhanced,width=16cm,height=14cm,colback=white,colframe=cfra,boxrule=2pt,arc=0mm,
left=6mm,right=6mm,top=0mm,bottom=4mm,sharp corners,
overlay={\fill[cfra] (frame.north west) rectangle ([yshift=-1.25cm]frame.north east);
\node[anchor=west,text=white,font=\sffamily\bfseries\LARGE] at ([xshift=5mm,yshift=-0.62cm]frame.north west) {Question 34 C};
\node[anchor=east,text=white,font=\sffamily\itshape\large] at ([xshift=-5mm,yshift=-0.62cm]frame.north east) {Fractions};}]
\vspace*{1.75cm}
\begin{minipage}[t][6.9cm][c]{\linewidth}\LARGE\centering
\begin{minipage}{0.95\linewidth}\raggedright\linespread{1.25}\selectfont Combien vaut $\dfrac{7}{10}$ de $60$~?\end{minipage}
\end{minipage}

\vfill
\begin{tcolorbox}[enhanced,nobeforeafter,width=7.1cm,height=2.1cm,valign=center,colback=cfra!6,colframe=cfra,arc=2.5mm,boxrule=0.9pt,left=3mm]{\Large\bfseries\color{cfra}a)}\quad\begin{minipage}[c]{5.6cm}\raggedright\Large $42$\end{minipage}\end{tcolorbox}\hfill\begin{tcolorbox}[enhanced,nobeforeafter,width=7.1cm,height=2.1cm,valign=center,colback=cfra!6,colframe=cfra,arc=2.5mm,boxrule=0.9pt,left=3mm]{\Large\bfseries\color{cfra}b)}\quad\begin{minipage}[c]{5.6cm}\raggedright\Large $6$\end{minipage}\end{tcolorbox}\par\vspace{2.5mm}
\begin{tcolorbox}[enhanced,nobeforeafter,width=7.1cm,height=2.1cm,valign=center,colback=cfra!6,colframe=cfra,arc=2.5mm,boxrule=0.9pt,left=3mm]{\Large\bfseries\color{cfra}c)}\quad\begin{minipage}[c]{5.6cm}\raggedright\Large $18$\end{minipage}\end{tcolorbox}\hfill\begin{tcolorbox}[enhanced,nobeforeafter,width=7.1cm,height=2.1cm,valign=center,colback=cfra!6,colframe=cfra,arc=2.5mm,boxrule=0.9pt,left=3mm]{\Large\bfseries\color{cfra}d)}\quad\begin{minipage}[c]{5.6cm}\raggedright\Large $420$\end{minipage}\end{tcolorbox}
\end{tcolorbox}
\newpage
\begin{tcolorbox}[enhanced,width=16cm,height=14cm,colback=white,colframe=cfra,boxrule=2pt,arc=0mm,
left=6mm,right=6mm,top=0mm,bottom=4mm,sharp corners,
overlay={\fill[cfra] (frame.north west) rectangle ([yshift=-1.25cm]frame.north east);
\node[anchor=west,text=white,font=\sffamily\bfseries\LARGE] at ([xshift=5mm,yshift=-0.62cm]frame.north west) {Question 35 C};
\node[anchor=east,text=white,font=\sffamily\itshape\large] at ([xshift=-5mm,yshift=-0.62cm]frame.north east) {Fractions};}]
\vspace*{1.75cm}
\begin{minipage}[t][6.9cm][c]{\linewidth}\LARGE\centering
\begin{minipage}{0.95\linewidth}\raggedright\linespread{1.25}\selectfont Compléter~: $\dfrac{4}{5}=\dfrac{?}{30}$\end{minipage}
\end{minipage}

\vfill
\begin{tcolorbox}[enhanced,nobeforeafter,width=7.1cm,height=2.1cm,valign=center,colback=cfra!6,colframe=cfra,arc=2.5mm,boxrule=0.9pt,left=3mm]{\Large\bfseries\color{cfra}a)}\quad\begin{minipage}[c]{5.6cm}\raggedright\Large $6$\end{minipage}\end{tcolorbox}\hfill\begin{tcolorbox}[enhanced,nobeforeafter,width=7.1cm,height=2.1cm,valign=center,colback=cfra!6,colframe=cfra,arc=2.5mm,boxrule=0.9pt,left=3mm]{\Large\bfseries\color{cfra}b)}\quad\begin{minipage}[c]{5.6cm}\raggedright\Large $29$\end{minipage}\end{tcolorbox}\par\vspace{2.5mm}
\begin{tcolorbox}[enhanced,nobeforeafter,width=7.1cm,height=2.1cm,valign=center,colback=cfra!6,colframe=cfra,arc=2.5mm,boxrule=0.9pt,left=3mm]{\Large\bfseries\color{cfra}c)}\quad\begin{minipage}[c]{5.6cm}\raggedright\Large $120$\end{minipage}\end{tcolorbox}\hfill\begin{tcolorbox}[enhanced,nobeforeafter,width=7.1cm,height=2.1cm,valign=center,colback=cfra!6,colframe=cfra,arc=2.5mm,boxrule=0.9pt,left=3mm]{\Large\bfseries\color{cfra}d)}\quad\begin{minipage}[c]{5.6cm}\raggedright\Large $24$\end{minipage}\end{tcolorbox}
\end{tcolorbox}
\newpage
\begin{tcolorbox}[enhanced,width=16cm,height=14cm,colback=white,colframe=cfra,boxrule=2pt,arc=0mm,
left=6mm,right=6mm,top=0mm,bottom=4mm,sharp corners,
overlay={\fill[cfra] (frame.north west) rectangle ([yshift=-1.25cm]frame.north east);
\node[anchor=west,text=white,font=\sffamily\bfseries\LARGE] at ([xshift=5mm,yshift=-0.62cm]frame.north west) {Question 36 C};
\node[anchor=east,text=white,font=\sffamily\itshape\large] at ([xshift=-5mm,yshift=-0.62cm]frame.north east) {Fractions};}]
\vspace*{1.75cm}
\begin{minipage}[t][6.9cm][c]{\linewidth}\LARGE\centering
\begin{minipage}{0.95\linewidth}\raggedright\linespread{1.25}\selectfont Compléter~: $\dfrac{3}{8}=\dfrac{?}{24}$\end{minipage}
\end{minipage}

\vfill
\begin{tcolorbox}[enhanced,nobeforeafter,width=7.1cm,height=2.1cm,valign=center,colback=cfra!6,colframe=cfra,arc=2.5mm,boxrule=0.9pt,left=3mm]{\Large\bfseries\color{cfra}a)}\quad\begin{minipage}[c]{5.6cm}\raggedright\Large $9$\end{minipage}\end{tcolorbox}\hfill\begin{tcolorbox}[enhanced,nobeforeafter,width=7.1cm,height=2.1cm,valign=center,colback=cfra!6,colframe=cfra,arc=2.5mm,boxrule=0.9pt,left=3mm]{\Large\bfseries\color{cfra}b)}\quad\begin{minipage}[c]{5.6cm}\raggedright\Large $3$\end{minipage}\end{tcolorbox}\par\vspace{2.5mm}
\begin{tcolorbox}[enhanced,nobeforeafter,width=7.1cm,height=2.1cm,valign=center,colback=cfra!6,colframe=cfra,arc=2.5mm,boxrule=0.9pt,left=3mm]{\Large\bfseries\color{cfra}c)}\quad\begin{minipage}[c]{5.6cm}\raggedright\Large $72$\end{minipage}\end{tcolorbox}\hfill\begin{tcolorbox}[enhanced,nobeforeafter,width=7.1cm,height=2.1cm,valign=center,colback=cfra!6,colframe=cfra,arc=2.5mm,boxrule=0.9pt,left=3mm]{\Large\bfseries\color{cfra}d)}\quad\begin{minipage}[c]{5.6cm}\raggedright\Large $19$\end{minipage}\end{tcolorbox}
\end{tcolorbox}
\newpage
\begin{tcolorbox}[enhanced,width=16cm,height=14cm,colback=white,colframe=cfra,boxrule=2pt,arc=0mm,
left=6mm,right=6mm,top=0mm,bottom=4mm,sharp corners,
overlay={\fill[cfra] (frame.north west) rectangle ([yshift=-1.25cm]frame.north east);
\node[anchor=west,text=white,font=\sffamily\bfseries\LARGE] at ([xshift=5mm,yshift=-0.62cm]frame.north west) {Question 37 C};
\node[anchor=east,text=white,font=\sffamily\itshape\large] at ([xshift=-5mm,yshift=-0.62cm]frame.north east) {Fractions};}]
\vspace*{1.75cm}
\begin{minipage}[t][6.9cm][c]{\linewidth}\LARGE\centering
\begin{minipage}{0.95\linewidth}\raggedright\linespread{1.25}\selectfont Quelle est l'écriture décimale de $\dfrac{9}{20}$~?\end{minipage}
\end{minipage}

\vfill
\begin{tcolorbox}[enhanced,nobeforeafter,width=7.1cm,height=2.1cm,valign=center,colback=cfra!6,colframe=cfra,arc=2.5mm,boxrule=0.9pt,left=3mm]{\Large\bfseries\color{cfra}a)}\quad\begin{minipage}[c]{5.6cm}\raggedright\Large $0{,}45$\end{minipage}\end{tcolorbox}\hfill\begin{tcolorbox}[enhanced,nobeforeafter,width=7.1cm,height=2.1cm,valign=center,colback=cfra!6,colframe=cfra,arc=2.5mm,boxrule=0.9pt,left=3mm]{\Large\bfseries\color{cfra}b)}\quad\begin{minipage}[c]{5.6cm}\raggedright\Large $\dfrac{20}{9}$\end{minipage}\end{tcolorbox}\par\vspace{2.5mm}
\begin{tcolorbox}[enhanced,nobeforeafter,width=7.1cm,height=2.1cm,valign=center,colback=cfra!6,colframe=cfra,arc=2.5mm,boxrule=0.9pt,left=3mm]{\Large\bfseries\color{cfra}c)}\quad\begin{minipage}[c]{5.6cm}\raggedright\Large $9{,}20$\end{minipage}\end{tcolorbox}\hfill\begin{tcolorbox}[enhanced,nobeforeafter,width=7.1cm,height=2.1cm,valign=center,colback=cfra!6,colframe=cfra,arc=2.5mm,boxrule=0.9pt,left=3mm]{\Large\bfseries\color{cfra}d)}\quad\begin{minipage}[c]{5.6cm}\raggedright\Large $4{,}5$\end{minipage}\end{tcolorbox}
\end{tcolorbox}
\newpage
\begin{tcolorbox}[enhanced,width=16cm,height=14cm,colback=white,colframe=cfra,boxrule=2pt,arc=0mm,
left=6mm,right=6mm,top=0mm,bottom=4mm,sharp corners,
overlay={\fill[cfra] (frame.north west) rectangle ([yshift=-1.25cm]frame.north east);
\node[anchor=west,text=white,font=\sffamily\bfseries\LARGE] at ([xshift=5mm,yshift=-0.62cm]frame.north west) {Question 38 C};
\node[anchor=east,text=white,font=\sffamily\itshape\large] at ([xshift=-5mm,yshift=-0.62cm]frame.north east) {Fractions};}]
\vspace*{1.75cm}
\begin{minipage}[t][6.9cm][c]{\linewidth}\LARGE\centering
\begin{minipage}{0.95\linewidth}\raggedright\linespread{1.25}\selectfont Quelle est l'écriture décimale de $\dfrac{3}{25}$~?\end{minipage}
\end{minipage}

\vfill
\begin{tcolorbox}[enhanced,nobeforeafter,width=7.1cm,height=2.1cm,valign=center,colback=cfra!6,colframe=cfra,arc=2.5mm,boxrule=0.9pt,left=3mm]{\Large\bfseries\color{cfra}a)}\quad\begin{minipage}[c]{5.6cm}\raggedright\Large $\dfrac{25}{3}$\end{minipage}\end{tcolorbox}\hfill\begin{tcolorbox}[enhanced,nobeforeafter,width=7.1cm,height=2.1cm,valign=center,colback=cfra!6,colframe=cfra,arc=2.5mm,boxrule=0.9pt,left=3mm]{\Large\bfseries\color{cfra}b)}\quad\begin{minipage}[c]{5.6cm}\raggedright\Large $1{,}2$\end{minipage}\end{tcolorbox}\par\vspace{2.5mm}
\begin{tcolorbox}[enhanced,nobeforeafter,width=7.1cm,height=2.1cm,valign=center,colback=cfra!6,colframe=cfra,arc=2.5mm,boxrule=0.9pt,left=3mm]{\Large\bfseries\color{cfra}c)}\quad\begin{minipage}[c]{5.6cm}\raggedright\Large $3{,}25$\end{minipage}\end{tcolorbox}\hfill\begin{tcolorbox}[enhanced,nobeforeafter,width=7.1cm,height=2.1cm,valign=center,colback=cfra!6,colframe=cfra,arc=2.5mm,boxrule=0.9pt,left=3mm]{\Large\bfseries\color{cfra}d)}\quad\begin{minipage}[c]{5.6cm}\raggedright\Large $0{,}12$\end{minipage}\end{tcolorbox}
\end{tcolorbox}
\newpage
\begin{tcolorbox}[enhanced,width=16cm,height=14cm,colback=white,colframe=cfra,boxrule=2pt,arc=0mm,
left=6mm,right=6mm,top=0mm,bottom=4mm,sharp corners,
overlay={\fill[cfra] (frame.north west) rectangle ([yshift=-1.25cm]frame.north east);
\node[anchor=west,text=white,font=\sffamily\bfseries\LARGE] at ([xshift=5mm,yshift=-0.62cm]frame.north west) {Question 39 C};
\node[anchor=east,text=white,font=\sffamily\itshape\large] at ([xshift=-5mm,yshift=-0.62cm]frame.north east) {Fractions};}]
\vspace*{1.75cm}
\begin{minipage}[t][6.9cm][c]{\linewidth}\LARGE\centering
\begin{minipage}{0.95\linewidth}\raggedright\linespread{1.25}\selectfont Laquelle de ces fractions est plus grande que $1$~?\end{minipage}
\end{minipage}

\vfill
\begin{tcolorbox}[enhanced,nobeforeafter,width=7.1cm,height=2.1cm,valign=center,colback=cfra!6,colframe=cfra,arc=2.5mm,boxrule=0.9pt,left=3mm]{\Large\bfseries\color{cfra}a)}\quad\begin{minipage}[c]{5.6cm}\raggedright\Large $\dfrac{12}{13}$\end{minipage}\end{tcolorbox}\hfill\begin{tcolorbox}[enhanced,nobeforeafter,width=7.1cm,height=2.1cm,valign=center,colback=cfra!6,colframe=cfra,arc=2.5mm,boxrule=0.9pt,left=3mm]{\Large\bfseries\color{cfra}b)}\quad\begin{minipage}[c]{5.6cm}\raggedright\Large $\dfrac{9}{9}$\end{minipage}\end{tcolorbox}\par\vspace{2.5mm}
\begin{tcolorbox}[enhanced,nobeforeafter,width=7.1cm,height=2.1cm,valign=center,colback=cfra!6,colframe=cfra,arc=2.5mm,boxrule=0.9pt,left=3mm]{\Large\bfseries\color{cfra}c)}\quad\begin{minipage}[c]{5.6cm}\raggedright\Large $\dfrac{1}{3}$\end{minipage}\end{tcolorbox}\hfill\begin{tcolorbox}[enhanced,nobeforeafter,width=7.1cm,height=2.1cm,valign=center,colback=cfra!6,colframe=cfra,arc=2.5mm,boxrule=0.9pt,left=3mm]{\Large\bfseries\color{cfra}d)}\quad\begin{minipage}[c]{5.6cm}\raggedright\Large $\dfrac{13}{12}$\end{minipage}\end{tcolorbox}
\end{tcolorbox}
\newpage
\begin{tcolorbox}[enhanced,width=16cm,height=14cm,colback=white,colframe=cfra,boxrule=2pt,arc=0mm,
left=6mm,right=6mm,top=0mm,bottom=4mm,sharp corners,
overlay={\fill[cfra] (frame.north west) rectangle ([yshift=-1.25cm]frame.north east);
\node[anchor=west,text=white,font=\sffamily\bfseries\LARGE] at ([xshift=5mm,yshift=-0.62cm]frame.north west) {Question 40 C};
\node[anchor=east,text=white,font=\sffamily\itshape\large] at ([xshift=-5mm,yshift=-0.62cm]frame.north east) {Fractions};}]
\vspace*{1.75cm}
\begin{minipage}[t][6.9cm][c]{\linewidth}\LARGE\centering
\begin{minipage}{0.95\linewidth}\raggedright\linespread{1.25}\selectfont Laquelle de ces fractions est plus grande que $1$~?\end{minipage}
\end{minipage}

\vfill
\begin{tcolorbox}[enhanced,nobeforeafter,width=7.1cm,height=2.1cm,valign=center,colback=cfra!6,colframe=cfra,arc=2.5mm,boxrule=0.9pt,left=3mm]{\Large\bfseries\color{cfra}a)}\quad\begin{minipage}[c]{5.6cm}\raggedright\Large $\dfrac{8}{7}$\end{minipage}\end{tcolorbox}\hfill\begin{tcolorbox}[enhanced,nobeforeafter,width=7.1cm,height=2.1cm,valign=center,colback=cfra!6,colframe=cfra,arc=2.5mm,boxrule=0.9pt,left=3mm]{\Large\bfseries\color{cfra}b)}\quad\begin{minipage}[c]{5.6cm}\raggedright\Large $\dfrac{5}{5}$\end{minipage}\end{tcolorbox}\par\vspace{2.5mm}
\begin{tcolorbox}[enhanced,nobeforeafter,width=7.1cm,height=2.1cm,valign=center,colback=cfra!6,colframe=cfra,arc=2.5mm,boxrule=0.9pt,left=3mm]{\Large\bfseries\color{cfra}c)}\quad\begin{minipage}[c]{5.6cm}\raggedright\Large $\dfrac{2}{5}$\end{minipage}\end{tcolorbox}\hfill\begin{tcolorbox}[enhanced,nobeforeafter,width=7.1cm,height=2.1cm,valign=center,colback=cfra!6,colframe=cfra,arc=2.5mm,boxrule=0.9pt,left=3mm]{\Large\bfseries\color{cfra}d)}\quad\begin{minipage}[c]{5.6cm}\raggedright\Large $\dfrac{7}{8}$\end{minipage}\end{tcolorbox}
\end{tcolorbox}
\newpage
\begin{tcolorbox}[enhanced,width=16cm,height=14cm,colback=white,colframe=cfra,boxrule=2pt,arc=0mm,
left=6mm,right=6mm,top=0mm,bottom=4mm,sharp corners,
overlay={\fill[cfra] (frame.north west) rectangle ([yshift=-1.25cm]frame.north east);
\node[anchor=west,text=white,font=\sffamily\bfseries\LARGE] at ([xshift=5mm,yshift=-0.62cm]frame.north west) {Question 41 C};
\node[anchor=east,text=white,font=\sffamily\itshape\large] at ([xshift=-5mm,yshift=-0.62cm]frame.north east) {Fractions};}]
\vspace*{1.75cm}
\begin{minipage}[t][6.9cm][c]{\linewidth}\LARGE\centering
\begin{minipage}{0.95\linewidth}\raggedright\linespread{1.25}\selectfont Quel nombre, multiplié par $6$, donne $1$~?\end{minipage}
\end{minipage}

\vfill
\begin{tcolorbox}[enhanced,nobeforeafter,width=7.1cm,height=2.1cm,valign=center,colback=cfra!6,colframe=cfra,arc=2.5mm,boxrule=0.9pt,left=3mm]{\Large\bfseries\color{cfra}a)}\quad\begin{minipage}[c]{5.6cm}\raggedright\Large $5$\end{minipage}\end{tcolorbox}\hfill\begin{tcolorbox}[enhanced,nobeforeafter,width=7.1cm,height=2.1cm,valign=center,colback=cfra!6,colframe=cfra,arc=2.5mm,boxrule=0.9pt,left=3mm]{\Large\bfseries\color{cfra}b)}\quad\begin{minipage}[c]{5.6cm}\raggedright\Large $6$\end{minipage}\end{tcolorbox}\par\vspace{2.5mm}
\begin{tcolorbox}[enhanced,nobeforeafter,width=7.1cm,height=2.1cm,valign=center,colback=cfra!6,colframe=cfra,arc=2.5mm,boxrule=0.9pt,left=3mm]{\Large\bfseries\color{cfra}c)}\quad\begin{minipage}[c]{5.6cm}\raggedright\Large $\dfrac{1}{6}$\end{minipage}\end{tcolorbox}\hfill\begin{tcolorbox}[enhanced,nobeforeafter,width=7.1cm,height=2.1cm,valign=center,colback=cfra!6,colframe=cfra,arc=2.5mm,boxrule=0.9pt,left=3mm]{\Large\bfseries\color{cfra}d)}\quad\begin{minipage}[c]{5.6cm}\raggedright\Large $\dfrac{6}{1}$\end{minipage}\end{tcolorbox}
\end{tcolorbox}
\newpage
\begin{tcolorbox}[enhanced,width=16cm,height=14cm,colback=white,colframe=cfra,boxrule=2pt,arc=0mm,
left=6mm,right=6mm,top=0mm,bottom=4mm,sharp corners,
overlay={\fill[cfra] (frame.north west) rectangle ([yshift=-1.25cm]frame.north east);
\node[anchor=west,text=white,font=\sffamily\bfseries\LARGE] at ([xshift=5mm,yshift=-0.62cm]frame.north west) {Question 42 C};
\node[anchor=east,text=white,font=\sffamily\itshape\large] at ([xshift=-5mm,yshift=-0.62cm]frame.north east) {Fractions};}]
\vspace*{1.75cm}
\begin{minipage}[t][6.9cm][c]{\linewidth}\LARGE\centering
\begin{minipage}{0.95\linewidth}\raggedright\linespread{1.25}\selectfont Quel nombre, multiplié par $10$, donne $7$~?\end{minipage}
\end{minipage}

\vfill
\begin{tcolorbox}[enhanced,nobeforeafter,width=7.1cm,height=2.1cm,valign=center,colback=cfra!6,colframe=cfra,arc=2.5mm,boxrule=0.9pt,left=3mm]{\Large\bfseries\color{cfra}a)}\quad\begin{minipage}[c]{5.6cm}\raggedright\Large $3$\end{minipage}\end{tcolorbox}\hfill\begin{tcolorbox}[enhanced,nobeforeafter,width=7.1cm,height=2.1cm,valign=center,colback=cfra!6,colframe=cfra,arc=2.5mm,boxrule=0.9pt,left=3mm]{\Large\bfseries\color{cfra}b)}\quad\begin{minipage}[c]{5.6cm}\raggedright\Large $70$\end{minipage}\end{tcolorbox}\par\vspace{2.5mm}
\begin{tcolorbox}[enhanced,nobeforeafter,width=7.1cm,height=2.1cm,valign=center,colback=cfra!6,colframe=cfra,arc=2.5mm,boxrule=0.9pt,left=3mm]{\Large\bfseries\color{cfra}c)}\quad\begin{minipage}[c]{5.6cm}\raggedright\Large $\dfrac{10}{7}$\end{minipage}\end{tcolorbox}\hfill\begin{tcolorbox}[enhanced,nobeforeafter,width=7.1cm,height=2.1cm,valign=center,colback=cfra!6,colframe=cfra,arc=2.5mm,boxrule=0.9pt,left=3mm]{\Large\bfseries\color{cfra}d)}\quad\begin{minipage}[c]{5.6cm}\raggedright\Large $\dfrac{7}{10}$\end{minipage}\end{tcolorbox}
\end{tcolorbox}
\newpage
\begin{tcolorbox}[enhanced,width=16cm,height=14cm,colback=white,colframe=cfra,boxrule=2pt,arc=0mm,
left=6mm,right=6mm,top=0mm,bottom=4mm,sharp corners,
overlay={\fill[cfra] (frame.north west) rectangle ([yshift=-1.25cm]frame.north east);
\node[anchor=west,text=white,font=\sffamily\bfseries\LARGE] at ([xshift=5mm,yshift=-0.62cm]frame.north west) {Question 43 C};
\node[anchor=east,text=white,font=\sffamily\itshape\large] at ([xshift=-5mm,yshift=-0.62cm]frame.north east) {Fractions};}]
\vspace*{1.75cm}
\begin{minipage}[t][6.9cm][c]{\linewidth}\LARGE\centering
\begin{minipage}{0.95\linewidth}\raggedright\linespread{1.25}\selectfont Que vaut $\dfrac{3}{8}+\dfrac{4}{8}$~?\end{minipage}
\end{minipage}

\vfill
\begin{tcolorbox}[enhanced,nobeforeafter,width=7.1cm,height=2.1cm,valign=center,colback=cfra!6,colframe=cfra,arc=2.5mm,boxrule=0.9pt,left=3mm]{\Large\bfseries\color{cfra}a)}\quad\begin{minipage}[c]{5.6cm}\raggedright\Large $\dfrac{12}{8}$\end{minipage}\end{tcolorbox}\hfill\begin{tcolorbox}[enhanced,nobeforeafter,width=7.1cm,height=2.1cm,valign=center,colback=cfra!6,colframe=cfra,arc=2.5mm,boxrule=0.9pt,left=3mm]{\Large\bfseries\color{cfra}b)}\quad\begin{minipage}[c]{5.6cm}\raggedright\Large $\dfrac{7}{16}$\end{minipage}\end{tcolorbox}\par\vspace{2.5mm}
\begin{tcolorbox}[enhanced,nobeforeafter,width=7.1cm,height=2.1cm,valign=center,colback=cfra!6,colframe=cfra,arc=2.5mm,boxrule=0.9pt,left=3mm]{\Large\bfseries\color{cfra}c)}\quad\begin{minipage}[c]{5.6cm}\raggedright\Large $\dfrac{7}{8}$\end{minipage}\end{tcolorbox}\hfill\begin{tcolorbox}[enhanced,nobeforeafter,width=7.1cm,height=2.1cm,valign=center,colback=cfra!6,colframe=cfra,arc=2.5mm,boxrule=0.9pt,left=3mm]{\Large\bfseries\color{cfra}d)}\quad\begin{minipage}[c]{5.6cm}\raggedright\Large $\dfrac{7}{64}$\end{minipage}\end{tcolorbox}
\end{tcolorbox}
\newpage
\begin{tcolorbox}[enhanced,width=16cm,height=14cm,colback=white,colframe=cfra,boxrule=2pt,arc=0mm,
left=6mm,right=6mm,top=0mm,bottom=4mm,sharp corners,
overlay={\fill[cfra] (frame.north west) rectangle ([yshift=-1.25cm]frame.north east);
\node[anchor=west,text=white,font=\sffamily\bfseries\LARGE] at ([xshift=5mm,yshift=-0.62cm]frame.north west) {Question 44 C};
\node[anchor=east,text=white,font=\sffamily\itshape\large] at ([xshift=-5mm,yshift=-0.62cm]frame.north east) {Fractions};}]
\vspace*{1.75cm}
\begin{minipage}[t][6.9cm][c]{\linewidth}\LARGE\centering
\begin{minipage}{0.95\linewidth}\raggedright\linespread{1.25}\selectfont Que vaut $\dfrac{2}{13}+\dfrac{9}{13}$~?\end{minipage}
\end{minipage}

\vfill
\begin{tcolorbox}[enhanced,nobeforeafter,width=7.1cm,height=2.1cm,valign=center,colback=cfra!6,colframe=cfra,arc=2.5mm,boxrule=0.9pt,left=3mm]{\Large\bfseries\color{cfra}a)}\quad\begin{minipage}[c]{5.6cm}\raggedright\Large $\dfrac{18}{13}$\end{minipage}\end{tcolorbox}\hfill\begin{tcolorbox}[enhanced,nobeforeafter,width=7.1cm,height=2.1cm,valign=center,colback=cfra!6,colframe=cfra,arc=2.5mm,boxrule=0.9pt,left=3mm]{\Large\bfseries\color{cfra}b)}\quad\begin{minipage}[c]{5.6cm}\raggedright\Large $\dfrac{11}{13}$\end{minipage}\end{tcolorbox}\par\vspace{2.5mm}
\begin{tcolorbox}[enhanced,nobeforeafter,width=7.1cm,height=2.1cm,valign=center,colback=cfra!6,colframe=cfra,arc=2.5mm,boxrule=0.9pt,left=3mm]{\Large\bfseries\color{cfra}c)}\quad\begin{minipage}[c]{5.6cm}\raggedright\Large $\dfrac{11}{169}$\end{minipage}\end{tcolorbox}\hfill\begin{tcolorbox}[enhanced,nobeforeafter,width=7.1cm,height=2.1cm,valign=center,colback=cfra!6,colframe=cfra,arc=2.5mm,boxrule=0.9pt,left=3mm]{\Large\bfseries\color{cfra}d)}\quad\begin{minipage}[c]{5.6cm}\raggedright\Large $\dfrac{11}{26}$\end{minipage}\end{tcolorbox}
\end{tcolorbox}
\newpage
\begin{tcolorbox}[enhanced,width=16cm,height=14cm,colback=white,colframe=cfra,boxrule=2pt,arc=0mm,
left=6mm,right=6mm,top=0mm,bottom=4mm,sharp corners,
overlay={\fill[cfra] (frame.north west) rectangle ([yshift=-1.25cm]frame.north east);
\node[anchor=west,text=white,font=\sffamily\bfseries\LARGE] at ([xshift=5mm,yshift=-0.62cm]frame.north west) {Question 45 C};
\node[anchor=east,text=white,font=\sffamily\itshape\large] at ([xshift=-5mm,yshift=-0.62cm]frame.north east) {Fractions};}]
\vspace*{1.75cm}
\begin{minipage}[t][6.9cm][c]{\linewidth}\LARGE\centering
\begin{minipage}{0.95\linewidth}\raggedright\linespread{1.25}\selectfont Un gâteau est coupé en $5$ parts égales. On mange $2$ parts. Quelle fraction du gâteau reste-t-il~?\end{minipage}
\end{minipage}

\vfill
\begin{tcolorbox}[enhanced,nobeforeafter,width=7.1cm,height=2.1cm,valign=center,colback=cfra!6,colframe=cfra,arc=2.5mm,boxrule=0.9pt,left=3mm]{\Large\bfseries\color{cfra}a)}\quad\begin{minipage}[c]{5.6cm}\raggedright\Large $\dfrac{2}{5}$\end{minipage}\end{tcolorbox}\hfill\begin{tcolorbox}[enhanced,nobeforeafter,width=7.1cm,height=2.1cm,valign=center,colback=cfra!6,colframe=cfra,arc=2.5mm,boxrule=0.9pt,left=3mm]{\Large\bfseries\color{cfra}b)}\quad\begin{minipage}[c]{5.6cm}\raggedright\Large $\dfrac{3}{5}$\end{minipage}\end{tcolorbox}\par\vspace{2.5mm}
\begin{tcolorbox}[enhanced,nobeforeafter,width=7.1cm,height=2.1cm,valign=center,colback=cfra!6,colframe=cfra,arc=2.5mm,boxrule=0.9pt,left=3mm]{\Large\bfseries\color{cfra}c)}\quad\begin{minipage}[c]{5.6cm}\raggedright\Large $\dfrac{2}{3}$\end{minipage}\end{tcolorbox}\hfill\begin{tcolorbox}[enhanced,nobeforeafter,width=7.1cm,height=2.1cm,valign=center,colback=cfra!6,colframe=cfra,arc=2.5mm,boxrule=0.9pt,left=3mm]{\Large\bfseries\color{cfra}d)}\quad\begin{minipage}[c]{5.6cm}\raggedright\Large $\dfrac{3}{2}$\end{minipage}\end{tcolorbox}
\end{tcolorbox}
\newpage
\begin{tcolorbox}[enhanced,width=16cm,height=14cm,colback=white,colframe=cfra,boxrule=2pt,arc=0mm,
left=6mm,right=6mm,top=0mm,bottom=4mm,sharp corners,
overlay={\fill[cfra] (frame.north west) rectangle ([yshift=-1.25cm]frame.north east);
\node[anchor=west,text=white,font=\sffamily\bfseries\LARGE] at ([xshift=5mm,yshift=-0.62cm]frame.north west) {Question 46 C};
\node[anchor=east,text=white,font=\sffamily\itshape\large] at ([xshift=-5mm,yshift=-0.62cm]frame.north east) {Fractions};}]
\vspace*{1.75cm}
\begin{minipage}[t][6.9cm][c]{\linewidth}\LARGE\centering
\begin{minipage}{0.95\linewidth}\raggedright\linespread{1.25}\selectfont Un gâteau est coupé en $9$ parts égales. On mange $4$ parts. Quelle fraction du gâteau reste-t-il~?\end{minipage}
\end{minipage}

\vfill
\begin{tcolorbox}[enhanced,nobeforeafter,width=7.1cm,height=2.1cm,valign=center,colback=cfra!6,colframe=cfra,arc=2.5mm,boxrule=0.9pt,left=3mm]{\Large\bfseries\color{cfra}a)}\quad\begin{minipage}[c]{5.6cm}\raggedright\Large $\dfrac{5}{9}$\end{minipage}\end{tcolorbox}\hfill\begin{tcolorbox}[enhanced,nobeforeafter,width=7.1cm,height=2.1cm,valign=center,colback=cfra!6,colframe=cfra,arc=2.5mm,boxrule=0.9pt,left=3mm]{\Large\bfseries\color{cfra}b)}\quad\begin{minipage}[c]{5.6cm}\raggedright\Large $\dfrac{5}{4}$\end{minipage}\end{tcolorbox}\par\vspace{2.5mm}
\begin{tcolorbox}[enhanced,nobeforeafter,width=7.1cm,height=2.1cm,valign=center,colback=cfra!6,colframe=cfra,arc=2.5mm,boxrule=0.9pt,left=3mm]{\Large\bfseries\color{cfra}c)}\quad\begin{minipage}[c]{5.6cm}\raggedright\Large $\dfrac{4}{9}$\end{minipage}\end{tcolorbox}\hfill\begin{tcolorbox}[enhanced,nobeforeafter,width=7.1cm,height=2.1cm,valign=center,colback=cfra!6,colframe=cfra,arc=2.5mm,boxrule=0.9pt,left=3mm]{\Large\bfseries\color{cfra}d)}\quad\begin{minipage}[c]{5.6cm}\raggedright\Large $\dfrac{4}{5}$\end{minipage}\end{tcolorbox}
\end{tcolorbox}
\newpage
\begin{tcolorbox}[enhanced,width=16cm,height=14cm,colback=white,colframe=cfra,boxrule=2pt,arc=0mm,
left=6mm,right=6mm,top=0mm,bottom=4mm,sharp corners,
overlay={\fill[cfra] (frame.north west) rectangle ([yshift=-1.25cm]frame.north east);
\node[anchor=west,text=white,font=\sffamily\bfseries\LARGE] at ([xshift=5mm,yshift=-0.62cm]frame.north west) {Question 47 C};
\node[anchor=east,text=white,font=\sffamily\itshape\large] at ([xshift=-5mm,yshift=-0.62cm]frame.north east) {Fractions};}]
\vspace*{1.75cm}
\begin{minipage}[t][6.9cm][c]{\linewidth}\LARGE\centering
\begin{minipage}{0.95\linewidth}\raggedright\linespread{1.25}\selectfont Entre quels nombres entiers consécutifs se trouve $\dfrac{50}{7}$~?\end{minipage}
\end{minipage}

\vfill
\begin{tcolorbox}[enhanced,nobeforeafter,width=7.1cm,height=2.1cm,valign=center,colback=cfra!6,colframe=cfra,arc=2.5mm,boxrule=0.9pt,left=3mm]{\Large\bfseries\color{cfra}a)}\quad\begin{minipage}[c]{5.6cm}\raggedright\Large $8\text{ et }9$\end{minipage}\end{tcolorbox}\hfill\begin{tcolorbox}[enhanced,nobeforeafter,width=7.1cm,height=2.1cm,valign=center,colback=cfra!6,colframe=cfra,arc=2.5mm,boxrule=0.9pt,left=3mm]{\Large\bfseries\color{cfra}b)}\quad\begin{minipage}[c]{5.6cm}\raggedright\Large $6\text{ et }7$\end{minipage}\end{tcolorbox}\par\vspace{2.5mm}
\begin{tcolorbox}[enhanced,nobeforeafter,width=7.1cm,height=2.1cm,valign=center,colback=cfra!6,colframe=cfra,arc=2.5mm,boxrule=0.9pt,left=3mm]{\Large\bfseries\color{cfra}c)}\quad\begin{minipage}[c]{5.6cm}\raggedright\Large $9\text{ et }10$\end{minipage}\end{tcolorbox}\hfill\begin{tcolorbox}[enhanced,nobeforeafter,width=7.1cm,height=2.1cm,valign=center,colback=cfra!6,colframe=cfra,arc=2.5mm,boxrule=0.9pt,left=3mm]{\Large\bfseries\color{cfra}d)}\quad\begin{minipage}[c]{5.6cm}\raggedright\Large $7\text{ et }8$\end{minipage}\end{tcolorbox}
\end{tcolorbox}
\newpage
\begin{tcolorbox}[enhanced,width=16cm,height=14cm,colback=white,colframe=cfra,boxrule=2pt,arc=0mm,
left=6mm,right=6mm,top=0mm,bottom=4mm,sharp corners,
overlay={\fill[cfra] (frame.north west) rectangle ([yshift=-1.25cm]frame.north east);
\node[anchor=west,text=white,font=\sffamily\bfseries\LARGE] at ([xshift=5mm,yshift=-0.62cm]frame.north west) {Question 48 C};
\node[anchor=east,text=white,font=\sffamily\itshape\large] at ([xshift=-5mm,yshift=-0.62cm]frame.north east) {Fractions};}]
\vspace*{1.75cm}
\begin{minipage}[t][6.9cm][c]{\linewidth}\LARGE\centering
\begin{minipage}{0.95\linewidth}\raggedright\linespread{1.25}\selectfont Entre quels nombres entiers consécutifs se trouve $\dfrac{25}{3}$~?\end{minipage}
\end{minipage}

\vfill
\begin{tcolorbox}[enhanced,nobeforeafter,width=7.1cm,height=2.1cm,valign=center,colback=cfra!6,colframe=cfra,arc=2.5mm,boxrule=0.9pt,left=3mm]{\Large\bfseries\color{cfra}a)}\quad\begin{minipage}[c]{5.6cm}\raggedright\Large $10\text{ et }11$\end{minipage}\end{tcolorbox}\hfill\begin{tcolorbox}[enhanced,nobeforeafter,width=7.1cm,height=2.1cm,valign=center,colback=cfra!6,colframe=cfra,arc=2.5mm,boxrule=0.9pt,left=3mm]{\Large\bfseries\color{cfra}b)}\quad\begin{minipage}[c]{5.6cm}\raggedright\Large $8\text{ et }9$\end{minipage}\end{tcolorbox}\par\vspace{2.5mm}
\begin{tcolorbox}[enhanced,nobeforeafter,width=7.1cm,height=2.1cm,valign=center,colback=cfra!6,colframe=cfra,arc=2.5mm,boxrule=0.9pt,left=3mm]{\Large\bfseries\color{cfra}c)}\quad\begin{minipage}[c]{5.6cm}\raggedright\Large $9\text{ et }10$\end{minipage}\end{tcolorbox}\hfill\begin{tcolorbox}[enhanced,nobeforeafter,width=7.1cm,height=2.1cm,valign=center,colback=cfra!6,colframe=cfra,arc=2.5mm,boxrule=0.9pt,left=3mm]{\Large\bfseries\color{cfra}d)}\quad\begin{minipage}[c]{5.6cm}\raggedright\Large $7\text{ et }8$\end{minipage}\end{tcolorbox}
\end{tcolorbox}
\newpage
\begin{tcolorbox}[enhanced,width=16cm,height=14cm,colback=white,colframe=cgeo,boxrule=2pt,arc=0mm,
left=6mm,right=6mm,top=0mm,bottom=4mm,sharp corners,
overlay={\fill[cgeo] (frame.north west) rectangle ([yshift=-1.25cm]frame.north east);
\node[anchor=west,text=white,font=\sffamily\bfseries\LARGE] at ([xshift=5mm,yshift=-0.62cm]frame.north west) {Question 49 C};
\node[anchor=east,text=white,font=\sffamily\itshape\large] at ([xshift=-5mm,yshift=-0.62cm]frame.north east) {Géométrie};}]
\vspace*{1.75cm}
\begin{minipage}[t][6.9cm][c]{\linewidth}\LARGE\centering
\begin{minipage}{0.95\linewidth}\raggedright\linespread{1.25}\selectfont Un angle mesure $89^\circ$. Quelle est sa nature~?\end{minipage}
\end{minipage}

\vfill
\begin{tcolorbox}[enhanced,nobeforeafter,width=7.1cm,height=2.1cm,valign=center,colback=cgeo!6,colframe=cgeo,arc=2.5mm,boxrule=0.9pt,left=3mm]{\Large\bfseries\color{cgeo}a)}\quad\begin{minipage}[c]{5.6cm}\raggedright\Large obtus\end{minipage}\end{tcolorbox}\hfill\begin{tcolorbox}[enhanced,nobeforeafter,width=7.1cm,height=2.1cm,valign=center,colback=cgeo!6,colframe=cgeo,arc=2.5mm,boxrule=0.9pt,left=3mm]{\Large\bfseries\color{cgeo}b)}\quad\begin{minipage}[c]{5.6cm}\raggedright\Large plat\end{minipage}\end{tcolorbox}\par\vspace{2.5mm}
\begin{tcolorbox}[enhanced,nobeforeafter,width=7.1cm,height=2.1cm,valign=center,colback=cgeo!6,colframe=cgeo,arc=2.5mm,boxrule=0.9pt,left=3mm]{\Large\bfseries\color{cgeo}c)}\quad\begin{minipage}[c]{5.6cm}\raggedright\Large aigu\end{minipage}\end{tcolorbox}\hfill\begin{tcolorbox}[enhanced,nobeforeafter,width=7.1cm,height=2.1cm,valign=center,colback=cgeo!6,colframe=cgeo,arc=2.5mm,boxrule=0.9pt,left=3mm]{\Large\bfseries\color{cgeo}d)}\quad\begin{minipage}[c]{5.6cm}\raggedright\Large droit\end{minipage}\end{tcolorbox}
\end{tcolorbox}
\newpage
\begin{tcolorbox}[enhanced,width=16cm,height=14cm,colback=white,colframe=cgeo,boxrule=2pt,arc=0mm,
left=6mm,right=6mm,top=0mm,bottom=4mm,sharp corners,
overlay={\fill[cgeo] (frame.north west) rectangle ([yshift=-1.25cm]frame.north east);
\node[anchor=west,text=white,font=\sffamily\bfseries\LARGE] at ([xshift=5mm,yshift=-0.62cm]frame.north west) {Question 50 C};
\node[anchor=east,text=white,font=\sffamily\itshape\large] at ([xshift=-5mm,yshift=-0.62cm]frame.north east) {Géométrie};}]
\vspace*{1.75cm}
\begin{minipage}[t][6.9cm][c]{\linewidth}\LARGE\centering
\begin{minipage}{0.95\linewidth}\raggedright\linespread{1.25}\selectfont Un angle mesure $91^\circ$. Quelle est sa nature~?\end{minipage}
\end{minipage}

\vfill
\begin{tcolorbox}[enhanced,nobeforeafter,width=7.1cm,height=2.1cm,valign=center,colback=cgeo!6,colframe=cgeo,arc=2.5mm,boxrule=0.9pt,left=3mm]{\Large\bfseries\color{cgeo}a)}\quad\begin{minipage}[c]{5.6cm}\raggedright\Large plat\end{minipage}\end{tcolorbox}\hfill\begin{tcolorbox}[enhanced,nobeforeafter,width=7.1cm,height=2.1cm,valign=center,colback=cgeo!6,colframe=cgeo,arc=2.5mm,boxrule=0.9pt,left=3mm]{\Large\bfseries\color{cgeo}b)}\quad\begin{minipage}[c]{5.6cm}\raggedright\Large obtus\end{minipage}\end{tcolorbox}\par\vspace{2.5mm}
\begin{tcolorbox}[enhanced,nobeforeafter,width=7.1cm,height=2.1cm,valign=center,colback=cgeo!6,colframe=cgeo,arc=2.5mm,boxrule=0.9pt,left=3mm]{\Large\bfseries\color{cgeo}c)}\quad\begin{minipage}[c]{5.6cm}\raggedright\Large aigu\end{minipage}\end{tcolorbox}\hfill\begin{tcolorbox}[enhanced,nobeforeafter,width=7.1cm,height=2.1cm,valign=center,colback=cgeo!6,colframe=cgeo,arc=2.5mm,boxrule=0.9pt,left=3mm]{\Large\bfseries\color{cgeo}d)}\quad\begin{minipage}[c]{5.6cm}\raggedright\Large droit\end{minipage}\end{tcolorbox}
\end{tcolorbox}
\newpage
\begin{tcolorbox}[enhanced,width=16cm,height=14cm,colback=white,colframe=cgeo,boxrule=2pt,arc=0mm,
left=6mm,right=6mm,top=0mm,bottom=4mm,sharp corners,
overlay={\fill[cgeo] (frame.north west) rectangle ([yshift=-1.25cm]frame.north east);
\node[anchor=west,text=white,font=\sffamily\bfseries\LARGE] at ([xshift=5mm,yshift=-0.62cm]frame.north west) {Question 51 C};
\node[anchor=east,text=white,font=\sffamily\itshape\large] at ([xshift=-5mm,yshift=-0.62cm]frame.north east) {Géométrie};}]
\vspace*{1.75cm}
\begin{minipage}[t][6.9cm][c]{\linewidth}\LARGE\centering
\begin{minipage}{0.95\linewidth}\raggedright\linespread{1.25}\selectfont L'angle droit $\widehat{xOy}$ est partagé par la demi-droite $[Oz)$ en deux angles. $\widehat{xOz}=61^\circ$. Combien mesure $\widehat{zOy}$~?\end{minipage}
\end{minipage}

\vfill
\begin{tcolorbox}[enhanced,nobeforeafter,width=7.1cm,height=2.1cm,valign=center,colback=cgeo!6,colframe=cgeo,arc=2.5mm,boxrule=0.9pt,left=3mm]{\Large\bfseries\color{cgeo}a)}\quad\begin{minipage}[c]{5.6cm}\raggedright\Large $119^\circ$\end{minipage}\end{tcolorbox}\hfill\begin{tcolorbox}[enhanced,nobeforeafter,width=7.1cm,height=2.1cm,valign=center,colback=cgeo!6,colframe=cgeo,arc=2.5mm,boxrule=0.9pt,left=3mm]{\Large\bfseries\color{cgeo}b)}\quad\begin{minipage}[c]{5.6cm}\raggedright\Large $29^\circ$\end{minipage}\end{tcolorbox}\par\vspace{2.5mm}
\begin{tcolorbox}[enhanced,nobeforeafter,width=7.1cm,height=2.1cm,valign=center,colback=cgeo!6,colframe=cgeo,arc=2.5mm,boxrule=0.9pt,left=3mm]{\Large\bfseries\color{cgeo}c)}\quad\begin{minipage}[c]{5.6cm}\raggedright\Large $151^\circ$\end{minipage}\end{tcolorbox}\hfill\begin{tcolorbox}[enhanced,nobeforeafter,width=7.1cm,height=2.1cm,valign=center,colback=cgeo!6,colframe=cgeo,arc=2.5mm,boxrule=0.9pt,left=3mm]{\Large\bfseries\color{cgeo}d)}\quad\begin{minipage}[c]{5.6cm}\raggedright\Large $61^\circ$\end{minipage}\end{tcolorbox}
\end{tcolorbox}
\newpage
\begin{tcolorbox}[enhanced,width=16cm,height=14cm,colback=white,colframe=cgeo,boxrule=2pt,arc=0mm,
left=6mm,right=6mm,top=0mm,bottom=4mm,sharp corners,
overlay={\fill[cgeo] (frame.north west) rectangle ([yshift=-1.25cm]frame.north east);
\node[anchor=west,text=white,font=\sffamily\bfseries\LARGE] at ([xshift=5mm,yshift=-0.62cm]frame.north west) {Question 52 C};
\node[anchor=east,text=white,font=\sffamily\itshape\large] at ([xshift=-5mm,yshift=-0.62cm]frame.north east) {Géométrie};}]
\vspace*{1.75cm}
\begin{minipage}[t][6.9cm][c]{\linewidth}\LARGE\centering
\begin{minipage}{0.95\linewidth}\raggedright\linespread{1.25}\selectfont L'angle droit $\widehat{xOy}$ est partagé par la demi-droite $[Oz)$ en deux angles. $\widehat{xOz}=13^\circ$. Combien mesure $\widehat{zOy}$~?\end{minipage}
\end{minipage}

\vfill
\begin{tcolorbox}[enhanced,nobeforeafter,width=7.1cm,height=2.1cm,valign=center,colback=cgeo!6,colframe=cgeo,arc=2.5mm,boxrule=0.9pt,left=3mm]{\Large\bfseries\color{cgeo}a)}\quad\begin{minipage}[c]{5.6cm}\raggedright\Large $77^\circ$\end{minipage}\end{tcolorbox}\hfill\begin{tcolorbox}[enhanced,nobeforeafter,width=7.1cm,height=2.1cm,valign=center,colback=cgeo!6,colframe=cgeo,arc=2.5mm,boxrule=0.9pt,left=3mm]{\Large\bfseries\color{cgeo}b)}\quad\begin{minipage}[c]{5.6cm}\raggedright\Large $103^\circ$\end{minipage}\end{tcolorbox}\par\vspace{2.5mm}
\begin{tcolorbox}[enhanced,nobeforeafter,width=7.1cm,height=2.1cm,valign=center,colback=cgeo!6,colframe=cgeo,arc=2.5mm,boxrule=0.9pt,left=3mm]{\Large\bfseries\color{cgeo}c)}\quad\begin{minipage}[c]{5.6cm}\raggedright\Large $13^\circ$\end{minipage}\end{tcolorbox}\hfill\begin{tcolorbox}[enhanced,nobeforeafter,width=7.1cm,height=2.1cm,valign=center,colback=cgeo!6,colframe=cgeo,arc=2.5mm,boxrule=0.9pt,left=3mm]{\Large\bfseries\color{cgeo}d)}\quad\begin{minipage}[c]{5.6cm}\raggedright\Large $167^\circ$\end{minipage}\end{tcolorbox}
\end{tcolorbox}
\newpage
\begin{tcolorbox}[enhanced,width=16cm,height=14cm,colback=white,colframe=cgeo,boxrule=2pt,arc=0mm,
left=6mm,right=6mm,top=0mm,bottom=4mm,sharp corners,
overlay={\fill[cgeo] (frame.north west) rectangle ([yshift=-1.25cm]frame.north east);
\node[anchor=west,text=white,font=\sffamily\bfseries\LARGE] at ([xshift=5mm,yshift=-0.62cm]frame.north west) {Question 53 C};
\node[anchor=east,text=white,font=\sffamily\itshape\large] at ([xshift=-5mm,yshift=-0.62cm]frame.north east) {Géométrie};}]
\vspace*{1.75cm}
\begin{minipage}[t][6.9cm][c]{\linewidth}\LARGE\centering
\begin{minipage}{0.95\linewidth}\raggedright\linespread{1.25}\selectfont Les points $A$, $O$ et $B$ sont alignés dans cet ordre et $\widehat{AOC}=100^\circ$. Combien mesure $\widehat{COB}$~?\end{minipage}
\end{minipage}

\vfill
\begin{tcolorbox}[enhanced,nobeforeafter,width=7.1cm,height=2.1cm,valign=center,colback=cgeo!6,colframe=cgeo,arc=2.5mm,boxrule=0.9pt,left=3mm]{\Large\bfseries\color{cgeo}a)}\quad\begin{minipage}[c]{5.6cm}\raggedright\Large $260^\circ$\end{minipage}\end{tcolorbox}\hfill\begin{tcolorbox}[enhanced,nobeforeafter,width=7.1cm,height=2.1cm,valign=center,colback=cgeo!6,colframe=cgeo,arc=2.5mm,boxrule=0.9pt,left=3mm]{\Large\bfseries\color{cgeo}b)}\quad\begin{minipage}[c]{5.6cm}\raggedright\Large $10^\circ$\end{minipage}\end{tcolorbox}\par\vspace{2.5mm}
\begin{tcolorbox}[enhanced,nobeforeafter,width=7.1cm,height=2.1cm,valign=center,colback=cgeo!6,colframe=cgeo,arc=2.5mm,boxrule=0.9pt,left=3mm]{\Large\bfseries\color{cgeo}c)}\quad\begin{minipage}[c]{5.6cm}\raggedright\Large $80^\circ$\end{minipage}\end{tcolorbox}\hfill\begin{tcolorbox}[enhanced,nobeforeafter,width=7.1cm,height=2.1cm,valign=center,colback=cgeo!6,colframe=cgeo,arc=2.5mm,boxrule=0.9pt,left=3mm]{\Large\bfseries\color{cgeo}d)}\quad\begin{minipage}[c]{5.6cm}\raggedright\Large $100^\circ$\end{minipage}\end{tcolorbox}
\end{tcolorbox}
\newpage
\begin{tcolorbox}[enhanced,width=16cm,height=14cm,colback=white,colframe=cgeo,boxrule=2pt,arc=0mm,
left=6mm,right=6mm,top=0mm,bottom=4mm,sharp corners,
overlay={\fill[cgeo] (frame.north west) rectangle ([yshift=-1.25cm]frame.north east);
\node[anchor=west,text=white,font=\sffamily\bfseries\LARGE] at ([xshift=5mm,yshift=-0.62cm]frame.north west) {Question 54 C};
\node[anchor=east,text=white,font=\sffamily\itshape\large] at ([xshift=-5mm,yshift=-0.62cm]frame.north east) {Géométrie};}]
\vspace*{1.75cm}
\begin{minipage}[t][6.9cm][c]{\linewidth}\LARGE\centering
\begin{minipage}{0.95\linewidth}\raggedright\linespread{1.25}\selectfont Les points $A$, $O$ et $B$ sont alignés dans cet ordre et $\widehat{AOC}=77^\circ$. Combien mesure $\widehat{COB}$~?\end{minipage}
\end{minipage}

\vfill
\begin{tcolorbox}[enhanced,nobeforeafter,width=7.1cm,height=2.1cm,valign=center,colback=cgeo!6,colframe=cgeo,arc=2.5mm,boxrule=0.9pt,left=3mm]{\Large\bfseries\color{cgeo}a)}\quad\begin{minipage}[c]{5.6cm}\raggedright\Large $77^\circ$\end{minipage}\end{tcolorbox}\hfill\begin{tcolorbox}[enhanced,nobeforeafter,width=7.1cm,height=2.1cm,valign=center,colback=cgeo!6,colframe=cgeo,arc=2.5mm,boxrule=0.9pt,left=3mm]{\Large\bfseries\color{cgeo}b)}\quad\begin{minipage}[c]{5.6cm}\raggedright\Large $283^\circ$\end{minipage}\end{tcolorbox}\par\vspace{2.5mm}
\begin{tcolorbox}[enhanced,nobeforeafter,width=7.1cm,height=2.1cm,valign=center,colback=cgeo!6,colframe=cgeo,arc=2.5mm,boxrule=0.9pt,left=3mm]{\Large\bfseries\color{cgeo}c)}\quad\begin{minipage}[c]{5.6cm}\raggedright\Large $13^\circ$\end{minipage}\end{tcolorbox}\hfill\begin{tcolorbox}[enhanced,nobeforeafter,width=7.1cm,height=2.1cm,valign=center,colback=cgeo!6,colframe=cgeo,arc=2.5mm,boxrule=0.9pt,left=3mm]{\Large\bfseries\color{cgeo}d)}\quad\begin{minipage}[c]{5.6cm}\raggedright\Large $103^\circ$\end{minipage}\end{tcolorbox}
\end{tcolorbox}
\newpage
\begin{tcolorbox}[enhanced,width=16cm,height=14cm,colback=white,colframe=cgeo,boxrule=2pt,arc=0mm,
left=6mm,right=6mm,top=0mm,bottom=4mm,sharp corners,
overlay={\fill[cgeo] (frame.north west) rectangle ([yshift=-1.25cm]frame.north east);
\node[anchor=west,text=white,font=\sffamily\bfseries\LARGE] at ([xshift=5mm,yshift=-0.62cm]frame.north west) {Question 55 C};
\node[anchor=east,text=white,font=\sffamily\itshape\large] at ([xshift=-5mm,yshift=-0.62cm]frame.north east) {Géométrie};}]
\vspace*{1.75cm}
\begin{minipage}[t][6.9cm][c]{\linewidth}\LARGE\centering
\begin{minipage}{0.95\linewidth}\raggedright\linespread{1.25}\selectfont Combien d'axes de symétrie possède un triangle isocèle qui n'est pas équilatéral~?\end{minipage}
\end{minipage}

\vfill
\begin{tcolorbox}[enhanced,nobeforeafter,width=7.1cm,height=2.1cm,valign=center,colback=cgeo!6,colframe=cgeo,arc=2.5mm,boxrule=0.9pt,left=3mm]{\Large\bfseries\color{cgeo}a)}\quad\begin{minipage}[c]{5.6cm}\raggedright\Large 2\end{minipage}\end{tcolorbox}\hfill\begin{tcolorbox}[enhanced,nobeforeafter,width=7.1cm,height=2.1cm,valign=center,colback=cgeo!6,colframe=cgeo,arc=2.5mm,boxrule=0.9pt,left=3mm]{\Large\bfseries\color{cgeo}b)}\quad\begin{minipage}[c]{5.6cm}\raggedright\Large 4\end{minipage}\end{tcolorbox}\par\vspace{2.5mm}
\begin{tcolorbox}[enhanced,nobeforeafter,width=7.1cm,height=2.1cm,valign=center,colback=cgeo!6,colframe=cgeo,arc=2.5mm,boxrule=0.9pt,left=3mm]{\Large\bfseries\color{cgeo}c)}\quad\begin{minipage}[c]{5.6cm}\raggedright\Large 1\end{minipage}\end{tcolorbox}\hfill\begin{tcolorbox}[enhanced,nobeforeafter,width=7.1cm,height=2.1cm,valign=center,colback=cgeo!6,colframe=cgeo,arc=2.5mm,boxrule=0.9pt,left=3mm]{\Large\bfseries\color{cgeo}d)}\quad\begin{minipage}[c]{5.6cm}\raggedright\Large 3\end{minipage}\end{tcolorbox}
\end{tcolorbox}
\newpage
\begin{tcolorbox}[enhanced,width=16cm,height=14cm,colback=white,colframe=cgeo,boxrule=2pt,arc=0mm,
left=6mm,right=6mm,top=0mm,bottom=4mm,sharp corners,
overlay={\fill[cgeo] (frame.north west) rectangle ([yshift=-1.25cm]frame.north east);
\node[anchor=west,text=white,font=\sffamily\bfseries\LARGE] at ([xshift=5mm,yshift=-0.62cm]frame.north west) {Question 56 C};
\node[anchor=east,text=white,font=\sffamily\itshape\large] at ([xshift=-5mm,yshift=-0.62cm]frame.north east) {Géométrie};}]
\vspace*{1.75cm}
\begin{minipage}[t][6.9cm][c]{\linewidth}\LARGE\centering
\begin{minipage}{0.95\linewidth}\raggedright\linespread{1.25}\selectfont Combien d'axes de symétrie possède un triangle quelconque~?\end{minipage}
\end{minipage}

\vfill
\begin{tcolorbox}[enhanced,nobeforeafter,width=7.1cm,height=2.1cm,valign=center,colback=cgeo!6,colframe=cgeo,arc=2.5mm,boxrule=0.9pt,left=3mm]{\Large\bfseries\color{cgeo}a)}\quad\begin{minipage}[c]{5.6cm}\raggedright\Large 1\end{minipage}\end{tcolorbox}\hfill\begin{tcolorbox}[enhanced,nobeforeafter,width=7.1cm,height=2.1cm,valign=center,colback=cgeo!6,colframe=cgeo,arc=2.5mm,boxrule=0.9pt,left=3mm]{\Large\bfseries\color{cgeo}b)}\quad\begin{minipage}[c]{5.6cm}\raggedright\Large 2\end{minipage}\end{tcolorbox}\par\vspace{2.5mm}
\begin{tcolorbox}[enhanced,nobeforeafter,width=7.1cm,height=2.1cm,valign=center,colback=cgeo!6,colframe=cgeo,arc=2.5mm,boxrule=0.9pt,left=3mm]{\Large\bfseries\color{cgeo}c)}\quad\begin{minipage}[c]{5.6cm}\raggedright\Large 0\end{minipage}\end{tcolorbox}\hfill\begin{tcolorbox}[enhanced,nobeforeafter,width=7.1cm,height=2.1cm,valign=center,colback=cgeo!6,colframe=cgeo,arc=2.5mm,boxrule=0.9pt,left=3mm]{\Large\bfseries\color{cgeo}d)}\quad\begin{minipage}[c]{5.6cm}\raggedright\Large 4\end{minipage}\end{tcolorbox}
\end{tcolorbox}
\newpage
\begin{tcolorbox}[enhanced,width=16cm,height=14cm,colback=white,colframe=cgeo,boxrule=2pt,arc=0mm,
left=6mm,right=6mm,top=0mm,bottom=4mm,sharp corners,
overlay={\fill[cgeo] (frame.north west) rectangle ([yshift=-1.25cm]frame.north east);
\node[anchor=west,text=white,font=\sffamily\bfseries\LARGE] at ([xshift=5mm,yshift=-0.62cm]frame.north west) {Question 57 C};
\node[anchor=east,text=white,font=\sffamily\itshape\large] at ([xshift=-5mm,yshift=-0.62cm]frame.north east) {Géométrie};}]
\vspace*{1.75cm}
\begin{minipage}[t][6.9cm][c]{\linewidth}\LARGE\centering
\begin{minipage}{0.95\linewidth}\raggedright\linespread{1.25}\selectfont Un cercle a un diamètre de $20$ cm. Quelle est la longueur de son rayon~?\end{minipage}
\end{minipage}

\vfill
\begin{tcolorbox}[enhanced,nobeforeafter,width=7.1cm,height=2.1cm,valign=center,colback=cgeo!6,colframe=cgeo,arc=2.5mm,boxrule=0.9pt,left=3mm]{\Large\bfseries\color{cgeo}a)}\quad\begin{minipage}[c]{5.6cm}\raggedright\Large $40~\text{cm}$\end{minipage}\end{tcolorbox}\hfill\begin{tcolorbox}[enhanced,nobeforeafter,width=7.1cm,height=2.1cm,valign=center,colback=cgeo!6,colframe=cgeo,arc=2.5mm,boxrule=0.9pt,left=3mm]{\Large\bfseries\color{cgeo}b)}\quad\begin{minipage}[c]{5.6cm}\raggedright\Large $22~\text{cm}$\end{minipage}\end{tcolorbox}\par\vspace{2.5mm}
\begin{tcolorbox}[enhanced,nobeforeafter,width=7.1cm,height=2.1cm,valign=center,colback=cgeo!6,colframe=cgeo,arc=2.5mm,boxrule=0.9pt,left=3mm]{\Large\bfseries\color{cgeo}c)}\quad\begin{minipage}[c]{5.6cm}\raggedright\Large $20~\text{cm}$\end{minipage}\end{tcolorbox}\hfill\begin{tcolorbox}[enhanced,nobeforeafter,width=7.1cm,height=2.1cm,valign=center,colback=cgeo!6,colframe=cgeo,arc=2.5mm,boxrule=0.9pt,left=3mm]{\Large\bfseries\color{cgeo}d)}\quad\begin{minipage}[c]{5.6cm}\raggedright\Large $10~\text{cm}$\end{minipage}\end{tcolorbox}
\end{tcolorbox}
\newpage
\begin{tcolorbox}[enhanced,width=16cm,height=14cm,colback=white,colframe=cgeo,boxrule=2pt,arc=0mm,
left=6mm,right=6mm,top=0mm,bottom=4mm,sharp corners,
overlay={\fill[cgeo] (frame.north west) rectangle ([yshift=-1.25cm]frame.north east);
\node[anchor=west,text=white,font=\sffamily\bfseries\LARGE] at ([xshift=5mm,yshift=-0.62cm]frame.north west) {Question 58 C};
\node[anchor=east,text=white,font=\sffamily\itshape\large] at ([xshift=-5mm,yshift=-0.62cm]frame.north east) {Géométrie};}]
\vspace*{1.75cm}
\begin{minipage}[t][6.9cm][c]{\linewidth}\LARGE\centering
\begin{minipage}{0.95\linewidth}\raggedright\linespread{1.25}\selectfont Un cercle a un rayon de $8$ cm. Quelle est la longueur de son diamètre~?\end{minipage}
\end{minipage}

\vfill
\begin{tcolorbox}[enhanced,nobeforeafter,width=7.1cm,height=2.1cm,valign=center,colback=cgeo!6,colframe=cgeo,arc=2.5mm,boxrule=0.9pt,left=3mm]{\Large\bfseries\color{cgeo}a)}\quad\begin{minipage}[c]{5.6cm}\raggedright\Large $4~\text{cm}$\end{minipage}\end{tcolorbox}\hfill\begin{tcolorbox}[enhanced,nobeforeafter,width=7.1cm,height=2.1cm,valign=center,colback=cgeo!6,colframe=cgeo,arc=2.5mm,boxrule=0.9pt,left=3mm]{\Large\bfseries\color{cgeo}b)}\quad\begin{minipage}[c]{5.6cm}\raggedright\Large $16~\text{cm}$\end{minipage}\end{tcolorbox}\par\vspace{2.5mm}
\begin{tcolorbox}[enhanced,nobeforeafter,width=7.1cm,height=2.1cm,valign=center,colback=cgeo!6,colframe=cgeo,arc=2.5mm,boxrule=0.9pt,left=3mm]{\Large\bfseries\color{cgeo}c)}\quad\begin{minipage}[c]{5.6cm}\raggedright\Large $10~\text{cm}$\end{minipage}\end{tcolorbox}\hfill\begin{tcolorbox}[enhanced,nobeforeafter,width=7.1cm,height=2.1cm,valign=center,colback=cgeo!6,colframe=cgeo,arc=2.5mm,boxrule=0.9pt,left=3mm]{\Large\bfseries\color{cgeo}d)}\quad\begin{minipage}[c]{5.6cm}\raggedright\Large $8~\text{cm}$\end{minipage}\end{tcolorbox}
\end{tcolorbox}
\newpage
\begin{tcolorbox}[enhanced,width=16cm,height=14cm,colback=white,colframe=cgeo,boxrule=2pt,arc=0mm,
left=6mm,right=6mm,top=0mm,bottom=4mm,sharp corners,
overlay={\fill[cgeo] (frame.north west) rectangle ([yshift=-1.25cm]frame.north east);
\node[anchor=west,text=white,font=\sffamily\bfseries\LARGE] at ([xshift=5mm,yshift=-0.62cm]frame.north west) {Question 59 C};
\node[anchor=east,text=white,font=\sffamily\itshape\large] at ([xshift=-5mm,yshift=-0.62cm]frame.north east) {Géométrie};}]
\vspace*{1.75cm}
\begin{minipage}[t][6.9cm][c]{\linewidth}\LARGE\centering
\begin{minipage}{0.95\linewidth}\raggedright\linespread{1.25}\selectfont La droite $(m)$ est perpendiculaire à $(p)$ et la droite $(n)$ est perpendiculaire à $(p)$. Que peut-on dire des droites $(m)$ et $(n)$~?\end{minipage}
\end{minipage}

\vfill
\begin{tcolorbox}[enhanced,nobeforeafter,width=7.1cm,height=2.1cm,valign=center,colback=cgeo!6,colframe=cgeo,arc=2.5mm,boxrule=0.9pt,left=3mm]{\Large\bfseries\color{cgeo}a)}\quad\begin{minipage}[c]{5.6cm}\raggedright\Large on ne peut pas savoir\end{minipage}\end{tcolorbox}\hfill\begin{tcolorbox}[enhanced,nobeforeafter,width=7.1cm,height=2.1cm,valign=center,colback=cgeo!6,colframe=cgeo,arc=2.5mm,boxrule=0.9pt,left=3mm]{\Large\bfseries\color{cgeo}b)}\quad\begin{minipage}[c]{5.6cm}\raggedright\Large parallèles\end{minipage}\end{tcolorbox}\par\vspace{2.5mm}
\begin{tcolorbox}[enhanced,nobeforeafter,width=7.1cm,height=2.1cm,valign=center,colback=cgeo!6,colframe=cgeo,arc=2.5mm,boxrule=0.9pt,left=3mm]{\Large\bfseries\color{cgeo}c)}\quad\begin{minipage}[c]{5.6cm}\raggedright\large sécantes, non perpendiculaires\end{minipage}\end{tcolorbox}\hfill\begin{tcolorbox}[enhanced,nobeforeafter,width=7.1cm,height=2.1cm,valign=center,colback=cgeo!6,colframe=cgeo,arc=2.5mm,boxrule=0.9pt,left=3mm]{\Large\bfseries\color{cgeo}d)}\quad\begin{minipage}[c]{5.6cm}\raggedright\Large perpendiculaires\end{minipage}\end{tcolorbox}
\end{tcolorbox}
\newpage
\begin{tcolorbox}[enhanced,width=16cm,height=14cm,colback=white,colframe=cgeo,boxrule=2pt,arc=0mm,
left=6mm,right=6mm,top=0mm,bottom=4mm,sharp corners,
overlay={\fill[cgeo] (frame.north west) rectangle ([yshift=-1.25cm]frame.north east);
\node[anchor=west,text=white,font=\sffamily\bfseries\LARGE] at ([xshift=5mm,yshift=-0.62cm]frame.north west) {Question 60 C};
\node[anchor=east,text=white,font=\sffamily\itshape\large] at ([xshift=-5mm,yshift=-0.62cm]frame.north east) {Géométrie};}]
\vspace*{1.75cm}
\begin{minipage}[t][6.9cm][c]{\linewidth}\LARGE\centering
\begin{minipage}{0.95\linewidth}\raggedright\linespread{1.25}\selectfont La droite $(x)$ est parallèle à $(z)$ et la droite $(y)$ est parallèle à $(z)$. Que peut-on dire des droites $(x)$ et $(y)$~?\end{minipage}
\end{minipage}

\vfill
\begin{tcolorbox}[enhanced,nobeforeafter,width=7.1cm,height=2.1cm,valign=center,colback=cgeo!6,colframe=cgeo,arc=2.5mm,boxrule=0.9pt,left=3mm]{\Large\bfseries\color{cgeo}a)}\quad\begin{minipage}[c]{5.6cm}\raggedright\Large on ne peut pas savoir\end{minipage}\end{tcolorbox}\hfill\begin{tcolorbox}[enhanced,nobeforeafter,width=7.1cm,height=2.1cm,valign=center,colback=cgeo!6,colframe=cgeo,arc=2.5mm,boxrule=0.9pt,left=3mm]{\Large\bfseries\color{cgeo}b)}\quad\begin{minipage}[c]{5.6cm}\raggedright\large sécantes, non perpendiculaires\end{minipage}\end{tcolorbox}\par\vspace{2.5mm}
\begin{tcolorbox}[enhanced,nobeforeafter,width=7.1cm,height=2.1cm,valign=center,colback=cgeo!6,colframe=cgeo,arc=2.5mm,boxrule=0.9pt,left=3mm]{\Large\bfseries\color{cgeo}c)}\quad\begin{minipage}[c]{5.6cm}\raggedright\Large parallèles\end{minipage}\end{tcolorbox}\hfill\begin{tcolorbox}[enhanced,nobeforeafter,width=7.1cm,height=2.1cm,valign=center,colback=cgeo!6,colframe=cgeo,arc=2.5mm,boxrule=0.9pt,left=3mm]{\Large\bfseries\color{cgeo}d)}\quad\begin{minipage}[c]{5.6cm}\raggedright\Large perpendiculaires\end{minipage}\end{tcolorbox}
\end{tcolorbox}
\newpage
\begin{tcolorbox}[enhanced,width=16cm,height=14cm,colback=white,colframe=cgeo,boxrule=2pt,arc=0mm,
left=6mm,right=6mm,top=0mm,bottom=4mm,sharp corners,
overlay={\fill[cgeo] (frame.north west) rectangle ([yshift=-1.25cm]frame.north east);
\node[anchor=west,text=white,font=\sffamily\bfseries\LARGE] at ([xshift=5mm,yshift=-0.62cm]frame.north west) {Question 61 C};
\node[anchor=east,text=white,font=\sffamily\itshape\large] at ([xshift=-5mm,yshift=-0.62cm]frame.north east) {Géométrie};}]
\vspace*{1.75cm}
\begin{minipage}[t][6.9cm][c]{\linewidth}\LARGE\centering
\begin{minipage}{0.95\linewidth}\raggedright\linespread{1.25}\selectfont Un triangle a deux côtés de $8$ cm et un troisième de $3$ cm. Quelle est sa nature~?\end{minipage}
\end{minipage}

\vfill
\begin{tcolorbox}[enhanced,nobeforeafter,width=7.1cm,height=2.1cm,valign=center,colback=cgeo!6,colframe=cgeo,arc=2.5mm,boxrule=0.9pt,left=3mm]{\Large\bfseries\color{cgeo}a)}\quad\begin{minipage}[c]{5.6cm}\raggedright\Large rectangle\end{minipage}\end{tcolorbox}\hfill\begin{tcolorbox}[enhanced,nobeforeafter,width=7.1cm,height=2.1cm,valign=center,colback=cgeo!6,colframe=cgeo,arc=2.5mm,boxrule=0.9pt,left=3mm]{\Large\bfseries\color{cgeo}b)}\quad\begin{minipage}[c]{5.6cm}\raggedright\Large équilatéral\end{minipage}\end{tcolorbox}\par\vspace{2.5mm}
\begin{tcolorbox}[enhanced,nobeforeafter,width=7.1cm,height=2.1cm,valign=center,colback=cgeo!6,colframe=cgeo,arc=2.5mm,boxrule=0.9pt,left=3mm]{\Large\bfseries\color{cgeo}c)}\quad\begin{minipage}[c]{5.6cm}\raggedright\Large isocèle\end{minipage}\end{tcolorbox}\hfill\begin{tcolorbox}[enhanced,nobeforeafter,width=7.1cm,height=2.1cm,valign=center,colback=cgeo!6,colframe=cgeo,arc=2.5mm,boxrule=0.9pt,left=3mm]{\Large\bfseries\color{cgeo}d)}\quad\begin{minipage}[c]{5.6cm}\raggedright\Large quelconque\end{minipage}\end{tcolorbox}
\end{tcolorbox}
\newpage
\begin{tcolorbox}[enhanced,width=16cm,height=14cm,colback=white,colframe=cgeo,boxrule=2pt,arc=0mm,
left=6mm,right=6mm,top=0mm,bottom=4mm,sharp corners,
overlay={\fill[cgeo] (frame.north west) rectangle ([yshift=-1.25cm]frame.north east);
\node[anchor=west,text=white,font=\sffamily\bfseries\LARGE] at ([xshift=5mm,yshift=-0.62cm]frame.north west) {Question 62 C};
\node[anchor=east,text=white,font=\sffamily\itshape\large] at ([xshift=-5mm,yshift=-0.62cm]frame.north east) {Géométrie};}]
\vspace*{1.75cm}
\begin{minipage}[t][6.9cm][c]{\linewidth}\LARGE\centering
\begin{minipage}{0.95\linewidth}\raggedright\linespread{1.25}\selectfont Un triangle a un angle droit et deux côtés de l'angle droit de même longueur. Quelle est sa nature~?\end{minipage}
\end{minipage}

\vfill
\begin{tcolorbox}[enhanced,nobeforeafter,width=7.1cm,height=2.1cm,valign=center,colback=cgeo!6,colframe=cgeo,arc=2.5mm,boxrule=0.9pt,left=3mm]{\Large\bfseries\color{cgeo}a)}\quad\begin{minipage}[c]{5.6cm}\raggedright\Large isocèle\end{minipage}\end{tcolorbox}\hfill\begin{tcolorbox}[enhanced,nobeforeafter,width=7.1cm,height=2.1cm,valign=center,colback=cgeo!6,colframe=cgeo,arc=2.5mm,boxrule=0.9pt,left=3mm]{\Large\bfseries\color{cgeo}b)}\quad\begin{minipage}[c]{5.6cm}\raggedright\Large rectangle\end{minipage}\end{tcolorbox}\par\vspace{2.5mm}
\begin{tcolorbox}[enhanced,nobeforeafter,width=7.1cm,height=2.1cm,valign=center,colback=cgeo!6,colframe=cgeo,arc=2.5mm,boxrule=0.9pt,left=3mm]{\Large\bfseries\color{cgeo}c)}\quad\begin{minipage}[c]{5.6cm}\raggedright\Large isocèle rectangle\end{minipage}\end{tcolorbox}\hfill\begin{tcolorbox}[enhanced,nobeforeafter,width=7.1cm,height=2.1cm,valign=center,colback=cgeo!6,colframe=cgeo,arc=2.5mm,boxrule=0.9pt,left=3mm]{\Large\bfseries\color{cgeo}d)}\quad\begin{minipage}[c]{5.6cm}\raggedright\Large équilatéral\end{minipage}\end{tcolorbox}
\end{tcolorbox}
\newpage
\begin{tcolorbox}[enhanced,width=16cm,height=14cm,colback=white,colframe=cgeo,boxrule=2pt,arc=0mm,
left=6mm,right=6mm,top=0mm,bottom=4mm,sharp corners,
overlay={\fill[cgeo] (frame.north west) rectangle ([yshift=-1.25cm]frame.north east);
\node[anchor=west,text=white,font=\sffamily\bfseries\LARGE] at ([xshift=5mm,yshift=-0.62cm]frame.north west) {Question 63 C};
\node[anchor=east,text=white,font=\sffamily\itshape\large] at ([xshift=-5mm,yshift=-0.62cm]frame.north east) {Géométrie};}]
\vspace*{1.75cm}
\begin{minipage}[t][6.9cm][c]{\linewidth}\LARGE\centering
\begin{minipage}{0.95\linewidth}\raggedright\linespread{1.25}\selectfont Un quadrilatère a des diagonales de même longueur qui se coupent en leur milieu. Quelle est sa nature la plus précise~?\end{minipage}
\end{minipage}

\vfill
\begin{tcolorbox}[enhanced,nobeforeafter,width=7.1cm,height=2.1cm,valign=center,colback=cgeo!6,colframe=cgeo,arc=2.5mm,boxrule=0.9pt,left=3mm]{\Large\bfseries\color{cgeo}a)}\quad\begin{minipage}[c]{5.6cm}\raggedright\Large un rectangle\end{minipage}\end{tcolorbox}\hfill\begin{tcolorbox}[enhanced,nobeforeafter,width=7.1cm,height=2.1cm,valign=center,colback=cgeo!6,colframe=cgeo,arc=2.5mm,boxrule=0.9pt,left=3mm]{\Large\bfseries\color{cgeo}b)}\quad\begin{minipage}[c]{5.6cm}\raggedright\Large un carré\end{minipage}\end{tcolorbox}\par\vspace{2.5mm}
\begin{tcolorbox}[enhanced,nobeforeafter,width=7.1cm,height=2.1cm,valign=center,colback=cgeo!6,colframe=cgeo,arc=2.5mm,boxrule=0.9pt,left=3mm]{\Large\bfseries\color{cgeo}c)}\quad\begin{minipage}[c]{5.6cm}\raggedright\Large un losange\end{minipage}\end{tcolorbox}\hfill\begin{tcolorbox}[enhanced,nobeforeafter,width=7.1cm,height=2.1cm,valign=center,colback=cgeo!6,colframe=cgeo,arc=2.5mm,boxrule=0.9pt,left=3mm]{\Large\bfseries\color{cgeo}d)}\quad\begin{minipage}[c]{5.6cm}\raggedright\Large un triangle\end{minipage}\end{tcolorbox}
\end{tcolorbox}
\newpage
\begin{tcolorbox}[enhanced,width=16cm,height=14cm,colback=white,colframe=cgeo,boxrule=2pt,arc=0mm,
left=6mm,right=6mm,top=0mm,bottom=4mm,sharp corners,
overlay={\fill[cgeo] (frame.north west) rectangle ([yshift=-1.25cm]frame.north east);
\node[anchor=west,text=white,font=\sffamily\bfseries\LARGE] at ([xshift=5mm,yshift=-0.62cm]frame.north west) {Question 64 C};
\node[anchor=east,text=white,font=\sffamily\itshape\large] at ([xshift=-5mm,yshift=-0.62cm]frame.north east) {Géométrie};}]
\vspace*{1.75cm}
\begin{minipage}[t][6.9cm][c]{\linewidth}\LARGE\centering
\begin{minipage}{0.95\linewidth}\raggedright\linespread{1.25}\selectfont Un quadrilatère a des diagonales perpendiculaires qui se coupent en leur milieu. Quelle est sa nature la plus précise~?\end{minipage}
\end{minipage}

\vfill
\begin{tcolorbox}[enhanced,nobeforeafter,width=7.1cm,height=2.1cm,valign=center,colback=cgeo!6,colframe=cgeo,arc=2.5mm,boxrule=0.9pt,left=3mm]{\Large\bfseries\color{cgeo}a)}\quad\begin{minipage}[c]{5.6cm}\raggedright\Large un carré\end{minipage}\end{tcolorbox}\hfill\begin{tcolorbox}[enhanced,nobeforeafter,width=7.1cm,height=2.1cm,valign=center,colback=cgeo!6,colframe=cgeo,arc=2.5mm,boxrule=0.9pt,left=3mm]{\Large\bfseries\color{cgeo}b)}\quad\begin{minipage}[c]{5.6cm}\raggedright\Large un triangle\end{minipage}\end{tcolorbox}\par\vspace{2.5mm}
\begin{tcolorbox}[enhanced,nobeforeafter,width=7.1cm,height=2.1cm,valign=center,colback=cgeo!6,colframe=cgeo,arc=2.5mm,boxrule=0.9pt,left=3mm]{\Large\bfseries\color{cgeo}c)}\quad\begin{minipage}[c]{5.6cm}\raggedright\Large un rectangle\end{minipage}\end{tcolorbox}\hfill\begin{tcolorbox}[enhanced,nobeforeafter,width=7.1cm,height=2.1cm,valign=center,colback=cgeo!6,colframe=cgeo,arc=2.5mm,boxrule=0.9pt,left=3mm]{\Large\bfseries\color{cgeo}d)}\quad\begin{minipage}[c]{5.6cm}\raggedright\Large un losange\end{minipage}\end{tcolorbox}
\end{tcolorbox}
\newpage
\begin{tcolorbox}[enhanced,width=16cm,height=14cm,colback=white,colframe=cmes,boxrule=2pt,arc=0mm,
left=6mm,right=6mm,top=0mm,bottom=4mm,sharp corners,
overlay={\fill[cmes] (frame.north west) rectangle ([yshift=-1.25cm]frame.north east);
\node[anchor=west,text=white,font=\sffamily\bfseries\LARGE] at ([xshift=5mm,yshift=-0.62cm]frame.north west) {Question 65 C};
\node[anchor=east,text=white,font=\sffamily\itshape\large] at ([xshift=-5mm,yshift=-0.62cm]frame.north east) {Grandeurs et mesures};}]
\vspace*{1.75cm}
\begin{minipage}[t][6.9cm][c]{\linewidth}\LARGE\centering
\begin{minipage}{0.95\linewidth}\raggedright\linespread{1.25}\selectfont Un rectangle a une longueur de $11$ cm et une largeur de $2{,}5$ cm. Quel est son périmètre~?\end{minipage}
\end{minipage}

\vfill
\begin{tcolorbox}[enhanced,nobeforeafter,width=7.1cm,height=2.1cm,valign=center,colback=cmes!6,colframe=cmes,arc=2.5mm,boxrule=0.9pt,left=3mm]{\Large\bfseries\color{cmes}a)}\quad\begin{minipage}[c]{5.6cm}\raggedright\Large $27{,}5~\text{cm}$\end{minipage}\end{tcolorbox}\hfill\begin{tcolorbox}[enhanced,nobeforeafter,width=7.1cm,height=2.1cm,valign=center,colback=cmes!6,colframe=cmes,arc=2.5mm,boxrule=0.9pt,left=3mm]{\Large\bfseries\color{cmes}b)}\quad\begin{minipage}[c]{5.6cm}\raggedright\Large $27~\text{cm}$\end{minipage}\end{tcolorbox}\par\vspace{2.5mm}
\begin{tcolorbox}[enhanced,nobeforeafter,width=7.1cm,height=2.1cm,valign=center,colback=cmes!6,colframe=cmes,arc=2.5mm,boxrule=0.9pt,left=3mm]{\Large\bfseries\color{cmes}c)}\quad\begin{minipage}[c]{5.6cm}\raggedright\Large $24{,}5~\text{cm}$\end{minipage}\end{tcolorbox}\hfill\begin{tcolorbox}[enhanced,nobeforeafter,width=7.1cm,height=2.1cm,valign=center,colback=cmes!6,colframe=cmes,arc=2.5mm,boxrule=0.9pt,left=3mm]{\Large\bfseries\color{cmes}d)}\quad\begin{minipage}[c]{5.6cm}\raggedright\Large $13{,}5~\text{cm}$\end{minipage}\end{tcolorbox}
\end{tcolorbox}
\newpage
\begin{tcolorbox}[enhanced,width=16cm,height=14cm,colback=white,colframe=cmes,boxrule=2pt,arc=0mm,
left=6mm,right=6mm,top=0mm,bottom=4mm,sharp corners,
overlay={\fill[cmes] (frame.north west) rectangle ([yshift=-1.25cm]frame.north east);
\node[anchor=west,text=white,font=\sffamily\bfseries\LARGE] at ([xshift=5mm,yshift=-0.62cm]frame.north west) {Question 66 C};
\node[anchor=east,text=white,font=\sffamily\itshape\large] at ([xshift=-5mm,yshift=-0.62cm]frame.north east) {Grandeurs et mesures};}]
\vspace*{1.75cm}
\begin{minipage}[t][6.9cm][c]{\linewidth}\LARGE\centering
\begin{minipage}{0.95\linewidth}\raggedright\linespread{1.25}\selectfont Un rectangle a une longueur de $7$ cm et une largeur de $3$ cm. Quel est son périmètre~?\end{minipage}
\end{minipage}

\vfill
\begin{tcolorbox}[enhanced,nobeforeafter,width=7.1cm,height=2.1cm,valign=center,colback=cmes!6,colframe=cmes,arc=2.5mm,boxrule=0.9pt,left=3mm]{\Large\bfseries\color{cmes}a)}\quad\begin{minipage}[c]{5.6cm}\raggedright\Large $10~\text{cm}$\end{minipage}\end{tcolorbox}\hfill\begin{tcolorbox}[enhanced,nobeforeafter,width=7.1cm,height=2.1cm,valign=center,colback=cmes!6,colframe=cmes,arc=2.5mm,boxrule=0.9pt,left=3mm]{\Large\bfseries\color{cmes}b)}\quad\begin{minipage}[c]{5.6cm}\raggedright\Large $17~\text{cm}$\end{minipage}\end{tcolorbox}\par\vspace{2.5mm}
\begin{tcolorbox}[enhanced,nobeforeafter,width=7.1cm,height=2.1cm,valign=center,colback=cmes!6,colframe=cmes,arc=2.5mm,boxrule=0.9pt,left=3mm]{\Large\bfseries\color{cmes}c)}\quad\begin{minipage}[c]{5.6cm}\raggedright\Large $20~\text{cm}$\end{minipage}\end{tcolorbox}\hfill\begin{tcolorbox}[enhanced,nobeforeafter,width=7.1cm,height=2.1cm,valign=center,colback=cmes!6,colframe=cmes,arc=2.5mm,boxrule=0.9pt,left=3mm]{\Large\bfseries\color{cmes}d)}\quad\begin{minipage}[c]{5.6cm}\raggedright\Large $21~\text{cm}$\end{minipage}\end{tcolorbox}
\end{tcolorbox}
\newpage
\begin{tcolorbox}[enhanced,width=16cm,height=14cm,colback=white,colframe=cmes,boxrule=2pt,arc=0mm,
left=6mm,right=6mm,top=0mm,bottom=4mm,sharp corners,
overlay={\fill[cmes] (frame.north west) rectangle ([yshift=-1.25cm]frame.north east);
\node[anchor=west,text=white,font=\sffamily\bfseries\LARGE] at ([xshift=5mm,yshift=-0.62cm]frame.north west) {Question 67 C};
\node[anchor=east,text=white,font=\sffamily\itshape\large] at ([xshift=-5mm,yshift=-0.62cm]frame.north east) {Grandeurs et mesures};}]
\vspace*{1.75cm}
\begin{minipage}[t][6.9cm][c]{\linewidth}\LARGE\centering
\begin{minipage}{0.95\linewidth}\raggedright\linespread{1.25}\selectfont Un rectangle a une longueur de $11$ cm et une largeur de $2{,}5$ cm. Quelle est son aire~?\end{minipage}
\end{minipage}

\vfill
\begin{tcolorbox}[enhanced,nobeforeafter,width=7.1cm,height=2.1cm,valign=center,colback=cmes!6,colframe=cmes,arc=2.5mm,boxrule=0.9pt,left=3mm]{\Large\bfseries\color{cmes}a)}\quad\begin{minipage}[c]{5.6cm}\raggedright\Large $27{,}5~\text{cm}^2$\end{minipage}\end{tcolorbox}\hfill\begin{tcolorbox}[enhanced,nobeforeafter,width=7.1cm,height=2.1cm,valign=center,colback=cmes!6,colframe=cmes,arc=2.5mm,boxrule=0.9pt,left=3mm]{\Large\bfseries\color{cmes}b)}\quad\begin{minipage}[c]{5.6cm}\raggedright\Large $27~\text{cm}^2$\end{minipage}\end{tcolorbox}\par\vspace{2.5mm}
\begin{tcolorbox}[enhanced,nobeforeafter,width=7.1cm,height=2.1cm,valign=center,colback=cmes!6,colframe=cmes,arc=2.5mm,boxrule=0.9pt,left=3mm]{\Large\bfseries\color{cmes}c)}\quad\begin{minipage}[c]{5.6cm}\raggedright\Large $13{,}75~\text{cm}^2$\end{minipage}\end{tcolorbox}\hfill\begin{tcolorbox}[enhanced,nobeforeafter,width=7.1cm,height=2.1cm,valign=center,colback=cmes!6,colframe=cmes,arc=2.5mm,boxrule=0.9pt,left=3mm]{\Large\bfseries\color{cmes}d)}\quad\begin{minipage}[c]{5.6cm}\raggedright\Large $13{,}5~\text{cm}^2$\end{minipage}\end{tcolorbox}
\end{tcolorbox}
\newpage
\begin{tcolorbox}[enhanced,width=16cm,height=14cm,colback=white,colframe=cmes,boxrule=2pt,arc=0mm,
left=6mm,right=6mm,top=0mm,bottom=4mm,sharp corners,
overlay={\fill[cmes] (frame.north west) rectangle ([yshift=-1.25cm]frame.north east);
\node[anchor=west,text=white,font=\sffamily\bfseries\LARGE] at ([xshift=5mm,yshift=-0.62cm]frame.north west) {Question 68 C};
\node[anchor=east,text=white,font=\sffamily\itshape\large] at ([xshift=-5mm,yshift=-0.62cm]frame.north east) {Grandeurs et mesures};}]
\vspace*{1.75cm}
\begin{minipage}[t][6.9cm][c]{\linewidth}\LARGE\centering
\begin{minipage}{0.95\linewidth}\raggedright\linespread{1.25}\selectfont Un rectangle a une longueur de $7$ cm et une largeur de $3$ cm. Quelle est son aire~?\end{minipage}
\end{minipage}

\vfill
\begin{tcolorbox}[enhanced,nobeforeafter,width=7.1cm,height=2.1cm,valign=center,colback=cmes!6,colframe=cmes,arc=2.5mm,boxrule=0.9pt,left=3mm]{\Large\bfseries\color{cmes}a)}\quad\begin{minipage}[c]{5.6cm}\raggedright\Large $10{,}5~\text{cm}^2$\end{minipage}\end{tcolorbox}\hfill\begin{tcolorbox}[enhanced,nobeforeafter,width=7.1cm,height=2.1cm,valign=center,colback=cmes!6,colframe=cmes,arc=2.5mm,boxrule=0.9pt,left=3mm]{\Large\bfseries\color{cmes}b)}\quad\begin{minipage}[c]{5.6cm}\raggedright\Large $10~\text{cm}^2$\end{minipage}\end{tcolorbox}\par\vspace{2.5mm}
\begin{tcolorbox}[enhanced,nobeforeafter,width=7.1cm,height=2.1cm,valign=center,colback=cmes!6,colframe=cmes,arc=2.5mm,boxrule=0.9pt,left=3mm]{\Large\bfseries\color{cmes}c)}\quad\begin{minipage}[c]{5.6cm}\raggedright\Large $21~\text{cm}^2$\end{minipage}\end{tcolorbox}\hfill\begin{tcolorbox}[enhanced,nobeforeafter,width=7.1cm,height=2.1cm,valign=center,colback=cmes!6,colframe=cmes,arc=2.5mm,boxrule=0.9pt,left=3mm]{\Large\bfseries\color{cmes}d)}\quad\begin{minipage}[c]{5.6cm}\raggedright\Large $20~\text{cm}^2$\end{minipage}\end{tcolorbox}
\end{tcolorbox}
\newpage
\begin{tcolorbox}[enhanced,width=16cm,height=14cm,colback=white,colframe=cmes,boxrule=2pt,arc=0mm,
left=6mm,right=6mm,top=0mm,bottom=4mm,sharp corners,
overlay={\fill[cmes] (frame.north west) rectangle ([yshift=-1.25cm]frame.north east);
\node[anchor=west,text=white,font=\sffamily\bfseries\LARGE] at ([xshift=5mm,yshift=-0.62cm]frame.north west) {Question 69 C};
\node[anchor=east,text=white,font=\sffamily\itshape\large] at ([xshift=-5mm,yshift=-0.62cm]frame.north east) {Grandeurs et mesures};}]
\vspace*{1.75cm}
\begin{minipage}[t][6.9cm][c]{\linewidth}\LARGE\centering
\begin{minipage}{0.95\linewidth}\raggedright\linespread{1.25}\selectfont Un triangle est rectangle ; les deux côtés de l'angle droit mesurent $5$ cm et $3$ cm. Quelle est son aire~?\end{minipage}
\end{minipage}

\vfill
\begin{tcolorbox}[enhanced,nobeforeafter,width=7.1cm,height=2.1cm,valign=center,colback=cmes!6,colframe=cmes,arc=2.5mm,boxrule=0.9pt,left=3mm]{\Large\bfseries\color{cmes}a)}\quad\begin{minipage}[c]{5.6cm}\raggedright\Large $7{,}5~\text{cm}^2$\end{minipage}\end{tcolorbox}\hfill\begin{tcolorbox}[enhanced,nobeforeafter,width=7.1cm,height=2.1cm,valign=center,colback=cmes!6,colframe=cmes,arc=2.5mm,boxrule=0.9pt,left=3mm]{\Large\bfseries\color{cmes}b)}\quad\begin{minipage}[c]{5.6cm}\raggedright\Large $15~\text{cm}^2$\end{minipage}\end{tcolorbox}\par\vspace{2.5mm}
\begin{tcolorbox}[enhanced,nobeforeafter,width=7.1cm,height=2.1cm,valign=center,colback=cmes!6,colframe=cmes,arc=2.5mm,boxrule=0.9pt,left=3mm]{\Large\bfseries\color{cmes}c)}\quad\begin{minipage}[c]{5.6cm}\raggedright\Large $8~\text{cm}^2$\end{minipage}\end{tcolorbox}\hfill\begin{tcolorbox}[enhanced,nobeforeafter,width=7.1cm,height=2.1cm,valign=center,colback=cmes!6,colframe=cmes,arc=2.5mm,boxrule=0.9pt,left=3mm]{\Large\bfseries\color{cmes}d)}\quad\begin{minipage}[c]{5.6cm}\raggedright\Large $16~\text{cm}^2$\end{minipage}\end{tcolorbox}
\end{tcolorbox}
\newpage
\begin{tcolorbox}[enhanced,width=16cm,height=14cm,colback=white,colframe=cmes,boxrule=2pt,arc=0mm,
left=6mm,right=6mm,top=0mm,bottom=4mm,sharp corners,
overlay={\fill[cmes] (frame.north west) rectangle ([yshift=-1.25cm]frame.north east);
\node[anchor=west,text=white,font=\sffamily\bfseries\LARGE] at ([xshift=5mm,yshift=-0.62cm]frame.north west) {Question 70 C};
\node[anchor=east,text=white,font=\sffamily\itshape\large] at ([xshift=-5mm,yshift=-0.62cm]frame.north east) {Grandeurs et mesures};}]
\vspace*{1.75cm}
\begin{minipage}[t][6.9cm][c]{\linewidth}\LARGE\centering
\begin{minipage}{0.95\linewidth}\raggedright\linespread{1.25}\selectfont Un triangle est rectangle ; les deux côtés de l'angle droit mesurent $14$ cm et $6$ cm. Quelle est son aire~?\end{minipage}
\end{minipage}

\vfill
\begin{tcolorbox}[enhanced,nobeforeafter,width=7.1cm,height=2.1cm,valign=center,colback=cmes!6,colframe=cmes,arc=2.5mm,boxrule=0.9pt,left=3mm]{\Large\bfseries\color{cmes}a)}\quad\begin{minipage}[c]{5.6cm}\raggedright\Large $40~\text{cm}^2$\end{minipage}\end{tcolorbox}\hfill\begin{tcolorbox}[enhanced,nobeforeafter,width=7.1cm,height=2.1cm,valign=center,colback=cmes!6,colframe=cmes,arc=2.5mm,boxrule=0.9pt,left=3mm]{\Large\bfseries\color{cmes}b)}\quad\begin{minipage}[c]{5.6cm}\raggedright\Large $20~\text{cm}^2$\end{minipage}\end{tcolorbox}\par\vspace{2.5mm}
\begin{tcolorbox}[enhanced,nobeforeafter,width=7.1cm,height=2.1cm,valign=center,colback=cmes!6,colframe=cmes,arc=2.5mm,boxrule=0.9pt,left=3mm]{\Large\bfseries\color{cmes}c)}\quad\begin{minipage}[c]{5.6cm}\raggedright\Large $84~\text{cm}^2$\end{minipage}\end{tcolorbox}\hfill\begin{tcolorbox}[enhanced,nobeforeafter,width=7.1cm,height=2.1cm,valign=center,colback=cmes!6,colframe=cmes,arc=2.5mm,boxrule=0.9pt,left=3mm]{\Large\bfseries\color{cmes}d)}\quad\begin{minipage}[c]{5.6cm}\raggedright\Large $42~\text{cm}^2$\end{minipage}\end{tcolorbox}
\end{tcolorbox}
\newpage
\begin{tcolorbox}[enhanced,width=16cm,height=14cm,colback=white,colframe=cmes,boxrule=2pt,arc=0mm,
left=6mm,right=6mm,top=0mm,bottom=4mm,sharp corners,
overlay={\fill[cmes] (frame.north west) rectangle ([yshift=-1.25cm]frame.north east);
\node[anchor=west,text=white,font=\sffamily\bfseries\LARGE] at ([xshift=5mm,yshift=-0.62cm]frame.north west) {Question 71 C};
\node[anchor=east,text=white,font=\sffamily\itshape\large] at ([xshift=-5mm,yshift=-0.62cm]frame.north east) {Grandeurs et mesures};}]
\vspace*{1.75cm}
\begin{minipage}[t][6.9cm][c]{\linewidth}\LARGE\centering
\begin{minipage}{0.95\linewidth}\raggedright\linespread{1.25}\selectfont Compléter~: $65~\text{cm}=\;?\;\text{m}$\end{minipage}
\end{minipage}

\vfill
\begin{tcolorbox}[enhanced,nobeforeafter,width=7.1cm,height=2.1cm,valign=center,colback=cmes!6,colframe=cmes,arc=2.5mm,boxrule=0.9pt,left=3mm]{\Large\bfseries\color{cmes}a)}\quad\begin{minipage}[c]{5.6cm}\raggedright\Large $6{,}5~\text{m}$\end{minipage}\end{tcolorbox}\hfill\begin{tcolorbox}[enhanced,nobeforeafter,width=7.1cm,height=2.1cm,valign=center,colback=cmes!6,colframe=cmes,arc=2.5mm,boxrule=0.9pt,left=3mm]{\Large\bfseries\color{cmes}b)}\quad\begin{minipage}[c]{5.6cm}\raggedright\Large $0{,}0065~\text{m}$\end{minipage}\end{tcolorbox}\par\vspace{2.5mm}
\begin{tcolorbox}[enhanced,nobeforeafter,width=7.1cm,height=2.1cm,valign=center,colback=cmes!6,colframe=cmes,arc=2.5mm,boxrule=0.9pt,left=3mm]{\Large\bfseries\color{cmes}c)}\quad\begin{minipage}[c]{5.6cm}\raggedright\Large $0{,}65~\text{m}$\end{minipage}\end{tcolorbox}\hfill\begin{tcolorbox}[enhanced,nobeforeafter,width=7.1cm,height=2.1cm,valign=center,colback=cmes!6,colframe=cmes,arc=2.5mm,boxrule=0.9pt,left=3mm]{\Large\bfseries\color{cmes}d)}\quad\begin{minipage}[c]{5.6cm}\raggedright\Large $0{,}065~\text{m}$\end{minipage}\end{tcolorbox}
\end{tcolorbox}
\newpage
\begin{tcolorbox}[enhanced,width=16cm,height=14cm,colback=white,colframe=cmes,boxrule=2pt,arc=0mm,
left=6mm,right=6mm,top=0mm,bottom=4mm,sharp corners,
overlay={\fill[cmes] (frame.north west) rectangle ([yshift=-1.25cm]frame.north east);
\node[anchor=west,text=white,font=\sffamily\bfseries\LARGE] at ([xshift=5mm,yshift=-0.62cm]frame.north west) {Question 72 C};
\node[anchor=east,text=white,font=\sffamily\itshape\large] at ([xshift=-5mm,yshift=-0.62cm]frame.north east) {Grandeurs et mesures};}]
\vspace*{1.75cm}
\begin{minipage}[t][6.9cm][c]{\linewidth}\LARGE\centering
\begin{minipage}{0.95\linewidth}\raggedright\linespread{1.25}\selectfont Compléter~: $7~\text{dm}=\;?\;\text{mm}$\end{minipage}
\end{minipage}

\vfill
\begin{tcolorbox}[enhanced,nobeforeafter,width=7.1cm,height=2.1cm,valign=center,colback=cmes!6,colframe=cmes,arc=2.5mm,boxrule=0.9pt,left=3mm]{\Large\bfseries\color{cmes}a)}\quad\begin{minipage}[c]{5.6cm}\raggedright\Large $70\,000~\text{mm}$\end{minipage}\end{tcolorbox}\hfill\begin{tcolorbox}[enhanced,nobeforeafter,width=7.1cm,height=2.1cm,valign=center,colback=cmes!6,colframe=cmes,arc=2.5mm,boxrule=0.9pt,left=3mm]{\Large\bfseries\color{cmes}b)}\quad\begin{minipage}[c]{5.6cm}\raggedright\Large $700~\text{mm}$\end{minipage}\end{tcolorbox}\par\vspace{2.5mm}
\begin{tcolorbox}[enhanced,nobeforeafter,width=7.1cm,height=2.1cm,valign=center,colback=cmes!6,colframe=cmes,arc=2.5mm,boxrule=0.9pt,left=3mm]{\Large\bfseries\color{cmes}c)}\quad\begin{minipage}[c]{5.6cm}\raggedright\Large $7\,000~\text{mm}$\end{minipage}\end{tcolorbox}\hfill\begin{tcolorbox}[enhanced,nobeforeafter,width=7.1cm,height=2.1cm,valign=center,colback=cmes!6,colframe=cmes,arc=2.5mm,boxrule=0.9pt,left=3mm]{\Large\bfseries\color{cmes}d)}\quad\begin{minipage}[c]{5.6cm}\raggedright\Large $70~\text{mm}$\end{minipage}\end{tcolorbox}
\end{tcolorbox}
\newpage
\begin{tcolorbox}[enhanced,width=16cm,height=14cm,colback=white,colframe=cmes,boxrule=2pt,arc=0mm,
left=6mm,right=6mm,top=0mm,bottom=4mm,sharp corners,
overlay={\fill[cmes] (frame.north west) rectangle ([yshift=-1.25cm]frame.north east);
\node[anchor=west,text=white,font=\sffamily\bfseries\LARGE] at ([xshift=5mm,yshift=-0.62cm]frame.north west) {Question 73 C};
\node[anchor=east,text=white,font=\sffamily\itshape\large] at ([xshift=-5mm,yshift=-0.62cm]frame.north east) {Grandeurs et mesures};}]
\vspace*{1.75cm}
\begin{minipage}[t][6.9cm][c]{\linewidth}\LARGE\centering
\begin{minipage}{0.95\linewidth}\raggedright\linespread{1.25}\selectfont Compléter~: $0{,}3~\text{kg}=\;?\;\text{g}$\end{minipage}
\end{minipage}

\vfill
\begin{tcolorbox}[enhanced,nobeforeafter,width=7.1cm,height=2.1cm,valign=center,colback=cmes!6,colframe=cmes,arc=2.5mm,boxrule=0.9pt,left=3mm]{\Large\bfseries\color{cmes}a)}\quad\begin{minipage}[c]{5.6cm}\raggedright\Large $30~\text{g}$\end{minipage}\end{tcolorbox}\hfill\begin{tcolorbox}[enhanced,nobeforeafter,width=7.1cm,height=2.1cm,valign=center,colback=cmes!6,colframe=cmes,arc=2.5mm,boxrule=0.9pt,left=3mm]{\Large\bfseries\color{cmes}b)}\quad\begin{minipage}[c]{5.6cm}\raggedright\Large $300~\text{g}$\end{minipage}\end{tcolorbox}\par\vspace{2.5mm}
\begin{tcolorbox}[enhanced,nobeforeafter,width=7.1cm,height=2.1cm,valign=center,colback=cmes!6,colframe=cmes,arc=2.5mm,boxrule=0.9pt,left=3mm]{\Large\bfseries\color{cmes}c)}\quad\begin{minipage}[c]{5.6cm}\raggedright\Large $30\,000~\text{g}$\end{minipage}\end{tcolorbox}\hfill\begin{tcolorbox}[enhanced,nobeforeafter,width=7.1cm,height=2.1cm,valign=center,colback=cmes!6,colframe=cmes,arc=2.5mm,boxrule=0.9pt,left=3mm]{\Large\bfseries\color{cmes}d)}\quad\begin{minipage}[c]{5.6cm}\raggedright\Large $3\,000~\text{g}$\end{minipage}\end{tcolorbox}
\end{tcolorbox}
\newpage
\begin{tcolorbox}[enhanced,width=16cm,height=14cm,colback=white,colframe=cmes,boxrule=2pt,arc=0mm,
left=6mm,right=6mm,top=0mm,bottom=4mm,sharp corners,
overlay={\fill[cmes] (frame.north west) rectangle ([yshift=-1.25cm]frame.north east);
\node[anchor=west,text=white,font=\sffamily\bfseries\LARGE] at ([xshift=5mm,yshift=-0.62cm]frame.north west) {Question 74 C};
\node[anchor=east,text=white,font=\sffamily\itshape\large] at ([xshift=-5mm,yshift=-0.62cm]frame.north east) {Grandeurs et mesures};}]
\vspace*{1.75cm}
\begin{minipage}[t][6.9cm][c]{\linewidth}\LARGE\centering
\begin{minipage}{0.95\linewidth}\raggedright\linespread{1.25}\selectfont Compléter~: $8~\text{g}=\;?\;\text{cg}$\end{minipage}
\end{minipage}

\vfill
\begin{tcolorbox}[enhanced,nobeforeafter,width=7.1cm,height=2.1cm,valign=center,colback=cmes!6,colframe=cmes,arc=2.5mm,boxrule=0.9pt,left=3mm]{\Large\bfseries\color{cmes}a)}\quad\begin{minipage}[c]{5.6cm}\raggedright\Large $800~\text{cg}$\end{minipage}\end{tcolorbox}\hfill\begin{tcolorbox}[enhanced,nobeforeafter,width=7.1cm,height=2.1cm,valign=center,colback=cmes!6,colframe=cmes,arc=2.5mm,boxrule=0.9pt,left=3mm]{\Large\bfseries\color{cmes}b)}\quad\begin{minipage}[c]{5.6cm}\raggedright\Large $80\,000~\text{cg}$\end{minipage}\end{tcolorbox}\par\vspace{2.5mm}
\begin{tcolorbox}[enhanced,nobeforeafter,width=7.1cm,height=2.1cm,valign=center,colback=cmes!6,colframe=cmes,arc=2.5mm,boxrule=0.9pt,left=3mm]{\Large\bfseries\color{cmes}c)}\quad\begin{minipage}[c]{5.6cm}\raggedright\Large $80~\text{cg}$\end{minipage}\end{tcolorbox}\hfill\begin{tcolorbox}[enhanced,nobeforeafter,width=7.1cm,height=2.1cm,valign=center,colback=cmes!6,colframe=cmes,arc=2.5mm,boxrule=0.9pt,left=3mm]{\Large\bfseries\color{cmes}d)}\quad\begin{minipage}[c]{5.6cm}\raggedright\Large $8\,000~\text{cg}$\end{minipage}\end{tcolorbox}
\end{tcolorbox}
\newpage
\begin{tcolorbox}[enhanced,width=16cm,height=14cm,colback=white,colframe=cmes,boxrule=2pt,arc=0mm,
left=6mm,right=6mm,top=0mm,bottom=4mm,sharp corners,
overlay={\fill[cmes] (frame.north west) rectangle ([yshift=-1.25cm]frame.north east);
\node[anchor=west,text=white,font=\sffamily\bfseries\LARGE] at ([xshift=5mm,yshift=-0.62cm]frame.north west) {Question 75 C};
\node[anchor=east,text=white,font=\sffamily\itshape\large] at ([xshift=-5mm,yshift=-0.62cm]frame.north east) {Grandeurs et mesures};}]
\vspace*{1.75cm}
\begin{minipage}[t][6.9cm][c]{\linewidth}\LARGE\centering
\begin{minipage}{0.95\linewidth}\raggedright\linespread{1.25}\selectfont Combien de minutes y a-t-il dans 4 h 20 min~?\end{minipage}
\end{minipage}

\vfill
\begin{tcolorbox}[enhanced,nobeforeafter,width=7.1cm,height=2.1cm,valign=center,colback=cmes!6,colframe=cmes,arc=2.5mm,boxrule=0.9pt,left=3mm]{\Large\bfseries\color{cmes}a)}\quad\begin{minipage}[c]{5.6cm}\raggedright\Large $260~\text{min}$\end{minipage}\end{tcolorbox}\hfill\begin{tcolorbox}[enhanced,nobeforeafter,width=7.1cm,height=2.1cm,valign=center,colback=cmes!6,colframe=cmes,arc=2.5mm,boxrule=0.9pt,left=3mm]{\Large\bfseries\color{cmes}b)}\quad\begin{minipage}[c]{5.6cm}\raggedright\Large $220~\text{min}$\end{minipage}\end{tcolorbox}\par\vspace{2.5mm}
\begin{tcolorbox}[enhanced,nobeforeafter,width=7.1cm,height=2.1cm,valign=center,colback=cmes!6,colframe=cmes,arc=2.5mm,boxrule=0.9pt,left=3mm]{\Large\bfseries\color{cmes}c)}\quad\begin{minipage}[c]{5.6cm}\raggedright\Large $24~\text{min}$\end{minipage}\end{tcolorbox}\hfill\begin{tcolorbox}[enhanced,nobeforeafter,width=7.1cm,height=2.1cm,valign=center,colback=cmes!6,colframe=cmes,arc=2.5mm,boxrule=0.9pt,left=3mm]{\Large\bfseries\color{cmes}d)}\quad\begin{minipage}[c]{5.6cm}\raggedright\Large $420~\text{min}$\end{minipage}\end{tcolorbox}
\end{tcolorbox}
\newpage
\begin{tcolorbox}[enhanced,width=16cm,height=14cm,colback=white,colframe=cmes,boxrule=2pt,arc=0mm,
left=6mm,right=6mm,top=0mm,bottom=4mm,sharp corners,
overlay={\fill[cmes] (frame.north west) rectangle ([yshift=-1.25cm]frame.north east);
\node[anchor=west,text=white,font=\sffamily\bfseries\LARGE] at ([xshift=5mm,yshift=-0.62cm]frame.north west) {Question 76 C};
\node[anchor=east,text=white,font=\sffamily\itshape\large] at ([xshift=-5mm,yshift=-0.62cm]frame.north east) {Grandeurs et mesures};}]
\vspace*{1.75cm}
\begin{minipage}[t][6.9cm][c]{\linewidth}\LARGE\centering
\begin{minipage}{0.95\linewidth}\raggedright\linespread{1.25}\selectfont Combien de minutes y a-t-il dans 2 h 50 min~?\end{minipage}
\end{minipage}

\vfill
\begin{tcolorbox}[enhanced,nobeforeafter,width=7.1cm,height=2.1cm,valign=center,colback=cmes!6,colframe=cmes,arc=2.5mm,boxrule=0.9pt,left=3mm]{\Large\bfseries\color{cmes}a)}\quad\begin{minipage}[c]{5.6cm}\raggedright\Large $250~\text{min}$\end{minipage}\end{tcolorbox}\hfill\begin{tcolorbox}[enhanced,nobeforeafter,width=7.1cm,height=2.1cm,valign=center,colback=cmes!6,colframe=cmes,arc=2.5mm,boxrule=0.9pt,left=3mm]{\Large\bfseries\color{cmes}b)}\quad\begin{minipage}[c]{5.6cm}\raggedright\Large $70~\text{min}$\end{minipage}\end{tcolorbox}\par\vspace{2.5mm}
\begin{tcolorbox}[enhanced,nobeforeafter,width=7.1cm,height=2.1cm,valign=center,colback=cmes!6,colframe=cmes,arc=2.5mm,boxrule=0.9pt,left=3mm]{\Large\bfseries\color{cmes}c)}\quad\begin{minipage}[c]{5.6cm}\raggedright\Large $52~\text{min}$\end{minipage}\end{tcolorbox}\hfill\begin{tcolorbox}[enhanced,nobeforeafter,width=7.1cm,height=2.1cm,valign=center,colback=cmes!6,colframe=cmes,arc=2.5mm,boxrule=0.9pt,left=3mm]{\Large\bfseries\color{cmes}d)}\quad\begin{minipage}[c]{5.6cm}\raggedright\Large $170~\text{min}$\end{minipage}\end{tcolorbox}
\end{tcolorbox}
\newpage
\begin{tcolorbox}[enhanced,width=16cm,height=14cm,colback=white,colframe=cmes,boxrule=2pt,arc=0mm,
left=6mm,right=6mm,top=0mm,bottom=4mm,sharp corners,
overlay={\fill[cmes] (frame.north west) rectangle ([yshift=-1.25cm]frame.north east);
\node[anchor=west,text=white,font=\sffamily\bfseries\LARGE] at ([xshift=5mm,yshift=-0.62cm]frame.north west) {Question 77 C};
\node[anchor=east,text=white,font=\sffamily\itshape\large] at ([xshift=-5mm,yshift=-0.62cm]frame.north east) {Grandeurs et mesures};}]
\vspace*{1.75cm}
\begin{minipage}[t][6.9cm][c]{\linewidth}\LARGE\centering
\begin{minipage}{0.95\linewidth}\raggedright\linespread{1.25}\selectfont En prenant $\pi\approx3{,}14$, quel est le périmètre d'un cercle de diamètre $8$ cm~?\end{minipage}
\end{minipage}

\vfill
\begin{tcolorbox}[enhanced,nobeforeafter,width=7.1cm,height=2.1cm,valign=center,colback=cmes!6,colframe=cmes,arc=2.5mm,boxrule=0.9pt,left=3mm]{\Large\bfseries\color{cmes}a)}\quad\begin{minipage}[c]{5.6cm}\raggedright\Large $25{,}12~\text{cm}$\end{minipage}\end{tcolorbox}\hfill\begin{tcolorbox}[enhanced,nobeforeafter,width=7.1cm,height=2.1cm,valign=center,colback=cmes!6,colframe=cmes,arc=2.5mm,boxrule=0.9pt,left=3mm]{\Large\bfseries\color{cmes}b)}\quad\begin{minipage}[c]{5.6cm}\raggedright\Large $24~\text{cm}$\end{minipage}\end{tcolorbox}\par\vspace{2.5mm}
\begin{tcolorbox}[enhanced,nobeforeafter,width=7.1cm,height=2.1cm,valign=center,colback=cmes!6,colframe=cmes,arc=2.5mm,boxrule=0.9pt,left=3mm]{\Large\bfseries\color{cmes}c)}\quad\begin{minipage}[c]{5.6cm}\raggedright\Large $12{,}56~\text{cm}$\end{minipage}\end{tcolorbox}\hfill\begin{tcolorbox}[enhanced,nobeforeafter,width=7.1cm,height=2.1cm,valign=center,colback=cmes!6,colframe=cmes,arc=2.5mm,boxrule=0.9pt,left=3mm]{\Large\bfseries\color{cmes}d)}\quad\begin{minipage}[c]{5.6cm}\raggedright\Large $50{,}24~\text{cm}$\end{minipage}\end{tcolorbox}
\end{tcolorbox}
\newpage
\begin{tcolorbox}[enhanced,width=16cm,height=14cm,colback=white,colframe=cmes,boxrule=2pt,arc=0mm,
left=6mm,right=6mm,top=0mm,bottom=4mm,sharp corners,
overlay={\fill[cmes] (frame.north west) rectangle ([yshift=-1.25cm]frame.north east);
\node[anchor=west,text=white,font=\sffamily\bfseries\LARGE] at ([xshift=5mm,yshift=-0.62cm]frame.north west) {Question 78 C};
\node[anchor=east,text=white,font=\sffamily\itshape\large] at ([xshift=-5mm,yshift=-0.62cm]frame.north east) {Grandeurs et mesures};}]
\vspace*{1.75cm}
\begin{minipage}[t][6.9cm][c]{\linewidth}\LARGE\centering
\begin{minipage}{0.95\linewidth}\raggedright\linespread{1.25}\selectfont En prenant $\pi\approx3{,}14$, quel est le périmètre d'un cercle de diamètre $15$ cm~?\end{minipage}
\end{minipage}

\vfill
\begin{tcolorbox}[enhanced,nobeforeafter,width=7.1cm,height=2.1cm,valign=center,colback=cmes!6,colframe=cmes,arc=2.5mm,boxrule=0.9pt,left=3mm]{\Large\bfseries\color{cmes}a)}\quad\begin{minipage}[c]{5.6cm}\raggedright\Large $23{,}55~\text{cm}$\end{minipage}\end{tcolorbox}\hfill\begin{tcolorbox}[enhanced,nobeforeafter,width=7.1cm,height=2.1cm,valign=center,colback=cmes!6,colframe=cmes,arc=2.5mm,boxrule=0.9pt,left=3mm]{\Large\bfseries\color{cmes}b)}\quad\begin{minipage}[c]{5.6cm}\raggedright\Large $94{,}2~\text{cm}$\end{minipage}\end{tcolorbox}\par\vspace{2.5mm}
\begin{tcolorbox}[enhanced,nobeforeafter,width=7.1cm,height=2.1cm,valign=center,colback=cmes!6,colframe=cmes,arc=2.5mm,boxrule=0.9pt,left=3mm]{\Large\bfseries\color{cmes}c)}\quad\begin{minipage}[c]{5.6cm}\raggedright\Large $176{,}625~\text{cm}$\end{minipage}\end{tcolorbox}\hfill\begin{tcolorbox}[enhanced,nobeforeafter,width=7.1cm,height=2.1cm,valign=center,colback=cmes!6,colframe=cmes,arc=2.5mm,boxrule=0.9pt,left=3mm]{\Large\bfseries\color{cmes}d)}\quad\begin{minipage}[c]{5.6cm}\raggedright\Large $47{,}1~\text{cm}$\end{minipage}\end{tcolorbox}
\end{tcolorbox}
\newpage
\begin{tcolorbox}[enhanced,width=16cm,height=14cm,colback=white,colframe=cmes,boxrule=2pt,arc=0mm,
left=6mm,right=6mm,top=0mm,bottom=4mm,sharp corners,
overlay={\fill[cmes] (frame.north west) rectangle ([yshift=-1.25cm]frame.north east);
\node[anchor=west,text=white,font=\sffamily\bfseries\LARGE] at ([xshift=5mm,yshift=-0.62cm]frame.north west) {Question 79 C};
\node[anchor=east,text=white,font=\sffamily\itshape\large] at ([xshift=-5mm,yshift=-0.62cm]frame.north east) {Grandeurs et mesures};}]
\vspace*{1.75cm}
\begin{minipage}[t][6.9cm][c]{\linewidth}\LARGE\centering
\begin{minipage}{0.95\linewidth}\raggedright\linespread{1.25}\selectfont Un pavé droit mesure $2{,}5$ cm de longueur, $4$ cm de largeur et $10$ cm de hauteur. Quel est son volume~?\end{minipage}
\end{minipage}

\vfill
\begin{tcolorbox}[enhanced,nobeforeafter,width=7.1cm,height=2.1cm,valign=center,colback=cmes!6,colframe=cmes,arc=2.5mm,boxrule=0.9pt,left=3mm]{\Large\bfseries\color{cmes}a)}\quad\begin{minipage}[c]{5.6cm}\raggedright\Large $150~\text{cm}^3$\end{minipage}\end{tcolorbox}\hfill\begin{tcolorbox}[enhanced,nobeforeafter,width=7.1cm,height=2.1cm,valign=center,colback=cmes!6,colframe=cmes,arc=2.5mm,boxrule=0.9pt,left=3mm]{\Large\bfseries\color{cmes}b)}\quad\begin{minipage}[c]{5.6cm}\raggedright\Large $10~\text{cm}^3$\end{minipage}\end{tcolorbox}\par\vspace{2.5mm}
\begin{tcolorbox}[enhanced,nobeforeafter,width=7.1cm,height=2.1cm,valign=center,colback=cmes!6,colframe=cmes,arc=2.5mm,boxrule=0.9pt,left=3mm]{\Large\bfseries\color{cmes}c)}\quad\begin{minipage}[c]{5.6cm}\raggedright\Large $100~\text{cm}^3$\end{minipage}\end{tcolorbox}\hfill\begin{tcolorbox}[enhanced,nobeforeafter,width=7.1cm,height=2.1cm,valign=center,colback=cmes!6,colframe=cmes,arc=2.5mm,boxrule=0.9pt,left=3mm]{\Large\bfseries\color{cmes}d)}\quad\begin{minipage}[c]{5.6cm}\raggedright\Large $16{,}5~\text{cm}^3$\end{minipage}\end{tcolorbox}
\end{tcolorbox}
\newpage
\begin{tcolorbox}[enhanced,width=16cm,height=14cm,colback=white,colframe=cmes,boxrule=2pt,arc=0mm,
left=6mm,right=6mm,top=0mm,bottom=4mm,sharp corners,
overlay={\fill[cmes] (frame.north west) rectangle ([yshift=-1.25cm]frame.north east);
\node[anchor=west,text=white,font=\sffamily\bfseries\LARGE] at ([xshift=5mm,yshift=-0.62cm]frame.north west) {Question 80 C};
\node[anchor=east,text=white,font=\sffamily\itshape\large] at ([xshift=-5mm,yshift=-0.62cm]frame.north east) {Grandeurs et mesures};}]
\vspace*{1.75cm}
\begin{minipage}[t][6.9cm][c]{\linewidth}\LARGE\centering
\begin{minipage}{0.95\linewidth}\raggedright\linespread{1.25}\selectfont Un pavé droit mesure $6$ cm de longueur, $6$ cm de largeur et $6$ cm de hauteur. Quel est son volume~?\end{minipage}
\end{minipage}

\vfill
\begin{tcolorbox}[enhanced,nobeforeafter,width=7.1cm,height=2.1cm,valign=center,colback=cmes!6,colframe=cmes,arc=2.5mm,boxrule=0.9pt,left=3mm]{\Large\bfseries\color{cmes}a)}\quad\begin{minipage}[c]{5.6cm}\raggedright\Large $36~\text{cm}^3$\end{minipage}\end{tcolorbox}\hfill\begin{tcolorbox}[enhanced,nobeforeafter,width=7.1cm,height=2.1cm,valign=center,colback=cmes!6,colframe=cmes,arc=2.5mm,boxrule=0.9pt,left=3mm]{\Large\bfseries\color{cmes}b)}\quad\begin{minipage}[c]{5.6cm}\raggedright\Large $108~\text{cm}^3$\end{minipage}\end{tcolorbox}\par\vspace{2.5mm}
\begin{tcolorbox}[enhanced,nobeforeafter,width=7.1cm,height=2.1cm,valign=center,colback=cmes!6,colframe=cmes,arc=2.5mm,boxrule=0.9pt,left=3mm]{\Large\bfseries\color{cmes}c)}\quad\begin{minipage}[c]{5.6cm}\raggedright\Large $18~\text{cm}^3$\end{minipage}\end{tcolorbox}\hfill\begin{tcolorbox}[enhanced,nobeforeafter,width=7.1cm,height=2.1cm,valign=center,colback=cmes!6,colframe=cmes,arc=2.5mm,boxrule=0.9pt,left=3mm]{\Large\bfseries\color{cmes}d)}\quad\begin{minipage}[c]{5.6cm}\raggedright\Large $216~\text{cm}^3$\end{minipage}\end{tcolorbox}
\end{tcolorbox}
\newpage
\begin{tcolorbox}[enhanced,width=16cm,height=14cm,colback=white,colframe=cpro,boxrule=2pt,arc=0mm,
left=6mm,right=6mm,top=0mm,bottom=4mm,sharp corners,
overlay={\fill[cpro] (frame.north west) rectangle ([yshift=-1.25cm]frame.north east);
\node[anchor=west,text=white,font=\sffamily\bfseries\LARGE] at ([xshift=5mm,yshift=-0.62cm]frame.north west) {Question 81 C};
\node[anchor=east,text=white,font=\sffamily\itshape\large] at ([xshift=-5mm,yshift=-0.62cm]frame.north east) {Proportionnalité et données};}]
\vspace*{1.75cm}
\begin{minipage}[t][6.9cm][c]{\linewidth}\LARGE\centering
\begin{minipage}{0.95\linewidth}\raggedright\linespread{1.25}\selectfont Pour $2$ personnes, il faut $150$ g de farine. Combien en faut-il pour $7$ personnes~?\end{minipage}
\end{minipage}

\vfill
\begin{tcolorbox}[enhanced,nobeforeafter,width=7.1cm,height=2.1cm,valign=center,colback=cpro!6,colframe=cpro,arc=2.5mm,boxrule=0.9pt,left=3mm]{\Large\bfseries\color{cpro}a)}\quad\begin{minipage}[c]{5.6cm}\raggedright\Large $155~\text{g}$\end{minipage}\end{tcolorbox}\hfill\begin{tcolorbox}[enhanced,nobeforeafter,width=7.1cm,height=2.1cm,valign=center,colback=cpro!6,colframe=cpro,arc=2.5mm,boxrule=0.9pt,left=3mm]{\Large\bfseries\color{cpro}b)}\quad\begin{minipage}[c]{5.6cm}\raggedright\Large $1\,050~\text{g}$\end{minipage}\end{tcolorbox}\par\vspace{2.5mm}
\begin{tcolorbox}[enhanced,nobeforeafter,width=7.1cm,height=2.1cm,valign=center,colback=cpro!6,colframe=cpro,arc=2.5mm,boxrule=0.9pt,left=3mm]{\Large\bfseries\color{cpro}c)}\quad\begin{minipage}[c]{5.6cm}\raggedright\Large $75~\text{g}$\end{minipage}\end{tcolorbox}\hfill\begin{tcolorbox}[enhanced,nobeforeafter,width=7.1cm,height=2.1cm,valign=center,colback=cpro!6,colframe=cpro,arc=2.5mm,boxrule=0.9pt,left=3mm]{\Large\bfseries\color{cpro}d)}\quad\begin{minipage}[c]{5.6cm}\raggedright\Large $525~\text{g}$\end{minipage}\end{tcolorbox}
\end{tcolorbox}
\newpage
\begin{tcolorbox}[enhanced,width=16cm,height=14cm,colback=white,colframe=cpro,boxrule=2pt,arc=0mm,
left=6mm,right=6mm,top=0mm,bottom=4mm,sharp corners,
overlay={\fill[cpro] (frame.north west) rectangle ([yshift=-1.25cm]frame.north east);
\node[anchor=west,text=white,font=\sffamily\bfseries\LARGE] at ([xshift=5mm,yshift=-0.62cm]frame.north west) {Question 82 C};
\node[anchor=east,text=white,font=\sffamily\itshape\large] at ([xshift=-5mm,yshift=-0.62cm]frame.north east) {Proportionnalité et données};}]
\vspace*{1.75cm}
\begin{minipage}[t][6.9cm][c]{\linewidth}\LARGE\centering
\begin{minipage}{0.95\linewidth}\raggedright\linespread{1.25}\selectfont Pour $6$ personnes, il faut $90$ g de farine. Combien en faut-il pour $4$ personnes~?\end{minipage}
\end{minipage}

\vfill
\begin{tcolorbox}[enhanced,nobeforeafter,width=7.1cm,height=2.1cm,valign=center,colback=cpro!6,colframe=cpro,arc=2.5mm,boxrule=0.9pt,left=3mm]{\Large\bfseries\color{cpro}a)}\quad\begin{minipage}[c]{5.6cm}\raggedright\Large $88~\text{g}$\end{minipage}\end{tcolorbox}\hfill\begin{tcolorbox}[enhanced,nobeforeafter,width=7.1cm,height=2.1cm,valign=center,colback=cpro!6,colframe=cpro,arc=2.5mm,boxrule=0.9pt,left=3mm]{\Large\bfseries\color{cpro}b)}\quad\begin{minipage}[c]{5.6cm}\raggedright\Large $360~\text{g}$\end{minipage}\end{tcolorbox}\par\vspace{2.5mm}
\begin{tcolorbox}[enhanced,nobeforeafter,width=7.1cm,height=2.1cm,valign=center,colback=cpro!6,colframe=cpro,arc=2.5mm,boxrule=0.9pt,left=3mm]{\Large\bfseries\color{cpro}c)}\quad\begin{minipage}[c]{5.6cm}\raggedright\Large $15~\text{g}$\end{minipage}\end{tcolorbox}\hfill\begin{tcolorbox}[enhanced,nobeforeafter,width=7.1cm,height=2.1cm,valign=center,colback=cpro!6,colframe=cpro,arc=2.5mm,boxrule=0.9pt,left=3mm]{\Large\bfseries\color{cpro}d)}\quad\begin{minipage}[c]{5.6cm}\raggedright\Large $60~\text{g}$\end{minipage}\end{tcolorbox}
\end{tcolorbox}
\newpage
\begin{tcolorbox}[enhanced,width=16cm,height=14cm,colback=white,colframe=cpro,boxrule=2pt,arc=0mm,
left=6mm,right=6mm,top=0mm,bottom=4mm,sharp corners,
overlay={\fill[cpro] (frame.north west) rectangle ([yshift=-1.25cm]frame.north east);
\node[anchor=west,text=white,font=\sffamily\bfseries\LARGE] at ([xshift=5mm,yshift=-0.62cm]frame.north west) {Question 83 C};
\node[anchor=east,text=white,font=\sffamily\itshape\large] at ([xshift=-5mm,yshift=-0.62cm]frame.north east) {Proportionnalité et données};}]
\vspace*{1.75cm}
\begin{minipage}[t][6.9cm][c]{\linewidth}\LARGE\centering
\begin{minipage}{0.95\linewidth}\raggedright\linespread{1.25}\selectfont $8$ cahiers coûtent $1\,600$ F. Combien coûte un cahier~?\end{minipage}
\end{minipage}

\vfill
\begin{tcolorbox}[enhanced,nobeforeafter,width=7.1cm,height=2.1cm,valign=center,colback=cpro!6,colframe=cpro,arc=2.5mm,boxrule=0.9pt,left=3mm]{\Large\bfseries\color{cpro}a)}\quad\begin{minipage}[c]{5.6cm}\raggedright\Large $208~\text{F}$\end{minipage}\end{tcolorbox}\hfill\begin{tcolorbox}[enhanced,nobeforeafter,width=7.1cm,height=2.1cm,valign=center,colback=cpro!6,colframe=cpro,arc=2.5mm,boxrule=0.9pt,left=3mm]{\Large\bfseries\color{cpro}b)}\quad\begin{minipage}[c]{5.6cm}\raggedright\Large $1\,592~\text{F}$\end{minipage}\end{tcolorbox}\par\vspace{2.5mm}
\begin{tcolorbox}[enhanced,nobeforeafter,width=7.1cm,height=2.1cm,valign=center,colback=cpro!6,colframe=cpro,arc=2.5mm,boxrule=0.9pt,left=3mm]{\Large\bfseries\color{cpro}c)}\quad\begin{minipage}[c]{5.6cm}\raggedright\Large $200~\text{F}$\end{minipage}\end{tcolorbox}\hfill\begin{tcolorbox}[enhanced,nobeforeafter,width=7.1cm,height=2.1cm,valign=center,colback=cpro!6,colframe=cpro,arc=2.5mm,boxrule=0.9pt,left=3mm]{\Large\bfseries\color{cpro}d)}\quad\begin{minipage}[c]{5.6cm}\raggedright\Large $400~\text{F}$\end{minipage}\end{tcolorbox}
\end{tcolorbox}
\newpage
\begin{tcolorbox}[enhanced,width=16cm,height=14cm,colback=white,colframe=cpro,boxrule=2pt,arc=0mm,
left=6mm,right=6mm,top=0mm,bottom=4mm,sharp corners,
overlay={\fill[cpro] (frame.north west) rectangle ([yshift=-1.25cm]frame.north east);
\node[anchor=west,text=white,font=\sffamily\bfseries\LARGE] at ([xshift=5mm,yshift=-0.62cm]frame.north west) {Question 84 C};
\node[anchor=east,text=white,font=\sffamily\itshape\large] at ([xshift=-5mm,yshift=-0.62cm]frame.north east) {Proportionnalité et données};}]
\vspace*{1.75cm}
\begin{minipage}[t][6.9cm][c]{\linewidth}\LARGE\centering
\begin{minipage}{0.95\linewidth}\raggedright\linespread{1.25}\selectfont $12$ cahiers coûtent $900$ F. Combien coûte un cahier~?\end{minipage}
\end{minipage}

\vfill
\begin{tcolorbox}[enhanced,nobeforeafter,width=7.1cm,height=2.1cm,valign=center,colback=cpro!6,colframe=cpro,arc=2.5mm,boxrule=0.9pt,left=3mm]{\Large\bfseries\color{cpro}a)}\quad\begin{minipage}[c]{5.6cm}\raggedright\Large $888~\text{F}$\end{minipage}\end{tcolorbox}\hfill\begin{tcolorbox}[enhanced,nobeforeafter,width=7.1cm,height=2.1cm,valign=center,colback=cpro!6,colframe=cpro,arc=2.5mm,boxrule=0.9pt,left=3mm]{\Large\bfseries\color{cpro}b)}\quad\begin{minipage}[c]{5.6cm}\raggedright\Large $75~\text{F}$\end{minipage}\end{tcolorbox}\par\vspace{2.5mm}
\begin{tcolorbox}[enhanced,nobeforeafter,width=7.1cm,height=2.1cm,valign=center,colback=cpro!6,colframe=cpro,arc=2.5mm,boxrule=0.9pt,left=3mm]{\Large\bfseries\color{cpro}c)}\quad\begin{minipage}[c]{5.6cm}\raggedright\Large $87~\text{F}$\end{minipage}\end{tcolorbox}\hfill\begin{tcolorbox}[enhanced,nobeforeafter,width=7.1cm,height=2.1cm,valign=center,colback=cpro!6,colframe=cpro,arc=2.5mm,boxrule=0.9pt,left=3mm]{\Large\bfseries\color{cpro}d)}\quad\begin{minipage}[c]{5.6cm}\raggedright\Large $150~\text{F}$\end{minipage}\end{tcolorbox}
\end{tcolorbox}
\newpage
\begin{tcolorbox}[enhanced,width=16cm,height=14cm,colback=white,colframe=cpro,boxrule=2pt,arc=0mm,
left=6mm,right=6mm,top=0mm,bottom=4mm,sharp corners,
overlay={\fill[cpro] (frame.north west) rectangle ([yshift=-1.25cm]frame.north east);
\node[anchor=west,text=white,font=\sffamily\bfseries\LARGE] at ([xshift=5mm,yshift=-0.62cm]frame.north west) {Question 85 C};
\node[anchor=east,text=white,font=\sffamily\itshape\large] at ([xshift=-5mm,yshift=-0.62cm]frame.north east) {Proportionnalité et données};}]
\vspace*{1.75cm}
\begin{minipage}[t][6.9cm][c]{\linewidth}\LARGE\centering
\begin{minipage}{0.95\linewidth}\raggedright\linespread{1.25}\selectfont Que vaut $75\,\%$ de $40$~?\end{minipage}
\end{minipage}

\vfill
\begin{tcolorbox}[enhanced,nobeforeafter,width=7.1cm,height=2.1cm,valign=center,colback=cpro!6,colframe=cpro,arc=2.5mm,boxrule=0.9pt,left=3mm]{\Large\bfseries\color{cpro}a)}\quad\begin{minipage}[c]{5.6cm}\raggedright\Large $40$\end{minipage}\end{tcolorbox}\hfill\begin{tcolorbox}[enhanced,nobeforeafter,width=7.1cm,height=2.1cm,valign=center,colback=cpro!6,colframe=cpro,arc=2.5mm,boxrule=0.9pt,left=3mm]{\Large\bfseries\color{cpro}b)}\quad\begin{minipage}[c]{5.6cm}\raggedright\Large $30$\end{minipage}\end{tcolorbox}\par\vspace{2.5mm}
\begin{tcolorbox}[enhanced,nobeforeafter,width=7.1cm,height=2.1cm,valign=center,colback=cpro!6,colframe=cpro,arc=2.5mm,boxrule=0.9pt,left=3mm]{\Large\bfseries\color{cpro}c)}\quad\begin{minipage}[c]{5.6cm}\raggedright\Large $60$\end{minipage}\end{tcolorbox}\hfill\begin{tcolorbox}[enhanced,nobeforeafter,width=7.1cm,height=2.1cm,valign=center,colback=cpro!6,colframe=cpro,arc=2.5mm,boxrule=0.9pt,left=3mm]{\Large\bfseries\color{cpro}d)}\quad\begin{minipage}[c]{5.6cm}\raggedright\Large $10$\end{minipage}\end{tcolorbox}
\end{tcolorbox}
\newpage
\begin{tcolorbox}[enhanced,width=16cm,height=14cm,colback=white,colframe=cpro,boxrule=2pt,arc=0mm,
left=6mm,right=6mm,top=0mm,bottom=4mm,sharp corners,
overlay={\fill[cpro] (frame.north west) rectangle ([yshift=-1.25cm]frame.north east);
\node[anchor=west,text=white,font=\sffamily\bfseries\LARGE] at ([xshift=5mm,yshift=-0.62cm]frame.north west) {Question 86 C};
\node[anchor=east,text=white,font=\sffamily\itshape\large] at ([xshift=-5mm,yshift=-0.62cm]frame.north east) {Proportionnalité et données};}]
\vspace*{1.75cm}
\begin{minipage}[t][6.9cm][c]{\linewidth}\LARGE\centering
\begin{minipage}{0.95\linewidth}\raggedright\linespread{1.25}\selectfont Que vaut $50\,\%$ de $1\,300$~?\end{minipage}
\end{minipage}

\vfill
\begin{tcolorbox}[enhanced,nobeforeafter,width=7.1cm,height=2.1cm,valign=center,colback=cpro!6,colframe=cpro,arc=2.5mm,boxrule=0.9pt,left=3mm]{\Large\bfseries\color{cpro}a)}\quad\begin{minipage}[c]{5.6cm}\raggedright\Large $65$\end{minipage}\end{tcolorbox}\hfill\begin{tcolorbox}[enhanced,nobeforeafter,width=7.1cm,height=2.1cm,valign=center,colback=cpro!6,colframe=cpro,arc=2.5mm,boxrule=0.9pt,left=3mm]{\Large\bfseries\color{cpro}b)}\quad\begin{minipage}[c]{5.6cm}\raggedright\Large $650$\end{minipage}\end{tcolorbox}\par\vspace{2.5mm}
\begin{tcolorbox}[enhanced,nobeforeafter,width=7.1cm,height=2.1cm,valign=center,colback=cpro!6,colframe=cpro,arc=2.5mm,boxrule=0.9pt,left=3mm]{\Large\bfseries\color{cpro}c)}\quad\begin{minipage}[c]{5.6cm}\raggedright\Large $26$\end{minipage}\end{tcolorbox}\hfill\begin{tcolorbox}[enhanced,nobeforeafter,width=7.1cm,height=2.1cm,valign=center,colback=cpro!6,colframe=cpro,arc=2.5mm,boxrule=0.9pt,left=3mm]{\Large\bfseries\color{cpro}d)}\quad\begin{minipage}[c]{5.6cm}\raggedright\Large $660$\end{minipage}\end{tcolorbox}
\end{tcolorbox}
\newpage
\begin{tcolorbox}[enhanced,width=16cm,height=14cm,colback=white,colframe=cpro,boxrule=2pt,arc=0mm,
left=6mm,right=6mm,top=0mm,bottom=4mm,sharp corners,
overlay={\fill[cpro] (frame.north west) rectangle ([yshift=-1.25cm]frame.north east);
\node[anchor=west,text=white,font=\sffamily\bfseries\LARGE] at ([xshift=5mm,yshift=-0.62cm]frame.north west) {Question 87 C};
\node[anchor=east,text=white,font=\sffamily\itshape\large] at ([xshift=-5mm,yshift=-0.62cm]frame.north east) {Proportionnalité et données};}]
\vspace*{1.75cm}
\begin{minipage}[t][6.9cm][c]{\linewidth}\LARGE\centering
\begin{minipage}{0.95\linewidth}\raggedright\linespread{1.25}\selectfont Sur un plan à l'échelle $\dfrac{1}{10\,000}$, une longueur mesure $2$ cm. Quelle est la longueur réelle~?\end{minipage}
\end{minipage}

\vfill
\begin{tcolorbox}[enhanced,nobeforeafter,width=7.1cm,height=2.1cm,valign=center,colback=cpro!6,colframe=cpro,arc=2.5mm,boxrule=0.9pt,left=3mm]{\Large\bfseries\color{cpro}a)}\quad\begin{minipage}[c]{5.6cm}\raggedright\Large $200~\text{m}$\end{minipage}\end{tcolorbox}\hfill\begin{tcolorbox}[enhanced,nobeforeafter,width=7.1cm,height=2.1cm,valign=center,colback=cpro!6,colframe=cpro,arc=2.5mm,boxrule=0.9pt,left=3mm]{\Large\bfseries\color{cpro}b)}\quad\begin{minipage}[c]{5.6cm}\raggedright\Large $20~\text{m}$\end{minipage}\end{tcolorbox}\par\vspace{2.5mm}
\begin{tcolorbox}[enhanced,nobeforeafter,width=7.1cm,height=2.1cm,valign=center,colback=cpro!6,colframe=cpro,arc=2.5mm,boxrule=0.9pt,left=3mm]{\Large\bfseries\color{cpro}c)}\quad\begin{minipage}[c]{5.6cm}\raggedright\Large $2\,000~\text{m}$\end{minipage}\end{tcolorbox}\hfill\begin{tcolorbox}[enhanced,nobeforeafter,width=7.1cm,height=2.1cm,valign=center,colback=cpro!6,colframe=cpro,arc=2.5mm,boxrule=0.9pt,left=3mm]{\Large\bfseries\color{cpro}d)}\quad\begin{minipage}[c]{5.6cm}\raggedright\Large $20\,000~\text{m}$\end{minipage}\end{tcolorbox}
\end{tcolorbox}
\newpage
\begin{tcolorbox}[enhanced,width=16cm,height=14cm,colback=white,colframe=cpro,boxrule=2pt,arc=0mm,
left=6mm,right=6mm,top=0mm,bottom=4mm,sharp corners,
overlay={\fill[cpro] (frame.north west) rectangle ([yshift=-1.25cm]frame.north east);
\node[anchor=west,text=white,font=\sffamily\bfseries\LARGE] at ([xshift=5mm,yshift=-0.62cm]frame.north west) {Question 88 C};
\node[anchor=east,text=white,font=\sffamily\itshape\large] at ([xshift=-5mm,yshift=-0.62cm]frame.north east) {Proportionnalité et données};}]
\vspace*{1.75cm}
\begin{minipage}[t][6.9cm][c]{\linewidth}\LARGE\centering
\begin{minipage}{0.95\linewidth}\raggedright\linespread{1.25}\selectfont Sur un plan à l'échelle $\dfrac{1}{250}$, une longueur mesure $4$ cm. Quelle est la longueur réelle~?\end{minipage}
\end{minipage}

\vfill
\begin{tcolorbox}[enhanced,nobeforeafter,width=7.1cm,height=2.1cm,valign=center,colback=cpro!6,colframe=cpro,arc=2.5mm,boxrule=0.9pt,left=3mm]{\Large\bfseries\color{cpro}a)}\quad\begin{minipage}[c]{5.6cm}\raggedright\Large $100~\text{m}$\end{minipage}\end{tcolorbox}\hfill\begin{tcolorbox}[enhanced,nobeforeafter,width=7.1cm,height=2.1cm,valign=center,colback=cpro!6,colframe=cpro,arc=2.5mm,boxrule=0.9pt,left=3mm]{\Large\bfseries\color{cpro}b)}\quad\begin{minipage}[c]{5.6cm}\raggedright\Large $1~\text{m}$\end{minipage}\end{tcolorbox}\par\vspace{2.5mm}
\begin{tcolorbox}[enhanced,nobeforeafter,width=7.1cm,height=2.1cm,valign=center,colback=cpro!6,colframe=cpro,arc=2.5mm,boxrule=0.9pt,left=3mm]{\Large\bfseries\color{cpro}c)}\quad\begin{minipage}[c]{5.6cm}\raggedright\Large $1\,000~\text{m}$\end{minipage}\end{tcolorbox}\hfill\begin{tcolorbox}[enhanced,nobeforeafter,width=7.1cm,height=2.1cm,valign=center,colback=cpro!6,colframe=cpro,arc=2.5mm,boxrule=0.9pt,left=3mm]{\Large\bfseries\color{cpro}d)}\quad\begin{minipage}[c]{5.6cm}\raggedright\Large $10~\text{m}$\end{minipage}\end{tcolorbox}
\end{tcolorbox}
\newpage
\begin{tcolorbox}[enhanced,width=16cm,height=14cm,colback=white,colframe=cpro,boxrule=2pt,arc=0mm,
left=6mm,right=6mm,top=0mm,bottom=4mm,sharp corners,
overlay={\fill[cpro] (frame.north west) rectangle ([yshift=-1.25cm]frame.north east);
\node[anchor=west,text=white,font=\sffamily\bfseries\LARGE] at ([xshift=5mm,yshift=-0.62cm]frame.north west) {Question 89 C};
\node[anchor=east,text=white,font=\sffamily\itshape\large] at ([xshift=-5mm,yshift=-0.62cm]frame.north east) {Proportionnalité et données};}]
\vspace*{1.75cm}
\begin{minipage}[t][6.9cm][c]{\linewidth}\Large\centering
\begin{minipage}{0.95\linewidth}\raggedright\linespread{1.25}\selectfont Ce tableau est un tableau de proportionnalité~: $\begin{array}{|c|c|c|}\hline 3 & 5 & 10\\\hline 21 & 35 & 70\\\hline\end{array}$. Par quel nombre multiplie-t-on la 1re ligne pour obtenir la 2e~?\end{minipage}
\end{minipage}

\vfill
\begin{tcolorbox}[enhanced,nobeforeafter,width=7.1cm,height=2.1cm,valign=center,colback=cpro!6,colframe=cpro,arc=2.5mm,boxrule=0.9pt,left=3mm]{\Large\bfseries\color{cpro}a)}\quad\begin{minipage}[c]{5.6cm}\raggedright\Large $7$\end{minipage}\end{tcolorbox}\hfill\begin{tcolorbox}[enhanced,nobeforeafter,width=7.1cm,height=2.1cm,valign=center,colback=cpro!6,colframe=cpro,arc=2.5mm,boxrule=0.9pt,left=3mm]{\Large\bfseries\color{cpro}b)}\quad\begin{minipage}[c]{5.6cm}\raggedright\Large $\dfrac{1}{7}$\end{minipage}\end{tcolorbox}\par\vspace{2.5mm}
\begin{tcolorbox}[enhanced,nobeforeafter,width=7.1cm,height=2.1cm,valign=center,colback=cpro!6,colframe=cpro,arc=2.5mm,boxrule=0.9pt,left=3mm]{\Large\bfseries\color{cpro}c)}\quad\begin{minipage}[c]{5.6cm}\raggedright\Large $21$\end{minipage}\end{tcolorbox}\hfill\begin{tcolorbox}[enhanced,nobeforeafter,width=7.1cm,height=2.1cm,valign=center,colback=cpro!6,colframe=cpro,arc=2.5mm,boxrule=0.9pt,left=3mm]{\Large\bfseries\color{cpro}d)}\quad\begin{minipage}[c]{5.6cm}\raggedright\Large $8$\end{minipage}\end{tcolorbox}
\end{tcolorbox}
\newpage
\begin{tcolorbox}[enhanced,width=16cm,height=14cm,colback=white,colframe=cpro,boxrule=2pt,arc=0mm,
left=6mm,right=6mm,top=0mm,bottom=4mm,sharp corners,
overlay={\fill[cpro] (frame.north west) rectangle ([yshift=-1.25cm]frame.north east);
\node[anchor=west,text=white,font=\sffamily\bfseries\LARGE] at ([xshift=5mm,yshift=-0.62cm]frame.north west) {Question 90 C};
\node[anchor=east,text=white,font=\sffamily\itshape\large] at ([xshift=-5mm,yshift=-0.62cm]frame.north east) {Proportionnalité et données};}]
\vspace*{1.75cm}
\begin{minipage}[t][6.9cm][c]{\linewidth}\Large\centering
\begin{minipage}{0.95\linewidth}\raggedright\linespread{1.25}\selectfont Ce tableau est un tableau de proportionnalité~: $\begin{array}{|c|c|c|}\hline 2 & 6 & 8\\\hline 3 & 9 & 12\\\hline\end{array}$. Par quel nombre multiplie-t-on la 1re ligne pour obtenir la 2e~?\end{minipage}
\end{minipage}

\vfill
\begin{tcolorbox}[enhanced,nobeforeafter,width=7.1cm,height=2.1cm,valign=center,colback=cpro!6,colframe=cpro,arc=2.5mm,boxrule=0.9pt,left=3mm]{\Large\bfseries\color{cpro}a)}\quad\begin{minipage}[c]{5.6cm}\raggedright\Large $3$\end{minipage}\end{tcolorbox}\hfill\begin{tcolorbox}[enhanced,nobeforeafter,width=7.1cm,height=2.1cm,valign=center,colback=cpro!6,colframe=cpro,arc=2.5mm,boxrule=0.9pt,left=3mm]{\Large\bfseries\color{cpro}b)}\quad\begin{minipage}[c]{5.6cm}\raggedright\Large $2{,}5$\end{minipage}\end{tcolorbox}\par\vspace{2.5mm}
\begin{tcolorbox}[enhanced,nobeforeafter,width=7.1cm,height=2.1cm,valign=center,colback=cpro!6,colframe=cpro,arc=2.5mm,boxrule=0.9pt,left=3mm]{\Large\bfseries\color{cpro}c)}\quad\begin{minipage}[c]{5.6cm}\raggedright\Large $1{,}5$\end{minipage}\end{tcolorbox}\hfill\begin{tcolorbox}[enhanced,nobeforeafter,width=7.1cm,height=2.1cm,valign=center,colback=cpro!6,colframe=cpro,arc=2.5mm,boxrule=0.9pt,left=3mm]{\Large\bfseries\color{cpro}d)}\quad\begin{minipage}[c]{5.6cm}\raggedright\Large $\dfrac{2}{3}$\end{minipage}\end{tcolorbox}
\end{tcolorbox}
\newpage
\begin{tcolorbox}[enhanced,width=16cm,height=14cm,colback=white,colframe=cpro,boxrule=2pt,arc=0mm,
left=6mm,right=6mm,top=0mm,bottom=4mm,sharp corners,
overlay={\fill[cpro] (frame.north west) rectangle ([yshift=-1.25cm]frame.north east);
\node[anchor=west,text=white,font=\sffamily\bfseries\LARGE] at ([xshift=5mm,yshift=-0.62cm]frame.north west) {Question 91 C};
\node[anchor=east,text=white,font=\sffamily\itshape\large] at ([xshift=-5mm,yshift=-0.62cm]frame.north east) {Proportionnalité et données};}]
\vspace*{1.75cm}
\begin{minipage}[t][6.9cm][c]{\linewidth}\Large\centering
\begin{minipage}{0.95\linewidth}\raggedright\linespread{1.25}\selectfont Dans une classe, les élèves viennent au collège ainsi~: $5$ à pied, $9$ en bus, $4$ à vélo et $6$ en voiture. Combien y a-t-il d'élèves dans la classe~?\end{minipage}
\end{minipage}

\vfill
\begin{tcolorbox}[enhanced,nobeforeafter,width=7.1cm,height=2.1cm,valign=center,colback=cpro!6,colframe=cpro,arc=2.5mm,boxrule=0.9pt,left=3mm]{\Large\bfseries\color{cpro}a)}\quad\begin{minipage}[c]{5.6cm}\raggedright\Large $24$\end{minipage}\end{tcolorbox}\hfill\begin{tcolorbox}[enhanced,nobeforeafter,width=7.1cm,height=2.1cm,valign=center,colback=cpro!6,colframe=cpro,arc=2.5mm,boxrule=0.9pt,left=3mm]{\Large\bfseries\color{cpro}b)}\quad\begin{minipage}[c]{5.6cm}\raggedright\Large $25$\end{minipage}\end{tcolorbox}\par\vspace{2.5mm}
\begin{tcolorbox}[enhanced,nobeforeafter,width=7.1cm,height=2.1cm,valign=center,colback=cpro!6,colframe=cpro,arc=2.5mm,boxrule=0.9pt,left=3mm]{\Large\bfseries\color{cpro}c)}\quad\begin{minipage}[c]{5.6cm}\raggedright\Large $9$\end{minipage}\end{tcolorbox}\hfill\begin{tcolorbox}[enhanced,nobeforeafter,width=7.1cm,height=2.1cm,valign=center,colback=cpro!6,colframe=cpro,arc=2.5mm,boxrule=0.9pt,left=3mm]{\Large\bfseries\color{cpro}d)}\quad\begin{minipage}[c]{5.6cm}\raggedright\Large $20$\end{minipage}\end{tcolorbox}
\end{tcolorbox}
\newpage
\begin{tcolorbox}[enhanced,width=16cm,height=14cm,colback=white,colframe=cpro,boxrule=2pt,arc=0mm,
left=6mm,right=6mm,top=0mm,bottom=4mm,sharp corners,
overlay={\fill[cpro] (frame.north west) rectangle ([yshift=-1.25cm]frame.north east);
\node[anchor=west,text=white,font=\sffamily\bfseries\LARGE] at ([xshift=5mm,yshift=-0.62cm]frame.north west) {Question 92 C};
\node[anchor=east,text=white,font=\sffamily\itshape\large] at ([xshift=-5mm,yshift=-0.62cm]frame.north east) {Proportionnalité et données};}]
\vspace*{1.75cm}
\begin{minipage}[t][6.9cm][c]{\linewidth}\LARGE\centering
\begin{minipage}{0.95\linewidth}\raggedright\linespread{1.25}\selectfont Dans une classe, les élèves viennent au collège ainsi~: $13$ à pied, $2$ en bus et $10$ à vélo. Combien y a-t-il d'élèves dans la classe~?\end{minipage}
\end{minipage}

\vfill
\begin{tcolorbox}[enhanced,nobeforeafter,width=7.1cm,height=2.1cm,valign=center,colback=cpro!6,colframe=cpro,arc=2.5mm,boxrule=0.9pt,left=3mm]{\Large\bfseries\color{cpro}a)}\quad\begin{minipage}[c]{5.6cm}\raggedright\Large $26$\end{minipage}\end{tcolorbox}\hfill\begin{tcolorbox}[enhanced,nobeforeafter,width=7.1cm,height=2.1cm,valign=center,colback=cpro!6,colframe=cpro,arc=2.5mm,boxrule=0.9pt,left=3mm]{\Large\bfseries\color{cpro}b)}\quad\begin{minipage}[c]{5.6cm}\raggedright\Large $23$\end{minipage}\end{tcolorbox}\par\vspace{2.5mm}
\begin{tcolorbox}[enhanced,nobeforeafter,width=7.1cm,height=2.1cm,valign=center,colback=cpro!6,colframe=cpro,arc=2.5mm,boxrule=0.9pt,left=3mm]{\Large\bfseries\color{cpro}c)}\quad\begin{minipage}[c]{5.6cm}\raggedright\Large $13$\end{minipage}\end{tcolorbox}\hfill\begin{tcolorbox}[enhanced,nobeforeafter,width=7.1cm,height=2.1cm,valign=center,colback=cpro!6,colframe=cpro,arc=2.5mm,boxrule=0.9pt,left=3mm]{\Large\bfseries\color{cpro}d)}\quad\begin{minipage}[c]{5.6cm}\raggedright\Large $25$\end{minipage}\end{tcolorbox}
\end{tcolorbox}
\newpage
\begin{tcolorbox}[enhanced,width=16cm,height=14cm,colback=white,colframe=cpro,boxrule=2pt,arc=0mm,
left=6mm,right=6mm,top=0mm,bottom=4mm,sharp corners,
overlay={\fill[cpro] (frame.north west) rectangle ([yshift=-1.25cm]frame.north east);
\node[anchor=west,text=white,font=\sffamily\bfseries\LARGE] at ([xshift=5mm,yshift=-0.62cm]frame.north west) {Question 93 C};
\node[anchor=east,text=white,font=\sffamily\itshape\large] at ([xshift=-5mm,yshift=-0.62cm]frame.north east) {Proportionnalité et données};}]
\vspace*{1.75cm}
\begin{minipage}[t][6.9cm][c]{\linewidth}\LARGE\centering
\begin{minipage}{0.95\linewidth}\raggedright\linespread{1.25}\selectfont Ce tableau est-il un tableau de proportionnalité~? $\begin{array}{|c|c|}\hline 6 & 9\\\hline 8 & 12\\\hline\end{array}$\end{minipage}
\end{minipage}

\vfill
\begin{tcolorbox}[enhanced,nobeforeafter,width=7.1cm,height=2.1cm,valign=center,colback=cpro!6,colframe=cpro,arc=2.5mm,boxrule=0.9pt,left=3mm]{\Large\bfseries\color{cpro}a)}\quad\begin{minipage}[c]{5.6cm}\raggedright\Large oui\end{minipage}\end{tcolorbox}\hfill\begin{tcolorbox}[enhanced,nobeforeafter,width=7.1cm,height=2.1cm,valign=center,colback=cpro!6,colframe=cpro,arc=2.5mm,boxrule=0.9pt,left=3mm]{\Large\bfseries\color{cpro}b)}\quad\begin{minipage}[c]{5.6cm}\raggedright\Large on ne peut pas savoir\end{minipage}\end{tcolorbox}\par\vspace{2.5mm}
\begin{tcolorbox}[enhanced,nobeforeafter,width=7.1cm,height=2.1cm,valign=center,colback=cpro!6,colframe=cpro,arc=2.5mm,boxrule=0.9pt,left=3mm]{\Large\bfseries\color{cpro}c)}\quad\begin{minipage}[c]{5.6cm}\raggedright\Large il manque une colonne\end{minipage}\end{tcolorbox}\hfill\begin{tcolorbox}[enhanced,nobeforeafter,width=7.1cm,height=2.1cm,valign=center,colback=cpro!6,colframe=cpro,arc=2.5mm,boxrule=0.9pt,left=3mm]{\Large\bfseries\color{cpro}d)}\quad\begin{minipage}[c]{5.6cm}\raggedright\Large non\end{minipage}\end{tcolorbox}
\end{tcolorbox}
\newpage
\begin{tcolorbox}[enhanced,width=16cm,height=14cm,colback=white,colframe=cpro,boxrule=2pt,arc=0mm,
left=6mm,right=6mm,top=0mm,bottom=4mm,sharp corners,
overlay={\fill[cpro] (frame.north west) rectangle ([yshift=-1.25cm]frame.north east);
\node[anchor=west,text=white,font=\sffamily\bfseries\LARGE] at ([xshift=5mm,yshift=-0.62cm]frame.north west) {Question 94 C};
\node[anchor=east,text=white,font=\sffamily\itshape\large] at ([xshift=-5mm,yshift=-0.62cm]frame.north east) {Proportionnalité et données};}]
\vspace*{1.75cm}
\begin{minipage}[t][6.9cm][c]{\linewidth}\LARGE\centering
\begin{minipage}{0.95\linewidth}\raggedright\linespread{1.25}\selectfont Ce tableau est-il un tableau de proportionnalité~? $\begin{array}{|c|c|}\hline 2 & 8\\\hline 5 & 18\\\hline\end{array}$\end{minipage}
\end{minipage}

\vfill
\begin{tcolorbox}[enhanced,nobeforeafter,width=7.1cm,height=2.1cm,valign=center,colback=cpro!6,colframe=cpro,arc=2.5mm,boxrule=0.9pt,left=3mm]{\Large\bfseries\color{cpro}a)}\quad\begin{minipage}[c]{5.6cm}\raggedright\Large non\end{minipage}\end{tcolorbox}\hfill\begin{tcolorbox}[enhanced,nobeforeafter,width=7.1cm,height=2.1cm,valign=center,colback=cpro!6,colframe=cpro,arc=2.5mm,boxrule=0.9pt,left=3mm]{\Large\bfseries\color{cpro}b)}\quad\begin{minipage}[c]{5.6cm}\raggedright\Large on ne peut pas savoir\end{minipage}\end{tcolorbox}\par\vspace{2.5mm}
\begin{tcolorbox}[enhanced,nobeforeafter,width=7.1cm,height=2.1cm,valign=center,colback=cpro!6,colframe=cpro,arc=2.5mm,boxrule=0.9pt,left=3mm]{\Large\bfseries\color{cpro}c)}\quad\begin{minipage}[c]{5.6cm}\raggedright\Large il manque une colonne\end{minipage}\end{tcolorbox}\hfill\begin{tcolorbox}[enhanced,nobeforeafter,width=7.1cm,height=2.1cm,valign=center,colback=cpro!6,colframe=cpro,arc=2.5mm,boxrule=0.9pt,left=3mm]{\Large\bfseries\color{cpro}d)}\quad\begin{minipage}[c]{5.6cm}\raggedright\Large oui\end{minipage}\end{tcolorbox}
\end{tcolorbox}
\newpage
\begin{tcolorbox}[enhanced,width=16cm,height=14cm,colback=white,colframe=cpro,boxrule=2pt,arc=0mm,
left=6mm,right=6mm,top=0mm,bottom=4mm,sharp corners,
overlay={\fill[cpro] (frame.north west) rectangle ([yshift=-1.25cm]frame.north east);
\node[anchor=west,text=white,font=\sffamily\bfseries\LARGE] at ([xshift=5mm,yshift=-0.62cm]frame.north west) {Question 95 C};
\node[anchor=east,text=white,font=\sffamily\itshape\large] at ([xshift=-5mm,yshift=-0.62cm]frame.north east) {Proportionnalité et données};}]
\vspace*{1.75cm}
\begin{minipage}[t][6.9cm][c]{\linewidth}\LARGE\centering
\begin{minipage}{0.95\linewidth}\raggedright\linespread{1.25}\selectfont Un robinet remplit $40$ L en $8$ min. Quel volume remplit-il en $3$ min, au même débit~?\end{minipage}
\end{minipage}

\vfill
\begin{tcolorbox}[enhanced,nobeforeafter,width=7.1cm,height=2.1cm,valign=center,colback=cpro!6,colframe=cpro,arc=2.5mm,boxrule=0.9pt,left=3mm]{\Large\bfseries\color{cpro}a)}\quad\begin{minipage}[c]{5.6cm}\raggedright\Large $35~\text{L}$\end{minipage}\end{tcolorbox}\hfill\begin{tcolorbox}[enhanced,nobeforeafter,width=7.1cm,height=2.1cm,valign=center,colback=cpro!6,colframe=cpro,arc=2.5mm,boxrule=0.9pt,left=3mm]{\Large\bfseries\color{cpro}b)}\quad\begin{minipage}[c]{5.6cm}\raggedright\Large $15~\text{L}$\end{minipage}\end{tcolorbox}\par\vspace{2.5mm}
\begin{tcolorbox}[enhanced,nobeforeafter,width=7.1cm,height=2.1cm,valign=center,colback=cpro!6,colframe=cpro,arc=2.5mm,boxrule=0.9pt,left=3mm]{\Large\bfseries\color{cpro}c)}\quad\begin{minipage}[c]{5.6cm}\raggedright\Large $120~\text{L}$\end{minipage}\end{tcolorbox}\hfill\begin{tcolorbox}[enhanced,nobeforeafter,width=7.1cm,height=2.1cm,valign=center,colback=cpro!6,colframe=cpro,arc=2.5mm,boxrule=0.9pt,left=3mm]{\Large\bfseries\color{cpro}d)}\quad\begin{minipage}[c]{5.6cm}\raggedright\Large $5~\text{L}$\end{minipage}\end{tcolorbox}
\end{tcolorbox}
\newpage
\begin{tcolorbox}[enhanced,width=16cm,height=14cm,colback=white,colframe=cpro,boxrule=2pt,arc=0mm,
left=6mm,right=6mm,top=0mm,bottom=4mm,sharp corners,
overlay={\fill[cpro] (frame.north west) rectangle ([yshift=-1.25cm]frame.north east);
\node[anchor=west,text=white,font=\sffamily\bfseries\LARGE] at ([xshift=5mm,yshift=-0.62cm]frame.north west) {Question 96 C};
\node[anchor=east,text=white,font=\sffamily\itshape\large] at ([xshift=-5mm,yshift=-0.62cm]frame.north east) {Proportionnalité et données};}]
\vspace*{1.75cm}
\begin{minipage}[t][6.9cm][c]{\linewidth}\LARGE\centering
\begin{minipage}{0.95\linewidth}\raggedright\linespread{1.25}\selectfont Un robinet remplit $14$ L en $4$ min. Quel volume remplit-il en $6$ min, au même débit~?\end{minipage}
\end{minipage}

\vfill
\begin{tcolorbox}[enhanced,nobeforeafter,width=7.1cm,height=2.1cm,valign=center,colback=cpro!6,colframe=cpro,arc=2.5mm,boxrule=0.9pt,left=3mm]{\Large\bfseries\color{cpro}a)}\quad\begin{minipage}[c]{5.6cm}\raggedright\Large $21~\text{L}$\end{minipage}\end{tcolorbox}\hfill\begin{tcolorbox}[enhanced,nobeforeafter,width=7.1cm,height=2.1cm,valign=center,colback=cpro!6,colframe=cpro,arc=2.5mm,boxrule=0.9pt,left=3mm]{\Large\bfseries\color{cpro}b)}\quad\begin{minipage}[c]{5.6cm}\raggedright\Large $3{,}5~\text{L}$\end{minipage}\end{tcolorbox}\par\vspace{2.5mm}
\begin{tcolorbox}[enhanced,nobeforeafter,width=7.1cm,height=2.1cm,valign=center,colback=cpro!6,colframe=cpro,arc=2.5mm,boxrule=0.9pt,left=3mm]{\Large\bfseries\color{cpro}c)}\quad\begin{minipage}[c]{5.6cm}\raggedright\Large $16~\text{L}$\end{minipage}\end{tcolorbox}\hfill\begin{tcolorbox}[enhanced,nobeforeafter,width=7.1cm,height=2.1cm,valign=center,colback=cpro!6,colframe=cpro,arc=2.5mm,boxrule=0.9pt,left=3mm]{\Large\bfseries\color{cpro}d)}\quad\begin{minipage}[c]{5.6cm}\raggedright\Large $84~\text{L}$\end{minipage}\end{tcolorbox}
\end{tcolorbox}
\newpage
\begin{tcolorbox}[enhanced,width=16cm,height=14cm,colback=white,colframe=cnum,boxrule=2pt,arc=0mm,
left=6mm,right=6mm,top=0mm,bottom=4mm,sharp corners,
overlay={\fill[cnum] (frame.north west) rectangle ([yshift=-1.25cm]frame.north east);
\node[anchor=west,text=white,font=\sffamily\bfseries\LARGE] at ([xshift=5mm,yshift=-0.62cm]frame.north west) {Question 1 D};
\node[anchor=east,text=white,font=\sffamily\itshape\large] at ([xshift=-5mm,yshift=-0.62cm]frame.north east) {Nombres entiers et décimaux};}]
\vspace*{1.75cm}
\begin{minipage}[t][6.9cm][c]{\linewidth}\LARGE\centering
\begin{minipage}{0.95\linewidth}\raggedright\linespread{1.25}\selectfont Dans le nombre $3{,}509$, quel est le chiffre des centièmes~?\end{minipage}
\end{minipage}

\vfill
\begin{tcolorbox}[enhanced,nobeforeafter,width=7.1cm,height=2.1cm,valign=center,colback=cnum!6,colframe=cnum,arc=2.5mm,boxrule=0.9pt,left=3mm]{\Large\bfseries\color{cnum}a)}\quad\begin{minipage}[c]{5.6cm}\raggedright\Large $0$\end{minipage}\end{tcolorbox}\hfill\begin{tcolorbox}[enhanced,nobeforeafter,width=7.1cm,height=2.1cm,valign=center,colback=cnum!6,colframe=cnum,arc=2.5mm,boxrule=0.9pt,left=3mm]{\Large\bfseries\color{cnum}b)}\quad\begin{minipage}[c]{5.6cm}\raggedright\Large $5$\end{minipage}\end{tcolorbox}\par\vspace{2.5mm}
\begin{tcolorbox}[enhanced,nobeforeafter,width=7.1cm,height=2.1cm,valign=center,colback=cnum!6,colframe=cnum,arc=2.5mm,boxrule=0.9pt,left=3mm]{\Large\bfseries\color{cnum}c)}\quad\begin{minipage}[c]{5.6cm}\raggedright\Large $3$\end{minipage}\end{tcolorbox}\hfill\begin{tcolorbox}[enhanced,nobeforeafter,width=7.1cm,height=2.1cm,valign=center,colback=cnum!6,colframe=cnum,arc=2.5mm,boxrule=0.9pt,left=3mm]{\Large\bfseries\color{cnum}d)}\quad\begin{minipage}[c]{5.6cm}\raggedright\Large $9$\end{minipage}\end{tcolorbox}
\end{tcolorbox}
\newpage
\begin{tcolorbox}[enhanced,width=16cm,height=14cm,colback=white,colframe=cnum,boxrule=2pt,arc=0mm,
left=6mm,right=6mm,top=0mm,bottom=4mm,sharp corners,
overlay={\fill[cnum] (frame.north west) rectangle ([yshift=-1.25cm]frame.north east);
\node[anchor=west,text=white,font=\sffamily\bfseries\LARGE] at ([xshift=5mm,yshift=-0.62cm]frame.north west) {Question 2 D};
\node[anchor=east,text=white,font=\sffamily\itshape\large] at ([xshift=-5mm,yshift=-0.62cm]frame.north east) {Nombres entiers et décimaux};}]
\vspace*{1.75cm}
\begin{minipage}[t][6.9cm][c]{\linewidth}\LARGE\centering
\begin{minipage}{0.95\linewidth}\raggedright\linespread{1.25}\selectfont Dans le nombre $74{,}16$, quel est le chiffre des unités~?\end{minipage}
\end{minipage}

\vfill
\begin{tcolorbox}[enhanced,nobeforeafter,width=7.1cm,height=2.1cm,valign=center,colback=cnum!6,colframe=cnum,arc=2.5mm,boxrule=0.9pt,left=3mm]{\Large\bfseries\color{cnum}a)}\quad\begin{minipage}[c]{5.6cm}\raggedright\Large $7$\end{minipage}\end{tcolorbox}\hfill\begin{tcolorbox}[enhanced,nobeforeafter,width=7.1cm,height=2.1cm,valign=center,colback=cnum!6,colframe=cnum,arc=2.5mm,boxrule=0.9pt,left=3mm]{\Large\bfseries\color{cnum}b)}\quad\begin{minipage}[c]{5.6cm}\raggedright\Large $1$\end{minipage}\end{tcolorbox}\par\vspace{2.5mm}
\begin{tcolorbox}[enhanced,nobeforeafter,width=7.1cm,height=2.1cm,valign=center,colback=cnum!6,colframe=cnum,arc=2.5mm,boxrule=0.9pt,left=3mm]{\Large\bfseries\color{cnum}c)}\quad\begin{minipage}[c]{5.6cm}\raggedright\Large $4$\end{minipage}\end{tcolorbox}\hfill\begin{tcolorbox}[enhanced,nobeforeafter,width=7.1cm,height=2.1cm,valign=center,colback=cnum!6,colframe=cnum,arc=2.5mm,boxrule=0.9pt,left=3mm]{\Large\bfseries\color{cnum}d)}\quad\begin{minipage}[c]{5.6cm}\raggedright\Large $6$\end{minipage}\end{tcolorbox}
\end{tcolorbox}
\newpage
\begin{tcolorbox}[enhanced,width=16cm,height=14cm,colback=white,colframe=cnum,boxrule=2pt,arc=0mm,
left=6mm,right=6mm,top=0mm,bottom=4mm,sharp corners,
overlay={\fill[cnum] (frame.north west) rectangle ([yshift=-1.25cm]frame.north east);
\node[anchor=west,text=white,font=\sffamily\bfseries\LARGE] at ([xshift=5mm,yshift=-0.62cm]frame.north west) {Question 3 D};
\node[anchor=east,text=white,font=\sffamily\itshape\large] at ([xshift=-5mm,yshift=-0.62cm]frame.north east) {Nombres entiers et décimaux};}]
\vspace*{1.75cm}
\begin{minipage}[t][6.9cm][c]{\linewidth}\LARGE\centering
\begin{minipage}{0.95\linewidth}\raggedright\linespread{1.25}\selectfont Quelle est l'écriture décimale de $\dfrac{4507}{100}$~?\end{minipage}
\end{minipage}

\vfill
\begin{tcolorbox}[enhanced,nobeforeafter,width=7.1cm,height=2.1cm,valign=center,colback=cnum!6,colframe=cnum,arc=2.5mm,boxrule=0.9pt,left=3mm]{\Large\bfseries\color{cnum}a)}\quad\begin{minipage}[c]{5.6cm}\raggedright\Large $450{,}7$\end{minipage}\end{tcolorbox}\hfill\begin{tcolorbox}[enhanced,nobeforeafter,width=7.1cm,height=2.1cm,valign=center,colback=cnum!6,colframe=cnum,arc=2.5mm,boxrule=0.9pt,left=3mm]{\Large\bfseries\color{cnum}b)}\quad\begin{minipage}[c]{5.6cm}\raggedright\Large $45{,}07$\end{minipage}\end{tcolorbox}\par\vspace{2.5mm}
\begin{tcolorbox}[enhanced,nobeforeafter,width=7.1cm,height=2.1cm,valign=center,colback=cnum!6,colframe=cnum,arc=2.5mm,boxrule=0.9pt,left=3mm]{\Large\bfseries\color{cnum}c)}\quad\begin{minipage}[c]{5.6cm}\raggedright\Large $4{,}507$\end{minipage}\end{tcolorbox}\hfill\begin{tcolorbox}[enhanced,nobeforeafter,width=7.1cm,height=2.1cm,valign=center,colback=cnum!6,colframe=cnum,arc=2.5mm,boxrule=0.9pt,left=3mm]{\Large\bfseries\color{cnum}d)}\quad\begin{minipage}[c]{5.6cm}\raggedright\Large $4\,507$\end{minipage}\end{tcolorbox}
\end{tcolorbox}
\newpage
\begin{tcolorbox}[enhanced,width=16cm,height=14cm,colback=white,colframe=cnum,boxrule=2pt,arc=0mm,
left=6mm,right=6mm,top=0mm,bottom=4mm,sharp corners,
overlay={\fill[cnum] (frame.north west) rectangle ([yshift=-1.25cm]frame.north east);
\node[anchor=west,text=white,font=\sffamily\bfseries\LARGE] at ([xshift=5mm,yshift=-0.62cm]frame.north west) {Question 4 D};
\node[anchor=east,text=white,font=\sffamily\itshape\large] at ([xshift=-5mm,yshift=-0.62cm]frame.north east) {Nombres entiers et décimaux};}]
\vspace*{1.75cm}
\begin{minipage}[t][6.9cm][c]{\linewidth}\LARGE\centering
\begin{minipage}{0.95\linewidth}\raggedright\linespread{1.25}\selectfont Quelle est l'écriture décimale de $\dfrac{36}{100}$~?\end{minipage}
\end{minipage}

\vfill
\begin{tcolorbox}[enhanced,nobeforeafter,width=7.1cm,height=2.1cm,valign=center,colback=cnum!6,colframe=cnum,arc=2.5mm,boxrule=0.9pt,left=3mm]{\Large\bfseries\color{cnum}a)}\quad\begin{minipage}[c]{5.6cm}\raggedright\Large $0{,}036$\end{minipage}\end{tcolorbox}\hfill\begin{tcolorbox}[enhanced,nobeforeafter,width=7.1cm,height=2.1cm,valign=center,colback=cnum!6,colframe=cnum,arc=2.5mm,boxrule=0.9pt,left=3mm]{\Large\bfseries\color{cnum}b)}\quad\begin{minipage}[c]{5.6cm}\raggedright\Large $3{,}6$\end{minipage}\end{tcolorbox}\par\vspace{2.5mm}
\begin{tcolorbox}[enhanced,nobeforeafter,width=7.1cm,height=2.1cm,valign=center,colback=cnum!6,colframe=cnum,arc=2.5mm,boxrule=0.9pt,left=3mm]{\Large\bfseries\color{cnum}c)}\quad\begin{minipage}[c]{5.6cm}\raggedright\Large $0{,}36$\end{minipage}\end{tcolorbox}\hfill\begin{tcolorbox}[enhanced,nobeforeafter,width=7.1cm,height=2.1cm,valign=center,colback=cnum!6,colframe=cnum,arc=2.5mm,boxrule=0.9pt,left=3mm]{\Large\bfseries\color{cnum}d)}\quad\begin{minipage}[c]{5.6cm}\raggedright\Large $36$\end{minipage}\end{tcolorbox}
\end{tcolorbox}
\newpage
\begin{tcolorbox}[enhanced,width=16cm,height=14cm,colback=white,colframe=cnum,boxrule=2pt,arc=0mm,
left=6mm,right=6mm,top=0mm,bottom=4mm,sharp corners,
overlay={\fill[cnum] (frame.north west) rectangle ([yshift=-1.25cm]frame.north east);
\node[anchor=west,text=white,font=\sffamily\bfseries\LARGE] at ([xshift=5mm,yshift=-0.62cm]frame.north west) {Question 5 D};
\node[anchor=east,text=white,font=\sffamily\itshape\large] at ([xshift=-5mm,yshift=-0.62cm]frame.north east) {Nombres entiers et décimaux};}]
\vspace*{1.75cm}
\begin{minipage}[t][6.9cm][c]{\linewidth}\LARGE\centering
\begin{minipage}{0.95\linewidth}\raggedright\linespread{1.25}\selectfont Quel est le plus grand de ces quatre nombres~: $10{,}1\ ;\ 10{,}01\ ;\ 10{,}11\ ;\ 10{,}011$~?\end{minipage}
\end{minipage}

\vfill
\begin{tcolorbox}[enhanced,nobeforeafter,width=7.1cm,height=2.1cm,valign=center,colback=cnum!6,colframe=cnum,arc=2.5mm,boxrule=0.9pt,left=3mm]{\Large\bfseries\color{cnum}a)}\quad\begin{minipage}[c]{5.6cm}\raggedright\Large $10{,}1$\end{minipage}\end{tcolorbox}\hfill\begin{tcolorbox}[enhanced,nobeforeafter,width=7.1cm,height=2.1cm,valign=center,colback=cnum!6,colframe=cnum,arc=2.5mm,boxrule=0.9pt,left=3mm]{\Large\bfseries\color{cnum}b)}\quad\begin{minipage}[c]{5.6cm}\raggedright\Large $10{,}01$\end{minipage}\end{tcolorbox}\par\vspace{2.5mm}
\begin{tcolorbox}[enhanced,nobeforeafter,width=7.1cm,height=2.1cm,valign=center,colback=cnum!6,colframe=cnum,arc=2.5mm,boxrule=0.9pt,left=3mm]{\Large\bfseries\color{cnum}c)}\quad\begin{minipage}[c]{5.6cm}\raggedright\Large $10{,}11$\end{minipage}\end{tcolorbox}\hfill\begin{tcolorbox}[enhanced,nobeforeafter,width=7.1cm,height=2.1cm,valign=center,colback=cnum!6,colframe=cnum,arc=2.5mm,boxrule=0.9pt,left=3mm]{\Large\bfseries\color{cnum}d)}\quad\begin{minipage}[c]{5.6cm}\raggedright\Large $10{,}011$\end{minipage}\end{tcolorbox}
\end{tcolorbox}
\newpage
\begin{tcolorbox}[enhanced,width=16cm,height=14cm,colback=white,colframe=cnum,boxrule=2pt,arc=0mm,
left=6mm,right=6mm,top=0mm,bottom=4mm,sharp corners,
overlay={\fill[cnum] (frame.north west) rectangle ([yshift=-1.25cm]frame.north east);
\node[anchor=west,text=white,font=\sffamily\bfseries\LARGE] at ([xshift=5mm,yshift=-0.62cm]frame.north west) {Question 6 D};
\node[anchor=east,text=white,font=\sffamily\itshape\large] at ([xshift=-5mm,yshift=-0.62cm]frame.north east) {Nombres entiers et décimaux};}]
\vspace*{1.75cm}
\begin{minipage}[t][6.9cm][c]{\linewidth}\LARGE\centering
\begin{minipage}{0.95\linewidth}\raggedright\linespread{1.25}\selectfont Quel est le plus petit de ces quatre nombres~: $0{,}4\ ;\ 0{,}39\ ;\ 0{,}401\ ;\ 0{,}0499$~?\end{minipage}
\end{minipage}

\vfill
\begin{tcolorbox}[enhanced,nobeforeafter,width=7.1cm,height=2.1cm,valign=center,colback=cnum!6,colframe=cnum,arc=2.5mm,boxrule=0.9pt,left=3mm]{\Large\bfseries\color{cnum}a)}\quad\begin{minipage}[c]{5.6cm}\raggedright\Large $0{,}4$\end{minipage}\end{tcolorbox}\hfill\begin{tcolorbox}[enhanced,nobeforeafter,width=7.1cm,height=2.1cm,valign=center,colback=cnum!6,colframe=cnum,arc=2.5mm,boxrule=0.9pt,left=3mm]{\Large\bfseries\color{cnum}b)}\quad\begin{minipage}[c]{5.6cm}\raggedright\Large $0{,}401$\end{minipage}\end{tcolorbox}\par\vspace{2.5mm}
\begin{tcolorbox}[enhanced,nobeforeafter,width=7.1cm,height=2.1cm,valign=center,colback=cnum!6,colframe=cnum,arc=2.5mm,boxrule=0.9pt,left=3mm]{\Large\bfseries\color{cnum}c)}\quad\begin{minipage}[c]{5.6cm}\raggedright\Large $0{,}0499$\end{minipage}\end{tcolorbox}\hfill\begin{tcolorbox}[enhanced,nobeforeafter,width=7.1cm,height=2.1cm,valign=center,colback=cnum!6,colframe=cnum,arc=2.5mm,boxrule=0.9pt,left=3mm]{\Large\bfseries\color{cnum}d)}\quad\begin{minipage}[c]{5.6cm}\raggedright\Large $0{,}39$\end{minipage}\end{tcolorbox}
\end{tcolorbox}
\newpage
\begin{tcolorbox}[enhanced,width=16cm,height=14cm,colback=white,colframe=cnum,boxrule=2pt,arc=0mm,
left=6mm,right=6mm,top=0mm,bottom=4mm,sharp corners,
overlay={\fill[cnum] (frame.north west) rectangle ([yshift=-1.25cm]frame.north east);
\node[anchor=west,text=white,font=\sffamily\bfseries\LARGE] at ([xshift=5mm,yshift=-0.62cm]frame.north west) {Question 7 D};
\node[anchor=east,text=white,font=\sffamily\itshape\large] at ([xshift=-5mm,yshift=-0.62cm]frame.north east) {Nombres entiers et décimaux};}]
\vspace*{1.75cm}
\begin{minipage}[t][6.9cm][c]{\linewidth}\LARGE\centering
\begin{minipage}{0.95\linewidth}\raggedright\linespread{1.25}\selectfont Quel est l'arrondi au dixième de $3{,}333$~?\end{minipage}
\end{minipage}

\vfill
\begin{tcolorbox}[enhanced,nobeforeafter,width=7.1cm,height=2.1cm,valign=center,colback=cnum!6,colframe=cnum,arc=2.5mm,boxrule=0.9pt,left=3mm]{\Large\bfseries\color{cnum}a)}\quad\begin{minipage}[c]{5.6cm}\raggedright\Large $3{,}33$\end{minipage}\end{tcolorbox}\hfill\begin{tcolorbox}[enhanced,nobeforeafter,width=7.1cm,height=2.1cm,valign=center,colback=cnum!6,colframe=cnum,arc=2.5mm,boxrule=0.9pt,left=3mm]{\Large\bfseries\color{cnum}b)}\quad\begin{minipage}[c]{5.6cm}\raggedright\Large $3{,}3$\end{minipage}\end{tcolorbox}\par\vspace{2.5mm}
\begin{tcolorbox}[enhanced,nobeforeafter,width=7.1cm,height=2.1cm,valign=center,colback=cnum!6,colframe=cnum,arc=2.5mm,boxrule=0.9pt,left=3mm]{\Large\bfseries\color{cnum}c)}\quad\begin{minipage}[c]{5.6cm}\raggedright\Large $3$\end{minipage}\end{tcolorbox}\hfill\begin{tcolorbox}[enhanced,nobeforeafter,width=7.1cm,height=2.1cm,valign=center,colback=cnum!6,colframe=cnum,arc=2.5mm,boxrule=0.9pt,left=3mm]{\Large\bfseries\color{cnum}d)}\quad\begin{minipage}[c]{5.6cm}\raggedright\Large $3{,}4$\end{minipage}\end{tcolorbox}
\end{tcolorbox}
\newpage
\begin{tcolorbox}[enhanced,width=16cm,height=14cm,colback=white,colframe=cnum,boxrule=2pt,arc=0mm,
left=6mm,right=6mm,top=0mm,bottom=4mm,sharp corners,
overlay={\fill[cnum] (frame.north west) rectangle ([yshift=-1.25cm]frame.north east);
\node[anchor=west,text=white,font=\sffamily\bfseries\LARGE] at ([xshift=5mm,yshift=-0.62cm]frame.north west) {Question 8 D};
\node[anchor=east,text=white,font=\sffamily\itshape\large] at ([xshift=-5mm,yshift=-0.62cm]frame.north east) {Nombres entiers et décimaux};}]
\vspace*{1.75cm}
\begin{minipage}[t][6.9cm][c]{\linewidth}\LARGE\centering
\begin{minipage}{0.95\linewidth}\raggedright\linespread{1.25}\selectfont Quel est l'arrondi au dixième de $14{,}75$~?\end{minipage}
\end{minipage}

\vfill
\begin{tcolorbox}[enhanced,nobeforeafter,width=7.1cm,height=2.1cm,valign=center,colback=cnum!6,colframe=cnum,arc=2.5mm,boxrule=0.9pt,left=3mm]{\Large\bfseries\color{cnum}a)}\quad\begin{minipage}[c]{5.6cm}\raggedright\Large $14{,}7$\end{minipage}\end{tcolorbox}\hfill\begin{tcolorbox}[enhanced,nobeforeafter,width=7.1cm,height=2.1cm,valign=center,colback=cnum!6,colframe=cnum,arc=2.5mm,boxrule=0.9pt,left=3mm]{\Large\bfseries\color{cnum}b)}\quad\begin{minipage}[c]{5.6cm}\raggedright\Large $15$\end{minipage}\end{tcolorbox}\par\vspace{2.5mm}
\begin{tcolorbox}[enhanced,nobeforeafter,width=7.1cm,height=2.1cm,valign=center,colback=cnum!6,colframe=cnum,arc=2.5mm,boxrule=0.9pt,left=3mm]{\Large\bfseries\color{cnum}c)}\quad\begin{minipage}[c]{5.6cm}\raggedright\Large $14{,}75$\end{minipage}\end{tcolorbox}\hfill\begin{tcolorbox}[enhanced,nobeforeafter,width=7.1cm,height=2.1cm,valign=center,colback=cnum!6,colframe=cnum,arc=2.5mm,boxrule=0.9pt,left=3mm]{\Large\bfseries\color{cnum}d)}\quad\begin{minipage}[c]{5.6cm}\raggedright\Large $14{,}8$\end{minipage}\end{tcolorbox}
\end{tcolorbox}
\newpage
\begin{tcolorbox}[enhanced,width=16cm,height=14cm,colback=white,colframe=cnum,boxrule=2pt,arc=0mm,
left=6mm,right=6mm,top=0mm,bottom=4mm,sharp corners,
overlay={\fill[cnum] (frame.north west) rectangle ([yshift=-1.25cm]frame.north east);
\node[anchor=west,text=white,font=\sffamily\bfseries\LARGE] at ([xshift=5mm,yshift=-0.62cm]frame.north west) {Question 9 D};
\node[anchor=east,text=white,font=\sffamily\itshape\large] at ([xshift=-5mm,yshift=-0.62cm]frame.north east) {Nombres entiers et décimaux};}]
\vspace*{1.75cm}
\begin{minipage}[t][6.9cm][c]{\linewidth}\LARGE\centering
\begin{minipage}{0.95\linewidth}\raggedright\linespread{1.25}\selectfont Que vaut $7{,}09\times100$~?\end{minipage}
\end{minipage}

\vfill
\begin{tcolorbox}[enhanced,nobeforeafter,width=7.1cm,height=2.1cm,valign=center,colback=cnum!6,colframe=cnum,arc=2.5mm,boxrule=0.9pt,left=3mm]{\Large\bfseries\color{cnum}a)}\quad\begin{minipage}[c]{5.6cm}\raggedright\Large $709$\end{minipage}\end{tcolorbox}\hfill\begin{tcolorbox}[enhanced,nobeforeafter,width=7.1cm,height=2.1cm,valign=center,colback=cnum!6,colframe=cnum,arc=2.5mm,boxrule=0.9pt,left=3mm]{\Large\bfseries\color{cnum}b)}\quad\begin{minipage}[c]{5.6cm}\raggedright\Large $70\,900$\end{minipage}\end{tcolorbox}\par\vspace{2.5mm}
\begin{tcolorbox}[enhanced,nobeforeafter,width=7.1cm,height=2.1cm,valign=center,colback=cnum!6,colframe=cnum,arc=2.5mm,boxrule=0.9pt,left=3mm]{\Large\bfseries\color{cnum}c)}\quad\begin{minipage}[c]{5.6cm}\raggedright\Large $70{,}9$\end{minipage}\end{tcolorbox}\hfill\begin{tcolorbox}[enhanced,nobeforeafter,width=7.1cm,height=2.1cm,valign=center,colback=cnum!6,colframe=cnum,arc=2.5mm,boxrule=0.9pt,left=3mm]{\Large\bfseries\color{cnum}d)}\quad\begin{minipage}[c]{5.6cm}\raggedright\Large $7\,090$\end{minipage}\end{tcolorbox}
\end{tcolorbox}
\newpage
\begin{tcolorbox}[enhanced,width=16cm,height=14cm,colback=white,colframe=cnum,boxrule=2pt,arc=0mm,
left=6mm,right=6mm,top=0mm,bottom=4mm,sharp corners,
overlay={\fill[cnum] (frame.north west) rectangle ([yshift=-1.25cm]frame.north east);
\node[anchor=west,text=white,font=\sffamily\bfseries\LARGE] at ([xshift=5mm,yshift=-0.62cm]frame.north west) {Question 10 D};
\node[anchor=east,text=white,font=\sffamily\itshape\large] at ([xshift=-5mm,yshift=-0.62cm]frame.north east) {Nombres entiers et décimaux};}]
\vspace*{1.75cm}
\begin{minipage}[t][6.9cm][c]{\linewidth}\LARGE\centering
\begin{minipage}{0.95\linewidth}\raggedright\linespread{1.25}\selectfont Que vaut $0{,}015\times1\,000$~?\end{minipage}
\end{minipage}

\vfill
\begin{tcolorbox}[enhanced,nobeforeafter,width=7.1cm,height=2.1cm,valign=center,colback=cnum!6,colframe=cnum,arc=2.5mm,boxrule=0.9pt,left=3mm]{\Large\bfseries\color{cnum}a)}\quad\begin{minipage}[c]{5.6cm}\raggedright\Large $1\,500$\end{minipage}\end{tcolorbox}\hfill\begin{tcolorbox}[enhanced,nobeforeafter,width=7.1cm,height=2.1cm,valign=center,colback=cnum!6,colframe=cnum,arc=2.5mm,boxrule=0.9pt,left=3mm]{\Large\bfseries\color{cnum}b)}\quad\begin{minipage}[c]{5.6cm}\raggedright\Large $15$\end{minipage}\end{tcolorbox}\par\vspace{2.5mm}
\begin{tcolorbox}[enhanced,nobeforeafter,width=7.1cm,height=2.1cm,valign=center,colback=cnum!6,colframe=cnum,arc=2.5mm,boxrule=0.9pt,left=3mm]{\Large\bfseries\color{cnum}c)}\quad\begin{minipage}[c]{5.6cm}\raggedright\Large $150$\end{minipage}\end{tcolorbox}\hfill\begin{tcolorbox}[enhanced,nobeforeafter,width=7.1cm,height=2.1cm,valign=center,colback=cnum!6,colframe=cnum,arc=2.5mm,boxrule=0.9pt,left=3mm]{\Large\bfseries\color{cnum}d)}\quad\begin{minipage}[c]{5.6cm}\raggedright\Large $1{,}5$\end{minipage}\end{tcolorbox}
\end{tcolorbox}
\newpage
\begin{tcolorbox}[enhanced,width=16cm,height=14cm,colback=white,colframe=cnum,boxrule=2pt,arc=0mm,
left=6mm,right=6mm,top=0mm,bottom=4mm,sharp corners,
overlay={\fill[cnum] (frame.north west) rectangle ([yshift=-1.25cm]frame.north east);
\node[anchor=west,text=white,font=\sffamily\bfseries\LARGE] at ([xshift=5mm,yshift=-0.62cm]frame.north west) {Question 11 D};
\node[anchor=east,text=white,font=\sffamily\itshape\large] at ([xshift=-5mm,yshift=-0.62cm]frame.north east) {Nombres entiers et décimaux};}]
\vspace*{1.75cm}
\begin{minipage}[t][6.9cm][c]{\linewidth}\LARGE\centering
\begin{minipage}{0.95\linewidth}\raggedright\linespread{1.25}\selectfont Que vaut $4{,}2\div100$~?\end{minipage}
\end{minipage}

\vfill
\begin{tcolorbox}[enhanced,nobeforeafter,width=7.1cm,height=2.1cm,valign=center,colback=cnum!6,colframe=cnum,arc=2.5mm,boxrule=0.9pt,left=3mm]{\Large\bfseries\color{cnum}a)}\quad\begin{minipage}[c]{5.6cm}\raggedright\Large $0{,}042$\end{minipage}\end{tcolorbox}\hfill\begin{tcolorbox}[enhanced,nobeforeafter,width=7.1cm,height=2.1cm,valign=center,colback=cnum!6,colframe=cnum,arc=2.5mm,boxrule=0.9pt,left=3mm]{\Large\bfseries\color{cnum}b)}\quad\begin{minipage}[c]{5.6cm}\raggedright\Large $0{,}42$\end{minipage}\end{tcolorbox}\par\vspace{2.5mm}
\begin{tcolorbox}[enhanced,nobeforeafter,width=7.1cm,height=2.1cm,valign=center,colback=cnum!6,colframe=cnum,arc=2.5mm,boxrule=0.9pt,left=3mm]{\Large\bfseries\color{cnum}c)}\quad\begin{minipage}[c]{5.6cm}\raggedright\Large $0{,}0042$\end{minipage}\end{tcolorbox}\hfill\begin{tcolorbox}[enhanced,nobeforeafter,width=7.1cm,height=2.1cm,valign=center,colback=cnum!6,colframe=cnum,arc=2.5mm,boxrule=0.9pt,left=3mm]{\Large\bfseries\color{cnum}d)}\quad\begin{minipage}[c]{5.6cm}\raggedright\Large $420$\end{minipage}\end{tcolorbox}
\end{tcolorbox}
\newpage
\begin{tcolorbox}[enhanced,width=16cm,height=14cm,colback=white,colframe=cnum,boxrule=2pt,arc=0mm,
left=6mm,right=6mm,top=0mm,bottom=4mm,sharp corners,
overlay={\fill[cnum] (frame.north west) rectangle ([yshift=-1.25cm]frame.north east);
\node[anchor=west,text=white,font=\sffamily\bfseries\LARGE] at ([xshift=5mm,yshift=-0.62cm]frame.north west) {Question 12 D};
\node[anchor=east,text=white,font=\sffamily\itshape\large] at ([xshift=-5mm,yshift=-0.62cm]frame.north east) {Nombres entiers et décimaux};}]
\vspace*{1.75cm}
\begin{minipage}[t][6.9cm][c]{\linewidth}\LARGE\centering
\begin{minipage}{0.95\linewidth}\raggedright\linespread{1.25}\selectfont Que vaut $73\div10$~?\end{minipage}
\end{minipage}

\vfill
\begin{tcolorbox}[enhanced,nobeforeafter,width=7.1cm,height=2.1cm,valign=center,colback=cnum!6,colframe=cnum,arc=2.5mm,boxrule=0.9pt,left=3mm]{\Large\bfseries\color{cnum}a)}\quad\begin{minipage}[c]{5.6cm}\raggedright\Large $0{,}73$\end{minipage}\end{tcolorbox}\hfill\begin{tcolorbox}[enhanced,nobeforeafter,width=7.1cm,height=2.1cm,valign=center,colback=cnum!6,colframe=cnum,arc=2.5mm,boxrule=0.9pt,left=3mm]{\Large\bfseries\color{cnum}b)}\quad\begin{minipage}[c]{5.6cm}\raggedright\Large $7{,}3$\end{minipage}\end{tcolorbox}\par\vspace{2.5mm}
\begin{tcolorbox}[enhanced,nobeforeafter,width=7.1cm,height=2.1cm,valign=center,colback=cnum!6,colframe=cnum,arc=2.5mm,boxrule=0.9pt,left=3mm]{\Large\bfseries\color{cnum}c)}\quad\begin{minipage}[c]{5.6cm}\raggedright\Large $73$\end{minipage}\end{tcolorbox}\hfill\begin{tcolorbox}[enhanced,nobeforeafter,width=7.1cm,height=2.1cm,valign=center,colback=cnum!6,colframe=cnum,arc=2.5mm,boxrule=0.9pt,left=3mm]{\Large\bfseries\color{cnum}d)}\quad\begin{minipage}[c]{5.6cm}\raggedright\Large $730$\end{minipage}\end{tcolorbox}
\end{tcolorbox}
\newpage
\begin{tcolorbox}[enhanced,width=16cm,height=14cm,colback=white,colframe=cnum,boxrule=2pt,arc=0mm,
left=6mm,right=6mm,top=0mm,bottom=4mm,sharp corners,
overlay={\fill[cnum] (frame.north west) rectangle ([yshift=-1.25cm]frame.north east);
\node[anchor=west,text=white,font=\sffamily\bfseries\LARGE] at ([xshift=5mm,yshift=-0.62cm]frame.north west) {Question 13 D};
\node[anchor=east,text=white,font=\sffamily\itshape\large] at ([xshift=-5mm,yshift=-0.62cm]frame.north east) {Nombres entiers et décimaux};}]
\vspace*{1.75cm}
\begin{minipage}[t][6.9cm][c]{\linewidth}\LARGE\centering
\begin{minipage}{0.95\linewidth}\raggedright\linespread{1.25}\selectfont Comment s'écrit en chiffres le nombre «~un million deux cent cinquante mille~»~?\end{minipage}
\end{minipage}

\vfill
\begin{tcolorbox}[enhanced,nobeforeafter,width=7.1cm,height=2.1cm,valign=center,colback=cnum!6,colframe=cnum,arc=2.5mm,boxrule=0.9pt,left=3mm]{\Large\bfseries\color{cnum}a)}\quad\begin{minipage}[c]{5.6cm}\raggedright\Large $1\,025\,000$\end{minipage}\end{tcolorbox}\hfill\begin{tcolorbox}[enhanced,nobeforeafter,width=7.1cm,height=2.1cm,valign=center,colback=cnum!6,colframe=cnum,arc=2.5mm,boxrule=0.9pt,left=3mm]{\Large\bfseries\color{cnum}b)}\quad\begin{minipage}[c]{5.6cm}\raggedright\Large $1\,205\,000$\end{minipage}\end{tcolorbox}\par\vspace{2.5mm}
\begin{tcolorbox}[enhanced,nobeforeafter,width=7.1cm,height=2.1cm,valign=center,colback=cnum!6,colframe=cnum,arc=2.5mm,boxrule=0.9pt,left=3mm]{\Large\bfseries\color{cnum}c)}\quad\begin{minipage}[c]{5.6cm}\raggedright\Large $125\,000$\end{minipage}\end{tcolorbox}\hfill\begin{tcolorbox}[enhanced,nobeforeafter,width=7.1cm,height=2.1cm,valign=center,colback=cnum!6,colframe=cnum,arc=2.5mm,boxrule=0.9pt,left=3mm]{\Large\bfseries\color{cnum}d)}\quad\begin{minipage}[c]{5.6cm}\raggedright\Large $1\,250\,000$\end{minipage}\end{tcolorbox}
\end{tcolorbox}
\newpage
\begin{tcolorbox}[enhanced,width=16cm,height=14cm,colback=white,colframe=cnum,boxrule=2pt,arc=0mm,
left=6mm,right=6mm,top=0mm,bottom=4mm,sharp corners,
overlay={\fill[cnum] (frame.north west) rectangle ([yshift=-1.25cm]frame.north east);
\node[anchor=west,text=white,font=\sffamily\bfseries\LARGE] at ([xshift=5mm,yshift=-0.62cm]frame.north west) {Question 14 D};
\node[anchor=east,text=white,font=\sffamily\itshape\large] at ([xshift=-5mm,yshift=-0.62cm]frame.north east) {Nombres entiers et décimaux};}]
\vspace*{1.75cm}
\begin{minipage}[t][6.9cm][c]{\linewidth}\LARGE\centering
\begin{minipage}{0.95\linewidth}\raggedright\linespread{1.25}\selectfont Comment s'écrit en chiffres le nombre «~quarante-huit mille neuf cents~»~?\end{minipage}
\end{minipage}

\vfill
\begin{tcolorbox}[enhanced,nobeforeafter,width=7.1cm,height=2.1cm,valign=center,colback=cnum!6,colframe=cnum,arc=2.5mm,boxrule=0.9pt,left=3mm]{\Large\bfseries\color{cnum}a)}\quad\begin{minipage}[c]{5.6cm}\raggedright\Large $489\,000$\end{minipage}\end{tcolorbox}\hfill\begin{tcolorbox}[enhanced,nobeforeafter,width=7.1cm,height=2.1cm,valign=center,colback=cnum!6,colframe=cnum,arc=2.5mm,boxrule=0.9pt,left=3mm]{\Large\bfseries\color{cnum}b)}\quad\begin{minipage}[c]{5.6cm}\raggedright\Large $4\,890$\end{minipage}\end{tcolorbox}\par\vspace{2.5mm}
\begin{tcolorbox}[enhanced,nobeforeafter,width=7.1cm,height=2.1cm,valign=center,colback=cnum!6,colframe=cnum,arc=2.5mm,boxrule=0.9pt,left=3mm]{\Large\bfseries\color{cnum}c)}\quad\begin{minipage}[c]{5.6cm}\raggedright\Large $48\,090$\end{minipage}\end{tcolorbox}\hfill\begin{tcolorbox}[enhanced,nobeforeafter,width=7.1cm,height=2.1cm,valign=center,colback=cnum!6,colframe=cnum,arc=2.5mm,boxrule=0.9pt,left=3mm]{\Large\bfseries\color{cnum}d)}\quad\begin{minipage}[c]{5.6cm}\raggedright\Large $48\,900$\end{minipage}\end{tcolorbox}
\end{tcolorbox}
\newpage
\begin{tcolorbox}[enhanced,width=16cm,height=14cm,colback=white,colframe=cnum,boxrule=2pt,arc=0mm,
left=6mm,right=6mm,top=0mm,bottom=4mm,sharp corners,
overlay={\fill[cnum] (frame.north west) rectangle ([yshift=-1.25cm]frame.north east);
\node[anchor=west,text=white,font=\sffamily\bfseries\LARGE] at ([xshift=5mm,yshift=-0.62cm]frame.north west) {Question 15 D};
\node[anchor=east,text=white,font=\sffamily\itshape\large] at ([xshift=-5mm,yshift=-0.62cm]frame.north east) {Nombres entiers et décimaux};}]
\vspace*{1.75cm}
\begin{minipage}[t][6.9cm][c]{\linewidth}\LARGE\centering
\begin{minipage}{0.95\linewidth}\raggedright\linespread{1.25}\selectfont Quel nombre est compris entre $1{,}4$ et $1{,}5$~?\end{minipage}
\end{minipage}

\vfill
\begin{tcolorbox}[enhanced,nobeforeafter,width=7.1cm,height=2.1cm,valign=center,colback=cnum!6,colframe=cnum,arc=2.5mm,boxrule=0.9pt,left=3mm]{\Large\bfseries\color{cnum}a)}\quad\begin{minipage}[c]{5.6cm}\raggedright\Large $1{,}52$\end{minipage}\end{tcolorbox}\hfill\begin{tcolorbox}[enhanced,nobeforeafter,width=7.1cm,height=2.1cm,valign=center,colback=cnum!6,colframe=cnum,arc=2.5mm,boxrule=0.9pt,left=3mm]{\Large\bfseries\color{cnum}b)}\quad\begin{minipage}[c]{5.6cm}\raggedright\Large $1{,}42$\end{minipage}\end{tcolorbox}\par\vspace{2.5mm}
\begin{tcolorbox}[enhanced,nobeforeafter,width=7.1cm,height=2.1cm,valign=center,colback=cnum!6,colframe=cnum,arc=2.5mm,boxrule=0.9pt,left=3mm]{\Large\bfseries\color{cnum}c)}\quad\begin{minipage}[c]{5.6cm}\raggedright\Large $1{,}04$\end{minipage}\end{tcolorbox}\hfill\begin{tcolorbox}[enhanced,nobeforeafter,width=7.1cm,height=2.1cm,valign=center,colback=cnum!6,colframe=cnum,arc=2.5mm,boxrule=0.9pt,left=3mm]{\Large\bfseries\color{cnum}d)}\quad\begin{minipage}[c]{5.6cm}\raggedright\Large $1{,}39$\end{minipage}\end{tcolorbox}
\end{tcolorbox}
\newpage
\begin{tcolorbox}[enhanced,width=16cm,height=14cm,colback=white,colframe=cnum,boxrule=2pt,arc=0mm,
left=6mm,right=6mm,top=0mm,bottom=4mm,sharp corners,
overlay={\fill[cnum] (frame.north west) rectangle ([yshift=-1.25cm]frame.north east);
\node[anchor=west,text=white,font=\sffamily\bfseries\LARGE] at ([xshift=5mm,yshift=-0.62cm]frame.north west) {Question 16 D};
\node[anchor=east,text=white,font=\sffamily\itshape\large] at ([xshift=-5mm,yshift=-0.62cm]frame.north east) {Nombres entiers et décimaux};}]
\vspace*{1.75cm}
\begin{minipage}[t][6.9cm][c]{\linewidth}\LARGE\centering
\begin{minipage}{0.95\linewidth}\raggedright\linespread{1.25}\selectfont Quel nombre est compris entre $0{,}1$ et $0{,}2$~?\end{minipage}
\end{minipage}

\vfill
\begin{tcolorbox}[enhanced,nobeforeafter,width=7.1cm,height=2.1cm,valign=center,colback=cnum!6,colframe=cnum,arc=2.5mm,boxrule=0.9pt,left=3mm]{\Large\bfseries\color{cnum}a)}\quad\begin{minipage}[c]{5.6cm}\raggedright\Large $0{,}015$\end{minipage}\end{tcolorbox}\hfill\begin{tcolorbox}[enhanced,nobeforeafter,width=7.1cm,height=2.1cm,valign=center,colback=cnum!6,colframe=cnum,arc=2.5mm,boxrule=0.9pt,left=3mm]{\Large\bfseries\color{cnum}b)}\quad\begin{minipage}[c]{5.6cm}\raggedright\Large $0{,}25$\end{minipage}\end{tcolorbox}\par\vspace{2.5mm}
\begin{tcolorbox}[enhanced,nobeforeafter,width=7.1cm,height=2.1cm,valign=center,colback=cnum!6,colframe=cnum,arc=2.5mm,boxrule=0.9pt,left=3mm]{\Large\bfseries\color{cnum}c)}\quad\begin{minipage}[c]{5.6cm}\raggedright\Large $0{,}15$\end{minipage}\end{tcolorbox}\hfill\begin{tcolorbox}[enhanced,nobeforeafter,width=7.1cm,height=2.1cm,valign=center,colback=cnum!6,colframe=cnum,arc=2.5mm,boxrule=0.9pt,left=3mm]{\Large\bfseries\color{cnum}d)}\quad\begin{minipage}[c]{5.6cm}\raggedright\Large $0{,}21$\end{minipage}\end{tcolorbox}
\end{tcolorbox}
\newpage
\begin{tcolorbox}[enhanced,width=16cm,height=14cm,colback=white,colframe=ccal,boxrule=2pt,arc=0mm,
left=6mm,right=6mm,top=0mm,bottom=4mm,sharp corners,
overlay={\fill[ccal] (frame.north west) rectangle ([yshift=-1.25cm]frame.north east);
\node[anchor=west,text=white,font=\sffamily\bfseries\LARGE] at ([xshift=5mm,yshift=-0.62cm]frame.north west) {Question 17 D};
\node[anchor=east,text=white,font=\sffamily\itshape\large] at ([xshift=-5mm,yshift=-0.62cm]frame.north east) {Calculs};}]
\vspace*{1.75cm}
\begin{minipage}[t][6.9cm][c]{\linewidth}\LARGE\centering
\begin{minipage}{0.95\linewidth}\raggedright\linespread{1.25}\selectfont Que vaut $6{,}08+3{,}9$~?\end{minipage}
\end{minipage}

\vfill
\begin{tcolorbox}[enhanced,nobeforeafter,width=7.1cm,height=2.1cm,valign=center,colback=ccal!6,colframe=ccal,arc=2.5mm,boxrule=0.9pt,left=3mm]{\Large\bfseries\color{ccal}a)}\quad\begin{minipage}[c]{5.6cm}\raggedright\Large $10{,}08$\end{minipage}\end{tcolorbox}\hfill\begin{tcolorbox}[enhanced,nobeforeafter,width=7.1cm,height=2.1cm,valign=center,colback=ccal!6,colframe=ccal,arc=2.5mm,boxrule=0.9pt,left=3mm]{\Large\bfseries\color{ccal}b)}\quad\begin{minipage}[c]{5.6cm}\raggedright\Large $9{,}98$\end{minipage}\end{tcolorbox}\par\vspace{2.5mm}
\begin{tcolorbox}[enhanced,nobeforeafter,width=7.1cm,height=2.1cm,valign=center,colback=ccal!6,colframe=ccal,arc=2.5mm,boxrule=0.9pt,left=3mm]{\Large\bfseries\color{ccal}c)}\quad\begin{minipage}[c]{5.6cm}\raggedright\Large $10{,}98$\end{minipage}\end{tcolorbox}\hfill\begin{tcolorbox}[enhanced,nobeforeafter,width=7.1cm,height=2.1cm,valign=center,colback=ccal!6,colframe=ccal,arc=2.5mm,boxrule=0.9pt,left=3mm]{\Large\bfseries\color{ccal}d)}\quad\begin{minipage}[c]{5.6cm}\raggedright\Large $8{,}98$\end{minipage}\end{tcolorbox}
\end{tcolorbox}
\newpage
\begin{tcolorbox}[enhanced,width=16cm,height=14cm,colback=white,colframe=ccal,boxrule=2pt,arc=0mm,
left=6mm,right=6mm,top=0mm,bottom=4mm,sharp corners,
overlay={\fill[ccal] (frame.north west) rectangle ([yshift=-1.25cm]frame.north east);
\node[anchor=west,text=white,font=\sffamily\bfseries\LARGE] at ([xshift=5mm,yshift=-0.62cm]frame.north west) {Question 18 D};
\node[anchor=east,text=white,font=\sffamily\itshape\large] at ([xshift=-5mm,yshift=-0.62cm]frame.north east) {Calculs};}]
\vspace*{1.75cm}
\begin{minipage}[t][6.9cm][c]{\linewidth}\LARGE\centering
\begin{minipage}{0.95\linewidth}\raggedright\linespread{1.25}\selectfont Que vaut $45+2{,}35$~?\end{minipage}
\end{minipage}

\vfill
\begin{tcolorbox}[enhanced,nobeforeafter,width=7.1cm,height=2.1cm,valign=center,colback=ccal!6,colframe=ccal,arc=2.5mm,boxrule=0.9pt,left=3mm]{\Large\bfseries\color{ccal}a)}\quad\begin{minipage}[c]{5.6cm}\raggedright\Large $48{,}35$\end{minipage}\end{tcolorbox}\hfill\begin{tcolorbox}[enhanced,nobeforeafter,width=7.1cm,height=2.1cm,valign=center,colback=ccal!6,colframe=ccal,arc=2.5mm,boxrule=0.9pt,left=3mm]{\Large\bfseries\color{ccal}b)}\quad\begin{minipage}[c]{5.6cm}\raggedright\Large $46{,}35$\end{minipage}\end{tcolorbox}\par\vspace{2.5mm}
\begin{tcolorbox}[enhanced,nobeforeafter,width=7.1cm,height=2.1cm,valign=center,colback=ccal!6,colframe=ccal,arc=2.5mm,boxrule=0.9pt,left=3mm]{\Large\bfseries\color{ccal}c)}\quad\begin{minipage}[c]{5.6cm}\raggedright\Large $47{,}45$\end{minipage}\end{tcolorbox}\hfill\begin{tcolorbox}[enhanced,nobeforeafter,width=7.1cm,height=2.1cm,valign=center,colback=ccal!6,colframe=ccal,arc=2.5mm,boxrule=0.9pt,left=3mm]{\Large\bfseries\color{ccal}d)}\quad\begin{minipage}[c]{5.6cm}\raggedright\Large $47{,}35$\end{minipage}\end{tcolorbox}
\end{tcolorbox}
\newpage
\begin{tcolorbox}[enhanced,width=16cm,height=14cm,colback=white,colframe=ccal,boxrule=2pt,arc=0mm,
left=6mm,right=6mm,top=0mm,bottom=4mm,sharp corners,
overlay={\fill[ccal] (frame.north west) rectangle ([yshift=-1.25cm]frame.north east);
\node[anchor=west,text=white,font=\sffamily\bfseries\LARGE] at ([xshift=5mm,yshift=-0.62cm]frame.north west) {Question 19 D};
\node[anchor=east,text=white,font=\sffamily\itshape\large] at ([xshift=-5mm,yshift=-0.62cm]frame.north east) {Calculs};}]
\vspace*{1.75cm}
\begin{minipage}[t][6.9cm][c]{\linewidth}\LARGE\centering
\begin{minipage}{0.95\linewidth}\raggedright\linespread{1.25}\selectfont Que vaut $4-0{,}85$~?\end{minipage}
\end{minipage}

\vfill
\begin{tcolorbox}[enhanced,nobeforeafter,width=7.1cm,height=2.1cm,valign=center,colback=ccal!6,colframe=ccal,arc=2.5mm,boxrule=0.9pt,left=3mm]{\Large\bfseries\color{ccal}a)}\quad\begin{minipage}[c]{5.6cm}\raggedright\Large $3{,}25$\end{minipage}\end{tcolorbox}\hfill\begin{tcolorbox}[enhanced,nobeforeafter,width=7.1cm,height=2.1cm,valign=center,colback=ccal!6,colframe=ccal,arc=2.5mm,boxrule=0.9pt,left=3mm]{\Large\bfseries\color{ccal}b)}\quad\begin{minipage}[c]{5.6cm}\raggedright\Large $3{,}05$\end{minipage}\end{tcolorbox}\par\vspace{2.5mm}
\begin{tcolorbox}[enhanced,nobeforeafter,width=7.1cm,height=2.1cm,valign=center,colback=ccal!6,colframe=ccal,arc=2.5mm,boxrule=0.9pt,left=3mm]{\Large\bfseries\color{ccal}c)}\quad\begin{minipage}[c]{5.6cm}\raggedright\Large $3{,}15$\end{minipage}\end{tcolorbox}\hfill\begin{tcolorbox}[enhanced,nobeforeafter,width=7.1cm,height=2.1cm,valign=center,colback=ccal!6,colframe=ccal,arc=2.5mm,boxrule=0.9pt,left=3mm]{\Large\bfseries\color{ccal}d)}\quad\begin{minipage}[c]{5.6cm}\raggedright\Large $4{,}15$\end{minipage}\end{tcolorbox}
\end{tcolorbox}
\newpage
\begin{tcolorbox}[enhanced,width=16cm,height=14cm,colback=white,colframe=ccal,boxrule=2pt,arc=0mm,
left=6mm,right=6mm,top=0mm,bottom=4mm,sharp corners,
overlay={\fill[ccal] (frame.north west) rectangle ([yshift=-1.25cm]frame.north east);
\node[anchor=west,text=white,font=\sffamily\bfseries\LARGE] at ([xshift=5mm,yshift=-0.62cm]frame.north west) {Question 20 D};
\node[anchor=east,text=white,font=\sffamily\itshape\large] at ([xshift=-5mm,yshift=-0.62cm]frame.north east) {Calculs};}]
\vspace*{1.75cm}
\begin{minipage}[t][6.9cm][c]{\linewidth}\LARGE\centering
\begin{minipage}{0.95\linewidth}\raggedright\linespread{1.25}\selectfont Que vaut $13{,}06-5{,}9$~?\end{minipage}
\end{minipage}

\vfill
\begin{tcolorbox}[enhanced,nobeforeafter,width=7.1cm,height=2.1cm,valign=center,colback=ccal!6,colframe=ccal,arc=2.5mm,boxrule=0.9pt,left=3mm]{\Large\bfseries\color{ccal}a)}\quad\begin{minipage}[c]{5.6cm}\raggedright\Large $7{,}26$\end{minipage}\end{tcolorbox}\hfill\begin{tcolorbox}[enhanced,nobeforeafter,width=7.1cm,height=2.1cm,valign=center,colback=ccal!6,colframe=ccal,arc=2.5mm,boxrule=0.9pt,left=3mm]{\Large\bfseries\color{ccal}b)}\quad\begin{minipage}[c]{5.6cm}\raggedright\Large $8{,}16$\end{minipage}\end{tcolorbox}\par\vspace{2.5mm}
\begin{tcolorbox}[enhanced,nobeforeafter,width=7.1cm,height=2.1cm,valign=center,colback=ccal!6,colframe=ccal,arc=2.5mm,boxrule=0.9pt,left=3mm]{\Large\bfseries\color{ccal}c)}\quad\begin{minipage}[c]{5.6cm}\raggedright\Large $7{,}06$\end{minipage}\end{tcolorbox}\hfill\begin{tcolorbox}[enhanced,nobeforeafter,width=7.1cm,height=2.1cm,valign=center,colback=ccal!6,colframe=ccal,arc=2.5mm,boxrule=0.9pt,left=3mm]{\Large\bfseries\color{ccal}d)}\quad\begin{minipage}[c]{5.6cm}\raggedright\Large $7{,}16$\end{minipage}\end{tcolorbox}
\end{tcolorbox}
\newpage
\begin{tcolorbox}[enhanced,width=16cm,height=14cm,colback=white,colframe=ccal,boxrule=2pt,arc=0mm,
left=6mm,right=6mm,top=0mm,bottom=4mm,sharp corners,
overlay={\fill[ccal] (frame.north west) rectangle ([yshift=-1.25cm]frame.north east);
\node[anchor=west,text=white,font=\sffamily\bfseries\LARGE] at ([xshift=5mm,yshift=-0.62cm]frame.north west) {Question 21 D};
\node[anchor=east,text=white,font=\sffamily\itshape\large] at ([xshift=-5mm,yshift=-0.62cm]frame.north east) {Calculs};}]
\vspace*{1.75cm}
\begin{minipage}[t][6.9cm][c]{\linewidth}\LARGE\centering
\begin{minipage}{0.95\linewidth}\raggedright\linespread{1.25}\selectfont Que vaut $12\times0{,}75$~?\end{minipage}
\end{minipage}

\vfill
\begin{tcolorbox}[enhanced,nobeforeafter,width=7.1cm,height=2.1cm,valign=center,colback=ccal!6,colframe=ccal,arc=2.5mm,boxrule=0.9pt,left=3mm]{\Large\bfseries\color{ccal}a)}\quad\begin{minipage}[c]{5.6cm}\raggedright\Large $12{,}75$\end{minipage}\end{tcolorbox}\hfill\begin{tcolorbox}[enhanced,nobeforeafter,width=7.1cm,height=2.1cm,valign=center,colback=ccal!6,colframe=ccal,arc=2.5mm,boxrule=0.9pt,left=3mm]{\Large\bfseries\color{ccal}b)}\quad\begin{minipage}[c]{5.6cm}\raggedright\Large $9$\end{minipage}\end{tcolorbox}\par\vspace{2.5mm}
\begin{tcolorbox}[enhanced,nobeforeafter,width=7.1cm,height=2.1cm,valign=center,colback=ccal!6,colframe=ccal,arc=2.5mm,boxrule=0.9pt,left=3mm]{\Large\bfseries\color{ccal}c)}\quad\begin{minipage}[c]{5.6cm}\raggedright\Large $90$\end{minipage}\end{tcolorbox}\hfill\begin{tcolorbox}[enhanced,nobeforeafter,width=7.1cm,height=2.1cm,valign=center,colback=ccal!6,colframe=ccal,arc=2.5mm,boxrule=0.9pt,left=3mm]{\Large\bfseries\color{ccal}d)}\quad\begin{minipage}[c]{5.6cm}\raggedright\Large $0{,}9$\end{minipage}\end{tcolorbox}
\end{tcolorbox}
\newpage
\begin{tcolorbox}[enhanced,width=16cm,height=14cm,colback=white,colframe=ccal,boxrule=2pt,arc=0mm,
left=6mm,right=6mm,top=0mm,bottom=4mm,sharp corners,
overlay={\fill[ccal] (frame.north west) rectangle ([yshift=-1.25cm]frame.north east);
\node[anchor=west,text=white,font=\sffamily\bfseries\LARGE] at ([xshift=5mm,yshift=-0.62cm]frame.north west) {Question 22 D};
\node[anchor=east,text=white,font=\sffamily\itshape\large] at ([xshift=-5mm,yshift=-0.62cm]frame.north east) {Calculs};}]
\vspace*{1.75cm}
\begin{minipage}[t][6.9cm][c]{\linewidth}\LARGE\centering
\begin{minipage}{0.95\linewidth}\raggedright\linespread{1.25}\selectfont Que vaut $0{,}6\times0{,}07$~?\end{minipage}
\end{minipage}

\vfill
\begin{tcolorbox}[enhanced,nobeforeafter,width=7.1cm,height=2.1cm,valign=center,colback=ccal!6,colframe=ccal,arc=2.5mm,boxrule=0.9pt,left=3mm]{\Large\bfseries\color{ccal}a)}\quad\begin{minipage}[c]{5.6cm}\raggedright\Large $0{,}67$\end{minipage}\end{tcolorbox}\hfill\begin{tcolorbox}[enhanced,nobeforeafter,width=7.1cm,height=2.1cm,valign=center,colback=ccal!6,colframe=ccal,arc=2.5mm,boxrule=0.9pt,left=3mm]{\Large\bfseries\color{ccal}b)}\quad\begin{minipage}[c]{5.6cm}\raggedright\Large $0{,}0042$\end{minipage}\end{tcolorbox}\par\vspace{2.5mm}
\begin{tcolorbox}[enhanced,nobeforeafter,width=7.1cm,height=2.1cm,valign=center,colback=ccal!6,colframe=ccal,arc=2.5mm,boxrule=0.9pt,left=3mm]{\Large\bfseries\color{ccal}c)}\quad\begin{minipage}[c]{5.6cm}\raggedright\Large $0{,}42$\end{minipage}\end{tcolorbox}\hfill\begin{tcolorbox}[enhanced,nobeforeafter,width=7.1cm,height=2.1cm,valign=center,colback=ccal!6,colframe=ccal,arc=2.5mm,boxrule=0.9pt,left=3mm]{\Large\bfseries\color{ccal}d)}\quad\begin{minipage}[c]{5.6cm}\raggedright\Large $0{,}042$\end{minipage}\end{tcolorbox}
\end{tcolorbox}
\newpage
\begin{tcolorbox}[enhanced,width=16cm,height=14cm,colback=white,colframe=ccal,boxrule=2pt,arc=0mm,
left=6mm,right=6mm,top=0mm,bottom=4mm,sharp corners,
overlay={\fill[ccal] (frame.north west) rectangle ([yshift=-1.25cm]frame.north east);
\node[anchor=west,text=white,font=\sffamily\bfseries\LARGE] at ([xshift=5mm,yshift=-0.62cm]frame.north west) {Question 23 D};
\node[anchor=east,text=white,font=\sffamily\itshape\large] at ([xshift=-5mm,yshift=-0.62cm]frame.north east) {Calculs};}]
\vspace*{1.75cm}
\begin{minipage}[t][6.9cm][c]{\linewidth}\LARGE\centering
\begin{minipage}{0.95\linewidth}\raggedright\linespread{1.25}\selectfont Dans la division euclidienne de $365$ par $7$, quels sont le quotient $q$ et le reste $r$~?\end{minipage}
\end{minipage}

\vfill
\begin{tcolorbox}[enhanced,nobeforeafter,width=7.1cm,height=2.1cm,valign=center,colback=ccal!6,colframe=ccal,arc=2.5mm,boxrule=0.9pt,left=3mm]{\Large\bfseries\color{ccal}a)}\quad\begin{minipage}[c]{5.6cm}\raggedright\Large $q=51\text{ et }r=8$\end{minipage}\end{tcolorbox}\hfill\begin{tcolorbox}[enhanced,nobeforeafter,width=7.1cm,height=2.1cm,valign=center,colback=ccal!6,colframe=ccal,arc=2.5mm,boxrule=0.9pt,left=3mm]{\Large\bfseries\color{ccal}b)}\quad\begin{minipage}[c]{5.6cm}\raggedright\Large $q=52\text{ et }r=1$\end{minipage}\end{tcolorbox}\par\vspace{2.5mm}
\begin{tcolorbox}[enhanced,nobeforeafter,width=7.1cm,height=2.1cm,valign=center,colback=ccal!6,colframe=ccal,arc=2.5mm,boxrule=0.9pt,left=3mm]{\Large\bfseries\color{ccal}c)}\quad\begin{minipage}[c]{5.6cm}\raggedright\Large $q=52\text{ et }r=2$\end{minipage}\end{tcolorbox}\hfill\begin{tcolorbox}[enhanced,nobeforeafter,width=7.1cm,height=2.1cm,valign=center,colback=ccal!6,colframe=ccal,arc=2.5mm,boxrule=0.9pt,left=3mm]{\Large\bfseries\color{ccal}d)}\quad\begin{minipage}[c]{5.6cm}\raggedright\Large $q=53\text{ et }r=1$\end{minipage}\end{tcolorbox}
\end{tcolorbox}
\newpage
\begin{tcolorbox}[enhanced,width=16cm,height=14cm,colback=white,colframe=ccal,boxrule=2pt,arc=0mm,
left=6mm,right=6mm,top=0mm,bottom=4mm,sharp corners,
overlay={\fill[ccal] (frame.north west) rectangle ([yshift=-1.25cm]frame.north east);
\node[anchor=west,text=white,font=\sffamily\bfseries\LARGE] at ([xshift=5mm,yshift=-0.62cm]frame.north west) {Question 24 D};
\node[anchor=east,text=white,font=\sffamily\itshape\large] at ([xshift=-5mm,yshift=-0.62cm]frame.north east) {Calculs};}]
\vspace*{1.75cm}
\begin{minipage}[t][6.9cm][c]{\linewidth}\LARGE\centering
\begin{minipage}{0.95\linewidth}\raggedright\linespread{1.25}\selectfont Dans la division euclidienne de $1\,000$ par $24$, quels sont le quotient $q$ et le reste $r$~?\end{minipage}
\end{minipage}

\vfill
\begin{tcolorbox}[enhanced,nobeforeafter,width=7.1cm,height=2.1cm,valign=center,colback=ccal!6,colframe=ccal,arc=2.5mm,boxrule=0.9pt,left=3mm]{\Large\bfseries\color{ccal}a)}\quad\begin{minipage}[c]{5.6cm}\raggedright\Large $q=41\text{ et }r=17$\end{minipage}\end{tcolorbox}\hfill\begin{tcolorbox}[enhanced,nobeforeafter,width=7.1cm,height=2.1cm,valign=center,colback=ccal!6,colframe=ccal,arc=2.5mm,boxrule=0.9pt,left=3mm]{\Large\bfseries\color{ccal}b)}\quad\begin{minipage}[c]{5.6cm}\raggedright\Large $q=40\text{ et }r=40$\end{minipage}\end{tcolorbox}\par\vspace{2.5mm}
\begin{tcolorbox}[enhanced,nobeforeafter,width=7.1cm,height=2.1cm,valign=center,colback=ccal!6,colframe=ccal,arc=2.5mm,boxrule=0.9pt,left=3mm]{\Large\bfseries\color{ccal}c)}\quad\begin{minipage}[c]{5.6cm}\raggedright\Large $q=42\text{ et }r=16$\end{minipage}\end{tcolorbox}\hfill\begin{tcolorbox}[enhanced,nobeforeafter,width=7.1cm,height=2.1cm,valign=center,colback=ccal!6,colframe=ccal,arc=2.5mm,boxrule=0.9pt,left=3mm]{\Large\bfseries\color{ccal}d)}\quad\begin{minipage}[c]{5.6cm}\raggedright\Large $q=41\text{ et }r=16$\end{minipage}\end{tcolorbox}
\end{tcolorbox}
\newpage
\begin{tcolorbox}[enhanced,width=16cm,height=14cm,colback=white,colframe=ccal,boxrule=2pt,arc=0mm,
left=6mm,right=6mm,top=0mm,bottom=4mm,sharp corners,
overlay={\fill[ccal] (frame.north west) rectangle ([yshift=-1.25cm]frame.north east);
\node[anchor=west,text=white,font=\sffamily\bfseries\LARGE] at ([xshift=5mm,yshift=-0.62cm]frame.north west) {Question 25 D};
\node[anchor=east,text=white,font=\sffamily\itshape\large] at ([xshift=-5mm,yshift=-0.62cm]frame.north east) {Calculs};}]
\vspace*{1.75cm}
\begin{minipage}[t][6.9cm][c]{\linewidth}\LARGE\centering
\begin{minipage}{0.95\linewidth}\raggedright\linespread{1.25}\selectfont Lequel de ces nombres est un multiple de $11$~?\end{minipage}
\end{minipage}

\vfill
\begin{tcolorbox}[enhanced,nobeforeafter,width=7.1cm,height=2.1cm,valign=center,colback=ccal!6,colframe=ccal,arc=2.5mm,boxrule=0.9pt,left=3mm]{\Large\bfseries\color{ccal}a)}\quad\begin{minipage}[c]{5.6cm}\raggedright\Large $133$\end{minipage}\end{tcolorbox}\hfill\begin{tcolorbox}[enhanced,nobeforeafter,width=7.1cm,height=2.1cm,valign=center,colback=ccal!6,colframe=ccal,arc=2.5mm,boxrule=0.9pt,left=3mm]{\Large\bfseries\color{ccal}b)}\quad\begin{minipage}[c]{5.6cm}\raggedright\Large $143$\end{minipage}\end{tcolorbox}\par\vspace{2.5mm}
\begin{tcolorbox}[enhanced,nobeforeafter,width=7.1cm,height=2.1cm,valign=center,colback=ccal!6,colframe=ccal,arc=2.5mm,boxrule=0.9pt,left=3mm]{\Large\bfseries\color{ccal}c)}\quad\begin{minipage}[c]{5.6cm}\raggedright\Large $141$\end{minipage}\end{tcolorbox}\hfill\begin{tcolorbox}[enhanced,nobeforeafter,width=7.1cm,height=2.1cm,valign=center,colback=ccal!6,colframe=ccal,arc=2.5mm,boxrule=0.9pt,left=3mm]{\Large\bfseries\color{ccal}d)}\quad\begin{minipage}[c]{5.6cm}\raggedright\Large $153$\end{minipage}\end{tcolorbox}
\end{tcolorbox}
\newpage
\begin{tcolorbox}[enhanced,width=16cm,height=14cm,colback=white,colframe=ccal,boxrule=2pt,arc=0mm,
left=6mm,right=6mm,top=0mm,bottom=4mm,sharp corners,
overlay={\fill[ccal] (frame.north west) rectangle ([yshift=-1.25cm]frame.north east);
\node[anchor=west,text=white,font=\sffamily\bfseries\LARGE] at ([xshift=5mm,yshift=-0.62cm]frame.north west) {Question 26 D};
\node[anchor=east,text=white,font=\sffamily\itshape\large] at ([xshift=-5mm,yshift=-0.62cm]frame.north east) {Calculs};}]
\vspace*{1.75cm}
\begin{minipage}[t][6.9cm][c]{\linewidth}\LARGE\centering
\begin{minipage}{0.95\linewidth}\raggedright\linespread{1.25}\selectfont Lequel de ces nombres est un multiple de $15$~?\end{minipage}
\end{minipage}

\vfill
\begin{tcolorbox}[enhanced,nobeforeafter,width=7.1cm,height=2.1cm,valign=center,colback=ccal!6,colframe=ccal,arc=2.5mm,boxrule=0.9pt,left=3mm]{\Large\bfseries\color{ccal}a)}\quad\begin{minipage}[c]{5.6cm}\raggedright\Large $165$\end{minipage}\end{tcolorbox}\hfill\begin{tcolorbox}[enhanced,nobeforeafter,width=7.1cm,height=2.1cm,valign=center,colback=ccal!6,colframe=ccal,arc=2.5mm,boxrule=0.9pt,left=3mm]{\Large\bfseries\color{ccal}b)}\quad\begin{minipage}[c]{5.6cm}\raggedright\Large $175$\end{minipage}\end{tcolorbox}\par\vspace{2.5mm}
\begin{tcolorbox}[enhanced,nobeforeafter,width=7.1cm,height=2.1cm,valign=center,colback=ccal!6,colframe=ccal,arc=2.5mm,boxrule=0.9pt,left=3mm]{\Large\bfseries\color{ccal}c)}\quad\begin{minipage}[c]{5.6cm}\raggedright\Large $155$\end{minipage}\end{tcolorbox}\hfill\begin{tcolorbox}[enhanced,nobeforeafter,width=7.1cm,height=2.1cm,valign=center,colback=ccal!6,colframe=ccal,arc=2.5mm,boxrule=0.9pt,left=3mm]{\Large\bfseries\color{ccal}d)}\quad\begin{minipage}[c]{5.6cm}\raggedright\Large $160$\end{minipage}\end{tcolorbox}
\end{tcolorbox}
\newpage
\begin{tcolorbox}[enhanced,width=16cm,height=14cm,colback=white,colframe=ccal,boxrule=2pt,arc=0mm,
left=6mm,right=6mm,top=0mm,bottom=4mm,sharp corners,
overlay={\fill[ccal] (frame.north west) rectangle ([yshift=-1.25cm]frame.north east);
\node[anchor=west,text=white,font=\sffamily\bfseries\LARGE] at ([xshift=5mm,yshift=-0.62cm]frame.north west) {Question 27 D};
\node[anchor=east,text=white,font=\sffamily\itshape\large] at ([xshift=-5mm,yshift=-0.62cm]frame.north east) {Calculs};}]
\vspace*{1.75cm}
\begin{minipage}[t][6.9cm][c]{\linewidth}\LARGE\centering
\begin{minipage}{0.95\linewidth}\raggedright\linespread{1.25}\selectfont Lequel de ces nombres est divisible par $9$~?\end{minipage}
\end{minipage}

\vfill
\begin{tcolorbox}[enhanced,nobeforeafter,width=7.1cm,height=2.1cm,valign=center,colback=ccal!6,colframe=ccal,arc=2.5mm,boxrule=0.9pt,left=3mm]{\Large\bfseries\color{ccal}a)}\quad\begin{minipage}[c]{5.6cm}\raggedright\Large $821$\end{minipage}\end{tcolorbox}\hfill\begin{tcolorbox}[enhanced,nobeforeafter,width=7.1cm,height=2.1cm,valign=center,colback=ccal!6,colframe=ccal,arc=2.5mm,boxrule=0.9pt,left=3mm]{\Large\bfseries\color{ccal}b)}\quad\begin{minipage}[c]{5.6cm}\raggedright\Large $819$\end{minipage}\end{tcolorbox}\par\vspace{2.5mm}
\begin{tcolorbox}[enhanced,nobeforeafter,width=7.1cm,height=2.1cm,valign=center,colback=ccal!6,colframe=ccal,arc=2.5mm,boxrule=0.9pt,left=3mm]{\Large\bfseries\color{ccal}c)}\quad\begin{minipage}[c]{5.6cm}\raggedright\Large $829$\end{minipage}\end{tcolorbox}\hfill\begin{tcolorbox}[enhanced,nobeforeafter,width=7.1cm,height=2.1cm,valign=center,colback=ccal!6,colframe=ccal,arc=2.5mm,boxrule=0.9pt,left=3mm]{\Large\bfseries\color{ccal}d)}\quad\begin{minipage}[c]{5.6cm}\raggedright\Large $812$\end{minipage}\end{tcolorbox}
\end{tcolorbox}
\newpage
\begin{tcolorbox}[enhanced,width=16cm,height=14cm,colback=white,colframe=ccal,boxrule=2pt,arc=0mm,
left=6mm,right=6mm,top=0mm,bottom=4mm,sharp corners,
overlay={\fill[ccal] (frame.north west) rectangle ([yshift=-1.25cm]frame.north east);
\node[anchor=west,text=white,font=\sffamily\bfseries\LARGE] at ([xshift=5mm,yshift=-0.62cm]frame.north west) {Question 28 D};
\node[anchor=east,text=white,font=\sffamily\itshape\large] at ([xshift=-5mm,yshift=-0.62cm]frame.north east) {Calculs};}]
\vspace*{1.75cm}
\begin{minipage}[t][6.9cm][c]{\linewidth}\LARGE\centering
\begin{minipage}{0.95\linewidth}\raggedright\linespread{1.25}\selectfont Lequel de ces nombres est divisible par $5$~?\end{minipage}
\end{minipage}

\vfill
\begin{tcolorbox}[enhanced,nobeforeafter,width=7.1cm,height=2.1cm,valign=center,colback=ccal!6,colframe=ccal,arc=2.5mm,boxrule=0.9pt,left=3mm]{\Large\bfseries\color{ccal}a)}\quad\begin{minipage}[c]{5.6cm}\raggedright\Large $2\,097$\end{minipage}\end{tcolorbox}\hfill\begin{tcolorbox}[enhanced,nobeforeafter,width=7.1cm,height=2.1cm,valign=center,colback=ccal!6,colframe=ccal,arc=2.5mm,boxrule=0.9pt,left=3mm]{\Large\bfseries\color{ccal}b)}\quad\begin{minipage}[c]{5.6cm}\raggedright\Large $2\,709$\end{minipage}\end{tcolorbox}\par\vspace{2.5mm}
\begin{tcolorbox}[enhanced,nobeforeafter,width=7.1cm,height=2.1cm,valign=center,colback=ccal!6,colframe=ccal,arc=2.5mm,boxrule=0.9pt,left=3mm]{\Large\bfseries\color{ccal}c)}\quad\begin{minipage}[c]{5.6cm}\raggedright\Large $7\,902$\end{minipage}\end{tcolorbox}\hfill\begin{tcolorbox}[enhanced,nobeforeafter,width=7.1cm,height=2.1cm,valign=center,colback=ccal!6,colframe=ccal,arc=2.5mm,boxrule=0.9pt,left=3mm]{\Large\bfseries\color{ccal}d)}\quad\begin{minipage}[c]{5.6cm}\raggedright\Large $2\,790$\end{minipage}\end{tcolorbox}
\end{tcolorbox}
\newpage
\begin{tcolorbox}[enhanced,width=16cm,height=14cm,colback=white,colframe=ccal,boxrule=2pt,arc=0mm,
left=6mm,right=6mm,top=0mm,bottom=4mm,sharp corners,
overlay={\fill[ccal] (frame.north west) rectangle ([yshift=-1.25cm]frame.north east);
\node[anchor=west,text=white,font=\sffamily\bfseries\LARGE] at ([xshift=5mm,yshift=-0.62cm]frame.north west) {Question 29 D};
\node[anchor=east,text=white,font=\sffamily\itshape\large] at ([xshift=-5mm,yshift=-0.62cm]frame.north east) {Calculs};}]
\vspace*{1.75cm}
\begin{minipage}[t][6.9cm][c]{\linewidth}\LARGE\centering
\begin{minipage}{0.95\linewidth}\raggedright\linespread{1.25}\selectfont Que vaut $2+6\times5-1$~?\end{minipage}
\end{minipage}

\vfill
\begin{tcolorbox}[enhanced,nobeforeafter,width=7.1cm,height=2.1cm,valign=center,colback=ccal!6,colframe=ccal,arc=2.5mm,boxrule=0.9pt,left=3mm]{\Large\bfseries\color{ccal}a)}\quad\begin{minipage}[c]{5.6cm}\raggedright\Large $39$\end{minipage}\end{tcolorbox}\hfill\begin{tcolorbox}[enhanced,nobeforeafter,width=7.1cm,height=2.1cm,valign=center,colback=ccal!6,colframe=ccal,arc=2.5mm,boxrule=0.9pt,left=3mm]{\Large\bfseries\color{ccal}b)}\quad\begin{minipage}[c]{5.6cm}\raggedright\Large $29$\end{minipage}\end{tcolorbox}\par\vspace{2.5mm}
\begin{tcolorbox}[enhanced,nobeforeafter,width=7.1cm,height=2.1cm,valign=center,colback=ccal!6,colframe=ccal,arc=2.5mm,boxrule=0.9pt,left=3mm]{\Large\bfseries\color{ccal}c)}\quad\begin{minipage}[c]{5.6cm}\raggedright\Large $31$\end{minipage}\end{tcolorbox}\hfill\begin{tcolorbox}[enhanced,nobeforeafter,width=7.1cm,height=2.1cm,valign=center,colback=ccal!6,colframe=ccal,arc=2.5mm,boxrule=0.9pt,left=3mm]{\Large\bfseries\color{ccal}d)}\quad\begin{minipage}[c]{5.6cm}\raggedright\Large $40$\end{minipage}\end{tcolorbox}
\end{tcolorbox}
\newpage
\begin{tcolorbox}[enhanced,width=16cm,height=14cm,colback=white,colframe=ccal,boxrule=2pt,arc=0mm,
left=6mm,right=6mm,top=0mm,bottom=4mm,sharp corners,
overlay={\fill[ccal] (frame.north west) rectangle ([yshift=-1.25cm]frame.north east);
\node[anchor=west,text=white,font=\sffamily\bfseries\LARGE] at ([xshift=5mm,yshift=-0.62cm]frame.north west) {Question 30 D};
\node[anchor=east,text=white,font=\sffamily\itshape\large] at ([xshift=-5mm,yshift=-0.62cm]frame.north east) {Calculs};}]
\vspace*{1.75cm}
\begin{minipage}[t][6.9cm][c]{\linewidth}\LARGE\centering
\begin{minipage}{0.95\linewidth}\raggedright\linespread{1.25}\selectfont Que vaut $(15-3)\div4$~?\end{minipage}
\end{minipage}

\vfill
\begin{tcolorbox}[enhanced,nobeforeafter,width=7.1cm,height=2.1cm,valign=center,colback=ccal!6,colframe=ccal,arc=2.5mm,boxrule=0.9pt,left=3mm]{\Large\bfseries\color{ccal}a)}\quad\begin{minipage}[c]{5.6cm}\raggedright\Large $12$\end{minipage}\end{tcolorbox}\hfill\begin{tcolorbox}[enhanced,nobeforeafter,width=7.1cm,height=2.1cm,valign=center,colback=ccal!6,colframe=ccal,arc=2.5mm,boxrule=0.9pt,left=3mm]{\Large\bfseries\color{ccal}b)}\quad\begin{minipage}[c]{5.6cm}\raggedright\Large $9$\end{minipage}\end{tcolorbox}\par\vspace{2.5mm}
\begin{tcolorbox}[enhanced,nobeforeafter,width=7.1cm,height=2.1cm,valign=center,colback=ccal!6,colframe=ccal,arc=2.5mm,boxrule=0.9pt,left=3mm]{\Large\bfseries\color{ccal}c)}\quad\begin{minipage}[c]{5.6cm}\raggedright\Large $3$\end{minipage}\end{tcolorbox}\hfill\begin{tcolorbox}[enhanced,nobeforeafter,width=7.1cm,height=2.1cm,valign=center,colback=ccal!6,colframe=ccal,arc=2.5mm,boxrule=0.9pt,left=3mm]{\Large\bfseries\color{ccal}d)}\quad\begin{minipage}[c]{5.6cm}\raggedright\Large $14$\end{minipage}\end{tcolorbox}
\end{tcolorbox}
\newpage
\begin{tcolorbox}[enhanced,width=16cm,height=14cm,colback=white,colframe=ccal,boxrule=2pt,arc=0mm,
left=6mm,right=6mm,top=0mm,bottom=4mm,sharp corners,
overlay={\fill[ccal] (frame.north west) rectangle ([yshift=-1.25cm]frame.north east);
\node[anchor=west,text=white,font=\sffamily\bfseries\LARGE] at ([xshift=5mm,yshift=-0.62cm]frame.north west) {Question 31 D};
\node[anchor=east,text=white,font=\sffamily\itshape\large] at ([xshift=-5mm,yshift=-0.62cm]frame.north east) {Calculs};}]
\vspace*{1.75cm}
\begin{minipage}[t][6.9cm][c]{\linewidth}\LARGE\centering
\begin{minipage}{0.95\linewidth}\raggedright\linespread{1.25}\selectfont Que vaut $1\div4$~?\end{minipage}
\end{minipage}

\vfill
\begin{tcolorbox}[enhanced,nobeforeafter,width=7.1cm,height=2.1cm,valign=center,colback=ccal!6,colframe=ccal,arc=2.5mm,boxrule=0.9pt,left=3mm]{\Large\bfseries\color{ccal}a)}\quad\begin{minipage}[c]{5.6cm}\raggedright\Large $2{,}5$\end{minipage}\end{tcolorbox}\hfill\begin{tcolorbox}[enhanced,nobeforeafter,width=7.1cm,height=2.1cm,valign=center,colback=ccal!6,colframe=ccal,arc=2.5mm,boxrule=0.9pt,left=3mm]{\Large\bfseries\color{ccal}b)}\quad\begin{minipage}[c]{5.6cm}\raggedright\Large $4$\end{minipage}\end{tcolorbox}\par\vspace{2.5mm}
\begin{tcolorbox}[enhanced,nobeforeafter,width=7.1cm,height=2.1cm,valign=center,colback=ccal!6,colframe=ccal,arc=2.5mm,boxrule=0.9pt,left=3mm]{\Large\bfseries\color{ccal}c)}\quad\begin{minipage}[c]{5.6cm}\raggedright\Large $0{,}25$\end{minipage}\end{tcolorbox}\hfill\begin{tcolorbox}[enhanced,nobeforeafter,width=7.1cm,height=2.1cm,valign=center,colback=ccal!6,colframe=ccal,arc=2.5mm,boxrule=0.9pt,left=3mm]{\Large\bfseries\color{ccal}d)}\quad\begin{minipage}[c]{5.6cm}\raggedright\Large $0{,}1$\end{minipage}\end{tcolorbox}
\end{tcolorbox}
\newpage
\begin{tcolorbox}[enhanced,width=16cm,height=14cm,colback=white,colframe=ccal,boxrule=2pt,arc=0mm,
left=6mm,right=6mm,top=0mm,bottom=4mm,sharp corners,
overlay={\fill[ccal] (frame.north west) rectangle ([yshift=-1.25cm]frame.north east);
\node[anchor=west,text=white,font=\sffamily\bfseries\LARGE] at ([xshift=5mm,yshift=-0.62cm]frame.north west) {Question 32 D};
\node[anchor=east,text=white,font=\sffamily\itshape\large] at ([xshift=-5mm,yshift=-0.62cm]frame.north east) {Calculs};}]
\vspace*{1.75cm}
\begin{minipage}[t][6.9cm][c]{\linewidth}\LARGE\centering
\begin{minipage}{0.95\linewidth}\raggedright\linespread{1.25}\selectfont Que vaut $34\div5$~?\end{minipage}
\end{minipage}

\vfill
\begin{tcolorbox}[enhanced,nobeforeafter,width=7.1cm,height=2.1cm,valign=center,colback=ccal!6,colframe=ccal,arc=2.5mm,boxrule=0.9pt,left=3mm]{\Large\bfseries\color{ccal}a)}\quad\begin{minipage}[c]{5.6cm}\raggedright\Large $6{,}4$\end{minipage}\end{tcolorbox}\hfill\begin{tcolorbox}[enhanced,nobeforeafter,width=7.1cm,height=2.1cm,valign=center,colback=ccal!6,colframe=ccal,arc=2.5mm,boxrule=0.9pt,left=3mm]{\Large\bfseries\color{ccal}b)}\quad\begin{minipage}[c]{5.6cm}\raggedright\Large $68$\end{minipage}\end{tcolorbox}\par\vspace{2.5mm}
\begin{tcolorbox}[enhanced,nobeforeafter,width=7.1cm,height=2.1cm,valign=center,colback=ccal!6,colframe=ccal,arc=2.5mm,boxrule=0.9pt,left=3mm]{\Large\bfseries\color{ccal}c)}\quad\begin{minipage}[c]{5.6cm}\raggedright\Large $\dfrac{5}{34}$\end{minipage}\end{tcolorbox}\hfill\begin{tcolorbox}[enhanced,nobeforeafter,width=7.1cm,height=2.1cm,valign=center,colback=ccal!6,colframe=ccal,arc=2.5mm,boxrule=0.9pt,left=3mm]{\Large\bfseries\color{ccal}d)}\quad\begin{minipage}[c]{5.6cm}\raggedright\Large $6{,}8$\end{minipage}\end{tcolorbox}
\end{tcolorbox}
\newpage
\begin{tcolorbox}[enhanced,width=16cm,height=14cm,colback=white,colframe=cfra,boxrule=2pt,arc=0mm,
left=6mm,right=6mm,top=0mm,bottom=4mm,sharp corners,
overlay={\fill[cfra] (frame.north west) rectangle ([yshift=-1.25cm]frame.north east);
\node[anchor=west,text=white,font=\sffamily\bfseries\LARGE] at ([xshift=5mm,yshift=-0.62cm]frame.north west) {Question 33 D};
\node[anchor=east,text=white,font=\sffamily\itshape\large] at ([xshift=-5mm,yshift=-0.62cm]frame.north east) {Fractions};}]
\vspace*{1.75cm}
\begin{minipage}[t][6.9cm][c]{\linewidth}\LARGE\centering
\begin{minipage}{0.95\linewidth}\raggedright\linespread{1.25}\selectfont Combien vaut $\dfrac{4}{9}$ de $36$~?\end{minipage}
\end{minipage}

\vfill
\begin{tcolorbox}[enhanced,nobeforeafter,width=7.1cm,height=2.1cm,valign=center,colback=cfra!6,colframe=cfra,arc=2.5mm,boxrule=0.9pt,left=3mm]{\Large\bfseries\color{cfra}a)}\quad\begin{minipage}[c]{5.6cm}\raggedright\Large $20$\end{minipage}\end{tcolorbox}\hfill\begin{tcolorbox}[enhanced,nobeforeafter,width=7.1cm,height=2.1cm,valign=center,colback=cfra!6,colframe=cfra,arc=2.5mm,boxrule=0.9pt,left=3mm]{\Large\bfseries\color{cfra}b)}\quad\begin{minipage}[c]{5.6cm}\raggedright\Large $4$\end{minipage}\end{tcolorbox}\par\vspace{2.5mm}
\begin{tcolorbox}[enhanced,nobeforeafter,width=7.1cm,height=2.1cm,valign=center,colback=cfra!6,colframe=cfra,arc=2.5mm,boxrule=0.9pt,left=3mm]{\Large\bfseries\color{cfra}c)}\quad\begin{minipage}[c]{5.6cm}\raggedright\Large $144$\end{minipage}\end{tcolorbox}\hfill\begin{tcolorbox}[enhanced,nobeforeafter,width=7.1cm,height=2.1cm,valign=center,colback=cfra!6,colframe=cfra,arc=2.5mm,boxrule=0.9pt,left=3mm]{\Large\bfseries\color{cfra}d)}\quad\begin{minipage}[c]{5.6cm}\raggedright\Large $16$\end{minipage}\end{tcolorbox}
\end{tcolorbox}
\newpage
\begin{tcolorbox}[enhanced,width=16cm,height=14cm,colback=white,colframe=cfra,boxrule=2pt,arc=0mm,
left=6mm,right=6mm,top=0mm,bottom=4mm,sharp corners,
overlay={\fill[cfra] (frame.north west) rectangle ([yshift=-1.25cm]frame.north east);
\node[anchor=west,text=white,font=\sffamily\bfseries\LARGE] at ([xshift=5mm,yshift=-0.62cm]frame.north west) {Question 34 D};
\node[anchor=east,text=white,font=\sffamily\itshape\large] at ([xshift=-5mm,yshift=-0.62cm]frame.north east) {Fractions};}]
\vspace*{1.75cm}
\begin{minipage}[t][6.9cm][c]{\linewidth}\LARGE\centering
\begin{minipage}{0.95\linewidth}\raggedright\linespread{1.25}\selectfont Combien vaut $\dfrac{3}{5}$ de $45$~?\end{minipage}
\end{minipage}

\vfill
\begin{tcolorbox}[enhanced,nobeforeafter,width=7.1cm,height=2.1cm,valign=center,colback=cfra!6,colframe=cfra,arc=2.5mm,boxrule=0.9pt,left=3mm]{\Large\bfseries\color{cfra}a)}\quad\begin{minipage}[c]{5.6cm}\raggedright\Large $18$\end{minipage}\end{tcolorbox}\hfill\begin{tcolorbox}[enhanced,nobeforeafter,width=7.1cm,height=2.1cm,valign=center,colback=cfra!6,colframe=cfra,arc=2.5mm,boxrule=0.9pt,left=3mm]{\Large\bfseries\color{cfra}b)}\quad\begin{minipage}[c]{5.6cm}\raggedright\Large $135$\end{minipage}\end{tcolorbox}\par\vspace{2.5mm}
\begin{tcolorbox}[enhanced,nobeforeafter,width=7.1cm,height=2.1cm,valign=center,colback=cfra!6,colframe=cfra,arc=2.5mm,boxrule=0.9pt,left=3mm]{\Large\bfseries\color{cfra}c)}\quad\begin{minipage}[c]{5.6cm}\raggedright\Large $27$\end{minipage}\end{tcolorbox}\hfill\begin{tcolorbox}[enhanced,nobeforeafter,width=7.1cm,height=2.1cm,valign=center,colback=cfra!6,colframe=cfra,arc=2.5mm,boxrule=0.9pt,left=3mm]{\Large\bfseries\color{cfra}d)}\quad\begin{minipage}[c]{5.6cm}\raggedright\Large $9$\end{minipage}\end{tcolorbox}
\end{tcolorbox}
\newpage
\begin{tcolorbox}[enhanced,width=16cm,height=14cm,colback=white,colframe=cfra,boxrule=2pt,arc=0mm,
left=6mm,right=6mm,top=0mm,bottom=4mm,sharp corners,
overlay={\fill[cfra] (frame.north west) rectangle ([yshift=-1.25cm]frame.north east);
\node[anchor=west,text=white,font=\sffamily\bfseries\LARGE] at ([xshift=5mm,yshift=-0.62cm]frame.north west) {Question 35 D};
\node[anchor=east,text=white,font=\sffamily\itshape\large] at ([xshift=-5mm,yshift=-0.62cm]frame.north east) {Fractions};}]
\vspace*{1.75cm}
\begin{minipage}[t][6.9cm][c]{\linewidth}\LARGE\centering
\begin{minipage}{0.95\linewidth}\raggedright\linespread{1.25}\selectfont Compléter~: $\dfrac{7}{9}=\dfrac{?}{27}$\end{minipage}
\end{minipage}

\vfill
\begin{tcolorbox}[enhanced,nobeforeafter,width=7.1cm,height=2.1cm,valign=center,colback=cfra!6,colframe=cfra,arc=2.5mm,boxrule=0.9pt,left=3mm]{\Large\bfseries\color{cfra}a)}\quad\begin{minipage}[c]{5.6cm}\raggedright\Large $25$\end{minipage}\end{tcolorbox}\hfill\begin{tcolorbox}[enhanced,nobeforeafter,width=7.1cm,height=2.1cm,valign=center,colback=cfra!6,colframe=cfra,arc=2.5mm,boxrule=0.9pt,left=3mm]{\Large\bfseries\color{cfra}b)}\quad\begin{minipage}[c]{5.6cm}\raggedright\Large $3$\end{minipage}\end{tcolorbox}\par\vspace{2.5mm}
\begin{tcolorbox}[enhanced,nobeforeafter,width=7.1cm,height=2.1cm,valign=center,colback=cfra!6,colframe=cfra,arc=2.5mm,boxrule=0.9pt,left=3mm]{\Large\bfseries\color{cfra}c)}\quad\begin{minipage}[c]{5.6cm}\raggedright\Large $189$\end{minipage}\end{tcolorbox}\hfill\begin{tcolorbox}[enhanced,nobeforeafter,width=7.1cm,height=2.1cm,valign=center,colback=cfra!6,colframe=cfra,arc=2.5mm,boxrule=0.9pt,left=3mm]{\Large\bfseries\color{cfra}d)}\quad\begin{minipage}[c]{5.6cm}\raggedright\Large $21$\end{minipage}\end{tcolorbox}
\end{tcolorbox}
\newpage
\begin{tcolorbox}[enhanced,width=16cm,height=14cm,colback=white,colframe=cfra,boxrule=2pt,arc=0mm,
left=6mm,right=6mm,top=0mm,bottom=4mm,sharp corners,
overlay={\fill[cfra] (frame.north west) rectangle ([yshift=-1.25cm]frame.north east);
\node[anchor=west,text=white,font=\sffamily\bfseries\LARGE] at ([xshift=5mm,yshift=-0.62cm]frame.north west) {Question 36 D};
\node[anchor=east,text=white,font=\sffamily\itshape\large] at ([xshift=-5mm,yshift=-0.62cm]frame.north east) {Fractions};}]
\vspace*{1.75cm}
\begin{minipage}[t][6.9cm][c]{\linewidth}\LARGE\centering
\begin{minipage}{0.95\linewidth}\raggedright\linespread{1.25}\selectfont Compléter~: $\dfrac{2}{11}=\dfrac{?}{44}$\end{minipage}
\end{minipage}

\vfill
\begin{tcolorbox}[enhanced,nobeforeafter,width=7.1cm,height=2.1cm,valign=center,colback=cfra!6,colframe=cfra,arc=2.5mm,boxrule=0.9pt,left=3mm]{\Large\bfseries\color{cfra}a)}\quad\begin{minipage}[c]{5.6cm}\raggedright\Large $8$\end{minipage}\end{tcolorbox}\hfill\begin{tcolorbox}[enhanced,nobeforeafter,width=7.1cm,height=2.1cm,valign=center,colback=cfra!6,colframe=cfra,arc=2.5mm,boxrule=0.9pt,left=3mm]{\Large\bfseries\color{cfra}b)}\quad\begin{minipage}[c]{5.6cm}\raggedright\Large $35$\end{minipage}\end{tcolorbox}\par\vspace{2.5mm}
\begin{tcolorbox}[enhanced,nobeforeafter,width=7.1cm,height=2.1cm,valign=center,colback=cfra!6,colframe=cfra,arc=2.5mm,boxrule=0.9pt,left=3mm]{\Large\bfseries\color{cfra}c)}\quad\begin{minipage}[c]{5.6cm}\raggedright\Large $88$\end{minipage}\end{tcolorbox}\hfill\begin{tcolorbox}[enhanced,nobeforeafter,width=7.1cm,height=2.1cm,valign=center,colback=cfra!6,colframe=cfra,arc=2.5mm,boxrule=0.9pt,left=3mm]{\Large\bfseries\color{cfra}d)}\quad\begin{minipage}[c]{5.6cm}\raggedright\Large $4$\end{minipage}\end{tcolorbox}
\end{tcolorbox}
\newpage
\begin{tcolorbox}[enhanced,width=16cm,height=14cm,colback=white,colframe=cfra,boxrule=2pt,arc=0mm,
left=6mm,right=6mm,top=0mm,bottom=4mm,sharp corners,
overlay={\fill[cfra] (frame.north west) rectangle ([yshift=-1.25cm]frame.north east);
\node[anchor=west,text=white,font=\sffamily\bfseries\LARGE] at ([xshift=5mm,yshift=-0.62cm]frame.north west) {Question 37 D};
\node[anchor=east,text=white,font=\sffamily\itshape\large] at ([xshift=-5mm,yshift=-0.62cm]frame.north east) {Fractions};}]
\vspace*{1.75cm}
\begin{minipage}[t][6.9cm][c]{\linewidth}\LARGE\centering
\begin{minipage}{0.95\linewidth}\raggedright\linespread{1.25}\selectfont Quelle est l'écriture décimale de $\dfrac{11}{4}$~?\end{minipage}
\end{minipage}

\vfill
\begin{tcolorbox}[enhanced,nobeforeafter,width=7.1cm,height=2.1cm,valign=center,colback=cfra!6,colframe=cfra,arc=2.5mm,boxrule=0.9pt,left=3mm]{\Large\bfseries\color{cfra}a)}\quad\begin{minipage}[c]{5.6cm}\raggedright\Large $27{,}5$\end{minipage}\end{tcolorbox}\hfill\begin{tcolorbox}[enhanced,nobeforeafter,width=7.1cm,height=2.1cm,valign=center,colback=cfra!6,colframe=cfra,arc=2.5mm,boxrule=0.9pt,left=3mm]{\Large\bfseries\color{cfra}b)}\quad\begin{minipage}[c]{5.6cm}\raggedright\Large $\dfrac{4}{11}$\end{minipage}\end{tcolorbox}\par\vspace{2.5mm}
\begin{tcolorbox}[enhanced,nobeforeafter,width=7.1cm,height=2.1cm,valign=center,colback=cfra!6,colframe=cfra,arc=2.5mm,boxrule=0.9pt,left=3mm]{\Large\bfseries\color{cfra}c)}\quad\begin{minipage}[c]{5.6cm}\raggedright\Large $2{,}75$\end{minipage}\end{tcolorbox}\hfill\begin{tcolorbox}[enhanced,nobeforeafter,width=7.1cm,height=2.1cm,valign=center,colback=cfra!6,colframe=cfra,arc=2.5mm,boxrule=0.9pt,left=3mm]{\Large\bfseries\color{cfra}d)}\quad\begin{minipage}[c]{5.6cm}\raggedright\Large $11{,}4$\end{minipage}\end{tcolorbox}
\end{tcolorbox}
\newpage
\begin{tcolorbox}[enhanced,width=16cm,height=14cm,colback=white,colframe=cfra,boxrule=2pt,arc=0mm,
left=6mm,right=6mm,top=0mm,bottom=4mm,sharp corners,
overlay={\fill[cfra] (frame.north west) rectangle ([yshift=-1.25cm]frame.north east);
\node[anchor=west,text=white,font=\sffamily\bfseries\LARGE] at ([xshift=5mm,yshift=-0.62cm]frame.north west) {Question 38 D};
\node[anchor=east,text=white,font=\sffamily\itshape\large] at ([xshift=-5mm,yshift=-0.62cm]frame.north east) {Fractions};}]
\vspace*{1.75cm}
\begin{minipage}[t][6.9cm][c]{\linewidth}\LARGE\centering
\begin{minipage}{0.95\linewidth}\raggedright\linespread{1.25}\selectfont Quelle est l'écriture décimale de $\dfrac{13}{10}$~?\end{minipage}
\end{minipage}

\vfill
\begin{tcolorbox}[enhanced,nobeforeafter,width=7.1cm,height=2.1cm,valign=center,colback=cfra!6,colframe=cfra,arc=2.5mm,boxrule=0.9pt,left=3mm]{\Large\bfseries\color{cfra}a)}\quad\begin{minipage}[c]{5.6cm}\raggedright\Large $13{,}10$\end{minipage}\end{tcolorbox}\hfill\begin{tcolorbox}[enhanced,nobeforeafter,width=7.1cm,height=2.1cm,valign=center,colback=cfra!6,colframe=cfra,arc=2.5mm,boxrule=0.9pt,left=3mm]{\Large\bfseries\color{cfra}b)}\quad\begin{minipage}[c]{5.6cm}\raggedright\Large $1{,}3$\end{minipage}\end{tcolorbox}\par\vspace{2.5mm}
\begin{tcolorbox}[enhanced,nobeforeafter,width=7.1cm,height=2.1cm,valign=center,colback=cfra!6,colframe=cfra,arc=2.5mm,boxrule=0.9pt,left=3mm]{\Large\bfseries\color{cfra}c)}\quad\begin{minipage}[c]{5.6cm}\raggedright\Large $\dfrac{10}{13}$\end{minipage}\end{tcolorbox}\hfill\begin{tcolorbox}[enhanced,nobeforeafter,width=7.1cm,height=2.1cm,valign=center,colback=cfra!6,colframe=cfra,arc=2.5mm,boxrule=0.9pt,left=3mm]{\Large\bfseries\color{cfra}d)}\quad\begin{minipage}[c]{5.6cm}\raggedright\Large $13$\end{minipage}\end{tcolorbox}
\end{tcolorbox}
\newpage
\begin{tcolorbox}[enhanced,width=16cm,height=14cm,colback=white,colframe=cfra,boxrule=2pt,arc=0mm,
left=6mm,right=6mm,top=0mm,bottom=4mm,sharp corners,
overlay={\fill[cfra] (frame.north west) rectangle ([yshift=-1.25cm]frame.north east);
\node[anchor=west,text=white,font=\sffamily\bfseries\LARGE] at ([xshift=5mm,yshift=-0.62cm]frame.north west) {Question 39 D};
\node[anchor=east,text=white,font=\sffamily\itshape\large] at ([xshift=-5mm,yshift=-0.62cm]frame.north east) {Fractions};}]
\vspace*{1.75cm}
\begin{minipage}[t][6.9cm][c]{\linewidth}\LARGE\centering
\begin{minipage}{0.95\linewidth}\raggedright\linespread{1.25}\selectfont Laquelle de ces fractions est plus grande que $1$~?\end{minipage}
\end{minipage}

\vfill
\begin{tcolorbox}[enhanced,nobeforeafter,width=7.1cm,height=2.1cm,valign=center,colback=cfra!6,colframe=cfra,arc=2.5mm,boxrule=0.9pt,left=3mm]{\Large\bfseries\color{cfra}a)}\quad\begin{minipage}[c]{5.6cm}\raggedright\Large $\dfrac{4}{7}$\end{minipage}\end{tcolorbox}\hfill\begin{tcolorbox}[enhanced,nobeforeafter,width=7.1cm,height=2.1cm,valign=center,colback=cfra!6,colframe=cfra,arc=2.5mm,boxrule=0.9pt,left=3mm]{\Large\bfseries\color{cfra}b)}\quad\begin{minipage}[c]{5.6cm}\raggedright\Large $\dfrac{2}{2}$\end{minipage}\end{tcolorbox}\par\vspace{2.5mm}
\begin{tcolorbox}[enhanced,nobeforeafter,width=7.1cm,height=2.1cm,valign=center,colback=cfra!6,colframe=cfra,arc=2.5mm,boxrule=0.9pt,left=3mm]{\Large\bfseries\color{cfra}c)}\quad\begin{minipage}[c]{5.6cm}\raggedright\Large $\dfrac{9}{10}$\end{minipage}\end{tcolorbox}\hfill\begin{tcolorbox}[enhanced,nobeforeafter,width=7.1cm,height=2.1cm,valign=center,colback=cfra!6,colframe=cfra,arc=2.5mm,boxrule=0.9pt,left=3mm]{\Large\bfseries\color{cfra}d)}\quad\begin{minipage}[c]{5.6cm}\raggedright\Large $\dfrac{10}{9}$\end{minipage}\end{tcolorbox}
\end{tcolorbox}
\newpage
\begin{tcolorbox}[enhanced,width=16cm,height=14cm,colback=white,colframe=cfra,boxrule=2pt,arc=0mm,
left=6mm,right=6mm,top=0mm,bottom=4mm,sharp corners,
overlay={\fill[cfra] (frame.north west) rectangle ([yshift=-1.25cm]frame.north east);
\node[anchor=west,text=white,font=\sffamily\bfseries\LARGE] at ([xshift=5mm,yshift=-0.62cm]frame.north west) {Question 40 D};
\node[anchor=east,text=white,font=\sffamily\itshape\large] at ([xshift=-5mm,yshift=-0.62cm]frame.north east) {Fractions};}]
\vspace*{1.75cm}
\begin{minipage}[t][6.9cm][c]{\linewidth}\LARGE\centering
\begin{minipage}{0.95\linewidth}\raggedright\linespread{1.25}\selectfont Laquelle de ces fractions est plus grande que $1$~?\end{minipage}
\end{minipage}

\vfill
\begin{tcolorbox}[enhanced,nobeforeafter,width=7.1cm,height=2.1cm,valign=center,colback=cfra!6,colframe=cfra,arc=2.5mm,boxrule=0.9pt,left=3mm]{\Large\bfseries\color{cfra}a)}\quad\begin{minipage}[c]{5.6cm}\raggedright\Large $\dfrac{15}{4}$\end{minipage}\end{tcolorbox}\hfill\begin{tcolorbox}[enhanced,nobeforeafter,width=7.1cm,height=2.1cm,valign=center,colback=cfra!6,colframe=cfra,arc=2.5mm,boxrule=0.9pt,left=3mm]{\Large\bfseries\color{cfra}b)}\quad\begin{minipage}[c]{5.6cm}\raggedright\Large $\dfrac{8}{8}$\end{minipage}\end{tcolorbox}\par\vspace{2.5mm}
\begin{tcolorbox}[enhanced,nobeforeafter,width=7.1cm,height=2.1cm,valign=center,colback=cfra!6,colframe=cfra,arc=2.5mm,boxrule=0.9pt,left=3mm]{\Large\bfseries\color{cfra}c)}\quad\begin{minipage}[c]{5.6cm}\raggedright\Large $\dfrac{4}{15}$\end{minipage}\end{tcolorbox}\hfill\begin{tcolorbox}[enhanced,nobeforeafter,width=7.1cm,height=2.1cm,valign=center,colback=cfra!6,colframe=cfra,arc=2.5mm,boxrule=0.9pt,left=3mm]{\Large\bfseries\color{cfra}d)}\quad\begin{minipage}[c]{5.6cm}\raggedright\Large $\dfrac{3}{8}$\end{minipage}\end{tcolorbox}
\end{tcolorbox}
\newpage
\begin{tcolorbox}[enhanced,width=16cm,height=14cm,colback=white,colframe=cfra,boxrule=2pt,arc=0mm,
left=6mm,right=6mm,top=0mm,bottom=4mm,sharp corners,
overlay={\fill[cfra] (frame.north west) rectangle ([yshift=-1.25cm]frame.north east);
\node[anchor=west,text=white,font=\sffamily\bfseries\LARGE] at ([xshift=5mm,yshift=-0.62cm]frame.north west) {Question 41 D};
\node[anchor=east,text=white,font=\sffamily\itshape\large] at ([xshift=-5mm,yshift=-0.62cm]frame.north east) {Fractions};}]
\vspace*{1.75cm}
\begin{minipage}[t][6.9cm][c]{\linewidth}\LARGE\centering
\begin{minipage}{0.95\linewidth}\raggedright\linespread{1.25}\selectfont Quel nombre, multiplié par $11$, donne $4$~?\end{minipage}
\end{minipage}

\vfill
\begin{tcolorbox}[enhanced,nobeforeafter,width=7.1cm,height=2.1cm,valign=center,colback=cfra!6,colframe=cfra,arc=2.5mm,boxrule=0.9pt,left=3mm]{\Large\bfseries\color{cfra}a)}\quad\begin{minipage}[c]{5.6cm}\raggedright\Large $\dfrac{4}{11}$\end{minipage}\end{tcolorbox}\hfill\begin{tcolorbox}[enhanced,nobeforeafter,width=7.1cm,height=2.1cm,valign=center,colback=cfra!6,colframe=cfra,arc=2.5mm,boxrule=0.9pt,left=3mm]{\Large\bfseries\color{cfra}b)}\quad\begin{minipage}[c]{5.6cm}\raggedright\Large $7$\end{minipage}\end{tcolorbox}\par\vspace{2.5mm}
\begin{tcolorbox}[enhanced,nobeforeafter,width=7.1cm,height=2.1cm,valign=center,colback=cfra!6,colframe=cfra,arc=2.5mm,boxrule=0.9pt,left=3mm]{\Large\bfseries\color{cfra}c)}\quad\begin{minipage}[c]{5.6cm}\raggedright\Large $\dfrac{11}{4}$\end{minipage}\end{tcolorbox}\hfill\begin{tcolorbox}[enhanced,nobeforeafter,width=7.1cm,height=2.1cm,valign=center,colback=cfra!6,colframe=cfra,arc=2.5mm,boxrule=0.9pt,left=3mm]{\Large\bfseries\color{cfra}d)}\quad\begin{minipage}[c]{5.6cm}\raggedright\Large $44$\end{minipage}\end{tcolorbox}
\end{tcolorbox}
\newpage
\begin{tcolorbox}[enhanced,width=16cm,height=14cm,colback=white,colframe=cfra,boxrule=2pt,arc=0mm,
left=6mm,right=6mm,top=0mm,bottom=4mm,sharp corners,
overlay={\fill[cfra] (frame.north west) rectangle ([yshift=-1.25cm]frame.north east);
\node[anchor=west,text=white,font=\sffamily\bfseries\LARGE] at ([xshift=5mm,yshift=-0.62cm]frame.north west) {Question 42 D};
\node[anchor=east,text=white,font=\sffamily\itshape\large] at ([xshift=-5mm,yshift=-0.62cm]frame.north east) {Fractions};}]
\vspace*{1.75cm}
\begin{minipage}[t][6.9cm][c]{\linewidth}\LARGE\centering
\begin{minipage}{0.95\linewidth}\raggedright\linespread{1.25}\selectfont Quel nombre, multiplié par $5$, donne $9$~?\end{minipage}
\end{minipage}

\vfill
\begin{tcolorbox}[enhanced,nobeforeafter,width=7.1cm,height=2.1cm,valign=center,colback=cfra!6,colframe=cfra,arc=2.5mm,boxrule=0.9pt,left=3mm]{\Large\bfseries\color{cfra}a)}\quad\begin{minipage}[c]{5.6cm}\raggedright\Large $45$\end{minipage}\end{tcolorbox}\hfill\begin{tcolorbox}[enhanced,nobeforeafter,width=7.1cm,height=2.1cm,valign=center,colback=cfra!6,colframe=cfra,arc=2.5mm,boxrule=0.9pt,left=3mm]{\Large\bfseries\color{cfra}b)}\quad\begin{minipage}[c]{5.6cm}\raggedright\Large $\dfrac{5}{9}$\end{minipage}\end{tcolorbox}\par\vspace{2.5mm}
\begin{tcolorbox}[enhanced,nobeforeafter,width=7.1cm,height=2.1cm,valign=center,colback=cfra!6,colframe=cfra,arc=2.5mm,boxrule=0.9pt,left=3mm]{\Large\bfseries\color{cfra}c)}\quad\begin{minipage}[c]{5.6cm}\raggedright\Large $4$\end{minipage}\end{tcolorbox}\hfill\begin{tcolorbox}[enhanced,nobeforeafter,width=7.1cm,height=2.1cm,valign=center,colback=cfra!6,colframe=cfra,arc=2.5mm,boxrule=0.9pt,left=3mm]{\Large\bfseries\color{cfra}d)}\quad\begin{minipage}[c]{5.6cm}\raggedright\Large $\dfrac{9}{5}$\end{minipage}\end{tcolorbox}
\end{tcolorbox}
\newpage
\begin{tcolorbox}[enhanced,width=16cm,height=14cm,colback=white,colframe=cfra,boxrule=2pt,arc=0mm,
left=6mm,right=6mm,top=0mm,bottom=4mm,sharp corners,
overlay={\fill[cfra] (frame.north west) rectangle ([yshift=-1.25cm]frame.north east);
\node[anchor=west,text=white,font=\sffamily\bfseries\LARGE] at ([xshift=5mm,yshift=-0.62cm]frame.north west) {Question 43 D};
\node[anchor=east,text=white,font=\sffamily\itshape\large] at ([xshift=-5mm,yshift=-0.62cm]frame.north east) {Fractions};}]
\vspace*{1.75cm}
\begin{minipage}[t][6.9cm][c]{\linewidth}\LARGE\centering
\begin{minipage}{0.95\linewidth}\raggedright\linespread{1.25}\selectfont Que vaut $\dfrac{7}{15}+\dfrac{5}{15}$~?\end{minipage}
\end{minipage}

\vfill
\begin{tcolorbox}[enhanced,nobeforeafter,width=7.1cm,height=2.1cm,valign=center,colback=cfra!6,colframe=cfra,arc=2.5mm,boxrule=0.9pt,left=3mm]{\Large\bfseries\color{cfra}a)}\quad\begin{minipage}[c]{5.6cm}\raggedright\Large $\dfrac{12}{15}$\end{minipage}\end{tcolorbox}\hfill\begin{tcolorbox}[enhanced,nobeforeafter,width=7.1cm,height=2.1cm,valign=center,colback=cfra!6,colframe=cfra,arc=2.5mm,boxrule=0.9pt,left=3mm]{\Large\bfseries\color{cfra}b)}\quad\begin{minipage}[c]{5.6cm}\raggedright\Large $\dfrac{35}{15}$\end{minipage}\end{tcolorbox}\par\vspace{2.5mm}
\begin{tcolorbox}[enhanced,nobeforeafter,width=7.1cm,height=2.1cm,valign=center,colback=cfra!6,colframe=cfra,arc=2.5mm,boxrule=0.9pt,left=3mm]{\Large\bfseries\color{cfra}c)}\quad\begin{minipage}[c]{5.6cm}\raggedright\Large $\dfrac{12}{225}$\end{minipage}\end{tcolorbox}\hfill\begin{tcolorbox}[enhanced,nobeforeafter,width=7.1cm,height=2.1cm,valign=center,colback=cfra!6,colframe=cfra,arc=2.5mm,boxrule=0.9pt,left=3mm]{\Large\bfseries\color{cfra}d)}\quad\begin{minipage}[c]{5.6cm}\raggedright\Large $\dfrac{12}{30}$\end{minipage}\end{tcolorbox}
\end{tcolorbox}
\newpage
\begin{tcolorbox}[enhanced,width=16cm,height=14cm,colback=white,colframe=cfra,boxrule=2pt,arc=0mm,
left=6mm,right=6mm,top=0mm,bottom=4mm,sharp corners,
overlay={\fill[cfra] (frame.north west) rectangle ([yshift=-1.25cm]frame.north east);
\node[anchor=west,text=white,font=\sffamily\bfseries\LARGE] at ([xshift=5mm,yshift=-0.62cm]frame.north west) {Question 44 D};
\node[anchor=east,text=white,font=\sffamily\itshape\large] at ([xshift=-5mm,yshift=-0.62cm]frame.north east) {Fractions};}]
\vspace*{1.75cm}
\begin{minipage}[t][6.9cm][c]{\linewidth}\LARGE\centering
\begin{minipage}{0.95\linewidth}\raggedright\linespread{1.25}\selectfont Que vaut $\dfrac{6}{12}+\dfrac{1}{12}$~?\end{minipage}
\end{minipage}

\vfill
\begin{tcolorbox}[enhanced,nobeforeafter,width=7.1cm,height=2.1cm,valign=center,colback=cfra!6,colframe=cfra,arc=2.5mm,boxrule=0.9pt,left=3mm]{\Large\bfseries\color{cfra}a)}\quad\begin{minipage}[c]{5.6cm}\raggedright\Large $\dfrac{7}{144}$\end{minipage}\end{tcolorbox}\hfill\begin{tcolorbox}[enhanced,nobeforeafter,width=7.1cm,height=2.1cm,valign=center,colback=cfra!6,colframe=cfra,arc=2.5mm,boxrule=0.9pt,left=3mm]{\Large\bfseries\color{cfra}b)}\quad\begin{minipage}[c]{5.6cm}\raggedright\Large $\dfrac{6}{12}$\end{minipage}\end{tcolorbox}\par\vspace{2.5mm}
\begin{tcolorbox}[enhanced,nobeforeafter,width=7.1cm,height=2.1cm,valign=center,colback=cfra!6,colframe=cfra,arc=2.5mm,boxrule=0.9pt,left=3mm]{\Large\bfseries\color{cfra}c)}\quad\begin{minipage}[c]{5.6cm}\raggedright\Large $\dfrac{7}{12}$\end{minipage}\end{tcolorbox}\hfill\begin{tcolorbox}[enhanced,nobeforeafter,width=7.1cm,height=2.1cm,valign=center,colback=cfra!6,colframe=cfra,arc=2.5mm,boxrule=0.9pt,left=3mm]{\Large\bfseries\color{cfra}d)}\quad\begin{minipage}[c]{5.6cm}\raggedright\Large $\dfrac{7}{24}$\end{minipage}\end{tcolorbox}
\end{tcolorbox}
\newpage
\begin{tcolorbox}[enhanced,width=16cm,height=14cm,colback=white,colframe=cfra,boxrule=2pt,arc=0mm,
left=6mm,right=6mm,top=0mm,bottom=4mm,sharp corners,
overlay={\fill[cfra] (frame.north west) rectangle ([yshift=-1.25cm]frame.north east);
\node[anchor=west,text=white,font=\sffamily\bfseries\LARGE] at ([xshift=5mm,yshift=-0.62cm]frame.north west) {Question 45 D};
\node[anchor=east,text=white,font=\sffamily\itshape\large] at ([xshift=-5mm,yshift=-0.62cm]frame.north east) {Fractions};}]
\vspace*{1.75cm}
\begin{minipage}[t][6.9cm][c]{\linewidth}\LARGE\centering
\begin{minipage}{0.95\linewidth}\raggedright\linespread{1.25}\selectfont Un gâteau est coupé en $4$ parts égales. On mange $3$ parts. Quelle fraction du gâteau reste-t-il~?\end{minipage}
\end{minipage}

\vfill
\begin{tcolorbox}[enhanced,nobeforeafter,width=7.1cm,height=2.1cm,valign=center,colback=cfra!6,colframe=cfra,arc=2.5mm,boxrule=0.9pt,left=3mm]{\Large\bfseries\color{cfra}a)}\quad\begin{minipage}[c]{5.6cm}\raggedright\Large $\dfrac{3}{4}$\end{minipage}\end{tcolorbox}\hfill\begin{tcolorbox}[enhanced,nobeforeafter,width=7.1cm,height=2.1cm,valign=center,colback=cfra!6,colframe=cfra,arc=2.5mm,boxrule=0.9pt,left=3mm]{\Large\bfseries\color{cfra}b)}\quad\begin{minipage}[c]{5.6cm}\raggedright\Large $\dfrac{1}{3}$\end{minipage}\end{tcolorbox}\par\vspace{2.5mm}
\begin{tcolorbox}[enhanced,nobeforeafter,width=7.1cm,height=2.1cm,valign=center,colback=cfra!6,colframe=cfra,arc=2.5mm,boxrule=0.9pt,left=3mm]{\Large\bfseries\color{cfra}c)}\quad\begin{minipage}[c]{5.6cm}\raggedright\Large $\dfrac{3}{1}$\end{minipage}\end{tcolorbox}\hfill\begin{tcolorbox}[enhanced,nobeforeafter,width=7.1cm,height=2.1cm,valign=center,colback=cfra!6,colframe=cfra,arc=2.5mm,boxrule=0.9pt,left=3mm]{\Large\bfseries\color{cfra}d)}\quad\begin{minipage}[c]{5.6cm}\raggedright\Large $\dfrac{1}{4}$\end{minipage}\end{tcolorbox}
\end{tcolorbox}
\newpage
\begin{tcolorbox}[enhanced,width=16cm,height=14cm,colback=white,colframe=cfra,boxrule=2pt,arc=0mm,
left=6mm,right=6mm,top=0mm,bottom=4mm,sharp corners,
overlay={\fill[cfra] (frame.north west) rectangle ([yshift=-1.25cm]frame.north east);
\node[anchor=west,text=white,font=\sffamily\bfseries\LARGE] at ([xshift=5mm,yshift=-0.62cm]frame.north west) {Question 46 D};
\node[anchor=east,text=white,font=\sffamily\itshape\large] at ([xshift=-5mm,yshift=-0.62cm]frame.north east) {Fractions};}]
\vspace*{1.75cm}
\begin{minipage}[t][6.9cm][c]{\linewidth}\LARGE\centering
\begin{minipage}{0.95\linewidth}\raggedright\linespread{1.25}\selectfont Un gâteau est coupé en $7$ parts égales. On mange $2$ parts. Quelle fraction du gâteau reste-t-il~?\end{minipage}
\end{minipage}

\vfill
\begin{tcolorbox}[enhanced,nobeforeafter,width=7.1cm,height=2.1cm,valign=center,colback=cfra!6,colframe=cfra,arc=2.5mm,boxrule=0.9pt,left=3mm]{\Large\bfseries\color{cfra}a)}\quad\begin{minipage}[c]{5.6cm}\raggedright\Large $\dfrac{5}{7}$\end{minipage}\end{tcolorbox}\hfill\begin{tcolorbox}[enhanced,nobeforeafter,width=7.1cm,height=2.1cm,valign=center,colback=cfra!6,colframe=cfra,arc=2.5mm,boxrule=0.9pt,left=3mm]{\Large\bfseries\color{cfra}b)}\quad\begin{minipage}[c]{5.6cm}\raggedright\Large $\dfrac{5}{2}$\end{minipage}\end{tcolorbox}\par\vspace{2.5mm}
\begin{tcolorbox}[enhanced,nobeforeafter,width=7.1cm,height=2.1cm,valign=center,colback=cfra!6,colframe=cfra,arc=2.5mm,boxrule=0.9pt,left=3mm]{\Large\bfseries\color{cfra}c)}\quad\begin{minipage}[c]{5.6cm}\raggedright\Large $\dfrac{2}{7}$\end{minipage}\end{tcolorbox}\hfill\begin{tcolorbox}[enhanced,nobeforeafter,width=7.1cm,height=2.1cm,valign=center,colback=cfra!6,colframe=cfra,arc=2.5mm,boxrule=0.9pt,left=3mm]{\Large\bfseries\color{cfra}d)}\quad\begin{minipage}[c]{5.6cm}\raggedright\Large $\dfrac{2}{5}$\end{minipage}\end{tcolorbox}
\end{tcolorbox}
\newpage
\begin{tcolorbox}[enhanced,width=16cm,height=14cm,colback=white,colframe=cfra,boxrule=2pt,arc=0mm,
left=6mm,right=6mm,top=0mm,bottom=4mm,sharp corners,
overlay={\fill[cfra] (frame.north west) rectangle ([yshift=-1.25cm]frame.north east);
\node[anchor=west,text=white,font=\sffamily\bfseries\LARGE] at ([xshift=5mm,yshift=-0.62cm]frame.north west) {Question 47 D};
\node[anchor=east,text=white,font=\sffamily\itshape\large] at ([xshift=-5mm,yshift=-0.62cm]frame.north east) {Fractions};}]
\vspace*{1.75cm}
\begin{minipage}[t][6.9cm][c]{\linewidth}\LARGE\centering
\begin{minipage}{0.95\linewidth}\raggedright\linespread{1.25}\selectfont Entre quels nombres entiers consécutifs se trouve $\dfrac{11}{10}$~?\end{minipage}
\end{minipage}

\vfill
\begin{tcolorbox}[enhanced,nobeforeafter,width=7.1cm,height=2.1cm,valign=center,colback=cfra!6,colframe=cfra,arc=2.5mm,boxrule=0.9pt,left=3mm]{\Large\bfseries\color{cfra}a)}\quad\begin{minipage}[c]{5.6cm}\raggedright\Large $2\text{ et }3$\end{minipage}\end{tcolorbox}\hfill\begin{tcolorbox}[enhanced,nobeforeafter,width=7.1cm,height=2.1cm,valign=center,colback=cfra!6,colframe=cfra,arc=2.5mm,boxrule=0.9pt,left=3mm]{\Large\bfseries\color{cfra}b)}\quad\begin{minipage}[c]{5.6cm}\raggedright\Large $1\text{ et }2$\end{minipage}\end{tcolorbox}\par\vspace{2.5mm}
\begin{tcolorbox}[enhanced,nobeforeafter,width=7.1cm,height=2.1cm,valign=center,colback=cfra!6,colframe=cfra,arc=2.5mm,boxrule=0.9pt,left=3mm]{\Large\bfseries\color{cfra}c)}\quad\begin{minipage}[c]{5.6cm}\raggedright\Large $0\text{ et }1$\end{minipage}\end{tcolorbox}\hfill\begin{tcolorbox}[enhanced,nobeforeafter,width=7.1cm,height=2.1cm,valign=center,colback=cfra!6,colframe=cfra,arc=2.5mm,boxrule=0.9pt,left=3mm]{\Large\bfseries\color{cfra}d)}\quad\begin{minipage}[c]{5.6cm}\raggedright\Large $3\text{ et }4$\end{minipage}\end{tcolorbox}
\end{tcolorbox}
\newpage
\begin{tcolorbox}[enhanced,width=16cm,height=14cm,colback=white,colframe=cfra,boxrule=2pt,arc=0mm,
left=6mm,right=6mm,top=0mm,bottom=4mm,sharp corners,
overlay={\fill[cfra] (frame.north west) rectangle ([yshift=-1.25cm]frame.north east);
\node[anchor=west,text=white,font=\sffamily\bfseries\LARGE] at ([xshift=5mm,yshift=-0.62cm]frame.north west) {Question 48 D};
\node[anchor=east,text=white,font=\sffamily\itshape\large] at ([xshift=-5mm,yshift=-0.62cm]frame.north east) {Fractions};}]
\vspace*{1.75cm}
\begin{minipage}[t][6.9cm][c]{\linewidth}\LARGE\centering
\begin{minipage}{0.95\linewidth}\raggedright\linespread{1.25}\selectfont Entre quels nombres entiers consécutifs se trouve $\dfrac{62}{9}$~?\end{minipage}
\end{minipage}

\vfill
\begin{tcolorbox}[enhanced,nobeforeafter,width=7.1cm,height=2.1cm,valign=center,colback=cfra!6,colframe=cfra,arc=2.5mm,boxrule=0.9pt,left=3mm]{\Large\bfseries\color{cfra}a)}\quad\begin{minipage}[c]{5.6cm}\raggedright\Large $8\text{ et }9$\end{minipage}\end{tcolorbox}\hfill\begin{tcolorbox}[enhanced,nobeforeafter,width=7.1cm,height=2.1cm,valign=center,colback=cfra!6,colframe=cfra,arc=2.5mm,boxrule=0.9pt,left=3mm]{\Large\bfseries\color{cfra}b)}\quad\begin{minipage}[c]{5.6cm}\raggedright\Large $5\text{ et }6$\end{minipage}\end{tcolorbox}\par\vspace{2.5mm}
\begin{tcolorbox}[enhanced,nobeforeafter,width=7.1cm,height=2.1cm,valign=center,colback=cfra!6,colframe=cfra,arc=2.5mm,boxrule=0.9pt,left=3mm]{\Large\bfseries\color{cfra}c)}\quad\begin{minipage}[c]{5.6cm}\raggedright\Large $6\text{ et }7$\end{minipage}\end{tcolorbox}\hfill\begin{tcolorbox}[enhanced,nobeforeafter,width=7.1cm,height=2.1cm,valign=center,colback=cfra!6,colframe=cfra,arc=2.5mm,boxrule=0.9pt,left=3mm]{\Large\bfseries\color{cfra}d)}\quad\begin{minipage}[c]{5.6cm}\raggedright\Large $7\text{ et }8$\end{minipage}\end{tcolorbox}
\end{tcolorbox}
\newpage
\begin{tcolorbox}[enhanced,width=16cm,height=14cm,colback=white,colframe=cgeo,boxrule=2pt,arc=0mm,
left=6mm,right=6mm,top=0mm,bottom=4mm,sharp corners,
overlay={\fill[cgeo] (frame.north west) rectangle ([yshift=-1.25cm]frame.north east);
\node[anchor=west,text=white,font=\sffamily\bfseries\LARGE] at ([xshift=5mm,yshift=-0.62cm]frame.north west) {Question 49 D};
\node[anchor=east,text=white,font=\sffamily\itshape\large] at ([xshift=-5mm,yshift=-0.62cm]frame.north east) {Géométrie};}]
\vspace*{1.75cm}
\begin{minipage}[t][6.9cm][c]{\linewidth}\LARGE\centering
\begin{minipage}{0.95\linewidth}\raggedright\linespread{1.25}\selectfont Un angle mesure $160^\circ$. Quelle est sa nature~?\end{minipage}
\end{minipage}

\vfill
\begin{tcolorbox}[enhanced,nobeforeafter,width=7.1cm,height=2.1cm,valign=center,colback=cgeo!6,colframe=cgeo,arc=2.5mm,boxrule=0.9pt,left=3mm]{\Large\bfseries\color{cgeo}a)}\quad\begin{minipage}[c]{5.6cm}\raggedright\Large plat\end{minipage}\end{tcolorbox}\hfill\begin{tcolorbox}[enhanced,nobeforeafter,width=7.1cm,height=2.1cm,valign=center,colback=cgeo!6,colframe=cgeo,arc=2.5mm,boxrule=0.9pt,left=3mm]{\Large\bfseries\color{cgeo}b)}\quad\begin{minipage}[c]{5.6cm}\raggedright\Large aigu\end{minipage}\end{tcolorbox}\par\vspace{2.5mm}
\begin{tcolorbox}[enhanced,nobeforeafter,width=7.1cm,height=2.1cm,valign=center,colback=cgeo!6,colframe=cgeo,arc=2.5mm,boxrule=0.9pt,left=3mm]{\Large\bfseries\color{cgeo}c)}\quad\begin{minipage}[c]{5.6cm}\raggedright\Large obtus\end{minipage}\end{tcolorbox}\hfill\begin{tcolorbox}[enhanced,nobeforeafter,width=7.1cm,height=2.1cm,valign=center,colback=cgeo!6,colframe=cgeo,arc=2.5mm,boxrule=0.9pt,left=3mm]{\Large\bfseries\color{cgeo}d)}\quad\begin{minipage}[c]{5.6cm}\raggedright\Large droit\end{minipage}\end{tcolorbox}
\end{tcolorbox}
\newpage
\begin{tcolorbox}[enhanced,width=16cm,height=14cm,colback=white,colframe=cgeo,boxrule=2pt,arc=0mm,
left=6mm,right=6mm,top=0mm,bottom=4mm,sharp corners,
overlay={\fill[cgeo] (frame.north west) rectangle ([yshift=-1.25cm]frame.north east);
\node[anchor=west,text=white,font=\sffamily\bfseries\LARGE] at ([xshift=5mm,yshift=-0.62cm]frame.north west) {Question 50 D};
\node[anchor=east,text=white,font=\sffamily\itshape\large] at ([xshift=-5mm,yshift=-0.62cm]frame.north east) {Géométrie};}]
\vspace*{1.75cm}
\begin{minipage}[t][6.9cm][c]{\linewidth}\LARGE\centering
\begin{minipage}{0.95\linewidth}\raggedright\linespread{1.25}\selectfont Un angle mesure $15^\circ$. Quelle est sa nature~?\end{minipage}
\end{minipage}

\vfill
\begin{tcolorbox}[enhanced,nobeforeafter,width=7.1cm,height=2.1cm,valign=center,colback=cgeo!6,colframe=cgeo,arc=2.5mm,boxrule=0.9pt,left=3mm]{\Large\bfseries\color{cgeo}a)}\quad\begin{minipage}[c]{5.6cm}\raggedright\Large obtus\end{minipage}\end{tcolorbox}\hfill\begin{tcolorbox}[enhanced,nobeforeafter,width=7.1cm,height=2.1cm,valign=center,colback=cgeo!6,colframe=cgeo,arc=2.5mm,boxrule=0.9pt,left=3mm]{\Large\bfseries\color{cgeo}b)}\quad\begin{minipage}[c]{5.6cm}\raggedright\Large plat\end{minipage}\end{tcolorbox}\par\vspace{2.5mm}
\begin{tcolorbox}[enhanced,nobeforeafter,width=7.1cm,height=2.1cm,valign=center,colback=cgeo!6,colframe=cgeo,arc=2.5mm,boxrule=0.9pt,left=3mm]{\Large\bfseries\color{cgeo}c)}\quad\begin{minipage}[c]{5.6cm}\raggedright\Large droit\end{minipage}\end{tcolorbox}\hfill\begin{tcolorbox}[enhanced,nobeforeafter,width=7.1cm,height=2.1cm,valign=center,colback=cgeo!6,colframe=cgeo,arc=2.5mm,boxrule=0.9pt,left=3mm]{\Large\bfseries\color{cgeo}d)}\quad\begin{minipage}[c]{5.6cm}\raggedright\Large aigu\end{minipage}\end{tcolorbox}
\end{tcolorbox}
\newpage
\begin{tcolorbox}[enhanced,width=16cm,height=14cm,colback=white,colframe=cgeo,boxrule=2pt,arc=0mm,
left=6mm,right=6mm,top=0mm,bottom=4mm,sharp corners,
overlay={\fill[cgeo] (frame.north west) rectangle ([yshift=-1.25cm]frame.north east);
\node[anchor=west,text=white,font=\sffamily\bfseries\LARGE] at ([xshift=5mm,yshift=-0.62cm]frame.north west) {Question 51 D};
\node[anchor=east,text=white,font=\sffamily\itshape\large] at ([xshift=-5mm,yshift=-0.62cm]frame.north east) {Géométrie};}]
\vspace*{1.75cm}
\begin{minipage}[t][6.9cm][c]{\linewidth}\LARGE\centering
\begin{minipage}{0.95\linewidth}\raggedright\linespread{1.25}\selectfont L'angle droit $\widehat{xOy}$ est partagé par la demi-droite $[Oz)$ en deux angles. $\widehat{xOz}=50^\circ$. Combien mesure $\widehat{zOy}$~?\end{minipage}
\end{minipage}

\vfill
\begin{tcolorbox}[enhanced,nobeforeafter,width=7.1cm,height=2.1cm,valign=center,colback=cgeo!6,colframe=cgeo,arc=2.5mm,boxrule=0.9pt,left=3mm]{\Large\bfseries\color{cgeo}a)}\quad\begin{minipage}[c]{5.6cm}\raggedright\Large $40^\circ$\end{minipage}\end{tcolorbox}\hfill\begin{tcolorbox}[enhanced,nobeforeafter,width=7.1cm,height=2.1cm,valign=center,colback=cgeo!6,colframe=cgeo,arc=2.5mm,boxrule=0.9pt,left=3mm]{\Large\bfseries\color{cgeo}b)}\quad\begin{minipage}[c]{5.6cm}\raggedright\Large $130^\circ$\end{minipage}\end{tcolorbox}\par\vspace{2.5mm}
\begin{tcolorbox}[enhanced,nobeforeafter,width=7.1cm,height=2.1cm,valign=center,colback=cgeo!6,colframe=cgeo,arc=2.5mm,boxrule=0.9pt,left=3mm]{\Large\bfseries\color{cgeo}c)}\quad\begin{minipage}[c]{5.6cm}\raggedright\Large $140^\circ$\end{minipage}\end{tcolorbox}\hfill\begin{tcolorbox}[enhanced,nobeforeafter,width=7.1cm,height=2.1cm,valign=center,colback=cgeo!6,colframe=cgeo,arc=2.5mm,boxrule=0.9pt,left=3mm]{\Large\bfseries\color{cgeo}d)}\quad\begin{minipage}[c]{5.6cm}\raggedright\Large $50^\circ$\end{minipage}\end{tcolorbox}
\end{tcolorbox}
\newpage
\begin{tcolorbox}[enhanced,width=16cm,height=14cm,colback=white,colframe=cgeo,boxrule=2pt,arc=0mm,
left=6mm,right=6mm,top=0mm,bottom=4mm,sharp corners,
overlay={\fill[cgeo] (frame.north west) rectangle ([yshift=-1.25cm]frame.north east);
\node[anchor=west,text=white,font=\sffamily\bfseries\LARGE] at ([xshift=5mm,yshift=-0.62cm]frame.north west) {Question 52 D};
\node[anchor=east,text=white,font=\sffamily\itshape\large] at ([xshift=-5mm,yshift=-0.62cm]frame.north east) {Géométrie};}]
\vspace*{1.75cm}
\begin{minipage}[t][6.9cm][c]{\linewidth}\LARGE\centering
\begin{minipage}{0.95\linewidth}\raggedright\linespread{1.25}\selectfont L'angle droit $\widehat{xOy}$ est partagé par la demi-droite $[Oz)$ en deux angles. $\widehat{xOz}=84^\circ$. Combien mesure $\widehat{zOy}$~?\end{minipage}
\end{minipage}

\vfill
\begin{tcolorbox}[enhanced,nobeforeafter,width=7.1cm,height=2.1cm,valign=center,colback=cgeo!6,colframe=cgeo,arc=2.5mm,boxrule=0.9pt,left=3mm]{\Large\bfseries\color{cgeo}a)}\quad\begin{minipage}[c]{5.6cm}\raggedright\Large $174^\circ$\end{minipage}\end{tcolorbox}\hfill\begin{tcolorbox}[enhanced,nobeforeafter,width=7.1cm,height=2.1cm,valign=center,colback=cgeo!6,colframe=cgeo,arc=2.5mm,boxrule=0.9pt,left=3mm]{\Large\bfseries\color{cgeo}b)}\quad\begin{minipage}[c]{5.6cm}\raggedright\Large $84^\circ$\end{minipage}\end{tcolorbox}\par\vspace{2.5mm}
\begin{tcolorbox}[enhanced,nobeforeafter,width=7.1cm,height=2.1cm,valign=center,colback=cgeo!6,colframe=cgeo,arc=2.5mm,boxrule=0.9pt,left=3mm]{\Large\bfseries\color{cgeo}c)}\quad\begin{minipage}[c]{5.6cm}\raggedright\Large $96^\circ$\end{minipage}\end{tcolorbox}\hfill\begin{tcolorbox}[enhanced,nobeforeafter,width=7.1cm,height=2.1cm,valign=center,colback=cgeo!6,colframe=cgeo,arc=2.5mm,boxrule=0.9pt,left=3mm]{\Large\bfseries\color{cgeo}d)}\quad\begin{minipage}[c]{5.6cm}\raggedright\Large $6^\circ$\end{minipage}\end{tcolorbox}
\end{tcolorbox}
\newpage
\begin{tcolorbox}[enhanced,width=16cm,height=14cm,colback=white,colframe=cgeo,boxrule=2pt,arc=0mm,
left=6mm,right=6mm,top=0mm,bottom=4mm,sharp corners,
overlay={\fill[cgeo] (frame.north west) rectangle ([yshift=-1.25cm]frame.north east);
\node[anchor=west,text=white,font=\sffamily\bfseries\LARGE] at ([xshift=5mm,yshift=-0.62cm]frame.north west) {Question 53 D};
\node[anchor=east,text=white,font=\sffamily\itshape\large] at ([xshift=-5mm,yshift=-0.62cm]frame.north east) {Géométrie};}]
\vspace*{1.75cm}
\begin{minipage}[t][6.9cm][c]{\linewidth}\LARGE\centering
\begin{minipage}{0.95\linewidth}\raggedright\linespread{1.25}\selectfont Les points $A$, $O$ et $B$ sont alignés dans cet ordre et $\widehat{AOC}=152^\circ$. Combien mesure $\widehat{COB}$~?\end{minipage}
\end{minipage}

\vfill
\begin{tcolorbox}[enhanced,nobeforeafter,width=7.1cm,height=2.1cm,valign=center,colback=cgeo!6,colframe=cgeo,arc=2.5mm,boxrule=0.9pt,left=3mm]{\Large\bfseries\color{cgeo}a)}\quad\begin{minipage}[c]{5.6cm}\raggedright\Large $28^\circ$\end{minipage}\end{tcolorbox}\hfill\begin{tcolorbox}[enhanced,nobeforeafter,width=7.1cm,height=2.1cm,valign=center,colback=cgeo!6,colframe=cgeo,arc=2.5mm,boxrule=0.9pt,left=3mm]{\Large\bfseries\color{cgeo}b)}\quad\begin{minipage}[c]{5.6cm}\raggedright\Large $62^\circ$\end{minipage}\end{tcolorbox}\par\vspace{2.5mm}
\begin{tcolorbox}[enhanced,nobeforeafter,width=7.1cm,height=2.1cm,valign=center,colback=cgeo!6,colframe=cgeo,arc=2.5mm,boxrule=0.9pt,left=3mm]{\Large\bfseries\color{cgeo}c)}\quad\begin{minipage}[c]{5.6cm}\raggedright\Large $208^\circ$\end{minipage}\end{tcolorbox}\hfill\begin{tcolorbox}[enhanced,nobeforeafter,width=7.1cm,height=2.1cm,valign=center,colback=cgeo!6,colframe=cgeo,arc=2.5mm,boxrule=0.9pt,left=3mm]{\Large\bfseries\color{cgeo}d)}\quad\begin{minipage}[c]{5.6cm}\raggedright\Large $152^\circ$\end{minipage}\end{tcolorbox}
\end{tcolorbox}
\newpage
\begin{tcolorbox}[enhanced,width=16cm,height=14cm,colback=white,colframe=cgeo,boxrule=2pt,arc=0mm,
left=6mm,right=6mm,top=0mm,bottom=4mm,sharp corners,
overlay={\fill[cgeo] (frame.north west) rectangle ([yshift=-1.25cm]frame.north east);
\node[anchor=west,text=white,font=\sffamily\bfseries\LARGE] at ([xshift=5mm,yshift=-0.62cm]frame.north west) {Question 54 D};
\node[anchor=east,text=white,font=\sffamily\itshape\large] at ([xshift=-5mm,yshift=-0.62cm]frame.north east) {Géométrie};}]
\vspace*{1.75cm}
\begin{minipage}[t][6.9cm][c]{\linewidth}\LARGE\centering
\begin{minipage}{0.95\linewidth}\raggedright\linespread{1.25}\selectfont Les points $A$, $O$ et $B$ sont alignés dans cet ordre et $\widehat{AOC}=49^\circ$. Combien mesure $\widehat{COB}$~?\end{minipage}
\end{minipage}

\vfill
\begin{tcolorbox}[enhanced,nobeforeafter,width=7.1cm,height=2.1cm,valign=center,colback=cgeo!6,colframe=cgeo,arc=2.5mm,boxrule=0.9pt,left=3mm]{\Large\bfseries\color{cgeo}a)}\quad\begin{minipage}[c]{5.6cm}\raggedright\Large $41^\circ$\end{minipage}\end{tcolorbox}\hfill\begin{tcolorbox}[enhanced,nobeforeafter,width=7.1cm,height=2.1cm,valign=center,colback=cgeo!6,colframe=cgeo,arc=2.5mm,boxrule=0.9pt,left=3mm]{\Large\bfseries\color{cgeo}b)}\quad\begin{minipage}[c]{5.6cm}\raggedright\Large $131^\circ$\end{minipage}\end{tcolorbox}\par\vspace{2.5mm}
\begin{tcolorbox}[enhanced,nobeforeafter,width=7.1cm,height=2.1cm,valign=center,colback=cgeo!6,colframe=cgeo,arc=2.5mm,boxrule=0.9pt,left=3mm]{\Large\bfseries\color{cgeo}c)}\quad\begin{minipage}[c]{5.6cm}\raggedright\Large $311^\circ$\end{minipage}\end{tcolorbox}\hfill\begin{tcolorbox}[enhanced,nobeforeafter,width=7.1cm,height=2.1cm,valign=center,colback=cgeo!6,colframe=cgeo,arc=2.5mm,boxrule=0.9pt,left=3mm]{\Large\bfseries\color{cgeo}d)}\quad\begin{minipage}[c]{5.6cm}\raggedright\Large $49^\circ$\end{minipage}\end{tcolorbox}
\end{tcolorbox}
\newpage
\begin{tcolorbox}[enhanced,width=16cm,height=14cm,colback=white,colframe=cgeo,boxrule=2pt,arc=0mm,
left=6mm,right=6mm,top=0mm,bottom=4mm,sharp corners,
overlay={\fill[cgeo] (frame.north west) rectangle ([yshift=-1.25cm]frame.north east);
\node[anchor=west,text=white,font=\sffamily\bfseries\LARGE] at ([xshift=5mm,yshift=-0.62cm]frame.north west) {Question 55 D};
\node[anchor=east,text=white,font=\sffamily\itshape\large] at ([xshift=-5mm,yshift=-0.62cm]frame.north east) {Géométrie};}]
\vspace*{1.75cm}
\begin{minipage}[t][6.9cm][c]{\linewidth}\LARGE\centering
\begin{minipage}{0.95\linewidth}\raggedright\linespread{1.25}\selectfont Combien d'axes de symétrie possède un hexagone régulier~?\end{minipage}
\end{minipage}

\vfill
\begin{tcolorbox}[enhanced,nobeforeafter,width=7.1cm,height=2.1cm,valign=center,colback=cgeo!6,colframe=cgeo,arc=2.5mm,boxrule=0.9pt,left=3mm]{\Large\bfseries\color{cgeo}a)}\quad\begin{minipage}[c]{5.6cm}\raggedright\Large 6\end{minipage}\end{tcolorbox}\hfill\begin{tcolorbox}[enhanced,nobeforeafter,width=7.1cm,height=2.1cm,valign=center,colback=cgeo!6,colframe=cgeo,arc=2.5mm,boxrule=0.9pt,left=3mm]{\Large\bfseries\color{cgeo}b)}\quad\begin{minipage}[c]{5.6cm}\raggedright\Large 2\end{minipage}\end{tcolorbox}\par\vspace{2.5mm}
\begin{tcolorbox}[enhanced,nobeforeafter,width=7.1cm,height=2.1cm,valign=center,colback=cgeo!6,colframe=cgeo,arc=2.5mm,boxrule=0.9pt,left=3mm]{\Large\bfseries\color{cgeo}c)}\quad\begin{minipage}[c]{5.6cm}\raggedright\Large 4\end{minipage}\end{tcolorbox}\hfill\begin{tcolorbox}[enhanced,nobeforeafter,width=7.1cm,height=2.1cm,valign=center,colback=cgeo!6,colframe=cgeo,arc=2.5mm,boxrule=0.9pt,left=3mm]{\Large\bfseries\color{cgeo}d)}\quad\begin{minipage}[c]{5.6cm}\raggedright\Large 1\end{minipage}\end{tcolorbox}
\end{tcolorbox}
\newpage
\begin{tcolorbox}[enhanced,width=16cm,height=14cm,colback=white,colframe=cgeo,boxrule=2pt,arc=0mm,
left=6mm,right=6mm,top=0mm,bottom=4mm,sharp corners,
overlay={\fill[cgeo] (frame.north west) rectangle ([yshift=-1.25cm]frame.north east);
\node[anchor=west,text=white,font=\sffamily\bfseries\LARGE] at ([xshift=5mm,yshift=-0.62cm]frame.north west) {Question 56 D};
\node[anchor=east,text=white,font=\sffamily\itshape\large] at ([xshift=-5mm,yshift=-0.62cm]frame.north east) {Géométrie};}]
\vspace*{1.75cm}
\begin{minipage}[t][6.9cm][c]{\linewidth}\LARGE\centering
\begin{minipage}{0.95\linewidth}\raggedright\linespread{1.25}\selectfont Combien d'axes de symétrie possède un cercle~?\end{minipage}
\end{minipage}

\vfill
\begin{tcolorbox}[enhanced,nobeforeafter,width=7.1cm,height=2.1cm,valign=center,colback=cgeo!6,colframe=cgeo,arc=2.5mm,boxrule=0.9pt,left=3mm]{\Large\bfseries\color{cgeo}a)}\quad\begin{minipage}[c]{5.6cm}\raggedright\Large 2\end{minipage}\end{tcolorbox}\hfill\begin{tcolorbox}[enhanced,nobeforeafter,width=7.1cm,height=2.1cm,valign=center,colback=cgeo!6,colframe=cgeo,arc=2.5mm,boxrule=0.9pt,left=3mm]{\Large\bfseries\color{cgeo}b)}\quad\begin{minipage}[c]{5.6cm}\raggedright\Large une infinité\end{minipage}\end{tcolorbox}\par\vspace{2.5mm}
\begin{tcolorbox}[enhanced,nobeforeafter,width=7.1cm,height=2.1cm,valign=center,colback=cgeo!6,colframe=cgeo,arc=2.5mm,boxrule=0.9pt,left=3mm]{\Large\bfseries\color{cgeo}c)}\quad\begin{minipage}[c]{5.6cm}\raggedright\Large 1\end{minipage}\end{tcolorbox}\hfill\begin{tcolorbox}[enhanced,nobeforeafter,width=7.1cm,height=2.1cm,valign=center,colback=cgeo!6,colframe=cgeo,arc=2.5mm,boxrule=0.9pt,left=3mm]{\Large\bfseries\color{cgeo}d)}\quad\begin{minipage}[c]{5.6cm}\raggedright\Large 4\end{minipage}\end{tcolorbox}
\end{tcolorbox}
\newpage
\begin{tcolorbox}[enhanced,width=16cm,height=14cm,colback=white,colframe=cgeo,boxrule=2pt,arc=0mm,
left=6mm,right=6mm,top=0mm,bottom=4mm,sharp corners,
overlay={\fill[cgeo] (frame.north west) rectangle ([yshift=-1.25cm]frame.north east);
\node[anchor=west,text=white,font=\sffamily\bfseries\LARGE] at ([xshift=5mm,yshift=-0.62cm]frame.north west) {Question 57 D};
\node[anchor=east,text=white,font=\sffamily\itshape\large] at ([xshift=-5mm,yshift=-0.62cm]frame.north east) {Géométrie};}]
\vspace*{1.75cm}
\begin{minipage}[t][6.9cm][c]{\linewidth}\LARGE\centering
\begin{minipage}{0.95\linewidth}\raggedright\linespread{1.25}\selectfont Un cercle a un diamètre de $7$ cm. Quelle est la longueur de son rayon~?\end{minipage}
\end{minipage}

\vfill
\begin{tcolorbox}[enhanced,nobeforeafter,width=7.1cm,height=2.1cm,valign=center,colback=cgeo!6,colframe=cgeo,arc=2.5mm,boxrule=0.9pt,left=3mm]{\Large\bfseries\color{cgeo}a)}\quad\begin{minipage}[c]{5.6cm}\raggedright\Large $7~\text{cm}$\end{minipage}\end{tcolorbox}\hfill\begin{tcolorbox}[enhanced,nobeforeafter,width=7.1cm,height=2.1cm,valign=center,colback=cgeo!6,colframe=cgeo,arc=2.5mm,boxrule=0.9pt,left=3mm]{\Large\bfseries\color{cgeo}b)}\quad\begin{minipage}[c]{5.6cm}\raggedright\Large $3{,}5~\text{cm}$\end{minipage}\end{tcolorbox}\par\vspace{2.5mm}
\begin{tcolorbox}[enhanced,nobeforeafter,width=7.1cm,height=2.1cm,valign=center,colback=cgeo!6,colframe=cgeo,arc=2.5mm,boxrule=0.9pt,left=3mm]{\Large\bfseries\color{cgeo}c)}\quad\begin{minipage}[c]{5.6cm}\raggedright\Large $9~\text{cm}$\end{minipage}\end{tcolorbox}\hfill\begin{tcolorbox}[enhanced,nobeforeafter,width=7.1cm,height=2.1cm,valign=center,colback=cgeo!6,colframe=cgeo,arc=2.5mm,boxrule=0.9pt,left=3mm]{\Large\bfseries\color{cgeo}d)}\quad\begin{minipage}[c]{5.6cm}\raggedright\Large $14~\text{cm}$\end{minipage}\end{tcolorbox}
\end{tcolorbox}
\newpage
\begin{tcolorbox}[enhanced,width=16cm,height=14cm,colback=white,colframe=cgeo,boxrule=2pt,arc=0mm,
left=6mm,right=6mm,top=0mm,bottom=4mm,sharp corners,
overlay={\fill[cgeo] (frame.north west) rectangle ([yshift=-1.25cm]frame.north east);
\node[anchor=west,text=white,font=\sffamily\bfseries\LARGE] at ([xshift=5mm,yshift=-0.62cm]frame.north west) {Question 58 D};
\node[anchor=east,text=white,font=\sffamily\itshape\large] at ([xshift=-5mm,yshift=-0.62cm]frame.north east) {Géométrie};}]
\vspace*{1.75cm}
\begin{minipage}[t][6.9cm][c]{\linewidth}\LARGE\centering
\begin{minipage}{0.95\linewidth}\raggedright\linespread{1.25}\selectfont Un cercle a un rayon de $12{,}5$ cm. Quelle est la longueur de son diamètre~?\end{minipage}
\end{minipage}

\vfill
\begin{tcolorbox}[enhanced,nobeforeafter,width=7.1cm,height=2.1cm,valign=center,colback=cgeo!6,colframe=cgeo,arc=2.5mm,boxrule=0.9pt,left=3mm]{\Large\bfseries\color{cgeo}a)}\quad\begin{minipage}[c]{5.6cm}\raggedright\Large $6{,}25~\text{cm}$\end{minipage}\end{tcolorbox}\hfill\begin{tcolorbox}[enhanced,nobeforeafter,width=7.1cm,height=2.1cm,valign=center,colback=cgeo!6,colframe=cgeo,arc=2.5mm,boxrule=0.9pt,left=3mm]{\Large\bfseries\color{cgeo}b)}\quad\begin{minipage}[c]{5.6cm}\raggedright\Large $25~\text{cm}$\end{minipage}\end{tcolorbox}\par\vspace{2.5mm}
\begin{tcolorbox}[enhanced,nobeforeafter,width=7.1cm,height=2.1cm,valign=center,colback=cgeo!6,colframe=cgeo,arc=2.5mm,boxrule=0.9pt,left=3mm]{\Large\bfseries\color{cgeo}c)}\quad\begin{minipage}[c]{5.6cm}\raggedright\Large $12{,}5~\text{cm}$\end{minipage}\end{tcolorbox}\hfill\begin{tcolorbox}[enhanced,nobeforeafter,width=7.1cm,height=2.1cm,valign=center,colback=cgeo!6,colframe=cgeo,arc=2.5mm,boxrule=0.9pt,left=3mm]{\Large\bfseries\color{cgeo}d)}\quad\begin{minipage}[c]{5.6cm}\raggedright\Large $14{,}5~\text{cm}$\end{minipage}\end{tcolorbox}
\end{tcolorbox}
\newpage
\begin{tcolorbox}[enhanced,width=16cm,height=14cm,colback=white,colframe=cgeo,boxrule=2pt,arc=0mm,
left=6mm,right=6mm,top=0mm,bottom=4mm,sharp corners,
overlay={\fill[cgeo] (frame.north west) rectangle ([yshift=-1.25cm]frame.north east);
\node[anchor=west,text=white,font=\sffamily\bfseries\LARGE] at ([xshift=5mm,yshift=-0.62cm]frame.north west) {Question 59 D};
\node[anchor=east,text=white,font=\sffamily\itshape\large] at ([xshift=-5mm,yshift=-0.62cm]frame.north east) {Géométrie};}]
\vspace*{1.75cm}
\begin{minipage}[t][6.9cm][c]{\linewidth}\LARGE\centering
\begin{minipage}{0.95\linewidth}\raggedright\linespread{1.25}\selectfont La droite $(r)$ est parallèle à $(t)$ et la droite $(s)$ est perpendiculaire à $(t)$. Que peut-on dire des droites $(r)$ et $(s)$~?\end{minipage}
\end{minipage}

\vfill
\begin{tcolorbox}[enhanced,nobeforeafter,width=7.1cm,height=2.1cm,valign=center,colback=cgeo!6,colframe=cgeo,arc=2.5mm,boxrule=0.9pt,left=3mm]{\Large\bfseries\color{cgeo}a)}\quad\begin{minipage}[c]{5.6cm}\raggedright\Large parallèles\end{minipage}\end{tcolorbox}\hfill\begin{tcolorbox}[enhanced,nobeforeafter,width=7.1cm,height=2.1cm,valign=center,colback=cgeo!6,colframe=cgeo,arc=2.5mm,boxrule=0.9pt,left=3mm]{\Large\bfseries\color{cgeo}b)}\quad\begin{minipage}[c]{5.6cm}\raggedright\large sécantes, non perpendiculaires\end{minipage}\end{tcolorbox}\par\vspace{2.5mm}
\begin{tcolorbox}[enhanced,nobeforeafter,width=7.1cm,height=2.1cm,valign=center,colback=cgeo!6,colframe=cgeo,arc=2.5mm,boxrule=0.9pt,left=3mm]{\Large\bfseries\color{cgeo}c)}\quad\begin{minipage}[c]{5.6cm}\raggedright\Large on ne peut pas savoir\end{minipage}\end{tcolorbox}\hfill\begin{tcolorbox}[enhanced,nobeforeafter,width=7.1cm,height=2.1cm,valign=center,colback=cgeo!6,colframe=cgeo,arc=2.5mm,boxrule=0.9pt,left=3mm]{\Large\bfseries\color{cgeo}d)}\quad\begin{minipage}[c]{5.6cm}\raggedright\Large perpendiculaires\end{minipage}\end{tcolorbox}
\end{tcolorbox}
\newpage
\begin{tcolorbox}[enhanced,width=16cm,height=14cm,colback=white,colframe=cgeo,boxrule=2pt,arc=0mm,
left=6mm,right=6mm,top=0mm,bottom=4mm,sharp corners,
overlay={\fill[cgeo] (frame.north west) rectangle ([yshift=-1.25cm]frame.north east);
\node[anchor=west,text=white,font=\sffamily\bfseries\LARGE] at ([xshift=5mm,yshift=-0.62cm]frame.north west) {Question 60 D};
\node[anchor=east,text=white,font=\sffamily\itshape\large] at ([xshift=-5mm,yshift=-0.62cm]frame.north east) {Géométrie};}]
\vspace*{1.75cm}
\begin{minipage}[t][6.9cm][c]{\linewidth}\LARGE\centering
\begin{minipage}{0.95\linewidth}\raggedright\linespread{1.25}\selectfont La droite $(k)$ est perpendiculaire à $(h)$ et la droite $(\ell)$ est parallèle à $(h)$. Que peut-on dire des droites $(k)$ et $(\ell)$~?\end{minipage}
\end{minipage}

\vfill
\begin{tcolorbox}[enhanced,nobeforeafter,width=7.1cm,height=2.1cm,valign=center,colback=cgeo!6,colframe=cgeo,arc=2.5mm,boxrule=0.9pt,left=3mm]{\Large\bfseries\color{cgeo}a)}\quad\begin{minipage}[c]{5.6cm}\raggedright\Large perpendiculaires\end{minipage}\end{tcolorbox}\hfill\begin{tcolorbox}[enhanced,nobeforeafter,width=7.1cm,height=2.1cm,valign=center,colback=cgeo!6,colframe=cgeo,arc=2.5mm,boxrule=0.9pt,left=3mm]{\Large\bfseries\color{cgeo}b)}\quad\begin{minipage}[c]{5.6cm}\raggedright\Large on ne peut pas savoir\end{minipage}\end{tcolorbox}\par\vspace{2.5mm}
\begin{tcolorbox}[enhanced,nobeforeafter,width=7.1cm,height=2.1cm,valign=center,colback=cgeo!6,colframe=cgeo,arc=2.5mm,boxrule=0.9pt,left=3mm]{\Large\bfseries\color{cgeo}c)}\quad\begin{minipage}[c]{5.6cm}\raggedright\large sécantes, non perpendiculaires\end{minipage}\end{tcolorbox}\hfill\begin{tcolorbox}[enhanced,nobeforeafter,width=7.1cm,height=2.1cm,valign=center,colback=cgeo!6,colframe=cgeo,arc=2.5mm,boxrule=0.9pt,left=3mm]{\Large\bfseries\color{cgeo}d)}\quad\begin{minipage}[c]{5.6cm}\raggedright\Large parallèles\end{minipage}\end{tcolorbox}
\end{tcolorbox}
\newpage
\begin{tcolorbox}[enhanced,width=16cm,height=14cm,colback=white,colframe=cgeo,boxrule=2pt,arc=0mm,
left=6mm,right=6mm,top=0mm,bottom=4mm,sharp corners,
overlay={\fill[cgeo] (frame.north west) rectangle ([yshift=-1.25cm]frame.north east);
\node[anchor=west,text=white,font=\sffamily\bfseries\LARGE] at ([xshift=5mm,yshift=-0.62cm]frame.north west) {Question 61 D};
\node[anchor=east,text=white,font=\sffamily\itshape\large] at ([xshift=-5mm,yshift=-0.62cm]frame.north east) {Géométrie};}]
\vspace*{1.75cm}
\begin{minipage}[t][6.9cm][c]{\linewidth}\LARGE\centering
\begin{minipage}{0.95\linewidth}\raggedright\linespread{1.25}\selectfont Un triangle a des côtés de $9$ cm, $9$ cm et $9$ cm. Quelle est sa nature~?\end{minipage}
\end{minipage}

\vfill
\begin{tcolorbox}[enhanced,nobeforeafter,width=7.1cm,height=2.1cm,valign=center,colback=cgeo!6,colframe=cgeo,arc=2.5mm,boxrule=0.9pt,left=3mm]{\Large\bfseries\color{cgeo}a)}\quad\begin{minipage}[c]{5.6cm}\raggedright\Large quelconque\end{minipage}\end{tcolorbox}\hfill\begin{tcolorbox}[enhanced,nobeforeafter,width=7.1cm,height=2.1cm,valign=center,colback=cgeo!6,colframe=cgeo,arc=2.5mm,boxrule=0.9pt,left=3mm]{\Large\bfseries\color{cgeo}b)}\quad\begin{minipage}[c]{5.6cm}\raggedright\Large rectangle\end{minipage}\end{tcolorbox}\par\vspace{2.5mm}
\begin{tcolorbox}[enhanced,nobeforeafter,width=7.1cm,height=2.1cm,valign=center,colback=cgeo!6,colframe=cgeo,arc=2.5mm,boxrule=0.9pt,left=3mm]{\Large\bfseries\color{cgeo}c)}\quad\begin{minipage}[c]{5.6cm}\raggedright\Large isocèle\end{minipage}\end{tcolorbox}\hfill\begin{tcolorbox}[enhanced,nobeforeafter,width=7.1cm,height=2.1cm,valign=center,colback=cgeo!6,colframe=cgeo,arc=2.5mm,boxrule=0.9pt,left=3mm]{\Large\bfseries\color{cgeo}d)}\quad\begin{minipage}[c]{5.6cm}\raggedright\Large équilatéral\end{minipage}\end{tcolorbox}
\end{tcolorbox}
\newpage
\begin{tcolorbox}[enhanced,width=16cm,height=14cm,colback=white,colframe=cgeo,boxrule=2pt,arc=0mm,
left=6mm,right=6mm,top=0mm,bottom=4mm,sharp corners,
overlay={\fill[cgeo] (frame.north west) rectangle ([yshift=-1.25cm]frame.north east);
\node[anchor=west,text=white,font=\sffamily\bfseries\LARGE] at ([xshift=5mm,yshift=-0.62cm]frame.north west) {Question 62 D};
\node[anchor=east,text=white,font=\sffamily\itshape\large] at ([xshift=-5mm,yshift=-0.62cm]frame.north east) {Géométrie};}]
\vspace*{1.75cm}
\begin{minipage}[t][6.9cm][c]{\linewidth}\LARGE\centering
\begin{minipage}{0.95\linewidth}\raggedright\linespread{1.25}\selectfont Un triangle a des côtés de $3$ cm, $5$ cm et $6$ cm, et aucun angle droit. Quelle est sa nature~?\end{minipage}
\end{minipage}

\vfill
\begin{tcolorbox}[enhanced,nobeforeafter,width=7.1cm,height=2.1cm,valign=center,colback=cgeo!6,colframe=cgeo,arc=2.5mm,boxrule=0.9pt,left=3mm]{\Large\bfseries\color{cgeo}a)}\quad\begin{minipage}[c]{5.6cm}\raggedright\Large isocèle\end{minipage}\end{tcolorbox}\hfill\begin{tcolorbox}[enhanced,nobeforeafter,width=7.1cm,height=2.1cm,valign=center,colback=cgeo!6,colframe=cgeo,arc=2.5mm,boxrule=0.9pt,left=3mm]{\Large\bfseries\color{cgeo}b)}\quad\begin{minipage}[c]{5.6cm}\raggedright\Large quelconque\end{minipage}\end{tcolorbox}\par\vspace{2.5mm}
\begin{tcolorbox}[enhanced,nobeforeafter,width=7.1cm,height=2.1cm,valign=center,colback=cgeo!6,colframe=cgeo,arc=2.5mm,boxrule=0.9pt,left=3mm]{\Large\bfseries\color{cgeo}c)}\quad\begin{minipage}[c]{5.6cm}\raggedright\Large rectangle\end{minipage}\end{tcolorbox}\hfill\begin{tcolorbox}[enhanced,nobeforeafter,width=7.1cm,height=2.1cm,valign=center,colback=cgeo!6,colframe=cgeo,arc=2.5mm,boxrule=0.9pt,left=3mm]{\Large\bfseries\color{cgeo}d)}\quad\begin{minipage}[c]{5.6cm}\raggedright\Large équilatéral\end{minipage}\end{tcolorbox}
\end{tcolorbox}
\newpage
\begin{tcolorbox}[enhanced,width=16cm,height=14cm,colback=white,colframe=cgeo,boxrule=2pt,arc=0mm,
left=6mm,right=6mm,top=0mm,bottom=4mm,sharp corners,
overlay={\fill[cgeo] (frame.north west) rectangle ([yshift=-1.25cm]frame.north east);
\node[anchor=west,text=white,font=\sffamily\bfseries\LARGE] at ([xshift=5mm,yshift=-0.62cm]frame.north west) {Question 63 D};
\node[anchor=east,text=white,font=\sffamily\itshape\large] at ([xshift=-5mm,yshift=-0.62cm]frame.north east) {Géométrie};}]
\vspace*{1.75cm}
\begin{minipage}[t][6.9cm][c]{\linewidth}\LARGE\centering
\begin{minipage}{0.95\linewidth}\raggedright\linespread{1.25}\selectfont Un quadrilatère a des diagonales de même longueur, perpendiculaires, qui se coupent en leur milieu. Quelle est sa nature la plus précise~?\end{minipage}
\end{minipage}

\vfill
\begin{tcolorbox}[enhanced,nobeforeafter,width=7.1cm,height=2.1cm,valign=center,colback=cgeo!6,colframe=cgeo,arc=2.5mm,boxrule=0.9pt,left=3mm]{\Large\bfseries\color{cgeo}a)}\quad\begin{minipage}[c]{5.6cm}\raggedright\Large un triangle\end{minipage}\end{tcolorbox}\hfill\begin{tcolorbox}[enhanced,nobeforeafter,width=7.1cm,height=2.1cm,valign=center,colback=cgeo!6,colframe=cgeo,arc=2.5mm,boxrule=0.9pt,left=3mm]{\Large\bfseries\color{cgeo}b)}\quad\begin{minipage}[c]{5.6cm}\raggedright\Large un rectangle\end{minipage}\end{tcolorbox}\par\vspace{2.5mm}
\begin{tcolorbox}[enhanced,nobeforeafter,width=7.1cm,height=2.1cm,valign=center,colback=cgeo!6,colframe=cgeo,arc=2.5mm,boxrule=0.9pt,left=3mm]{\Large\bfseries\color{cgeo}c)}\quad\begin{minipage}[c]{5.6cm}\raggedright\Large un carré\end{minipage}\end{tcolorbox}\hfill\begin{tcolorbox}[enhanced,nobeforeafter,width=7.1cm,height=2.1cm,valign=center,colback=cgeo!6,colframe=cgeo,arc=2.5mm,boxrule=0.9pt,left=3mm]{\Large\bfseries\color{cgeo}d)}\quad\begin{minipage}[c]{5.6cm}\raggedright\Large un losange\end{minipage}\end{tcolorbox}
\end{tcolorbox}
\newpage
\begin{tcolorbox}[enhanced,width=16cm,height=14cm,colback=white,colframe=cgeo,boxrule=2pt,arc=0mm,
left=6mm,right=6mm,top=0mm,bottom=4mm,sharp corners,
overlay={\fill[cgeo] (frame.north west) rectangle ([yshift=-1.25cm]frame.north east);
\node[anchor=west,text=white,font=\sffamily\bfseries\LARGE] at ([xshift=5mm,yshift=-0.62cm]frame.north west) {Question 64 D};
\node[anchor=east,text=white,font=\sffamily\itshape\large] at ([xshift=-5mm,yshift=-0.62cm]frame.north east) {Géométrie};}]
\vspace*{1.75cm}
\begin{minipage}[t][6.9cm][c]{\linewidth}\LARGE\centering
\begin{minipage}{0.95\linewidth}\raggedright\linespread{1.25}\selectfont Un quadrilatère a ses côtés opposés de même longueur et quatre angles droits. Quelle est sa nature la plus précise~?\end{minipage}
\end{minipage}

\vfill
\begin{tcolorbox}[enhanced,nobeforeafter,width=7.1cm,height=2.1cm,valign=center,colback=cgeo!6,colframe=cgeo,arc=2.5mm,boxrule=0.9pt,left=3mm]{\Large\bfseries\color{cgeo}a)}\quad\begin{minipage}[c]{5.6cm}\raggedright\Large un triangle\end{minipage}\end{tcolorbox}\hfill\begin{tcolorbox}[enhanced,nobeforeafter,width=7.1cm,height=2.1cm,valign=center,colback=cgeo!6,colframe=cgeo,arc=2.5mm,boxrule=0.9pt,left=3mm]{\Large\bfseries\color{cgeo}b)}\quad\begin{minipage}[c]{5.6cm}\raggedright\Large un losange\end{minipage}\end{tcolorbox}\par\vspace{2.5mm}
\begin{tcolorbox}[enhanced,nobeforeafter,width=7.1cm,height=2.1cm,valign=center,colback=cgeo!6,colframe=cgeo,arc=2.5mm,boxrule=0.9pt,left=3mm]{\Large\bfseries\color{cgeo}c)}\quad\begin{minipage}[c]{5.6cm}\raggedright\Large un rectangle\end{minipage}\end{tcolorbox}\hfill\begin{tcolorbox}[enhanced,nobeforeafter,width=7.1cm,height=2.1cm,valign=center,colback=cgeo!6,colframe=cgeo,arc=2.5mm,boxrule=0.9pt,left=3mm]{\Large\bfseries\color{cgeo}d)}\quad\begin{minipage}[c]{5.6cm}\raggedright\Large un carré\end{minipage}\end{tcolorbox}
\end{tcolorbox}
\newpage
\begin{tcolorbox}[enhanced,width=16cm,height=14cm,colback=white,colframe=cmes,boxrule=2pt,arc=0mm,
left=6mm,right=6mm,top=0mm,bottom=4mm,sharp corners,
overlay={\fill[cmes] (frame.north west) rectangle ([yshift=-1.25cm]frame.north east);
\node[anchor=west,text=white,font=\sffamily\bfseries\LARGE] at ([xshift=5mm,yshift=-0.62cm]frame.north west) {Question 65 D};
\node[anchor=east,text=white,font=\sffamily\itshape\large] at ([xshift=-5mm,yshift=-0.62cm]frame.north east) {Grandeurs et mesures};}]
\vspace*{1.75cm}
\begin{minipage}[t][6.9cm][c]{\linewidth}\LARGE\centering
\begin{minipage}{0.95\linewidth}\raggedright\linespread{1.25}\selectfont Un rectangle a une longueur de $20$ cm et une largeur de $11$ cm. Quel est son périmètre~?\end{minipage}
\end{minipage}

\vfill
\begin{tcolorbox}[enhanced,nobeforeafter,width=7.1cm,height=2.1cm,valign=center,colback=cmes!6,colframe=cmes,arc=2.5mm,boxrule=0.9pt,left=3mm]{\Large\bfseries\color{cmes}a)}\quad\begin{minipage}[c]{5.6cm}\raggedright\Large $62~\text{cm}$\end{minipage}\end{tcolorbox}\hfill\begin{tcolorbox}[enhanced,nobeforeafter,width=7.1cm,height=2.1cm,valign=center,colback=cmes!6,colframe=cmes,arc=2.5mm,boxrule=0.9pt,left=3mm]{\Large\bfseries\color{cmes}b)}\quad\begin{minipage}[c]{5.6cm}\raggedright\Large $31~\text{cm}$\end{minipage}\end{tcolorbox}\par\vspace{2.5mm}
\begin{tcolorbox}[enhanced,nobeforeafter,width=7.1cm,height=2.1cm,valign=center,colback=cmes!6,colframe=cmes,arc=2.5mm,boxrule=0.9pt,left=3mm]{\Large\bfseries\color{cmes}c)}\quad\begin{minipage}[c]{5.6cm}\raggedright\Large $51~\text{cm}$\end{minipage}\end{tcolorbox}\hfill\begin{tcolorbox}[enhanced,nobeforeafter,width=7.1cm,height=2.1cm,valign=center,colback=cmes!6,colframe=cmes,arc=2.5mm,boxrule=0.9pt,left=3mm]{\Large\bfseries\color{cmes}d)}\quad\begin{minipage}[c]{5.6cm}\raggedright\Large $220~\text{cm}$\end{minipage}\end{tcolorbox}
\end{tcolorbox}
\newpage
\begin{tcolorbox}[enhanced,width=16cm,height=14cm,colback=white,colframe=cmes,boxrule=2pt,arc=0mm,
left=6mm,right=6mm,top=0mm,bottom=4mm,sharp corners,
overlay={\fill[cmes] (frame.north west) rectangle ([yshift=-1.25cm]frame.north east);
\node[anchor=west,text=white,font=\sffamily\bfseries\LARGE] at ([xshift=5mm,yshift=-0.62cm]frame.north west) {Question 66 D};
\node[anchor=east,text=white,font=\sffamily\itshape\large] at ([xshift=-5mm,yshift=-0.62cm]frame.north east) {Grandeurs et mesures};}]
\vspace*{1.75cm}
\begin{minipage}[t][6.9cm][c]{\linewidth}\LARGE\centering
\begin{minipage}{0.95\linewidth}\raggedright\linespread{1.25}\selectfont Un rectangle a une longueur de $9{,}5$ cm et une largeur de $6$ cm. Quel est son périmètre~?\end{minipage}
\end{minipage}

\vfill
\begin{tcolorbox}[enhanced,nobeforeafter,width=7.1cm,height=2.1cm,valign=center,colback=cmes!6,colframe=cmes,arc=2.5mm,boxrule=0.9pt,left=3mm]{\Large\bfseries\color{cmes}a)}\quad\begin{minipage}[c]{5.6cm}\raggedright\Large $15{,}5~\text{cm}$\end{minipage}\end{tcolorbox}\hfill\begin{tcolorbox}[enhanced,nobeforeafter,width=7.1cm,height=2.1cm,valign=center,colback=cmes!6,colframe=cmes,arc=2.5mm,boxrule=0.9pt,left=3mm]{\Large\bfseries\color{cmes}b)}\quad\begin{minipage}[c]{5.6cm}\raggedright\Large $31~\text{cm}$\end{minipage}\end{tcolorbox}\par\vspace{2.5mm}
\begin{tcolorbox}[enhanced,nobeforeafter,width=7.1cm,height=2.1cm,valign=center,colback=cmes!6,colframe=cmes,arc=2.5mm,boxrule=0.9pt,left=3mm]{\Large\bfseries\color{cmes}c)}\quad\begin{minipage}[c]{5.6cm}\raggedright\Large $57~\text{cm}$\end{minipage}\end{tcolorbox}\hfill\begin{tcolorbox}[enhanced,nobeforeafter,width=7.1cm,height=2.1cm,valign=center,colback=cmes!6,colframe=cmes,arc=2.5mm,boxrule=0.9pt,left=3mm]{\Large\bfseries\color{cmes}d)}\quad\begin{minipage}[c]{5.6cm}\raggedright\Large $25~\text{cm}$\end{minipage}\end{tcolorbox}
\end{tcolorbox}
\newpage
\begin{tcolorbox}[enhanced,width=16cm,height=14cm,colback=white,colframe=cmes,boxrule=2pt,arc=0mm,
left=6mm,right=6mm,top=0mm,bottom=4mm,sharp corners,
overlay={\fill[cmes] (frame.north west) rectangle ([yshift=-1.25cm]frame.north east);
\node[anchor=west,text=white,font=\sffamily\bfseries\LARGE] at ([xshift=5mm,yshift=-0.62cm]frame.north west) {Question 67 D};
\node[anchor=east,text=white,font=\sffamily\itshape\large] at ([xshift=-5mm,yshift=-0.62cm]frame.north east) {Grandeurs et mesures};}]
\vspace*{1.75cm}
\begin{minipage}[t][6.9cm][c]{\linewidth}\LARGE\centering
\begin{minipage}{0.95\linewidth}\raggedright\linespread{1.25}\selectfont Un rectangle a une longueur de $20$ cm et une largeur de $11$ cm. Quelle est son aire~?\end{minipage}
\end{minipage}

\vfill
\begin{tcolorbox}[enhanced,nobeforeafter,width=7.1cm,height=2.1cm,valign=center,colback=cmes!6,colframe=cmes,arc=2.5mm,boxrule=0.9pt,left=3mm]{\Large\bfseries\color{cmes}a)}\quad\begin{minipage}[c]{5.6cm}\raggedright\Large $110~\text{cm}^2$\end{minipage}\end{tcolorbox}\hfill\begin{tcolorbox}[enhanced,nobeforeafter,width=7.1cm,height=2.1cm,valign=center,colback=cmes!6,colframe=cmes,arc=2.5mm,boxrule=0.9pt,left=3mm]{\Large\bfseries\color{cmes}b)}\quad\begin{minipage}[c]{5.6cm}\raggedright\Large $62~\text{cm}^2$\end{minipage}\end{tcolorbox}\par\vspace{2.5mm}
\begin{tcolorbox}[enhanced,nobeforeafter,width=7.1cm,height=2.1cm,valign=center,colback=cmes!6,colframe=cmes,arc=2.5mm,boxrule=0.9pt,left=3mm]{\Large\bfseries\color{cmes}c)}\quad\begin{minipage}[c]{5.6cm}\raggedright\Large $220~\text{cm}^2$\end{minipage}\end{tcolorbox}\hfill\begin{tcolorbox}[enhanced,nobeforeafter,width=7.1cm,height=2.1cm,valign=center,colback=cmes!6,colframe=cmes,arc=2.5mm,boxrule=0.9pt,left=3mm]{\Large\bfseries\color{cmes}d)}\quad\begin{minipage}[c]{5.6cm}\raggedright\Large $31~\text{cm}^2$\end{minipage}\end{tcolorbox}
\end{tcolorbox}
\newpage
\begin{tcolorbox}[enhanced,width=16cm,height=14cm,colback=white,colframe=cmes,boxrule=2pt,arc=0mm,
left=6mm,right=6mm,top=0mm,bottom=4mm,sharp corners,
overlay={\fill[cmes] (frame.north west) rectangle ([yshift=-1.25cm]frame.north east);
\node[anchor=west,text=white,font=\sffamily\bfseries\LARGE] at ([xshift=5mm,yshift=-0.62cm]frame.north west) {Question 68 D};
\node[anchor=east,text=white,font=\sffamily\itshape\large] at ([xshift=-5mm,yshift=-0.62cm]frame.north east) {Grandeurs et mesures};}]
\vspace*{1.75cm}
\begin{minipage}[t][6.9cm][c]{\linewidth}\LARGE\centering
\begin{minipage}{0.95\linewidth}\raggedright\linespread{1.25}\selectfont Un rectangle a une longueur de $9{,}5$ cm et une largeur de $6$ cm. Quelle est son aire~?\end{minipage}
\end{minipage}

\vfill
\begin{tcolorbox}[enhanced,nobeforeafter,width=7.1cm,height=2.1cm,valign=center,colback=cmes!6,colframe=cmes,arc=2.5mm,boxrule=0.9pt,left=3mm]{\Large\bfseries\color{cmes}a)}\quad\begin{minipage}[c]{5.6cm}\raggedright\Large $15{,}5~\text{cm}^2$\end{minipage}\end{tcolorbox}\hfill\begin{tcolorbox}[enhanced,nobeforeafter,width=7.1cm,height=2.1cm,valign=center,colback=cmes!6,colframe=cmes,arc=2.5mm,boxrule=0.9pt,left=3mm]{\Large\bfseries\color{cmes}b)}\quad\begin{minipage}[c]{5.6cm}\raggedright\Large $57~\text{cm}^2$\end{minipage}\end{tcolorbox}\par\vspace{2.5mm}
\begin{tcolorbox}[enhanced,nobeforeafter,width=7.1cm,height=2.1cm,valign=center,colback=cmes!6,colframe=cmes,arc=2.5mm,boxrule=0.9pt,left=3mm]{\Large\bfseries\color{cmes}c)}\quad\begin{minipage}[c]{5.6cm}\raggedright\Large $31~\text{cm}^2$\end{minipage}\end{tcolorbox}\hfill\begin{tcolorbox}[enhanced,nobeforeafter,width=7.1cm,height=2.1cm,valign=center,colback=cmes!6,colframe=cmes,arc=2.5mm,boxrule=0.9pt,left=3mm]{\Large\bfseries\color{cmes}d)}\quad\begin{minipage}[c]{5.6cm}\raggedright\Large $28{,}5~\text{cm}^2$\end{minipage}\end{tcolorbox}
\end{tcolorbox}
\newpage
\begin{tcolorbox}[enhanced,width=16cm,height=14cm,colback=white,colframe=cmes,boxrule=2pt,arc=0mm,
left=6mm,right=6mm,top=0mm,bottom=4mm,sharp corners,
overlay={\fill[cmes] (frame.north west) rectangle ([yshift=-1.25cm]frame.north east);
\node[anchor=west,text=white,font=\sffamily\bfseries\LARGE] at ([xshift=5mm,yshift=-0.62cm]frame.north west) {Question 69 D};
\node[anchor=east,text=white,font=\sffamily\itshape\large] at ([xshift=-5mm,yshift=-0.62cm]frame.north east) {Grandeurs et mesures};}]
\vspace*{1.75cm}
\begin{minipage}[t][6.9cm][c]{\linewidth}\LARGE\centering
\begin{minipage}{0.95\linewidth}\raggedright\linespread{1.25}\selectfont Un triangle est rectangle ; les deux côtés de l'angle droit mesurent $9$ cm et $8$ cm. Quelle est son aire~?\end{minipage}
\end{minipage}

\vfill
\begin{tcolorbox}[enhanced,nobeforeafter,width=7.1cm,height=2.1cm,valign=center,colback=cmes!6,colframe=cmes,arc=2.5mm,boxrule=0.9pt,left=3mm]{\Large\bfseries\color{cmes}a)}\quad\begin{minipage}[c]{5.6cm}\raggedright\Large $34~\text{cm}^2$\end{minipage}\end{tcolorbox}\hfill\begin{tcolorbox}[enhanced,nobeforeafter,width=7.1cm,height=2.1cm,valign=center,colback=cmes!6,colframe=cmes,arc=2.5mm,boxrule=0.9pt,left=3mm]{\Large\bfseries\color{cmes}b)}\quad\begin{minipage}[c]{5.6cm}\raggedright\Large $36~\text{cm}^2$\end{minipage}\end{tcolorbox}\par\vspace{2.5mm}
\begin{tcolorbox}[enhanced,nobeforeafter,width=7.1cm,height=2.1cm,valign=center,colback=cmes!6,colframe=cmes,arc=2.5mm,boxrule=0.9pt,left=3mm]{\Large\bfseries\color{cmes}c)}\quad\begin{minipage}[c]{5.6cm}\raggedright\Large $17~\text{cm}^2$\end{minipage}\end{tcolorbox}\hfill\begin{tcolorbox}[enhanced,nobeforeafter,width=7.1cm,height=2.1cm,valign=center,colback=cmes!6,colframe=cmes,arc=2.5mm,boxrule=0.9pt,left=3mm]{\Large\bfseries\color{cmes}d)}\quad\begin{minipage}[c]{5.6cm}\raggedright\Large $72~\text{cm}^2$\end{minipage}\end{tcolorbox}
\end{tcolorbox}
\newpage
\begin{tcolorbox}[enhanced,width=16cm,height=14cm,colback=white,colframe=cmes,boxrule=2pt,arc=0mm,
left=6mm,right=6mm,top=0mm,bottom=4mm,sharp corners,
overlay={\fill[cmes] (frame.north west) rectangle ([yshift=-1.25cm]frame.north east);
\node[anchor=west,text=white,font=\sffamily\bfseries\LARGE] at ([xshift=5mm,yshift=-0.62cm]frame.north west) {Question 70 D};
\node[anchor=east,text=white,font=\sffamily\itshape\large] at ([xshift=-5mm,yshift=-0.62cm]frame.north east) {Grandeurs et mesures};}]
\vspace*{1.75cm}
\begin{minipage}[t][6.9cm][c]{\linewidth}\LARGE\centering
\begin{minipage}{0.95\linewidth}\raggedright\linespread{1.25}\selectfont Un triangle est rectangle ; les deux côtés de l'angle droit mesurent $7$ cm et $4$ cm. Quelle est son aire~?\end{minipage}
\end{minipage}

\vfill
\begin{tcolorbox}[enhanced,nobeforeafter,width=7.1cm,height=2.1cm,valign=center,colback=cmes!6,colframe=cmes,arc=2.5mm,boxrule=0.9pt,left=3mm]{\Large\bfseries\color{cmes}a)}\quad\begin{minipage}[c]{5.6cm}\raggedright\Large $22~\text{cm}^2$\end{minipage}\end{tcolorbox}\hfill\begin{tcolorbox}[enhanced,nobeforeafter,width=7.1cm,height=2.1cm,valign=center,colback=cmes!6,colframe=cmes,arc=2.5mm,boxrule=0.9pt,left=3mm]{\Large\bfseries\color{cmes}b)}\quad\begin{minipage}[c]{5.6cm}\raggedright\Large $14~\text{cm}^2$\end{minipage}\end{tcolorbox}\par\vspace{2.5mm}
\begin{tcolorbox}[enhanced,nobeforeafter,width=7.1cm,height=2.1cm,valign=center,colback=cmes!6,colframe=cmes,arc=2.5mm,boxrule=0.9pt,left=3mm]{\Large\bfseries\color{cmes}c)}\quad\begin{minipage}[c]{5.6cm}\raggedright\Large $28~\text{cm}^2$\end{minipage}\end{tcolorbox}\hfill\begin{tcolorbox}[enhanced,nobeforeafter,width=7.1cm,height=2.1cm,valign=center,colback=cmes!6,colframe=cmes,arc=2.5mm,boxrule=0.9pt,left=3mm]{\Large\bfseries\color{cmes}d)}\quad\begin{minipage}[c]{5.6cm}\raggedright\Large $11~\text{cm}^2$\end{minipage}\end{tcolorbox}
\end{tcolorbox}
\newpage
\begin{tcolorbox}[enhanced,width=16cm,height=14cm,colback=white,colframe=cmes,boxrule=2pt,arc=0mm,
left=6mm,right=6mm,top=0mm,bottom=4mm,sharp corners,
overlay={\fill[cmes] (frame.north west) rectangle ([yshift=-1.25cm]frame.north east);
\node[anchor=west,text=white,font=\sffamily\bfseries\LARGE] at ([xshift=5mm,yshift=-0.62cm]frame.north west) {Question 71 D};
\node[anchor=east,text=white,font=\sffamily\itshape\large] at ([xshift=-5mm,yshift=-0.62cm]frame.north east) {Grandeurs et mesures};}]
\vspace*{1.75cm}
\begin{minipage}[t][6.9cm][c]{\linewidth}\LARGE\centering
\begin{minipage}{0.95\linewidth}\raggedright\linespread{1.25}\selectfont Compléter~: $2{,}3~\text{m}=\;?\;\text{mm}$\end{minipage}
\end{minipage}

\vfill
\begin{tcolorbox}[enhanced,nobeforeafter,width=7.1cm,height=2.1cm,valign=center,colback=cmes!6,colframe=cmes,arc=2.5mm,boxrule=0.9pt,left=3mm]{\Large\bfseries\color{cmes}a)}\quad\begin{minipage}[c]{5.6cm}\raggedright\Large $23\,000~\text{mm}$\end{minipage}\end{tcolorbox}\hfill\begin{tcolorbox}[enhanced,nobeforeafter,width=7.1cm,height=2.1cm,valign=center,colback=cmes!6,colframe=cmes,arc=2.5mm,boxrule=0.9pt,left=3mm]{\Large\bfseries\color{cmes}b)}\quad\begin{minipage}[c]{5.6cm}\raggedright\Large $230\,000~\text{mm}$\end{minipage}\end{tcolorbox}\par\vspace{2.5mm}
\begin{tcolorbox}[enhanced,nobeforeafter,width=7.1cm,height=2.1cm,valign=center,colback=cmes!6,colframe=cmes,arc=2.5mm,boxrule=0.9pt,left=3mm]{\Large\bfseries\color{cmes}c)}\quad\begin{minipage}[c]{5.6cm}\raggedright\Large $2\,300~\text{mm}$\end{minipage}\end{tcolorbox}\hfill\begin{tcolorbox}[enhanced,nobeforeafter,width=7.1cm,height=2.1cm,valign=center,colback=cmes!6,colframe=cmes,arc=2.5mm,boxrule=0.9pt,left=3mm]{\Large\bfseries\color{cmes}d)}\quad\begin{minipage}[c]{5.6cm}\raggedright\Large $230~\text{mm}$\end{minipage}\end{tcolorbox}
\end{tcolorbox}
\newpage
\begin{tcolorbox}[enhanced,width=16cm,height=14cm,colback=white,colframe=cmes,boxrule=2pt,arc=0mm,
left=6mm,right=6mm,top=0mm,bottom=4mm,sharp corners,
overlay={\fill[cmes] (frame.north west) rectangle ([yshift=-1.25cm]frame.north east);
\node[anchor=west,text=white,font=\sffamily\bfseries\LARGE] at ([xshift=5mm,yshift=-0.62cm]frame.north west) {Question 72 D};
\node[anchor=east,text=white,font=\sffamily\itshape\large] at ([xshift=-5mm,yshift=-0.62cm]frame.north east) {Grandeurs et mesures};}]
\vspace*{1.75cm}
\begin{minipage}[t][6.9cm][c]{\linewidth}\LARGE\centering
\begin{minipage}{0.95\linewidth}\raggedright\linespread{1.25}\selectfont Compléter~: $5\,000~\text{mm}=\;?\;\text{m}$\end{minipage}
\end{minipage}

\vfill
\begin{tcolorbox}[enhanced,nobeforeafter,width=7.1cm,height=2.1cm,valign=center,colback=cmes!6,colframe=cmes,arc=2.5mm,boxrule=0.9pt,left=3mm]{\Large\bfseries\color{cmes}a)}\quad\begin{minipage}[c]{5.6cm}\raggedright\Large $50~\text{m}$\end{minipage}\end{tcolorbox}\hfill\begin{tcolorbox}[enhanced,nobeforeafter,width=7.1cm,height=2.1cm,valign=center,colback=cmes!6,colframe=cmes,arc=2.5mm,boxrule=0.9pt,left=3mm]{\Large\bfseries\color{cmes}b)}\quad\begin{minipage}[c]{5.6cm}\raggedright\Large $5~\text{m}$\end{minipage}\end{tcolorbox}\par\vspace{2.5mm}
\begin{tcolorbox}[enhanced,nobeforeafter,width=7.1cm,height=2.1cm,valign=center,colback=cmes!6,colframe=cmes,arc=2.5mm,boxrule=0.9pt,left=3mm]{\Large\bfseries\color{cmes}c)}\quad\begin{minipage}[c]{5.6cm}\raggedright\Large $0{,}5~\text{m}$\end{minipage}\end{tcolorbox}\hfill\begin{tcolorbox}[enhanced,nobeforeafter,width=7.1cm,height=2.1cm,valign=center,colback=cmes!6,colframe=cmes,arc=2.5mm,boxrule=0.9pt,left=3mm]{\Large\bfseries\color{cmes}d)}\quad\begin{minipage}[c]{5.6cm}\raggedright\Large $500~\text{m}$\end{minipage}\end{tcolorbox}
\end{tcolorbox}
\newpage
\begin{tcolorbox}[enhanced,width=16cm,height=14cm,colback=white,colframe=cmes,boxrule=2pt,arc=0mm,
left=6mm,right=6mm,top=0mm,bottom=4mm,sharp corners,
overlay={\fill[cmes] (frame.north west) rectangle ([yshift=-1.25cm]frame.north east);
\node[anchor=west,text=white,font=\sffamily\bfseries\LARGE] at ([xshift=5mm,yshift=-0.62cm]frame.north west) {Question 73 D};
\node[anchor=east,text=white,font=\sffamily\itshape\large] at ([xshift=-5mm,yshift=-0.62cm]frame.north east) {Grandeurs et mesures};}]
\vspace*{1.75cm}
\begin{minipage}[t][6.9cm][c]{\linewidth}\LARGE\centering
\begin{minipage}{0.95\linewidth}\raggedright\linespread{1.25}\selectfont Compléter~: $250~\text{mg}=\;?\;\text{g}$\end{minipage}
\end{minipage}

\vfill
\begin{tcolorbox}[enhanced,nobeforeafter,width=7.1cm,height=2.1cm,valign=center,colback=cmes!6,colframe=cmes,arc=2.5mm,boxrule=0.9pt,left=3mm]{\Large\bfseries\color{cmes}a)}\quad\begin{minipage}[c]{5.6cm}\raggedright\Large $25~\text{g}$\end{minipage}\end{tcolorbox}\hfill\begin{tcolorbox}[enhanced,nobeforeafter,width=7.1cm,height=2.1cm,valign=center,colback=cmes!6,colframe=cmes,arc=2.5mm,boxrule=0.9pt,left=3mm]{\Large\bfseries\color{cmes}b)}\quad\begin{minipage}[c]{5.6cm}\raggedright\Large $0{,}025~\text{g}$\end{minipage}\end{tcolorbox}\par\vspace{2.5mm}
\begin{tcolorbox}[enhanced,nobeforeafter,width=7.1cm,height=2.1cm,valign=center,colback=cmes!6,colframe=cmes,arc=2.5mm,boxrule=0.9pt,left=3mm]{\Large\bfseries\color{cmes}c)}\quad\begin{minipage}[c]{5.6cm}\raggedright\Large $0{,}25~\text{g}$\end{minipage}\end{tcolorbox}\hfill\begin{tcolorbox}[enhanced,nobeforeafter,width=7.1cm,height=2.1cm,valign=center,colback=cmes!6,colframe=cmes,arc=2.5mm,boxrule=0.9pt,left=3mm]{\Large\bfseries\color{cmes}d)}\quad\begin{minipage}[c]{5.6cm}\raggedright\Large $2{,}5~\text{g}$\end{minipage}\end{tcolorbox}
\end{tcolorbox}
\newpage
\begin{tcolorbox}[enhanced,width=16cm,height=14cm,colback=white,colframe=cmes,boxrule=2pt,arc=0mm,
left=6mm,right=6mm,top=0mm,bottom=4mm,sharp corners,
overlay={\fill[cmes] (frame.north west) rectangle ([yshift=-1.25cm]frame.north east);
\node[anchor=west,text=white,font=\sffamily\bfseries\LARGE] at ([xshift=5mm,yshift=-0.62cm]frame.north west) {Question 74 D};
\node[anchor=east,text=white,font=\sffamily\itshape\large] at ([xshift=-5mm,yshift=-0.62cm]frame.north east) {Grandeurs et mesures};}]
\vspace*{1.75cm}
\begin{minipage}[t][6.9cm][c]{\linewidth}\LARGE\centering
\begin{minipage}{0.95\linewidth}\raggedright\linespread{1.25}\selectfont Compléter~: $4{,}05~\text{kg}=\;?\;\text{g}$\end{minipage}
\end{minipage}

\vfill
\begin{tcolorbox}[enhanced,nobeforeafter,width=7.1cm,height=2.1cm,valign=center,colback=cmes!6,colframe=cmes,arc=2.5mm,boxrule=0.9pt,left=3mm]{\Large\bfseries\color{cmes}a)}\quad\begin{minipage}[c]{5.6cm}\raggedright\Large $4\,050~\text{g}$\end{minipage}\end{tcolorbox}\hfill\begin{tcolorbox}[enhanced,nobeforeafter,width=7.1cm,height=2.1cm,valign=center,colback=cmes!6,colframe=cmes,arc=2.5mm,boxrule=0.9pt,left=3mm]{\Large\bfseries\color{cmes}b)}\quad\begin{minipage}[c]{5.6cm}\raggedright\Large $405~\text{g}$\end{minipage}\end{tcolorbox}\par\vspace{2.5mm}
\begin{tcolorbox}[enhanced,nobeforeafter,width=7.1cm,height=2.1cm,valign=center,colback=cmes!6,colframe=cmes,arc=2.5mm,boxrule=0.9pt,left=3mm]{\Large\bfseries\color{cmes}c)}\quad\begin{minipage}[c]{5.6cm}\raggedright\Large $405\,000~\text{g}$\end{minipage}\end{tcolorbox}\hfill\begin{tcolorbox}[enhanced,nobeforeafter,width=7.1cm,height=2.1cm,valign=center,colback=cmes!6,colframe=cmes,arc=2.5mm,boxrule=0.9pt,left=3mm]{\Large\bfseries\color{cmes}d)}\quad\begin{minipage}[c]{5.6cm}\raggedright\Large $40\,500~\text{g}$\end{minipage}\end{tcolorbox}
\end{tcolorbox}
\newpage
\begin{tcolorbox}[enhanced,width=16cm,height=14cm,colback=white,colframe=cmes,boxrule=2pt,arc=0mm,
left=6mm,right=6mm,top=0mm,bottom=4mm,sharp corners,
overlay={\fill[cmes] (frame.north west) rectangle ([yshift=-1.25cm]frame.north east);
\node[anchor=west,text=white,font=\sffamily\bfseries\LARGE] at ([xshift=5mm,yshift=-0.62cm]frame.north west) {Question 75 D};
\node[anchor=east,text=white,font=\sffamily\itshape\large] at ([xshift=-5mm,yshift=-0.62cm]frame.north east) {Grandeurs et mesures};}]
\vspace*{1.75cm}
\begin{minipage}[t][6.9cm][c]{\linewidth}\LARGE\centering
\begin{minipage}{0.95\linewidth}\raggedright\linespread{1.25}\selectfont Combien de minutes y a-t-il dans 5 h~?\end{minipage}
\end{minipage}

\vfill
\begin{tcolorbox}[enhanced,nobeforeafter,width=7.1cm,height=2.1cm,valign=center,colback=cmes!6,colframe=cmes,arc=2.5mm,boxrule=0.9pt,left=3mm]{\Large\bfseries\color{cmes}a)}\quad\begin{minipage}[c]{5.6cm}\raggedright\Large $300~\text{min}$\end{minipage}\end{tcolorbox}\hfill\begin{tcolorbox}[enhanced,nobeforeafter,width=7.1cm,height=2.1cm,valign=center,colback=cmes!6,colframe=cmes,arc=2.5mm,boxrule=0.9pt,left=3mm]{\Large\bfseries\color{cmes}b)}\quad\begin{minipage}[c]{5.6cm}\raggedright\Large $5~\text{min}$\end{minipage}\end{tcolorbox}\par\vspace{2.5mm}
\begin{tcolorbox}[enhanced,nobeforeafter,width=7.1cm,height=2.1cm,valign=center,colback=cmes!6,colframe=cmes,arc=2.5mm,boxrule=0.9pt,left=3mm]{\Large\bfseries\color{cmes}c)}\quad\begin{minipage}[c]{5.6cm}\raggedright\Large $330~\text{min}$\end{minipage}\end{tcolorbox}\hfill\begin{tcolorbox}[enhanced,nobeforeafter,width=7.1cm,height=2.1cm,valign=center,colback=cmes!6,colframe=cmes,arc=2.5mm,boxrule=0.9pt,left=3mm]{\Large\bfseries\color{cmes}d)}\quad\begin{minipage}[c]{5.6cm}\raggedright\Large $500~\text{min}$\end{minipage}\end{tcolorbox}
\end{tcolorbox}
\newpage
\begin{tcolorbox}[enhanced,width=16cm,height=14cm,colback=white,colframe=cmes,boxrule=2pt,arc=0mm,
left=6mm,right=6mm,top=0mm,bottom=4mm,sharp corners,
overlay={\fill[cmes] (frame.north west) rectangle ([yshift=-1.25cm]frame.north east);
\node[anchor=west,text=white,font=\sffamily\bfseries\LARGE] at ([xshift=5mm,yshift=-0.62cm]frame.north west) {Question 76 D};
\node[anchor=east,text=white,font=\sffamily\itshape\large] at ([xshift=-5mm,yshift=-0.62cm]frame.north east) {Grandeurs et mesures};}]
\vspace*{1.75cm}
\begin{minipage}[t][6.9cm][c]{\linewidth}\LARGE\centering
\begin{minipage}{0.95\linewidth}\raggedright\linespread{1.25}\selectfont Combien de minutes y a-t-il dans 3 h 40 min~?\end{minipage}
\end{minipage}

\vfill
\begin{tcolorbox}[enhanced,nobeforeafter,width=7.1cm,height=2.1cm,valign=center,colback=cmes!6,colframe=cmes,arc=2.5mm,boxrule=0.9pt,left=3mm]{\Large\bfseries\color{cmes}a)}\quad\begin{minipage}[c]{5.6cm}\raggedright\Large $43~\text{min}$\end{minipage}\end{tcolorbox}\hfill\begin{tcolorbox}[enhanced,nobeforeafter,width=7.1cm,height=2.1cm,valign=center,colback=cmes!6,colframe=cmes,arc=2.5mm,boxrule=0.9pt,left=3mm]{\Large\bfseries\color{cmes}b)}\quad\begin{minipage}[c]{5.6cm}\raggedright\Large $140~\text{min}$\end{minipage}\end{tcolorbox}\par\vspace{2.5mm}
\begin{tcolorbox}[enhanced,nobeforeafter,width=7.1cm,height=2.1cm,valign=center,colback=cmes!6,colframe=cmes,arc=2.5mm,boxrule=0.9pt,left=3mm]{\Large\bfseries\color{cmes}c)}\quad\begin{minipage}[c]{5.6cm}\raggedright\Large $340~\text{min}$\end{minipage}\end{tcolorbox}\hfill\begin{tcolorbox}[enhanced,nobeforeafter,width=7.1cm,height=2.1cm,valign=center,colback=cmes!6,colframe=cmes,arc=2.5mm,boxrule=0.9pt,left=3mm]{\Large\bfseries\color{cmes}d)}\quad\begin{minipage}[c]{5.6cm}\raggedright\Large $220~\text{min}$\end{minipage}\end{tcolorbox}
\end{tcolorbox}
\newpage
\begin{tcolorbox}[enhanced,width=16cm,height=14cm,colback=white,colframe=cmes,boxrule=2pt,arc=0mm,
left=6mm,right=6mm,top=0mm,bottom=4mm,sharp corners,
overlay={\fill[cmes] (frame.north west) rectangle ([yshift=-1.25cm]frame.north east);
\node[anchor=west,text=white,font=\sffamily\bfseries\LARGE] at ([xshift=5mm,yshift=-0.62cm]frame.north west) {Question 77 D};
\node[anchor=east,text=white,font=\sffamily\itshape\large] at ([xshift=-5mm,yshift=-0.62cm]frame.north east) {Grandeurs et mesures};}]
\vspace*{1.75cm}
\begin{minipage}[t][6.9cm][c]{\linewidth}\LARGE\centering
\begin{minipage}{0.95\linewidth}\raggedright\linespread{1.25}\selectfont En prenant $\pi\approx3{,}14$, quel est le périmètre d'un cercle de diamètre $30$ cm~?\end{minipage}
\end{minipage}

\vfill
\begin{tcolorbox}[enhanced,nobeforeafter,width=7.1cm,height=2.1cm,valign=center,colback=cmes!6,colframe=cmes,arc=2.5mm,boxrule=0.9pt,left=3mm]{\Large\bfseries\color{cmes}a)}\quad\begin{minipage}[c]{5.6cm}\raggedright\Large $47{,}1~\text{cm}$\end{minipage}\end{tcolorbox}\hfill\begin{tcolorbox}[enhanced,nobeforeafter,width=7.1cm,height=2.1cm,valign=center,colback=cmes!6,colframe=cmes,arc=2.5mm,boxrule=0.9pt,left=3mm]{\Large\bfseries\color{cmes}b)}\quad\begin{minipage}[c]{5.6cm}\raggedright\Large $188{,}4~\text{cm}$\end{minipage}\end{tcolorbox}\par\vspace{2.5mm}
\begin{tcolorbox}[enhanced,nobeforeafter,width=7.1cm,height=2.1cm,valign=center,colback=cmes!6,colframe=cmes,arc=2.5mm,boxrule=0.9pt,left=3mm]{\Large\bfseries\color{cmes}c)}\quad\begin{minipage}[c]{5.6cm}\raggedright\Large $706{,}5~\text{cm}$\end{minipage}\end{tcolorbox}\hfill\begin{tcolorbox}[enhanced,nobeforeafter,width=7.1cm,height=2.1cm,valign=center,colback=cmes!6,colframe=cmes,arc=2.5mm,boxrule=0.9pt,left=3mm]{\Large\bfseries\color{cmes}d)}\quad\begin{minipage}[c]{5.6cm}\raggedright\Large $94{,}2~\text{cm}$\end{minipage}\end{tcolorbox}
\end{tcolorbox}
\newpage
\begin{tcolorbox}[enhanced,width=16cm,height=14cm,colback=white,colframe=cmes,boxrule=2pt,arc=0mm,
left=6mm,right=6mm,top=0mm,bottom=4mm,sharp corners,
overlay={\fill[cmes] (frame.north west) rectangle ([yshift=-1.25cm]frame.north east);
\node[anchor=west,text=white,font=\sffamily\bfseries\LARGE] at ([xshift=5mm,yshift=-0.62cm]frame.north west) {Question 78 D};
\node[anchor=east,text=white,font=\sffamily\itshape\large] at ([xshift=-5mm,yshift=-0.62cm]frame.north east) {Grandeurs et mesures};}]
\vspace*{1.75cm}
\begin{minipage}[t][6.9cm][c]{\linewidth}\LARGE\centering
\begin{minipage}{0.95\linewidth}\raggedright\linespread{1.25}\selectfont En prenant $\pi\approx3{,}14$, quel est le périmètre d'un cercle de diamètre $2{,}5$ cm~?\end{minipage}
\end{minipage}

\vfill
\begin{tcolorbox}[enhanced,nobeforeafter,width=7.1cm,height=2.1cm,valign=center,colback=cmes!6,colframe=cmes,arc=2.5mm,boxrule=0.9pt,left=3mm]{\Large\bfseries\color{cmes}a)}\quad\begin{minipage}[c]{5.6cm}\raggedright\Large $15{,}7~\text{cm}$\end{minipage}\end{tcolorbox}\hfill\begin{tcolorbox}[enhanced,nobeforeafter,width=7.1cm,height=2.1cm,valign=center,colback=cmes!6,colframe=cmes,arc=2.5mm,boxrule=0.9pt,left=3mm]{\Large\bfseries\color{cmes}b)}\quad\begin{minipage}[c]{5.6cm}\raggedright\Large $3{,}925~\text{cm}$\end{minipage}\end{tcolorbox}\par\vspace{2.5mm}
\begin{tcolorbox}[enhanced,nobeforeafter,width=7.1cm,height=2.1cm,valign=center,colback=cmes!6,colframe=cmes,arc=2.5mm,boxrule=0.9pt,left=3mm]{\Large\bfseries\color{cmes}c)}\quad\begin{minipage}[c]{5.6cm}\raggedright\Large $7{,}85~\text{cm}$\end{minipage}\end{tcolorbox}\hfill\begin{tcolorbox}[enhanced,nobeforeafter,width=7.1cm,height=2.1cm,valign=center,colback=cmes!6,colframe=cmes,arc=2.5mm,boxrule=0.9pt,left=3mm]{\Large\bfseries\color{cmes}d)}\quad\begin{minipage}[c]{5.6cm}\raggedright\Large $4{,}90625~\text{cm}$\end{minipage}\end{tcolorbox}
\end{tcolorbox}
\newpage
\begin{tcolorbox}[enhanced,width=16cm,height=14cm,colback=white,colframe=cmes,boxrule=2pt,arc=0mm,
left=6mm,right=6mm,top=0mm,bottom=4mm,sharp corners,
overlay={\fill[cmes] (frame.north west) rectangle ([yshift=-1.25cm]frame.north east);
\node[anchor=west,text=white,font=\sffamily\bfseries\LARGE] at ([xshift=5mm,yshift=-0.62cm]frame.north west) {Question 79 D};
\node[anchor=east,text=white,font=\sffamily\itshape\large] at ([xshift=-5mm,yshift=-0.62cm]frame.north east) {Grandeurs et mesures};}]
\vspace*{1.75cm}
\begin{minipage}[t][6.9cm][c]{\linewidth}\LARGE\centering
\begin{minipage}{0.95\linewidth}\raggedright\linespread{1.25}\selectfont Un pavé droit mesure $12$ cm de longueur, $5$ cm de largeur et $1{,}5$ cm de hauteur. Quel est son volume~?\end{minipage}
\end{minipage}

\vfill
\begin{tcolorbox}[enhanced,nobeforeafter,width=7.1cm,height=2.1cm,valign=center,colback=cmes!6,colframe=cmes,arc=2.5mm,boxrule=0.9pt,left=3mm]{\Large\bfseries\color{cmes}a)}\quad\begin{minipage}[c]{5.6cm}\raggedright\Large $60~\text{cm}^3$\end{minipage}\end{tcolorbox}\hfill\begin{tcolorbox}[enhanced,nobeforeafter,width=7.1cm,height=2.1cm,valign=center,colback=cmes!6,colframe=cmes,arc=2.5mm,boxrule=0.9pt,left=3mm]{\Large\bfseries\color{cmes}b)}\quad\begin{minipage}[c]{5.6cm}\raggedright\Large $18{,}5~\text{cm}^3$\end{minipage}\end{tcolorbox}\par\vspace{2.5mm}
\begin{tcolorbox}[enhanced,nobeforeafter,width=7.1cm,height=2.1cm,valign=center,colback=cmes!6,colframe=cmes,arc=2.5mm,boxrule=0.9pt,left=3mm]{\Large\bfseries\color{cmes}c)}\quad\begin{minipage}[c]{5.6cm}\raggedright\Large $171~\text{cm}^3$\end{minipage}\end{tcolorbox}\hfill\begin{tcolorbox}[enhanced,nobeforeafter,width=7.1cm,height=2.1cm,valign=center,colback=cmes!6,colframe=cmes,arc=2.5mm,boxrule=0.9pt,left=3mm]{\Large\bfseries\color{cmes}d)}\quad\begin{minipage}[c]{5.6cm}\raggedright\Large $90~\text{cm}^3$\end{minipage}\end{tcolorbox}
\end{tcolorbox}
\newpage
\begin{tcolorbox}[enhanced,width=16cm,height=14cm,colback=white,colframe=cmes,boxrule=2pt,arc=0mm,
left=6mm,right=6mm,top=0mm,bottom=4mm,sharp corners,
overlay={\fill[cmes] (frame.north west) rectangle ([yshift=-1.25cm]frame.north east);
\node[anchor=west,text=white,font=\sffamily\bfseries\LARGE] at ([xshift=5mm,yshift=-0.62cm]frame.north west) {Question 80 D};
\node[anchor=east,text=white,font=\sffamily\itshape\large] at ([xshift=-5mm,yshift=-0.62cm]frame.north east) {Grandeurs et mesures};}]
\vspace*{1.75cm}
\begin{minipage}[t][6.9cm][c]{\linewidth}\LARGE\centering
\begin{minipage}{0.95\linewidth}\raggedright\linespread{1.25}\selectfont Un pavé droit mesure $9$ cm de longueur, $2$ cm de largeur et $3$ cm de hauteur. Quel est son volume~?\end{minipage}
\end{minipage}

\vfill
\begin{tcolorbox}[enhanced,nobeforeafter,width=7.1cm,height=2.1cm,valign=center,colback=cmes!6,colframe=cmes,arc=2.5mm,boxrule=0.9pt,left=3mm]{\Large\bfseries\color{cmes}a)}\quad\begin{minipage}[c]{5.6cm}\raggedright\Large $14~\text{cm}^3$\end{minipage}\end{tcolorbox}\hfill\begin{tcolorbox}[enhanced,nobeforeafter,width=7.1cm,height=2.1cm,valign=center,colback=cmes!6,colframe=cmes,arc=2.5mm,boxrule=0.9pt,left=3mm]{\Large\bfseries\color{cmes}b)}\quad\begin{minipage}[c]{5.6cm}\raggedright\Large $102~\text{cm}^3$\end{minipage}\end{tcolorbox}\par\vspace{2.5mm}
\begin{tcolorbox}[enhanced,nobeforeafter,width=7.1cm,height=2.1cm,valign=center,colback=cmes!6,colframe=cmes,arc=2.5mm,boxrule=0.9pt,left=3mm]{\Large\bfseries\color{cmes}c)}\quad\begin{minipage}[c]{5.6cm}\raggedright\Large $54~\text{cm}^3$\end{minipage}\end{tcolorbox}\hfill\begin{tcolorbox}[enhanced,nobeforeafter,width=7.1cm,height=2.1cm,valign=center,colback=cmes!6,colframe=cmes,arc=2.5mm,boxrule=0.9pt,left=3mm]{\Large\bfseries\color{cmes}d)}\quad\begin{minipage}[c]{5.6cm}\raggedright\Large $18~\text{cm}^3$\end{minipage}\end{tcolorbox}
\end{tcolorbox}
\newpage
\begin{tcolorbox}[enhanced,width=16cm,height=14cm,colback=white,colframe=cpro,boxrule=2pt,arc=0mm,
left=6mm,right=6mm,top=0mm,bottom=4mm,sharp corners,
overlay={\fill[cpro] (frame.north west) rectangle ([yshift=-1.25cm]frame.north east);
\node[anchor=west,text=white,font=\sffamily\bfseries\LARGE] at ([xshift=5mm,yshift=-0.62cm]frame.north west) {Question 81 D};
\node[anchor=east,text=white,font=\sffamily\itshape\large] at ([xshift=-5mm,yshift=-0.62cm]frame.north east) {Proportionnalité et données};}]
\vspace*{1.75cm}
\begin{minipage}[t][6.9cm][c]{\linewidth}\LARGE\centering
\begin{minipage}{0.95\linewidth}\raggedright\linespread{1.25}\selectfont Pour $10$ personnes, il faut $750$ g de farine. Combien en faut-il pour $4$ personnes~?\end{minipage}
\end{minipage}

\vfill
\begin{tcolorbox}[enhanced,nobeforeafter,width=7.1cm,height=2.1cm,valign=center,colback=cpro!6,colframe=cpro,arc=2.5mm,boxrule=0.9pt,left=3mm]{\Large\bfseries\color{cpro}a)}\quad\begin{minipage}[c]{5.6cm}\raggedright\Large $75~\text{g}$\end{minipage}\end{tcolorbox}\hfill\begin{tcolorbox}[enhanced,nobeforeafter,width=7.1cm,height=2.1cm,valign=center,colback=cpro!6,colframe=cpro,arc=2.5mm,boxrule=0.9pt,left=3mm]{\Large\bfseries\color{cpro}b)}\quad\begin{minipage}[c]{5.6cm}\raggedright\Large $3\,000~\text{g}$\end{minipage}\end{tcolorbox}\par\vspace{2.5mm}
\begin{tcolorbox}[enhanced,nobeforeafter,width=7.1cm,height=2.1cm,valign=center,colback=cpro!6,colframe=cpro,arc=2.5mm,boxrule=0.9pt,left=3mm]{\Large\bfseries\color{cpro}c)}\quad\begin{minipage}[c]{5.6cm}\raggedright\Large $744~\text{g}$\end{minipage}\end{tcolorbox}\hfill\begin{tcolorbox}[enhanced,nobeforeafter,width=7.1cm,height=2.1cm,valign=center,colback=cpro!6,colframe=cpro,arc=2.5mm,boxrule=0.9pt,left=3mm]{\Large\bfseries\color{cpro}d)}\quad\begin{minipage}[c]{5.6cm}\raggedright\Large $300~\text{g}$\end{minipage}\end{tcolorbox}
\end{tcolorbox}
\newpage
\begin{tcolorbox}[enhanced,width=16cm,height=14cm,colback=white,colframe=cpro,boxrule=2pt,arc=0mm,
left=6mm,right=6mm,top=0mm,bottom=4mm,sharp corners,
overlay={\fill[cpro] (frame.north west) rectangle ([yshift=-1.25cm]frame.north east);
\node[anchor=west,text=white,font=\sffamily\bfseries\LARGE] at ([xshift=5mm,yshift=-0.62cm]frame.north west) {Question 82 D};
\node[anchor=east,text=white,font=\sffamily\itshape\large] at ([xshift=-5mm,yshift=-0.62cm]frame.north east) {Proportionnalité et données};}]
\vspace*{1.75cm}
\begin{minipage}[t][6.9cm][c]{\linewidth}\LARGE\centering
\begin{minipage}{0.95\linewidth}\raggedright\linespread{1.25}\selectfont Pour $4$ personnes, il faut $120$ g de farine. Combien en faut-il pour $10$ personnes~?\end{minipage}
\end{minipage}

\vfill
\begin{tcolorbox}[enhanced,nobeforeafter,width=7.1cm,height=2.1cm,valign=center,colback=cpro!6,colframe=cpro,arc=2.5mm,boxrule=0.9pt,left=3mm]{\Large\bfseries\color{cpro}a)}\quad\begin{minipage}[c]{5.6cm}\raggedright\Large $30~\text{g}$\end{minipage}\end{tcolorbox}\hfill\begin{tcolorbox}[enhanced,nobeforeafter,width=7.1cm,height=2.1cm,valign=center,colback=cpro!6,colframe=cpro,arc=2.5mm,boxrule=0.9pt,left=3mm]{\Large\bfseries\color{cpro}b)}\quad\begin{minipage}[c]{5.6cm}\raggedright\Large $126~\text{g}$\end{minipage}\end{tcolorbox}\par\vspace{2.5mm}
\begin{tcolorbox}[enhanced,nobeforeafter,width=7.1cm,height=2.1cm,valign=center,colback=cpro!6,colframe=cpro,arc=2.5mm,boxrule=0.9pt,left=3mm]{\Large\bfseries\color{cpro}c)}\quad\begin{minipage}[c]{5.6cm}\raggedright\Large $1\,200~\text{g}$\end{minipage}\end{tcolorbox}\hfill\begin{tcolorbox}[enhanced,nobeforeafter,width=7.1cm,height=2.1cm,valign=center,colback=cpro!6,colframe=cpro,arc=2.5mm,boxrule=0.9pt,left=3mm]{\Large\bfseries\color{cpro}d)}\quad\begin{minipage}[c]{5.6cm}\raggedright\Large $300~\text{g}$\end{minipage}\end{tcolorbox}
\end{tcolorbox}
\newpage
\begin{tcolorbox}[enhanced,width=16cm,height=14cm,colback=white,colframe=cpro,boxrule=2pt,arc=0mm,
left=6mm,right=6mm,top=0mm,bottom=4mm,sharp corners,
overlay={\fill[cpro] (frame.north west) rectangle ([yshift=-1.25cm]frame.north east);
\node[anchor=west,text=white,font=\sffamily\bfseries\LARGE] at ([xshift=5mm,yshift=-0.62cm]frame.north west) {Question 83 D};
\node[anchor=east,text=white,font=\sffamily\itshape\large] at ([xshift=-5mm,yshift=-0.62cm]frame.north east) {Proportionnalité et données};}]
\vspace*{1.75cm}
\begin{minipage}[t][6.9cm][c]{\linewidth}\LARGE\centering
\begin{minipage}{0.95\linewidth}\raggedright\linespread{1.25}\selectfont $7$ cahiers coûtent $2\,450$ F. Combien coûte un cahier~?\end{minipage}
\end{minipage}

\vfill
\begin{tcolorbox}[enhanced,nobeforeafter,width=7.1cm,height=2.1cm,valign=center,colback=cpro!6,colframe=cpro,arc=2.5mm,boxrule=0.9pt,left=3mm]{\Large\bfseries\color{cpro}a)}\quad\begin{minipage}[c]{5.6cm}\raggedright\Large $2\,443~\text{F}$\end{minipage}\end{tcolorbox}\hfill\begin{tcolorbox}[enhanced,nobeforeafter,width=7.1cm,height=2.1cm,valign=center,colback=cpro!6,colframe=cpro,arc=2.5mm,boxrule=0.9pt,left=3mm]{\Large\bfseries\color{cpro}b)}\quad\begin{minipage}[c]{5.6cm}\raggedright\Large $700~\text{F}$\end{minipage}\end{tcolorbox}\par\vspace{2.5mm}
\begin{tcolorbox}[enhanced,nobeforeafter,width=7.1cm,height=2.1cm,valign=center,colback=cpro!6,colframe=cpro,arc=2.5mm,boxrule=0.9pt,left=3mm]{\Large\bfseries\color{cpro}c)}\quad\begin{minipage}[c]{5.6cm}\raggedright\Large $350~\text{F}$\end{minipage}\end{tcolorbox}\hfill\begin{tcolorbox}[enhanced,nobeforeafter,width=7.1cm,height=2.1cm,valign=center,colback=cpro!6,colframe=cpro,arc=2.5mm,boxrule=0.9pt,left=3mm]{\Large\bfseries\color{cpro}d)}\quad\begin{minipage}[c]{5.6cm}\raggedright\Large $357~\text{F}$\end{minipage}\end{tcolorbox}
\end{tcolorbox}
\newpage
\begin{tcolorbox}[enhanced,width=16cm,height=14cm,colback=white,colframe=cpro,boxrule=2pt,arc=0mm,
left=6mm,right=6mm,top=0mm,bottom=4mm,sharp corners,
overlay={\fill[cpro] (frame.north west) rectangle ([yshift=-1.25cm]frame.north east);
\node[anchor=west,text=white,font=\sffamily\bfseries\LARGE] at ([xshift=5mm,yshift=-0.62cm]frame.north west) {Question 84 D};
\node[anchor=east,text=white,font=\sffamily\itshape\large] at ([xshift=-5mm,yshift=-0.62cm]frame.north east) {Proportionnalité et données};}]
\vspace*{1.75cm}
\begin{minipage}[t][6.9cm][c]{\linewidth}\LARGE\centering
\begin{minipage}{0.95\linewidth}\raggedright\linespread{1.25}\selectfont $9$ cahiers coûtent $1\,080$ F. Combien coûte un cahier~?\end{minipage}
\end{minipage}

\vfill
\begin{tcolorbox}[enhanced,nobeforeafter,width=7.1cm,height=2.1cm,valign=center,colback=cpro!6,colframe=cpro,arc=2.5mm,boxrule=0.9pt,left=3mm]{\Large\bfseries\color{cpro}a)}\quad\begin{minipage}[c]{5.6cm}\raggedright\Large $120~\text{F}$\end{minipage}\end{tcolorbox}\hfill\begin{tcolorbox}[enhanced,nobeforeafter,width=7.1cm,height=2.1cm,valign=center,colback=cpro!6,colframe=cpro,arc=2.5mm,boxrule=0.9pt,left=3mm]{\Large\bfseries\color{cpro}b)}\quad\begin{minipage}[c]{5.6cm}\raggedright\Large $1\,071~\text{F}$\end{minipage}\end{tcolorbox}\par\vspace{2.5mm}
\begin{tcolorbox}[enhanced,nobeforeafter,width=7.1cm,height=2.1cm,valign=center,colback=cpro!6,colframe=cpro,arc=2.5mm,boxrule=0.9pt,left=3mm]{\Large\bfseries\color{cpro}c)}\quad\begin{minipage}[c]{5.6cm}\raggedright\Large $129~\text{F}$\end{minipage}\end{tcolorbox}\hfill\begin{tcolorbox}[enhanced,nobeforeafter,width=7.1cm,height=2.1cm,valign=center,colback=cpro!6,colframe=cpro,arc=2.5mm,boxrule=0.9pt,left=3mm]{\Large\bfseries\color{cpro}d)}\quad\begin{minipage}[c]{5.6cm}\raggedright\Large $240~\text{F}$\end{minipage}\end{tcolorbox}
\end{tcolorbox}
\newpage
\begin{tcolorbox}[enhanced,width=16cm,height=14cm,colback=white,colframe=cpro,boxrule=2pt,arc=0mm,
left=6mm,right=6mm,top=0mm,bottom=4mm,sharp corners,
overlay={\fill[cpro] (frame.north west) rectangle ([yshift=-1.25cm]frame.north east);
\node[anchor=west,text=white,font=\sffamily\bfseries\LARGE] at ([xshift=5mm,yshift=-0.62cm]frame.north west) {Question 85 D};
\node[anchor=east,text=white,font=\sffamily\itshape\large] at ([xshift=-5mm,yshift=-0.62cm]frame.north east) {Proportionnalité et données};}]
\vspace*{1.75cm}
\begin{minipage}[t][6.9cm][c]{\linewidth}\LARGE\centering
\begin{minipage}{0.95\linewidth}\raggedright\linespread{1.25}\selectfont Que vaut $25\,\%$ de $48$~?\end{minipage}
\end{minipage}

\vfill
\begin{tcolorbox}[enhanced,nobeforeafter,width=7.1cm,height=2.1cm,valign=center,colback=cpro!6,colframe=cpro,arc=2.5mm,boxrule=0.9pt,left=3mm]{\Large\bfseries\color{cpro}a)}\quad\begin{minipage}[c]{5.6cm}\raggedright\Large $12$\end{minipage}\end{tcolorbox}\hfill\begin{tcolorbox}[enhanced,nobeforeafter,width=7.1cm,height=2.1cm,valign=center,colback=cpro!6,colframe=cpro,arc=2.5mm,boxrule=0.9pt,left=3mm]{\Large\bfseries\color{cpro}b)}\quad\begin{minipage}[c]{5.6cm}\raggedright\Large $36$\end{minipage}\end{tcolorbox}\par\vspace{2.5mm}
\begin{tcolorbox}[enhanced,nobeforeafter,width=7.1cm,height=2.1cm,valign=center,colback=cpro!6,colframe=cpro,arc=2.5mm,boxrule=0.9pt,left=3mm]{\Large\bfseries\color{cpro}c)}\quad\begin{minipage}[c]{5.6cm}\raggedright\Large $24$\end{minipage}\end{tcolorbox}\hfill\begin{tcolorbox}[enhanced,nobeforeafter,width=7.1cm,height=2.1cm,valign=center,colback=cpro!6,colframe=cpro,arc=2.5mm,boxrule=0.9pt,left=3mm]{\Large\bfseries\color{cpro}d)}\quad\begin{minipage}[c]{5.6cm}\raggedright\Large $22$\end{minipage}\end{tcolorbox}
\end{tcolorbox}
\newpage
\begin{tcolorbox}[enhanced,width=16cm,height=14cm,colback=white,colframe=cpro,boxrule=2pt,arc=0mm,
left=6mm,right=6mm,top=0mm,bottom=4mm,sharp corners,
overlay={\fill[cpro] (frame.north west) rectangle ([yshift=-1.25cm]frame.north east);
\node[anchor=west,text=white,font=\sffamily\bfseries\LARGE] at ([xshift=5mm,yshift=-0.62cm]frame.north west) {Question 86 D};
\node[anchor=east,text=white,font=\sffamily\itshape\large] at ([xshift=-5mm,yshift=-0.62cm]frame.north east) {Proportionnalité et données};}]
\vspace*{1.75cm}
\begin{minipage}[t][6.9cm][c]{\linewidth}\LARGE\centering
\begin{minipage}{0.95\linewidth}\raggedright\linespread{1.25}\selectfont Que vaut $10\,\%$ de $72$~?\end{minipage}
\end{minipage}

\vfill
\begin{tcolorbox}[enhanced,nobeforeafter,width=7.1cm,height=2.1cm,valign=center,colback=cpro!6,colframe=cpro,arc=2.5mm,boxrule=0.9pt,left=3mm]{\Large\bfseries\color{cpro}a)}\quad\begin{minipage}[c]{5.6cm}\raggedright\Large $0{,}72$\end{minipage}\end{tcolorbox}\hfill\begin{tcolorbox}[enhanced,nobeforeafter,width=7.1cm,height=2.1cm,valign=center,colback=cpro!6,colframe=cpro,arc=2.5mm,boxrule=0.9pt,left=3mm]{\Large\bfseries\color{cpro}b)}\quad\begin{minipage}[c]{5.6cm}\raggedright\Large $17{,}2$\end{minipage}\end{tcolorbox}\par\vspace{2.5mm}
\begin{tcolorbox}[enhanced,nobeforeafter,width=7.1cm,height=2.1cm,valign=center,colback=cpro!6,colframe=cpro,arc=2.5mm,boxrule=0.9pt,left=3mm]{\Large\bfseries\color{cpro}c)}\quad\begin{minipage}[c]{5.6cm}\raggedright\Large $7{,}2$\end{minipage}\end{tcolorbox}\hfill\begin{tcolorbox}[enhanced,nobeforeafter,width=7.1cm,height=2.1cm,valign=center,colback=cpro!6,colframe=cpro,arc=2.5mm,boxrule=0.9pt,left=3mm]{\Large\bfseries\color{cpro}d)}\quad\begin{minipage}[c]{5.6cm}\raggedright\Large $64{,}8$\end{minipage}\end{tcolorbox}
\end{tcolorbox}
\newpage
\begin{tcolorbox}[enhanced,width=16cm,height=14cm,colback=white,colframe=cpro,boxrule=2pt,arc=0mm,
left=6mm,right=6mm,top=0mm,bottom=4mm,sharp corners,
overlay={\fill[cpro] (frame.north west) rectangle ([yshift=-1.25cm]frame.north east);
\node[anchor=west,text=white,font=\sffamily\bfseries\LARGE] at ([xshift=5mm,yshift=-0.62cm]frame.north west) {Question 87 D};
\node[anchor=east,text=white,font=\sffamily\itshape\large] at ([xshift=-5mm,yshift=-0.62cm]frame.north east) {Proportionnalité et données};}]
\vspace*{1.75cm}
\begin{minipage}[t][6.9cm][c]{\linewidth}\LARGE\centering
\begin{minipage}{0.95\linewidth}\raggedright\linespread{1.25}\selectfont Sur un plan à l'échelle $\dfrac{1}{2\,000}$, une longueur mesure $7$ cm. Quelle est la longueur réelle~?\end{minipage}
\end{minipage}

\vfill
\begin{tcolorbox}[enhanced,nobeforeafter,width=7.1cm,height=2.1cm,valign=center,colback=cpro!6,colframe=cpro,arc=2.5mm,boxrule=0.9pt,left=3mm]{\Large\bfseries\color{cpro}a)}\quad\begin{minipage}[c]{5.6cm}\raggedright\Large $14~\text{m}$\end{minipage}\end{tcolorbox}\hfill\begin{tcolorbox}[enhanced,nobeforeafter,width=7.1cm,height=2.1cm,valign=center,colback=cpro!6,colframe=cpro,arc=2.5mm,boxrule=0.9pt,left=3mm]{\Large\bfseries\color{cpro}b)}\quad\begin{minipage}[c]{5.6cm}\raggedright\Large $14\,000~\text{m}$\end{minipage}\end{tcolorbox}\par\vspace{2.5mm}
\begin{tcolorbox}[enhanced,nobeforeafter,width=7.1cm,height=2.1cm,valign=center,colback=cpro!6,colframe=cpro,arc=2.5mm,boxrule=0.9pt,left=3mm]{\Large\bfseries\color{cpro}c)}\quad\begin{minipage}[c]{5.6cm}\raggedright\Large $1\,400~\text{m}$\end{minipage}\end{tcolorbox}\hfill\begin{tcolorbox}[enhanced,nobeforeafter,width=7.1cm,height=2.1cm,valign=center,colback=cpro!6,colframe=cpro,arc=2.5mm,boxrule=0.9pt,left=3mm]{\Large\bfseries\color{cpro}d)}\quad\begin{minipage}[c]{5.6cm}\raggedright\Large $140~\text{m}$\end{minipage}\end{tcolorbox}
\end{tcolorbox}
\newpage
\begin{tcolorbox}[enhanced,width=16cm,height=14cm,colback=white,colframe=cpro,boxrule=2pt,arc=0mm,
left=6mm,right=6mm,top=0mm,bottom=4mm,sharp corners,
overlay={\fill[cpro] (frame.north west) rectangle ([yshift=-1.25cm]frame.north east);
\node[anchor=west,text=white,font=\sffamily\bfseries\LARGE] at ([xshift=5mm,yshift=-0.62cm]frame.north west) {Question 88 D};
\node[anchor=east,text=white,font=\sffamily\itshape\large] at ([xshift=-5mm,yshift=-0.62cm]frame.north east) {Proportionnalité et données};}]
\vspace*{1.75cm}
\begin{minipage}[t][6.9cm][c]{\linewidth}\LARGE\centering
\begin{minipage}{0.95\linewidth}\raggedright\linespread{1.25}\selectfont Sur un plan à l'échelle $\dfrac{1}{100}$, une longueur mesure $12{,}5$ cm. Quelle est la longueur réelle~?\end{minipage}
\end{minipage}

\vfill
\begin{tcolorbox}[enhanced,nobeforeafter,width=7.1cm,height=2.1cm,valign=center,colback=cpro!6,colframe=cpro,arc=2.5mm,boxrule=0.9pt,left=3mm]{\Large\bfseries\color{cpro}a)}\quad\begin{minipage}[c]{5.6cm}\raggedright\Large $125~\text{m}$\end{minipage}\end{tcolorbox}\hfill\begin{tcolorbox}[enhanced,nobeforeafter,width=7.1cm,height=2.1cm,valign=center,colback=cpro!6,colframe=cpro,arc=2.5mm,boxrule=0.9pt,left=3mm]{\Large\bfseries\color{cpro}b)}\quad\begin{minipage}[c]{5.6cm}\raggedright\Large $1{,}25~\text{m}$\end{minipage}\end{tcolorbox}\par\vspace{2.5mm}
\begin{tcolorbox}[enhanced,nobeforeafter,width=7.1cm,height=2.1cm,valign=center,colback=cpro!6,colframe=cpro,arc=2.5mm,boxrule=0.9pt,left=3mm]{\Large\bfseries\color{cpro}c)}\quad\begin{minipage}[c]{5.6cm}\raggedright\Large $12{,}5~\text{m}$\end{minipage}\end{tcolorbox}\hfill\begin{tcolorbox}[enhanced,nobeforeafter,width=7.1cm,height=2.1cm,valign=center,colback=cpro!6,colframe=cpro,arc=2.5mm,boxrule=0.9pt,left=3mm]{\Large\bfseries\color{cpro}d)}\quad\begin{minipage}[c]{5.6cm}\raggedright\Large $1\,250~\text{m}$\end{minipage}\end{tcolorbox}
\end{tcolorbox}
\newpage
\begin{tcolorbox}[enhanced,width=16cm,height=14cm,colback=white,colframe=cpro,boxrule=2pt,arc=0mm,
left=6mm,right=6mm,top=0mm,bottom=4mm,sharp corners,
overlay={\fill[cpro] (frame.north west) rectangle ([yshift=-1.25cm]frame.north east);
\node[anchor=west,text=white,font=\sffamily\bfseries\LARGE] at ([xshift=5mm,yshift=-0.62cm]frame.north west) {Question 89 D};
\node[anchor=east,text=white,font=\sffamily\itshape\large] at ([xshift=-5mm,yshift=-0.62cm]frame.north east) {Proportionnalité et données};}]
\vspace*{1.75cm}
\begin{minipage}[t][6.9cm][c]{\linewidth}\Large\centering
\begin{minipage}{0.95\linewidth}\raggedright\linespread{1.25}\selectfont Ce tableau est un tableau de proportionnalité~: $\begin{array}{|c|c|c|}\hline 1 & 2 & 4\\\hline 9 & 18 & 36\\\hline\end{array}$. Par quel nombre multiplie-t-on la 1re ligne pour obtenir la 2e~?\end{minipage}
\end{minipage}

\vfill
\begin{tcolorbox}[enhanced,nobeforeafter,width=7.1cm,height=2.1cm,valign=center,colback=cpro!6,colframe=cpro,arc=2.5mm,boxrule=0.9pt,left=3mm]{\Large\bfseries\color{cpro}a)}\quad\begin{minipage}[c]{5.6cm}\raggedright\Large $8$\end{minipage}\end{tcolorbox}\hfill\begin{tcolorbox}[enhanced,nobeforeafter,width=7.1cm,height=2.1cm,valign=center,colback=cpro!6,colframe=cpro,arc=2.5mm,boxrule=0.9pt,left=3mm]{\Large\bfseries\color{cpro}b)}\quad\begin{minipage}[c]{5.6cm}\raggedright\Large $10$\end{minipage}\end{tcolorbox}\par\vspace{2.5mm}
\begin{tcolorbox}[enhanced,nobeforeafter,width=7.1cm,height=2.1cm,valign=center,colback=cpro!6,colframe=cpro,arc=2.5mm,boxrule=0.9pt,left=3mm]{\Large\bfseries\color{cpro}c)}\quad\begin{minipage}[c]{5.6cm}\raggedright\Large $9$\end{minipage}\end{tcolorbox}\hfill\begin{tcolorbox}[enhanced,nobeforeafter,width=7.1cm,height=2.1cm,valign=center,colback=cpro!6,colframe=cpro,arc=2.5mm,boxrule=0.9pt,left=3mm]{\Large\bfseries\color{cpro}d)}\quad\begin{minipage}[c]{5.6cm}\raggedright\Large $\dfrac{1}{9}$\end{minipage}\end{tcolorbox}
\end{tcolorbox}
\newpage
\begin{tcolorbox}[enhanced,width=16cm,height=14cm,colback=white,colframe=cpro,boxrule=2pt,arc=0mm,
left=6mm,right=6mm,top=0mm,bottom=4mm,sharp corners,
overlay={\fill[cpro] (frame.north west) rectangle ([yshift=-1.25cm]frame.north east);
\node[anchor=west,text=white,font=\sffamily\bfseries\LARGE] at ([xshift=5mm,yshift=-0.62cm]frame.north west) {Question 90 D};
\node[anchor=east,text=white,font=\sffamily\itshape\large] at ([xshift=-5mm,yshift=-0.62cm]frame.north east) {Proportionnalité et données};}]
\vspace*{1.75cm}
\begin{minipage}[t][6.9cm][c]{\linewidth}\Large\centering
\begin{minipage}{0.95\linewidth}\raggedright\linespread{1.25}\selectfont Ce tableau est un tableau de proportionnalité~: $\begin{array}{|c|c|c|}\hline 2 & 3 & 5\\\hline 24 & 36 & 60\\\hline\end{array}$. Par quel nombre multiplie-t-on la 1re ligne pour obtenir la 2e~?\end{minipage}
\end{minipage}

\vfill
\begin{tcolorbox}[enhanced,nobeforeafter,width=7.1cm,height=2.1cm,valign=center,colback=cpro!6,colframe=cpro,arc=2.5mm,boxrule=0.9pt,left=3mm]{\Large\bfseries\color{cpro}a)}\quad\begin{minipage}[c]{5.6cm}\raggedright\Large $13$\end{minipage}\end{tcolorbox}\hfill\begin{tcolorbox}[enhanced,nobeforeafter,width=7.1cm,height=2.1cm,valign=center,colback=cpro!6,colframe=cpro,arc=2.5mm,boxrule=0.9pt,left=3mm]{\Large\bfseries\color{cpro}b)}\quad\begin{minipage}[c]{5.6cm}\raggedright\Large $12$\end{minipage}\end{tcolorbox}\par\vspace{2.5mm}
\begin{tcolorbox}[enhanced,nobeforeafter,width=7.1cm,height=2.1cm,valign=center,colback=cpro!6,colframe=cpro,arc=2.5mm,boxrule=0.9pt,left=3mm]{\Large\bfseries\color{cpro}c)}\quad\begin{minipage}[c]{5.6cm}\raggedright\Large $24$\end{minipage}\end{tcolorbox}\hfill\begin{tcolorbox}[enhanced,nobeforeafter,width=7.1cm,height=2.1cm,valign=center,colback=cpro!6,colframe=cpro,arc=2.5mm,boxrule=0.9pt,left=3mm]{\Large\bfseries\color{cpro}d)}\quad\begin{minipage}[c]{5.6cm}\raggedright\Large $\dfrac{1}{12}$\end{minipage}\end{tcolorbox}
\end{tcolorbox}
\newpage
\begin{tcolorbox}[enhanced,width=16cm,height=14cm,colback=white,colframe=cpro,boxrule=2pt,arc=0mm,
left=6mm,right=6mm,top=0mm,bottom=4mm,sharp corners,
overlay={\fill[cpro] (frame.north west) rectangle ([yshift=-1.25cm]frame.north east);
\node[anchor=west,text=white,font=\sffamily\bfseries\LARGE] at ([xshift=5mm,yshift=-0.62cm]frame.north west) {Question 91 D};
\node[anchor=east,text=white,font=\sffamily\itshape\large] at ([xshift=-5mm,yshift=-0.62cm]frame.north east) {Proportionnalité et données};}]
\vspace*{1.75cm}
\begin{minipage}[t][6.9cm][c]{\linewidth}\Large\centering
\begin{minipage}{0.95\linewidth}\raggedright\linespread{1.25}\selectfont Dans une classe, les élèves viennent au collège ainsi~: $11$ à pied, $16$ en bus, $6$ à vélo et $2$ en voiture. Combien y a-t-il d'élèves dans la classe~?\end{minipage}
\end{minipage}

\vfill
\begin{tcolorbox}[enhanced,nobeforeafter,width=7.1cm,height=2.1cm,valign=center,colback=cpro!6,colframe=cpro,arc=2.5mm,boxrule=0.9pt,left=3mm]{\Large\bfseries\color{cpro}a)}\quad\begin{minipage}[c]{5.6cm}\raggedright\Large $16$\end{minipage}\end{tcolorbox}\hfill\begin{tcolorbox}[enhanced,nobeforeafter,width=7.1cm,height=2.1cm,valign=center,colback=cpro!6,colframe=cpro,arc=2.5mm,boxrule=0.9pt,left=3mm]{\Large\bfseries\color{cpro}b)}\quad\begin{minipage}[c]{5.6cm}\raggedright\Large $35$\end{minipage}\end{tcolorbox}\par\vspace{2.5mm}
\begin{tcolorbox}[enhanced,nobeforeafter,width=7.1cm,height=2.1cm,valign=center,colback=cpro!6,colframe=cpro,arc=2.5mm,boxrule=0.9pt,left=3mm]{\Large\bfseries\color{cpro}c)}\quad\begin{minipage}[c]{5.6cm}\raggedright\Large $36$\end{minipage}\end{tcolorbox}\hfill\begin{tcolorbox}[enhanced,nobeforeafter,width=7.1cm,height=2.1cm,valign=center,colback=cpro!6,colframe=cpro,arc=2.5mm,boxrule=0.9pt,left=3mm]{\Large\bfseries\color{cpro}d)}\quad\begin{minipage}[c]{5.6cm}\raggedright\Large $33$\end{minipage}\end{tcolorbox}
\end{tcolorbox}
\newpage
\begin{tcolorbox}[enhanced,width=16cm,height=14cm,colback=white,colframe=cpro,boxrule=2pt,arc=0mm,
left=6mm,right=6mm,top=0mm,bottom=4mm,sharp corners,
overlay={\fill[cpro] (frame.north west) rectangle ([yshift=-1.25cm]frame.north east);
\node[anchor=west,text=white,font=\sffamily\bfseries\LARGE] at ([xshift=5mm,yshift=-0.62cm]frame.north west) {Question 92 D};
\node[anchor=east,text=white,font=\sffamily\itshape\large] at ([xshift=-5mm,yshift=-0.62cm]frame.north east) {Proportionnalité et données};}]
\vspace*{1.75cm}
\begin{minipage}[t][6.9cm][c]{\linewidth}\LARGE\centering
\begin{minipage}{0.95\linewidth}\raggedright\linespread{1.25}\selectfont Dans une classe, les élèves viennent au collège ainsi~: $9$ à pied, $18$ en bus et $7$ à vélo. Combien y a-t-il d'élèves dans la classe~?\end{minipage}
\end{minipage}

\vfill
\begin{tcolorbox}[enhanced,nobeforeafter,width=7.1cm,height=2.1cm,valign=center,colback=cpro!6,colframe=cpro,arc=2.5mm,boxrule=0.9pt,left=3mm]{\Large\bfseries\color{cpro}a)}\quad\begin{minipage}[c]{5.6cm}\raggedright\Large $18$\end{minipage}\end{tcolorbox}\hfill\begin{tcolorbox}[enhanced,nobeforeafter,width=7.1cm,height=2.1cm,valign=center,colback=cpro!6,colframe=cpro,arc=2.5mm,boxrule=0.9pt,left=3mm]{\Large\bfseries\color{cpro}b)}\quad\begin{minipage}[c]{5.6cm}\raggedright\Large $34$\end{minipage}\end{tcolorbox}\par\vspace{2.5mm}
\begin{tcolorbox}[enhanced,nobeforeafter,width=7.1cm,height=2.1cm,valign=center,colback=cpro!6,colframe=cpro,arc=2.5mm,boxrule=0.9pt,left=3mm]{\Large\bfseries\color{cpro}c)}\quad\begin{minipage}[c]{5.6cm}\raggedright\Large $35$\end{minipage}\end{tcolorbox}\hfill\begin{tcolorbox}[enhanced,nobeforeafter,width=7.1cm,height=2.1cm,valign=center,colback=cpro!6,colframe=cpro,arc=2.5mm,boxrule=0.9pt,left=3mm]{\Large\bfseries\color{cpro}d)}\quad\begin{minipage}[c]{5.6cm}\raggedright\Large $27$\end{minipage}\end{tcolorbox}
\end{tcolorbox}
\newpage
\begin{tcolorbox}[enhanced,width=16cm,height=14cm,colback=white,colframe=cpro,boxrule=2pt,arc=0mm,
left=6mm,right=6mm,top=0mm,bottom=4mm,sharp corners,
overlay={\fill[cpro] (frame.north west) rectangle ([yshift=-1.25cm]frame.north east);
\node[anchor=west,text=white,font=\sffamily\bfseries\LARGE] at ([xshift=5mm,yshift=-0.62cm]frame.north west) {Question 93 D};
\node[anchor=east,text=white,font=\sffamily\itshape\large] at ([xshift=-5mm,yshift=-0.62cm]frame.north east) {Proportionnalité et données};}]
\vspace*{1.75cm}
\begin{minipage}[t][6.9cm][c]{\linewidth}\LARGE\centering
\begin{minipage}{0.95\linewidth}\raggedright\linespread{1.25}\selectfont Ce tableau est-il un tableau de proportionnalité~? $\begin{array}{|c|c|}\hline 10 & 15\\\hline 4 & 6\\\hline\end{array}$\end{minipage}
\end{minipage}

\vfill
\begin{tcolorbox}[enhanced,nobeforeafter,width=7.1cm,height=2.1cm,valign=center,colback=cpro!6,colframe=cpro,arc=2.5mm,boxrule=0.9pt,left=3mm]{\Large\bfseries\color{cpro}a)}\quad\begin{minipage}[c]{5.6cm}\raggedright\Large oui\end{minipage}\end{tcolorbox}\hfill\begin{tcolorbox}[enhanced,nobeforeafter,width=7.1cm,height=2.1cm,valign=center,colback=cpro!6,colframe=cpro,arc=2.5mm,boxrule=0.9pt,left=3mm]{\Large\bfseries\color{cpro}b)}\quad\begin{minipage}[c]{5.6cm}\raggedright\Large il manque une colonne\end{minipage}\end{tcolorbox}\par\vspace{2.5mm}
\begin{tcolorbox}[enhanced,nobeforeafter,width=7.1cm,height=2.1cm,valign=center,colback=cpro!6,colframe=cpro,arc=2.5mm,boxrule=0.9pt,left=3mm]{\Large\bfseries\color{cpro}c)}\quad\begin{minipage}[c]{5.6cm}\raggedright\Large on ne peut pas savoir\end{minipage}\end{tcolorbox}\hfill\begin{tcolorbox}[enhanced,nobeforeafter,width=7.1cm,height=2.1cm,valign=center,colback=cpro!6,colframe=cpro,arc=2.5mm,boxrule=0.9pt,left=3mm]{\Large\bfseries\color{cpro}d)}\quad\begin{minipage}[c]{5.6cm}\raggedright\Large non\end{minipage}\end{tcolorbox}
\end{tcolorbox}
\newpage
\begin{tcolorbox}[enhanced,width=16cm,height=14cm,colback=white,colframe=cpro,boxrule=2pt,arc=0mm,
left=6mm,right=6mm,top=0mm,bottom=4mm,sharp corners,
overlay={\fill[cpro] (frame.north west) rectangle ([yshift=-1.25cm]frame.north east);
\node[anchor=west,text=white,font=\sffamily\bfseries\LARGE] at ([xshift=5mm,yshift=-0.62cm]frame.north west) {Question 94 D};
\node[anchor=east,text=white,font=\sffamily\itshape\large] at ([xshift=-5mm,yshift=-0.62cm]frame.north east) {Proportionnalité et données};}]
\vspace*{1.75cm}
\begin{minipage}[t][6.9cm][c]{\linewidth}\LARGE\centering
\begin{minipage}{0.95\linewidth}\raggedright\linespread{1.25}\selectfont Ce tableau est-il un tableau de proportionnalité~? $\begin{array}{|c|c|}\hline 3 & 9\\\hline 7 & 20\\\hline\end{array}$\end{minipage}
\end{minipage}

\vfill
\begin{tcolorbox}[enhanced,nobeforeafter,width=7.1cm,height=2.1cm,valign=center,colback=cpro!6,colframe=cpro,arc=2.5mm,boxrule=0.9pt,left=3mm]{\Large\bfseries\color{cpro}a)}\quad\begin{minipage}[c]{5.6cm}\raggedright\Large non\end{minipage}\end{tcolorbox}\hfill\begin{tcolorbox}[enhanced,nobeforeafter,width=7.1cm,height=2.1cm,valign=center,colback=cpro!6,colframe=cpro,arc=2.5mm,boxrule=0.9pt,left=3mm]{\Large\bfseries\color{cpro}b)}\quad\begin{minipage}[c]{5.6cm}\raggedright\Large on ne peut pas savoir\end{minipage}\end{tcolorbox}\par\vspace{2.5mm}
\begin{tcolorbox}[enhanced,nobeforeafter,width=7.1cm,height=2.1cm,valign=center,colback=cpro!6,colframe=cpro,arc=2.5mm,boxrule=0.9pt,left=3mm]{\Large\bfseries\color{cpro}c)}\quad\begin{minipage}[c]{5.6cm}\raggedright\Large oui\end{minipage}\end{tcolorbox}\hfill\begin{tcolorbox}[enhanced,nobeforeafter,width=7.1cm,height=2.1cm,valign=center,colback=cpro!6,colframe=cpro,arc=2.5mm,boxrule=0.9pt,left=3mm]{\Large\bfseries\color{cpro}d)}\quad\begin{minipage}[c]{5.6cm}\raggedright\Large il manque une colonne\end{minipage}\end{tcolorbox}
\end{tcolorbox}
\newpage
\begin{tcolorbox}[enhanced,width=16cm,height=14cm,colback=white,colframe=cpro,boxrule=2pt,arc=0mm,
left=6mm,right=6mm,top=0mm,bottom=4mm,sharp corners,
overlay={\fill[cpro] (frame.north west) rectangle ([yshift=-1.25cm]frame.north east);
\node[anchor=west,text=white,font=\sffamily\bfseries\LARGE] at ([xshift=5mm,yshift=-0.62cm]frame.north west) {Question 95 D};
\node[anchor=east,text=white,font=\sffamily\itshape\large] at ([xshift=-5mm,yshift=-0.62cm]frame.north east) {Proportionnalité et données};}]
\vspace*{1.75cm}
\begin{minipage}[t][6.9cm][c]{\linewidth}\LARGE\centering
\begin{minipage}{0.95\linewidth}\raggedright\linespread{1.25}\selectfont Un robinet remplit $30$ L en $10$ min. Quel volume remplit-il en $4$ min, au même débit~?\end{minipage}
\end{minipage}

\vfill
\begin{tcolorbox}[enhanced,nobeforeafter,width=7.1cm,height=2.1cm,valign=center,colback=cpro!6,colframe=cpro,arc=2.5mm,boxrule=0.9pt,left=3mm]{\Large\bfseries\color{cpro}a)}\quad\begin{minipage}[c]{5.6cm}\raggedright\Large $12~\text{L}$\end{minipage}\end{tcolorbox}\hfill\begin{tcolorbox}[enhanced,nobeforeafter,width=7.1cm,height=2.1cm,valign=center,colback=cpro!6,colframe=cpro,arc=2.5mm,boxrule=0.9pt,left=3mm]{\Large\bfseries\color{cpro}b)}\quad\begin{minipage}[c]{5.6cm}\raggedright\Large $120~\text{L}$\end{minipage}\end{tcolorbox}\par\vspace{2.5mm}
\begin{tcolorbox}[enhanced,nobeforeafter,width=7.1cm,height=2.1cm,valign=center,colback=cpro!6,colframe=cpro,arc=2.5mm,boxrule=0.9pt,left=3mm]{\Large\bfseries\color{cpro}c)}\quad\begin{minipage}[c]{5.6cm}\raggedright\Large $3~\text{L}$\end{minipage}\end{tcolorbox}\hfill\begin{tcolorbox}[enhanced,nobeforeafter,width=7.1cm,height=2.1cm,valign=center,colback=cpro!6,colframe=cpro,arc=2.5mm,boxrule=0.9pt,left=3mm]{\Large\bfseries\color{cpro}d)}\quad\begin{minipage}[c]{5.6cm}\raggedright\Large $24~\text{L}$\end{minipage}\end{tcolorbox}
\end{tcolorbox}
\newpage
\begin{tcolorbox}[enhanced,width=16cm,height=14cm,colback=white,colframe=cpro,boxrule=2pt,arc=0mm,
left=6mm,right=6mm,top=0mm,bottom=4mm,sharp corners,
overlay={\fill[cpro] (frame.north west) rectangle ([yshift=-1.25cm]frame.north east);
\node[anchor=west,text=white,font=\sffamily\bfseries\LARGE] at ([xshift=5mm,yshift=-0.62cm]frame.north west) {Question 96 D};
\node[anchor=east,text=white,font=\sffamily\itshape\large] at ([xshift=-5mm,yshift=-0.62cm]frame.north east) {Proportionnalité et données};}]
\vspace*{1.75cm}
\begin{minipage}[t][6.9cm][c]{\linewidth}\LARGE\centering
\begin{minipage}{0.95\linewidth}\raggedright\linespread{1.25}\selectfont Un robinet remplit $7$ L en $2$ min. Quel volume remplit-il en $10$ min, au même débit~?\end{minipage}
\end{minipage}

\vfill
\begin{tcolorbox}[enhanced,nobeforeafter,width=7.1cm,height=2.1cm,valign=center,colback=cpro!6,colframe=cpro,arc=2.5mm,boxrule=0.9pt,left=3mm]{\Large\bfseries\color{cpro}a)}\quad\begin{minipage}[c]{5.6cm}\raggedright\Large $15~\text{L}$\end{minipage}\end{tcolorbox}\hfill\begin{tcolorbox}[enhanced,nobeforeafter,width=7.1cm,height=2.1cm,valign=center,colback=cpro!6,colframe=cpro,arc=2.5mm,boxrule=0.9pt,left=3mm]{\Large\bfseries\color{cpro}b)}\quad\begin{minipage}[c]{5.6cm}\raggedright\Large $35~\text{L}$\end{minipage}\end{tcolorbox}\par\vspace{2.5mm}
\begin{tcolorbox}[enhanced,nobeforeafter,width=7.1cm,height=2.1cm,valign=center,colback=cpro!6,colframe=cpro,arc=2.5mm,boxrule=0.9pt,left=3mm]{\Large\bfseries\color{cpro}c)}\quad\begin{minipage}[c]{5.6cm}\raggedright\Large $70~\text{L}$\end{minipage}\end{tcolorbox}\hfill\begin{tcolorbox}[enhanced,nobeforeafter,width=7.1cm,height=2.1cm,valign=center,colback=cpro!6,colframe=cpro,arc=2.5mm,boxrule=0.9pt,left=3mm]{\Large\bfseries\color{cpro}d)}\quad\begin{minipage}[c]{5.6cm}\raggedright\Large $3{,}5~\text{L}$\end{minipage}\end{tcolorbox}
\end{tcolorbox}
\newpage
\end{document}